\documentclass{article}
\usepackage{amsmath,amssymb,amsthm}
\IfFileExists{hyperref.sty}{%
  \IfFileExists{color.sty}{\usepackage{color}\definecolor{paperlinkblue}{rgb}{0.10,0.25,0.42}%
    \usepackage[colorlinks=true,linkcolor=paperlinkblue]{hyperref}}{\usepackage[hidelinks]{hyperref}}%
}{}
\ifdefined\hyperlink
  \newcommand{\paperanchor}[1]{\hypertarget{#1}{}}%
  \newcommand{\paperlink}[2]{\hyperlink{#1}{\begin{NoHyper}#2\end{NoHyper}}}%
\else
  \ifdefined\pdfextension
    \ifnum\outputmode>0
      \def\paperdest{\pdfextension dest}\def\paperstart{\pdfextension startlink}%
      \def\paperend{\pdfextension endlink}\def\papercolor{\pdfextension colorstack}%
    \fi
  \else\ifdefined\pdfstartlink
    \ifnum\pdfoutput>0
      \let\paperdest\pdfdest\let\paperstart\pdfstartlink\let\paperend\pdfendlink\let\papercolor\pdfcolorstack
    \fi
  \fi\fi
  \ifdefined\paperdest
    \newcommand{\paperanchor}[1]{\paperdest name{#1} xyz}%
    \newcommand{\paperlink}[2]{\leavevmode\paperstart attr{/Border [0 0 0]} goto name{#1}%
      \papercolor 0 push {0.10 0.25 0.42 rg 0.10 0.25 0.42 RG}#2\papercolor 0 pop\paperend}%
  \else\ifdefined\XeTeXversion
    \newcommand{\paperanchor}[1]{\special{pdf:dest (#1) [@thispage /XYZ null @ypos null]}}%
    \newcommand{\paperlink}[2]{\leavevmode\special{pdf:bann << /Subtype /Link /Border [0 0 0] /A << /S /GoTo /D (#1) >> >>}%
      \special{color push rgb 0.10 0.25 0.42}#2\special{color pop}\special{pdf:eann}}%
  \else
    \PackageError{lean-to-latex}{No PDF hyperlink backend available}{Use XeLaTeX, pdfLaTeX, LuaLaTeX, or install hyperref.}%
  \fi\fi
\fi
\newcommand{\paperlineanchor}[1]{\raisebox{10pt}[0pt][0pt]{\paperanchor{#1}}}
\DeclareRobustCommand{\paperresultname}[1]{\paperlineanchor{theorem:\arabic{theorem}}#1}
\newtheoremstyle{paperresult}{8pt}{5pt}{\normalfont\raggedright}{}{\bfseries}{.}{\newline}{}
\theoremstyle{paperresult}
\newtheorem{theorem}{\paperresultname{Theorem}}
\newcounter{paperpremise}
\newtheorem{lemma}[theorem]{\paperresultname{Lemma}}
\newtheorem{corollary}[theorem]{\paperresultname{Corollary}}
\newcounter{definingequation}
\newenvironment{proofstep}{\par\smallskip\ifdim\leftskip<3em\relax\advance\leftskip by 1em\relax\fi\noindent\ignorespaces}{\par\smallskip}
\setlength{\textwidth}{6.5in}
\setlength{\oddsidemargin}{0in}
\AddToHook{env/proof/begin}{\raggedright}
\setlength{\emergencystretch}{\linewidth}
\title{Matchings in graphs}
\author{}
\date{}
\begin{document}
\maketitle

The inclusion \(\operatorname{incl}_{G,H}\) of graphs on the same vertex set fixes each vertex and carries each edge of \(G\) to the same edge in \(H\).

\section*{Definitions}
\paragraph{Is matching.} \par\noindent\refstepcounter{definingequation}\label{defeq:1}\paperlineanchor{defeq:1}\textup{(\paperlink{defeq:1}{D1})} Let \(V\) be a type. Let \(G\) be a simple graph on \(V\). Let \(M\) be a subgraph of \(G\). Then \((M\text{ is a matching}) = (\forall v\, :\, V,\; v \in V(M) \implies (\exists! w\, :\, V,\; v \sim_{M} w))\).

\paragraph{To edge.} \par\noindent\refstepcounter{definingequation}\label{defeq:2}\paperlineanchor{defeq:2}\textup{(\paperlink{defeq:2}{D2})} Let \(V\) be a type. Let \(G\) be a simple graph on \(V\). Let \(M\) be a subgraph of \(G\). Let \(h\) be a proof of \(M\text{ is a matching}\). Let \(v\) be an element of \(V(M)\). Then \(\operatorname{to\,Edge}(v) = (\operatorname{Sym2\mathord{:}mk}((v,\allowbreak{} \text{the fixed choice from }\{x\in V \mid v \sim_{M} x \land (\forall y\, :\, V,\; v \sim_{M} y \implies y = x)\}))\, :\, \operatorname{edge\,Set}(M))\).

\paragraph{Is perfect matching.} \par\noindent\refstepcounter{definingequation}\label{defeq:3}\paperlineanchor{defeq:3}\textup{(\paperlink{defeq:3}{D3})} Let \(V\) be a type. Let \(G\) be a simple graph on \(V\). Let \(M\) be a subgraph of \(G\). Then \((M\text{ is a perfect matching}) = (M\text{ is a matching} \land \operatorname{Is\,Spanning}(M))\).

\paragraph{Is matching free.} \par\noindent\refstepcounter{definingequation}\label{defeq:4}\paperlineanchor{defeq:4}\textup{(\paperlink{defeq:4}{D4})} Let \(V\) be a type. Let \(G\) be a simple graph on \(V\). Then \(\operatorname{Is\,Matching\,Free}(G) = (\forall M\, :\, \operatorname{Subgraph}(G),\; \neg (M\text{ is a perfect matching}))\).

\paragraph{Is cycles.} \par\noindent\refstepcounter{definingequation}\label{defeq:5}\paperlineanchor{defeq:5}\textup{(\paperlink{defeq:5}{D5})} Let \(V\) be a type. Let \(G\) be a simple graph on \(V\). Then \((\text{every nonisolated vertex of }G\text{ has degree two}) = (\forall v\, :\, V,\; N_{G}(v)\ne\varnothing \implies \operatorname{ncard}(N_{G}(v)) = 2)\).

\paragraph{Is alternating.} \par\noindent\refstepcounter{definingequation}\label{defeq:6}\paperlineanchor{defeq:6}\textup{(\paperlink{defeq:6}{D6})} Let \(V\) be a type. Let \(G\) and \(G^{'}\) be simple graphs on \(V\). Then \((G\text{ alternates with respect to }G^{'}) = (\forall v\, :\, V,\; \forall w\, :\, V,\; \forall w^{'}\, :\, V,\; w \ne w^{'} \implies v \sim_{G} w \implies v \sim_{G} w^{'} \implies (v \sim_{G^{'}} w \iff \neg (v \sim_{G^{'}} w^{'})))\).

\subsection*{Prerequisites}
\begin{description}
\item[\paperlineanchor{premise:4165736f702e4275696c74696e52756c65732e6e6f745f696e74726f}\paperlink{premise:4165736f702e4275696c74696e52756c65732e6e6f745f696e74726f}{P1.}] \refstepcounter{paperpremise}\label{premise:4165736f702e4275696c74696e52756c65732e6e6f745f696e74726f}\textbf{Not intro.} Let \(P\) be a proposition. Assume \(P \implies \mathrm{False}\). Then \(\neg P\).
\item[\paperlineanchor{premise:416e642e6c656674}\paperlink{premise:416e642e6c656674}{P2.}] \refstepcounter{paperpremise}\label{premise:416e642e6c656674}\textbf{Left.} Let \(a\) and \(b\) be propositions. Assume \(a \land b\). Then \(a\).
\item[\paperlineanchor{premise:416e642e7269676874}\paperlink{premise:416e642e7269676874}{P3.}] \refstepcounter{paperpremise}\label{premise:416e642e7269676874}\textbf{Right.} Let \(a\) and \(b\) be propositions. Assume \(a \land b\). Then \(b\).
\item[\paperlineanchor{premise:43617264696e616c2e6571}\paperlink{premise:43617264696e616c2e6571}{P4.}] \refstepcounter{paperpremise}\label{premise:43617264696e616c2e6571}\textbf{Equals.} Let \(\alpha\) and \(\beta\) be types. Then \(\operatorname{Cardinal\mathord{:}mk}(\alpha) = \operatorname{Cardinal\mathord{:}mk}(\beta) \iff \operatorname{Bij}(\alpha,\allowbreak{} \beta)\text{ is nonempty}\).
\item[\paperlineanchor{premise:436c6173736963616c2e6e6f745f6e6f74}\paperlink{premise:436c6173736963616c2e6e6f745f6e6f74}{P5.}] \refstepcounter{paperpremise}\label{premise:436c6173736963616c2e6e6f745f6e6f74}\textbf{Not not.} Let \(a\) be a proposition. Then \(\neg \left(\neg a\right) \iff a\).
\item[\paperlineanchor{premise:446563696461626c652e6279436f6e74726164696374696f6e}\paperlink{premise:446563696461626c652e6279436f6e74726164696374696f6e}{P6.}] \refstepcounter{paperpremise}\label{premise:446563696461626c652e6279436f6e74726164696374696f6e}\textbf{By contradiction.} Let \(p\) be a proposition. Assume \(p\) is decidable. Assume \(\neg p \implies \mathrm{False}\). Then \(p\).
\item[\paperlineanchor{premise:4469736a6f696e742e6e655f6f665f6d656d}\paperlink{premise:4469736a6f696e742e6e655f6f665f6d656d}{P7.}] \refstepcounter{paperpremise}\label{premise:4469736a6f696e742e6e655f6f665f6d656d}\textbf{Differs from of membership.} Let \(\alpha\) be a type. Let \(s\) and \(t\) be subsets of \(\alpha\). Assume \(\operatorname{Disjoint}\left(s, t\right)\). Let \(a_{1}\) be an element of \(\alpha\). Assume \(a_{1} \in s\). Let \(b\) be an element of \(\alpha\). Assume \(b \in t\). Then \(a_{1} \ne b\).
\item[\paperlineanchor{premise:4469736a6f696e742e73796d6d}\paperlink{premise:4469736a6f696e742e73796d6d}{P8.}] \refstepcounter{paperpremise}\label{premise:4469736a6f696e742e73796d6d}\textbf{Symm.} Let \(\alpha\) be a type. Equip \(\alpha\) with a partial order. Fix a structure \(\mathit{inst}_{1}\) of type \(\operatorname{Order\,Bot}\left(\alpha\right)\). Let \(a\) and \(b\) be elements of \(\alpha\). Assume \(\operatorname{Disjoint}\left(a, b\right)\). Then \(\operatorname{Disjoint}\left(b, a\right)\).
\item[\paperlineanchor{premise:45717569762e6170706c795f73796d6d5f6170706c79}\paperlink{premise:45717569762e6170706c795f73796d6d5f6170706c79}{P9.}] \refstepcounter{paperpremise}\label{premise:45717569762e6170706c795f73796d6d5f6170706c79}\textbf{Apply symm apply.} Let \(\alpha\) be a sort. Let \(\beta\) be a sort. Let \(e\) be a bijection from \(\alpha\) to \(\beta\). Let \(x\) be an element of \(\beta\). Then \(e\left(\left(e^{-1}\right)\left(x\right)\right) = x\).
\item[\paperlineanchor{premise:45717569762e73796d6d5f6170706c795f6170706c79}\paperlink{premise:45717569762e73796d6d5f6170706c795f6170706c79}{P10.}] \refstepcounter{paperpremise}\label{premise:45717569762e73796d6d5f6170706c795f6170706c79}\textbf{Symm apply apply.} Let \(\alpha\) be a sort. Let \(\beta\) be a sort. Let \(e\) be a bijection from \(\alpha\) to \(\beta\). Let \(x\) be an element of \(\alpha\). Then \(\left(e^{-1}\right)\left(e\left(x\right)\right) = x\).
\item[\paperlineanchor{premise:4578697374732e63686f6f73655f73706563}\paperlink{premise:4578697374732e63686f6f73655f73706563}{P11.}] \refstepcounter{paperpremise}\label{premise:4578697374732e63686f6f73655f73706563}\textbf{Choose spec.} Let \(\alpha\) be a sort. Let \(p\) be a function of type \(\alpha \to \mathsf{Prop}\). Let \(P\) be a proof of \(\exists a\, :\, \alpha,\; p\left(a\right)\). Then \(p\left(\text{the fixed choice from }\{x\in \alpha \mid p(x)\}\right)\).
\item[\paperlineanchor{premise:457869737473556e697175652e756e69717565}\paperlink{premise:457869737473556e697175652e756e69717565}{P12.}] \refstepcounter{paperpremise}\label{premise:457869737473556e697175652e756e69717565}\textbf{Unique.} Let \(\alpha\) be a sort. Let \(p\) be a function of type \(\alpha \to \mathsf{Prop}\). Assume \(\exists! x\, :\, \alpha,\; p\left(x\right)\). Let \(y_{1}\) and \(y_{2}\) be elements of \(\alpha\). Assume \(p\left(y_{1}\right)\). Assume \(p\left(y_{2}\right)\). Then \(y_{1} = y_{2}\).
\item[\paperlineanchor{premise:46696e6974652e6578697374735f6c655f6d6178696d616c}\paperlink{premise:46696e6974652e6578697374735f6c655f6d6178696d616c}{P13.}] \refstepcounter{paperpremise}\label{premise:46696e6974652e6578697374735f6c655f6d6178696d616c}\textbf{Exists at most maximal.} Let \(\alpha\) be a type. Equip \(\alpha\) with a partial order. Let \(a\) be an element of \(\alpha\). Let \(p\) be a function of type \(\alpha \to \mathsf{Prop}\). Assume \(\alpha\) is finite. Assume \(p\left(a\right)\). Then \(\exists b\, :\, \alpha,\; a \le b \land \operatorname{Maximal}\left(p, b\right)\).
\item[\paperlineanchor{premise:46696e6974652e6f665f66696e74797065}\paperlink{premise:46696e6974652e6f665f66696e74797065}{P14.}] \refstepcounter{paperpremise}\label{premise:46696e6974652e6f665f66696e74797065}\textbf{Of fintype.} Let \(\alpha\) be a type. Fix a finite enumeration of \(\alpha\). Then \(\operatorname{Finite}\left(\alpha\right)\).
\item[\paperlineanchor{premise:46696e7365742e636172645f696d6167655f6f665f696e6a656374697665}\paperlink{premise:46696e7365742e636172645f696d6167655f6f665f696e6a656374697665}{P15.}] \refstepcounter{paperpremise}\label{premise:46696e7365742e636172645f696d6167655f6f665f696e6a656374697665}\textbf{Cardinality image of injective.} Let \(\alpha\) and \(\beta\) be types. Let \(f\) be a function of type \(\alpha \to \beta\). Assume equality on \(\beta\) is decidable. Let \(s\) be a finite subset of \(\alpha\). Assume \(f\text{ is injective}\). Then \(\left|f\left[s\right]\right| = \left|s\right|\).
\item[\paperlineanchor{premise:46696e7365742e6d656d5f66696c746572}\paperlink{premise:46696e7365742e6d656d5f66696c746572}{P16.}] \refstepcounter{paperpremise}\label{premise:46696e7365742e6d656d5f66696c746572}\textbf{Membership filter.} Let \(\alpha\) be a type. Let \(p\) be a function of type \(\alpha \to \mathsf{Prop}\). Assume the predicate \(p\) is decidable. Let \(s\) be a finite subset of \(\alpha\). Let \(a\) be an element of \(\alpha\). Then \(a \in \{x\in s \mid p(x)\} \iff a \in s \land p\left(a\right)\).
\item[\paperlineanchor{premise:46696e7365742e6d656d5f756e6976}\paperlink{premise:46696e7365742e6d656d5f756e6976}{P17.}] \refstepcounter{paperpremise}\label{premise:46696e7365742e6d656d5f756e6976}\textbf{Membership in the full set.} Let \(\alpha\) be a type. Fix a finite enumeration of \(\alpha\). Let \(x\) be an element of \(\alpha\). Then \(x \in \alpha\).
\item[\paperlineanchor{premise:46696e7365742e73756d5f636f6e6772}\paperlink{premise:46696e7365742e73756d5f636f6e6772}{P18.}] \refstepcounter{paperpremise}\label{premise:46696e7365742e73756d5f636f6e6772}\textbf{Equality of sums with equal summands.} Let \(\iota\) and \(M\) be types. Let \(s_{1}\) and \(s_{2}\) be finite subsets of \(\iota\). Equip \(M\) with a commutative additive monoid structure. Let \(f\) and \(g\) be functions of type \(\iota \to M\). Assume \(s_{1} = s_{2}\). Assume \(\forall x\, :\, \iota,\; x \in s_{2} \implies f\left(x\right) = g\left(x\right)\). Then \(\sum_{i\in s_{1}}\left(f(i)\right) = \sum_{i\in s_{2}}\left(g(i)\right)\).
\item[\paperlineanchor{premise:46696e7365742e73756d5f636f6e7374}\paperlink{premise:46696e7365742e73756d5f636f6e7374}{P19.}] \refstepcounter{paperpremise}\label{premise:46696e7365742e73756d5f636f6e7374}\textbf{The sum of a constant.} Let \(\iota\) and \(M\) be types. Let \(s\) be a finite subset of \(\iota\). Equip \(M\) with a commutative additive monoid structure. Let \(b\) be an element of \(M\). Then \(\sum_{x\in s}\left(b\right) = \sum_{k=1}^{\left|s\right|}\left(b\right)\).
\item[\paperlineanchor{premise:46696e747970652e636172645f6f6646696e736574}\paperlink{premise:46696e747970652e636172645f6f6646696e736574}{P20.}] \refstepcounter{paperpremise}\label{premise:46696e747970652e636172645f6f6646696e736574}\textbf{Cardinality of finset.} Let \(\alpha\) be a type. Let \(p\) be a subset of \(\alpha\). Let \(s\) be a finite subset of \(\alpha\). Let \(H\) be a proof of \(\forall x\, :\, \alpha,\; x \in s \iff x \in p\). Then \(\left|p\right| = \left|s\right|\).
\item[\paperlineanchor{premise:46696e747970652e696e73744661737453756273696e676c65746f6e}\paperlink{premise:46696e747970652e696e73744661737453756273696e676c65746f6e}{P21.}] \refstepcounter{paperpremise}\label{premise:46696e747970652e696e73744661737453756273696e676c65746f6e}\textbf{Inst fast subsingleton.} Let \(\alpha\) be a type. Then \(\operatorname{Fast\,Subsingleton}\left(\operatorname{Fintype}\left(\alpha\right)\right)\).
\item[\paperlineanchor{premise:4861735375627365742e5375627365742e616e746973796d6d}\paperlink{premise:4861735375627365742e5375627365742e616e746973796d6d}{P22.}] \refstepcounter{paperpremise}\label{premise:4861735375627365742e5375627365742e616e746973796d6d}\textbf{Antisymm.} Let \(\alpha\) be a type. Equip \(\alpha\) with a subset relation. Let \(a\) and \(b\) be elements of \(\alpha\). Let \(\mathit{inst}_{1}\) be a proof of \(\operatorname{Is\,Antisymm}(\alpha,\allowbreak{} \mathit{x1} \mapsto \mathit{x2} \mapsto \mathit{x1} \subseteq \mathit{x2})\). Assume \(a \subseteq b\). Assume \(b \subseteq a\). Then \(a = b\).
\item[\paperlineanchor{premise:4966662e6d70}\paperlink{premise:4966662e6d70}{P23.}] \refstepcounter{paperpremise}\label{premise:4966662e6d70}\textbf{Mp.} Let \(a\) and \(b\) be propositions. Assume \(a \iff b\). Assume \(a\). Then \(b\).
\item[\paperlineanchor{premise:4966662e6d7072}\paperlink{premise:4966662e6d7072}{P24.}] \refstepcounter{paperpremise}\label{premise:4966662e6d7072}\textbf{Mpr.} Let \(a\) and \(b\) be propositions. Assume \(a \iff b\). Assume \(b\). Then \(a\).
\item[\paperlineanchor{premise:4966662e6f665f6571}\paperlink{premise:4966662e6f665f6571}{P25.}] \refstepcounter{paperpremise}\label{premise:4966662e6f665f6571}\textbf{Of equals.} Let \(a\) and \(b\) be propositions. Assume \(a = b\). Then \(a \iff b\).
\item[\paperlineanchor{premise:496e742e6164645f6f6e655f6c655f6f665f6c74}\paperlink{premise:496e742e6164645f6f6e655f6c655f6f665f6c74}{P26.}] \refstepcounter{paperpremise}\label{premise:496e742e6164645f6f6e655f6c655f6f665f6c74}\textbf{Sum one at most of less than.} Let \(a\) and \(b\) be integers. Assume \(a < b\). Then \(a + 1 \le b\).
\item[\paperlineanchor{premise:496e742e7375625f65715f7a65726f5f6f665f6571}\paperlink{premise:496e742e7375625f65715f7a65726f5f6f665f6571}{P27.}] \refstepcounter{paperpremise}\label{premise:496e742e7375625f65715f7a65726f5f6f665f6571}\textbf{Difference equals zero of equals.} Let \(a\) and \(b\) be integers. Assume \(a = b\). Then \(a - b = 0\).
\item[\paperlineanchor{premise:496e742e7375625f6e6f6e6e65675f6f665f6c65}\paperlink{premise:496e742e7375625f6e6f6e6e65675f6f665f6c65}{P28.}] \refstepcounter{paperpremise}\label{premise:496e742e7375625f6e6f6e6e65675f6f665f6c65}\textbf{Difference nonnegative of at most.} Let \(a\) and \(b\) be integers. Assume \(b \le a\). Then \(0 \le a - b\).
\item[\paperlineanchor{premise:4c65616e2e4f6d6567612e436f6e73747261696e742e616464457175616c6974795f736174}\paperlink{premise:4c65616e2e4f6d6567612e436f6e73747261696e742e616464457175616c6974795f736174}{P29.}] \refstepcounter{paperpremise}\label{premise:4c65616e2e4f6d6567612e436f6e73747261696e742e616464457175616c6974795f736174}\textbf{Sum equality sat.} Let \(c\) be an integer. Let \(x\) and \(y\) be elements of \(\operatorname{Coeffs}\). Assume \(c + \operatorname{dot}\left(x, y\right) = 0\). Then \(\operatorname{sat^{\prime}}(\operatorname{Lean\mathord{:}Omega\mathord{:}Constraint\mathord{:}mk}(\operatorname{Option\mathord{:}some}(-c),\allowbreak{} \operatorname{Option\mathord{:}some}(-c)),\allowbreak{} x,\allowbreak{} y) = \mathrm{true}\).
\item[\paperlineanchor{premise:4c65616e2e4f6d6567612e436f6e73747261696e742e616464496e657175616c6974795f736174}\paperlink{premise:4c65616e2e4f6d6567612e436f6e73747261696e742e616464496e657175616c6974795f736174}{P30.}] \refstepcounter{paperpremise}\label{premise:4c65616e2e4f6d6567612e436f6e73747261696e742e616464496e657175616c6974795f736174}\textbf{Sum inequality sat.} Let \(c\) be an integer. Let \(x\) and \(y\) be elements of \(\operatorname{Coeffs}\). Assume \(c + \operatorname{dot}\left(x, y\right) \ge 0\). Then \(\operatorname{sat^{\prime}}(\operatorname{Lean\mathord{:}Omega\mathord{:}Constraint\mathord{:}mk}(\operatorname{Option\mathord{:}some}(-c),\allowbreak{} \operatorname{none}),\allowbreak{} x,\allowbreak{} y) = \mathrm{true}\).
\item[\paperlineanchor{premise:4c65616e2e4f6d6567612e436f6e73747261696e742e636f6d62696e655f73617427}\paperlink{premise:4c65616e2e4f6d6567612e436f6e73747261696e742e636f6d62696e655f73617427}{P31.}] \refstepcounter{paperpremise}\label{premise:4c65616e2e4f6d6567612e436f6e73747261696e742e636f6d62696e655f73617427}\textbf{Combine sat'.} Let \(s\) and \(t\) be elements of \(\operatorname{Constraint}\). Let \(x\) and \(y\) be elements of \(\operatorname{Coeffs}\). Assume \(\operatorname{sat^{\prime}}\left(s, x, y\right) = \mathrm{true}\). Assume \(\operatorname{sat^{\prime}}\left(t, x, y\right) = \mathrm{true}\). Then \(\operatorname{sat^{\prime}}\left(\operatorname{combine}\left(s, t\right), x, y\right) = \mathrm{true}\).
\item[\paperlineanchor{premise:4c65616e2e4f6d6567612e436f6e73747261696e742e6e6f745f736174275f6f665f6973496d706f737369626c65}\paperlink{premise:4c65616e2e4f6d6567612e436f6e73747261696e742e6e6f745f736174275f6f665f6973496d706f737369626c65}{P32.}] \refstepcounter{paperpremise}\label{premise:4c65616e2e4f6d6567612e436f6e73747261696e742e6e6f745f736174275f6f665f6973496d706f737369626c65}\textbf{Not sat' of is impossible.} Let \(c\) be an element of \(\operatorname{Constraint}\). Assume \(\operatorname{is\,Impossible}\left(c\right) = \mathrm{true}\). Let \(x\) and \(y\) be elements of \(\operatorname{Coeffs}\). Then \(\neg \left(\operatorname{sat^{\prime}}\left(c, x, y\right) = \mathrm{true}\right)\).
\item[\paperlineanchor{premise:4c65616e2e4f6d6567612e496e742e6164645f636f6e6772}\paperlink{premise:4c65616e2e4f6d6567612e496e742e6164645f636f6e6772}{P33.}] \refstepcounter{paperpremise}\label{premise:4c65616e2e4f6d6567612e496e742e6164645f636f6e6772}\textbf{Sum congruence.} Let \(a\), \(b\), \(c\), and \(d\) be integers. Assume \(a = b\). Assume \(c = d\). Then \(a + c = b + d\).
\item[\paperlineanchor{premise:4c65616e2e4f6d6567612e496e742e6f664e61745f636f6e6772}\paperlink{premise:4c65616e2e4f6d6567612e496e742e6f664e61745f636f6e6772}{P34.}] \refstepcounter{paperpremise}\label{premise:4c65616e2e4f6d6567612e496e742e6f664e61745f636f6e6772}\textbf{Of nat congruence.} Let \(a\) and \(b\) be natural numbers. Assume \(a = b\). Then \(a = b\).
\item[\paperlineanchor{premise:4c65616e2e4f6d6567612e496e742e6f664e61745f6c655f6f665f6c65}\paperlink{premise:4c65616e2e4f6d6567612e496e742e6f664e61745f6c655f6f665f6c65}{P35.}] \refstepcounter{paperpremise}\label{premise:4c65616e2e4f6d6567612e496e742e6f664e61745f6c655f6f665f6c65}\textbf{Of nat at most of at most.} Let \(x\) and \(y\) be natural numbers. Assume \(x \le y\). Then \(x \le y\).
\item[\paperlineanchor{premise:4c65616e2e4f6d6567612e496e742e6f664e61745f6c745f6f665f6c74}\paperlink{premise:4c65616e2e4f6d6567612e496e742e6f664e61745f6c745f6f665f6c74}{P36.}] \refstepcounter{paperpremise}\label{premise:4c65616e2e4f6d6567612e496e742e6f664e61745f6c745f6f665f6c74}\textbf{Of nat less than of less than.} Let \(x\) and \(y\) be natural numbers. Assume \(x < y\). Then \(x < y\).
\item[\paperlineanchor{premise:4c65616e2e4f6d6567612e496e742e7375625f636f6e6772}\paperlink{premise:4c65616e2e4f6d6567612e496e742e7375625f636f6e6772}{P37.}] \refstepcounter{paperpremise}\label{premise:4c65616e2e4f6d6567612e496e742e7375625f636f6e6772}\textbf{Difference congruence.} Let \(a\), \(b\), \(c\), and \(d\) be integers. Assume \(a = b\). Assume \(c = d\). Then \(a - c = b - d\).
\item[\paperlineanchor{premise:4c65616e2e4f6d6567612e4c696e656172436f6d626f2e6164645f6576616c}\paperlink{premise:4c65616e2e4f6d6567612e4c696e656172436f6d626f2e6164645f6576616c}{P38.}] \refstepcounter{paperpremise}\label{premise:4c65616e2e4f6d6567612e4c696e656172436f6d626f2e6164645f6576616c}\textbf{Sum eval.} Let \(l_{1}\) and \(l_{2}\) be elements of \(\operatorname{Linear\,Combo}\). Let \(v\) be an element of \(\operatorname{Coeffs}\). Then \(\operatorname{eval}(l_{1} + l_{2},\allowbreak{} v) = \operatorname{eval}(l_{1},\allowbreak{} v) + \operatorname{eval}(l_{2},\allowbreak{} v)\).
\item[\paperlineanchor{premise:4c65616e2e4f6d6567612e4c696e656172436f6d626f2e636f6f7264696e6174655f6576616c5f30}\paperlink{premise:4c65616e2e4f6d6567612e4c696e656172436f6d626f2e636f6f7264696e6174655f6576616c5f30}{P39.}] \refstepcounter{paperpremise}\label{premise:4c65616e2e4f6d6567612e4c696e656172436f6d626f2e636f6f7264696e6174655f6576616c5f30}\textbf{Coordinate eval 0.} Let \(\mathit{a0}\) be an integer. Let \(t\) be a list of elements of \(\mathbb{Z}\). Then \(\operatorname{eval}(\operatorname{coordinate}(0),\allowbreak{} \operatorname{of\,List}(\langle \mathit{a0}\rangle \mathbin{\Vert} t)) = \mathit{a0}\).
\item[\paperlineanchor{premise:4c65616e2e4f6d6567612e4c696e656172436f6d626f2e636f6f7264696e6174655f6576616c5f31}\paperlink{premise:4c65616e2e4f6d6567612e4c696e656172436f6d626f2e636f6f7264696e6174655f6576616c5f31}{P40.}] \refstepcounter{paperpremise}\label{premise:4c65616e2e4f6d6567612e4c696e656172436f6d626f2e636f6f7264696e6174655f6576616c5f31}\textbf{Coordinate eval 1.} Let \(\mathit{a0}\) and \(\mathit{a1}\) be integers. Let \(t\) be a list of elements of \(\mathbb{Z}\). Then \(\operatorname{eval}(\operatorname{coordinate}(1),\allowbreak{} \operatorname{of\,List}(\langle \mathit{a0}\rangle \mathbin{\Vert} (\langle \mathit{a1}\rangle \mathbin{\Vert} t))) = \mathit{a1}\).
\item[\paperlineanchor{premise:4c65616e2e4f6d6567612e4c696e656172436f6d626f2e7375625f6576616c}\paperlink{premise:4c65616e2e4f6d6567612e4c696e656172436f6d626f2e7375625f6576616c}{P41.}] \refstepcounter{paperpremise}\label{premise:4c65616e2e4f6d6567612e4c696e656172436f6d626f2e7375625f6576616c}\textbf{Difference eval.} Let \(l_{1}\) and \(l_{2}\) be elements of \(\operatorname{Linear\,Combo}\). Let \(v\) be an element of \(\operatorname{Coeffs}\). Then \(\operatorname{eval}(l_{1} - l_{2},\allowbreak{} v) = \operatorname{eval}(l_{1},\allowbreak{} v) - \operatorname{eval}(l_{2},\allowbreak{} v)\).
\item[\paperlineanchor{premise:4c65616e2e4f6d6567612e616e645f6e6f745f6e6f745f6f665f6e6f745f6f72}\paperlink{premise:4c65616e2e4f6d6567612e616e645f6e6f745f6e6f745f6f665f6e6f745f6f72}{P42.}] \refstepcounter{paperpremise}\label{premise:4c65616e2e4f6d6567612e616e645f6e6f745f6e6f745f6f665f6e6f745f6f72}\textbf{And not not of not or.} Let \(p\) and \(q\) be propositions. Assume \(\neg \left(p \lor q\right)\). Then \(\neg p \land \neg q\).
\item[\paperlineanchor{premise:4c65616e2e4f6d6567612e636f6d626f5f73617427}\paperlink{premise:4c65616e2e4f6d6567612e636f6d626f5f73617427}{P43.}] \refstepcounter{paperpremise}\label{premise:4c65616e2e4f6d6567612e636f6d626f5f73617427}\textbf{Combo sat'.} Let \(s\) and \(t\) be elements of \(\operatorname{Constraint}\). Let \(a\) be an integer. Let \(x\) be an element of \(\operatorname{Coeffs}\). Let \(b\) be an integer. Let \(y\) and \(v\) be elements of \(\operatorname{Coeffs}\). Assume \(\operatorname{sat^{\prime}}\left(s, x, v\right) = \mathrm{true}\). Assume \(\operatorname{sat^{\prime}}\left(t, y, v\right) = \mathrm{true}\). Then \(\operatorname{sat^{\prime}}(\operatorname{Lean\mathord{:}Omega\mathord{:}Constraint\mathord{:}combo}(a,\allowbreak{} s,\allowbreak{} b,\allowbreak{} t),\allowbreak{} \operatorname{Lean\mathord{:}Omega\mathord{:}Coeffs\mathord{:}combo}(a,\allowbreak{} x,\allowbreak{} b,\allowbreak{} y),\allowbreak{} v) = \mathrm{true}\).
\item[\paperlineanchor{premise:4c65616e2e4f6d6567612e746964795f736174}\paperlink{premise:4c65616e2e4f6d6567612e746964795f736174}{P44.}] \refstepcounter{paperpremise}\label{premise:4c65616e2e4f6d6567612e746964795f736174}\textbf{Tidy sat.} Let \(s\) be an element of \(\operatorname{Constraint}\). Let \(x\) and \(v\) be elements of \(\operatorname{Coeffs}\). Assume \(\operatorname{sat^{\prime}}\left(s, x, v\right) = \mathrm{true}\). Then \(\operatorname{sat^{\prime}}(\operatorname{tidy\,Constraint}(s,\allowbreak{} x),\allowbreak{} \operatorname{tidy\,Coeffs}(s,\allowbreak{} x),\allowbreak{} v) = \mathrm{true}\).
\item[\paperlineanchor{premise:4c6973742e6d61705f69645f66756e}\paperlink{premise:4c6973742e6d61705f69645f66756e}{P45.}] \refstepcounter{paperpremise}\label{premise:4c6973742e6d61705f69645f66756e}\textbf{Map id fun.} Let \(\alpha\) be a type. Then \(\operatorname{List\mathord{:}map}\left(\operatorname{id}\right) = \operatorname{id}\).
\item[\paperlineanchor{premise:4d6174686c69622e5461637469632e507573684e65672e6e6f745f616e645f6571}\paperlink{premise:4d6174686c69622e5461637469632e507573684e65672e6e6f745f616e645f6571}{P46.}] \refstepcounter{paperpremise}\label{premise:4d6174686c69622e5461637469632e507573684e65672e6e6f745f616e645f6571}\textbf{Not and equals.} Let \(p\) and \(q\) be propositions. Then \(\left(\neg \left(p \land q\right)\right) = \left(p \implies \neg q\right)\).
\item[\paperlineanchor{premise:4d6174686c69622e5461637469632e507573684e65672e6e6f745f6578697374735f6571}\paperlink{premise:4d6174686c69622e5461637469632e507573684e65672e6e6f745f6578697374735f6571}{P47.}] \refstepcounter{paperpremise}\label{premise:4d6174686c69622e5461637469632e507573684e65672e6e6f745f6578697374735f6571}\textbf{Not exists equals.} Let \(\alpha\) be a sort. Let \(s\) be a function of type \(\alpha \to \mathsf{Prop}\). Then \((\neg (\exists x\, :\, \alpha,\; s(x))) = (\forall x\, :\, \alpha,\; \neg s(x))\).
\item[\paperlineanchor{premise:4d6178696d616c2e6e6f745f70726f705f6f665f6774}\paperlink{premise:4d6178696d616c2e6e6f745f70726f705f6f665f6774}{P48.}] \refstepcounter{paperpremise}\label{premise:4d6178696d616c2e6e6f745f70726f705f6f665f6774}\textbf{Not prop of greater than.} Let \(\alpha\) be a type. Let \(P\) be a function of type \(\alpha \to \mathsf{Prop}\). Let \(x\) and \(y\) be elements of \(\alpha\). Equip \(\alpha\) with a preorder. Assume \(\operatorname{Maximal}\left(P, x\right)\). Assume \(x < y\). Then \(\neg P\left(y\right)\).
\item[\paperlineanchor{premise:4d6178696d616c2e70726f70}\paperlink{premise:4d6178696d616c2e70726f70}{P49.}] \refstepcounter{paperpremise}\label{premise:4d6178696d616c2e70726f70}\textbf{Prop.} Let \(\alpha\) be a type. Fix a structure \(\mathit{inst}\) of type \(\operatorname{LE}\left(\alpha\right)\). Let \(P\) be a function of type \(\alpha \to \mathsf{Prop}\). Let \(x\) be an element of \(\alpha\). Assume \(\operatorname{Maximal}\left(P, x\right)\). Then \(P\left(x\right)\).
\item[\paperlineanchor{premise:4e61742e636172645f65715f66696e747970655f63617264}\paperlink{premise:4e61742e636172645f65715f66696e747970655f63617264}{P50.}] \refstepcounter{paperpremise}\label{premise:4e61742e636172645f65715f66696e747970655f63617264}\textbf{Cardinality equals fintype cardinality.} Let \(\alpha\) be a type. Fix a finite enumeration of \(\alpha\). Then \(\operatorname{Nat\mathord{:}card}\left(\alpha\right) = \left|\alpha\right|\).
\item[\paperlineanchor{premise:4e61742e6c655f6f665f6e6f745f6c74}\paperlink{premise:4e61742e6c655f6f665f6e6f745f6c74}{P51.}] \refstepcounter{paperpremise}\label{premise:4e61742e6c655f6f665f6e6f745f6c74}\textbf{At most of not less than.} Let \(a\) and \(b\) be natural numbers. Assume \(\neg \left(a < b\right)\). Then \(b \le a\).
\item[\paperlineanchor{premise:4e61742e6c745f6f725f67745f6f665f6e65}\paperlink{premise:4e61742e6c745f6f725f67745f6f665f6e65}{P52.}] \refstepcounter{paperpremise}\label{premise:4e61742e6c745f6f725f67745f6f665f6e65}\textbf{Less than or greater than of differs from.} Let \(a\) and \(b\) be natural numbers. Assume \(a \ne b\). Then \(a < b \lor a > b\).
\item[\paperlineanchor{premise:4e61742e6e6f745f6576656e5f6966665f6f6464}\paperlink{premise:4e61742e6e6f745f6576656e5f6966665f6f6464}{P53.}] \refstepcounter{paperpremise}\label{premise:4e61742e6e6f745f6576656e5f6966665f6f6464}\textbf{Not even if and only if odd.} Let \(n\) be a natural number. Then \(\neg \operatorname{Even}\left(n\right) \iff \operatorname{Odd}\left(n\right)\).
\item[\paperlineanchor{premise:4e652e73796d6d}\paperlink{premise:4e652e73796d6d}{P54.}] \refstepcounter{paperpremise}\label{premise:4e652e73796d6d}\textbf{Symm.} Let \(\alpha\) be a sort. Let \(a\) and \(b\) be elements of \(\alpha\). Assume \(a \ne b\). Then \(b \ne a\).
\item[\paperlineanchor{premise:4f72646572546f702e6c655f746f70}\paperlink{premise:4f72646572546f702e6c655f746f70}{P55.}] \refstepcounter{paperpremise}\label{premise:4f72646572546f702e6c655f746f70}\textbf{At most top.} Let \(\alpha\) be a type. Let \(\mathit{inst}\) be an element of \(\operatorname{LE}\left(\alpha\right)\). Fix a structure \(\mathit{self}\) of type \(\operatorname{Order\,Top}\left(\alpha\right)\). Let \(a\) be an element of \(\alpha\). Then \(a \le \operatorname{top}\).
\item[\paperlineanchor{premise:50726f642e6d6b2e696e6a4571}\paperlink{premise:50726f642e6d6b2e696e6a4571}{P56.}] \refstepcounter{paperpremise}\label{premise:50726f642e6d6b2e696e6a4571}\textbf{Injectivity equals.} Let \(\alpha\) and \(\beta\) be types. Let \(\mathit{fst}\) be an element of \(\alpha\). Let \(\mathit{snd}\) be an element of \(\beta\). Let \(\mathit{fst}_{1}\) be an element of \(\alpha\). Let \(\mathit{snd}_{1}\) be an element of \(\beta\). Then \(((\mathit{fst},\allowbreak{} \mathit{snd}) = (\mathit{fst}_{1},\allowbreak{} \mathit{snd}_{1})) = (\mathit{fst} = \mathit{fst}_{1} \land \mathit{snd} = \mathit{snd}_{1})\).
\item[\paperlineanchor{premise:52656c49736f2e696e6a656374697665}\paperlink{premise:52656c49736f2e696e6a656374697665}{P57.}] \refstepcounter{paperpremise}\label{premise:52656c49736f2e696e6a656374697665}\textbf{Injective.} Let \(\alpha\) and \(\beta\) be types. Let \(r\) be a relation from \(\alpha\) to \(\alpha\). Let \(s\) be a relation from \(\beta\) to \(\beta\). Let \(e\) be an element of \(\operatorname{Rel\,Iso}\left(r, s\right)\). Then \(\operatorname{DFun\,Like\mathord{:}coe}\left(e\right)\text{ is injective}\).
\item[\paperlineanchor{premise:52656c6174696f6e2e6d61705f69645f6964}\paperlink{premise:52656c6174696f6e2e6d61705f69645f6964}{P58.}] \refstepcounter{paperpremise}\label{premise:52656c6174696f6e2e6d61705f69645f6964}\textbf{Map id id.} Let \(\alpha\) and \(\beta\) be types. Let \(r\) be a relation from \(\alpha\) to \(\beta\). Then \(\operatorname{Map}\left(r, \operatorname{id}, \operatorname{id}\right) = r\).
\item[\paperlineanchor{premise:5365742e4e61742e636172645f636f655f7365745f6571}\paperlink{premise:5365742e4e61742e636172645f636f655f7365745f6571}{P59.}] \refstepcounter{paperpremise}\label{premise:5365742e4e61742e636172645f636f655f7365745f6571}\textbf{Cardinality coercion set equals.} Let \(\alpha\) be a type. Let \(s\) be a subset of \(\alpha\). Then \(\operatorname{Nat\mathord{:}card}\left(s\right) = \operatorname{ncard}(s)\).
\item[\paperlineanchor{premise:5365742e636173745f6e63617264}\paperlink{premise:5365742e636173745f6e63617264}{P60.}] \refstepcounter{paperpremise}\label{premise:5365742e636173745f6e63617264}\textbf{Cast ncard.} Let \(\alpha\) be a type. Let \(s\) be a subset of \(\alpha\). Assume \(\operatorname{Set\mathord{:}Finite}\left(s\right)\). Then \(\operatorname{cast}\left(\operatorname{ncard}(s)\right) = \operatorname{Cardinal\mathord{:}mk}\left(s\right)\).
\item[\paperlineanchor{premise:5365742e6469736a6f696e745f6c656674}\paperlink{premise:5365742e6469736a6f696e745f6c656674}{P61.}] \refstepcounter{paperpremise}\label{premise:5365742e6469736a6f696e745f6c656674}\textbf{Disjoint left.} Let \(\alpha\) be a type. Let \(s\) and \(t\) be subsets of \(\alpha\). Then \(\operatorname{Disjoint}(s,\allowbreak{} t) \iff (\forall a\, :\, \alpha,\; a \in s \implies a \notin t)\).
\item[\paperlineanchor{premise:5365742e65715f656d7074795f6966665f666f72616c6c5f6e6f745f6d656d}\paperlink{premise:5365742e65715f656d7074795f6966665f666f72616c6c5f6e6f745f6d656d}{P62.}] \refstepcounter{paperpremise}\label{premise:5365742e65715f656d7074795f6966665f666f72616c6c5f6e6f745f6d656d}\textbf{Equals empty if and only if forall not membership.} Let \(\alpha\) be a type. Let \(s\) be a subset of \(\alpha\). Then \(s = \varnothing \iff \left(\forall x\, :\, \alpha,\; x \notin s\right)\).
\item[\paperlineanchor{premise:5365742e6578697374735f6e655f6f665f6f6e655f6c745f6e63617264}\paperlink{premise:5365742e6578697374735f6e655f6f665f6f6e655f6c745f6e63617264}{P63.}] \refstepcounter{paperpremise}\label{premise:5365742e6578697374735f6e655f6f665f6f6e655f6c745f6e63617264}\textbf{Exists differs from of one less than ncard.} Let \(\alpha\) be a type. Let \(s\) be a subset of \(\alpha\). Assume \(1 < \operatorname{ncard}(s)\). Let \(a\) be an element of \(\alpha\). Then \(\exists b\, :\, \alpha,\; b \in s \land b \ne a\).
\item[\paperlineanchor{premise:5365742e6578697374735f756e696f6e5f6469736a6f696e745f6e636172645f65715f6f665f6576656e}\paperlink{premise:5365742e6578697374735f756e696f6e5f6469736a6f696e745f6e636172645f65715f6f665f6576656e}{P64.}] \refstepcounter{paperpremise}\label{premise:5365742e6578697374735f756e696f6e5f6469736a6f696e745f6e636172645f65715f6f665f6576656e}\textbf{Exists union disjoint ncard equals of even.} Let \(\alpha\) be a type. Let \(s\) be a subset of \(\alpha\). Assume \(\operatorname{Even}\left(\operatorname{ncard}(s)\right)\). Then \(\exists t\, :\, \mathcal{P}(\alpha),\; \exists u\, :\, \mathcal{P}(\alpha),\; t \cup u = s \land (\operatorname{Disjoint}(t,\allowbreak{} u) \land \operatorname{ncard}(t) = \operatorname{ncard}(u))\).
\item[\paperlineanchor{premise:5365742e657874}\paperlink{premise:5365742e657874}{P65.}] \refstepcounter{paperpremise}\label{premise:5365742e657874}\textbf{Ext.} Let \(\alpha\) be a type. Let \(a\) and \(b\) be subsets of \(\alpha\). Assume \(\forall x\, :\, \alpha,\; x \in a \iff x \in b\). Then \(a = b\).
\item[\paperlineanchor{premise:5365742e66696e6974655f756e696f6e}\paperlink{premise:5365742e66696e6974655f756e696f6e}{P66.}] \refstepcounter{paperpremise}\label{premise:5365742e66696e6974655f756e696f6e}\textbf{Finite union.} Let \(\alpha\) be a type. Let \(s\) and \(t\) be subsets of \(\alpha\). Then \(\operatorname{Set\mathord{:}Finite}(s \cup t) \iff \operatorname{Set\mathord{:}Finite}(s) \land \operatorname{Set\mathord{:}Finite}(t)\).
\item[\paperlineanchor{premise:5365742e696d6167655f636f6e6772}\paperlink{premise:5365742e696d6167655f636f6e6772}{P67.}] \refstepcounter{paperpremise}\label{premise:5365742e696d6167655f636f6e6772}\textbf{Image congruence.} Let \(\alpha\) and \(\beta\) be types. Let \(f\) and \(g\) be functions of type \(\alpha \to \beta\). Let \(s\) be a subset of \(\alpha\). Assume \(\forall a\, :\, \alpha,\; a \in s \implies f\left(a\right) = g\left(a\right)\). Then \(f\left[s\right] = g\left[s\right]\).
\item[\paperlineanchor{premise:5365742e696d6167655f696427}\paperlink{premise:5365742e696d6167655f696427}{P68.}] \refstepcounter{paperpremise}\label{premise:5365742e696d6167655f696427}\textbf{Image id'.} Let \(\alpha\) be a type. Let \(s\) be a subset of \(\alpha\). Then \(\left(x \mapsto x\right)\left[s\right] = s\).
\item[\paperlineanchor{premise:5365742e696e73744973416e746973796d6d537562736574}\paperlink{premise:5365742e696e73744973416e746973796d6d537562736574}{P69.}] \refstepcounter{paperpremise}\label{premise:5365742e696e73744973416e746973796d6d537562736574}\textbf{Inst is antisymm subset.} Let \(\alpha\) be a type. Then \(\operatorname{Is\,Antisymm}(\mathcal{P}(\alpha),\allowbreak{} \mathit{x1} \mapsto \mathit{x2} \mapsto \mathit{x1} \subseteq \mathit{x2})\).
\item[\paperlineanchor{premise:5365742e6d656d5f64696666}\paperlink{premise:5365742e6d656d5f64696666}{P70.}] \refstepcounter{paperpremise}\label{premise:5365742e6d656d5f64696666}\textbf{Membership diff.} Let \(\alpha\) be a type. Let \(s\) and \(t\) be subsets of \(\alpha\). Let \(x\) be an element of \(\alpha\). Then \(x \in s \setminus t \iff x \in s \land x \notin t\).
\item[\paperlineanchor{premise:5365742e6d656d5f69556e696f6e}\paperlink{premise:5365742e6d656d5f69556e696f6e}{P71.}] \refstepcounter{paperpremise}\label{premise:5365742e6d656d5f69556e696f6e}\textbf{Membership i union.} Let \(\alpha\) be a type. Let \(\iota\) be a sort. Let \(x\) be an element of \(\alpha\). Let \(s\) be a function of type \(\iota \to \mathcal{P}(\alpha)\). Then \(x \in \bigcup_{i\, :\, \iota} s\left(i\right) \iff \left(\exists i\, :\, \iota,\; x \in s\left(i\right)\right)\).
\item[\paperlineanchor{premise:5365742e6d656d5f696d616765}\paperlink{premise:5365742e6d656d5f696d616765}{P72.}] \refstepcounter{paperpremise}\label{premise:5365742e6d656d5f696d616765}\textbf{Membership image.} Let \(\alpha\) and \(\beta\) be types. Let \(f\) be a function of type \(\alpha \to \beta\). Let \(s\) be a subset of \(\alpha\). Let \(y\) be an element of \(\beta\). Then \(y \in f\left[s\right] \iff \left(\exists x\, :\, \alpha,\; x \in s \land f\left(x\right) = y\right)\).
\item[\paperlineanchor{premise:5365742e6d656d5f696e736572745f696666}\paperlink{premise:5365742e6d656d5f696e736572745f696666}{P73.}] \refstepcounter{paperpremise}\label{premise:5365742e6d656d5f696e736572745f696666}\textbf{Membership insert if and only if.} Let \(\alpha\) be a type. Let \(x\) and \(a\) be elements of \(\alpha\). Let \(s\) be a subset of \(\alpha\). Then \(x \in \{a\} \cup s \iff x = a \lor x \in s\).
\item[\paperlineanchor{premise:5365742e6d656d5f696e7465725f696666}\paperlink{premise:5365742e6d656d5f696e7465725f696666}{P74.}] \refstepcounter{paperpremise}\label{premise:5365742e6d656d5f696e7465725f696666}\textbf{Membership inter if and only if.} Let \(\alpha\) be a type. Let \(x\) be an element of \(\alpha\). Let \(a\) and \(b\) be subsets of \(\alpha\). Then \(x \in a \cap b \iff x \in a \land x \in b\).
\item[\paperlineanchor{premise:5365742e6d656d5f6f665f6d656d5f696d6167655f76616c}\paperlink{premise:5365742e6d656d5f6f665f6d656d5f696d6167655f76616c}{P75.}] \refstepcounter{paperpremise}\label{premise:5365742e6d656d5f6f665f6d656d5f696d6167655f76616c}\textbf{Membership of membership image val.} Let \(\alpha\) be a type. Let \(\beta\) be a subset of \(\alpha\). Let \(\gamma\) be a subset of \(\beta\). Let \(a\) be an element of \(\alpha\). Let \(\mathit{ha}\) be a proof of \(a \in \operatorname{val}\left[\gamma\right]\). Then \(\left(a\, :\, \beta\right) \in \gamma\).
\item[\paperlineanchor{premise:5365742e6d656d5f6f725f6d656d5f6f665f6d656d5f756e696f6e}\paperlink{premise:5365742e6d656d5f6f725f6d656d5f6f665f6d656d5f756e696f6e}{P76.}] \refstepcounter{paperpremise}\label{premise:5365742e6d656d5f6f725f6d656d5f6f665f6d656d5f756e696f6e}\textbf{Membership or membership of membership union.} Let \(\alpha\) be a type. Let \(x\) be an element of \(\alpha\). Let \(a\) and \(b\) be subsets of \(\alpha\). Assume \(x \in a \cup b\). Then \(x \in a \lor x \in b\).
\item[\paperlineanchor{premise:5365742e6d656d5f73696e676c65746f6e5f696666}\paperlink{premise:5365742e6d656d5f73696e676c65746f6e5f696666}{P77.}] \refstepcounter{paperpremise}\label{premise:5365742e6d656d5f73696e676c65746f6e5f696666}\textbf{Membership singleton if and only if.} Let \(\alpha\) be a type. Let \(a\) and \(b\) be elements of \(\alpha\). Then \(a \in \{b\} \iff a = b\).
\item[\paperlineanchor{premise:5365742e6d656d5f756e696f6e}\paperlink{premise:5365742e6d656d5f756e696f6e}{P78.}] \refstepcounter{paperpremise}\label{premise:5365742e6d656d5f756e696f6e}\textbf{Membership union.} Let \(\alpha\) be a type. Let \(x\) be an element of \(\alpha\). Let \(a\) and \(b\) be subsets of \(\alpha\). Then \(x \in a \cup b \iff x \in a \lor x \in b\).
\item[\paperlineanchor{premise:5365742e6d656d5f756e6976}\paperlink{premise:5365742e6d656d5f756e6976}{P79.}] \refstepcounter{paperpremise}\label{premise:5365742e6d656d5f756e6976}\textbf{Membership whole space.} Let \(\alpha\) be a type. Let \(x\) be an element of \(\alpha\). Then \(x \in \alpha\).
\item[\paperlineanchor{premise:5365742e6e636172645f65715f746f46696e7365745f63617264}\paperlink{premise:5365742e6e636172645f65715f746f46696e7365745f63617264}{P80.}] \refstepcounter{paperpremise}\label{premise:5365742e6e636172645f65715f746f46696e7365745f63617264}\textbf{Ncard equals to finset cardinality.} Let \(\alpha\) be a type. Let \(s\) be a subset of \(\alpha\). Let \(\mathit{hs}\) be a proof of \(\operatorname{Set\mathord{:}Finite}\left(s\right)\). Then \(\operatorname{ncard}(s) = \left|\operatorname{as\,finite\,set}(s)\right|\).
\item[\paperlineanchor{premise:5365742e6e636172645f65715f74776f}\paperlink{premise:5365742e6e636172645f65715f74776f}{P81.}] \refstepcounter{paperpremise}\label{premise:5365742e6e636172645f65715f74776f}\textbf{Ncard equals two.} Let \(\alpha\) be a type. Let \(s\) be a subset of \(\alpha\). Then \(\operatorname{ncard}(s) = 2 \iff (\exists x\, :\, \alpha,\; \exists y\, :\, \alpha,\; x \ne y \land s = \{x\} \cup \{y\})\).
\item[\paperlineanchor{premise:5365742e6e6f6e656d7074795f6f665f6d656d}\paperlink{premise:5365742e6e6f6e656d7074795f6f665f6d656d}{P82.}] \refstepcounter{paperpremise}\label{premise:5365742e6e6f6e656d7074795f6f665f6d656d}\textbf{Nonempty of membership.} Let \(\alpha\) be a type. Let \(s\) be a subset of \(\alpha\). Let \(x\) be an element of \(\alpha\). Assume \(x \in s\). Then \(s\ne\varnothing\).
\item[\paperlineanchor{premise:5365742e6e6f745f6e6f6e656d7074795f6966665f65715f656d707479}\paperlink{premise:5365742e6e6f745f6e6f6e656d7074795f6966665f65715f656d707479}{P83.}] \refstepcounter{paperpremise}\label{premise:5365742e6e6f745f6e6f6e656d7074795f6966665f65715f656d707479}\textbf{Not nonempty if and only if equals empty.} Let \(\alpha\) be a type. Let \(s\) be a subset of \(\alpha\). Then \(\neg \left(s\ne\varnothing\right) \iff s = \varnothing\).
\item[\paperlineanchor{premise:5365742e746f46696e697465}\paperlink{premise:5365742e746f46696e697465}{P84.}] \refstepcounter{paperpremise}\label{premise:5365742e746f46696e697465}\textbf{To finite.} Let \(\alpha\) be a type. Let \(s\) be a subset of \(\alpha\). Assume \(s\) is finite. Then \(\operatorname{Set\mathord{:}Finite}\left(s\right)\).
\item[\paperlineanchor{premise:5365742e746f46696e7365745f63617264}\paperlink{premise:5365742e746f46696e7365745f63617264}{P85.}] \refstepcounter{paperpremise}\label{premise:5365742e746f46696e7365745f63617264}\textbf{To finset cardinality.} Let \(\alpha\) be a type. Let \(s\) be a subset of \(\alpha\). Fix a finite enumeration of \(s\). Then \(\left|\operatorname{as\,finite\,set}(s)\right| = \left|s\right|\).
\item[\paperlineanchor{premise:5365742e746f46696e7365745f636f6e6772}\paperlink{premise:5365742e746f46696e7365745f636f6e6772}{P86.}] \refstepcounter{paperpremise}\label{premise:5365742e746f46696e7365745f636f6e6772}\textbf{To finset congruence.} Let \(\alpha\) be a type. Let \(s\) and \(t\) be subsets of \(\alpha\). Fix a finite enumeration of \(s\). Fix a finite enumeration of \(t\). Assume \(s = t\). Then \(\operatorname{as\,finite\,set}(s) = \operatorname{as\,finite\,set}(t)\).
\item[\paperlineanchor{premise:5365742e746f46696e7365745f696d616765}\paperlink{premise:5365742e746f46696e7365745f696d616765}{P87.}] \refstepcounter{paperpremise}\label{premise:5365742e746f46696e7365745f696d616765}\textbf{To finset image.} Let \(\alpha\) and \(\beta\) be types. Assume equality on \(\beta\) is decidable. Let \(f\) be a function of type \(\alpha \to \beta\). Let \(s\) be a subset of \(\alpha\). Fix a finite enumeration of \(s\). Fix a finite enumeration of \(f\left[s\right]\). Then \(\operatorname{as\,finite\,set}(f\left[s\right]) = f\left[\operatorname{as\,finite\,set}(s)\right]\).
\item[\paperlineanchor{premise:5365742e756e69765f696e746572}\paperlink{premise:5365742e756e69765f696e746572}{P88.}] \refstepcounter{paperpremise}\label{premise:5365742e756e69765f696e746572}\textbf{Whole space inter.} Let \(\alpha\) be a type. Let \(a\) be a subset of \(\alpha\). Then \(\alpha \cap a = a\).
\item[\paperlineanchor{premise:53696d706c6547726170682e41646a2e726561636861626c65}\paperlink{premise:53696d706c6547726170682e41646a2e726561636861626c65}{P89.}] \refstepcounter{paperpremise}\label{premise:53696d706c6547726170682e41646a2e726561636861626c65}\textbf{Reachable.} Let \(V\) be a type. Let \(G\) be a simple graph on \(V\). Let \(u\) and \(v\) be elements of \(V\). Assume \(u \sim_{G} v\). Then \(u\text{ can reach }v\text{ in }G\).
\item[\paperlineanchor{premise:53696d706c6547726170682e41646a2e73796d6d}\paperlink{premise:53696d706c6547726170682e41646a2e73796d6d}{P90.}] \refstepcounter{paperpremise}\label{premise:53696d706c6547726170682e41646a2e73796d6d}\textbf{Symm.} Let \(V\) be a type. Let \(G\) be a simple graph on \(V\). Let \(u\) and \(v\) be elements of \(V\). Assume \(u \sim_{G} v\). Then \(v \sim_{G} u\).
\item[\paperlineanchor{premise:53696d706c6547726170682e436f6e6e6563746564436f6d706f6e656e742e61646a5f7370616e6e696e67436f655f696e647563655f73757070}\paperlink{premise:53696d706c6547726170682e436f6e6e6563746564436f6d706f6e656e742e61646a5f7370616e6e696e67436f655f696e647563655f73757070}{P91.}] \refstepcounter{paperpremise}\label{premise:53696d706c6547726170682e436f6e6e6563746564436f6d706f6e656e742e61646a5f7370616e6e696e67436f655f696e647563655f73757070}\textbf{Adj spanning coercion induce supp.} Let \(V\) be a type. Let \(G\) be a simple graph on \(V\). Let \(v\) and \(w\) be elements of \(V\). Let \(c\) be an element of \(\operatorname{Connected\,Component}\left(G\right)\). Then \(v \sim_{\operatorname{Simple\,Graph\mathord{:}spanning\,Coe}(\operatorname{Simple\,Graph\mathord{:}induce}(\operatorname{supp}(c),\allowbreak{} G))} w \iff v \in \operatorname{supp}(c) \land v \sim_{G} w\).
\item[\paperlineanchor{premise:53696d706c6547726170682e436f6e6e6563746564436f6d706f6e656e742e6571}\paperlink{premise:53696d706c6547726170682e436f6e6e6563746564436f6d706f6e656e742e6571}{P92.}] \refstepcounter{paperpremise}\label{premise:53696d706c6547726170682e436f6e6e6563746564436f6d706f6e656e742e6571}\textbf{Equals.} Let \(V\) be a type. Let \(G\) be a simple graph on \(V\). Let \(v\) and \(w\) be elements of \(V\). Then \(\operatorname{connected\,Component\,Mk}(G,\allowbreak{} v) = \operatorname{connected\,Component\,Mk}(G,\allowbreak{} w) \iff v\text{ can reach }w\text{ in }G\).
\item[\paperlineanchor{premise:53696d706c6547726170682e436f6e6e6563746564436f6d706f6e656e742e6d656d5f636f655f737570705f6f665f61646a}\paperlink{premise:53696d706c6547726170682e436f6e6e6563746564436f6d706f6e656e742e6d656d5f636f655f737570705f6f665f61646a}{P93.}] \refstepcounter{paperpremise}\label{premise:53696d706c6547726170682e436f6e6e6563746564436f6d706f6e656e742e6d656d5f636f655f737570705f6f665f61646a}\textbf{Membership coercion supp of adj.} Let \(V\) be a type. Let \(G\) be a simple graph on \(V\). Let \(v\) and \(w\) be elements of \(V\). Let \(H\) be a subgraph of \(G\). Let \(c\) be an element of \(\operatorname{Connected\,Component}\left(\operatorname{graph}(H)\right)\). Assume \(v \in \operatorname{val}\left[\operatorname{Set\,Like\mathord{:}coe}\left(c\right)\right]\). Assume \(w \in V(H)\). Assume \(v \sim_{H} w\). Then \(w \in \operatorname{val}\left[\operatorname{Set\,Like\mathord{:}coe}\left(c\right)\right]\).
\item[\paperlineanchor{premise:53696d706c6547726170682e436f6e6e6563746564436f6d706f6e656e742e6d656d5f737570705f696666}\paperlink{premise:53696d706c6547726170682e436f6e6e6563746564436f6d706f6e656e742e6d656d5f737570705f696666}{P94.}] \refstepcounter{paperpremise}\label{premise:53696d706c6547726170682e436f6e6e6563746564436f6d706f6e656e742e6d656d5f737570705f696666}\textbf{Membership supp if and only if.} Let \(V\) be a type. Let \(G\) be a simple graph on \(V\). Let \(C\) be an element of \(\operatorname{Connected\,Component}\left(G\right)\). Let \(v\) be an element of \(V\). Then \(v \in \operatorname{supp}\left(C\right) \iff \operatorname{connected\,Component\,Mk}\left(G, v\right) = C\).
\item[\paperlineanchor{premise:53696d706c6547726170682e49736f2e73796d6d5f746f486f6d5f636f6d705f746f486f6d}\paperlink{premise:53696d706c6547726170682e49736f2e73796d6d5f746f486f6d5f636f6d705f746f486f6d}{P95.}] \refstepcounter{paperpremise}\label{premise:53696d706c6547726170682e49736f2e73796d6d5f746f486f6d5f636f6d705f746f486f6d}\textbf{Symm to hom comp to hom.} Let \(V\) and \(W\) be types. Let \(G\) be a simple graph on \(V\). Let \(G^{'}\) be a simple graph on \(W\). Let \(f\) be a graph isomorphism from \(G\) to \(G^{'}\). Then \(\operatorname{comp}(\operatorname{to\,Hom}(f^{-1}),\allowbreak{} \operatorname{to\,Hom}(f)) = \operatorname{Simple\,Graph\mathord{:}Hom\mathord{:}id}\).
\item[\paperlineanchor{premise:53696d706c6547726170682e526561636861626c652e73796d6d}\paperlink{premise:53696d706c6547726170682e526561636861626c652e73796d6d}{P96.}] \refstepcounter{paperpremise}\label{premise:53696d706c6547726170682e526561636861626c652e73796d6d}\textbf{Symm.} Let \(V\) be a type. Let \(G\) be a simple graph on \(V\). Let \(u\) and \(v\) be elements of \(V\). Assume \(u\text{ can reach }v\text{ in }G\). Then \(v\text{ can reach }u\text{ in }G\).
\item[\paperlineanchor{premise:53696d706c6547726170682e53756267726170682e41646a2e61646a5f737562}\paperlink{premise:53696d706c6547726170682e53756267726170682e41646a2e61646a5f737562}{P97.}] \refstepcounter{paperpremise}\label{premise:53696d706c6547726170682e53756267726170682e41646a2e61646a5f737562}\textbf{Adj difference.} Let \(V\) be a type. Let \(G\) be a simple graph on \(V\). Let \(H\) be a subgraph of \(G\). Let \(u\) and \(v\) be elements of \(V\). Assume \(u \sim_{H} v\). Then \(u \sim_{G} v\).
\item[\paperlineanchor{premise:53696d706c6547726170682e53756267726170682e41646a2e73796d6d}\paperlink{premise:53696d706c6547726170682e53756267726170682e41646a2e73796d6d}{P98.}] \refstepcounter{paperpremise}\label{premise:53696d706c6547726170682e53756267726170682e41646a2e73796d6d}\textbf{Symm.} Let \(V\) be a type. Let \(G\) be a simple graph on \(V\). Let \(G^{'}\) be a subgraph of \(G\). Let \(u\) and \(v\) be elements of \(V\). Assume \(u \sim_{G^{'}} v\). Then \(v \sim_{G^{'}} u\).
\item[\paperlineanchor{premise:53696d706c6547726170682e53756267726170682e436f6e6e65637465642e6578697374735f76657274735f65715f636f6e6e6563746564436f6d706f6e656e7453757070}\paperlink{premise:53696d706c6547726170682e53756267726170682e436f6e6e65637465642e6578697374735f76657274735f65715f636f6e6e6563746564436f6d706f6e656e7453757070}{P99.}] \refstepcounter{paperpremise}\label{premise:53696d706c6547726170682e53756267726170682e436f6e6e65637465642e6578697374735f76657274735f65715f636f6e6e6563746564436f6d706f6e656e7453757070}\textbf{Exists verts equals connected component supp.} Let \(V\) be a type. Let \(G\) be a simple graph on \(V\). Let \(H\) be a subgraph of \(G\). Assume \(\operatorname{Connected}\left(H\right)\). Assume \(\forall v\, :\, V,\; v \in V(H) \implies \left(\forall w\, :\, V,\; v \sim_{G} w \implies v \sim_{H} w\right)\). Then \(\exists c\, :\, \operatorname{Connected\,Component}\left(G\right),\; V(H) = \operatorname{supp}\left(c\right)\).
\item[\paperlineanchor{premise:53696d706c6547726170682e53756267726170682e49735370616e6e696e672e636172645f7665727473}\paperlink{premise:53696d706c6547726170682e53756267726170682e49735370616e6e696e672e636172645f7665727473}{P100.}] \refstepcounter{paperpremise}\label{premise:53696d706c6547726170682e53756267726170682e49735370616e6e696e672e636172645f7665727473}\textbf{Cardinality verts.} Let \(V\) be a type. Let \(G\) be a simple graph on \(V\). Fix a finite enumeration of \(V\). Let \(G^{'}\) be a subgraph of \(G\). Fix a finite enumeration of \(V(G^{'})\). Assume \(\operatorname{Is\,Spanning}\left(G^{'}\right)\). Then \(\left|\operatorname{as\,finite\,set}(V(G^{'}))\right| = \left|V\right|\).
\item[\paperlineanchor{premise:53696d706c6547726170682e53756267726170682e49735370616e6e696e672e76657274735f65715f756e6976}\paperlink{premise:53696d706c6547726170682e53756267726170682e49735370616e6e696e672e76657274735f65715f756e6976}{P101.}] \refstepcounter{paperpremise}\label{premise:53696d706c6547726170682e53756267726170682e49735370616e6e696e672e76657274735f65715f756e6976}\textbf{Verts equals whole space.} Let \(V\) be a type. Let \(G\) be a simple graph on \(V\). Let \(G^{'}\) be a subgraph of \(G\). Assume \(\operatorname{Is\,Spanning}\left(G^{'}\right)\). Then \(V(G^{'}) = V\).
\item[\paperlineanchor{premise:53696d706c6547726170682e53756267726170682e61646a5f636f6d6d}\paperlink{premise:53696d706c6547726170682e53756267726170682e61646a5f636f6d6d}{P102.}] \refstepcounter{paperpremise}\label{premise:53696d706c6547726170682e53756267726170682e61646a5f636f6d6d}\textbf{Adj commutation.} Let \(V\) be a type. Let \(G\) be a simple graph on \(V\). Let \(G^{'}\) be a subgraph of \(G\). Let \(v\) and \(w\) be elements of \(V\). Then \(v \sim_{G^{'}} w \iff w \sim_{G^{'}} v\).
\item[\paperlineanchor{premise:53696d706c6547726170682e53756267726170682e61646a5f6966665f6f665f6e65696768626f725365745f6571756976}\paperlink{premise:53696d706c6547726170682e53756267726170682e61646a5f6966665f6f665f6e65696768626f725365745f6571756976}{P103.}] \refstepcounter{paperpremise}\label{premise:53696d706c6547726170682e53756267726170682e61646a5f6966665f6f665f6e65696768626f725365745f6571756976}\textbf{Adj if and only if of neighbor set equiv.} Let \(V\) be a type. Let \(G\) be a simple graph on \(V\). Let \(v\) be an element of \(V\). Let \(H\) be a subgraph of \(G\). Let \(h\) be a bijection from \(N_{G}(v)\) to \(\operatorname{Simple\,Graph\mathord{:}Subgraph\mathord{:}neighbor\,Set}\left(H, v\right)\). Assume \(\operatorname{Set\mathord{:}Finite}\left(N_{G}(v)\right)\). Let \(w\) be an element of \(V\). Then \(v \sim_{H} w \iff v \sim_{G} w\).
\item[\paperlineanchor{premise:53696d706c6547726170682e53756267726170682e61646a5f737562}\paperlink{premise:53696d706c6547726170682e53756267726170682e61646a5f737562}{P104.}] \refstepcounter{paperpremise}\label{premise:53696d706c6547726170682e53756267726170682e61646a5f737562}\textbf{Adj difference.} Let \(V\) be a type. Let \(G\) be a simple graph on \(V\). Let \(\mathit{self}\) be a subgraph of \(G\). Let \(v\) and \(w\) be elements of \(V\). Assume \(v \sim_{\mathit{self}} w\). Then \(v \sim_{G} w\).
\item[\paperlineanchor{premise:53696d706c6547726170682e53756267726170682e636f6553756267726170685f61646a}\paperlink{premise:53696d706c6547726170682e53756267726170682e636f6553756267726170685f61646a}{P105.}] \refstepcounter{paperpremise}\label{premise:53696d706c6547726170682e53756267726170682e636f6553756267726170685f61646a}\textbf{Coercion subgraph adj.} Let \(V\) be a type. Let \(G\) be a simple graph on \(V\). Let \(G^{'}\) be a subgraph of \(G\). Let \(G^{''}\) be a subgraph of \(\operatorname{graph}(G^{'})\). Let \(v\) and \(w\) be elements of \(V\). Then \(v \sim_{\operatorname{coe\,Subgraph}(G^{''})} w \iff (\exists \mathit{hv}\text{ proving }(v \in V(G^{'})),\; \exists \mathit{hw}\text{ proving }(w \in V(G^{'})),\; (v\, :\, V(G^{'})) \sim_{G^{''}} (w\, :\, V(G^{'})))\).
\item[\paperlineanchor{premise:53696d706c6547726170682e53756267726170682e636f655f646567726565}\paperlink{premise:53696d706c6547726170682e53756267726170682e636f655f646567726565}{P106.}] \refstepcounter{paperpremise}\label{premise:53696d706c6547726170682e53756267726170682e636f655f646567726565}\textbf{Coercion degree.} Let \(V\) be a type. Let \(G\) be a simple graph on \(V\). Let \(G^{'}\) be a subgraph of \(G\). Let \(v\) be an element of \(V(G^{'})\). Fix a finite enumeration of \(N_{\operatorname{graph}(G^{'})}(v)\). Fix a finite enumeration of \(\operatorname{Simple\,Graph\mathord{:}Subgraph\mathord{:}neighbor\,Set}\left(G^{'}, v\right)\). Then \(\deg_{\operatorname{graph}(G^{'})}(v) = \deg_{G^{'}}(v)\).
\item[\paperlineanchor{premise:53696d706c6547726170682e53756267726170682e6465677265655f65715f6f6e655f6966665f756e697175655f61646a}\paperlink{premise:53696d706c6547726170682e53756267726170682e6465677265655f65715f6f6e655f6966665f756e697175655f61646a}{P107.}] \refstepcounter{paperpremise}\label{premise:53696d706c6547726170682e53756267726170682e6465677265655f65715f6f6e655f6966665f756e697175655f61646a}\textbf{Degree equals one if and only if unique adj.} Let \(V\) be a type. Let \(G\) be a simple graph on \(V\). Let \(G^{'}\) be a subgraph of \(G\). Let \(v\) be an element of \(V\). Fix a finite enumeration of \(\operatorname{Simple\,Graph\mathord{:}Subgraph\mathord{:}neighbor\,Set}\left(G^{'}, v\right)\). Then \(\deg_{G^{'}}(v) = 1 \iff \left(\exists! w\, :\, V,\; v \sim_{G^{'}} w\right)\).
\item[\paperlineanchor{premise:53696d706c6547726170682e53756267726170682e656467655f76657274}\paperlink{premise:53696d706c6547726170682e53756267726170682e656467655f76657274}{P108.}] \refstepcounter{paperpremise}\label{premise:53696d706c6547726170682e53756267726170682e656467655f76657274}\textbf{Edge vert.} Let \(V\) be a type. Let \(G\) be a simple graph on \(V\). Let \(\mathit{self}\) be a subgraph of \(G\). Let \(v\) and \(w\) be elements of \(V\). Assume \(v \sim_{\mathit{self}} w\). Then \(v \in V(\mathit{self})\).
\item[\paperlineanchor{premise:53696d706c6547726170682e53756267726170682e695375705f61646a}\paperlink{premise:53696d706c6547726170682e53756267726170682e695375705f61646a}{P109.}] \refstepcounter{paperpremise}\label{premise:53696d706c6547726170682e53756267726170682e695375705f61646a}\textbf{I sup adj.} Let \(\iota\) be a sort. Let \(V\) be a type. Let \(G\) be a simple graph on \(V\). Let \(a\) and \(b\) be elements of \(V\). Let \(f\) be a function of type \(\iota \to \operatorname{Subgraph}\left(G\right)\). Then \(a \sim_{\operatorname{i\,Sup}(i \mapsto f(i))} b \iff (\exists i\, :\, \iota,\; a \sim_{f(i)} b)\).
\item[\paperlineanchor{premise:53696d706c6547726170682e53756267726170682e6d61705f636f6d70}\paperlink{premise:53696d706c6547726170682e53756267726170682e6d61705f636f6d70}{P110.}] \refstepcounter{paperpremise}\label{premise:53696d706c6547726170682e53756267726170682e6d61705f636f6d70}\textbf{Map comp.} Let \(V\) and \(W\) be types. Let \(G\) be a simple graph on \(V\). Let \(G^{'}\) be a simple graph on \(W\). Let \(U\) be a type. Let \(G^{''}\) be a simple graph on \(U\). Let \(H\) be a subgraph of \(G\). Let \(f\) be a graph homomorphism from \(G\) to \(G^{'}\). Let \(g\) be a graph homomorphism from \(G^{'}\) to \(G^{''}\). Then \(\operatorname{Simple\,Graph\mathord{:}Subgraph\mathord{:}map}(\operatorname{comp}(g,\allowbreak{} f),\allowbreak{} H) = \operatorname{Simple\,Graph\mathord{:}Subgraph\mathord{:}map}(g,\allowbreak{} \operatorname{Simple\,Graph\mathord{:}Subgraph\mathord{:}map}(f,\allowbreak{} H))\).
\item[\paperlineanchor{premise:53696d706c6547726170682e53756267726170682e6d61705f6964}\paperlink{premise:53696d706c6547726170682e53756267726170682e6d61705f6964}{P111.}] \refstepcounter{paperpremise}\label{premise:53696d706c6547726170682e53756267726170682e6d61705f6964}\textbf{Map id.} Let \(V\) be a type. Let \(G\) be a simple graph on \(V\). Let \(H\) be a subgraph of \(G\). Then \(\operatorname{Simple\,Graph\mathord{:}Subgraph\mathord{:}map}(\operatorname{Simple\,Graph\mathord{:}Hom\mathord{:}id},\allowbreak{} H) = H\).
\item[\paperlineanchor{premise:53696d706c6547726170682e53756267726170682e6d656d5f737570706f7274}\paperlink{premise:53696d706c6547726170682e53756267726170682e6d656d5f737570706f7274}{P112.}] \refstepcounter{paperpremise}\label{premise:53696d706c6547726170682e53756267726170682e6d656d5f737570706f7274}\textbf{Membership support.} Let \(V\) be a type. Let \(G\) be a simple graph on \(V\). Let \(H\) be a subgraph of \(G\). Let \(v\) be an element of \(V\). Then \(v \in \operatorname{Simple\,Graph\mathord{:}Subgraph\mathord{:}support}(H) \iff (\exists w\, :\, V,\; v \sim_{H} w)\).
\item[\paperlineanchor{premise:53696d706c6547726170682e53756267726170682e7370616e6e696e67436f655f73756267726170684f6641646a}\paperlink{premise:53696d706c6547726170682e53756267726170682e7370616e6e696e67436f655f73756267726170684f6641646a}{P113.}] \refstepcounter{paperpremise}\label{premise:53696d706c6547726170682e53756267726170682e7370616e6e696e67436f655f73756267726170684f6641646a}\textbf{Spanning coercion subgraph of adj.} Let \(V\) be a type. Let \(G\) be a simple graph on \(V\). Let \(v\) and \(w\) be elements of \(V\). Let \(\mathit{hadj}\) be a proof of \(v \sim_{G} w\). Then \(\operatorname{spanning\,graph}(\operatorname{subgraph\,Of\,Adj}(G)) = \operatorname{from\,Edge\,Set}(\{\operatorname{Sym2\mathord{:}mk}((v,\allowbreak{} w))\})\).
\item[\paperlineanchor{premise:53696d706c6547726170682e53756267726170682e7375705f61646a}\paperlink{premise:53696d706c6547726170682e53756267726170682e7375705f61646a}{P114.}] \refstepcounter{paperpremise}\label{premise:53696d706c6547726170682e53756267726170682e7375705f61646a}\textbf{Sup adj.} Let \(V\) be a type. Let \(G\) be a simple graph on \(V\). Let \(G_{1}\) be a subgraph of \(G\). Let \(G_{2}\) be a subgraph of \(G\). Let \(a\) and \(b\) be elements of \(V\). Then \(a \sim_{\operatorname{max}\left(G_{1}, G_{2}\right)} b \iff a \sim_{G_{1}} b \lor a \sim_{G_{2}} b\).
\item[\paperlineanchor{premise:53696d706c6547726170682e53756267726170682e737570706f72745f7375627365745f7665727473}\paperlink{premise:53696d706c6547726170682e53756267726170682e737570706f72745f7375627365745f7665727473}{P115.}] \refstepcounter{paperpremise}\label{premise:53696d706c6547726170682e53756267726170682e737570706f72745f7375627365745f7665727473}\textbf{Support subset verts.} Let \(V\) be a type. Let \(G\) be a simple graph on \(V\). Let \(H\) be a subgraph of \(G\). Then \(\operatorname{Simple\,Graph\mathord{:}Subgraph\mathord{:}support}\left(H\right) \subseteq V(H)\).
\item[\paperlineanchor{premise:53696d706c6547726170682e53756267726170682e73796d6d}\paperlink{premise:53696d706c6547726170682e53756267726170682e73796d6d}{P116.}] \refstepcounter{paperpremise}\label{premise:53696d706c6547726170682e53756267726170682e73796d6d}\textbf{Symm.} Let \(V\) be a type. Let \(G\) be a simple graph on \(V\). Let \(\mathit{self}\) be a subgraph of \(G\). Then \(\operatorname{Symmetric}(\operatorname{Simple\,Graph\mathord{:}Subgraph\mathord{:}Adj}(\mathit{self}))\).
\item[\paperlineanchor{premise:53696d706c6547726170682e53756267726170682e746f705f61646a}\paperlink{premise:53696d706c6547726170682e53756267726170682e746f705f61646a}{P117.}] \refstepcounter{paperpremise}\label{premise:53696d706c6547726170682e53756267726170682e746f705f61646a}\textbf{Top adj.} Let \(V\) be a type. Let \(G\) be a simple graph on \(V\). Let \(a\) and \(b\) be elements of \(V\). Then \(a \sim_{\operatorname{full\,subgraph}(G)} b \iff a \sim_{G} b\).
\item[\paperlineanchor{premise:53696d706c6547726170682e53756267726170682e76657274735f636f655375626772617068}\paperlink{premise:53696d706c6547726170682e53756267726170682e76657274735f636f655375626772617068}{P118.}] \refstepcounter{paperpremise}\label{premise:53696d706c6547726170682e53756267726170682e76657274735f636f655375626772617068}\textbf{Verts coercion subgraph.} Let \(V\) be a type. Let \(G\) be a simple graph on \(V\). Let \(G^{'}\) be a subgraph of \(G\). Let \(G^{''}\) be a subgraph of \(\operatorname{graph}(G^{'})\). Then \(V(\operatorname{coe\,Subgraph}\left(G^{''}\right)) = \operatorname{val}\left[V(G^{''})\right]\).
\item[\paperlineanchor{premise:53696d706c6547726170682e53756267726170682e76657274735f69537570}\paperlink{premise:53696d706c6547726170682e53756267726170682e76657274735f69537570}{P119.}] \refstepcounter{paperpremise}\label{premise:53696d706c6547726170682e53756267726170682e76657274735f69537570}\textbf{Verts i sup.} Let \(\iota\) be a sort. Let \(V\) be a type. Let \(G\) be a simple graph on \(V\). Let \(f\) be a function of type \(\iota \to \operatorname{Subgraph}\left(G\right)\). Then \(V(\operatorname{i\,Sup}\left(i \mapsto f\left(i\right)\right)) = \bigcup_{i\, :\, \iota} V(f\left(i\right))\).
\item[\paperlineanchor{premise:53696d706c6547726170682e57616c6b2e49734379636c652e6e636172645f6e65696768626f725365745f746f53756267726170685f65715f74776f}\paperlink{premise:53696d706c6547726170682e57616c6b2e49734379636c652e6e636172645f6e65696768626f725365745f746f53756267726170685f65715f74776f}{P120.}] \refstepcounter{paperpremise}\label{premise:53696d706c6547726170682e57616c6b2e49734379636c652e6e636172645f6e65696768626f725365745f746f53756267726170685f65715f74776f}\textbf{Ncard neighbor set to subgraph equals two.} Let \(V\) be a type. Let \(G\) be a simple graph on \(V\). Let \(u\) and \(v\) be elements of \(V\). Let \(p\) be a closed walk in \(G\) based at \(u\). Assume \(p\text{ is a cycle}\). Assume \(p\text{ visits }v\). Then \(\operatorname{ncard}(\operatorname{Simple\,Graph\mathord{:}Subgraph\mathord{:}neighbor\,Set}(\operatorname{Simple\,Graph\mathord{:}Walk\mathord{:}to\,Subgraph}(p),\allowbreak{} v)) = 2\).
\item[\paperlineanchor{premise:53696d706c6547726170682e57616c6b2e49734379636c652e726f74617465}\paperlink{premise:53696d706c6547726170682e57616c6b2e49734379636c652e726f74617465}{P121.}] \refstepcounter{paperpremise}\label{premise:53696d706c6547726170682e57616c6b2e49734379636c652e726f74617465}\textbf{Rotate.} Let \(V\) be a type. Let \(G\) be a simple graph on \(V\). Assume equality on \(V\) is decidable. Let \(u\) and \(v\) be elements of \(V\). Let \(c\) be a closed walk in \(G\) based at \(v\). Assume \(c\text{ is a cycle}\). Let \(h\) be a proof of \(c\text{ visits }u\). Then \(\operatorname{Simple\,Graph\mathord{:}Walk\mathord{:}rotate}\left(c\right)\text{ is a cycle}\).
\item[\paperlineanchor{premise:53696d706c6547726170682e57616c6b2e4973506174682e6e636172645f6e65696768626f725365745f746f53756267726170685f696e7465726e616c5f65715f74776f}\paperlink{premise:53696d706c6547726170682e57616c6b2e4973506174682e6e636172645f6e65696768626f725365745f746f53756267726170685f696e7465726e616c5f65715f74776f}{P122.}] \refstepcounter{paperpremise}\label{premise:53696d706c6547726170682e57616c6b2e4973506174682e6e636172645f6e65696768626f725365745f746f53756267726170685f696e7465726e616c5f65715f74776f}\textbf{Ncard neighbor set to subgraph internal equals two.} Let \(V\) be a type. Let \(G\) be a simple graph on \(V\). Let \(v\) and \(u\) be elements of \(V\). Let \(i\) be a natural number. Let \(p\) be a walk in \(G\) from \(u\) to \(v\). Assume \(p\text{ has no repeated vertices}\). Assume \(i \ne 0\). Assume \(i < \operatorname{length}(p)\). Then \(\operatorname{ncard}(\operatorname{Simple\,Graph\mathord{:}Subgraph\mathord{:}neighbor\,Set}(\operatorname{Simple\,Graph\mathord{:}Walk\mathord{:}to\,Subgraph}(p),\allowbreak{} \operatorname{get\,Vert}(p,\allowbreak{} i))) = 2\).
\item[\paperlineanchor{premise:53696d706c6547726170682e57616c6b2e4973506174682e6f665f61646a}\paperlink{premise:53696d706c6547726170682e57616c6b2e4973506174682e6f665f61646a}{P123.}] \refstepcounter{paperpremise}\label{premise:53696d706c6547726170682e57616c6b2e4973506174682e6f665f61646a}\textbf{Of adj.} Let \(V\) be a type. Let \(G\) be a simple graph on \(V\). Let \(u\) and \(v\) be elements of \(V\). Let \(h\) be a proof of \(u \sim_{G} v\). Then \(\operatorname{to\,Walk}\text{ has no repeated vertices}\).
\item[\paperlineanchor{premise:53696d706c6547726170682e57616c6b2e4973506174682e736e645f6f665f746f53756267726170685f61646a}\paperlink{premise:53696d706c6547726170682e57616c6b2e4973506174682e736e645f6f665f746f53756267726170685f61646a}{P124.}] \refstepcounter{paperpremise}\label{premise:53696d706c6547726170682e57616c6b2e4973506174682e736e645f6f665f746f53756267726170685f61646a}\textbf{Snd of to subgraph adj.} Let \(V\) be a type. Let \(G\) be a simple graph on \(V\). Let \(u\), \(v\), and \(v^{'}\) be elements of \(V\). Let \(p\) be a walk in \(G\) from \(u\) to \(v\). Assume \(p\text{ has no repeated vertices}\). Assume \(u \sim_{\operatorname{Simple\,Graph\mathord{:}Walk\mathord{:}to\,Subgraph}\left(p\right)} v^{'}\). Then \(\operatorname{Simple\,Graph\mathord{:}Walk\mathord{:}snd}\left(p\right) = v^{'}\).
\item[\paperlineanchor{premise:53696d706c6547726170682e57616c6b2e61646a5f746f53756267726170685f6d61704c65}\paperlink{premise:53696d706c6547726170682e57616c6b2e61646a5f746f53756267726170685f6d61704c65}{P125.}] \refstepcounter{paperpremise}\label{premise:53696d706c6547726170682e57616c6b2e61646a5f746f53756267726170685f6d61704c65}\textbf{Adj to subgraph map at most.} Let \(V\) be a type. Let \(G\) be a simple graph on \(V\). Let \(u\) and \(v\) be elements of \(V\). Let \(G^{'}\) be a simple graph on \(V\). Let \(w\) and \(x\) be elements of \(V\). Let \(p\) be a walk in \(G\) from \(u\) to \(v\). Let \(h\) be a proof of \(G \subseteq G^{'}\). Then \(w \sim_{\operatorname{Simple\,Graph\mathord{:}Walk\mathord{:}to\,Subgraph}(\operatorname{map\,Le}(p))} x \iff w \sim_{\operatorname{Simple\,Graph\mathord{:}Walk\mathord{:}to\,Subgraph}(p)} x\).
\item[\paperlineanchor{premise:53696d706c6547726170682e57616c6b2e636f6e735f69734379636c655f696666}\paperlink{premise:53696d706c6547726170682e57616c6b2e636f6e735f69734379636c655f696666}{P126.}] \refstepcounter{paperpremise}\label{premise:53696d706c6547726170682e57616c6b2e636f6e735f69734379636c655f696666}\textbf{Cons is cycle if and only if.} Let \(V\) be a type. Let \(G\) be a simple graph on \(V\). Let \(u\) and \(v\) be elements of \(V\). Let \(p\) be a walk in \(G\) from \(v\) to \(u\). Let \(h\) be a proof of \(u \sim_{G} v\). Then \(\operatorname{Simple\,Graph\mathord{:}Walk\mathord{:}cons}(p)\text{ is a cycle} \iff p\text{ has no repeated vertices} \land \operatorname{Sym2\mathord{:}mk}((u,\allowbreak{} v)) \notin \operatorname{edges}(p)\).
\item[\paperlineanchor{premise:53696d706c6547726170682e57616c6b2e636f6e735f6973506174685f696666}\paperlink{premise:53696d706c6547726170682e57616c6b2e636f6e735f6973506174685f696666}{P127.}] \refstepcounter{paperpremise}\label{premise:53696d706c6547726170682e57616c6b2e636f6e735f6973506174685f696666}\textbf{Cons is path if and only if.} Let \(V\) be a type. Let \(G\) be a simple graph on \(V\). Let \(u\), \(v\), and \(w\) be elements of \(V\). Let \(h\) be a proof of \(u \sim_{G} v\). Let \(p\) be a walk in \(G\) from \(v\) to \(w\). Then \(\operatorname{Simple\,Graph\mathord{:}Walk\mathord{:}cons}(p)\text{ has no repeated vertices} \iff p\text{ has no repeated vertices} \land u \notin \operatorname{vertex\,sequence}(p)\).
\item[\paperlineanchor{premise:53696d706c6547726170682e57616c6b2e65646765735f6d6170}\paperlink{premise:53696d706c6547726170682e57616c6b2e65646765735f6d6170}{P128.}] \refstepcounter{paperpremise}\label{premise:53696d706c6547726170682e57616c6b2e65646765735f6d6170}\textbf{Edges map.} Let \(V\) and \(V^{'}\) be types. Let \(G\) be a simple graph on \(V\). Let \(G^{'}\) be a simple graph on \(V^{'}\). Let \(f\) be a graph homomorphism from \(G\) to \(G^{'}\). Let \(u\) and \(v\) be elements of \(V\). Let \(p\) be a walk in \(G\) from \(u\) to \(v\). Then \(\operatorname{edges}(\operatorname{image}_{f}(p)) = \langle \operatorname{Sym2\mathord{:}map}(\operatorname{DFun\,Like\mathord{:}coe}(f))(x) \mid x\in \operatorname{edges}(p)\rangle\).
\item[\paperlineanchor{premise:53696d706c6547726170682e57616c6b2e6673745f6d656d5f737570706f72745f6f665f6d656d5f6564676573}\paperlink{premise:53696d706c6547726170682e57616c6b2e6673745f6d656d5f737570706f72745f6f665f6d656d5f6564676573}{P129.}] \refstepcounter{paperpremise}\label{premise:53696d706c6547726170682e57616c6b2e6673745f6d656d5f737570706f72745f6f665f6d656d5f6564676573}\textbf{Fst membership support of membership edges.} Let \(V\) be a type. Let \(G\) be a simple graph on \(V\). Let \(t\), \(u\), \(v\), and \(w\) be elements of \(V\). Let \(p\) be a walk in \(G\) from \(v\) to \(w\). Assume \(\operatorname{Sym2\mathord{:}mk}\left(\left(t, u\right)\right) \in \operatorname{edges}\left(p\right)\). Then \(p\text{ visits }t\).
\item[\paperlineanchor{premise:53696d706c6547726170682e57616c6b2e676574566572745f6c656e677468}\paperlink{premise:53696d706c6547726170682e57616c6b2e676574566572745f6c656e677468}{P130.}] \refstepcounter{paperpremise}\label{premise:53696d706c6547726170682e57616c6b2e676574566572745f6c656e677468}\textbf{Get vert length.} Let \(V\) be a type. Let \(G\) be a simple graph on \(V\). Let \(u\) and \(v\) be elements of \(V\). Let \(w\) be a walk in \(G\) from \(u\) to \(v\). Then \(\operatorname{get\,Vert}\left(w, \operatorname{length}(w)\right) = v\).
\item[\paperlineanchor{premise:53696d706c6547726170682e57616c6b2e676574566572745f7a65726f}\paperlink{premise:53696d706c6547726170682e57616c6b2e676574566572745f7a65726f}{P131.}] \refstepcounter{paperpremise}\label{premise:53696d706c6547726170682e57616c6b2e676574566572745f7a65726f}\textbf{Get vert zero.} Let \(V\) be a type. Let \(G\) be a simple graph on \(V\). Let \(u\) and \(v\) be elements of \(V\). Let \(w\) be a walk in \(G\) from \(u\) to \(v\). Then \(\operatorname{get\,Vert}\left(w, 0\right) = u\).
\item[\paperlineanchor{premise:53696d706c6547726170682e57616c6b2e6973506174685f6966665f65715f6e696c}\paperlink{premise:53696d706c6547726170682e57616c6b2e6973506174685f6966665f65715f6e696c}{P132.}] \refstepcounter{paperpremise}\label{premise:53696d706c6547726170682e57616c6b2e6973506174685f6966665f65715f6e696c}\textbf{Is path if and only if equals nil.} Let \(V\) be a type. Let \(G\) be a simple graph on \(V\). Let \(u\) be an element of \(V\). Let \(p\) be a closed walk in \(G\) based at \(u\). Then \(p\text{ has no repeated vertices} \iff p = \operatorname{Simple\,Graph\mathord{:}Walk\mathord{:}nil}\).
\item[\paperlineanchor{premise:53696d706c6547726170682e57616c6b2e6d61704c655f697350617468}\paperlink{premise:53696d706c6547726170682e57616c6b2e6d61704c655f697350617468}{P133.}] \refstepcounter{paperpremise}\label{premise:53696d706c6547726170682e57616c6b2e6d61704c655f697350617468}\textbf{Map at most is path.} Let \(V\) be a type. Let \(G\) and \(G^{'}\) be simple graphs on \(V\). Let \(h\) be a proof of \(G \subseteq G^{'}\). Let \(u\) and \(v\) be elements of \(V\). Let \(p\) be a walk in \(G\) from \(u\) to \(v\). Then \(\operatorname{map\,Le}\left(p\right)\text{ has no repeated vertices} \iff p\text{ has no repeated vertices}\).
\item[\paperlineanchor{premise:53696d706c6547726170682e57616c6b2e6d656d5f737570706f72745f6966665f6578697374735f67657456657274}\paperlink{premise:53696d706c6547726170682e57616c6b2e6d656d5f737570706f72745f6966665f6578697374735f67657456657274}{P134.}] \refstepcounter{paperpremise}\label{premise:53696d706c6547726170682e57616c6b2e6d656d5f737570706f72745f6966665f6578697374735f67657456657274}\textbf{Membership support if and only if exists get vert.} Let \(V\) be a type. Let \(G\) be a simple graph on \(V\). Let \(u\), \(v\), and \(w\) be elements of \(V\). Let \(p\) be a walk in \(G\) from \(v\) to \(w\). Then \(p\text{ visits }u \iff (\exists n\, :\, \mathbb{N},\; \operatorname{get\,Vert}(p,\allowbreak{} n) = u \land n \le \operatorname{length}(p))\).
\item[\paperlineanchor{premise:53696d706c6547726170682e57616c6b2e6d656d5f76657274735f746f5375626772617068}\paperlink{premise:53696d706c6547726170682e57616c6b2e6d656d5f76657274735f746f5375626772617068}{P135.}] \refstepcounter{paperpremise}\label{premise:53696d706c6547726170682e57616c6b2e6d656d5f76657274735f746f5375626772617068}\textbf{Membership verts to subgraph.} Let \(V\) be a type. Let \(G\) be a simple graph on \(V\). Let \(u\), \(v\), and \(w\) be elements of \(V\). Let \(p\) be a walk in \(G\) from \(u\) to \(v\). Then \(w \in V(\operatorname{Simple\,Graph\mathord{:}Walk\mathord{:}to\,Subgraph}(p)) \iff p\text{ visits }w\).
\item[\paperlineanchor{premise:53696d706c6547726170682e57616c6b2e6e6f745f6e696c5f6f665f6e65}\paperlink{premise:53696d706c6547726170682e57616c6b2e6e6f745f6e696c5f6f665f6e65}{P136.}] \refstepcounter{paperpremise}\label{premise:53696d706c6547726170682e57616c6b2e6e6f745f6e696c5f6f665f6e65}\textbf{Not nil of differs from.} Let \(V\) be a type. Let \(G\) be a simple graph on \(V\). Let \(v\) and \(w\) be elements of \(V\). Let \(p\) be a walk in \(G\) from \(v\) to \(w\). Assume \(v \ne w\). Then \(\neg \operatorname{Nil}\left(p\right)\).
\item[\paperlineanchor{premise:53696d706c6547726170682e57616c6b2e746f53756267726170685f61646a5f736e64}\paperlink{premise:53696d706c6547726170682e57616c6b2e746f53756267726170685f61646a5f736e64}{P137.}] \refstepcounter{paperpremise}\label{premise:53696d706c6547726170682e57616c6b2e746f53756267726170685f61646a5f736e64}\textbf{To subgraph adj snd.} Let \(V\) be a type. Let \(G\) be a simple graph on \(V\). Let \(u\) and \(v\) be elements of \(V\). Let \(w\) be a walk in \(G\) from \(u\) to \(v\). Assume \(\neg \operatorname{Nil}\left(w\right)\). Then \(u \sim_{\operatorname{Simple\,Graph\mathord{:}Walk\mathord{:}to\,Subgraph}(w)} \operatorname{Simple\,Graph\mathord{:}Walk\mathord{:}snd}(w)\).
\item[\paperlineanchor{premise:53696d706c6547726170682e57616c6b2e746f53756267726170685f636f6e6e6563746564}\paperlink{premise:53696d706c6547726170682e57616c6b2e746f53756267726170685f636f6e6e6563746564}{P138.}] \refstepcounter{paperpremise}\label{premise:53696d706c6547726170682e57616c6b2e746f53756267726170685f636f6e6e6563746564}\textbf{To subgraph connected.} Let \(V\) be a type. Let \(G\) be a simple graph on \(V\). Let \(u\) and \(v\) be elements of \(V\). Let \(p\) be a walk in \(G\) from \(u\) to \(v\). Then \(\operatorname{Connected}(\operatorname{Simple\,Graph\mathord{:}Walk\mathord{:}to\,Subgraph}(p))\).
\item[\paperlineanchor{premise:53696d706c6547726170682e57616c6b2e746f53756267726170685f726f74617465}\paperlink{premise:53696d706c6547726170682e57616c6b2e746f53756267726170685f726f74617465}{P139.}] \refstepcounter{paperpremise}\label{premise:53696d706c6547726170682e57616c6b2e746f53756267726170685f726f74617465}\textbf{To subgraph rotate.} Let \(V\) be a type. Let \(G\) be a simple graph on \(V\). Let \(u\) and \(v\) be elements of \(V\). Assume equality on \(V\) is decidable. Let \(c\) be a closed walk in \(G\) based at \(v\). Let \(h\) be a proof of \(c\text{ visits }u\). Then \(\operatorname{Simple\,Graph\mathord{:}Walk\mathord{:}to\,Subgraph}(\operatorname{Simple\,Graph\mathord{:}Walk\mathord{:}rotate}(c)) = \operatorname{Simple\,Graph\mathord{:}Walk\mathord{:}to\,Subgraph}(c)\).
\item[\paperlineanchor{premise:53696d706c6547726170682e57616c6b2e76657274735f746f5375626772617068}\paperlink{premise:53696d706c6547726170682e57616c6b2e76657274735f746f5375626772617068}{P140.}] \refstepcounter{paperpremise}\label{premise:53696d706c6547726170682e57616c6b2e76657274735f746f5375626772617068}\textbf{Verts to subgraph.} Let \(V\) be a type. Let \(G\) be a simple graph on \(V\). Let \(u\) and \(v\) be elements of \(V\). Let \(p\) be a walk in \(G\) from \(u\) to \(v\). Then \(V(\operatorname{Simple\,Graph\mathord{:}Walk\mathord{:}to\,Subgraph}(p)) = \{w\in V \mid p\text{ visits }w\}\).
\item[\paperlineanchor{premise:53696d706c6547726170682e61646a5f616e645f726561636861626c655f64656c6574655f65646765735f6966665f6578697374735f6379636c65}\paperlink{premise:53696d706c6547726170682e61646a5f616e645f726561636861626c655f64656c6574655f65646765735f6966665f6578697374735f6379636c65}{P141.}] \refstepcounter{paperpremise}\label{premise:53696d706c6547726170682e61646a5f616e645f726561636861626c655f64656c6574655f65646765735f6966665f6578697374735f6379636c65}\textbf{Adj and reachable delete edges if and only if exists cycle.} Let \(V\) be a type. Let \(G\) be a simple graph on \(V\). Let \(v\) and \(w\) be elements of \(V\). Then \(v \sim_{G} w \land v\text{ can reach }w\text{ in }G \setminus \operatorname{from\,Edge\,Set}(\{\operatorname{Sym2\mathord{:}mk}((v,\allowbreak{} w))\}) \iff (\exists u\, :\, V,\; \exists p\, :\, \operatorname{Walk}_{G}(u,\allowbreak{} u),\; p\text{ is a cycle} \land \operatorname{Sym2\mathord{:}mk}((v,\allowbreak{} w)) \in \operatorname{edges}(p))\).
\item[\paperlineanchor{premise:53696d706c6547726170682e61646a5f636f6e67725f6f665f73796d32}\paperlink{premise:53696d706c6547726170682e61646a5f636f6e67725f6f665f73796d32}{P142.}] \refstepcounter{paperpremise}\label{premise:53696d706c6547726170682e61646a5f636f6e67725f6f665f73796d32}\textbf{Adj congruence of sym2.} Let \(V\) be a type. Let \(G\) be a simple graph on \(V\). Let \(u\), \(v\), \(w\), and \(x\) be elements of \(V\). Assume \(\operatorname{Sym2\mathord{:}mk}((u,\allowbreak{} v)) = \operatorname{Sym2\mathord{:}mk}((w,\allowbreak{} x))\). Then \(u \sim_{G} v \iff w \sim_{G} x\).
\item[\paperlineanchor{premise:53696d706c6547726170682e656467655f61646a}\paperlink{premise:53696d706c6547726170682e656467655f61646a}{P143.}] \refstepcounter{paperpremise}\label{premise:53696d706c6547726170682e656467655f61646a}\textbf{Edge adj.} Let \(V\) be a type. Let \(s\), \(t\), \(v\), and \(w\) be elements of \(V\). Then \(v \sim_{\operatorname{edge}(s,\allowbreak{} t)} w \iff (v = s \land w = t \lor v = t \land w = s) \land v \ne w\).
\item[\paperlineanchor{premise:53696d706c6547726170682e657874}\paperlink{premise:53696d706c6547726170682e657874}{P144.}] \refstepcounter{paperpremise}\label{premise:53696d706c6547726170682e657874}\textbf{Ext.} Let \(V\) be a type. Let \(x\) and \(y\) be simple graphs on \(V\). Assume \(\operatorname{Adj}_{x} = \operatorname{Adj}_{y}\). Then \(x = y\).
\item[\paperlineanchor{premise:53696d706c6547726170682e696e737446696e697465}\paperlink{premise:53696d706c6547726170682e696e737446696e697465}{P145.}] \refstepcounter{paperpremise}\label{premise:53696d706c6547726170682e696e737446696e697465}\textbf{Inst finite.} Let \(V\) be a type. Assume \(V\) is finite. Then \(\operatorname{Finite}\left(\operatorname{Simple\,Graph}(V)\right)\).
\item[\paperlineanchor{premise:53696d706c6547726170682e69727265666c}\paperlink{premise:53696d706c6547726170682e69727265666c}{P146.}] \refstepcounter{paperpremise}\label{premise:53696d706c6547726170682e69727265666c}\textbf{Irrefl.} Let \(V\) be a type. Let \(G\) be a simple graph on \(V\). Let \(v\) be an element of \(V\). Then \(\neg \left(v \sim_{G} v\right)\).
\item[\paperlineanchor{premise:53696d706c6547726170682e6d656d5f6e65696768626f72536574}\paperlink{premise:53696d706c6547726170682e6d656d5f6e65696768626f72536574}{P147.}] \refstepcounter{paperpremise}\label{premise:53696d706c6547726170682e6d656d5f6e65696768626f72536574}\textbf{Membership neighbor set.} Let \(V\) be a type. Let \(G\) be a simple graph on \(V\). Let \(v\) and \(w\) be elements of \(V\). Then \(w \in N_{G}(v) \iff v \sim_{G} w\).
\item[\paperlineanchor{premise:53696d706c6547726170682e73646966665f61646a}\paperlink{premise:53696d706c6547726170682e73646966665f61646a}{P148.}] \refstepcounter{paperpremise}\label{premise:53696d706c6547726170682e73646966665f61646a}\textbf{Sdiff adj.} Let \(V\) be a type. Let \(x\) and \(y\) be simple graphs on \(V\). Let \(v\) and \(w\) be elements of \(V\). Then \(v \sim_{x \setminus y} w \iff v \sim_{x} w \land \neg \left(v \sim_{y} w\right)\).
\item[\paperlineanchor{premise:53696d706c6547726170682e73696e676c65746f6e53756267726170685f6c655f696666}\paperlink{premise:53696d706c6547726170682e73696e676c65746f6e53756267726170685f6c655f696666}{P149.}] \refstepcounter{paperpremise}\label{premise:53696d706c6547726170682e73696e676c65746f6e53756267726170685f6c655f696666}\textbf{Singleton subgraph at most if and only if.} Let \(V\) be a type. Let \(G\) be a simple graph on \(V\). Let \(v\) be an element of \(V\). Let \(H\) be a subgraph of \(G\). Then \(\operatorname{singleton\,Subgraph}\left(G, v\right) \le H \iff v \in V(H)\).
\item[\paperlineanchor{premise:53696d706c6547726170682e73756267726170684f6641646a5f7665727473}\paperlink{premise:53696d706c6547726170682e73756267726170684f6641646a5f7665727473}{P150.}] \refstepcounter{paperpremise}\label{premise:53696d706c6547726170682e73756267726170684f6641646a5f7665727473}\textbf{Subgraph of adj verts.} Let \(V\) be a type. Let \(G\) be a simple graph on \(V\). Let \(v\) and \(w\) be elements of \(V\). Let \(\mathit{hvw}\) be a proof of \(v \sim_{G} w\). Then \(V(\operatorname{subgraph\,Of\,Adj}\left(G\right)) = \{v\} \cup \{w\}\).
\item[\paperlineanchor{premise:53696d706c6547726170682e73756d5f646567726565735f65715f74776963655f636172645f6564676573}\paperlink{premise:53696d706c6547726170682e73756d5f646567726565735f65715f74776963655f636172645f6564676573}{P151.}] \refstepcounter{paperpremise}\label{premise:53696d706c6547726170682e73756d5f646567726565735f65715f74776963655f636172645f6564676573}\textbf{Sum degrees equals twice cardinality edges.} Let \(V\) be a type. Let \(G\) be a simple graph on \(V\). Fix a finite enumeration of \(V\). Assume adjacency in \(G\) is decidable. Then \(\sum_{v\in V}\left(\deg_{G}(v)\right) = 2 \cdot \left|E(G)\right|\).
\item[\paperlineanchor{premise:53696d706c6547726170682e7375705f61646a}\paperlink{premise:53696d706c6547726170682e7375705f61646a}{P152.}] \refstepcounter{paperpremise}\label{premise:53696d706c6547726170682e7375705f61646a}\textbf{Sup adj.} Let \(V\) be a type. Let \(x\) and \(y\) be simple graphs on \(V\). Let \(v\) and \(w\) be elements of \(V\). Then \(v \sim_{x \cup y} w \iff v \sim_{x} w \lor v \sim_{y} w\).
\item[\paperlineanchor{premise:53696d706c6547726170682e7375705f656467655f6f665f61646a}\paperlink{premise:53696d706c6547726170682e7375705f656467655f6f665f61646a}{P153.}] \refstepcounter{paperpremise}\label{premise:53696d706c6547726170682e7375705f656467655f6f665f61646a}\textbf{Sup edge of adj.} Let \(V\) be a type. Let \(G\) be a simple graph on \(V\). Let \(s\) and \(t\) be elements of \(V\). Assume \(s \sim_{G} t\). Then \(G \cup \operatorname{edge}\left(s, t\right) = G\).
\item[\paperlineanchor{premise:53696d706c6547726170682e746f705f61646a}\paperlink{premise:53696d706c6547726170682e746f705f61646a}{P154.}] \refstepcounter{paperpremise}\label{premise:53696d706c6547726170682e746f705f61646a}\textbf{Top adj.} Let \(V\) be a type. Let \(v\) and \(w\) be elements of \(V\). Then \(v \sim_{\operatorname{complete\,graph}(V)} w \iff v \ne w\).
\item[\paperlineanchor{premise:537562747970652e636f655f657461}\paperlink{premise:537562747970652e636f655f657461}{P155.}] \refstepcounter{paperpremise}\label{premise:537562747970652e636f655f657461}\textbf{Coercion eta.} Let \(\alpha\) be a sort. Let \(p\) be a function of type \(\alpha \to \mathsf{Prop}\). Let \(a\) be an element of \(\{a_{1}\in \alpha \mid p\left(a_{1}\right)\}\). Let \(h\) be a proof of \(p\left(a\right)\). Then \(\left(a\, :\, \{a_{1}\in \alpha \mid p\left(a_{1}\right)\}\right) = a\).
\item[\paperlineanchor{premise:537562747970652e636f655f70726f70}\paperlink{premise:537562747970652e636f655f70726f70}{P156.}] \refstepcounter{paperpremise}\label{premise:537562747970652e636f655f70726f70}\textbf{Coercion prop.} Let \(\alpha\) be a type. Let \(S\) be a subset of \(\alpha\). Let \(a\) be an element of \(S\). Then \(a \in S\).
\item[\paperlineanchor{premise:537562747970652e66696e697465}\paperlink{premise:537562747970652e66696e697465}{P157.}] \refstepcounter{paperpremise}\label{premise:537562747970652e66696e697465}\textbf{Finite.} Let \(\alpha\) be a sort. Assume \(\alpha\) is finite. Let \(p\) be a function of type \(\alpha \to \mathsf{Prop}\). Then \(\operatorname{Finite}\left(\{x\in \alpha \mid p\left(x\right)\}\right)\).
\item[\paperlineanchor{premise:537562747970652e6d6b5f65715f6d6b}\paperlink{premise:537562747970652e6d6b5f65715f6d6b}{P158.}] \refstepcounter{paperpremise}\label{premise:537562747970652e6d6b5f65715f6d6b}\textbf{Mk equals mk.} Let \(\alpha\) be a sort. Let \(p\) be a function of type \(\alpha \to \mathsf{Prop}\). Let \(a\) be an element of \(\alpha\). Let \(h\) be a proof of \(p\left(a\right)\). Let \(a^{'}\) be an element of \(\alpha\). Let \(h^{'}\) be a proof of \(p\left(a^{'}\right)\). Then \((a\, :\, \{x\in \alpha \mid p(x)\}) = (a^{'}\, :\, \{x\in \alpha \mid p(x)\}) \iff a = a^{'}\).
\item[\paperlineanchor{premise:537562747970652e70726f70}\paperlink{premise:537562747970652e70726f70}{P159.}] \refstepcounter{paperpremise}\label{premise:537562747970652e70726f70}\textbf{Prop.} Let \(\alpha\) be a sort. Let \(p\) be a function of type \(\alpha \to \mathsf{Prop}\). Let \(x\) be an element of \(\{x_{1}\in \alpha \mid p(x_{1})\}\). Then \(p\left(x\right)\).
\item[\paperlineanchor{premise:537562747970652e70726f7065727479}\paperlink{premise:537562747970652e70726f7065727479}{P160.}] \refstepcounter{paperpremise}\label{premise:537562747970652e70726f7065727479}\textbf{Property.} Let \(\alpha\) be a sort. Let \(p\) be a function of type \(\alpha \to \mathsf{Prop}\). Let \(\mathit{self}\) be an element of \(\{x\in \alpha \mid p(x)\}\). Then \(p\left(\mathit{self}\right)\).
\item[\paperlineanchor{premise:537562747970652e76616c5f696e6a656374697665}\paperlink{premise:537562747970652e76616c5f696e6a656374697665}{P161.}] \refstepcounter{paperpremise}\label{premise:537562747970652e76616c5f696e6a656374697665}\textbf{Val injective.} Let \(\alpha\) be a sort. Let \(p\) be a function of type \(\alpha \to \mathsf{Prop}\). Then \(\operatorname{val}\text{ is injective}\).
\item[\paperlineanchor{premise:53796d322e6571}\paperlink{premise:53796d322e6571}{P162.}] \refstepcounter{paperpremise}\label{premise:53796d322e6571}\textbf{Equals.} Let \(\alpha\) be a type. Let \(p\) be an element of \(\alpha \times \alpha\). Let \(p^{'}\) be an element of \(\alpha \times \alpha\). Then \(\operatorname{Sym2\mathord{:}mk}(p) = \operatorname{Sym2\mathord{:}mk}(p^{'}) \iff \operatorname{Rel}(\alpha,\allowbreak{} p,\allowbreak{} p^{'})\).
\item[\paperlineanchor{premise:53796d322e65715f73776170}\paperlink{premise:53796d322e65715f73776170}{P163.}] \refstepcounter{paperpremise}\label{premise:53796d322e65715f73776170}\textbf{Equals swap.} Let \(\alpha\) be a type. Let \(a\) and \(b\) be elements of \(\alpha\). Then \(\operatorname{Sym2\mathord{:}mk}((a,\allowbreak{} b)) = \operatorname{Sym2\mathord{:}mk}((b,\allowbreak{} a))\).
\item[\paperlineanchor{premise:53796d322e6d61705f6964}\paperlink{premise:53796d322e6d61705f6964}{P164.}] \refstepcounter{paperpremise}\label{premise:53796d322e6d61705f6964}\textbf{Map id.} Let \(\alpha\) be a type. Then \(\operatorname{Sym2\mathord{:}map}\left(\operatorname{id}\right) = \operatorname{id}\).
\item[\paperlineanchor{premise:53796d322e72656c5f69666627}\paperlink{premise:53796d322e72656c5f69666627}{P165.}] \refstepcounter{paperpremise}\label{premise:53796d322e72656c5f69666627}\textbf{Rel iff'.} Let \(\alpha\) be a type. Let \(p\) be an element of \(\alpha \times \alpha\). Let \(q\) be an element of \(\alpha \times \alpha\). Then \(\operatorname{Rel}\left(\alpha, p, q\right) \iff p = q \lor p = \operatorname{swap}\left(q\right)\).
\item[\paperlineanchor{premise:616e645f66616c7365}\paperlink{premise:616e645f66616c7365}{P166.}] \refstepcounter{paperpremise}\label{premise:616e645f66616c7365}\textbf{And false.} Let \(p\) be a proposition. Then \(\left(p \land \mathrm{False}\right) = \mathrm{False}\).
\item[\paperlineanchor{premise:616e645f696d70}\paperlink{premise:616e645f696d70}{P167.}] \refstepcounter{paperpremise}\label{premise:616e645f696d70}\textbf{And imp.} Let \(a\), \(b\), and \(c\) be propositions. Then \(a \land b \implies c \iff a \implies b \implies c\).
\item[\paperlineanchor{premise:616e645f73656c66}\paperlink{premise:616e645f73656c66}{P168.}] \refstepcounter{paperpremise}\label{premise:616e645f73656c66}\textbf{And self.} Let \(p\) be a proposition. Then \(\left(p \land p\right) = p\).
\item[\paperlineanchor{premise:616e645f74727565}\paperlink{premise:616e645f74727565}{P169.}] \refstepcounter{paperpremise}\label{premise:616e645f74727565}\textbf{And true.} Let \(p\) be a proposition. Then \(\left(p \land \mathrm{True}\right) = p\).
\item[\paperlineanchor{premise:65715f66616c7365}\paperlink{premise:65715f66616c7365}{P170.}] \refstepcounter{paperpremise}\label{premise:65715f66616c7365}\textbf{Equals false.} Let \(p\) be a proposition. Assume \(\neg p\). Then \(p = \mathrm{False}\).
\item[\paperlineanchor{premise:65715f6f665f686571}\paperlink{premise:65715f6f665f686571}{P171.}] \refstepcounter{paperpremise}\label{premise:65715f6f665f686571}\textbf{Equals of heq.} Let \(\alpha\) be a sort. Let \(a\) and \(a^{'}\) be elements of \(\alpha\). Assume \(\left(a\, :\, \alpha\right)\text{ is heterogeneously equal to }\left(a^{'}\, :\, \alpha\right)\). Then \(a = a^{'}\).
\item[\paperlineanchor{premise:65715f73656c66}\paperlink{premise:65715f73656c66}{P172.}] \refstepcounter{paperpremise}\label{premise:65715f73656c66}\textbf{Equals self.} Let \(\alpha\) be a sort. Let \(a\) be an element of \(\alpha\). Then \(\left(a = a\right) = \mathrm{True}\).
\item[\paperlineanchor{premise:65715f74727565}\paperlink{premise:65715f74727565}{P173.}] \refstepcounter{paperpremise}\label{premise:65715f74727565}\textbf{Equals true.} Let \(p\) be a proposition. Assume \(p\). Then \(p = \mathrm{True}\).
\item[\paperlineanchor{premise:6578697374735f636f6e7374}\paperlink{premise:6578697374735f636f6e7374}{P174.}] \refstepcounter{paperpremise}\label{premise:6578697374735f636f6e7374}\textbf{Exists constant.} Let \(b\) be a proposition. Let \(\alpha\) be a sort. Assume \(\alpha\text{ is nonempty}\). Then \(\left(\exists x\, :\, \alpha,\; b\right) \iff b\).
\item[\paperlineanchor{premise:6578697374735f70726f705f636f6e6772}\paperlink{premise:6578697374735f70726f705f636f6e6772}{P175.}] \refstepcounter{paperpremise}\label{premise:6578697374735f70726f705f636f6e6772}\textbf{Exists prop congruence.} Let \(p\) and \(p^{'}\) be propositions. Let \(q\) be a function of type \(p \to \mathsf{Prop}\). Let \(q^{'}\) be a function of type \(p \to \mathsf{Prop}\). Assume \(\forall h\, :\, p,\; q\left(h\right) \iff q^{'}\left(h\right)\). Let \(\mathit{hp}\) be a proof of \(p \iff p^{'}\). Then \((\exists x\, :\, p,\; q(x)) \iff (\exists h\text{ proving }p^{'},\; q^{'}(\operatorname{Iff\mathord{:}mpr}(\mathit{hp},\allowbreak{} h)))\).
\item[\paperlineanchor{premise:6578697374735f747275655f6c656674}\paperlink{premise:6578697374735f747275655f6c656674}{P176.}] \refstepcounter{paperpremise}\label{premise:6578697374735f747275655f6c656674}\textbf{Exists true left.} Let \(p\) be a function of type \(\mathrm{True} \to \mathsf{Prop}\). Then \(\left(\exists x\, :\, \mathrm{True},\; p(x)\right) \iff p\left(\mathrm{trivial}\right)\).
\item[\paperlineanchor{premise:66616c73655f696666}\paperlink{premise:66616c73655f696666}{P177.}] \refstepcounter{paperpremise}\label{premise:66616c73655f696666}\textbf{False if and only if.} Let \(p\) be a proposition. Then \(\left(\mathrm{False} \iff p\right) = \left(\neg p\right)\).
\item[\paperlineanchor{premise:66616c73655f6f72}\paperlink{premise:66616c73655f6f72}{P178.}] \refstepcounter{paperpremise}\label{premise:66616c73655f6f72}\textbf{False or.} Let \(p\) be a proposition. Then \(\left(\mathrm{False} \lor p\right) = p\).
\item[\paperlineanchor{premise:666f72616c6c5f636f6e6772}\paperlink{premise:666f72616c6c5f636f6e6772}{P179.}] \refstepcounter{paperpremise}\label{premise:666f72616c6c5f636f6e6772}\textbf{Forall congruence.} Let \(\alpha\) be a sort. Let \(p\) and \(q\) be functions of type \(\alpha \to \mathsf{Prop}\). Assume \(\forall a\, :\, \alpha,\; p\left(a\right) = q\left(a\right)\). Then \((\forall a\, :\, \alpha,\; p(a)) = (\forall a\, :\, \alpha,\; q(a))\).
\item[\paperlineanchor{premise:66756e657874}\paperlink{premise:66756e657874}{P180.}] \refstepcounter{paperpremise}\label{premise:66756e657874}\textbf{Function extensionality.} Let \(\alpha\) be a sort. Let \(\beta\) be a function of type \(\alpha \to \mathsf{Sort}\). Let \(f\) be a function of type \(\Pi_{x\, :\, \alpha}\, \beta\left(x\right)\). Let \(g\) be a function of type \(\Pi_{x\, :\, \alpha}\, \beta\left(x\right)\). Assume \(\forall x\, :\, \alpha,\; f\left(x\right) = g\left(x\right)\). Then \(f = g\).
\item[\paperlineanchor{premise:6966665f66616c7365}\paperlink{premise:6966665f66616c7365}{P181.}] \refstepcounter{paperpremise}\label{premise:6966665f66616c7365}\textbf{If and only if false.} Let \(p\) be a proposition. Then \(\left(p \iff \mathrm{False}\right) = \left(\neg p\right)\).
\item[\paperlineanchor{premise:6966665f73656c66}\paperlink{premise:6966665f73656c66}{P182.}] \refstepcounter{paperpremise}\label{premise:6966665f73656c66}\textbf{If and only if self.} Let \(p\) be a proposition. Then \(\left(p \iff p\right) = \mathrm{True}\).
\item[\paperlineanchor{premise:6966665f74727565}\paperlink{premise:6966665f74727565}{P183.}] \refstepcounter{paperpremise}\label{premise:6966665f74727565}\textbf{If and only if true.} Let \(p\) be a proposition. Then \(\left(p \iff \mathrm{True}\right) = p\).
\item[\paperlineanchor{premise:696d705f66616c7365}\paperlink{premise:696d705f66616c7365}{P184.}] \refstepcounter{paperpremise}\label{premise:696d705f66616c7365}\textbf{Imp false.} Let \(a\) be a proposition. Then \(a \implies \mathrm{False} \iff \neg a\).
\item[\paperlineanchor{premise:696d705f6966665f6e6f745f6f72}\paperlink{premise:696d705f6966665f6e6f745f6f72}{P185.}] \refstepcounter{paperpremise}\label{premise:696d705f6966665f6e6f745f6f72}\textbf{Imp if and only if not or.} Let \(a\) and \(b\) be propositions. Then \(a \implies b \iff \neg a \lor b\).
\item[\paperlineanchor{premise:696d706c6965735f636f6e6772}\paperlink{premise:696d706c6965735f636f6e6772}{P186.}] \refstepcounter{paperpremise}\label{premise:696d706c6965735f636f6e6772}\textbf{Implies congruence.} Let \(p_{1}\) and \(p_{2}\) be sorts. Let \(q_{1}\) and \(q_{2}\) be sorts. Assume \(p_{1} = p_{2}\). Assume \(q_{1} = q_{2}\). Then \(\left(p_{1} \to q_{1}\right) = \left(p_{2} \to q_{2}\right)\).
\item[\paperlineanchor{premise:696d706c6965735f636f6e67725f637478}\paperlink{premise:696d706c6965735f636f6e67725f637478}{P187.}] \refstepcounter{paperpremise}\label{premise:696d706c6965735f636f6e67725f637478}\textbf{Implies congruence ctx.} Let \(p_{1}\), \(p_{2}\), \(q_{1}\), and \(q_{2}\) be propositions. Assume \(p_{1} = p_{2}\). Assume \(p_{2} \implies q_{1} = q_{2}\). Then \(\left(p_{1} \implies q_{1}\right) = \left(p_{2} \implies q_{2}\right)\).
\item[\paperlineanchor{premise:696d706c6965735f74727565}\paperlink{premise:696d706c6965735f74727565}{P188.}] \refstepcounter{paperpremise}\label{premise:696d706c6965735f74727565}\textbf{Implies true.} Let \(\alpha\) be a sort. Then \(\left(\forall a\, :\, \alpha,\; \mathrm{True}\right) = \mathrm{True}\).
\item[\paperlineanchor{premise:696e73744e6f6e656d7074794f66496e68616269746564}\paperlink{premise:696e73744e6f6e656d7074794f66496e68616269746564}{P189.}] \refstepcounter{paperpremise}\label{premise:696e73744e6f6e656d7074794f66496e68616269746564}\textbf{Inst nonempty of inhabited.} Let \(\alpha\) be a sort. Fix a structure \(\mathit{inst}\) of type \(\operatorname{Inhabited}\left(\alpha\right)\). Then \(\alpha\text{ is nonempty}\).
\item[\paperlineanchor{premise:6c655f6f665f6c655f6f665f6571}\paperlink{premise:6c655f6f665f6c655f6f665f6571}{P190.}] \refstepcounter{paperpremise}\label{premise:6c655f6f665f6c655f6f665f6571}\textbf{At most of at most of equals.} Let \(\alpha\) be a type. Let \(a\), \(b\), and \(c\) be elements of \(\alpha\). Fix a structure \(\mathit{inst}\) of type \(\operatorname{LE}\left(\alpha\right)\). Assume \(a \le b\). Assume \(b = c\). Then \(a \le c\).
\item[\paperlineanchor{premise:6c655f7265666c}\paperlink{premise:6c655f7265666c}{P191.}] \refstepcounter{paperpremise}\label{premise:6c655f7265666c}\textbf{At most refl.} Let \(\alpha\) be a type. Equip \(\alpha\) with a preorder. Let \(a\) be an element of \(\alpha\). Then \(a \le a\).
\item[\paperlineanchor{premise:6c655f7375705f7269676874}\paperlink{premise:6c655f7375705f7269676874}{P192.}] \refstepcounter{paperpremise}\label{premise:6c655f7375705f7269676874}\textbf{At most sup right.} Let \(\alpha\) be a type. Fix a structure \(\mathit{inst}\) of type \(\operatorname{Semilattice\,Sup}\left(\alpha\right)\). Let \(a\) and \(b\) be elements of \(\alpha\). Then \(b \le \operatorname{max}\left(a, b\right)\).
\item[\paperlineanchor{premise:6d756c5f6f6e65}\paperlink{premise:6d756c5f6f6e65}{P193.}] \refstepcounter{paperpremise}\label{premise:6d756c5f6f6e65}\textbf{Product one.} Let \(M\) be a type. Fix a structure \(\mathit{inst}\) of type \(\operatorname{Mul\,One\,Class}\left(M\right)\). Let \(a\) be an element of \(M\). Then \(a \cdot 1 = a\).
\item[\paperlineanchor{premise:6e6f745f616e64}\paperlink{premise:6e6f745f616e64}{P194.}] \refstepcounter{paperpremise}\label{premise:6e6f745f616e64}\textbf{Not and.} Let \(a\) and \(b\) be propositions. Then \(\neg \left(a \land b\right) \iff a \implies \neg b\).
\item[\paperlineanchor{premise:6e6f745f66616c73655f65715f74727565}\paperlink{premise:6e6f745f66616c73655f65715f74727565}{P195.}] \refstepcounter{paperpremise}\label{premise:6e6f745f66616c73655f65715f74727565}\textbf{Not false equals true.} \(\left(\neg \mathrm{False}\right) = \mathrm{True}\).
\item[\paperlineanchor{premise:6e6f745f666f72616c6c5f6e6f74}\paperlink{premise:6e6f745f666f72616c6c5f6e6f74}{P196.}] \refstepcounter{paperpremise}\label{premise:6e6f745f666f72616c6c5f6e6f74}\textbf{Not forall not.} Let \(\alpha\) be a sort. Let \(p\) be a function of type \(\alpha \to \mathsf{Prop}\). Then \(\neg (\forall x\, :\, \alpha,\; \neg p(x)) \iff (\exists x\, :\, \alpha,\; p(x))\).
\item[\paperlineanchor{premise:6e6f745f6f72}\paperlink{premise:6e6f745f6f72}{P197.}] \refstepcounter{paperpremise}\label{premise:6e6f745f6f72}\textbf{Not or.} Let \(p\) and \(q\) be propositions. Then \(\neg \left(p \lor q\right) \iff \neg p \land \neg q\).
\item[\paperlineanchor{premise:6e6f745f747275655f65715f66616c7365}\paperlink{premise:6e6f745f747275655f65715f66616c7365}{P198.}] \refstepcounter{paperpremise}\label{premise:6e6f745f747275655f65715f66616c7365}\textbf{Not true equals false.} \(\left(\neg \mathrm{True}\right) = \mathrm{False}\).
\item[\paperlineanchor{premise:6f665f6465636964655f65715f74727565}\paperlink{premise:6f665f6465636964655f65715f74727565}{P199.}] \refstepcounter{paperpremise}\label{premise:6f665f6465636964655f65715f74727565}\textbf{Of decide equals true.} Let \(p\) be a proposition. Assume \(p\) is decidable. Assume \(\operatorname{decide}\left(p\right) = \mathrm{true}\). Then \(p\).
\item[\paperlineanchor{premise:6f665f65715f74727565}\paperlink{premise:6f665f65715f74727565}{P200.}] \refstepcounter{paperpremise}\label{premise:6f665f65715f74727565}\textbf{Of equals true.} Let \(p\) be a proposition. Assume \(p = \mathrm{True}\). Then \(p\).
\item[\paperlineanchor{premise:6f725f66616c7365}\paperlink{premise:6f725f66616c7365}{P201.}] \refstepcounter{paperpremise}\label{premise:6f725f66616c7365}\textbf{Or false.} Let \(p\) be a proposition. Then \(\left(p \lor \mathrm{False}\right) = p\).
\item[\paperlineanchor{premise:6f725f74727565}\paperlink{premise:6f725f74727565}{P202.}] \refstepcounter{paperpremise}\label{premise:6f725f74727565}\textbf{Or true.} Let \(p\) be a proposition. Then \(\left(p \lor \mathrm{True}\right) = \mathrm{True}\).
\item[\paperlineanchor{premise:73646966665f6c655f7364696666}\paperlink{premise:73646966665f6c655f7364696666}{P203.}] \refstepcounter{paperpremise}\label{premise:73646966665f6c655f7364696666}\textbf{Sdiff at most sdiff.} Let \(\alpha\) be a type. Fix a structure \(\mathit{inst}\) of type \(\operatorname{Generalized\,Coheyting\,Algebra}\left(\alpha\right)\). Let \(a\), \(b\), \(c\), and \(d\) be elements of \(\alpha\). Assume \(a \le b\). Assume \(c \le d\). Then \(a \setminus d \le b \setminus c\).
\item[\paperlineanchor{premise:7375705f636f6d6d}\paperlink{premise:7375705f636f6d6d}{P204.}] \refstepcounter{paperpremise}\label{premise:7375705f636f6d6d}\textbf{Sup commutation.} Let \(\alpha\) be a type. Fix a structure \(\mathit{inst}\) of type \(\operatorname{Semilattice\,Sup}\left(\alpha\right)\). Let \(a\) and \(b\) be elements of \(\alpha\). Then \(\operatorname{max}\left(a, b\right) = \operatorname{max}\left(b, a\right)\).
\item[\paperlineanchor{premise:7375705f6f665f6c655f6c656674}\paperlink{premise:7375705f6f665f6c655f6c656674}{P205.}] \refstepcounter{paperpremise}\label{premise:7375705f6f665f6c655f6c656674}\textbf{Sup of at most left.} Let \(\alpha\) be a type. Fix a structure \(\mathit{inst}\) of type \(\operatorname{Semilattice\,Sup}\left(\alpha\right)\). Let \(a\) and \(b\) be elements of \(\alpha\). Assume \(b \le a\). Then \(\operatorname{max}\left(a, b\right) = a\).
\item[\paperlineanchor{premise:73796d6d446966665f636f6d6d}\paperlink{premise:73796d6d446966665f636f6d6d}{P206.}] \refstepcounter{paperpremise}\label{premise:73796d6d446966665f636f6d6d}\textbf{Symm diff commutation.} Let \(\alpha\) be a type. Fix a structure \(\mathit{inst}\) of type \(\operatorname{Generalized\,Coheyting\,Algebra}\left(\alpha\right)\). Let \(a\) and \(b\) be elements of \(\alpha\). Then \(\operatorname{symm\,Diff}\left(a, b\right) = \operatorname{symm\,Diff}\left(b, a\right)\).
\item[\paperlineanchor{premise:747275655f616e64}\paperlink{premise:747275655f616e64}{P207.}] \refstepcounter{paperpremise}\label{premise:747275655f616e64}\textbf{True and.} Let \(p\) be a proposition. Then \(\left(\mathrm{True} \land p\right) = p\).
\item[\paperlineanchor{premise:747275655f696666}\paperlink{premise:747275655f696666}{P208.}] \refstepcounter{paperpremise}\label{premise:747275655f696666}\textbf{True if and only if.} Let \(p\) be a proposition. Then \(\left(\mathrm{True} \iff p\right) = p\).
\item[\paperlineanchor{premise:747275655f6f72}\paperlink{premise:747275655f6f72}{P209.}] \refstepcounter{paperpremise}\label{premise:747275655f6f72}\textbf{True or.} Let \(p\) be a proposition. Then \(\left(\mathrm{True} \lor p\right) = \mathrm{True}\).
\item[\paperlineanchor{premise:74776f5f6d756c}\paperlink{premise:74776f5f6d756c}{P210.}] \refstepcounter{paperpremise}\label{premise:74776f5f6d756c}\textbf{Two product.} Let \(\alpha\) be a type. Fix a structure \(\mathit{inst}\) of type \(\operatorname{Non\,Assoc\,Semiring}\left(\alpha\right)\). Let \(n\) be an element of \(\alpha\). Then \(2 \cdot n = n + n\).
\end{description}

\section{Subgraph}
\begin{lemma}[{To edge equals}]\label{theorem:1}
Let \(V\) be a type. Let \(G\) be a simple graph on \(V\). Let \(M\) be a subgraph of \(G\). Let \(v\) and \(w\) be elements of \(V\). Let \(h\) be a proof of \(M\text{ is a matching}\). Let \(\mathit{hv}\) be a proof of \(v \in V(M)\). Let \(\mathit{hvw}\) be a proof of \(v \sim_{M} w\). Then
\begin{center}
\(\displaystyle \operatorname{to\,Edge}((v\, :\, V(M))) = (\operatorname{Sym2\mathord{:}mk}((v,\allowbreak{} w))\, :\, \operatorname{edge\,Set}(M)).\)
\end{center}
\end{lemma}

\begin{proof}
By definition, we obtain \(\operatorname{Sym2\mathord{:}mk}((v,\allowbreak{} \text{the fixed choice from }\{x\in V \mid v \sim_{M} x \land (\forall y\, :\, V,\; v \sim_{M} y \implies y = x)\})) = \operatorname{Sym2\mathord{:}mk}((v,\allowbreak{} \text{the fixed choice from }\{x\in V \mid v \sim_{M} x \land (\forall y\, :\, V,\; v \sim_{M} y \implies y = x)\}))\).
For the required implication, argue as follows.
\begin{proofstep}
Fix \(\mathit{snd}\) in \(V\).
Fix \(\mathit{snd}_{1}\) in \(V\).
Assume \(\mathit{snd} = \mathit{snd}_{1}\).
By definition, we obtain \(\left(v, \mathit{snd}\right) = \left(v, \mathit{snd}\right)\).
Use the hypothesis \(\mathit{snd} = \mathit{snd}_{1}\).
Substituting equals, we obtain \(\left(v, \mathit{snd}\right) = \left(v, \mathit{snd}_{1}\right)\).
\end{proofstep}
By definition, we obtain \(v = v\).
Use the hypothesis \(v \sim_{M} w\).
By \paperlink{premise:53696d706c6547726170682e53756267726170682e656467655f76657274}{Prerequisite~P\ref{premise:53696d706c6547726170682e53756267726170682e656467655f76657274}} with \(\mathit{self} = M\), we obtain \(v \in V(M)\).
Applying the hypothesis \(M\text{ is a matching}\), we obtain \(\exists! w_{1}\, :\, V,\; v \sim_{M} w_{1}\).
By \paperlink{premise:4578697374732e63686f6f73655f73706563}{Prerequisite~P\ref{premise:4578697374732e63686f6f73655f73706563}}, we obtain \(v \sim_{M} (\text{the fixed choice from }\{x\in V \mid v \sim_{M} x \land (\forall y\, :\, V,\; v \sim_{M} y \implies y = x)\}) \land (\forall y\, :\, V,\; v \sim_{M} y \implies y = (\text{the fixed choice from }\{x\in V \mid v \sim_{M} x \land (\forall y_{1}\, :\, V,\; v \sim_{M} y_{1} \implies y_{1} = x)\}))\).
Use the hypothesis \(v \sim_{M} w\).
Taking the second conjunct, we obtain \(w = (\text{the fixed choice from }\{x\in V \mid v \sim_{M} x \land (\forall y\, :\, V,\; v \sim_{M} y \implies y = x)\})\).
Reversing the equality, we obtain \((\text{the fixed choice from }\{x\in V \mid v \sim_{M} x \land (\forall y\, :\, V,\; v \sim_{M} y \implies y = x)\}) = w\).
Substituting equals, we obtain \((v,\allowbreak{} \text{the fixed choice from }\{x\in V \mid v \sim_{M} x \land (\forall y\, :\, V,\; v \sim_{M} y \implies y = x)\}) = (v,\allowbreak{} w)\).
Substituting equals, we obtain \(\operatorname{Sym2\mathord{:}mk}((v,\allowbreak{} \text{the fixed choice from }\{x\in V \mid v \sim_{M} x \land (\forall y\, :\, V,\; v \sim_{M} y \implies y = x)\})) = \operatorname{Sym2\mathord{:}mk}((v,\allowbreak{} w))\).
Rewriting with \paperlink{premise:537562747970652e6d6b5f65715f6d6b}{Prerequisite~P\ref{premise:537562747970652e6d6b5f65715f6d6b}} (using \(v \sim_{M} w\)), \paperlink{premise:416e642e6c656674}{Prerequisite~P\ref{premise:416e642e6c656674}}, \paperlink{premise:4578697374732e63686f6f73655f73706563}{Prerequisite~P\ref{premise:4578697374732e63686f6f73655f73706563}}, \(M\text{ is a matching}\) and \paperlink{premise:537562747970652e70726f7065727479}{Prerequisite~P\ref{premise:537562747970652e70726f7065727479}} in the preceding result gives \(\operatorname{to\,Edge}((v\, :\, V(M))) = (\operatorname{Sym2\mathord{:}mk}((v,\allowbreak{} w))\, :\, \operatorname{edge\,Set}(M))\).
\end{proof}

\begin{theorem}[{Surjective}]\label{theorem:2}
Let \(V\) be a type. Let \(G\) be a simple graph on \(V\). Let \(M\) be a subgraph of \(G\). Let \(h\) be a proof of \(M\text{ is a matching}\). Then \(\operatorname{to\,Edge}\text{ is surjective}\).
\end{theorem}

\begin{proof}
Fix \(b\) in \(\operatorname{edge\,Set}\left(M\right)\).
Consider the following case.
\begin{proofstep}
Fix \(\mathit{val}\) in \(\operatorname{Sym2}\left(V\right)\).
Assume \(\mathit{val} \in \operatorname{edge\,Set}\left(M\right)\).
Consider the following case.
\begin{proofstep}
Fix \(a\) in \(V \times V\).
Assume \(\operatorname{Quot\mathord{:}mk}(\operatorname{Rel}(V),\allowbreak{} a) \in \operatorname{edge\,Set}(M)\).
Consider the following case.
\begin{proofstep}
Fix \(x\) in \(V\).
Fix \(y\) in \(V\).
Assume \(\operatorname{Quot\mathord{:}mk}(\operatorname{Rel}(V),\allowbreak{} (x,\allowbreak{} y)) \in \operatorname{edge\,Set}(M)\).
Choose the witness \(\left(x\, :\, V(M)\right)\).
\par\smallskip\noindent\emph{Property.}\enspace
Use the hypothesis \(\operatorname{Quot\mathord{:}mk}(\operatorname{Rel}(V),\allowbreak{} (x,\allowbreak{} y)) \in \operatorname{edge\,Set}(M)\).
By \paperlink{premise:53696d706c6547726170682e53756267726170682e656467655f76657274}{Prerequisite~P\ref{premise:53696d706c6547726170682e53756267726170682e656467655f76657274}} with \(\mathit{self} = M\), we obtain \(\left(x_{1} \mapsto x_{1} \in V(M)\right)\left(x\right)\).
Use the hypothesis \(\operatorname{Quot\mathord{:}mk}(\operatorname{Rel}(V),\allowbreak{} (x,\allowbreak{} y)) \in \operatorname{edge\,Set}(M)\).
By \paperlink{premise:53696d706c6547726170682e53756267726170682e656467655f76657274}{Prerequisite~P\ref{premise:53696d706c6547726170682e53756267726170682e656467655f76657274}} with \(\mathit{self} = M\), we obtain \(x \in V(M)\).
By \paperlink{theorem:1}{Lemma~\ref{theorem:1}}, we obtain \(\operatorname{to\,Edge}((\pi_1((x,\allowbreak{} y))\, :\, V(M))) = (\operatorname{Sym2\mathord{:}mk}((\pi_1((x,\allowbreak{} y)),\allowbreak{} \pi_2((x,\allowbreak{} y))))\, :\, \operatorname{edge\,Set}(M))\).
This witness proves existence. Hence \(\exists a_{1}\, :\, V(M),\; \operatorname{to\,Edge}(a_{1}) = (\operatorname{Quot\mathord{:}mk}(\operatorname{Rel}(V),\allowbreak{} (x,\allowbreak{} y))\, :\, \operatorname{edge\,Set}(M))\).
\end{proofstep}
Use the hypothesis \(\operatorname{Quot\mathord{:}mk}(\operatorname{Rel}(V),\allowbreak{} a) \in \operatorname{edge\,Set}(M)\).
Every ordered pair has two components. Hence \(\exists a_{1}\, :\, V(M),\; \operatorname{to\,Edge}(a_{1}) = (\operatorname{Quot\mathord{:}mk}(\operatorname{Rel}(V),\allowbreak{} a)\, :\, \operatorname{edge\,Set}(M))\).
\end{proofstep}
Use the hypothesis \(\mathit{val} \in \operatorname{edge\,Set}\left(M\right)\).
Every element of the quotient has a representative. Hence \(\exists a\, :\, V(M),\; \operatorname{to\,Edge}(a) = (\mathit{val}\, :\, \operatorname{edge\,Set}(M))\).
\end{proofstep}
Using the underlying element and its defining property, we obtain \(\exists a\, :\, V(M),\; \operatorname{to\,Edge}\left(a\right) = b\).
\end{proof}

\begin{theorem}[{To edge equals to edge}]\label{theorem:3}
Let \(V\) be a type. Let \(G\) be a simple graph on \(V\). Let \(M\) be a subgraph of \(G\). Let \(v\) and \(w\) be elements of \(V\). Let \(h\) be a proof of \(M\text{ is a matching}\). Let \(\mathit{hv}\) be a proof of \(v \in V(M)\). Let \(\mathit{hw}\) be a proof of \(w \in V(M)\). Assume \(v \sim_{M} w\). Then
\begin{center}
\(\displaystyle \operatorname{to\,Edge}((v\, :\, V(M))) = \operatorname{to\,Edge}((w\, :\, V(M))).\)
\end{center}
\end{theorem}

\begin{proof}
Rewriting with \paperlink{premise:53796d322e65715f73776170}{Prerequisite~P\ref{premise:53796d322e65715f73776170}}, \paperlink{premise:537562747970652e6d6b5f65715f6d6b}{Prerequisite~P\ref{premise:537562747970652e6d6b5f65715f6d6b}} (using \(v \sim_{M} w\)), \paperlink{premise:53696d706c6547726170682e53756267726170682e73796d6d}{Prerequisite~P\ref{premise:53696d706c6547726170682e53756267726170682e73796d6d}} (using \(v \sim_{M} w\)), \paperlink{theorem:1}{Lemma~\ref{theorem:1}} (using \(M\text{ is a matching}\) and \(w \in V(M)\)) and \paperlink{theorem:1}{Lemma~\ref{theorem:1}} (using \(M\text{ is a matching}\) and \(v \in V(M)\) and \(v \sim_{M} w\)) gives \(\operatorname{to\,Edge}((v\, :\, V(M))) = \operatorname{to\,Edge}((w\, :\, V(M)))\).
\end{proof}

\begin{lemma}[{Map of at most}]\label{theorem:4}
Let \(V\) be a type. Let \(G\) and \(G^{'}\) be simple graphs on \(V\). Let \(M\) be a subgraph of \(G\). Assume \(M\text{ is a matching}\). Let \(h_{1}\) be a proof of \(G \subseteq G^{'}\). Then
\begin{center}
\(\displaystyle \operatorname{Simple\,Graph\mathord{:}Subgraph\mathord{:}map}(\operatorname{incl}_{G,G^{'}},\allowbreak{} M)\text{ is a matching}.\)
\end{center}
\end{lemma}

\begin{proof}
Fix \(v\) in \(V\).
Assume \(v \in V(\operatorname{Simple\,Graph\mathord{:}Subgraph\mathord{:}map}(\operatorname{incl}_{G,G^{'}},\allowbreak{} M))\).
By \paperlink{premise:5365742e6d656d5f696d616765}{Prerequisite~P\ref{premise:5365742e6d656d5f696d616765}} with \(f = \operatorname{DFun\,Like\mathord{:}coe}\left(\operatorname{incl}_{G,G^{'}}\right)\) and \(s = V(M)\) and \(y = v\), we obtain \(v \in \operatorname{DFun\,Like\mathord{:}coe}(\operatorname{incl}_{G,G^{'}})[V(M)] \iff (\exists x\, :\, V,\; x \in V(M) \land \operatorname{incl}_{G,G^{'}}(x) = v)\).
Use the hypothesis \(v \in V(\operatorname{Simple\,Graph\mathord{:}Subgraph\mathord{:}map}(\operatorname{incl}_{G,G^{'}},\allowbreak{} M))\).
By the forward implication, we obtain \(\exists x\, :\, V,\; x \in V(M) \land \operatorname{incl}_{G,G^{'}}(x) = v\).
Consider the following case.
\begin{proofstep}
Fix \(w\) in \(V\).
Assume \(w \in V(M) \land \operatorname{incl}_{G,G^{'}}(w) = v\).
Use the hypothesis \(w \in V(M) \land \operatorname{incl}_{G,G^{'}}(w) = v\).
Consider the following case.
\begin{proofstep}
Assume \(w \in V(M)\).
Assume \(\operatorname{incl}_{G,G^{'}}(w) = v\).
Use the hypothesis \(w \in V(M)\).
Applying the hypothesis \(M\text{ is a matching}\), we obtain \(\exists! w_{1}\, :\, V,\; w \sim_{M} w_{1}\).
Consider the following case.
\begin{proofstep}
Fix \(w_{1}\) in \(V\).
Assume \(w \sim_{M} w_{1} \land \left(\forall y\, :\, V,\; w \sim_{M} y \implies y = w_{1}\right)\).
Choose the witness \(w_{1}\).
The substitution uses the following argument.
\begin{proofstep}
Fix \(y\) in \(V\).
By \paperlink{premise:52656c6174696f6e2e6d61705f69645f6964}{Prerequisite~P\ref{premise:52656c6174696f6e2e6d61705f69645f6964}} with \(r = \operatorname{Simple\,Graph\mathord{:}Subgraph\mathord{:}Adj}\left(M\right)\), we obtain \(\operatorname{Map}(\operatorname{Simple\,Graph\mathord{:}Subgraph\mathord{:}Adj}(M),\allowbreak{} \operatorname{id},\allowbreak{} \operatorname{id}) = \operatorname{Simple\,Graph\mathord{:}Subgraph\mathord{:}Adj}(M)\).
Applying the same function to equal arguments, we obtain \(\operatorname{Map}(\operatorname{Simple\,Graph\mathord{:}Subgraph\mathord{:}Adj}(M),\allowbreak{} \operatorname{id},\allowbreak{} \operatorname{id},\allowbreak{} v,\allowbreak{} y) = (v \sim_{M} y)\).
By definition, we obtain \(\left(y = w_{1}\right) = \left(y = w_{1}\right)\).
By \paperlink{premise:696d706c6965735f636f6e6772}{Prerequisite~P\ref{premise:696d706c6965735f636f6e6772}} with \(p_{1} = \operatorname{Map}(\operatorname{Simple\,Graph\mathord{:}Subgraph\mathord{:}Adj}(M),\allowbreak{} \operatorname{id},\allowbreak{} \operatorname{id},\allowbreak{} v,\allowbreak{} y)\) and \(p_{2} = v \sim_{M} y\) and \(q_{1} = y = w_{1}\) and \(q_{2} = y = w_{1}\), we obtain \((\operatorname{Map}(\operatorname{Simple\,Graph\mathord{:}Subgraph\mathord{:}Adj}(M),\allowbreak{} \operatorname{id},\allowbreak{} \operatorname{id},\allowbreak{} v,\allowbreak{} y) \implies y = w_{1}) = (v \sim_{M} y \implies y = w_{1})\).
\end{proofstep}
Use the hypothesis \(w \sim_{M} w_{1} \land \left(\forall y\, :\, V,\; w \sim_{M} y \implies y = w_{1}\right)\).
Use the hypothesis \(\operatorname{incl}_{G,G^{'}}(w) = v\).
Substituting equals, we obtain \(v \sim_{M} w_{1} \land \left(\forall y\, :\, V,\; v \sim_{M} y \implies y = w_{1}\right)\).
Rewriting with \paperlink{premise:52656c6174696f6e2e6d61705f69645f6964}{Prerequisite~P\ref{premise:52656c6174696f6e2e6d61705f69645f6964}}, \paperlink{premise:666f72616c6c5f636f6e6772}{Prerequisite~P\ref{premise:666f72616c6c5f636f6e6772}} and the identity established above in the preceding result gives \(v \sim_{\operatorname{Simple\,Graph\mathord{:}Subgraph\mathord{:}map}(\operatorname{incl}_{G,G^{'}},\allowbreak{} M)} w_{1} \land (\forall y\, :\, V,\; v \sim_{\operatorname{Simple\,Graph\mathord{:}Subgraph\mathord{:}map}(\operatorname{incl}_{G,G^{'}},\allowbreak{} M)} y \implies y = w_{1})\).
This witness proves existence. Hence \(\exists x\, :\, V,\; v \sim_{\operatorname{Simple\,Graph\mathord{:}Subgraph\mathord{:}map}(\operatorname{incl}_{G,G^{'}},\allowbreak{} M)} x \land (\forall y\, :\, V,\; v \sim_{\operatorname{Simple\,Graph\mathord{:}Subgraph\mathord{:}map}(\operatorname{incl}_{G,G^{'}},\allowbreak{} M)} y \implies y = x)\).
\end{proofstep}
Using the existential hypothesis, we obtain \(\exists! w_{1}\, :\, V,\; v \sim_{\operatorname{Simple\,Graph\mathord{:}Subgraph\mathord{:}map}(\operatorname{incl}_{G,G^{'}},\allowbreak{} M)} w_{1}\).
\end{proofstep}
Using the two conjuncts, we obtain \(\exists! w_{1}\, :\, V,\; v \sim_{\operatorname{Simple\,Graph\mathord{:}Subgraph\mathord{:}map}(\operatorname{incl}_{G,G^{'}},\allowbreak{} M)} w_{1}\).
\end{proofstep}
Using the existential hypothesis, we obtain \(\exists! w\, :\, V,\; v \sim_{\operatorname{Simple\,Graph\mathord{:}Subgraph\mathord{:}map}(\operatorname{incl}_{G,G^{'}},\allowbreak{} M)} w\).
\end{proof}

\begin{lemma}[{Equals}]\label{theorem:5}
Let \(V\) be a type. Let \(G\) be a simple graph on \(V\). Let \(M\) be a subgraph of \(G\). Let \(u\), \(v\), and \(w\) be elements of \(V\). Assume \(M\text{ is a matching}\). Assume \(u \sim_{M} v\). Assume \(u \sim_{M} w\). Then \(v = w\).
\end{lemma}

\begin{proof}
Use the hypothesis \(u \sim_{M} v\).
By \paperlink{premise:53696d706c6547726170682e53756267726170682e656467655f76657274}{Prerequisite~P\ref{premise:53696d706c6547726170682e53756267726170682e656467655f76657274}} with \(\mathit{self} = M\), we obtain \(u \in V(M)\).
Applying the hypothesis \(M\text{ is a matching}\), we obtain \(\exists! w_{1}\, :\, V,\; u \sim_{M} w_{1}\).
Use the hypothesis \(u \sim_{M} v\).
Use the hypothesis \(u \sim_{M} w\).
By \paperlink{premise:457869737473556e697175652e756e69717565}{Prerequisite~P\ref{premise:457869737473556e697175652e756e69717565}}, we obtain \(v = w\).
\end{proof}

\begin{lemma}[{Equals}]\label{theorem:6}
Let \(V\) be a type. Let \(G\) be a simple graph on \(V\). Let \(M\) be a subgraph of \(G\). Let \(u\), \(v\), and \(w\) be elements of \(V\). Assume \(M\text{ is a matching}\). Assume \(u \sim_{M} w\). Assume \(v \sim_{M} w\). Then \(u = v\).
\end{lemma}

\begin{proof}
Use the hypothesis \(u \sim_{M} w\).
By \paperlink{premise:53696d706c6547726170682e53756267726170682e41646a2e73796d6d}{Prerequisite~P\ref{premise:53696d706c6547726170682e53756267726170682e41646a2e73796d6d}}, we obtain \(w \sim_{M} u\).
Use the hypothesis \(v \sim_{M} w\).
By \paperlink{premise:53696d706c6547726170682e53756267726170682e41646a2e73796d6d}{Prerequisite~P\ref{premise:53696d706c6547726170682e53756267726170682e41646a2e73796d6d}}, we obtain \(w \sim_{M} v\).
Apply \paperlink{theorem:5}{Lemma~\ref{theorem:5}}.
\end{proof}

\begin{theorem}[{Not adj left}]\label{theorem:7}
Let \(V\) be a type. Let \(G\) be a simple graph on \(V\). Let \(M\) be a subgraph of \(G\). Let \(u\), \(v\), and \(w\) be elements of \(V\). Assume \(M\text{ is a matching}\). Assume \(v \ne w\). Assume \(u \sim_{M} v\). Then \(\neg \left(u \sim_{M} w\right)\).
\end{theorem}

\begin{proof}
Assume \(u \sim_{M} w\).
By \paperlink{theorem:5}{Lemma~\ref{theorem:5}}, we obtain \(v = w\).
Applying the hypothesis \(v \ne w\), we obtain \(\mathrm{False}\).
\end{proof}

\begin{theorem}[{Not adj right}]\label{theorem:8}
Let \(V\) be a type. Let \(G\) be a simple graph on \(V\). Let \(M\) be a subgraph of \(G\). Let \(u\), \(v\), and \(w\) be elements of \(V\). Assume \(M\text{ is a matching}\). Assume \(u \ne v\). Assume \(u \sim_{M} w\). Then \(\neg \left(v \sim_{M} w\right)\).
\end{theorem}

\begin{proof}
Assume \(v \sim_{M} w\).
By \paperlink{theorem:6}{Lemma~\ref{theorem:6}}, we obtain \(u = v\).
Applying the hypothesis \(u \ne v\), we obtain \(\mathrm{False}\).
\end{proof}

\begin{theorem}[{Sup}]\label{theorem:9}
Let \(V\) be a type. Let \(G\) be a simple graph on \(V\). Let \(M\) be a subgraph of \(G\). Let \(M^{'}\) be a subgraph of \(G\). Assume \(M\text{ is a matching}\). Assume \(M^{'}\text{ is a matching}\). Assume \(\operatorname{Disjoint}(\operatorname{Simple\,Graph\mathord{:}Subgraph\mathord{:}support}(M),\allowbreak{} \operatorname{Simple\,Graph\mathord{:}Subgraph\mathord{:}support}(M^{'}))\). Then \(\operatorname{max}\left(M, M^{'}\right)\text{ is a matching}\).
\end{theorem}

\begin{proof}
Fix \(v\) in \(V\).
Assume \(v \in V(\operatorname{max}\left(M, M^{'}\right))\).
Fix \(N\) in \(\operatorname{Subgraph}\left(G\right)\).
Fix \(N^{'}\) in \(\operatorname{Subgraph}\left(G\right)\).
Assume \(N\text{ is a matching}\).
Assume \(\operatorname{Disjoint}(\operatorname{Simple\,Graph\mathord{:}Subgraph\mathord{:}support}(N),\allowbreak{} \operatorname{Simple\,Graph\mathord{:}Subgraph\mathord{:}support}(N^{'}))\).
Assume \(v \in V(N)\).
Use the hypothesis \(v \in V(N)\).
Applying the hypothesis \(N\text{ is a matching}\), we obtain \(\exists! w\, :\, V,\; v \sim_{N} w\).
Consider the following case.
\begin{proofstep}
Fix \(w\) in \(V\).
Assume \(v \sim_{N} w \land \left(\forall y\, :\, V,\; v \sim_{N} y \implies y = w\right)\).
Choose the witness \(w\).
By \paperlink{premise:53696d706c6547726170682e53756267726170682e7375705f61646a}{Prerequisite~P\ref{premise:53696d706c6547726170682e53756267726170682e7375705f61646a}}, we obtain \(v \sim_{\operatorname{max}\left(N, N^{'}\right)} w \iff v \sim_{N} w \lor v \sim_{N^{'}} w\).
Use the hypothesis \(v \sim_{N} w \land \left(\forall y\, :\, V,\; v \sim_{N} y \implies y = w\right)\).
Taking the first conjunct, we obtain \(v \sim_{N} w\).
The first alternative suffices. Hence \(v \sim_{N} w \lor v \sim_{N^{'}} w\).
By the reverse implication, we obtain \(v \sim_{\operatorname{max}\left(N, N^{'}\right)} w\).
For the required implication, argue as follows.
\begin{proofstep}
Fix \(y\) in \(V\).
Assume \(v \sim_{\operatorname{max}\left(N, N^{'}\right)} y\).
Use the hypothesis \(v \sim_{\operatorname{max}\left(N, N^{'}\right)} y\).
Consider the following case.
\begin{proofstep}
Assume \(v \sim_{N} y\).
Assume \(x \mapsto y_{1} \mapsto x = y_{1}\).
Use the hypothesis \(v \sim_{N} w \land \left(\forall y\, :\, V,\; v \sim_{N} y \implies y = w\right)\).
Use the hypothesis \(v \sim_{N} y\).
Taking the second conjunct, we obtain \(y = w\).
\end{proofstep}
Consider the following case.
\begin{proofstep}
Assume \(v \sim_{N^{'}} y\).
Assume \(x \mapsto y_{1} \mapsto x = y_{1}\).
The substitution uses the following argument.
\begin{proofstep}
Choose the witness \(w\).
Use the hypothesis \(v \sim_{N} w \land \left(\forall y\, :\, V,\; v \sim_{N} y \implies y = w\right)\).
Taking the first conjunct, we obtain \(v \sim_{N} w\).
This witness proves existence. Hence \(\exists w_{1}\, :\, V,\; v \sim_{N} w_{1}\).
\end{proofstep}
The substitution uses the following argument.
\begin{proofstep}
Choose the witness \(y\).
Use the hypothesis \(v \sim_{N^{'}} y\).
This witness proves existence. Hence \(\exists w_{1}\, :\, V,\; v \sim_{N^{'}} w_{1}\).
\end{proofstep}
By \paperlink{premise:5365742e6469736a6f696e745f6c656674}{Prerequisite~P\ref{premise:5365742e6469736a6f696e745f6c656674}}, we obtain \(\operatorname{Disjoint}(\operatorname{Simple\,Graph\mathord{:}Subgraph\mathord{:}support}(N),\allowbreak{} \operatorname{Simple\,Graph\mathord{:}Subgraph\mathord{:}support}(N^{'})) \iff (\forall a\, :\, V,\; a \in \operatorname{Simple\,Graph\mathord{:}Subgraph\mathord{:}support}(N) \implies a \notin \operatorname{Simple\,Graph\mathord{:}Subgraph\mathord{:}support}(N^{'}))\).
By propositional extensionality, we obtain \(\operatorname{Disjoint}(\operatorname{Simple\,Graph\mathord{:}Subgraph\mathord{:}support}(N),\allowbreak{} \operatorname{Simple\,Graph\mathord{:}Subgraph\mathord{:}support}(N^{'})) = (\forall a\, :\, V,\; a \in \operatorname{Simple\,Graph\mathord{:}Subgraph\mathord{:}support}(N) \implies a \notin \operatorname{Simple\,Graph\mathord{:}Subgraph\mathord{:}support}(N^{'}))\).
Applying the same function to equal arguments, we obtain \(\operatorname{Disjoint}(\operatorname{Simple\,Graph\mathord{:}Subgraph\mathord{:}support}(N),\allowbreak{} \operatorname{Simple\,Graph\mathord{:}Subgraph\mathord{:}support}(N^{'})) = (\forall a\, :\, V,\; a \in \operatorname{Simple\,Graph\mathord{:}Subgraph\mathord{:}support}(N) \implies a \notin \operatorname{Simple\,Graph\mathord{:}Subgraph\mathord{:}support}(N^{'}))\).
Use the hypothesis \(\operatorname{Disjoint}(\operatorname{Simple\,Graph\mathord{:}Subgraph\mathord{:}support}(N),\allowbreak{} \operatorname{Simple\,Graph\mathord{:}Subgraph\mathord{:}support}(N^{'}))\).
Substituting along the equality, we obtain \(v \in \operatorname{Simple\,Graph\mathord{:}Subgraph\mathord{:}support}(N) \implies v \notin \operatorname{Simple\,Graph\mathord{:}Subgraph\mathord{:}support}(N^{'})\).
Rewriting with \paperlink{premise:696d706c6965735f636f6e6772}{Prerequisite~P\ref{premise:696d706c6965735f636f6e6772}}, \paperlink{premise:4966662e6d7072}{Prerequisite~P\ref{premise:4966662e6d7072}}, \paperlink{premise:53696d706c6547726170682e53756267726170682e6d656d5f737570706f7274}{Prerequisite~P\ref{premise:53696d706c6547726170682e53756267726170682e6d656d5f737570706f7274}}, \paperlink{premise:6e6f745f747275655f65715f66616c7365}{Prerequisite~P\ref{premise:6e6f745f747275655f65715f66616c7365}}, \paperlink{premise:696d705f66616c7365}{Prerequisite~P\ref{premise:696d705f66616c7365}} and the two auxiliary results established above in the preceding result gives \(\mathrm{False}\).
The contradiction rules out this case. Hence \(y = w\).
\end{proofstep}
Use the hypothesis \(v \sim_{\operatorname{max}\left(N, N^{'}\right)} y\).
By definition, we obtain \(x \mapsto y_{1} \mapsto x = y_{1}\).
The two alternatives exhaust the cases. Hence \(y = w\).
\end{proofstep}
Combining the two conclusions, we obtain \(v \sim_{\operatorname{max}(N,\allowbreak{} N^{'})} w \land (\forall y\, :\, V,\; v \sim_{\operatorname{max}(N,\allowbreak{} N^{'})} y \implies y = w)\).
This witness proves existence. Hence \(\exists x\, :\, V,\; v \sim_{\operatorname{max}(N,\allowbreak{} N^{'})} x \land (\forall y\, :\, V,\; v \sim_{\operatorname{max}(N,\allowbreak{} N^{'})} y \implies y = x)\).
\end{proofstep}
Using the existential hypothesis, we obtain \(\exists! w\, :\, V,\; v \sim_{\operatorname{max}\left(N, N^{'}\right)} w\).
Use the hypothesis \(v \in V(\operatorname{max}\left(M, M^{'}\right))\).
By \paperlink{premise:5365742e6d656d5f6f725f6d656d5f6f665f6d656d5f756e696f6e}{Prerequisite~P\ref{premise:5365742e6d656d5f6f725f6d656d5f6f665f6d656d5f756e696f6e}}, we obtain \(v \in V(M) \lor v \in V(M^{'})\).
Consider the following case.
\begin{proofstep}
Assume \(v \in V(M)\).
Assume \(x \mapsto y \mapsto x = y\).
Use the hypothesis \(M\text{ is a matching}\).
Use the hypothesis \(\operatorname{Disjoint}(\operatorname{Simple\,Graph\mathord{:}Subgraph\mathord{:}support}(M),\allowbreak{} \operatorname{Simple\,Graph\mathord{:}Subgraph\mathord{:}support}(M^{'}))\).
Use the hypothesis \(v \in V(M)\).
Applying the hypothesis \(\forall N\, :\, \operatorname{Subgraph}(G),\; \forall N^{'}\, :\, \operatorname{Subgraph}(G),\; N\text{ is a matching} \implies \operatorname{Disjoint}(\operatorname{Simple\,Graph\mathord{:}Subgraph\mathord{:}support}(N),\allowbreak{} \operatorname{Simple\,Graph\mathord{:}Subgraph\mathord{:}support}(N^{'})) \implies v \in V(N) \implies (\exists! w\, :\, V,\; v \sim_{\operatorname{max}(N,\allowbreak{} N^{'})} w)\), we obtain \(\exists! w\, :\, V,\; v \sim_{\operatorname{max}\left(M, M^{'}\right)} w\).
\end{proofstep}
Consider the following case.
\begin{proofstep}
Assume \(v \in V(M^{'})\).
Assume \(x \mapsto y \mapsto x = y\).
Use the hypothesis \(M^{'}\text{ is a matching}\).
Use the hypothesis \(\operatorname{Disjoint}(\operatorname{Simple\,Graph\mathord{:}Subgraph\mathord{:}support}(M),\allowbreak{} \operatorname{Simple\,Graph\mathord{:}Subgraph\mathord{:}support}(M^{'}))\).
By \paperlink{premise:4469736a6f696e742e73796d6d}{Prerequisite~P\ref{premise:4469736a6f696e742e73796d6d}} with \(\alpha = \mathcal{P}(V)\), we obtain \(\operatorname{Disjoint}(\operatorname{Simple\,Graph\mathord{:}Subgraph\mathord{:}support}(M^{'}),\allowbreak{} \operatorname{Simple\,Graph\mathord{:}Subgraph\mathord{:}support}(M))\).
Use the hypothesis \(v \in V(M^{'})\).
Applying the hypothesis \(\forall N\, :\, \operatorname{Subgraph}(G),\; \forall N^{'}\, :\, \operatorname{Subgraph}(G),\; N\text{ is a matching} \implies \operatorname{Disjoint}(\operatorname{Simple\,Graph\mathord{:}Subgraph\mathord{:}support}(N),\allowbreak{} \operatorname{Simple\,Graph\mathord{:}Subgraph\mathord{:}support}(N^{'})) \implies v \in V(N) \implies (\exists! w\, :\, V,\; v \sim_{\operatorname{max}(N,\allowbreak{} N^{'})} w)\), we obtain \(\exists! w\, :\, V,\; v \sim_{\operatorname{max}\left(M^{'}, M\right)} w\).
Rewriting with \paperlink{premise:7375705f636f6d6d}{Prerequisite~P\ref{premise:7375705f636f6d6d}} in the preceding result gives \(\exists! w\, :\, V,\; v \sim_{\operatorname{max}\left(M, M^{'}\right)} w\).
\end{proofstep}
Use the hypothesis \(v \in V(\operatorname{max}\left(M, M^{'}\right))\).
By \paperlink{premise:5365742e6d656d5f6f725f6d656d5f6f665f6d656d5f756e696f6e}{Prerequisite~P\ref{premise:5365742e6d656d5f6f725f6d656d5f6f665f6d656d5f756e696f6e}}, we obtain \(v \in V(M) \lor v \in V(M^{'})\).
By definition, we obtain \(x \mapsto y \mapsto x = y\).
The two alternatives exhaust the cases. Hence \(\exists! w\, :\, V,\; v \sim_{\operatorname{max}\left(M, M^{'}\right)} w\).
\end{proof}

\begin{theorem}[{I sup}]\label{theorem:10}
Let \(V\) be a type. Let \(G\) be a simple graph on \(V\). Let \(\iota\) be a type. Let \(f\) be a function of type \(\iota \to \operatorname{Subgraph}\left(G\right)\). Assume \(\forall i\, :\, \iota,\; f\left(i\right)\text{ is a matching}\). Assume \(\operatorname{Pairwise}(i \mapsto j \mapsto \operatorname{Disjoint}(\operatorname{Simple\,Graph\mathord{:}Subgraph\mathord{:}support}(f(i)),\allowbreak{} \operatorname{Simple\,Graph\mathord{:}Subgraph\mathord{:}support}(f(j))))\). Then \(\operatorname{i\,Sup}\left(i \mapsto f\left(i\right)\right)\text{ is a matching}\).
\end{theorem}

\begin{proof}
Fix \(v\) in \(V\).
Assume \(v \in V(\operatorname{i\,Sup}\left(i \mapsto f\left(i\right)\right))\).
By \paperlink{premise:5365742e6d656d5f69556e696f6e}{Prerequisite~P\ref{premise:5365742e6d656d5f69556e696f6e}}, we obtain \(v \in \bigcup_{i\, :\, \iota} V(f(i)) \iff (\exists i\, :\, \iota,\; v \in V(f(i)))\).
Use the hypothesis \(v \in V(\operatorname{i\,Sup}\left(i \mapsto f\left(i\right)\right))\).
By \paperlink{premise:53696d706c6547726170682e53756267726170682e76657274735f69537570}{Prerequisite~P\ref{premise:53696d706c6547726170682e53756267726170682e76657274735f69537570}}, we obtain \(V(\operatorname{i\,Sup}\left(i \mapsto f\left(i\right)\right)) = \bigcup_{i\, :\, \iota} V(f\left(i\right))\).
Substituting equals, we obtain \(v \in \bigcup_{i\, :\, \iota} V(f\left(i\right))\).
By the forward implication, we obtain \(\exists i\, :\, \iota,\; v \in V(f\left(i\right))\).
Consider the following case.
\begin{proofstep}
Fix \(i\) in \(\iota\).
Assume \(v \in V(f\left(i\right))\).
Use the hypothesis \(v \in V(f\left(i\right))\).
Applying the hypothesis \(\forall i\, :\, \iota,\; f\left(i\right)\text{ is a matching}\), we obtain \(\exists! w\, :\, V,\; v \sim_{f\left(i\right)} w\).
Consider the following case.
\begin{proofstep}
Fix \(w\) in \(V\).
Assume \(v \sim_{f\left(i\right)} w \land \left(\forall y\, :\, V,\; v \sim_{f\left(i\right)} y \implies y = w\right)\).
Choose the witness \(w\).
By \paperlink{premise:53696d706c6547726170682e53756267726170682e695375705f61646a}{Prerequisite~P\ref{premise:53696d706c6547726170682e53756267726170682e695375705f61646a}}, we obtain \(v \sim_{\operatorname{i\,Sup}(i_{1} \mapsto f(i_{1}))} w \iff (\exists i_{1}\, :\, \iota,\; v \sim_{f(i_{1})} w)\).
Choose the witness \(i\).
Use the hypothesis \(v \sim_{f\left(i\right)} w \land \left(\forall y\, :\, V,\; v \sim_{f\left(i\right)} y \implies y = w\right)\).
Taking the first conjunct, we obtain \(v \sim_{f\left(i\right)} w\).
This witness proves existence. Hence \(\exists i_{1}\, :\, \iota,\; v \sim_{f\left(i_{1}\right)} w\).
By the reverse implication, we obtain \(v \sim_{\operatorname{i\,Sup}\left(i_{1} \mapsto f\left(i_{1}\right)\right)} w\).
For the required implication, argue as follows.
\begin{proofstep}
Fix \(y\) in \(V\).
Assume \(v \sim_{\operatorname{i\,Sup}\left(i_{1} \mapsto f\left(i_{1}\right)\right)} y\).
By \paperlink{premise:53696d706c6547726170682e53756267726170682e695375705f61646a}{Prerequisite~P\ref{premise:53696d706c6547726170682e53756267726170682e695375705f61646a}}, we obtain \(v \sim_{\operatorname{i\,Sup}(i_{1} \mapsto f(i_{1}))} y \iff (\exists i_{1}\, :\, \iota,\; v \sim_{f(i_{1})} y)\).
Use the hypothesis \(v \sim_{\operatorname{i\,Sup}\left(i_{1} \mapsto f\left(i_{1}\right)\right)} y\).
By the forward implication, we obtain \(\exists i_{1}\, :\, \iota,\; v \sim_{f\left(i_{1}\right)} y\).
Consider the following case.
\begin{proofstep}
Fix \(i^{'}\) in \(\iota\).
Assume \(v \sim_{f\left(i^{'}\right)} y\).
Distinguish whether \(i = i^{'}\) holds.
Consider the following case.
\begin{proofstep}
Assume \(i = i^{'}\).
Use the hypothesis \(v \sim_{f\left(i\right)} w \land \left(\forall y\, :\, V,\; v \sim_{f\left(i\right)} y \implies y = w\right)\).
Use the hypothesis \(v \sim_{f\left(i^{'}\right)} y\).
Use the hypothesis \(i = i^{'}\).
Reversing the equality, we obtain \(i^{'} = i\).
Substituting equals, we obtain \(v \sim_{f\left(i\right)} y\).
Taking the second conjunct, we obtain \(y = w\).
\end{proofstep}
Consider the following case.
\begin{proofstep}
Assume \(\neg \left(i = i^{'}\right)\).
Use the hypothesis \(\neg \left(i = i^{'}\right)\).
Applying the hypothesis \(\operatorname{Pairwise}(i \mapsto j \mapsto \operatorname{Disjoint}(\operatorname{Simple\,Graph\mathord{:}Subgraph\mathord{:}support}(f(i)),\allowbreak{} \operatorname{Simple\,Graph\mathord{:}Subgraph\mathord{:}support}(f(j))))\), we obtain \(\operatorname{Disjoint}(\operatorname{Simple\,Graph\mathord{:}Subgraph\mathord{:}support}(f(i)),\allowbreak{} \operatorname{Simple\,Graph\mathord{:}Subgraph\mathord{:}support}(f(i^{'})))\).
The substitution uses the following argument.
\begin{proofstep}
Choose the witness \(w\).
Use the hypothesis \(v \sim_{f\left(i\right)} w \land \left(\forall y\, :\, V,\; v \sim_{f\left(i\right)} y \implies y = w\right)\).
Taking the first conjunct, we obtain \(v \sim_{f\left(i\right)} w\).
This witness proves existence. Hence \(\exists w_{1}\, :\, V,\; v \sim_{f\left(i\right)} w_{1}\).
\end{proofstep}
The substitution uses the following argument.
\begin{proofstep}
Choose the witness \(y\).
Use the hypothesis \(v \sim_{f\left(i^{'}\right)} y\).
This witness proves existence. Hence \(\exists w_{1}\, :\, V,\; v \sim_{f\left(i^{'}\right)} w_{1}\).
\end{proofstep}
By \paperlink{premise:5365742e6469736a6f696e745f6c656674}{Prerequisite~P\ref{premise:5365742e6469736a6f696e745f6c656674}}, we obtain \(\operatorname{Disjoint}(\operatorname{Simple\,Graph\mathord{:}Subgraph\mathord{:}support}(f(i)),\allowbreak{} \operatorname{Simple\,Graph\mathord{:}Subgraph\mathord{:}support}(f(i^{'}))) \iff (\forall a\, :\, V,\; a \in \operatorname{Simple\,Graph\mathord{:}Subgraph\mathord{:}support}(f(i)) \implies a \notin \operatorname{Simple\,Graph\mathord{:}Subgraph\mathord{:}support}(f(i^{'})))\).
By propositional extensionality, we obtain \(\operatorname{Disjoint}(\operatorname{Simple\,Graph\mathord{:}Subgraph\mathord{:}support}(f(i)),\allowbreak{} \operatorname{Simple\,Graph\mathord{:}Subgraph\mathord{:}support}(f(i^{'}))) = (\forall a\, :\, V,\; a \in \operatorname{Simple\,Graph\mathord{:}Subgraph\mathord{:}support}(f(i)) \implies a \notin \operatorname{Simple\,Graph\mathord{:}Subgraph\mathord{:}support}(f(i^{'})))\).
Use the previously proved fact \(\operatorname{Disjoint}(\operatorname{Simple\,Graph\mathord{:}Subgraph\mathord{:}support}(f(i)),\allowbreak{} \operatorname{Simple\,Graph\mathord{:}Subgraph\mathord{:}support}(f(i^{'})))\).
Substituting along the equality, we obtain \(v \in \operatorname{Simple\,Graph\mathord{:}Subgraph\mathord{:}support}(f(i)) \implies v \notin \operatorname{Simple\,Graph\mathord{:}Subgraph\mathord{:}support}(f(i^{'}))\).
Rewriting with \paperlink{premise:696d706c6965735f636f6e6772}{Prerequisite~P\ref{premise:696d706c6965735f636f6e6772}}, \paperlink{premise:4966662e6d7072}{Prerequisite~P\ref{premise:4966662e6d7072}}, \paperlink{premise:53696d706c6547726170682e53756267726170682e6d656d5f737570706f7274}{Prerequisite~P\ref{premise:53696d706c6547726170682e53756267726170682e6d656d5f737570706f7274}}, \paperlink{premise:6e6f745f747275655f65715f66616c7365}{Prerequisite~P\ref{premise:6e6f745f747275655f65715f66616c7365}}, \paperlink{premise:696d705f66616c7365}{Prerequisite~P\ref{premise:696d705f66616c7365}} and the two auxiliary results established above in the preceding result gives \(\mathrm{False}\).
The contradiction rules out this case. Hence \(y = w\).
\end{proofstep}
The proposition or its negation holds. Hence \(y = w\).
\end{proofstep}
Using the existential hypothesis, we obtain \(y = w\).
\end{proofstep}
Combining the two conclusions, we obtain \(v \sim_{\operatorname{i\,Sup}(i_{1} \mapsto f(i_{1}))} w \land (\forall y\, :\, V,\; v \sim_{\operatorname{i\,Sup}(i_{1} \mapsto f(i_{1}))} y \implies y = w)\).
This witness proves existence. Hence \(\exists x\, :\, V,\; v \sim_{\operatorname{i\,Sup}(i_{1} \mapsto f(i_{1}))} x \land (\forall y\, :\, V,\; v \sim_{\operatorname{i\,Sup}(i_{1} \mapsto f(i_{1}))} y \implies y = x)\).
\end{proofstep}
Using the existential hypothesis, we obtain \(\exists! w\, :\, V,\; v \sim_{\operatorname{i\,Sup}\left(i_{1} \mapsto f\left(i_{1}\right)\right)} w\).
\end{proofstep}
Using the existential hypothesis, we obtain \(\exists! w\, :\, V,\; v \sim_{\operatorname{i\,Sup}\left(i \mapsto f\left(i\right)\right)} w\).
\end{proof}

\begin{theorem}[{Subgraph of adj}]\label{theorem:11}
Let \(V\) be a type. Let \(G\) be a simple graph on \(V\). Let \(v\) and \(w\) be elements of \(V\). Let \(h\) be a proof of \(v \sim_{G} w\). Then \(\operatorname{subgraph\,Of\,Adj}\left(G\right)\text{ is a matching}\).
\end{theorem}

\begin{proof}
Fix \(v_{1}\) in \(V\).
Assume \(v_{1} \in V(\operatorname{subgraph\,Of\,Adj}\left(G\right))\).
Use the hypothesis \(v_{1} \in V(\operatorname{subgraph\,Of\,Adj}\left(G\right))\).
Rewriting with \paperlink{premise:53696d706c6547726170682e73756267726170684f6641646a5f7665727473}{Prerequisite~P\ref{premise:53696d706c6547726170682e73756267726170684f6641646a5f7665727473}} (using \(v \sim_{G} w\)) in the preceding result gives \(v_{1} \in \{v\} \cup \{w\}\).
Rewriting with \paperlink{premise:5365742e6d656d5f696e736572745f696666}{Prerequisite~P\ref{premise:5365742e6d656d5f696e736572745f696666}} in the preceding result gives \(v_{1} = v \lor v_{1} \in \{w\}\).
Rewriting with \paperlink{premise:5365742e6d656d5f73696e676c65746f6e5f696666}{Prerequisite~P\ref{premise:5365742e6d656d5f73696e676c65746f6e5f696666}} in the preceding result gives \(v_{1} = v \lor v_{1} = w\).
Consider the following case.
\begin{proofstep}
Assume \(v_{1} = v\).
Assume \(x \mapsto y \mapsto x = y\).
Choose the witness \(w\).
For the required implication, argue as follows.
\begin{proofstep}
Assume \(v_{1} \sim_{G} w\).
The substitution uses the following argument.
\begin{proofstep}
Fix \(y\) in \(V\).
By \paperlink{premise:53796d322e6571}{Prerequisite~P\ref{premise:53796d322e6571}}, we obtain \(\operatorname{Sym2\mathord{:}mk}((v_{1},\allowbreak{} w)) = \operatorname{Sym2\mathord{:}mk}((v_{1},\allowbreak{} y)) \iff \operatorname{Rel}(V,\allowbreak{} (v_{1},\allowbreak{} w),\allowbreak{} (v_{1},\allowbreak{} y))\).
By propositional extensionality, we obtain \((\operatorname{Sym2\mathord{:}mk}((v_{1},\allowbreak{} w)) = \operatorname{Sym2\mathord{:}mk}((v_{1},\allowbreak{} y))) = \operatorname{Rel}(V,\allowbreak{} (v_{1},\allowbreak{} w),\allowbreak{} (v_{1},\allowbreak{} y))\).
By \paperlink{premise:53796d322e72656c5f69666627}{Prerequisite~P\ref{premise:53796d322e72656c5f69666627}}, we obtain \(\operatorname{Rel}(V,\allowbreak{} (v_{1},\allowbreak{} w),\allowbreak{} (v_{1},\allowbreak{} y)) \iff (v_{1},\allowbreak{} w) = (v_{1},\allowbreak{} y) \lor (v_{1},\allowbreak{} w) = \operatorname{swap}((v_{1},\allowbreak{} y))\).
By propositional extensionality, we obtain \(\operatorname{Rel}(V,\allowbreak{} (v_{1},\allowbreak{} w),\allowbreak{} (v_{1},\allowbreak{} y)) = ((v_{1},\allowbreak{} w) = (v_{1},\allowbreak{} y) \lor (v_{1},\allowbreak{} w) = \operatorname{swap}((v_{1},\allowbreak{} y)))\).
By \paperlink{premise:50726f642e6d6b2e696e6a4571}{Prerequisite~P\ref{premise:50726f642e6d6b2e696e6a4571}} with \(\mathit{fst} = v_{1}\) and \(\mathit{snd} = w\) and \(\mathit{fst}_{1} = v_{1}\) and \(\mathit{snd}_{1} = y\), we obtain \(\left(\left(v_{1}, w\right) = \left(v_{1}, y\right)\right) = \left(v_{1} = v_{1} \land w = y\right)\).
By \paperlink{premise:65715f73656c66}{Prerequisite~P\ref{premise:65715f73656c66}} with \(a = v_{1}\), we obtain \(\left(v_{1} = v_{1}\right) = \mathrm{True}\).
Applying the same function to equal arguments, we obtain \(\left(v_{1} = v_{1} \land w = y\right) = \left(\mathrm{True} \land w = y\right)\).
By \paperlink{premise:747275655f616e64}{Prerequisite~P\ref{premise:747275655f616e64}} with \(p = w = y\), we obtain \(\left(\mathrm{True} \land w = y\right) = \left(w = y\right)\).
Chaining the equalities, we obtain \(\left(v_{1} = v_{1} \land w = y\right) = \left(w = y\right)\).
Chaining the equalities, we obtain \(\left(\left(v_{1}, w\right) = \left(v_{1}, y\right)\right) = \left(w = y\right)\).
Applying the same function to equal arguments, we obtain \((y_{1} \mapsto (v_{1},\allowbreak{} w) = (v_{1},\allowbreak{} y) \lor y_{1}) = (y_{1} \mapsto w = y \lor y_{1})\).
By \paperlink{premise:50726f642e6d6b2e696e6a4571}{Prerequisite~P\ref{premise:50726f642e6d6b2e696e6a4571}} with \(\mathit{fst} = v_{1}\) and \(\mathit{snd} = w\) and \(\mathit{fst}_{1} = y\) and \(\mathit{snd}_{1} = v_{1}\), we obtain \(\left(\left(v_{1}, w\right) = \left(y, v_{1}\right)\right) = \left(v_{1} = y \land w = v_{1}\right)\).
By congruence, we obtain \(((v_{1},\allowbreak{} w) = (v_{1},\allowbreak{} y) \lor (v_{1},\allowbreak{} w) = (y,\allowbreak{} v_{1})) = (w = y \lor v_{1} = y \land w = v_{1})\).
Chaining the equalities, we obtain \(\operatorname{Rel}(V,\allowbreak{} (v_{1},\allowbreak{} w),\allowbreak{} (v_{1},\allowbreak{} y)) = (w = y \lor v_{1} = y \land w = v_{1})\).
Chaining the equalities, we obtain \((\operatorname{Sym2\mathord{:}mk}((v_{1},\allowbreak{} w)) = \operatorname{Sym2\mathord{:}mk}((v_{1},\allowbreak{} y))) = (w = y \lor v_{1} = y \land w = v_{1})\).
By definition, we obtain \(\left(y = w\right) = \left(y = w\right)\).
By \paperlink{premise:696d706c6965735f636f6e6772}{Prerequisite~P\ref{premise:696d706c6965735f636f6e6772}} with \(p_{1} = \operatorname{Sym2\mathord{:}mk}((v_{1},\allowbreak{} w)) = \operatorname{Sym2\mathord{:}mk}((v_{1},\allowbreak{} y))\) and \(p_{2} = w = y \lor v_{1} = y \land w = v_{1}\) and \(q_{1} = y = w\) and \(q_{2} = y = w\), we obtain \((\operatorname{Sym2\mathord{:}mk}((v_{1},\allowbreak{} w)) = \operatorname{Sym2\mathord{:}mk}((v_{1},\allowbreak{} y)) \implies y = w) = (w = y \lor v_{1} = y \land w = v_{1} \implies y = w)\).
\end{proofstep}
Fix \(y\) in \(V\).
Assume \(w = y \lor v_{1} = y \land w = v_{1}\).
Use the hypothesis \(w = y \lor v_{1} = y \land w = v_{1}\).
Consider the following case.
\begin{proofstep}
Assume \(w = y\).
Simplifying by definition gives \(w = w\).
Use the hypothesis \(w = y\).
Substituting equals, we obtain \(y = w\).
\end{proofstep}
Consider the following case.
\begin{proofstep}
Assume \(v_{1} = y \land w = v_{1}\).
Use the hypothesis \(v_{1} = y \land w = v_{1}\).
Taking the first conjunct gives \(v_{1} = y\).
Use the hypothesis \(v_{1} = y \land w = v_{1}\).
Taking the second conjunct gives \(w = v_{1}\).
Use the hypothesis \(v_{1} = y \land w = v_{1}\).
Taking the first conjunct gives \(v_{1} = y\).
Use the hypothesis \(v_{1} \sim_{G} w\).
Rewriting with \(v_{1} = y\), \(w = v_{1}\), \paperlink{premise:65715f66616c7365}{Prerequisite~P\ref{premise:65715f66616c7365}} and \paperlink{premise:53696d706c6547726170682e69727265666c}{Prerequisite~P\ref{premise:53696d706c6547726170682e69727265666c}} in the preceding result gives \(\mathrm{False}\).
The contradiction rules out this case. Hence \(y = w\).
\end{proofstep}
The two alternatives exhaust the cases. Hence \(y = w\).
Rewriting with \paperlink{premise:666f72616c6c5f636f6e6772}{Prerequisite~P\ref{premise:666f72616c6c5f636f6e6772}}, \paperlink{premise:747275655f616e64}{Prerequisite~P\ref{premise:747275655f616e64}} and the identity established above in the preceding result gives \(v_{1} \sim_{\operatorname{subgraph\,Of\,Adj}(G)} w \land (\forall y\, :\, V,\; v_{1} \sim_{\operatorname{subgraph\,Of\,Adj}(G)} y \implies y = w)\).
\end{proofstep}
Use the hypothesis \(v_{1} = v\).
Use the hypothesis \(v \sim_{G} w\).
Substituting equals, we obtain \(v_{1} \sim_{\operatorname{subgraph\,Of\,Adj}(G)} w \land (\forall y\, :\, V,\; v_{1} \sim_{\operatorname{subgraph\,Of\,Adj}(G)} y \implies y = w)\).
This witness proves existence. Hence \(\exists x\, :\, V,\; v_{1} \sim_{\operatorname{subgraph\,Of\,Adj}(G)} x \land (\forall y\, :\, V,\; v_{1} \sim_{\operatorname{subgraph\,Of\,Adj}(G)} y \implies y = x)\).
\end{proofstep}
Consider the following case.
\begin{proofstep}
Assume \(v_{1} = w\).
Assume \(x \mapsto y \mapsto x = y\).
Choose the witness \(v\).
For the required implication, argue as follows.
\begin{proofstep}
Assume \(v \sim_{G} v_{1}\).
The substitution uses the following argument.
\begin{proofstep}
Fix \(y\) in \(V\).
By \paperlink{premise:53796d322e6571}{Prerequisite~P\ref{premise:53796d322e6571}}, we obtain \(\operatorname{Sym2\mathord{:}mk}((v,\allowbreak{} v_{1})) = \operatorname{Sym2\mathord{:}mk}((v_{1},\allowbreak{} y)) \iff \operatorname{Rel}(V,\allowbreak{} (v,\allowbreak{} v_{1}),\allowbreak{} (v_{1},\allowbreak{} y))\).
By propositional extensionality, we obtain \((\operatorname{Sym2\mathord{:}mk}((v,\allowbreak{} v_{1})) = \operatorname{Sym2\mathord{:}mk}((v_{1},\allowbreak{} y))) = \operatorname{Rel}(V,\allowbreak{} (v,\allowbreak{} v_{1}),\allowbreak{} (v_{1},\allowbreak{} y))\).
By \paperlink{premise:53796d322e72656c5f69666627}{Prerequisite~P\ref{premise:53796d322e72656c5f69666627}}, we obtain \(\operatorname{Rel}(V,\allowbreak{} (v,\allowbreak{} v_{1}),\allowbreak{} (v_{1},\allowbreak{} y)) \iff (v,\allowbreak{} v_{1}) = (v_{1},\allowbreak{} y) \lor (v,\allowbreak{} v_{1}) = \operatorname{swap}((v_{1},\allowbreak{} y))\).
By propositional extensionality, we obtain \(\operatorname{Rel}(V,\allowbreak{} (v,\allowbreak{} v_{1}),\allowbreak{} (v_{1},\allowbreak{} y)) = ((v,\allowbreak{} v_{1}) = (v_{1},\allowbreak{} y) \lor (v,\allowbreak{} v_{1}) = \operatorname{swap}((v_{1},\allowbreak{} y)))\).
By \paperlink{premise:50726f642e6d6b2e696e6a4571}{Prerequisite~P\ref{premise:50726f642e6d6b2e696e6a4571}} with \(\mathit{fst} = v\) and \(\mathit{snd} = v_{1}\) and \(\mathit{fst}_{1} = v_{1}\) and \(\mathit{snd}_{1} = y\), we obtain \(\left(\left(v, v_{1}\right) = \left(v_{1}, y\right)\right) = \left(v = v_{1} \land v_{1} = y\right)\).
Applying the same function to equal arguments, we obtain \((y_{1} \mapsto (v,\allowbreak{} v_{1}) = (v_{1},\allowbreak{} y) \lor y_{1}) = (y_{1} \mapsto v = v_{1} \land v_{1} = y \lor y_{1})\).
By \paperlink{premise:50726f642e6d6b2e696e6a4571}{Prerequisite~P\ref{premise:50726f642e6d6b2e696e6a4571}} with \(\mathit{fst} = v\) and \(\mathit{snd} = v_{1}\) and \(\mathit{fst}_{1} = y\) and \(\mathit{snd}_{1} = v_{1}\), we obtain \(\left(\left(v, v_{1}\right) = \left(y, v_{1}\right)\right) = \left(v = y \land v_{1} = v_{1}\right)\).
By \paperlink{premise:65715f73656c66}{Prerequisite~P\ref{premise:65715f73656c66}} with \(a = v_{1}\), we obtain \(\left(v_{1} = v_{1}\right) = \mathrm{True}\).
Applying the same function to equal arguments, we obtain \(\left(v = y \land v_{1} = v_{1}\right) = \left(v = y \land \mathrm{True}\right)\).
By \paperlink{premise:616e645f74727565}{Prerequisite~P\ref{premise:616e645f74727565}} with \(p = v = y\), we obtain \(\left(v = y \land \mathrm{True}\right) = \left(v = y\right)\).
Chaining the equalities, we obtain \(\left(v = y \land v_{1} = v_{1}\right) = \left(v = y\right)\).
Chaining the equalities, we obtain \(\left(\left(v, v_{1}\right) = \left(y, v_{1}\right)\right) = \left(v = y\right)\).
By congruence, we obtain \(((v,\allowbreak{} v_{1}) = (v_{1},\allowbreak{} y) \lor (v,\allowbreak{} v_{1}) = (y,\allowbreak{} v_{1})) = (v = v_{1} \land v_{1} = y \lor v = y)\).
Chaining the equalities, we obtain \(\operatorname{Rel}(V,\allowbreak{} (v,\allowbreak{} v_{1}),\allowbreak{} (v_{1},\allowbreak{} y)) = (v = v_{1} \land v_{1} = y \lor v = y)\).
Chaining the equalities, we obtain \((\operatorname{Sym2\mathord{:}mk}((v,\allowbreak{} v_{1})) = \operatorname{Sym2\mathord{:}mk}((v_{1},\allowbreak{} y))) = (v = v_{1} \land v_{1} = y \lor v = y)\).
By definition, we obtain \(\left(y = v\right) = \left(y = v\right)\).
By \paperlink{premise:696d706c6965735f636f6e6772}{Prerequisite~P\ref{premise:696d706c6965735f636f6e6772}} with \(p_{1} = \operatorname{Sym2\mathord{:}mk}((v,\allowbreak{} v_{1})) = \operatorname{Sym2\mathord{:}mk}((v_{1},\allowbreak{} y))\) and \(p_{2} = v = v_{1} \land v_{1} = y \lor v = y\) and \(q_{1} = y = v\) and \(q_{2} = y = v\), we obtain \((\operatorname{Sym2\mathord{:}mk}((v,\allowbreak{} v_{1})) = \operatorname{Sym2\mathord{:}mk}((v_{1},\allowbreak{} y)) \implies y = v) = (v = v_{1} \land v_{1} = y \lor v = y \implies y = v)\).
\end{proofstep}
Fix \(y\) in \(V\).
Assume \(v = v_{1} \land v_{1} = y \lor v = y\).
Use the hypothesis \(v = v_{1} \land v_{1} = y \lor v = y\).
Consider the following case.
\begin{proofstep}
Assume \(v = v_{1} \land v_{1} = y\).
Use the hypothesis \(v = v_{1} \land v_{1} = y\).
Taking the first conjunct gives \(v = v_{1}\).
Use the hypothesis \(v = v_{1} \land v_{1} = y\).
Taking the second conjunct gives \(v_{1} = y\).
Use the hypothesis \(v = v_{1} \land v_{1} = y\).
Taking the second conjunct gives \(v_{1} = y\).
Use the hypothesis \(v \sim_{G} v_{1}\).
Rewriting with \(v = v_{1}\), \(v_{1} = y\), \paperlink{premise:65715f66616c7365}{Prerequisite~P\ref{premise:65715f66616c7365}} and \paperlink{premise:53696d706c6547726170682e69727265666c}{Prerequisite~P\ref{premise:53696d706c6547726170682e69727265666c}} in the preceding result gives \(\mathrm{False}\).
The contradiction rules out this case. Hence \(y = v\).
\end{proofstep}
Consider the following case.
\begin{proofstep}
Assume \(v = y\).
Simplifying by definition gives \(v = v\).
Use the hypothesis \(v = y\).
Substituting equals, we obtain \(y = v\).
\end{proofstep}
The two alternatives exhaust the cases. Hence \(y = v\).
Rewriting with \paperlink{premise:53796d322e6571}{Prerequisite~P\ref{premise:53796d322e6571}}, \paperlink{premise:53796d322e72656c5f69666627}{Prerequisite~P\ref{premise:53796d322e72656c5f69666627}}, \paperlink{premise:50726f642e6d6b2e696e6a4571}{Prerequisite~P\ref{premise:50726f642e6d6b2e696e6a4571}}, \paperlink{premise:6f725f74727565}{Prerequisite~P\ref{premise:6f725f74727565}}, \paperlink{premise:666f72616c6c5f636f6e6772}{Prerequisite~P\ref{premise:666f72616c6c5f636f6e6772}}, \paperlink{premise:747275655f616e64}{Prerequisite~P\ref{premise:747275655f616e64}} and the identity established above in the preceding result gives \(v_{1} \sim_{\operatorname{subgraph\,Of\,Adj}(G)} v \land (\forall y\, :\, V,\; v_{1} \sim_{\operatorname{subgraph\,Of\,Adj}(G)} y \implies y = v)\).
\end{proofstep}
Use the hypothesis \(v_{1} = w\).
Use the hypothesis \(v \sim_{G} w\).
Substituting equals, we obtain \(v_{1} \sim_{\operatorname{subgraph\,Of\,Adj}(G)} v \land (\forall y\, :\, V,\; v_{1} \sim_{\operatorname{subgraph\,Of\,Adj}(G)} y \implies y = v)\).
This witness proves existence. Hence \(\exists x\, :\, V,\; v_{1} \sim_{\operatorname{subgraph\,Of\,Adj}(G)} x \land (\forall y\, :\, V,\; v_{1} \sim_{\operatorname{subgraph\,Of\,Adj}(G)} y \implies y = x)\).
\end{proofstep}
Use the hypothesis \(v_{1} \in V(\operatorname{subgraph\,Of\,Adj}\left(G\right))\).
Rewriting with \paperlink{premise:53696d706c6547726170682e73756267726170684f6641646a5f7665727473}{Prerequisite~P\ref{premise:53696d706c6547726170682e73756267726170684f6641646a5f7665727473}} (using \(v \sim_{G} w\)) in the preceding result gives \(v_{1} \in \{v\} \cup \{w\}\).
Rewriting with \paperlink{premise:5365742e6d656d5f696e736572745f696666}{Prerequisite~P\ref{premise:5365742e6d656d5f696e736572745f696666}} in the preceding result gives \(v_{1} = v \lor v_{1} \in \{w\}\).
Rewriting with \paperlink{premise:5365742e6d656d5f73696e676c65746f6e5f696666}{Prerequisite~P\ref{premise:5365742e6d656d5f73696e676c65746f6e5f696666}} in the preceding result gives \(v_{1} = v \lor v_{1} = w\).
By definition, we obtain \(x \mapsto y \mapsto x = y\).
The two alternatives exhaust the cases. Hence \(\exists! w_{1}\, :\, V,\; v_{1} \sim_{\operatorname{subgraph\,Of\,Adj}\left(G\right)} w_{1}\).
\end{proof}

\begin{theorem}[{Coercion subgraph}]\label{theorem:12}
Let \(V\) be a type. Let \(G\) be a simple graph on \(V\). Let \(G^{'}\) be a subgraph of \(G\). Let \(M\) be a subgraph of \(\operatorname{graph}(G^{'})\). Assume \(M\text{ is a matching}\). Then \(\operatorname{coe\,Subgraph}\left(M\right)\text{ is a matching}\).
\end{theorem}

\begin{proof}
Fix \(v\) in \(V\).
Assume \(v \in V(\operatorname{coe\,Subgraph}\left(M\right))\).
Use the hypothesis \(v \in V(\operatorname{coe\,Subgraph}\left(M\right))\).
Using \paperlink{premise:53696d706c6547726170682e53756267726170682e76657274735f636f655375626772617068}{Prerequisite~P\ref{premise:53696d706c6547726170682e53756267726170682e76657274735f636f655375626772617068}} with \(G^{''} = M\) in reverse gives \(\operatorname{val}\left[V(M)\right] = V(\operatorname{coe\,Subgraph}\left(M\right))\).
Reversing the equality, we obtain \(V(\operatorname{coe\,Subgraph}\left(M\right)) = \operatorname{val}\left[V(M)\right]\).
Substituting equals, we obtain \(v \in \operatorname{val}\left[V(M)\right]\).
By \paperlink{premise:5365742e6d656d5f6f665f6d656d5f696d6167655f76616c}{Prerequisite~P\ref{premise:5365742e6d656d5f6f665f6d656d5f696d6167655f76616c}}, we obtain \(\left(v\, :\, V(G^{'})\right) \in V(M)\).
Applying the hypothesis \(M\text{ is a matching}\), we obtain \(\exists! w\, :\, V(G^{'}),\; \left(v\, :\, V(G^{'})\right) \sim_{M} w\).
Consider the following case.
\begin{proofstep}
Fix \(w\) in \(V(G^{'})\).
Assume \((v\, :\, V(G^{'})) \sim_{M} w \land (\forall y\, :\, V(G^{'}),\; (v\, :\, V(G^{'})) \sim_{M} y \implies y = w)\).
Choose the witness \(\operatorname{val}(w)\).
By \paperlink{premise:5365742e6d656d5f696d616765}{Prerequisite~P\ref{premise:5365742e6d656d5f696d616765}} with \(\alpha = V(G^{'})\) and \(f = \operatorname{val}\) and \(s = V(M)\) and \(y = v\), we obtain \(v \in \operatorname{val}\left[V(M)\right] \iff \left(\exists x\, :\, V(G^{'}),\; x \in V(M) \land x = v\right)\).
Use the hypothesis \(v \in V(\operatorname{coe\,Subgraph}\left(M\right))\).
Using \paperlink{premise:53696d706c6547726170682e53756267726170682e76657274735f636f655375626772617068}{Prerequisite~P\ref{premise:53696d706c6547726170682e53756267726170682e76657274735f636f655375626772617068}} with \(G^{''} = M\) in reverse gives \(\operatorname{val}\left[V(M)\right] = V(\operatorname{coe\,Subgraph}\left(M\right))\).
Reversing the equality, we obtain \(V(\operatorname{coe\,Subgraph}\left(M\right)) = \operatorname{val}\left[V(M)\right]\).
Substituting equals, we obtain \(v \in \operatorname{val}\left[V(M)\right]\).
By the forward implication, we obtain \(\exists x\, :\, V(G^{'}),\; x \in V(M) \land x = v\).
Consider the following case.
\begin{proofstep}
Fix \(v_{1}\) in \(V(G^{'})\).
Assume \(v_{1} \in V(M) \land \operatorname{val}(v_{1}) = v\).
The substitution uses the following argument.
\begin{proofstep}
Assume \(v \in V(G^{'})\).
For the required implication, argue as follows.
\begin{proofstep}
Assume \(\operatorname{val}(w) \in V(G^{'})\).
Use the hypothesis \(\operatorname{val}(w) \in V(G^{'})\).
By \paperlink{premise:537562747970652e636f655f657461}{Prerequisite~P\ref{premise:537562747970652e636f655f657461}} with \(a = w\), we obtain \(\left(\operatorname{val}(w)\, :\, V(G^{'})\right) = w\).
Applying the same function to equal arguments, we obtain \(((v\, :\, V(G^{'})) \sim_{M} (\operatorname{val}(w)\, :\, V(G^{'}))) = ((v\, :\, V(G^{'})) \sim_{M} w)\).
By \paperlink{premise:4966662e6f665f6571}{Prerequisite~P\ref{premise:4966662e6f665f6571}}, we obtain \((v\, :\, V(G^{'})) \sim_{M} (\operatorname{val}(w)\, :\, V(G^{'})) \iff (v\, :\, V(G^{'})) \sim_{M} w\).
\end{proofstep}
By \paperlink{premise:537562747970652e636f655f70726f70}{Prerequisite~P\ref{premise:537562747970652e636f655f70726f70}} with \(a = w\), we obtain \(\operatorname{val}(w) \in V(G^{'})\).
By \paperlink{premise:65715f74727565}{Prerequisite~P\ref{premise:65715f74727565}}, we obtain \(\left(\operatorname{val}(w) \in V(G^{'})\right) = \mathrm{True}\).
By \paperlink{premise:4966662e6f665f6571}{Prerequisite~P\ref{premise:4966662e6f665f6571}}, we obtain \(\operatorname{val}(w) \in V(G^{'}) \iff \mathrm{True}\).
By \paperlink{premise:6578697374735f70726f705f636f6e6772}{Prerequisite~P\ref{premise:6578697374735f70726f705f636f6e6772}}, we obtain \((\exists \mathit{hw}_{1}\text{ proving }(\operatorname{val}(w) \in V(G^{'})),\; (v\, :\, V(G^{'})) \sim_{M} (\operatorname{val}(w)\, :\, V(G^{'}))) \iff \mathrm{True} \land (v\, :\, V(G^{'})) \sim_{M} w\).
By propositional extensionality, we obtain \((\exists \mathit{hw}_{1}\text{ proving }(\operatorname{val}(w) \in V(G^{'})),\; (v\, :\, V(G^{'})) \sim_{M} (\operatorname{val}(w)\, :\, V(G^{'}))) = (\mathrm{True} \land (v\, :\, V(G^{'})) \sim_{M} w)\).
By \paperlink{premise:696e73744e6f6e656d7074794f66496e68616269746564}{Prerequisite~P\ref{premise:696e73744e6f6e656d7074794f66496e68616269746564}} with \(\alpha = \mathrm{True}\), we obtain \(\mathrm{True}\text{ is nonempty}\).
By \paperlink{premise:6578697374735f636f6e7374}{Prerequisite~P\ref{premise:6578697374735f636f6e7374}} with \(\alpha = \mathrm{True}\), we obtain \((\exists x\, :\, \mathrm{True},\; (v\, :\, V(G^{'})) \sim_{M} w) \iff (v\, :\, V(G^{'})) \sim_{M} w\).
By propositional extensionality, we obtain \((\mathrm{True} \land (v\, :\, V(G^{'})) \sim_{M} w) = ((v\, :\, V(G^{'})) \sim_{M} w)\).
Chaining the equalities, we obtain \((\exists \mathit{hw}_{1}\text{ proving }(\operatorname{val}(w) \in V(G^{'})),\; (v\, :\, V(G^{'})) \sim_{M} (\operatorname{val}(w)\, :\, V(G^{'}))) = ((v\, :\, V(G^{'})) \sim_{M} w)\).
By \paperlink{premise:4966662e6f665f6571}{Prerequisite~P\ref{premise:4966662e6f665f6571}}, we obtain \((\exists \mathit{hw}_{1}\text{ proving }(\operatorname{val}(w) \in V(G^{'})),\; (v\, :\, V(G^{'})) \sim_{M} (\operatorname{val}(w)\, :\, V(G^{'}))) \iff (v\, :\, V(G^{'})) \sim_{M} w\).
\end{proofstep}
The construction requires evidence that \(v \in V(G^{'})\).
By \paperlink{premise:537562747970652e70726f7065727479}{Prerequisite~P\ref{premise:537562747970652e70726f7065727479}} with \(\mathit{self} = v_{1}\), we obtain \(\operatorname{val}(v_{1}) \in V(G^{'})\).
Use the hypothesis \(v_{1} \in V(M) \land \operatorname{val}(v_{1}) = v\).
Taking the second conjunct, we obtain \(\operatorname{val}(v_{1}) = v\).
Substituting equals, we obtain \(v \in V(G^{'})\).
Use the hypothesis \((v\, :\, V(G^{'})) \sim_{M} w \land (\forall y\, :\, V(G^{'}),\; (v\, :\, V(G^{'})) \sim_{M} y \implies y = w)\).
Taking the first conjunct, we obtain \(\left(v\, :\, V(G^{'})\right) \sim_{M} w\).
This witness proves existence. Hence \(\exists h\text{ proving }\left(v \in V(G^{'})\right),\; \left(v\, :\, V(G^{'})\right) \sim_{M} w\).
Rewriting with \paperlink{premise:53696d706c6547726170682e53756267726170682e636f6553756267726170685f61646a}{Prerequisite~P\ref{premise:53696d706c6547726170682e53756267726170682e636f6553756267726170685f61646a}}, \paperlink{premise:6578697374735f70726f705f636f6e6772}{Prerequisite~P\ref{premise:6578697374735f70726f705f636f6e6772}}, \paperlink{premise:4966662e6f665f6571}{Prerequisite~P\ref{premise:4966662e6f665f6571}} and the auxiliary result established above in the preceding result gives \(v \sim_{\operatorname{coe\,Subgraph}\left(M\right)} \operatorname{val}(w)\).
\end{proofstep}
Using the existential hypothesis, we obtain \(v \sim_{\operatorname{coe\,Subgraph}\left(M\right)} \operatorname{val}(w)\).
For the required implication, argue as follows.
\begin{proofstep}
Fix \(y\) in \(V\).
Assume \(v \sim_{\operatorname{coe\,Subgraph}\left(M\right)} y\).
By \paperlink{premise:53696d706c6547726170682e53756267726170682e636f6553756267726170685f61646a}{Prerequisite~P\ref{premise:53696d706c6547726170682e53756267726170682e636f6553756267726170685f61646a}} with \(G^{''} = M\) and \(y\) in place of its parameter \(w\), we obtain \(v \sim_{\operatorname{coe\,Subgraph}(M)} y \iff (\exists \mathit{hv}_{1}\text{ proving }(v \in V(G^{'})),\; \exists \mathit{hw}_{1}\text{ proving }(y \in V(G^{'})),\; (v\, :\, V(G^{'})) \sim_{M} (y\, :\, V(G^{'})))\).
Use the hypothesis \(v \sim_{\operatorname{coe\,Subgraph}\left(M\right)} y\).
By the forward implication, we obtain \(\exists \mathit{hv}_{1}\text{ proving }(v \in V(G^{'})),\; \exists \mathit{hw}_{1}\text{ proving }(y \in V(G^{'})),\; (v\, :\, V(G^{'})) \sim_{M} (y\, :\, V(G^{'}))\).
Consider the following case.
\begin{proofstep}
Assume \(v \in V(G^{'})\).
Assume \(\exists \mathit{hw}_{1}\text{ proving }(y \in V(G^{'})),\; (v\, :\, V(G^{'})) \sim_{M} (y\, :\, V(G^{'}))\).
Use the hypothesis \(\exists \mathit{hw}_{1}\text{ proving }(y \in V(G^{'})),\; (v\, :\, V(G^{'})) \sim_{M} (y\, :\, V(G^{'}))\).
Consider the following case.
\begin{proofstep}
Assume \(y \in V(G^{'})\).
Assume \(\left(v\, :\, V(G^{'})\right) \sim_{M} \left(y\, :\, V(G^{'})\right)\).
By definition, we obtain \(y = y\).
Rewriting with \paperlink{premise:416e642e7269676874}{Prerequisite~P\ref{premise:416e642e7269676874}} (using \((v\, :\, V(G^{'})) \sim_{M} w \land (\forall y\, :\, V(G^{'}),\; (v\, :\, V(G^{'})) \sim_{M} y \implies y = w)\) and \(\left(v\, :\, V(G^{'})\right) \sim_{M} \left(y\, :\, V(G^{'})\right)\)) in the preceding result gives \(y = \operatorname{val}(w)\).
\end{proofstep}
Using the existential hypothesis, we obtain \(y = \operatorname{val}(w)\).
\end{proofstep}
Using the existential hypothesis, we obtain \(y = \operatorname{val}(w)\).
\end{proofstep}
Combining the two conclusions, we obtain \(v \sim_{\operatorname{coe\,Subgraph}(M)} \operatorname{val}(w) \land (\forall y\, :\, V,\; v \sim_{\operatorname{coe\,Subgraph}(M)} y \implies y = \operatorname{val}(w))\).
This witness proves existence. Hence \(\exists x\, :\, V,\; v \sim_{\operatorname{coe\,Subgraph}(M)} x \land (\forall y\, :\, V,\; v \sim_{\operatorname{coe\,Subgraph}(M)} y \implies y = x)\).
\end{proofstep}
Using the existential hypothesis, we obtain \(\exists! w\, :\, V,\; v \sim_{\operatorname{coe\,Subgraph}\left(M\right)} w\).
\end{proof}

\begin{lemma}[{Exists}]\label{theorem:13}
Let \(V\) be a type. Let \(G\) be a simple graph on \(V\). Let \(s\) and \(t\) be subsets of \(V\). Assume \(\operatorname{Disjoint}\left(s, t\right)\). Let \(f\) be a bijection from \(s\) to \(t\). Assume \(\forall v\, :\, s,\; v \sim_{G} \operatorname{val}(f\left(v\right))\). Then
\begin{center}
\(\displaystyle \exists M\, :\, \operatorname{Subgraph}(G),\; V(M) = s \cup t \land M\text{ is a matching}.\)
\end{center}
\end{lemma}

\begin{proof}
Construct an element of \(\operatorname{Subgraph}\left(G\right)\) with verts \(s \cup t\) and adj \(v \mapsto w \mapsto (\exists h_{1}\text{ proving }(v \in s),\; \operatorname{val}(f((v\, :\, s))) = w) \lor (\exists h_{1}\text{ proving }(w \in s),\; \operatorname{val}(f((w\, :\, s))) = v)\).
\par\smallskip\noindent\emph{Adj difference.}\enspace
Fix \(v\) in \(V\).
Fix \(w\) in \(V\).
Assume \((\exists h_{1}\text{ proving }(v \in s),\; \operatorname{val}(f((v\, :\, s))) = w) \lor (\exists h_{1}\text{ proving }(w \in s),\; \operatorname{val}(f((w\, :\, s))) = v)\).
Use the hypothesis \((\exists h_{1}\text{ proving }(v \in s),\; \operatorname{val}(f((v\, :\, s))) = w) \lor (\exists h_{1}\text{ proving }(w \in s),\; \operatorname{val}(f((w\, :\, s))) = v)\).
Consider the following case.
\begin{proofstep}
Assume \(\exists h_{2}\text{ proving }(v \in s),\; \operatorname{val}(f((v\, :\, s))) = w\).
Use the hypothesis \(\exists h_{2}\text{ proving }(v \in s),\; \operatorname{val}(f((v\, :\, s))) = w\).
Consider the following case.
\begin{proofstep}
Assume \(v \in s\).
Assume \(\operatorname{val}(f\left(\left(v\, :\, s\right)\right)) = w\).
Applying the hypothesis \(\forall v\, :\, s,\; v \sim_{G} \operatorname{val}(f\left(v\right))\), we obtain \(v \sim_{G} \operatorname{val}(f\left(\left(v\, :\, s\right)\right))\).
Use the hypothesis \(\operatorname{val}(f\left(\left(v\, :\, s\right)\right)) = w\).
Substituting equals, we obtain \(v \sim_{G} w\).
\end{proofstep}
Using the existential hypothesis, we obtain \(v \sim_{G} w\).
\end{proofstep}
Consider the following case.
\begin{proofstep}
Assume \(\exists h_{2}\text{ proving }(w \in s),\; \operatorname{val}(f((w\, :\, s))) = v\).
Use the hypothesis \(\exists h_{2}\text{ proving }(w \in s),\; \operatorname{val}(f((w\, :\, s))) = v\).
Consider the following case.
\begin{proofstep}
Assume \(w \in s\).
Assume \(\operatorname{val}(f\left(\left(w\, :\, s\right)\right)) = v\).
Applying the hypothesis \(\forall v\, :\, s,\; v \sim_{G} \operatorname{val}(f\left(v\right))\), we obtain \(w \sim_{G} \operatorname{val}(f\left(\left(w\, :\, s\right)\right))\).
By \paperlink{premise:53696d706c6547726170682e41646a2e73796d6d}{Prerequisite~P\ref{premise:53696d706c6547726170682e41646a2e73796d6d}}, we obtain \(\operatorname{val}(f\left(\left(w\, :\, s\right)\right)) \sim_{G} w\).
Use the hypothesis \(\operatorname{val}(f\left(\left(w\, :\, s\right)\right)) = v\).
Substituting equals, we obtain \(v \sim_{G} w\).
\end{proofstep}
Using the existential hypothesis, we obtain \(v \sim_{G} w\).
\end{proofstep}
The two alternatives exhaust the cases. Hence \(v \sim_{G} w\).
\par\smallskip\noindent\emph{Edge vert.}\enspace
Fix \(v\) in \(V\).
Fix \(w\) in \(V\).
Assume \((\exists h_{1}\text{ proving }(v \in s),\; \operatorname{val}(f((v\, :\, s))) = w) \lor (\exists h_{1}\text{ proving }(w \in s),\; \operatorname{val}(f((w\, :\, s))) = v)\).
Use the hypothesis \((\exists h_{1}\text{ proving }(v \in s),\; \operatorname{val}(f((v\, :\, s))) = w) \lor (\exists h_{1}\text{ proving }(w \in s),\; \operatorname{val}(f((w\, :\, s))) = v)\).
Consider the following case.
\begin{proofstep}
Assume \(\exists h_{2}\text{ proving }(v \in s),\; \operatorname{val}(f((v\, :\, s))) = w\).
Use the hypothesis \(\exists h_{2}\text{ proving }(v \in s),\; \operatorname{val}(f((v\, :\, s))) = w\).
Consider the following case.
\begin{proofstep}
Assume \(v \in s\).
Assume \(\operatorname{val}(f\left(\left(v\, :\, s\right)\right)) = w\).
Rewriting with \(v \in s\) and \paperlink{premise:747275655f6f72}{Prerequisite~P\ref{premise:747275655f6f72}} gives \(v \in s \lor v \in t\).
\end{proofstep}
Using the existential hypothesis, we obtain \(v \in s \lor v \in t\).
\end{proofstep}
Consider the following case.
\begin{proofstep}
Assume \(\exists h_{2}\text{ proving }(w \in s),\; \operatorname{val}(f((w\, :\, s))) = v\).
Use the hypothesis \(\exists h_{2}\text{ proving }(w \in s),\; \operatorname{val}(f((w\, :\, s))) = v\).
Consider the following case.
\begin{proofstep}
Assume \(w \in s\).
Assume \(\operatorname{val}(f\left(\left(w\, :\, s\right)\right)) = v\).
Rewriting with \paperlink{premise:537562747970652e636f655f70726f70}{Prerequisite~P\ref{premise:537562747970652e636f655f70726f70}} and \paperlink{premise:6f725f74727565}{Prerequisite~P\ref{premise:6f725f74727565}} gives \(\operatorname{val}(f((w\, :\, s))) \in s \lor \operatorname{val}(f((w\, :\, s))) \in t\).
Use the hypothesis \(\operatorname{val}(f\left(\left(w\, :\, s\right)\right)) = v\).
Substituting equals, we obtain \(v \in s \lor v \in t\).
\end{proofstep}
Using the existential hypothesis, we obtain \(v \in s \lor v \in t\).
\end{proofstep}
The two alternatives exhaust the cases. Hence \(v \in s \lor v \in t\).
Rewriting with \paperlink{premise:5365742e6d656d5f756e696f6e}{Prerequisite~P\ref{premise:5365742e6d656d5f756e696f6e}} in the preceding result gives \(v \in s \cup t\).
\par\smallskip\noindent\emph{Symmetry.}\enspace
Fix \(x\) in \(V\).
Fix \(y\) in \(V\).
Assume \((\exists h_{1}\text{ proving }(x \in s),\; \operatorname{val}(f((x\, :\, s))) = y) \lor (\exists h_{1}\text{ proving }(y \in s),\; \operatorname{val}(f((y\, :\, s))) = x)\).
Use the hypothesis \((\exists h_{1}\text{ proving }(x \in s),\; \operatorname{val}(f((x\, :\, s))) = y) \lor (\exists h_{1}\text{ proving }(y \in s),\; \operatorname{val}(f((y\, :\, s))) = x)\).
Consider the following case.
\begin{proofstep}
Assume \(\exists h_{2}\text{ proving }(x \in s),\; \operatorname{val}(f((x\, :\, s))) = y\).
Rewriting with \(\exists h_{2}\text{ proving }(x \in s),\; \operatorname{val}(f((x\, :\, s))) = y\) and \paperlink{premise:6f725f74727565}{Prerequisite~P\ref{premise:6f725f74727565}} gives \((\exists h_{2}\text{ proving }(y \in s),\; \operatorname{val}(f((y\, :\, s))) = x) \lor (\exists h_{2}\text{ proving }(x \in s),\; \operatorname{val}(f((x\, :\, s))) = y)\).
\end{proofstep}
Consider the following case.
\begin{proofstep}
Assume \(\exists h_{2}\text{ proving }(y \in s),\; \operatorname{val}(f((y\, :\, s))) = x\).
Rewriting with \(\exists h_{2}\text{ proving }(y \in s),\; \operatorname{val}(f((y\, :\, s))) = x\) and \paperlink{premise:747275655f6f72}{Prerequisite~P\ref{premise:747275655f6f72}} gives \((\exists h_{2}\text{ proving }(y \in s),\; \operatorname{val}(f((y\, :\, s))) = x) \lor (\exists h_{2}\text{ proving }(x \in s),\; \operatorname{val}(f((x\, :\, s))) = y)\).
\end{proofstep}
The two alternatives exhaust the cases. Hence \((\exists h_{1}\text{ proving }(y \in s),\; \operatorname{val}(f((y\, :\, s))) = x) \lor (\exists h_{1}\text{ proving }(x \in s),\; \operatorname{val}(f((x\, :\, s))) = y)\).
The substitution uses the following argument.
\begin{proofstep}
Fix \(v\) in \(V\).
By \paperlink{premise:5365742e6d656d5f756e696f6e}{Prerequisite~P\ref{premise:5365742e6d656d5f756e696f6e}} with \(x = v\) and \(a = s\) and \(b = t\), we obtain \(v \in s \cup t \iff v \in s \lor v \in t\).
By propositional extensionality, we obtain \(\left(v \in s \cup t\right) = \left(v \in s \lor v \in t\right)\).
By definition, we obtain \((\exists! w\, :\, V,\; (\exists h_{1}\text{ proving }(v \in s),\; \operatorname{val}(f((v\, :\, s))) = w) \lor (\exists h_{1}\text{ proving }(w \in s),\; \operatorname{val}(f((w\, :\, s))) = v)) = (\exists! w\, :\, V,\; (\exists h_{1}\text{ proving }(v \in s),\; \operatorname{val}(f((v\, :\, s))) = w) \lor (\exists h_{1}\text{ proving }(w \in s),\; \operatorname{val}(f((w\, :\, s))) = v))\).
By \paperlink{premise:696d706c6965735f636f6e6772}{Prerequisite~P\ref{premise:696d706c6965735f636f6e6772}} with \(p_{1} = v \in s \cup t\) and \(p_{2} = v \in s \lor v \in t\) and \(q_{1} = \exists! w\, :\, V,\; (\exists h_{1}\text{ proving }(v \in s),\; \operatorname{val}(f((v\, :\, s))) = w) \lor (\exists h_{1}\text{ proving }(w \in s),\; \operatorname{val}(f((w\, :\, s))) = v)\) and \(q_{2} = \exists! w\, :\, V,\; (\exists h_{1}\text{ proving }(v \in s),\; \operatorname{val}(f((v\, :\, s))) = w) \lor (\exists h_{1}\text{ proving }(w \in s),\; \operatorname{val}(f((w\, :\, s))) = v)\), we obtain \((v \in s \cup t \implies (\exists! w\, :\, V,\; (\exists h_{1}\text{ proving }(v \in s),\; \operatorname{val}(f((v\, :\, s))) = w) \lor (\exists h_{1}\text{ proving }(w \in s),\; \operatorname{val}(f((w\, :\, s))) = v))) = (v \in s \lor v \in t \implies (\exists! w\, :\, V,\; (\exists h_{1}\text{ proving }(v \in s),\; \operatorname{val}(f((v\, :\, s))) = w) \lor (\exists h_{1}\text{ proving }(w \in s),\; \operatorname{val}(f((w\, :\, s))) = v)))\).
\end{proofstep}
Fix \(v\) in \(V\).
Assume \(v \in s \lor v \in t\).
Use the hypothesis \(v \in s \lor v \in t\).
Consider the following case.
\begin{proofstep}
Assume \(v \in s\).
Choose the witness \(\operatorname{val}(f\left(\left(v\, :\, s\right)\right))\).
The substitution uses the following argument.
\begin{proofstep}
Assume \(v \in s\).
By \paperlink{premise:65715f73656c66}{Prerequisite~P\ref{premise:65715f73656c66}} with \(\operatorname{val}(f\left(\left(v\, :\, s\right)\right))\) in place of its parameter \(a\), we obtain \((\operatorname{val}(f((v\, :\, s))) = \operatorname{val}(f((v\, :\, s)))) = \mathrm{True}\).
By \paperlink{premise:4966662e6f665f6571}{Prerequisite~P\ref{premise:4966662e6f665f6571}}, we obtain \(\operatorname{val}(f((v\, :\, s))) = \operatorname{val}(f((v\, :\, s))) \iff \mathrm{True}\).
\end{proofstep}
The substitution uses the following argument.
\begin{proofstep}
Fix \(y\) in \(V\).
For the required implication, argue as follows.
\begin{proofstep}
Assume \(v \in s\).
By definition, we obtain \((\operatorname{val}(f((v\, :\, s))) = y) = (\operatorname{val}(f((v\, :\, s))) = y)\).
By \paperlink{premise:4966662e6f665f6571}{Prerequisite~P\ref{premise:4966662e6f665f6571}}, we obtain \(\operatorname{val}(f((v\, :\, s))) = y \iff \operatorname{val}(f((v\, :\, s))) = y\).
\end{proofstep}
Use the hypothesis \(v \in s\).
By \paperlink{premise:65715f74727565}{Prerequisite~P\ref{premise:65715f74727565}}, we obtain \(\left(v \in s\right) = \mathrm{True}\).
By \paperlink{premise:4966662e6f665f6571}{Prerequisite~P\ref{premise:4966662e6f665f6571}}, we obtain \(v \in s \iff \mathrm{True}\).
By \paperlink{premise:6578697374735f70726f705f636f6e6772}{Prerequisite~P\ref{premise:6578697374735f70726f705f636f6e6772}}, we obtain \((\exists h_{1}\text{ proving }(v \in s),\; \operatorname{val}(f((v\, :\, s))) = y) \iff (\exists h_{1}\text{ proving }\mathrm{True},\; \operatorname{val}(f((v\, :\, s))) = y)\).
By propositional extensionality, we obtain \((\exists h_{1}\text{ proving }(v \in s),\; \operatorname{val}(f((v\, :\, s))) = y) = (\exists h_{1}\text{ proving }\mathrm{True},\; \operatorname{val}(f((v\, :\, s))) = y)\).
By \paperlink{premise:6578697374735f747275655f6c656674}{Prerequisite~P\ref{premise:6578697374735f747275655f6c656674}}, we obtain \((\exists h_{1}\text{ proving }\mathrm{True},\; \operatorname{val}(f((v\, :\, s))) = y) \iff (h_{1} \mapsto \operatorname{val}(f((v\, :\, s))) = y)(\mathrm{trivial})\).
By propositional extensionality, we obtain \((\exists h_{1}\text{ proving }\mathrm{True},\; \operatorname{val}(f((v\, :\, s))) = y) = (\operatorname{val}(f((v\, :\, s))) = y)\).
Chaining the equalities, we obtain \((\exists h_{1}\text{ proving }(v \in s),\; \operatorname{val}(f((v\, :\, s))) = y) = (\operatorname{val}(f((v\, :\, s))) = y)\).
Applying the same function to equal arguments, we obtain \(((\exists h_{1}\text{ proving }(v \in s),\; \operatorname{val}(f((v\, :\, s))) = y) \lor (\exists h_{1}\text{ proving }(y \in s),\; \operatorname{val}(f((y\, :\, s))) = v)) = (\operatorname{val}(f((v\, :\, s))) = y \lor (\exists h_{1}\text{ proving }(y \in s),\; \operatorname{val}(f((y\, :\, s))) = v))\).
By definition, we obtain \((y = \operatorname{val}(f((v\, :\, s)))) = (y = \operatorname{val}(f((v\, :\, s))))\).
By \paperlink{premise:696d706c6965735f636f6e6772}{Prerequisite~P\ref{premise:696d706c6965735f636f6e6772}} with \(p_{1} = (\exists h_{1}\text{ proving }(v \in s),\; \operatorname{val}(f((v\, :\, s))) = y) \lor (\exists h_{1}\text{ proving }(y \in s),\; \operatorname{val}(f((y\, :\, s))) = v)\) and \(p_{2} = \operatorname{val}(f((v\, :\, s))) = y \lor (\exists h_{1}\text{ proving }(y \in s),\; \operatorname{val}(f((y\, :\, s))) = v)\) and \(q_{1} = y = \operatorname{val}(f\left(\left(v\, :\, s\right)\right))\) and \(q_{2} = y = \operatorname{val}(f\left(\left(v\, :\, s\right)\right))\), we obtain \(((\exists h_{1}\text{ proving }(v \in s),\; \operatorname{val}(f((v\, :\, s))) = y) \lor (\exists h_{1}\text{ proving }(y \in s),\; \operatorname{val}(f((y\, :\, s))) = v) \implies y = \operatorname{val}(f((v\, :\, s)))) = (\operatorname{val}(f((v\, :\, s))) = y \lor (\exists h_{1}\text{ proving }(y \in s),\; \operatorname{val}(f((y\, :\, s))) = v) \implies y = \operatorname{val}(f((v\, :\, s))))\).
\end{proofstep}
Fix \(y\) in \(V\).
Assume \(\operatorname{val}(f((v\, :\, s))) = y \lor (\exists h_{1}\text{ proving }(y \in s),\; \operatorname{val}(f((y\, :\, s))) = v)\).
Use the hypothesis \(\operatorname{val}(f((v\, :\, s))) = y \lor (\exists h_{1}\text{ proving }(y \in s),\; \operatorname{val}(f((y\, :\, s))) = v)\).
Consider the following case.
\begin{proofstep}
Assume \(\operatorname{val}(f\left(\left(v\, :\, s\right)\right)) = y\).
By definition, we obtain \(\operatorname{val}(f((v\, :\, s))) = \operatorname{val}(f((v\, :\, s)))\).
Use the hypothesis \(\operatorname{val}(f\left(\left(v\, :\, s\right)\right)) = y\).
Substituting equals, we obtain \(y = \operatorname{val}(f\left(\left(v\, :\, s\right)\right))\).
\end{proofstep}
Consider the following case.
\begin{proofstep}
Assume \(\exists h_{2}\text{ proving }(y \in s),\; \operatorname{val}(f((y\, :\, s))) = v\).
Use the hypothesis \(\exists h_{2}\text{ proving }(y \in s),\; \operatorname{val}(f((y\, :\, s))) = v\).
Consider the following case.
\begin{proofstep}
Assume \(y \in s\).
Assume \(\operatorname{val}(f\left(\left(y\, :\, s\right)\right)) = v\).
For the required implication, argue as follows.
\begin{proofstep}
Assume \(\operatorname{val}(f\left(\left(y\, :\, s\right)\right)) \in s\).
Use the hypothesis \(\operatorname{Disjoint}\left(s, t\right)\).
Use the hypothesis \(\operatorname{val}(f\left(\left(y\, :\, s\right)\right)) \in s\).
By \paperlink{premise:537562747970652e636f655f70726f70}{Prerequisite~P\ref{premise:537562747970652e636f655f70726f70}} with \(f\left(\left(y\, :\, s\right)\right)\) in place of its parameter \(a\), we obtain \(\operatorname{val}(f\left(\left(y\, :\, s\right)\right)) \in t\).
By definition, we obtain \(\operatorname{val}(f((y\, :\, s))) = \operatorname{val}(f((y\, :\, s)))\).
By \paperlink{premise:4469736a6f696e742e6e655f6f665f6d656d}{Prerequisite~P\ref{premise:4469736a6f696e742e6e655f6f665f6d656d}}, we obtain \(\mathrm{False}\).
The contradiction rules out this case. Hence \(y = \operatorname{val}(f((\operatorname{val}(f((y\, :\, s)))\, :\, s)))\).
\end{proofstep}
Use the hypothesis \(\operatorname{val}(f\left(\left(y\, :\, s\right)\right)) = v\).
Use the hypothesis \(v \in s\).
Substituting equals, we obtain \(y = \operatorname{val}(f\left(\left(v\, :\, s\right)\right))\).
\end{proofstep}
Using the existential hypothesis, we obtain \(y = \operatorname{val}(f\left(\left(v\, :\, s\right)\right))\).
\end{proofstep}
The two alternatives exhaust the cases. Hence \(y = \operatorname{val}(f\left(\left(v\, :\, s\right)\right))\).
Rewriting with \paperlink{premise:6578697374735f70726f705f636f6e6772}{Prerequisite~P\ref{premise:6578697374735f70726f705f636f6e6772}}, \paperlink{premise:4966662e6f665f6571}{Prerequisite~P\ref{premise:4966662e6f665f6571}}, \(v \in s\), \paperlink{premise:6578697374735f636f6e7374}{Prerequisite~P\ref{premise:6578697374735f636f6e7374}}, \paperlink{premise:696e73744e6f6e656d7074794f66496e68616269746564}{Prerequisite~P\ref{premise:696e73744e6f6e656d7074794f66496e68616269746564}}, \paperlink{premise:747275655f6f72}{Prerequisite~P\ref{premise:747275655f6f72}}, \paperlink{premise:666f72616c6c5f636f6e6772}{Prerequisite~P\ref{premise:666f72616c6c5f636f6e6772}}, \paperlink{premise:747275655f616e64}{Prerequisite~P\ref{premise:747275655f616e64}}, the identity established above and the auxiliary result established above in the preceding result gives \(((\exists h_{1}\text{ proving }(v \in s),\; \operatorname{val}(f((v\, :\, s))) = \operatorname{val}(f((v\, :\, s)))) \lor (\exists h_{1}\text{ proving }(\operatorname{val}(f((v\, :\, s))) \in s),\; \operatorname{val}(f((\operatorname{val}(f((v\, :\, s)))\, :\, s))) = v)) \land (\forall y\, :\, V,\; (\exists h_{1}\text{ proving }(v \in s),\; \operatorname{val}(f((v\, :\, s))) = y) \lor (\exists h_{1}\text{ proving }(y \in s),\; \operatorname{val}(f((y\, :\, s))) = v) \implies y = \operatorname{val}(f((v\, :\, s))))\).
This witness proves existence. Hence \(\exists x\, :\, V,\; ((\exists h_{1}\text{ proving }(v \in s),\; \operatorname{val}(f((v\, :\, s))) = x) \lor (\exists h_{1}\text{ proving }(x \in s),\; \operatorname{val}(f((x\, :\, s))) = v)) \land (\forall y\, :\, V,\; (\exists h_{1}\text{ proving }(v \in s),\; \operatorname{val}(f((v\, :\, s))) = y) \lor (\exists h_{1}\text{ proving }(y \in s),\; \operatorname{val}(f((y\, :\, s))) = v) \implies y = x)\).
\end{proofstep}
Consider the following case.
\begin{proofstep}
Assume \(v \in t\).
Choose the witness \(\operatorname{val}(\left(f^{-1}\right)\left(\left(v\, :\, t\right)\right))\).
The substitution uses the following argument.
\begin{proofstep}
Assume \(\operatorname{val}(\left(f^{-1}\right)\left(\left(v\, :\, t\right)\right)) \in s\).
Use the hypothesis \(\operatorname{val}(\left(f^{-1}\right)\left(\left(v\, :\, t\right)\right)) \in s\).
By \paperlink{premise:537562747970652e636f655f657461}{Prerequisite~P\ref{premise:537562747970652e636f655f657461}} with \(\left(f^{-1}\right)\left(\left(v\, :\, t\right)\right)\) in place of its parameter \(a\), we obtain \((\operatorname{val}((f^{-1})((v\, :\, t)))\, :\, s) = (f^{-1})((v\, :\, t))\).
Applying the same function to equal arguments, we obtain \(f((\operatorname{val}((f^{-1})((v\, :\, t)))\, :\, s)) = f((f^{-1})((v\, :\, t)))\).
By \paperlink{premise:45717569762e6170706c795f73796d6d5f6170706c79}{Prerequisite~P\ref{premise:45717569762e6170706c795f73796d6d5f6170706c79}} with \(\alpha = s\) and \(\beta = t\) and \(e = f\) and \(x = \left(v\, :\, t\right)\), we obtain \(f\left(\left(f^{-1}\right)\left(\left(v\, :\, t\right)\right)\right) = \left(v\, :\, t\right)\).
Chaining the equalities, we obtain \(f((\operatorname{val}((f^{-1})((v\, :\, t)))\, :\, s)) = (v\, :\, t)\).
Applying the same function to equal arguments, we obtain \((\operatorname{val}(f((\operatorname{val}((f^{-1})((v\, :\, t)))\, :\, s))) = v) = (v = v)\).
By \paperlink{premise:65715f73656c66}{Prerequisite~P\ref{premise:65715f73656c66}} with \(v\) in place of its parameter \(a\), we obtain \(\left(v = v\right) = \mathrm{True}\).
Chaining the equalities, we obtain \((\operatorname{val}(f((\operatorname{val}((f^{-1})((v\, :\, t)))\, :\, s))) = v) = \mathrm{True}\).
By \paperlink{premise:4966662e6f665f6571}{Prerequisite~P\ref{premise:4966662e6f665f6571}}, we obtain \(\operatorname{val}(f((\operatorname{val}((f^{-1})((v\, :\, t)))\, :\, s))) = v \iff \mathrm{True}\).
\end{proofstep}
Fix \(y\) in \(V\).
Assume \((\exists h_{1}\text{ proving }(v \in s),\; \operatorname{val}(f((v\, :\, s))) = y) \lor (\exists h_{1}\text{ proving }(y \in s),\; \operatorname{val}(f((y\, :\, s))) = v)\).
Use the hypothesis \((\exists h_{1}\text{ proving }(v \in s),\; \operatorname{val}(f((v\, :\, s))) = y) \lor (\exists h_{1}\text{ proving }(y \in s),\; \operatorname{val}(f((y\, :\, s))) = v)\).
Consider the following case.
\begin{proofstep}
Assume \(\exists h_{2}\text{ proving }(v \in s),\; \operatorname{val}(f((v\, :\, s))) = y\).
Use the hypothesis \(\exists h_{2}\text{ proving }(v \in s),\; \operatorname{val}(f((v\, :\, s))) = y\).
Consider the following case.
\begin{proofstep}
Assume \(v \in s\).
Assume \(\operatorname{val}(f\left(\left(v\, :\, s\right)\right)) = y\).
Use the hypothesis \(\operatorname{Disjoint}\left(s, t\right)\).
Use the hypothesis \(v \in s\).
Use the hypothesis \(v \in t\).
By definition, we obtain \(v = v\).
By \paperlink{premise:4469736a6f696e742e6e655f6f665f6d656d}{Prerequisite~P\ref{premise:4469736a6f696e742e6e655f6f665f6d656d}}, we obtain \(\mathrm{False}\).
The contradiction rules out this case. Hence \(\operatorname{val}(f((v\, :\, s))) = \operatorname{val}((f^{-1})((v\, :\, t)))\).
Use the hypothesis \(\operatorname{val}(f\left(\left(v\, :\, s\right)\right)) = y\).
Substituting equals, we obtain \(y = \operatorname{val}(\left(f^{-1}\right)\left(\left(v\, :\, t\right)\right))\).
\end{proofstep}
Using the existential hypothesis, we obtain \(y = \operatorname{val}(\left(f^{-1}\right)\left(\left(v\, :\, t\right)\right))\).
\end{proofstep}
Consider the following case.
\begin{proofstep}
Assume \(\exists h_{2}\text{ proving }(y \in s),\; \operatorname{val}(f((y\, :\, s))) = v\).
Use the hypothesis \(\exists h_{2}\text{ proving }(y \in s),\; \operatorname{val}(f((y\, :\, s))) = v\).
Consider the following case.
\begin{proofstep}
Assume \(y \in s\).
Assume \(\operatorname{val}(f\left(\left(y\, :\, s\right)\right)) = v\).
For the required implication, argue as follows.
\begin{proofstep}
Assume \(\operatorname{val}(f\left(\left(y\, :\, s\right)\right)) \in t\).
Rewriting with \paperlink{premise:537562747970652e636f655f657461}{Prerequisite~P\ref{premise:537562747970652e636f655f657461}} (using \(\operatorname{val}(f\left(\left(y\, :\, s\right)\right)) \in t\)) and \paperlink{premise:45717569762e73796d6d5f6170706c795f6170706c79}{Prerequisite~P\ref{premise:45717569762e73796d6d5f6170706c795f6170706c79}} gives \(y = \operatorname{val}((f^{-1})((\operatorname{val}(f((y\, :\, s)))\, :\, t)))\).
\end{proofstep}
Use the hypothesis \(\operatorname{val}(f\left(\left(y\, :\, s\right)\right)) = v\).
Use the hypothesis \(v \in t\).
Substituting equals, we obtain \(y = \operatorname{val}(\left(f^{-1}\right)\left(\left(v\, :\, t\right)\right))\).
\end{proofstep}
Using the existential hypothesis, we obtain \(y = \operatorname{val}(\left(f^{-1}\right)\left(\left(v\, :\, t\right)\right))\).
\end{proofstep}
The two alternatives exhaust the cases. Hence \(y = \operatorname{val}(\left(f^{-1}\right)\left(\left(v\, :\, t\right)\right))\).
Rewriting with \paperlink{premise:6578697374735f70726f705f636f6e6772}{Prerequisite~P\ref{premise:6578697374735f70726f705f636f6e6772}}, \paperlink{premise:4966662e6f665f6571}{Prerequisite~P\ref{premise:4966662e6f665f6571}}, \paperlink{premise:537562747970652e636f655f70726f70}{Prerequisite~P\ref{premise:537562747970652e636f655f70726f70}}, \paperlink{premise:6578697374735f636f6e7374}{Prerequisite~P\ref{premise:6578697374735f636f6e7374}}, \paperlink{premise:696e73744e6f6e656d7074794f66496e68616269746564}{Prerequisite~P\ref{premise:696e73744e6f6e656d7074794f66496e68616269746564}}, \paperlink{premise:6f725f74727565}{Prerequisite~P\ref{premise:6f725f74727565}}, \paperlink{premise:747275655f616e64}{Prerequisite~P\ref{premise:747275655f616e64}} and the auxiliary result established above in the preceding result gives \(((\exists h_{1}\text{ proving }(v \in s),\; \operatorname{val}(f((v\, :\, s))) = \operatorname{val}((f^{-1})((v\, :\, t)))) \lor (\exists h_{1}\text{ proving }(\operatorname{val}((f^{-1})((v\, :\, t))) \in s),\; \operatorname{val}(f((\operatorname{val}((f^{-1})((v\, :\, t)))\, :\, s))) = v)) \land (\forall y\, :\, V,\; (\exists h_{1}\text{ proving }(v \in s),\; \operatorname{val}(f((v\, :\, s))) = y) \lor (\exists h_{1}\text{ proving }(y \in s),\; \operatorname{val}(f((y\, :\, s))) = v) \implies y = \operatorname{val}((f^{-1})((v\, :\, t))))\).
This witness proves existence. Hence \(\exists x\, :\, V,\; ((\exists h_{1}\text{ proving }(v \in s),\; \operatorname{val}(f((v\, :\, s))) = x) \lor (\exists h_{1}\text{ proving }(x \in s),\; \operatorname{val}(f((x\, :\, s))) = v)) \land (\forall y\, :\, V,\; (\exists h_{1}\text{ proving }(v \in s),\; \operatorname{val}(f((v\, :\, s))) = y) \lor (\exists h_{1}\text{ proving }(y \in s),\; \operatorname{val}(f((y\, :\, s))) = v) \implies y = x)\).
\end{proofstep}
The two alternatives exhaust the cases. Hence \(\exists! w\, :\, V,\; (\exists h_{1}\text{ proving }(v \in s),\; \operatorname{val}(f((v\, :\, s))) = w) \lor (\exists h_{1}\text{ proving }(w \in s),\; \operatorname{val}(f((w\, :\, s))) = v)\).
Rewriting with \paperlink{premise:666f72616c6c5f636f6e6772}{Prerequisite~P\ref{premise:666f72616c6c5f636f6e6772}}, \paperlink{premise:747275655f616e64}{Prerequisite~P\ref{premise:747275655f616e64}} and the identity established above in the preceding result gives \(V(\operatorname{Simple\,Graph\mathord{:}Subgraph\mathord{:}mk}(s \cup t,\allowbreak{} v \mapsto w \mapsto (\exists h_{1}\text{ proving }(v \in s),\; \operatorname{val}(f((v\, :\, s))) = w) \lor (\exists h_{1}\text{ proving }(w \in s),\; \operatorname{val}(f((w\, :\, s))) = v))) = s \cup t \land \operatorname{Simple\,Graph\mathord{:}Subgraph\mathord{:}mk}(s \cup t,\allowbreak{} v \mapsto w \mapsto (\exists h_{1}\text{ proving }(v \in s),\; \operatorname{val}(f((v\, :\, s))) = w) \lor (\exists h_{1}\text{ proving }(w \in s),\; \operatorname{val}(f((w\, :\, s))) = v))\text{ is a matching}\).
This witness proves existence. Hence \(\exists M\, :\, \operatorname{Subgraph}\left(G\right),\; V(M) = s \cup t \land M\text{ is a matching}\).
\end{proof}

\begin{lemma}[{Map}]\label{theorem:14}
Let \(V\) and \(W\) be types. Let \(G\) be a simple graph on \(V\). Let \(G^{'}\) be a simple graph on \(W\). Let \(M\) be a subgraph of \(G\). Let \(f\) be a graph homomorphism from \(G\) to \(G^{'}\). Assume \(\operatorname{DFun\,Like\mathord{:}coe}\left(f\right)\text{ is injective}\). Assume \(M\text{ is a matching}\). Then
\begin{center}
\(\displaystyle \operatorname{Simple\,Graph\mathord{:}Subgraph\mathord{:}map}(f,\allowbreak{} M)\text{ is a matching}.\)
\end{center}
\end{lemma}

\begin{proof}
Fix \(v\) in \(W\).
Assume \(v \in V(\operatorname{Simple\,Graph\mathord{:}Subgraph\mathord{:}map}\left(f, M\right))\).
Use the hypothesis \(v \in V(\operatorname{Simple\,Graph\mathord{:}Subgraph\mathord{:}map}\left(f, M\right))\).
Consider the following case.
\begin{proofstep}
Fix \(v_{1}\) in \(V\).
Assume \(v_{1} \in V(M) \land f(v_{1}) = v\).
Use the hypothesis \(v_{1} \in V(M) \land f(v_{1}) = v\).
Consider the following case.
\begin{proofstep}
Assume \(v_{1} \in V(M)\).
Assume \(f(v_{1}) = v\).
Use the hypothesis \(v_{1} \in V(M)\).
Applying the hypothesis \(M\text{ is a matching}\), we obtain \(\exists! w\, :\, V,\; v_{1} \sim_{M} w\).
Consider the following case.
\begin{proofstep}
Fix \(v^{'}\) in \(V\).
Assume \(v_{1} \sim_{M} v^{'} \land \left(\forall y\, :\, V,\; v_{1} \sim_{M} y \implies y = v^{'}\right)\).
Choose the witness \(f(v^{'})\).
Choose the witness \(v_{1}\).
Choose the witness \(v^{'}\).
Use the hypothesis \(v_{1} \sim_{M} v^{'} \land \left(\forall y\, :\, V,\; v_{1} \sim_{M} y \implies y = v^{'}\right)\).
Taking the first conjunct, we obtain \(v_{1} \sim_{M} v^{'}\).
By definition, we obtain \(f(v_{1}) = f(v_{1})\).
By definition, we obtain \(f(v^{'}) = f(v^{'})\).
Combining the two conclusions, we obtain \(f(v_{1}) = f(v_{1}) \land f(v^{'}) = f(v^{'})\).
Combining the two conclusions, we obtain \(v_{1} \sim_{M} v^{'} \land \left(f(v_{1}) = f(v_{1}) \land f(v^{'}) = f(v^{'})\right)\).
This witness proves existence. Hence \(\exists b\, :\, V,\; v_{1} \sim_{M} b \land \left(f(v_{1}) = f(v_{1}) \land f(b) = f(v^{'})\right)\).
This witness proves existence. Hence \(\exists a_{1}\, :\, V,\; \exists b\, :\, V,\; a_{1} \sim_{M} b \land (f(a_{1}) = f(v_{1}) \land f(b) = f(v^{'}))\).
For the required implication, argue as follows.
\begin{proofstep}
Fix \(y\) in \(W\).
Assume \(f(v_{1}) \sim_{\operatorname{Simple\,Graph\mathord{:}Subgraph\mathord{:}map}\left(f, M\right)} y\).
Use the hypothesis \(f(v_{1}) \sim_{\operatorname{Simple\,Graph\mathord{:}Subgraph\mathord{:}map}\left(f, M\right)} y\).
Consider the following case.
\begin{proofstep}
Fix \(w\) in \(V\).
Assume \(\exists b\, :\, V,\; w \sim_{M} b \land \left(f(w) = f(v_{1}) \land f(b) = y\right)\).
Use the hypothesis \(\exists b\, :\, V,\; w \sim_{M} b \land \left(f(w) = f(v_{1}) \land f(b) = y\right)\).
Consider the following case.
\begin{proofstep}
Fix \(w^{'}\) in \(V\).
Assume \(w \sim_{M} w^{'} \land \left(f(w) = f(v_{1}) \land f(w^{'}) = y\right)\).
Use the hypothesis \(w \sim_{M} w^{'} \land \left(f(w) = f(v_{1}) \land f(w^{'}) = y\right)\).
Consider the following case.
\begin{proofstep}
Assume \(w \sim_{M} w^{'}\).
Assume \(f(w) = f(v_{1}) \land f(w^{'}) = y\).
Use the hypothesis \(f(w) = f(v_{1}) \land f(w^{'}) = y\).
Consider the following case.
\begin{proofstep}
Assume \(f(w) = f(v_{1})\).
Assume \(f(w^{'}) = y\).
Use the hypothesis \(f(w) = f(v_{1})\).
Reversing the equality, we obtain \(f(v_{1}) = f(w)\).
Applying the hypothesis \(\operatorname{DFun\,Like\mathord{:}coe}\left(f\right)\text{ is injective}\), we obtain \(v_{1} = w\).
For the required implication, argue as follows.
\begin{proofstep}
Assume \(w = v_{1}\).
For the required implication, argue as follows.
\begin{proofstep}
Assume \(v_{1} \sim_{M} w^{'}\).
Assume \(f(v_{1}) = f(v_{1})\).
Assume \(v_{1} = v_{1}\).
Assume \(\text{the proof of }(v_{1} = v_{1})\text{ agrees under transport with the proof of }(v_{1} = v_{1})\).
Rewriting with \paperlink{premise:416e642e7269676874}{Prerequisite~P\ref{premise:416e642e7269676874}} (using \(v_{1} \sim_{M} v^{'} \land \left(\forall y\, :\, V,\; v_{1} \sim_{M} y \implies y = v^{'}\right)\) and \(v_{1} \sim_{M} w^{'}\)) gives \(f(w^{'}) = f(v^{'})\).
\end{proofstep}
Use the hypothesis \(w = v_{1}\).
Reversing the equality, we obtain \(v_{1} = w\).
Use the hypothesis \(w \sim_{M} w^{'}\).
Use the hypothesis \(f(w) = f(v_{1})\).
Use the hypothesis \(f(w) = f(v_{1})\).
Reversing the equality, we obtain \(f(v_{1}) = f(w)\).
Applying the hypothesis \(\operatorname{DFun\,Like\mathord{:}coe}\left(f\right)\text{ is injective}\), we obtain \(v_{1} = w\).
Substituting equals, we obtain \(\text{the proof of }(v_{1} = w)\text{ agrees under transport with the proof of }(v_{1} = v_{1}) \implies f(w^{'}) = f(v^{'})\).
\end{proofstep}
By definition, we obtain \(w = w\).
Use the hypothesis \(f(w) = f(v_{1})\).
Reversing the equality, we obtain \(f(v_{1}) = f(w)\).
Applying the hypothesis \(\operatorname{DFun\,Like\mathord{:}coe}\left(f\right)\text{ is injective}\), we obtain \(v_{1} = w\).
By reflexivity, we obtain \(\text{the proof of }(v_{1} = w)\text{ agrees under transport with the proof of }(v_{1} = w)\).
Substituting equals, we obtain \(f(w^{'}) = f(v^{'})\).
Use the hypothesis \(f(w^{'}) = y\).
Substituting equals, we obtain \(y = f(v^{'})\).
\end{proofstep}
Using the two conjuncts, we obtain \(y = f(v^{'})\).
\end{proofstep}
Using the two conjuncts, we obtain \(y = f(v^{'})\).
\end{proofstep}
Using the existential hypothesis, we obtain \(y = f(v^{'})\).
\end{proofstep}
Using the existential hypothesis, we obtain \(y = f(v^{'})\).
\end{proofstep}
Combining the two conclusions, we obtain \(f(v_{1}) \sim_{\operatorname{Simple\,Graph\mathord{:}Subgraph\mathord{:}map}(f,\allowbreak{} M)} f(v^{'}) \land (\forall y\, :\, W,\; f(v_{1}) \sim_{\operatorname{Simple\,Graph\mathord{:}Subgraph\mathord{:}map}(f,\allowbreak{} M)} y \implies y = f(v^{'}))\).
This witness proves existence. Hence \(\exists x\, :\, W,\; f(v_{1}) \sim_{\operatorname{Simple\,Graph\mathord{:}Subgraph\mathord{:}map}(f,\allowbreak{} M)} x \land (\forall y\, :\, W,\; f(v_{1}) \sim_{\operatorname{Simple\,Graph\mathord{:}Subgraph\mathord{:}map}(f,\allowbreak{} M)} y \implies y = x)\).
\end{proofstep}
Using the existential hypothesis, we obtain \(\exists! w\, :\, W,\; f(v_{1}) \sim_{\operatorname{Simple\,Graph\mathord{:}Subgraph\mathord{:}map}(f,\allowbreak{} M)} w\).
Use the hypothesis \(f(v_{1}) = v\).
Substituting equals, we obtain \(\exists! w\, :\, W,\; v \sim_{\operatorname{Simple\,Graph\mathord{:}Subgraph\mathord{:}map}(f,\allowbreak{} M)} w\).
\end{proofstep}
Using the two conjuncts, we obtain \(\exists! w\, :\, W,\; v \sim_{\operatorname{Simple\,Graph\mathord{:}Subgraph\mathord{:}map}(f,\allowbreak{} M)} w\).
\end{proofstep}
Using the existential hypothesis, we obtain \(\exists! w\, :\, W,\; v \sim_{\operatorname{Simple\,Graph\mathord{:}Subgraph\mathord{:}map}(f,\allowbreak{} M)} w\).
\end{proof}

\begin{theorem}[{Is matching map}]\label{theorem:15}
Let \(V\) and \(W\) be types. Let \(G\) be a simple graph on \(V\). Let \(G^{'}\) be a simple graph on \(W\). Let \(M\) be a subgraph of \(G\). Let \(f\) be a graph isomorphism from \(G\) to \(G^{'}\). Then
\begin{center}
\(\displaystyle \operatorname{Simple\,Graph\mathord{:}Subgraph\mathord{:}map}(\operatorname{to\,Hom}(f),\allowbreak{} M)\text{ is a matching} \iff M\text{ is a matching}.\)
\end{center}
\end{theorem}

\begin{proof}
Consider the following case.
\begin{proofstep}
Assume \(\operatorname{Simple\,Graph\mathord{:}Subgraph\mathord{:}map}(\operatorname{to\,Hom}(f),\allowbreak{} M)\text{ is a matching}\).
By \paperlink{premise:52656c49736f2e696e6a656374697665}{Prerequisite~P\ref{premise:52656c49736f2e696e6a656374697665}} with \(e = f^{-1}\), we obtain \(\operatorname{DFun\,Like\mathord{:}coe}\left(f^{-1}\right)\text{ is injective}\).
By \paperlink{theorem:14}{Lemma~\ref{theorem:14}} with \(\operatorname{to\,Hom}\left(f^{-1}\right)\) in place of its parameter \(f\), we obtain \(\operatorname{Simple\,Graph\mathord{:}Subgraph\mathord{:}map}(\operatorname{to\,Hom}(f^{-1}),\allowbreak{} \operatorname{Simple\,Graph\mathord{:}Subgraph\mathord{:}map}(\operatorname{to\,Hom}(f),\allowbreak{} M))\text{ is a matching}\).
Rewriting with \paperlink{premise:53696d706c6547726170682e53756267726170682e6d61705f636f6d70}{Prerequisite~P\ref{premise:53696d706c6547726170682e53756267726170682e6d61705f636f6d70}}, \paperlink{premise:53696d706c6547726170682e49736f2e73796d6d5f746f486f6d5f636f6d705f746f486f6d}{Prerequisite~P\ref{premise:53696d706c6547726170682e49736f2e73796d6d5f746f486f6d5f636f6d705f746f486f6d}} and \paperlink{premise:53696d706c6547726170682e53756267726170682e6d61705f6964}{Prerequisite~P\ref{premise:53696d706c6547726170682e53756267726170682e6d61705f6964}} in the preceding result gives \(M\text{ is a matching}\).
\end{proofstep}
By \paperlink{premise:52656c49736f2e696e6a656374697665}{Prerequisite~P\ref{premise:52656c49736f2e696e6a656374697665}} with \(e = f\), we obtain \(\operatorname{DFun\,Like\mathord{:}coe}\left(f\right)\text{ is injective}\).
By \paperlink{theorem:14}{Lemma~\ref{theorem:14}} with \(\operatorname{to\,Hom}\left(f\right)\) in place of its parameter \(f\), we obtain \(M\text{ is a matching} \implies \operatorname{Simple\,Graph\mathord{:}Subgraph\mathord{:}map}(\operatorname{to\,Hom}(f),\allowbreak{} M)\text{ is a matching}\).
Together the two implications give the equivalence. Hence \(\operatorname{Simple\,Graph\mathord{:}Subgraph\mathord{:}map}(\operatorname{to\,Hom}(f),\allowbreak{} M)\text{ is a matching} \iff M\text{ is a matching}\).
\end{proof}

\begin{theorem}[{Support equals verts}]\label{theorem:16}
Let \(V\) be a type. Let \(G\) be a simple graph on \(V\). Let \(M\) be a subgraph of \(G\). Assume \(M\text{ is a matching}\). Then \(\operatorname{Simple\,Graph\mathord{:}Subgraph\mathord{:}support}\left(M\right) = V(M)\).
\end{theorem}

\begin{proof}
By \paperlink{premise:5365742e696e73744973416e746973796d6d537562736574}{Prerequisite~P\ref{premise:5365742e696e73744973416e746973796d6d537562736574}}, we obtain \(\operatorname{Is\,Antisymm}(\mathcal{P}(V),\allowbreak{} \mathit{x1} \mapsto \mathit{x2} \mapsto \mathit{x1} \subseteq \mathit{x2})\).
By \paperlink{premise:53696d706c6547726170682e53756267726170682e737570706f72745f7375627365745f7665727473}{Prerequisite~P\ref{premise:53696d706c6547726170682e53756267726170682e737570706f72745f7375627365745f7665727473}} with \(H = M\), we obtain \(\operatorname{Simple\,Graph\mathord{:}Subgraph\mathord{:}support}\left(M\right) \subseteq V(M)\).
For the required implication, argue as follows.
\begin{proofstep}
Fix \(v\) in \(V\).
Assume \(v \in V(M)\).
Use the hypothesis \(v \in V(M)\).
Applying the hypothesis \(M\text{ is a matching}\), we obtain \(\exists! w\, :\, V,\; v \sim_{M} w\).
Consider the following case.
\begin{proofstep}
Fix \(w\) in \(V\).
Assume \(v \sim_{M} w \land \left(\forall y\, :\, V,\; v \sim_{M} y \implies y = w\right)\).
Use the hypothesis \(v \sim_{M} w \land \left(\forall y\, :\, V,\; v \sim_{M} y \implies y = w\right)\).
Consider the following case.
\begin{proofstep}
Assume \(v \sim_{M} w\).
Assume \(\forall y\, :\, V,\; v \sim_{M} y \implies y = w\).
Choose the witness \(w\).
Use the hypothesis \(v \sim_{M} w\).
This witness proves existence. Hence \(\exists y\, :\, V,\; v \sim_{M} y\).
\end{proofstep}
Using the two conjuncts, we obtain \(v \in \operatorname{Simple\,Graph\mathord{:}Subgraph\mathord{:}support}\left(M\right)\).
\end{proofstep}
Using the existential hypothesis, we obtain \(v \in \operatorname{Simple\,Graph\mathord{:}Subgraph\mathord{:}support}\left(M\right)\).
\end{proofstep}
By \paperlink{premise:4861735375627365742e5375627365742e616e746973796d6d}{Prerequisite~P\ref{premise:4861735375627365742e5375627365742e616e746973796d6d}} with \(\alpha = \mathcal{P}(V)\), we obtain \(\operatorname{Simple\,Graph\mathord{:}Subgraph\mathord{:}support}\left(M\right) = V(M)\).
\end{proof}

\begin{lemma}[{Is matching if and only if forall degree}]\label{theorem:17}
Let \(V\) be a type. Let \(G\) be a simple graph on \(V\). Let \(M\) be a subgraph of \(G\). Fix a structure \(\mathit{inst}\) of type \(\Pi_{v\, :\, V}\, \operatorname{Fintype}(\operatorname{Simple\,Graph\mathord{:}Subgraph\mathord{:}neighbor\,Set}(M,\allowbreak{} v))\). Then
\begin{center}
\(\displaystyle M\text{ is a matching} \iff (\forall v\, :\, V,\; v \in V(M) \implies \deg_{M}(v) = 1).\)
\end{center}
\end{lemma}

\begin{proof}
For the required implication, argue as follows.
\begin{proofstep}
Fix \(v\) in \(V\).
By definition, we obtain \(\left(v \in V(M)\right) = \left(v \in V(M)\right)\).
By \paperlink{premise:53696d706c6547726170682e53756267726170682e6465677265655f65715f6f6e655f6966665f756e697175655f61646a}{Prerequisite~P\ref{premise:53696d706c6547726170682e53756267726170682e6465677265655f65715f6f6e655f6966665f756e697175655f61646a}}, we obtain \(\deg_{M}(v) = 1 \iff \left(\exists! w\, :\, V,\; v \sim_{M} w\right)\).
By propositional extensionality, we obtain \(\left(\deg_{M}(v) = 1\right) = \left(\exists! w\, :\, V,\; v \sim_{M} w\right)\).
By \paperlink{premise:696d706c6965735f636f6e6772}{Prerequisite~P\ref{premise:696d706c6965735f636f6e6772}} with \(p_{1} = v \in V(M)\) and \(p_{2} = v \in V(M)\) and \(q_{1} = \deg_{M}(v) = 1\) and \(q_{2} = \exists! w\, :\, V,\; v \sim_{M} w\), we obtain \((v \in V(M) \implies \deg_{M}(v) = 1) = (v \in V(M) \implies (\exists! w\, :\, V,\; v \sim_{M} w))\).
\end{proofstep}
By \paperlink{premise:666f72616c6c5f636f6e6772}{Prerequisite~P\ref{premise:666f72616c6c5f636f6e6772}}, we obtain \((\forall a\, :\, V,\; a \in V(M) \implies \deg_{M}(a) = 1) = (\forall a\, :\, V,\; a \in V(M) \implies (\exists! w\, :\, V,\; a \sim_{M} w))\).
Applying the same function to equal arguments, we obtain \(((\forall v\, :\, V,\; v \in V(M) \implies (\exists! w\, :\, V,\; v \sim_{M} w)) \iff (\forall a\, :\, V,\; a \in V(M) \implies \deg_{M}(a) = 1)) = ((\forall v\, :\, V,\; v \in V(M) \implies (\exists! w\, :\, V,\; v \sim_{M} w)) \iff (\forall a\, :\, V,\; a \in V(M) \implies (\exists! w\, :\, V,\; a \sim_{M} w)))\).
By \paperlink{premise:6966665f73656c66}{Prerequisite~P\ref{premise:6966665f73656c66}} with \(p = \forall v\, :\, V,\; v \in V(M) \implies \left(\exists! w\, :\, V,\; v \sim_{M} w\right)\), we obtain \(((\forall v\, :\, V,\; v \in V(M) \implies (\exists! w\, :\, V,\; v \sim_{M} w)) \iff (\forall v\, :\, V,\; v \in V(M) \implies (\exists! w\, :\, V,\; v \sim_{M} w))) = \mathrm{True}\).
Chaining the equalities, we obtain \(((\forall v\, :\, V,\; v \in V(M) \implies (\exists! w\, :\, V,\; v \sim_{M} w)) \iff (\forall a\, :\, V,\; a \in V(M) \implies \deg_{M}(a) = 1)) = \mathrm{True}\).
By \paperlink{premise:6f665f65715f74727565}{Prerequisite~P\ref{premise:6f665f65715f74727565}}, we obtain \(M\text{ is a matching} \iff \left(\forall v\, :\, V,\; v \in V(M) \implies \deg_{M}(v) = 1\right)\).
\end{proof}

\begin{lemma}[{Even cardinality}]\label{theorem:18}
Let \(V\) be a type. Let \(G\) be a simple graph on \(V\). Let \(M\) be a subgraph of \(G\). Fix a finite enumeration of \(V(M)\). Assume \(M\text{ is a matching}\). Then \(\operatorname{Even}\left(\left|\operatorname{as\,finite\,set}(V(M))\right|\right)\).
\end{lemma}

\begin{proof}
Choose the witness \(\left|E(\operatorname{graph}(M))\right|\).
The substitution uses the following argument.
\begin{proofstep}
Fix \(x\) in \(V(M)\).
Assume \(x \in V(M)\).
By \paperlink{premise:53696d706c6547726170682e53756267726170682e636f655f646567726565}{Prerequisite~P\ref{premise:53696d706c6547726170682e53756267726170682e636f655f646567726565}} with \(G^{'} = M\) and \(v = x\), we obtain \(\deg_{\operatorname{graph}(M)}(x) = \deg_{M}(\operatorname{val}(x))\).
By \paperlink{theorem:17}{Lemma~\ref{theorem:17}}, we obtain \(M\text{ is a matching} \iff \left(\forall v\, :\, V,\; v \in V(M) \implies \deg_{M}(v) = 1\right)\).
By propositional extensionality, we obtain \(\left(M\text{ is a matching}\right) = \left(\forall v\, :\, V,\; v \in V(M) \implies \deg_{M}(v) = 1\right)\).
Applying the same function to equal arguments, we obtain \((x_{1} \mapsto x_{1})(M\text{ is a matching}) = (x_{1} \mapsto x_{1})(\forall v\, :\, V,\; v \in V(M) \implies \deg_{M}(v) = 1)\).
Use the hypothesis \(M\text{ is a matching}\).
Rewriting with \paperlink{premise:537562747970652e636f655f70726f70}{Prerequisite~P\ref{premise:537562747970652e636f655f70726f70}} gives \(\operatorname{val}(x) \in V(M)\).
Substituting along the equality, we obtain \(\deg_{M}(\operatorname{val}(x)) = 1\).
Chaining the equalities, we obtain \(\deg_{\operatorname{graph}(M)}(x) = 1\).
\end{proofstep}
Rewriting with \paperlink{premise:5365742e746f46696e7365745f63617264}{Prerequisite~P\ref{premise:5365742e746f46696e7365745f63617264}}, \paperlink{premise:46696e7365742e73756d5f636f6e6772}{Prerequisite~P\ref{premise:46696e7365742e73756d5f636f6e6772}}, \paperlink{premise:46696e7365742e73756d5f636f6e7374}{Prerequisite~P\ref{premise:46696e7365742e73756d5f636f6e7374}}, \paperlink{premise:6d756c5f6f6e65}{Prerequisite~P\ref{premise:6d756c5f6f6e65}} and the identity established above gives \(|\operatorname{as\,finite\,set}(V(M))| = \sum_{v\in V(M)}(\deg_{\operatorname{graph}(M)}(v))\).
By the preceding argument, we obtain \(|\operatorname{as\,finite\,set}(V(M))| = \sum_{v\in V(M)}(\deg_{\operatorname{graph}(M)}(v))\).
Rewriting with \paperlink{premise:65715f6f665f686571}{Prerequisite~P\ref{premise:65715f6f665f686571}}, \paperlink{premise:46696e7365742e73756d5f636f6e6772}{Prerequisite~P\ref{premise:46696e7365742e73756d5f636f6e6772}}, \paperlink{premise:46696e747970652e696e73744661737453756273696e676c65746f6e}{Prerequisite~P\ref{premise:46696e747970652e696e73744661737453756273696e676c65746f6e}} and substitution of equals in the preceding result gives \(|\operatorname{as\,finite\,set}(V(M))| = \sum_{v\in V(M)}(\deg_{\operatorname{graph}(M)}(v))\).
Rewriting with \paperlink{premise:53696d706c6547726170682e73756d5f646567726565735f65715f74776963655f636172645f6564676573}{Prerequisite~P\ref{premise:53696d706c6547726170682e73756d5f646567726565735f65715f74776963655f636172645f6564676573}} in the preceding result gives \(\left|\operatorname{as\,finite\,set}(V(M))\right| = 2 \cdot \left|E(\operatorname{graph}(M))\right|\).
Rewriting with \paperlink{premise:74776f5f6d756c}{Prerequisite~P\ref{premise:74776f5f6d756c}} in the preceding result gives \(|\operatorname{as\,finite\,set}(V(M))| = |E(\operatorname{graph}(M))| + |E(\operatorname{graph}(M))|\).
This witness proves existence. Hence \(\operatorname{Even}\left(\left|\operatorname{as\,finite\,set}(V(M))\right|\right)\).
\end{proof}

\begin{lemma}[{Is perfect matching if and only if}]\label{theorem:19}
Let \(V\) be a type. Let \(G\) be a simple graph on \(V\). Let \(M\) be a subgraph of \(G\). Then
\begin{center}
\(\displaystyle M\text{ is a perfect matching} \iff (\forall v\, :\, V,\; \exists! w\, :\, V,\; v \sim_{M} w).\)
\end{center}
\end{lemma}

\begin{proof}
Consider the following case.
\begin{proofstep}
Assume \(M\text{ is a perfect matching}\).
Use the hypothesis \(M\text{ is a perfect matching}\).
Consider the following case.
\begin{proofstep}
Assume \(M\text{ is a matching}\).
Assume \(\operatorname{Is\,Spanning}\left(M\right)\).
Fix \(v\) in \(V\).
Applying the hypothesis \(\operatorname{Is\,Spanning}\left(M\right)\), we obtain \(v \in V(M)\).
Applying the hypothesis \(M\text{ is a matching}\), we obtain \(\exists! w\, :\, V,\; v \sim_{M} w\).
\end{proofstep}
Using the two conjuncts, we obtain \(\forall v\, :\, V,\; \exists! w\, :\, V,\; v \sim_{M} w\).
\end{proofstep}
Consider the following case.
\begin{proofstep}
Assume \(\forall v\, :\, V,\; \exists! w\, :\, V,\; v \sim_{M} w\).
For the required implication, argue as follows.
\begin{proofstep}
Fix \(v\) in \(V\).
Assume \(v \in V(M)\).
Applying the hypothesis \(\forall v\, :\, V,\; \exists! w\, :\, V,\; v \sim_{M} w\), we obtain \(\exists! w\, :\, V,\; v \sim_{M} w\).
\end{proofstep}
For the required implication, argue as follows.
\begin{proofstep}
Fix \(v\) in \(V\).
Applying the hypothesis \(\forall v\, :\, V,\; \exists! w\, :\, V,\; v \sim_{M} w\), we obtain \(\exists! w\, :\, V,\; v \sim_{M} w\).
Consider the following case.
\begin{proofstep}
Fix \(w\) in \(V\).
Assume \(v \sim_{M} w \land \left(\forall y\, :\, V,\; v \sim_{M} y \implies y = w\right)\).
Use the hypothesis \(v \sim_{M} w \land \left(\forall y\, :\, V,\; v \sim_{M} y \implies y = w\right)\).
Consider the following case.
\begin{proofstep}
Assume \(v \sim_{M} w\).
Assume \(\forall y\, :\, V,\; v \sim_{M} y \implies y = w\).
Use the hypothesis \(v \sim_{M} w\).
By \paperlink{premise:53696d706c6547726170682e53756267726170682e656467655f76657274}{Prerequisite~P\ref{premise:53696d706c6547726170682e53756267726170682e656467655f76657274}} with \(\mathit{self} = M\), we obtain \(v \in V(M)\).
\end{proofstep}
Using the two conjuncts, we obtain \(v \in V(M)\).
\end{proofstep}
Using the existential hypothesis, we obtain \(v \in V(M)\).
\end{proofstep}
Combining the two conclusions, we obtain \(M\text{ is a perfect matching}\).
\end{proofstep}
Together the two implications give the equivalence. Hence \(M\text{ is a perfect matching} \iff \left(\forall v\, :\, V,\; \exists! w\, :\, V,\; v \sim_{M} w\right)\).
\end{proof}

\begin{theorem}[{Is perfect matching if and only if forall degree}]\label{theorem:20}
Let \(V\) be a type. Let \(G\) be a simple graph on \(V\). Let \(M\) be a subgraph of \(G\). Fix a structure \(\mathit{inst}\) of type \(\Pi_{v\, :\, V}\, \operatorname{Fintype}(\operatorname{Simple\,Graph\mathord{:}Subgraph\mathord{:}neighbor\,Set}(M,\allowbreak{} v))\). Then \(M\text{ is a perfect matching} \iff \left(\forall v\, :\, V,\; \deg_{M}(v) = 1\right)\).
\end{theorem}

\begin{proof}
The substitution uses the following argument.
\begin{proofstep}
Fix \(v\) in \(V\).
By \paperlink{premise:53696d706c6547726170682e53756267726170682e6465677265655f65715f6f6e655f6966665f756e697175655f61646a}{Prerequisite~P\ref{premise:53696d706c6547726170682e53756267726170682e6465677265655f65715f6f6e655f6966665f756e697175655f61646a}}, we obtain \(\deg_{M}(v) = 1 \iff \left(\exists! w\, :\, V,\; v \sim_{M} w\right)\).
By propositional extensionality, we obtain \(\left(\deg_{M}(v) = 1\right) = \left(\exists! w\, :\, V,\; v \sim_{M} w\right)\).
\end{proofstep}
Rewriting with \paperlink{theorem:19}{Lemma~\ref{theorem:19}}, \paperlink{premise:666f72616c6c5f636f6e6772}{Prerequisite~P\ref{premise:666f72616c6c5f636f6e6772}}, \paperlink{premise:6966665f73656c66}{Prerequisite~P\ref{premise:6966665f73656c66}} and the identity established above gives \(M\text{ is a perfect matching} \iff \left(\forall v\, :\, V,\; \deg_{M}(v) = 1\right)\).
\end{proof}

\begin{corollary}[{Even cardinality}]\label{theorem:21}
Let \(V\) be a type. Let \(G\) be a simple graph on \(V\). Let \(M\) be a subgraph of \(G\). Fix a finite enumeration of \(V\). Assume \(M\text{ is a perfect matching}\). Then \(\operatorname{Even}\left(\left|V\right|\right)\).
\end{corollary}

\begin{proof}
Use the hypothesis \(M\text{ is a perfect matching}\).
Taking the second conjunct, we obtain \(\operatorname{Is\,Spanning}\left(M\right)\).
By \paperlink{premise:53696d706c6547726170682e53756267726170682e49735370616e6e696e672e636172645f7665727473}{Prerequisite~P\ref{premise:53696d706c6547726170682e53756267726170682e49735370616e6e696e672e636172645f7665727473}}, we obtain \(\left|\operatorname{as\,finite\,set}(V(M))\right| = \left|V\right|\).
Applying the same function to equal arguments, we obtain \(\operatorname{Even}(|\operatorname{as\,finite\,set}(V(M))|) = \operatorname{Even}(|V|)\).
Use the hypothesis \(M\text{ is a perfect matching}\).
Taking the first conjunct, we obtain \(M\text{ is a matching}\).
By \paperlink{theorem:18}{Lemma~\ref{theorem:18}}, we obtain \(\operatorname{Even}\left(\left|\operatorname{as\,finite\,set}(V(M))\right|\right)\).
Substituting along the equality, we obtain \(\operatorname{Even}\left(\left|V\right|\right)\).
\end{proof}

\begin{lemma}[{Induce connected component}]\label{theorem:22}
Let \(V\) be a type. Let \(G\) be a simple graph on \(V\). Let \(M\) be a subgraph of \(G\). Assume \(M\text{ is a matching}\). Let \(c\) be an element of \(\operatorname{Connected\,Component}\left(G\right)\). Then
\begin{center}
\(\displaystyle \operatorname{Simple\,Graph\mathord{:}Subgraph\mathord{:}induce}(M,\allowbreak{} V(M) \cap \operatorname{supp}(c))\text{ is a matching}.\)
\end{center}
\end{lemma}

\begin{proof}
Fix \(v\) in \(V\).
Assume \(v \in V(\operatorname{Simple\,Graph\mathord{:}Subgraph\mathord{:}induce}(M,\allowbreak{} V(M) \cap \operatorname{supp}(c)))\).
Use the hypothesis \(v \in V(\operatorname{Simple\,Graph\mathord{:}Subgraph\mathord{:}induce}(M,\allowbreak{} V(M) \cap \operatorname{supp}(c)))\).
Consider the following case.
\begin{proofstep}
Assume \(v \in V(M)\).
Assume \(v \in \operatorname{supp}\left(c\right)\).
Use the hypothesis \(v \in V(M)\).
Applying the hypothesis \(M\text{ is a matching}\), we obtain \(\exists! w\, :\, V,\; v \sim_{M} w\).
Consider the following case.
\begin{proofstep}
Fix \(w\) in \(V\).
Assume \(v \sim_{M} w \land \left(\forall y\, :\, V,\; v \sim_{M} y \implies y = w\right)\).
Use the hypothesis \(v \sim_{M} w \land \left(\forall y\, :\, V,\; v \sim_{M} y \implies y = w\right)\).
Consider the following case.
\begin{proofstep}
Assume \(v \sim_{M} w\).
Assume \(\forall y\, :\, V,\; v \sim_{M} y \implies y = w\).
Choose the witness \(w\).
The substitution uses the following argument.
\begin{proofstep}
Fix \(y\) in \(V\).
By \paperlink{premise:5365742e6d656d5f696e7465725f696666}{Prerequisite~P\ref{premise:5365742e6d656d5f696e7465725f696666}} with \(x = v\) and \(a = V(M)\) and \(b = \operatorname{supp}\left(\operatorname{connected\,Component\,Mk}\left(G, v\right)\right)\), we obtain \(v \in V(M) \cap \operatorname{supp}(\operatorname{connected\,Component\,Mk}(G,\allowbreak{} v)) \iff v \in V(M) \land v \in \operatorname{supp}(\operatorname{connected\,Component\,Mk}(G,\allowbreak{} v))\).
By propositional extensionality, we obtain \((v \in V(M) \cap \operatorname{supp}(\operatorname{connected\,Component\,Mk}(G,\allowbreak{} v))) = (v \in V(M) \land v \in \operatorname{supp}(\operatorname{connected\,Component\,Mk}(G,\allowbreak{} v)))\).
Use the hypothesis \(v \in V(M)\).
By \paperlink{premise:65715f74727565}{Prerequisite~P\ref{premise:65715f74727565}}, we obtain \(\left(v \in V(M)\right) = \mathrm{True}\).
Applying the same function to equal arguments, we obtain \(\left(y_{1} \mapsto v \in V(M) \land y_{1}\right) = \left(y_{1} \mapsto \mathrm{True} \land y_{1}\right)\).
By \paperlink{premise:53696d706c6547726170682e436f6e6e6563746564436f6d706f6e656e742e6d656d5f737570705f696666}{Prerequisite~P\ref{premise:53696d706c6547726170682e436f6e6e6563746564436f6d706f6e656e742e6d656d5f737570705f696666}} with \(C = \operatorname{connected\,Component\,Mk}\left(G, v\right)\), we obtain \(v \in \operatorname{supp}(\operatorname{connected\,Component\,Mk}(G,\allowbreak{} v)) \iff \operatorname{connected\,Component\,Mk}(G,\allowbreak{} v) = \operatorname{connected\,Component\,Mk}(G,\allowbreak{} v)\).
By propositional extensionality, we obtain \((v \in \operatorname{supp}(\operatorname{connected\,Component\,Mk}(G,\allowbreak{} v))) = (\operatorname{connected\,Component\,Mk}(G,\allowbreak{} v) = \operatorname{connected\,Component\,Mk}(G,\allowbreak{} v))\).
By \paperlink{premise:65715f73656c66}{Prerequisite~P\ref{premise:65715f73656c66}} with \(\alpha = \operatorname{Connected\,Component}\left(G\right)\) and \(a = \operatorname{connected\,Component\,Mk}\left(G, v\right)\), we obtain \((\operatorname{connected\,Component\,Mk}(G,\allowbreak{} v) = \operatorname{connected\,Component\,Mk}(G,\allowbreak{} v)) = \mathrm{True}\).
Chaining the equalities, we obtain \((v \in \operatorname{supp}(\operatorname{connected\,Component\,Mk}(G,\allowbreak{} v))) = \mathrm{True}\).
By congruence, we obtain \((v \in V(M) \land v \in \operatorname{supp}(\operatorname{connected\,Component\,Mk}(G,\allowbreak{} v))) = (\mathrm{True} \land \mathrm{True})\).
By \paperlink{premise:616e645f73656c66}{Prerequisite~P\ref{premise:616e645f73656c66}} with \(p = \mathrm{True}\), we obtain \(\left(\mathrm{True} \land \mathrm{True}\right) = \mathrm{True}\).
Chaining the equalities, we obtain \((v \in V(M) \land v \in \operatorname{supp}(\operatorname{connected\,Component\,Mk}(G,\allowbreak{} v))) = \mathrm{True}\).
Chaining the equalities, we obtain \((v \in V(M) \cap \operatorname{supp}(\operatorname{connected\,Component\,Mk}(G,\allowbreak{} v))) = \mathrm{True}\).
Applying the same function to equal arguments, we obtain \((y_{1} \mapsto v \in V(M) \cap \operatorname{supp}(\operatorname{connected\,Component\,Mk}(G,\allowbreak{} v)) \land y_{1}) = (y_{1} \mapsto \mathrm{True} \land y_{1})\).
By \paperlink{premise:5365742e6d656d5f696e7465725f696666}{Prerequisite~P\ref{premise:5365742e6d656d5f696e7465725f696666}} with \(x = y\) and \(a = V(M)\) and \(b = \operatorname{supp}\left(\operatorname{connected\,Component\,Mk}\left(G, v\right)\right)\), we obtain \(y \in V(M) \cap \operatorname{supp}(\operatorname{connected\,Component\,Mk}(G,\allowbreak{} v)) \iff y \in V(M) \land y \in \operatorname{supp}(\operatorname{connected\,Component\,Mk}(G,\allowbreak{} v))\).
By propositional extensionality, we obtain \((y \in V(M) \cap \operatorname{supp}(\operatorname{connected\,Component\,Mk}(G,\allowbreak{} v))) = (y \in V(M) \land y \in \operatorname{supp}(\operatorname{connected\,Component\,Mk}(G,\allowbreak{} v)))\).
By \paperlink{premise:53696d706c6547726170682e436f6e6e6563746564436f6d706f6e656e742e6d656d5f737570705f696666}{Prerequisite~P\ref{premise:53696d706c6547726170682e436f6e6e6563746564436f6d706f6e656e742e6d656d5f737570705f696666}} with \(C = \operatorname{connected\,Component\,Mk}\left(G, v\right)\) and \(y\) in place of its parameter \(v\), we obtain \(y \in \operatorname{supp}(\operatorname{connected\,Component\,Mk}(G,\allowbreak{} v)) \iff \operatorname{connected\,Component\,Mk}(G,\allowbreak{} y) = \operatorname{connected\,Component\,Mk}(G,\allowbreak{} v)\).
By propositional extensionality, we obtain \((y \in \operatorname{supp}(\operatorname{connected\,Component\,Mk}(G,\allowbreak{} v))) = (\operatorname{connected\,Component\,Mk}(G,\allowbreak{} y) = \operatorname{connected\,Component\,Mk}(G,\allowbreak{} v))\).
By \paperlink{premise:53696d706c6547726170682e436f6e6e6563746564436f6d706f6e656e742e6571}{Prerequisite~P\ref{premise:53696d706c6547726170682e436f6e6e6563746564436f6d706f6e656e742e6571}}, we obtain \(\operatorname{connected\,Component\,Mk}(G,\allowbreak{} y) = \operatorname{connected\,Component\,Mk}(G,\allowbreak{} v) \iff y\text{ can reach }v\text{ in }G\).
By propositional extensionality, we obtain \((\operatorname{connected\,Component\,Mk}(G,\allowbreak{} y) = \operatorname{connected\,Component\,Mk}(G,\allowbreak{} v)) = (y\text{ can reach }v\text{ in }G)\).
Chaining the equalities, we obtain \((y \in \operatorname{supp}(\operatorname{connected\,Component\,Mk}(G,\allowbreak{} v))) = (y\text{ can reach }v\text{ in }G)\).
Applying the same function to equal arguments, we obtain \((y \in V(M) \land y \in \operatorname{supp}(\operatorname{connected\,Component\,Mk}(G,\allowbreak{} v))) = (y \in V(M) \land y\text{ can reach }v\text{ in }G)\).
Chaining the equalities, we obtain \((y \in V(M) \cap \operatorname{supp}(\operatorname{connected\,Component\,Mk}(G,\allowbreak{} v))) = (y \in V(M) \land y\text{ can reach }v\text{ in }G)\).
Applying the same function to equal arguments, we obtain \((y \in V(M) \cap \operatorname{supp}(\operatorname{connected\,Component\,Mk}(G,\allowbreak{} v)) \land v \sim_{M} y) = ((y \in V(M) \land y\text{ can reach }v\text{ in }G) \land v \sim_{M} y)\).
By congruence, we obtain \((v \in V(M) \cap \operatorname{supp}(\operatorname{connected\,Component\,Mk}(G,\allowbreak{} v)) \land (y \in V(M) \cap \operatorname{supp}(\operatorname{connected\,Component\,Mk}(G,\allowbreak{} v)) \land v \sim_{M} y)) = (\mathrm{True} \land ((y \in V(M) \land y\text{ can reach }v\text{ in }G) \land v \sim_{M} y))\).
By \paperlink{premise:747275655f616e64}{Prerequisite~P\ref{premise:747275655f616e64}} with \(p = \left(y \in V(M) \land y\text{ can reach }v\text{ in }G\right) \land v \sim_{M} y\), we obtain \((\mathrm{True} \land ((y \in V(M) \land y\text{ can reach }v\text{ in }G) \land v \sim_{M} y)) = ((y \in V(M) \land y\text{ can reach }v\text{ in }G) \land v \sim_{M} y)\).
Chaining the equalities, we obtain \((v \in V(M) \cap \operatorname{supp}(\operatorname{connected\,Component\,Mk}(G,\allowbreak{} v)) \land (y \in V(M) \cap \operatorname{supp}(\operatorname{connected\,Component\,Mk}(G,\allowbreak{} v)) \land v \sim_{M} y)) = ((y \in V(M) \land y\text{ can reach }v\text{ in }G) \land v \sim_{M} y)\).
By definition, we obtain \(\left(y = w\right) = \left(y = w\right)\).
By \paperlink{premise:696d706c6965735f636f6e6772}{Prerequisite~P\ref{premise:696d706c6965735f636f6e6772}} with \(p_{1} = v \in V(M) \cap \operatorname{supp}(\operatorname{connected\,Component\,Mk}(G,\allowbreak{} v)) \land (y \in V(M) \cap \operatorname{supp}(\operatorname{connected\,Component\,Mk}(G,\allowbreak{} v)) \land v \sim_{M} y)\) and \(p_{2} = \left(y \in V(M) \land y\text{ can reach }v\text{ in }G\right) \land v \sim_{M} y\) and \(q_{1} = y = w\) and \(q_{2} = y = w\), we obtain \((v \in V(M) \cap \operatorname{supp}(\operatorname{connected\,Component\,Mk}(G,\allowbreak{} v)) \land (y \in V(M) \cap \operatorname{supp}(\operatorname{connected\,Component\,Mk}(G,\allowbreak{} v)) \land v \sim_{M} y) \implies y = w) = ((y \in V(M) \land y\text{ can reach }v\text{ in }G) \land v \sim_{M} y \implies y = w)\).
By \paperlink{premise:616e645f696d70}{Prerequisite~P\ref{premise:616e645f696d70}}, we obtain \((y \in V(M) \land y\text{ can reach }v\text{ in }G) \land v \sim_{M} y \implies y = w \iff y \in V(M) \land y\text{ can reach }v\text{ in }G \implies v \sim_{M} y \implies y = w\).
By propositional extensionality, we obtain \(((y \in V(M) \land y\text{ can reach }v\text{ in }G) \land v \sim_{M} y \implies y = w) = (y \in V(M) \land y\text{ can reach }v\text{ in }G \implies v \sim_{M} y \implies y = w)\).
Chaining the equalities, we obtain \((v \in V(M) \cap \operatorname{supp}(\operatorname{connected\,Component\,Mk}(G,\allowbreak{} v)) \land (y \in V(M) \cap \operatorname{supp}(\operatorname{connected\,Component\,Mk}(G,\allowbreak{} v)) \land v \sim_{M} y) \implies y = w) = (y \in V(M) \land y\text{ can reach }v\text{ in }G \implies v \sim_{M} y \implies y = w)\).
By \paperlink{premise:616e645f696d70}{Prerequisite~P\ref{premise:616e645f696d70}}, we obtain \(y \in V(M) \land y\text{ can reach }v\text{ in }G \implies v \sim_{M} y \implies y = w \iff y \in V(M) \implies y\text{ can reach }v\text{ in }G \implies v \sim_{M} y \implies y = w\).
By propositional extensionality, we obtain \((y \in V(M) \land y\text{ can reach }v\text{ in }G \implies v \sim_{M} y \implies y = w) = (y \in V(M) \implies y\text{ can reach }v\text{ in }G \implies v \sim_{M} y \implies y = w)\).
Chaining the equalities, we obtain \((v \in V(M) \cap \operatorname{supp}(\operatorname{connected\,Component\,Mk}(G,\allowbreak{} v)) \land (y \in V(M) \cap \operatorname{supp}(\operatorname{connected\,Component\,Mk}(G,\allowbreak{} v)) \land v \sim_{M} y) \implies y = w) = (y \in V(M) \implies y\text{ can reach }v\text{ in }G \implies v \sim_{M} y \implies y = w)\).
\end{proofstep}
Fix \(y\) in \(V\).
Assume \(y \in V(M)\).
Assume \(y\text{ can reach }v\text{ in }G\).
Applying the hypothesis \(\forall y\, :\, V,\; v \sim_{M} y \implies y = w\), we obtain \(v \sim_{M} y \implies y = w\).
Rewriting with \paperlink{premise:5365742e6d656d5f696e7465725f696666}{Prerequisite~P\ref{premise:5365742e6d656d5f696e7465725f696666}}, \(v \in V(M)\), \paperlink{premise:53696d706c6547726170682e436f6e6e6563746564436f6d706f6e656e742e6d656d5f737570705f696666}{Prerequisite~P\ref{premise:53696d706c6547726170682e436f6e6e6563746564436f6d706f6e656e742e6d656d5f737570705f696666}}, \paperlink{premise:616e645f73656c66}{Prerequisite~P\ref{premise:616e645f73656c66}}, \paperlink{premise:53696d706c6547726170682e53756267726170682e656467655f76657274}{Prerequisite~P\ref{premise:53696d706c6547726170682e53756267726170682e656467655f76657274}}, \paperlink{premise:53696d706c6547726170682e53756267726170682e41646a2e73796d6d}{Prerequisite~P\ref{premise:53696d706c6547726170682e53756267726170682e41646a2e73796d6d}} (using \(v \sim_{M} w\)), \paperlink{premise:53696d706c6547726170682e436f6e6e6563746564436f6d706f6e656e742e6571}{Prerequisite~P\ref{premise:53696d706c6547726170682e436f6e6e6563746564436f6d706f6e656e742e6571}}, \paperlink{premise:53696d706c6547726170682e41646a2e726561636861626c65}{Prerequisite~P\ref{premise:53696d706c6547726170682e41646a2e726561636861626c65}}, \paperlink{premise:53696d706c6547726170682e41646a2e73796d6d}{Prerequisite~P\ref{premise:53696d706c6547726170682e41646a2e73796d6d}}, \paperlink{premise:53696d706c6547726170682e53756267726170682e61646a5f737562}{Prerequisite~P\ref{premise:53696d706c6547726170682e53756267726170682e61646a5f737562}} (using \(v \sim_{M} w\)), \(v \sim_{M} w\), \paperlink{premise:666f72616c6c5f636f6e6772}{Prerequisite~P\ref{premise:666f72616c6c5f636f6e6772}}, \paperlink{premise:747275655f616e64}{Prerequisite~P\ref{premise:747275655f616e64}} and the identity established above in the preceding result gives \(v \sim_{\operatorname{Simple\,Graph\mathord{:}Subgraph\mathord{:}induce}(M,\allowbreak{} V(M) \cap \operatorname{supp}(\operatorname{connected\,Component\,Mk}(G,\allowbreak{} v)))} w \land (\forall y\, :\, V,\; v \sim_{\operatorname{Simple\,Graph\mathord{:}Subgraph\mathord{:}induce}(M,\allowbreak{} V(M) \cap \operatorname{supp}(\operatorname{connected\,Component\,Mk}(G,\allowbreak{} v)))} y \implies y = w)\).
This witness proves existence. Hence \(\exists x\, :\, V,\; v \sim_{\operatorname{Simple\,Graph\mathord{:}Subgraph\mathord{:}induce}(M,\allowbreak{} V(M) \cap \operatorname{supp}(\operatorname{connected\,Component\,Mk}(G,\allowbreak{} v)))} x \land (\forall y\, :\, V,\; v \sim_{\operatorname{Simple\,Graph\mathord{:}Subgraph\mathord{:}induce}(M,\allowbreak{} V(M) \cap \operatorname{supp}(\operatorname{connected\,Component\,Mk}(G,\allowbreak{} v)))} y \implies y = x)\).
\end{proofstep}
Using the two conjuncts, we obtain \(\exists! w_{1}\, :\, V,\; v \sim_{\operatorname{Simple\,Graph\mathord{:}Subgraph\mathord{:}induce}(M,\allowbreak{} V(M) \cap \operatorname{supp}(\operatorname{connected\,Component\,Mk}(G,\allowbreak{} v)))} w_{1}\).
\end{proofstep}
Using the existential hypothesis, we obtain \(\exists! w\, :\, V,\; v \sim_{\operatorname{Simple\,Graph\mathord{:}Subgraph\mathord{:}induce}(M,\allowbreak{} V(M) \cap \operatorname{supp}(\operatorname{connected\,Component\,Mk}(G,\allowbreak{} v)))} w\).
Use the hypothesis \(v \in \operatorname{supp}\left(c\right)\).
Substituting equals, we obtain \(\exists! w\, :\, V,\; v \sim_{\operatorname{Simple\,Graph\mathord{:}Subgraph\mathord{:}induce}(M,\allowbreak{} V(M) \cap \operatorname{supp}(c))} w\).
\end{proofstep}
Using the two conjuncts, we obtain \(\exists! w\, :\, V,\; v \sim_{\operatorname{Simple\,Graph\mathord{:}Subgraph\mathord{:}induce}(M,\allowbreak{} V(M) \cap \operatorname{supp}(c))} w\).
\end{proof}

\begin{corollary}[{Induce connected component is matching}]\label{theorem:23}
Let \(V\) be a type. Let \(G\) be a simple graph on \(V\). Let \(M\) be a subgraph of \(G\). Assume \(M\text{ is a perfect matching}\). Let \(c\) be an element of \(\operatorname{Connected\,Component}\left(G\right)\). Then
\begin{center}
\(\displaystyle \operatorname{Simple\,Graph\mathord{:}Subgraph\mathord{:}induce}(M,\allowbreak{} \operatorname{supp}(c))\text{ is a matching}.\)
\end{center}
\end{corollary}

\begin{proof}
Use the hypothesis \(M\text{ is a perfect matching}\).
Taking the second conjunct, we obtain \(\operatorname{Is\,Spanning}\left(M\right)\).
By \paperlink{premise:53696d706c6547726170682e53756267726170682e49735370616e6e696e672e76657274735f65715f756e6976}{Prerequisite~P\ref{premise:53696d706c6547726170682e53756267726170682e49735370616e6e696e672e76657274735f65715f756e6976}}, we obtain \(V(M) = V\).
Applying the same function to equal arguments, we obtain \(V(M) \cap \operatorname{supp}\left(c\right) = V \cap \operatorname{supp}\left(c\right)\).
By \paperlink{premise:5365742e756e69765f696e746572}{Prerequisite~P\ref{premise:5365742e756e69765f696e746572}} with \(a = \operatorname{supp}\left(c\right)\), we obtain \(V \cap \operatorname{supp}\left(c\right) = \operatorname{supp}\left(c\right)\).
Chaining the equalities, we obtain \(V(M) \cap \operatorname{supp}\left(c\right) = \operatorname{supp}\left(c\right)\).
Applying the same function to equal arguments, we obtain \((\operatorname{Simple\,Graph\mathord{:}Subgraph\mathord{:}induce}(M,\allowbreak{} V(M) \cap \operatorname{supp}(c))\text{ is a matching}) = (\operatorname{Simple\,Graph\mathord{:}Subgraph\mathord{:}induce}(M,\allowbreak{} \operatorname{supp}(c))\text{ is a matching})\).
Use the hypothesis \(M\text{ is a perfect matching}\).
Taking the first conjunct, we obtain \(M\text{ is a matching}\).
By \paperlink{theorem:22}{Lemma~\ref{theorem:22}}, we obtain \(\operatorname{Simple\,Graph\mathord{:}Subgraph\mathord{:}induce}(M,\allowbreak{} V(M) \cap \operatorname{supp}(c))\text{ is a matching}\).
Substituting along the equality, we obtain \(\operatorname{Simple\,Graph\mathord{:}Subgraph\mathord{:}induce}(M,\allowbreak{} \operatorname{supp}(c))\text{ is a matching}\).
\end{proof}

\begin{theorem}[{To subgraph if and only if}]\label{theorem:24}
Let \(V\) be a type. Let \(G\) and \(G^{'}\) be simple graphs on \(V\). Let \(M\) be a subgraph of \(G\). Let \(h\) be a proof of \(\operatorname{spanning\,graph}(M) \subseteq G^{'}\). Then
\begin{center}
\(\displaystyle \operatorname{Simple\,Graph\mathord{:}to\,Subgraph}(\operatorname{spanning\,graph}(M))\text{ is a perfect matching} \iff M\text{ is a perfect matching}.\)
\end{center}
\end{theorem}

\begin{proof}
By \paperlink{theorem:19}{Lemma~\ref{theorem:19}}, we obtain \(\operatorname{Simple\,Graph\mathord{:}to\,Subgraph}(\operatorname{spanning\,graph}(M))\text{ is a perfect matching} \iff (\forall v\, :\, V,\; \exists! w\, :\, V,\; v \sim_{\operatorname{Simple\,Graph\mathord{:}to\,Subgraph}(\operatorname{spanning\,graph}(M))} w)\).
By propositional extensionality, we obtain \((\operatorname{Simple\,Graph\mathord{:}to\,Subgraph}(\operatorname{spanning\,graph}(M))\text{ is a perfect matching}) = (\forall v\, :\, V,\; \exists! w\, :\, V,\; v \sim_{\operatorname{Simple\,Graph\mathord{:}to\,Subgraph}(\operatorname{spanning\,graph}(M))} w)\).
Applying the same function to equal arguments, we obtain \((y \mapsto \operatorname{Simple\,Graph\mathord{:}to\,Subgraph}(\operatorname{spanning\,graph}(M))\text{ is a perfect matching} \iff y) = (y \mapsto (\forall v\, :\, V,\; \exists! w\, :\, V,\; v \sim_{\operatorname{Simple\,Graph\mathord{:}to\,Subgraph}(\operatorname{spanning\,graph}(M))} w) \iff y)\).
By \paperlink{theorem:19}{Lemma~\ref{theorem:19}}, we obtain \(M\text{ is a perfect matching} \iff \left(\forall v\, :\, V,\; \exists! w\, :\, V,\; v \sim_{M} w\right)\).
By propositional extensionality, we obtain \((M\text{ is a perfect matching}) = (\forall v\, :\, V,\; \exists! w\, :\, V,\; v \sim_{M} w)\).
By congruence, we obtain \((\operatorname{Simple\,Graph\mathord{:}to\,Subgraph}(\operatorname{spanning\,graph}(M))\text{ is a perfect matching} \iff M\text{ is a perfect matching}) = ((\forall v\, :\, V,\; \exists! w\, :\, V,\; v \sim_{\operatorname{Simple\,Graph\mathord{:}to\,Subgraph}(\operatorname{spanning\,graph}(M))} w) \iff (\forall v\, :\, V,\; \exists! w\, :\, V,\; v \sim_{M} w))\).
By \paperlink{premise:6966665f73656c66}{Prerequisite~P\ref{premise:6966665f73656c66}} with \(p = \forall v\, :\, V,\; \exists! w\, :\, V,\; v \sim_{M} w\), we obtain \(((\forall v\, :\, V,\; \exists! w\, :\, V,\; v \sim_{M} w) \iff (\forall v\, :\, V,\; \exists! w\, :\, V,\; v \sim_{M} w)) = \mathrm{True}\).
Chaining the equalities, we obtain \((\operatorname{Simple\,Graph\mathord{:}to\,Subgraph}(\operatorname{spanning\,graph}(M))\text{ is a perfect matching} \iff M\text{ is a perfect matching}) = \mathrm{True}\).
By \paperlink{premise:6f665f65715f74727565}{Prerequisite~P\ref{premise:6f665f65715f74727565}}, we obtain \(\operatorname{Simple\,Graph\mathord{:}to\,Subgraph}(\operatorname{spanning\,graph}(M))\text{ is a perfect matching} \iff M\text{ is a perfect matching}\).
\end{proof}

\section{Simple graphs}
\begin{theorem}[{Even if and only if exists is matching}]\label{theorem:25}
Let \(V\) be a type. Let \(G\) be a simple graph on \(V\). Let \(u\) be a subset of \(V\). Assume \(u\text{ is a clique in }G\). Assume \(\operatorname{Set\mathord{:}Finite}\left(u\right)\). Then
\begin{center}
\(\displaystyle \operatorname{Even}(\operatorname{ncard}(u)) \iff (\exists M\, :\, \operatorname{Subgraph}(G),\; V(M) = u \land M\text{ is a matching}).\)
\end{center}
\end{theorem}

\begin{proof}
Consider the following case.
\begin{proofstep}
Assume \(\operatorname{Even}\left(\operatorname{ncard}(u)\right)\).
Use the hypothesis \(\operatorname{Even}\left(\operatorname{ncard}(u)\right)\).
By \paperlink{premise:5365742e6578697374735f756e696f6e5f6469736a6f696e745f6e636172645f65715f6f665f6576656e}{Prerequisite~P\ref{premise:5365742e6578697374735f756e696f6e5f6469736a6f696e745f6e636172645f65715f6f665f6576656e}}, we obtain \(\exists t\, :\, \mathcal{P}(V),\; \exists u_{1}\, :\, \mathcal{P}(V),\; t \cup u_{1} = u \land (\operatorname{Disjoint}(t,\allowbreak{} u_{1}) \land \operatorname{ncard}(t) = \operatorname{ncard}(u_{1}))\).
Consider the following case.
\begin{proofstep}
Fix \(t\) in \(\mathcal{P}(V)\).
Assume \(\exists u_{1}\, :\, \mathcal{P}(V),\; t \cup u_{1} = u \land (\operatorname{Disjoint}(t,\allowbreak{} u_{1}) \land \operatorname{ncard}(t) = \operatorname{ncard}(u_{1}))\).
Use the hypothesis \(\exists u_{1}\, :\, \mathcal{P}(V),\; t \cup u_{1} = u \land (\operatorname{Disjoint}(t,\allowbreak{} u_{1}) \land \operatorname{ncard}(t) = \operatorname{ncard}(u_{1}))\).
Consider the following case.
\begin{proofstep}
Fix \(u_{1}\) in \(\mathcal{P}(V)\).
Assume \(t \cup u_{1} = u \land (\operatorname{Disjoint}(t,\allowbreak{} u_{1}) \land \operatorname{ncard}(t) = \operatorname{ncard}(u_{1}))\).
Use the hypothesis \(t \cup u_{1} = u \land (\operatorname{Disjoint}(t,\allowbreak{} u_{1}) \land \operatorname{ncard}(t) = \operatorname{ncard}(u_{1}))\).
Consider the following case.
\begin{proofstep}
Assume \(t \cup u_{1} = u\).
Assume \(\operatorname{Disjoint}\left(t, u_{1}\right) \land \operatorname{ncard}(t) = \operatorname{ncard}(u_{1})\).
For the required implication, argue as follows.
\begin{proofstep}
Assume \(t \cup u_{1}\text{ is a clique in }G\).
Assume \(\operatorname{Set\mathord{:}Finite}\left(t \cup u_{1}\right)\).
Assume \(\operatorname{Even}\left(\operatorname{ncard}(t \cup u_{1})\right)\).
Use the hypothesis \(\operatorname{Disjoint}\left(t, u_{1}\right) \land \operatorname{ncard}(t) = \operatorname{ncard}(u_{1})\).
Consider the following case.
\begin{proofstep}
Assume \(\operatorname{Disjoint}\left(t, u_{1}\right)\).
Assume \(\operatorname{ncard}(t) = \operatorname{ncard}(u_{1})\).
Use the hypothesis \(\operatorname{ncard}(t) = \operatorname{ncard}(u_{1})\).
Applying the same function to equal arguments, we obtain \(\operatorname{cast}(\operatorname{ncard}(t)) = \operatorname{cast}(\operatorname{ncard}(u_{1}))\).
Rewriting with \paperlink{premise:5365742e636173745f6e63617264}{Prerequisite~P\ref{premise:5365742e636173745f6e63617264}}, \paperlink{premise:416e642e7269676874}{Prerequisite~P\ref{premise:416e642e7269676874}}, \paperlink{premise:4966662e6d70}{Prerequisite~P\ref{premise:4966662e6d70}} (using \(\operatorname{Set\mathord{:}Finite}\left(t \cup u_{1}\right)\)) and \paperlink{premise:5365742e66696e6974655f756e696f6e}{Prerequisite~P\ref{premise:5365742e66696e6974655f756e696f6e}} in the preceding result gives \(\operatorname{cast}(\operatorname{ncard}(t)) = \operatorname{Cardinal\mathord{:}mk}(u_{1})\).
Rewriting with \paperlink{premise:5365742e636173745f6e63617264}{Prerequisite~P\ref{premise:5365742e636173745f6e63617264}}, \paperlink{premise:416e642e6c656674}{Prerequisite~P\ref{premise:416e642e6c656674}}, \paperlink{premise:4966662e6d70}{Prerequisite~P\ref{premise:4966662e6d70}} (using \(\operatorname{Set\mathord{:}Finite}\left(t \cup u_{1}\right)\)) and \paperlink{premise:5365742e66696e6974655f756e696f6e}{Prerequisite~P\ref{premise:5365742e66696e6974655f756e696f6e}} in the preceding result gives \(\operatorname{Cardinal\mathord{:}mk}\left(t\right) = \operatorname{Cardinal\mathord{:}mk}\left(u_{1}\right)\).
Rewriting with \paperlink{premise:43617264696e616c2e6571}{Prerequisite~P\ref{premise:43617264696e616c2e6571}} in the preceding result gives \(\operatorname{Bij}\left(t, u_{1}\right)\text{ is nonempty}\).
Consider the possible forms of \(\operatorname{Bij}\left(t, u_{1}\right)\text{ is nonempty}\).
\emph{Intro case.}
\begin{proofstep}
Fix \(f\) in \(\operatorname{Bij}\left(t, u_{1}\right)\).
For the required implication, argue as follows.
\begin{proofstep}
Fix \(v\) in \(t\).
Rewriting with \paperlink{premise:5365742e6d656d5f756e696f6e}{Prerequisite~P\ref{premise:5365742e6d656d5f756e696f6e}}, \paperlink{premise:537562747970652e636f655f70726f70}{Prerequisite~P\ref{premise:537562747970652e636f655f70726f70}} and \paperlink{premise:747275655f6f72}{Prerequisite~P\ref{premise:747275655f6f72}} gives \(\operatorname{val}(v) \in t \cup u_{1}\).
Rewriting with \paperlink{premise:5365742e6d656d5f756e696f6e}{Prerequisite~P\ref{premise:5365742e6d656d5f756e696f6e}}, \paperlink{premise:537562747970652e636f655f70726f70}{Prerequisite~P\ref{premise:537562747970652e636f655f70726f70}} and \paperlink{premise:6f725f74727565}{Prerequisite~P\ref{premise:6f725f74727565}} gives \(\operatorname{val}(f\left(v\right)) \in t \cup u_{1}\).
Use the hypothesis \(\operatorname{Disjoint}\left(t, u_{1}\right)\).
Rewriting with \paperlink{premise:537562747970652e636f655f70726f70}{Prerequisite~P\ref{premise:537562747970652e636f655f70726f70}} gives \(\operatorname{val}(v) \in t\).
Rewriting with \paperlink{premise:537562747970652e636f655f70726f70}{Prerequisite~P\ref{premise:537562747970652e636f655f70726f70}} gives \(\operatorname{val}(f\left(v\right)) \in u_{1}\).
By \paperlink{premise:4469736a6f696e742e6e655f6f665f6d656d}{Prerequisite~P\ref{premise:4469736a6f696e742e6e655f6f665f6d656d}}, we obtain \(\operatorname{val}(v) \ne \operatorname{val}(f\left(v\right))\).
Applying the hypothesis \(t \cup u_{1}\text{ is a clique in }G\), we obtain \(\operatorname{val}(v) \sim_{G} \operatorname{val}(f\left(v\right))\).
\end{proofstep}
By \paperlink{theorem:13}{Lemma~\ref{theorem:13}}, we obtain \(\exists M\, :\, \operatorname{Subgraph}\left(G\right),\; V(M) = t \cup u_{1} \land M\text{ is a matching}\).
\end{proofstep}
This proves \(\exists M\, :\, \operatorname{Subgraph}\left(G\right),\; V(M) = t \cup u_{1} \land M\text{ is a matching}\).
\end{proofstep}
Using the two conjuncts, we obtain \(\exists M\, :\, \operatorname{Subgraph}\left(G\right),\; V(M) = t \cup u_{1} \land M\text{ is a matching}\).
\end{proofstep}
Use the hypothesis \(t \cup u_{1} = u\).
Use the hypothesis \(u\text{ is a clique in }G\).
Use the hypothesis \(\operatorname{Set\mathord{:}Finite}\left(u\right)\).
Use the hypothesis \(\operatorname{Even}\left(\operatorname{ncard}(u)\right)\).
Substituting equals, we obtain \(\exists M\, :\, \operatorname{Subgraph}\left(G\right),\; V(M) = u \land M\text{ is a matching}\).
\end{proofstep}
Using the two conjuncts, we obtain \(\exists M\, :\, \operatorname{Subgraph}\left(G\right),\; V(M) = u \land M\text{ is a matching}\).
\end{proofstep}
Using the existential hypothesis, we obtain \(\exists M\, :\, \operatorname{Subgraph}\left(G\right),\; V(M) = u \land M\text{ is a matching}\).
\end{proofstep}
Using the existential hypothesis, we obtain \(\exists M\, :\, \operatorname{Subgraph}\left(G\right),\; V(M) = u \land M\text{ is a matching}\).
\end{proofstep}
Consider the following case.
\begin{proofstep}
Assume \(\exists M\, :\, \operatorname{Subgraph}\left(G\right),\; V(M) = u \land M\text{ is a matching}\).
Use the hypothesis \(\exists M\, :\, \operatorname{Subgraph}\left(G\right),\; V(M) = u \land M\text{ is a matching}\).
Consider the following case.
\begin{proofstep}
Fix \(M\) in \(\operatorname{Subgraph}\left(G\right)\).
Assume \(V(M) = u \land M\text{ is a matching}\).
Use the hypothesis \(V(M) = u \land M\text{ is a matching}\).
Consider the following case.
\begin{proofstep}
Assume \(V(M) = u\).
Assume \(M\text{ is a matching}\).
For the required implication, argue as follows.
\begin{proofstep}
Assume \(V(M)\text{ is a clique in }G\).
Assume \(\operatorname{Set\mathord{:}Finite}\left(V(M)\right)\).
By \paperlink{theorem:18}{Lemma~\ref{theorem:18}}, we obtain \(\operatorname{Even}\left(\left|\operatorname{as\,finite\,set}(V(M))\right|\right)\).
Rewriting with \paperlink{premise:5365742e6e636172645f65715f746f46696e7365745f63617264}{Prerequisite~P\ref{premise:5365742e6e636172645f65715f746f46696e7365745f63617264}} (using \(\operatorname{Set\mathord{:}Finite}\left(V(M)\right)\)) in the preceding result gives \(\operatorname{Even}\left(\operatorname{ncard}(V(M))\right)\).
\end{proofstep}
Use the hypothesis \(V(M) = u\).
Use the hypothesis \(u\text{ is a clique in }G\).
Use the hypothesis \(\operatorname{Set\mathord{:}Finite}\left(u\right)\).
Substituting equals, we obtain \(\operatorname{Even}\left(\operatorname{ncard}(u)\right)\).
\end{proofstep}
Using the two conjuncts, we obtain \(\operatorname{Even}\left(\operatorname{ncard}(u)\right)\).
\end{proofstep}
Using the existential hypothesis, we obtain \(\operatorname{Even}\left(\operatorname{ncard}(u)\right)\).
\end{proofstep}
Together the two implications give the equivalence. Hence \(\operatorname{Even}(\operatorname{ncard}(u)) \iff (\exists M\, :\, \operatorname{Subgraph}(G),\; V(M) = u \land M\text{ is a matching})\).
\end{proof}

\section{Finite}
\begin{theorem}[{Even cardinality}]\label{theorem:26}
Let \(V\) be a type. Let \(G\) be a simple graph on \(V\). Let \(M\) be a subgraph of \(G\). Fix a finite enumeration of \(V\). Assume equality on \(V\) is decidable. Assume adjacency in \(G\) is decidable. Let \(c\) be an element of \(\operatorname{Connected\,Component}\left(G\right)\). Assume \(M\text{ is a perfect matching}\). Then \(\operatorname{Even}\left(\left|\operatorname{supp}\left(c\right)\right|\right)\).
\end{theorem}

\begin{proof}
The substitution uses the following argument.
\begin{proofstep}
Fix \(x\) in \(V\).
By \paperlink{premise:46696e7365742e6d656d5f66696c746572}{Prerequisite~P\ref{premise:46696e7365742e6d656d5f66696c746572}} with \(s = V\), we obtain \(x \in \{x_{1}\in V \mid (y \mapsto y \in \operatorname{supp}(c))(x_{1})\} \iff x \in V \land x \in \operatorname{supp}(c)\).
By propositional extensionality, we obtain \((x \in \{x_{1}\in V \mid (y \mapsto y \in \operatorname{supp}(c))(x_{1})\}) = (x \in V \land x \in \operatorname{supp}(c))\).
By \paperlink{premise:46696e7365742e6d656d5f756e6976}{Prerequisite~P\ref{premise:46696e7365742e6d656d5f756e6976}}, we obtain \(x \in V\).
By \paperlink{premise:65715f74727565}{Prerequisite~P\ref{premise:65715f74727565}}, we obtain \(\left(x \in V\right) = \mathrm{True}\).
Applying the same function to equal arguments, we obtain \((x \in V \land x \in \operatorname{supp}(c)) = (\mathrm{True} \land x \in \operatorname{supp}(c))\).
By \paperlink{premise:747275655f616e64}{Prerequisite~P\ref{premise:747275655f616e64}} with \(p = x \in \operatorname{supp}\left(c\right)\), we obtain \((\mathrm{True} \land x \in \operatorname{supp}(c)) = (x \in \operatorname{supp}(c))\).
Chaining the equalities, we obtain \((x \in V \land x \in \operatorname{supp}(c)) = (x \in \operatorname{supp}(c))\).
Chaining the equalities, we obtain \((x \in \{x_{1}\in V \mid (y \mapsto y \in \operatorname{supp}(c))(x_{1})\}) = (x \in \operatorname{supp}(c))\).
Applying the same function to equal arguments, we obtain \((x \in \{x_{1}\in V \mid (y \mapsto y \in \operatorname{supp}(c))(x_{1})\} \iff x \in \operatorname{supp}(c)) = (x \in \operatorname{supp}(c) \iff x \in \operatorname{supp}(c))\).
By \paperlink{premise:6966665f73656c66}{Prerequisite~P\ref{premise:6966665f73656c66}} with \(p = x \in \operatorname{supp}\left(c\right)\), we obtain \((x \in \operatorname{supp}(c) \iff x \in \operatorname{supp}(c)) = \mathrm{True}\).
Chaining the equalities, we obtain \((x \in \{x_{1}\in V \mid (y \mapsto y \in \operatorname{supp}(c))(x_{1})\} \iff x \in \operatorname{supp}(c)) = \mathrm{True}\).
\end{proofstep}
The substitution uses the following argument.
\begin{proofstep}
Fix \(x\) in \(V\).
By \paperlink{premise:46696e7365742e6d656d5f66696c746572}{Prerequisite~P\ref{premise:46696e7365742e6d656d5f66696c746572}} with \(s = V\), we obtain \(x \in \{x_{1}\in V \mid (y \mapsto y \in \operatorname{supp}(c))(x_{1})\} \iff x \in V \land x \in \operatorname{supp}(c)\).
By propositional extensionality, we obtain \((x \in \{x_{1}\in V \mid (y \mapsto y \in \operatorname{supp}(c))(x_{1})\}) = (x \in V \land x \in \operatorname{supp}(c))\).
By \paperlink{premise:46696e7365742e6d656d5f756e6976}{Prerequisite~P\ref{premise:46696e7365742e6d656d5f756e6976}}, we obtain \(x \in V\).
By \paperlink{premise:65715f74727565}{Prerequisite~P\ref{premise:65715f74727565}}, we obtain \(\left(x \in V\right) = \mathrm{True}\).
Applying the same function to equal arguments, we obtain \((x \in V \land x \in \operatorname{supp}(c)) = (\mathrm{True} \land x \in \operatorname{supp}(c))\).
By \paperlink{premise:747275655f616e64}{Prerequisite~P\ref{premise:747275655f616e64}} with \(p = x \in \operatorname{supp}\left(c\right)\), we obtain \((\mathrm{True} \land x \in \operatorname{supp}(c)) = (x \in \operatorname{supp}(c))\).
Chaining the equalities, we obtain \((x \in V \land x \in \operatorname{supp}(c)) = (x \in \operatorname{supp}(c))\).
Chaining the equalities, we obtain \((x \in \{x_{1}\in V \mid (y \mapsto y \in \operatorname{supp}(c))(x_{1})\}) = (x \in \operatorname{supp}(c))\).
Applying the same function to equal arguments, we obtain \((x \in \{x_{1}\in V \mid (y \mapsto y \in \operatorname{supp}(c))(x_{1})\} \iff x \in \operatorname{supp}(c)) = (x \in \operatorname{supp}(c) \iff x \in \operatorname{supp}(c))\).
By \paperlink{premise:6966665f73656c66}{Prerequisite~P\ref{premise:6966665f73656c66}} with \(p = x \in \operatorname{supp}\left(c\right)\), we obtain \((x \in \operatorname{supp}(c) \iff x \in \operatorname{supp}(c)) = \mathrm{True}\).
Chaining the equalities, we obtain \((x \in \{x_{1}\in V \mid (y \mapsto y \in \operatorname{supp}(c))(x_{1})\} \iff x \in \operatorname{supp}(c)) = \mathrm{True}\).
\end{proofstep}
By \paperlink{theorem:23}{Corollary~\ref{theorem:23}}, we obtain \(\operatorname{Simple\,Graph\mathord{:}Subgraph\mathord{:}induce}(M,\allowbreak{} \operatorname{supp}(c))\text{ is a matching}\).
By \paperlink{theorem:18}{Lemma~\ref{theorem:18}}, we obtain \(\operatorname{Even}(|\operatorname{as\,finite\,set}(V(\operatorname{Simple\,Graph\mathord{:}Subgraph\mathord{:}induce}(M,\allowbreak{} \operatorname{supp}(c))))|)\).
Rewriting with \paperlink{premise:5365742e746f46696e7365745f63617264}{Prerequisite~P\ref{premise:5365742e746f46696e7365745f63617264}}, \paperlink{premise:46696e747970652e636172645f6f6646696e736574}{Prerequisite~P\ref{premise:46696e747970652e636172645f6f6646696e736574}}, \paperlink{premise:666f72616c6c5f636f6e6772}{Prerequisite~P\ref{premise:666f72616c6c5f636f6e6772}}, \paperlink{premise:696d706c6965735f74727565}{Prerequisite~P\ref{premise:696d706c6965735f74727565}} and the identity established above in the preceding result gives \(\operatorname{Even}(|\{x\in V \mid (y \mapsto y \in \operatorname{supp}(c))(x)\}|)\).
Rewriting with \paperlink{premise:46696e747970652e636172645f6f6646696e736574}{Prerequisite~P\ref{premise:46696e747970652e636172645f6f6646696e736574}}, \paperlink{premise:666f72616c6c5f636f6e6772}{Prerequisite~P\ref{premise:666f72616c6c5f636f6e6772}}, \paperlink{premise:696d706c6965735f74727565}{Prerequisite~P\ref{premise:696d706c6965735f74727565}} and the identity established above in the preceding result gives \(\operatorname{Even}\left(\left|\operatorname{supp}\left(c\right)\right|\right)\).
\end{proof}

\begin{theorem}[{Odd matches node outside}]\label{theorem:27}
Let \(V\) be a type. Let \(G\) be a simple graph on \(V\). Let \(M\) be a subgraph of \(G\). Assume \(V\) is finite. Let \(u\) be a subset of \(V\). Assume \(M\text{ is a perfect matching}\). Let \(c\) be an element of \(\operatorname{odd\,Components}(\operatorname{graph}(\operatorname{delete\,Verts}(\operatorname{full\,subgraph}(G),\allowbreak{} u)))\). Then
\begin{center}
\(\displaystyle \exists w\, :\, V,\; w \in u \land (\exists v\, :\, V(\operatorname{delete\,Verts}(\operatorname{full\,subgraph}(G),\allowbreak{} u)),\; v \sim_{M} w \land v \in \operatorname{supp}(c)).\)
\end{center}
\end{theorem}

\begin{proof}
Consider the following case.
\begin{proofstep}
Assume \(\neg (\exists w\, :\, V,\; w \in u \land (\exists v\, :\, V(\operatorname{delete\,Verts}(\operatorname{full\,subgraph}(G),\allowbreak{} u)),\; v \sim_{M} w \land v \in \operatorname{supp}(\operatorname{val}(c))))\).
Fix \(v\) in \(V\).
Assume \(v \in V(\operatorname{Simple\,Graph\mathord{:}Subgraph\mathord{:}induce}(M,\allowbreak{} \operatorname{val}[\operatorname{supp}(\operatorname{val}(c))]))\).
Use the hypothesis \(M\text{ is a perfect matching}\).
Use the hypothesis \(M\text{ is a perfect matching}\).
Taking the second conjunct, we obtain \(v \in V(M)\).
Taking the first conjunct, we obtain \(\exists! w\, :\, V,\; v \sim_{M} w\).
Consider the following case.
\begin{proofstep}
Fix \(w\) in \(V\).
Assume \(v \sim_{M} w \land \left(\forall y\, :\, V,\; v \sim_{M} y \implies y = w\right)\).
Use the hypothesis \(v \in V(\operatorname{Simple\,Graph\mathord{:}Subgraph\mathord{:}induce}(M,\allowbreak{} \operatorname{val}[\operatorname{supp}(\operatorname{val}(c))]))\).
Consider the following case.
\begin{proofstep}
Fix \(w_{1}\) in \(V(\operatorname{delete\,Verts}\left(\operatorname{full\,subgraph}(G), u\right))\).
Assume \(w_{1} \in \operatorname{supp}\left(\operatorname{val}(c)\right) \land \operatorname{val}(w_{1}) = v\).
Consider the following case.
\begin{proofstep}
Fix \(v^{'}\) in \(V\).
Assume \(v^{'} \in V(\operatorname{delete\,Verts}\left(\operatorname{full\,subgraph}(G), u\right))\).
Assume \((v^{'}\, :\, V(\operatorname{delete\,Verts}(\operatorname{full\,subgraph}(G),\allowbreak{} u))) \in \operatorname{supp}(\operatorname{val}(c)) \land v^{'} = v\).
Use the hypothesis \((v^{'}\, :\, V(\operatorname{delete\,Verts}(\operatorname{full\,subgraph}(G),\allowbreak{} u))) \in \operatorname{supp}(\operatorname{val}(c)) \land v^{'} = v\).
Consider the following case.
\begin{proofstep}
Assume \((v^{'}\, :\, V(\operatorname{delete\,Verts}(\operatorname{full\,subgraph}(G),\allowbreak{} u))) \in \operatorname{supp}(\operatorname{val}(c))\).
Assume \(v^{'} = v\).
For the required implication, argue as follows.
\begin{proofstep}
Assume \(v^{'} \sim_{M} w \land \left(\forall y\, :\, V,\; v^{'} \sim_{M} y \implies y = w\right)\).
Choose the witness \(w\).
Assume \(w \in u\).
By \paperlink{premise:4d6174686c69622e5461637469632e507573684e65672e6e6f745f6578697374735f6571}{Prerequisite~P\ref{premise:4d6174686c69622e5461637469632e507573684e65672e6e6f745f6578697374735f6571}} with \(s = w_{2} \mapsto w_{2} \in u \land (\exists v_{1}\, :\, V(\operatorname{delete\,Verts}(\operatorname{full\,subgraph}(G),\allowbreak{} u)),\; v_{1} \sim_{M} w_{2} \land v_{1} \in \operatorname{supp}(\operatorname{val}(c)))\), we obtain \((\neg (\exists x\, :\, V,\; x \in u \land (\exists v_{1}\, :\, V(\operatorname{delete\,Verts}(\operatorname{full\,subgraph}(G),\allowbreak{} u)),\; v_{1} \sim_{M} x \land v_{1} \in \operatorname{supp}(\operatorname{val}(c))))) = (\forall x\, :\, V,\; \neg (x \in u \land (\exists v_{1}\, :\, V(\operatorname{delete\,Verts}(\operatorname{full\,subgraph}(G),\allowbreak{} u)),\; v_{1} \sim_{M} x \land v_{1} \in \operatorname{supp}(\operatorname{val}(c)))))\).
For the required implication, argue as follows.
\begin{proofstep}
Fix \(w_{2}\) in \(V\).
By \paperlink{premise:4d6174686c69622e5461637469632e507573684e65672e6e6f745f616e645f6571}{Prerequisite~P\ref{premise:4d6174686c69622e5461637469632e507573684e65672e6e6f745f616e645f6571}} with \(p = w_{2} \in u\) and \(q = \exists v_{1}\, :\, V(\operatorname{delete\,Verts}(\operatorname{full\,subgraph}(G),\allowbreak{} u)),\; v_{1} \sim_{M} w_{2} \land v_{1} \in \operatorname{supp}(\operatorname{val}(c))\), we obtain \((\neg (w_{2} \in u \land (\exists v_{1}\, :\, V(\operatorname{delete\,Verts}(\operatorname{full\,subgraph}(G),\allowbreak{} u)),\; v_{1} \sim_{M} w_{2} \land v_{1} \in \operatorname{supp}(\operatorname{val}(c))))) = (w_{2} \in u \implies \neg (\exists v_{1}\, :\, V(\operatorname{delete\,Verts}(\operatorname{full\,subgraph}(G),\allowbreak{} u)),\; v_{1} \sim_{M} w_{2} \land v_{1} \in \operatorname{supp}(\operatorname{val}(c))))\).
By definition, we obtain \(\left(w_{2} \in u\right) = \left(w_{2} \in u\right)\).
By \paperlink{premise:4d6174686c69622e5461637469632e507573684e65672e6e6f745f6578697374735f6571}{Prerequisite~P\ref{premise:4d6174686c69622e5461637469632e507573684e65672e6e6f745f6578697374735f6571}} with \(\alpha = V(\operatorname{delete\,Verts}\left(\operatorname{full\,subgraph}(G), u\right))\) and \(s = v_{1} \mapsto v_{1} \sim_{M} w_{2} \land v_{1} \in \operatorname{supp}\left(\operatorname{val}(c)\right)\), we obtain \((\neg (\exists x\, :\, V(\operatorname{delete\,Verts}(\operatorname{full\,subgraph}(G),\allowbreak{} u)),\; x \sim_{M} w_{2} \land x \in \operatorname{supp}(\operatorname{val}(c)))) = (\forall x\, :\, V(\operatorname{delete\,Verts}(\operatorname{full\,subgraph}(G),\allowbreak{} u)),\; \neg (x \sim_{M} w_{2} \land x \in \operatorname{supp}(\operatorname{val}(c))))\).
For the required implication, argue as follows.
\begin{proofstep}
Fix \(v_{1}\) in \(V(\operatorname{delete\,Verts}\left(\operatorname{full\,subgraph}(G), u\right))\).
By \paperlink{premise:4d6174686c69622e5461637469632e507573684e65672e6e6f745f616e645f6571}{Prerequisite~P\ref{premise:4d6174686c69622e5461637469632e507573684e65672e6e6f745f616e645f6571}} with \(p = \operatorname{val}(v_{1}) \sim_{M} w_{2}\) and \(q = v_{1} \in \operatorname{supp}\left(\operatorname{val}(c)\right)\), we obtain \((\neg (\operatorname{val}(v_{1}) \sim_{M} w_{2} \land v_{1} \in \operatorname{supp}(\operatorname{val}(c)))) = (\operatorname{val}(v_{1}) \sim_{M} w_{2} \implies v_{1} \notin \operatorname{supp}(\operatorname{val}(c)))\).
\end{proofstep}
By \paperlink{premise:666f72616c6c5f636f6e6772}{Prerequisite~P\ref{premise:666f72616c6c5f636f6e6772}} with \(\alpha = V(\operatorname{delete\,Verts}\left(\operatorname{full\,subgraph}(G), u\right))\), we obtain \((\forall a\, :\, V(\operatorname{delete\,Verts}(\operatorname{full\,subgraph}(G),\allowbreak{} u)),\; \neg (a \sim_{M} w_{2} \land a \in \operatorname{supp}(\operatorname{val}(c)))) = (\forall a\, :\, V(\operatorname{delete\,Verts}(\operatorname{full\,subgraph}(G),\allowbreak{} u)),\; a \sim_{M} w_{2} \implies a \notin \operatorname{supp}(\operatorname{val}(c)))\).
Chaining the equalities, we obtain \((\neg (\exists x\, :\, V(\operatorname{delete\,Verts}(\operatorname{full\,subgraph}(G),\allowbreak{} u)),\; x \sim_{M} w_{2} \land x \in \operatorname{supp}(\operatorname{val}(c)))) = (\forall a\, :\, V(\operatorname{delete\,Verts}(\operatorname{full\,subgraph}(G),\allowbreak{} u)),\; a \sim_{M} w_{2} \implies a \notin \operatorname{supp}(\operatorname{val}(c)))\).
By \paperlink{premise:696d706c6965735f636f6e6772}{Prerequisite~P\ref{premise:696d706c6965735f636f6e6772}} with \(p_{1} = w_{2} \in u\) and \(p_{2} = w_{2} \in u\) and \(q_{1} = \neg (\exists x\, :\, V(\operatorname{delete\,Verts}(\operatorname{full\,subgraph}(G),\allowbreak{} u)),\; x \sim_{M} w_{2} \land x \in \operatorname{supp}(\operatorname{val}(c)))\) and \(q_{2} = \forall a\, :\, V(\operatorname{delete\,Verts}(\operatorname{full\,subgraph}(G),\allowbreak{} u)),\; a \sim_{M} w_{2} \implies a \notin \operatorname{supp}(\operatorname{val}(c))\), we obtain \((w_{2} \in u \implies \neg (\exists x\, :\, V(\operatorname{delete\,Verts}(\operatorname{full\,subgraph}(G),\allowbreak{} u)),\; x \sim_{M} w_{2} \land x \in \operatorname{supp}(\operatorname{val}(c)))) = (w_{2} \in u \implies (\forall a_{1}\, :\, V(\operatorname{delete\,Verts}(\operatorname{full\,subgraph}(G),\allowbreak{} u)),\; a_{1} \sim_{M} w_{2} \implies a_{1} \notin \operatorname{supp}(\operatorname{val}(c))))\).
Chaining the equalities, we obtain \((\neg (w_{2} \in u \land (\exists v_{1}\, :\, V(\operatorname{delete\,Verts}(\operatorname{full\,subgraph}(G),\allowbreak{} u)),\; v_{1} \sim_{M} w_{2} \land v_{1} \in \operatorname{supp}(\operatorname{val}(c))))) = (w_{2} \in u \implies (\forall a_{1}\, :\, V(\operatorname{delete\,Verts}(\operatorname{full\,subgraph}(G),\allowbreak{} u)),\; a_{1} \sim_{M} w_{2} \implies a_{1} \notin \operatorname{supp}(\operatorname{val}(c))))\).
\end{proofstep}
By \paperlink{premise:666f72616c6c5f636f6e6772}{Prerequisite~P\ref{premise:666f72616c6c5f636f6e6772}}, we obtain \((\forall a\, :\, V,\; \neg (a \in u \land (\exists v_{1}\, :\, V(\operatorname{delete\,Verts}(\operatorname{full\,subgraph}(G),\allowbreak{} u)),\; v_{1} \sim_{M} a \land v_{1} \in \operatorname{supp}(\operatorname{val}(c))))) = (\forall a\, :\, V,\; a \in u \implies (\forall a_{2}\, :\, V(\operatorname{delete\,Verts}(\operatorname{full\,subgraph}(G),\allowbreak{} u)),\; a_{2} \sim_{M} a \implies a_{2} \notin \operatorname{supp}(\operatorname{val}(c))))\).
Chaining the equalities, we obtain \((\neg (\exists x\, :\, V,\; x \in u \land (\exists v_{1}\, :\, V(\operatorname{delete\,Verts}(\operatorname{full\,subgraph}(G),\allowbreak{} u)),\; v_{1} \sim_{M} x \land v_{1} \in \operatorname{supp}(\operatorname{val}(c))))) = (\forall a\, :\, V,\; a \in u \implies (\forall a_{2}\, :\, V(\operatorname{delete\,Verts}(\operatorname{full\,subgraph}(G),\allowbreak{} u)),\; a_{2} \sim_{M} a \implies a_{2} \notin \operatorname{supp}(\operatorname{val}(c))))\).
Use the hypothesis \(\neg (\exists w\, :\, V,\; w \in u \land (\exists v\, :\, V(\operatorname{delete\,Verts}(\operatorname{full\,subgraph}(G),\allowbreak{} u)),\; v \sim_{M} w \land v \in \operatorname{supp}(\operatorname{val}(c))))\).
Use the hypothesis \(w \in u\).
Use the hypothesis \(v^{'} \sim_{M} w \land \left(\forall y\, :\, V,\; v^{'} \sim_{M} y \implies y = w\right)\).
Taking the first conjunct, we obtain \(v^{'} \sim_{M} w\).
Use the hypothesis \((v^{'}\, :\, V(\operatorname{delete\,Verts}(\operatorname{full\,subgraph}(G),\allowbreak{} u))) \in \operatorname{supp}(\operatorname{val}(c))\).
Substituting along the equality, we obtain \(\mathrm{False}\).
Choose the witness \(\left(v^{'}\, :\, V(\operatorname{delete\,Verts}\left(\operatorname{full\,subgraph}(G), u\right))\right)\).
\par\smallskip\noindent\emph{Property.}\enspace
Use the hypothesis \(v^{'} \in V(\operatorname{delete\,Verts}\left(\operatorname{full\,subgraph}(G), u\right))\), which simplifies to \((x \mapsto x \in V(\operatorname{delete\,Verts}(\operatorname{full\,subgraph}(G),\allowbreak{} u)))(v^{'})\).
Use the hypothesis \((v^{'}\, :\, V(\operatorname{delete\,Verts}(\operatorname{full\,subgraph}(G),\allowbreak{} u))) \in \operatorname{supp}(\operatorname{val}(c))\).
By definition, we obtain \(v^{'} = v^{'}\).
Combining the two conclusions, we obtain \((v^{'}\, :\, V(\operatorname{delete\,Verts}(\operatorname{full\,subgraph}(G),\allowbreak{} u))) \in \operatorname{supp}(\operatorname{val}(c)) \land v^{'} = v^{'}\).
This witness proves existence. Hence \(\exists a\, :\, V(\operatorname{delete\,Verts}(\operatorname{full\,subgraph}(G),\allowbreak{} u)),\; a \in \operatorname{supp}(\operatorname{val}(c)) \land a = v^{'}\).
Choose the witness \(\left(v^{'}\, :\, V(\operatorname{delete\,Verts}\left(\operatorname{full\,subgraph}(G), u\right))\right)\).
\par\smallskip\noindent\emph{Property.}\enspace
Use the hypothesis \(v^{'} \in V(\operatorname{delete\,Verts}\left(\operatorname{full\,subgraph}(G), u\right))\), which simplifies to \((x \mapsto x \in V(\operatorname{delete\,Verts}(\operatorname{full\,subgraph}(G),\allowbreak{} u)))(v^{'})\).
Use the hypothesis \((v^{'}\, :\, V(\operatorname{delete\,Verts}(\operatorname{full\,subgraph}(G),\allowbreak{} u))) \in \operatorname{supp}(\operatorname{val}(c))\).
By definition, we obtain \(v^{'} = v^{'}\).
Combining the two conclusions, we obtain \((v^{'}\, :\, V(\operatorname{delete\,Verts}(\operatorname{full\,subgraph}(G),\allowbreak{} u))) \in \operatorname{Set\,Like\mathord{:}coe}(\operatorname{val}(c)) \land v^{'} = v^{'}\).
This witness proves existence. Hence \(\exists a\, :\, V(\operatorname{delete\,Verts}(\operatorname{full\,subgraph}(G),\allowbreak{} u)),\; a \in \operatorname{Set\,Like\mathord{:}coe}(\operatorname{val}(c)) \land a = v^{'}\).
By the definition of truth, we obtain \(\mathrm{True}\).
Use the previously proved fact \(w \notin u\).
Combining the two conclusions, we obtain \(w \in V(\operatorname{full\,subgraph}(G)) \land w \notin u\).
Use the hypothesis \(v^{'} \in V(\operatorname{delete\,Verts}\left(\operatorname{full\,subgraph}(G), u\right))\).
Rewriting with \paperlink{premise:5365742e6d656d5f64696666}{Prerequisite~P\ref{premise:5365742e6d656d5f64696666}}, \paperlink{premise:5365742e6d656d5f756e6976}{Prerequisite~P\ref{premise:5365742e6d656d5f756e6976}} and \paperlink{premise:747275655f616e64}{Prerequisite~P\ref{premise:747275655f616e64}} in the preceding result gives \(v^{'} \notin u\).
Rewriting with \paperlink{premise:5365742e6d656d5f64696666}{Prerequisite~P\ref{premise:5365742e6d656d5f64696666}}, \paperlink{premise:5365742e6d656d5f756e6976}{Prerequisite~P\ref{premise:5365742e6d656d5f756e6976}}, \paperlink{premise:747275655f616e64}{Prerequisite~P\ref{premise:747275655f616e64}}, \paperlink{premise:65715f66616c7365}{Prerequisite~P\ref{premise:65715f66616c7365}} (using \(w \notin u\)), \paperlink{premise:6e6f745f66616c73655f65715f74727565}{Prerequisite~P\ref{premise:6e6f745f66616c73655f65715f74727565}}, \paperlink{premise:616e645f73656c66}{Prerequisite~P\ref{premise:616e645f73656c66}}, \paperlink{premise:53696d706c6547726170682e53756267726170682e746f705f61646a}{Prerequisite~P\ref{premise:53696d706c6547726170682e53756267726170682e746f705f61646a}}, \paperlink{premise:53696d706c6547726170682e53756267726170682e61646a5f737562}{Prerequisite~P\ref{premise:53696d706c6547726170682e53756267726170682e61646a5f737562}}, \paperlink{premise:416e642e6c656674}{Prerequisite~P\ref{premise:416e642e6c656674}} (using \(v^{'} \sim_{M} w \land \left(\forall y\, :\, V,\; v^{'} \sim_{M} y \implies y = w\right)\)) and \paperlink{premise:616e645f74727565}{Prerequisite~P\ref{premise:616e645f74727565}} in the preceding result gives \(v^{'} \sim_{\operatorname{delete\,Verts}\left(\operatorname{full\,subgraph}(G), u\right)} w\).
By \paperlink{premise:53696d706c6547726170682e436f6e6e6563746564436f6d706f6e656e742e6d656d5f636f655f737570705f6f665f61646a}{Prerequisite~P\ref{premise:53696d706c6547726170682e436f6e6e6563746564436f6d706f6e656e742e6d656d5f636f655f737570705f6f665f61646a}}, we obtain \(w \in \operatorname{val}\left[\operatorname{Set\,Like\mathord{:}coe}\left(\operatorname{val}(c)\right)\right]\).
Use the hypothesis \(v^{'} \sim_{M} w \land \left(\forall y\, :\, V,\; v^{'} \sim_{M} y \implies y = w\right)\).
Taking the first conjunct, we obtain \(v^{'} \sim_{M} w\).
Combining the two conclusions, we obtain \(w \in \operatorname{val}[\operatorname{supp}(\operatorname{val}(c))] \land v^{'} \sim_{M} w\).
Combining the two conclusions, we obtain \(v^{'} \in \operatorname{val}[\operatorname{supp}(\operatorname{val}(c))] \land (w \in \operatorname{val}[\operatorname{supp}(\operatorname{val}(c))] \land v^{'} \sim_{M} w)\).
For the required implication, argue as follows.
\begin{proofstep}
Fix \(x\) in \(V\).
Assume \(v^{'} \sim_{\operatorname{Simple\,Graph\mathord{:}Subgraph\mathord{:}induce}(M,\allowbreak{} \operatorname{val}[\operatorname{supp}(\operatorname{val}(c))])} x\).
Use the hypothesis \(v^{'} \sim_{M} w \land \left(\forall y\, :\, V,\; v^{'} \sim_{M} y \implies y = w\right)\).
Use the hypothesis \(v^{'} \sim_{\operatorname{Simple\,Graph\mathord{:}Subgraph\mathord{:}induce}(M,\allowbreak{} \operatorname{val}[\operatorname{supp}(\operatorname{val}(c))])} x\).
Taking the second conjunct, we obtain \(x \in \operatorname{val}[\operatorname{supp}(\operatorname{val}(c))] \land v^{'} \sim_{M} x\).
Taking the second conjunct, we obtain \(v^{'} \sim_{M} x\).
Taking the second conjunct, we obtain \(x = w\).
\end{proofstep}
Combining the two conclusions, we obtain \(v^{'} \sim_{\operatorname{Simple\,Graph\mathord{:}Subgraph\mathord{:}induce}(M,\allowbreak{} \operatorname{val}[\operatorname{supp}(\operatorname{val}(c))])} w \land (\forall y\, :\, V,\; v^{'} \sim_{\operatorname{Simple\,Graph\mathord{:}Subgraph\mathord{:}induce}(M,\allowbreak{} \operatorname{val}[\operatorname{supp}(\operatorname{val}(c))])} y \implies y = w)\).
This witness proves existence. Hence \(\exists x\, :\, V,\; v^{'} \sim_{\operatorname{Simple\,Graph\mathord{:}Subgraph\mathord{:}induce}(M,\allowbreak{} \operatorname{val}[\operatorname{supp}(\operatorname{val}(c))])} x \land (\forall y\, :\, V,\; v^{'} \sim_{\operatorname{Simple\,Graph\mathord{:}Subgraph\mathord{:}induce}(M,\allowbreak{} \operatorname{val}[\operatorname{supp}(\operatorname{val}(c))])} y \implies y = x)\).
\end{proofstep}
Use the hypothesis \(v^{'} = v\).
Use the hypothesis \(v \sim_{M} w \land \left(\forall y\, :\, V,\; v \sim_{M} y \implies y = w\right)\).
Substituting equals, we obtain \(\exists! w_{2}\, :\, V,\; v \sim_{\operatorname{Simple\,Graph\mathord{:}Subgraph\mathord{:}induce}(M,\allowbreak{} \operatorname{val}[\operatorname{supp}(\operatorname{val}(c))])} w_{2}\).
\end{proofstep}
Using the two conjuncts, we obtain \(\exists! w_{2}\, :\, V,\; v \sim_{\operatorname{Simple\,Graph\mathord{:}Subgraph\mathord{:}induce}(M,\allowbreak{} \operatorname{val}[\operatorname{supp}(\operatorname{val}(c))])} w_{2}\).
\end{proofstep}
Use the hypothesis \(w_{1} \in \operatorname{supp}\left(\operatorname{val}(c)\right) \land \operatorname{val}(w_{1}) = v\).
Using the underlying element and its defining property, we obtain \(\exists! w_{2}\, :\, V,\; v \sim_{\operatorname{Simple\,Graph\mathord{:}Subgraph\mathord{:}induce}(M,\allowbreak{} \operatorname{val}[\operatorname{supp}(\operatorname{val}(c))])} w_{2}\).
\end{proofstep}
Using the existential hypothesis, we obtain \(\exists! w_{1}\, :\, V,\; v \sim_{\operatorname{Simple\,Graph\mathord{:}Subgraph\mathord{:}induce}(M,\allowbreak{} \operatorname{val}[\operatorname{supp}(\operatorname{val}(c))])} w_{1}\).
\end{proofstep}
Using the existential hypothesis, we obtain \(\exists! w\, :\, V,\; v \sim_{\operatorname{Simple\,Graph\mathord{:}Subgraph\mathord{:}induce}(M,\allowbreak{} \operatorname{val}[\operatorname{supp}(\operatorname{val}(c))])} w\).
By \paperlink{premise:4e61742e6e6f745f6576656e5f6966665f6f6464}{Prerequisite~P\ref{premise:4e61742e6e6f745f6576656e5f6966665f6f6464}}, we obtain \(\neg \operatorname{Even}(\operatorname{ncard}(\operatorname{supp}(\operatorname{val}(c)))) \iff \operatorname{Odd}(\operatorname{ncard}(\operatorname{supp}(\operatorname{val}(c))))\).
By \paperlink{premise:537562747970652e70726f70}{Prerequisite~P\ref{premise:537562747970652e70726f70}} with \(\alpha = \operatorname{Connected\,Component}(\operatorname{graph}(\operatorname{delete\,Verts}(\operatorname{full\,subgraph}(G),\allowbreak{} u)))\) and \(x = c\), we obtain \(\operatorname{val}(c) \in \operatorname{odd\,Components}(\operatorname{graph}(\operatorname{delete\,Verts}(\operatorname{full\,subgraph}(G),\allowbreak{} u)))\).
The substitution uses the following argument.
\begin{proofstep}
Fix \(a\) in \(V(\operatorname{delete\,Verts}\left(\operatorname{full\,subgraph}(G), u\right))\).
Assume \(a \in \operatorname{supp}\left(\operatorname{val}(c)\right)\).
By definition, we obtain \(\operatorname{val}(a) = \operatorname{val}(a)\).
\end{proofstep}
By \paperlink{theorem:18}{Lemma~\ref{theorem:18}}, we obtain \(\operatorname{Even}(|\operatorname{as\,finite\,set}(V(\operatorname{Simple\,Graph\mathord{:}Subgraph\mathord{:}induce}(M,\allowbreak{} \operatorname{val}[\operatorname{supp}(\operatorname{val}(c))])))|)\).
Rewriting with \paperlink{premise:5365742e746f46696e7365745f636f6e6772}{Prerequisite~P\ref{premise:5365742e746f46696e7365745f636f6e6772}}, \paperlink{premise:5365742e696d6167655f636f6e6772}{Prerequisite~P\ref{premise:5365742e696d6167655f636f6e6772}}, \paperlink{premise:5365742e746f46696e7365745f696d616765}{Prerequisite~P\ref{premise:5365742e746f46696e7365745f696d616765}}, \paperlink{premise:46696e7365742e636172645f696d6167655f6f665f696e6a656374697665}{Prerequisite~P\ref{premise:46696e7365742e636172645f696d6167655f6f665f696e6a656374697665}}, \paperlink{premise:537562747970652e76616c5f696e6a656374697665}{Prerequisite~P\ref{premise:537562747970652e76616c5f696e6a656374697665}}, \paperlink{premise:5365742e746f46696e7365745f63617264}{Prerequisite~P\ref{premise:5365742e746f46696e7365745f63617264}} and the identity established above in the preceding result gives \(\operatorname{Even}\left(\left|\operatorname{supp}\left(\operatorname{val}(c)\right)\right|\right)\).
Rewriting with \paperlink{premise:5365742e4e61742e636172645f636f655f7365745f6571}{Prerequisite~P\ref{premise:5365742e4e61742e636172645f636f655f7365745f6571}} and \paperlink{premise:4e61742e636172645f65715f66696e747970655f63617264}{Prerequisite~P\ref{premise:4e61742e636172645f65715f66696e747970655f63617264}} in the preceding result gives \(\operatorname{Even}\left(\operatorname{ncard}(\operatorname{supp}\left(\operatorname{val}(c)\right))\right)\).
By the reverse implication, we obtain \(\mathrm{False}\).
\end{proofstep}
The contrary assumption leads to a contradiction. Hence \(\exists w\, :\, V,\; w \in u \land (\exists v\, :\, V(\operatorname{delete\,Verts}(\operatorname{full\,subgraph}(G),\allowbreak{} u)),\; v \sim_{M} w \land v \in \operatorname{supp}(\operatorname{val}(c)))\).
\end{proof}

\section{Simple graphs}
\begin{theorem}[{Mono}]\label{theorem:28}
Let \(V\) be a type. Let \(G\) and \(G^{'}\) be simple graphs on \(V\). Assume \(G \subseteq G^{'}\). Assume \(\operatorname{Is\,Matching\,Free}\left(G^{'}\right)\). Then \(\operatorname{Is\,Matching\,Free}\left(G\right)\).
\end{theorem}

\begin{proof}
Fix \(x\) in \(\operatorname{Subgraph}\left(G\right)\).
Assume \(x\text{ is a perfect matching}\).
Use the hypothesis \(x\text{ is a perfect matching}\).
Taking the first conjunct, we obtain \(x\text{ is a matching}\).
By \paperlink{theorem:4}{Lemma~\ref{theorem:4}}, we obtain \(\operatorname{Simple\,Graph\mathord{:}Subgraph\mathord{:}map}(\operatorname{incl}_{G,G^{'}},\allowbreak{} x)\text{ is a matching}\).
For the required implication, argue as follows.
\begin{proofstep}
Fix \(v\) in \(V\).
The substitution uses the following argument.
\begin{proofstep}
Fix \(a\) in \(V\).
Assume \(a \in V(x)\).
By definition, we obtain \(a = a\).
\end{proofstep}
Use the hypothesis \(x\text{ is a perfect matching}\).
Taking the second conjunct, we obtain \(v \in V(x)\).
Rewriting with \paperlink{premise:5365742e696d6167655f636f6e6772}{Prerequisite~P\ref{premise:5365742e696d6167655f636f6e6772}}, \paperlink{premise:5365742e696d6167655f696427}{Prerequisite~P\ref{premise:5365742e696d6167655f696427}} and the identity established above in the preceding result gives \(v \in V(\operatorname{Simple\,Graph\mathord{:}Subgraph\mathord{:}map}(\operatorname{incl}_{G,G^{'}},\allowbreak{} x))\).
\end{proofstep}
Combining the two conclusions, we obtain \(\operatorname{Simple\,Graph\mathord{:}Subgraph\mathord{:}map}(\operatorname{incl}_{G,G^{'}},\allowbreak{} x)\text{ is a matching} \land \operatorname{Is\,Spanning}(\operatorname{Simple\,Graph\mathord{:}Subgraph\mathord{:}map}(\operatorname{incl}_{G,G^{'}},\allowbreak{} x))\).
Applying the hypothesis \(\operatorname{Is\,Matching\,Free}\left(G^{'}\right)\), we obtain \(\mathrm{False}\).
\end{proof}

\begin{theorem}[{Exists maximal is matching free}]\label{theorem:29}
Let \(V\) be a type. Let \(G\) be a simple graph on \(V\). Assume \(V\) is finite. Assume \(\operatorname{Is\,Matching\,Free}\left(G\right)\). Then
\begin{center}
\(\displaystyle \exists \mathit{Gmax}\, :\, \operatorname{Simple\,Graph}(V),\; G \subseteq \mathit{Gmax} \land (\operatorname{Is\,Matching\,Free}(\mathit{Gmax}) \land (\forall G^{'}\, :\, \operatorname{Simple\,Graph}(V),\; G^{'} > \mathit{Gmax} \implies (\exists M\, :\, \operatorname{Subgraph}(G^{'}),\; M\text{ is a perfect matching}))).\)
\end{center}
\end{theorem}

\begin{proof}
The substitution uses the following argument.
\begin{proofstep}
Fix \(\mathit{Gmax}\) in \(\operatorname{Simple\,Graph}(V)\).
For the required implication, argue as follows.
\begin{proofstep}
Fix \(G^{'}\) in \(\operatorname{Simple\,Graph}(V)\).
By definition, we obtain \(\left(G^{'} > \mathit{Gmax}\right) = \left(G^{'} > \mathit{Gmax}\right)\).
By \paperlink{premise:6e6f745f666f72616c6c5f6e6f74}{Prerequisite~P\ref{premise:6e6f745f666f72616c6c5f6e6f74}} with \(\alpha = \operatorname{Subgraph}\left(G^{'}\right)\), we obtain \(\neg (\forall x\, :\, \operatorname{Subgraph}(G^{'}),\; \neg (x\text{ is a perfect matching})) \iff (\exists x\, :\, \operatorname{Subgraph}(G^{'}),\; x\text{ is a perfect matching})\).
By propositional extensionality, we obtain \((\neg (\forall x\, :\, \operatorname{Subgraph}(G^{'}),\; \neg (x\text{ is a perfect matching}))) = (\exists x\, :\, \operatorname{Subgraph}(G^{'}),\; x\text{ is a perfect matching})\).
Reversing the equality, we obtain \((\exists x\, :\, \operatorname{Subgraph}(G^{'}),\; x\text{ is a perfect matching}) = (\neg (\forall x\, :\, \operatorname{Subgraph}(G^{'}),\; \neg (x\text{ is a perfect matching})))\).
By \paperlink{premise:696d706c6965735f636f6e6772}{Prerequisite~P\ref{premise:696d706c6965735f636f6e6772}} with \(p_{1} = G^{'} > \mathit{Gmax}\) and \(p_{2} = G^{'} > \mathit{Gmax}\) and \(q_{1} = \exists x\, :\, \operatorname{Subgraph}\left(G^{'}\right),\; x\text{ is a perfect matching}\) and \(q_{2} = \neg (\forall x\, :\, \operatorname{Subgraph}(G^{'}),\; \neg (x\text{ is a perfect matching}))\), we obtain \((G^{'} > \mathit{Gmax} \implies (\exists x\, :\, \operatorname{Subgraph}(G^{'}),\; x\text{ is a perfect matching})) = (G^{'} > \mathit{Gmax} \implies \neg (\forall x\, :\, \operatorname{Subgraph}(G^{'}),\; \neg (x\text{ is a perfect matching})))\).
\end{proofstep}
By \paperlink{premise:666f72616c6c5f636f6e6772}{Prerequisite~P\ref{premise:666f72616c6c5f636f6e6772}} with \(\alpha = \operatorname{Simple\,Graph}(V)\), we obtain \((\forall a\, :\, \operatorname{Simple\,Graph}(V),\; a > \mathit{Gmax} \implies (\exists x\, :\, \operatorname{Subgraph}(a),\; x\text{ is a perfect matching})) = (\forall a\, :\, \operatorname{Simple\,Graph}(V),\; a > \mathit{Gmax} \implies \neg (\forall x\, :\, \operatorname{Subgraph}(a),\; \neg (x\text{ is a perfect matching})))\).
Applying the same function to equal arguments, we obtain \((G \subseteq \mathit{Gmax} \land (\operatorname{Is\,Matching\,Free}(\mathit{Gmax}) \land (\forall a\, :\, \operatorname{Simple\,Graph}(V),\; a > \mathit{Gmax} \implies (\exists x\, :\, \operatorname{Subgraph}(a),\; x\text{ is a perfect matching})))) = (G \subseteq \mathit{Gmax} \land (\operatorname{Is\,Matching\,Free}(\mathit{Gmax}) \land (\forall a\, :\, \operatorname{Simple\,Graph}(V),\; a > \mathit{Gmax} \implies \neg (\forall x\, :\, \operatorname{Subgraph}(a),\; \neg (x\text{ is a perfect matching})))))\).
\end{proofstep}
Use the hypothesis \(\operatorname{Finite}\left(V\right)\).
By \paperlink{premise:53696d706c6547726170682e696e737446696e697465}{Prerequisite~P\ref{premise:53696d706c6547726170682e696e737446696e697465}}, we obtain \(\operatorname{Finite}\left(\operatorname{Simple\,Graph}(V)\right)\).
Use the hypothesis \(\operatorname{Is\,Matching\,Free}\left(G\right)\).
By \paperlink{premise:46696e6974652e6578697374735f6c655f6d6178696d616c}{Prerequisite~P\ref{premise:46696e6974652e6578697374735f6c655f6d6178696d616c}} with \(\alpha = \operatorname{Simple\,Graph}(V)\), we obtain \(\exists b\, :\, \operatorname{Simple\,Graph}(V),\; G \subseteq b \land \operatorname{Maximal}(\operatorname{Is\,Matching\,Free},\allowbreak{} b)\).
Consider the following case.
\begin{proofstep}
Fix \(\mathit{Gmax}\) in \(\operatorname{Simple\,Graph}(V)\).
Assume \(G \subseteq \mathit{Gmax} \land \operatorname{Maximal}(\operatorname{Is\,Matching\,Free},\allowbreak{} \mathit{Gmax})\).
Choose the witness \(\mathit{Gmax}\).
Use the hypothesis \(G \subseteq \mathit{Gmax} \land \operatorname{Maximal}(\operatorname{Is\,Matching\,Free},\allowbreak{} \mathit{Gmax})\).
Taking the first conjunct, we obtain \(G \subseteq \mathit{Gmax}\).
Use the hypothesis \(G \subseteq \mathit{Gmax} \land \operatorname{Maximal}(\operatorname{Is\,Matching\,Free},\allowbreak{} \mathit{Gmax})\).
Taking the second conjunct, we obtain \(\operatorname{Maximal}\left(\operatorname{Is\,Matching\,Free}, \mathit{Gmax}\right)\).
By \paperlink{premise:4d6178696d616c2e70726f70}{Prerequisite~P\ref{premise:4d6178696d616c2e70726f70}} with \(\alpha = \operatorname{Simple\,Graph}(V)\), we obtain \(\operatorname{Is\,Matching\,Free}\left(\mathit{Gmax}\right)\).
For the required implication, argue as follows.
\begin{proofstep}
Fix \(x\) in \(\operatorname{Simple\,Graph}(V)\).
Assume \(x > \mathit{Gmax}\).
Use the hypothesis \(G \subseteq \mathit{Gmax} \land \operatorname{Maximal}(\operatorname{Is\,Matching\,Free},\allowbreak{} \mathit{Gmax})\).
Taking the second conjunct, we obtain \(\operatorname{Maximal}\left(\operatorname{Is\,Matching\,Free}, \mathit{Gmax}\right)\).
Use the hypothesis \(x > \mathit{Gmax}\).
By \paperlink{premise:4d6178696d616c2e6e6f745f70726f705f6f665f6774}{Prerequisite~P\ref{premise:4d6178696d616c2e6e6f745f70726f705f6f665f6774}} with \(\alpha = \operatorname{Simple\,Graph}(V)\), we obtain \(\neg \operatorname{Is\,Matching\,Free}\left(x\right)\).
\end{proofstep}
Combining the two conclusions, we obtain \(\operatorname{Is\,Matching\,Free}(\mathit{Gmax}) \land (\forall G^{'}\, :\, \operatorname{Simple\,Graph}(V),\; G^{'} > \mathit{Gmax} \implies \neg (\forall x\, :\, \operatorname{Subgraph}(G^{'}),\; \neg (x\text{ is a perfect matching})))\).
Combining the two conclusions, we obtain \(G \subseteq \mathit{Gmax} \land (\operatorname{Is\,Matching\,Free}(\mathit{Gmax}) \land (\forall G^{'}\, :\, \operatorname{Simple\,Graph}(V),\; G^{'} > \mathit{Gmax} \implies \neg (\forall x\, :\, \operatorname{Subgraph}(G^{'}),\; \neg (x\text{ is a perfect matching}))))\).
This witness proves existence. Hence \(\exists \mathit{Gmax}_{1}\, :\, \operatorname{Simple\,Graph}(V),\; G \subseteq \mathit{Gmax}_{1} \land (\operatorname{Is\,Matching\,Free}(\mathit{Gmax}_{1}) \land (\forall G^{'}\, :\, \operatorname{Simple\,Graph}(V),\; G^{'} > \mathit{Gmax}_{1} \implies \neg (\forall x\, :\, \operatorname{Subgraph}(G^{'}),\; \neg (x\text{ is a perfect matching}))))\).
\end{proofstep}
Using the existential hypothesis, we obtain \(\exists \mathit{Gmax}\, :\, \operatorname{Simple\,Graph}(V),\; G \subseteq \mathit{Gmax} \land (\operatorname{Is\,Matching\,Free}(\mathit{Gmax}) \land (\forall G^{'}\, :\, \operatorname{Simple\,Graph}(V),\; G^{'} > \mathit{Gmax} \implies \neg (\forall x\, :\, \operatorname{Subgraph}(G^{'}),\; \neg (x\text{ is a perfect matching}))))\).
Rewriting with \paperlink{premise:66756e657874}{Prerequisite~P\ref{premise:66756e657874}} and the identity established above in the preceding result gives \(\exists \mathit{Gmax}\, :\, \operatorname{Simple\,Graph}(V),\; G \subseteq \mathit{Gmax} \land (\operatorname{Is\,Matching\,Free}(\mathit{Gmax}) \land (\forall G^{'}\, :\, \operatorname{Simple\,Graph}(V),\; G^{'} > \mathit{Gmax} \implies (\exists M\, :\, \operatorname{Subgraph}(G^{'}),\; M\text{ is a perfect matching})))\).
\end{proof}

\begin{lemma}[{Other adj}]\label{theorem:30}
Let \(V\) be a type. Let \(G\) be a simple graph on \(V\). Let \(v\) and \(w\) be elements of \(V\). Assume \(\text{every nonisolated vertex of }G\text{ has degree two}\). Assume \(v \sim_{G} w\). Then \(\exists w^{'}\, :\, V,\; w \ne w^{'} \land v \sim_{G} w^{'}\).
\end{lemma}

\begin{proof}
The substitution uses the following argument.
\begin{proofstep}
Fix \(w^{'}\) in \(V\).
By \paperlink{premise:53696d706c6547726170682e6d656d5f6e65696768626f72536574}{Prerequisite~P\ref{premise:53696d706c6547726170682e6d656d5f6e65696768626f72536574}} with \(w^{'}\) in place of its parameter \(w\), we obtain \(w^{'} \in N_{G}(v) \iff v \sim_{G} w^{'}\).
By propositional extensionality, we obtain \(\left(w^{'} \in N_{G}(v)\right) = \left(v \sim_{G} w^{'}\right)\).
Reversing the equality, we obtain \(\left(v \sim_{G} w^{'}\right) = \left(w^{'} \in N_{G}(v)\right)\).
Applying the same function to equal arguments, we obtain \(\left(w \ne w^{'} \land v \sim_{G} w^{'}\right) = \left(w \ne w^{'} \land w^{'} \in N_{G}(v)\right)\).
\end{proofstep}
Choose the witness \(w\).
Use the hypothesis \(v \sim_{G} w\).
Rewriting with \paperlink{premise:53696d706c6547726170682e6d656d5f6e65696768626f72536574}{Prerequisite~P\ref{premise:53696d706c6547726170682e6d656d5f6e65696768626f72536574}} in the preceding result gives \(w \in N_{G}(v)\).
This witness proves existence. Hence \(\exists x\, :\, V,\; x \in N_{G}(v)\).
Applying the hypothesis \(\text{every nonisolated vertex of }G\text{ has degree two}\), we obtain \(\operatorname{ncard}(N_{G}(v)) = 2\).
Consider the following case.
\begin{proofstep}
Assume \(\neg \left(1 < \operatorname{ncard}(N_{G}(v))\right)\).
By definition, we obtain \(\mathrm{true} = \mathrm{true}\).
By the preceding argument, we obtain \(\operatorname{decide}(\operatorname{is\,Impossible}(\operatorname{Lean\mathord{:}Omega\mathord{:}Constraint\mathord{:}mk}(\operatorname{Option\mathord{:}some}(1),\allowbreak{} \operatorname{Option\mathord{:}some}(0))) = \mathrm{true}) = \mathrm{true}\).
By \paperlink{premise:6f665f6465636964655f65715f74727565}{Prerequisite~P\ref{premise:6f665f6465636964655f65715f74727565}}, we obtain \(\operatorname{is\,Impossible}(\operatorname{Lean\mathord{:}Omega\mathord{:}Constraint\mathord{:}mk}(\operatorname{Option\mathord{:}some}(1),\allowbreak{} \operatorname{Option\mathord{:}some}(0))) = \mathrm{true}\).
Using \paperlink{premise:4c65616e2e4f6d6567612e4c696e656172436f6d626f2e636f6f7264696e6174655f6576616c5f30}{Prerequisite~P\ref{premise:4c65616e2e4f6d6567612e4c696e656172436f6d626f2e636f6f7264696e6174655f6576616c5f30}} in reverse gives \(\operatorname{ncard}(N_{G}(v)) = \operatorname{eval}(\operatorname{coordinate}(0),\allowbreak{} \operatorname{of\,List}(\langle \operatorname{ncard}(N_{G}(v))\rangle))\).
By definition, we obtain \(2 = 2\).
By the preceding argument, we obtain \(2 = \operatorname{eval}(\operatorname{Lean\mathord{:}Omega\mathord{:}Linear\,Combo\mathord{:}mk}(2,\allowbreak{} \langle\,\rangle),\allowbreak{} \operatorname{of\,List}(\langle \operatorname{ncard}(N_{G}(v))\rangle))\).
By \paperlink{premise:4c65616e2e4f6d6567612e496e742e7375625f636f6e6772}{Prerequisite~P\ref{premise:4c65616e2e4f6d6567612e496e742e7375625f636f6e6772}}, we obtain \(\operatorname{ncard}(N_{G}(v)) - 2 = \operatorname{eval}(\operatorname{coordinate}(0),\allowbreak{} \operatorname{of\,List}(\langle \operatorname{ncard}(N_{G}(v))\rangle)) - \operatorname{eval}(\operatorname{Lean\mathord{:}Omega\mathord{:}Linear\,Combo\mathord{:}mk}(2,\allowbreak{} \langle\,\rangle),\allowbreak{} \operatorname{of\,List}(\langle \operatorname{ncard}(N_{G}(v))\rangle))\).
Using \paperlink{premise:4c65616e2e4f6d6567612e4c696e656172436f6d626f2e7375625f6576616c}{Prerequisite~P\ref{premise:4c65616e2e4f6d6567612e4c696e656172436f6d626f2e7375625f6576616c}} with \(l_{1} = \operatorname{Lean\mathord{:}Omega\mathord{:}Linear\,Combo\mathord{:}mk}(0,\allowbreak{} \langle 1\rangle)\) and \(l_{2} = \operatorname{Lean\mathord{:}Omega\mathord{:}Linear\,Combo\mathord{:}mk}(2,\allowbreak{} \langle\,\rangle)\) and \(\operatorname{of\,List}\left(\langle \operatorname{ncard}(N_{G}(v))\rangle\right)\) in place of its parameter \(v\) in reverse gives \(\operatorname{eval}(\operatorname{Lean\mathord{:}Omega\mathord{:}Linear\,Combo\mathord{:}mk}(0,\allowbreak{} \langle 1\rangle),\allowbreak{} \operatorname{of\,List}(\langle \operatorname{ncard}(N_{G}(v))\rangle)) - \operatorname{eval}(\operatorname{Lean\mathord{:}Omega\mathord{:}Linear\,Combo\mathord{:}mk}(2,\allowbreak{} \langle\,\rangle),\allowbreak{} \operatorname{of\,List}(\langle \operatorname{ncard}(N_{G}(v))\rangle)) = \operatorname{eval}(\operatorname{Lean\mathord{:}Omega\mathord{:}Linear\,Combo\mathord{:}mk}(0,\allowbreak{} \langle 1\rangle) - \operatorname{Lean\mathord{:}Omega\mathord{:}Linear\,Combo\mathord{:}mk}(2,\allowbreak{} \langle\,\rangle),\allowbreak{} \operatorname{of\,List}(\langle \operatorname{ncard}(N_{G}(v))\rangle))\).
Chaining the equalities, we obtain \(\operatorname{ncard}(N_{G}(v)) - 2 = \operatorname{eval}(\operatorname{Lean\mathord{:}Omega\mathord{:}Linear\,Combo\mathord{:}mk}(0,\allowbreak{} \langle 1\rangle) - \operatorname{Lean\mathord{:}Omega\mathord{:}Linear\,Combo\mathord{:}mk}(2,\allowbreak{} \langle\,\rangle),\allowbreak{} \operatorname{of\,List}(\langle \operatorname{ncard}(N_{G}(v))\rangle))\).
Reversing the equality, we obtain \(\operatorname{eval}(\operatorname{Lean\mathord{:}Omega\mathord{:}Linear\,Combo\mathord{:}mk}(0,\allowbreak{} \langle 1\rangle) - \operatorname{Lean\mathord{:}Omega\mathord{:}Linear\,Combo\mathord{:}mk}(2,\allowbreak{} \langle\,\rangle),\allowbreak{} \operatorname{of\,List}(\langle \operatorname{ncard}(N_{G}(v))\rangle)) = \operatorname{ncard}(N_{G}(v)) - 2\).
Use the previously proved fact \(\operatorname{ncard}(N_{G}(v)) = 2\).
By \paperlink{premise:4c65616e2e4f6d6567612e496e742e6f664e61745f636f6e6772}{Prerequisite~P\ref{premise:4c65616e2e4f6d6567612e496e742e6f664e61745f636f6e6772}}, we obtain \(\operatorname{ncard}(N_{G}(v)) = 2\).
By \paperlink{premise:496e742e7375625f65715f7a65726f5f6f665f6571}{Prerequisite~P\ref{premise:496e742e7375625f65715f7a65726f5f6f665f6571}}, we obtain \(\operatorname{ncard}(N_{G}(v)) - 2 = 0\).
Chaining the equalities, we obtain \(\operatorname{eval}(\operatorname{Lean\mathord{:}Omega\mathord{:}Linear\,Combo\mathord{:}mk}(0,\allowbreak{} \langle 1\rangle) - \operatorname{Lean\mathord{:}Omega\mathord{:}Linear\,Combo\mathord{:}mk}(2,\allowbreak{} \langle\,\rangle),\allowbreak{} \operatorname{of\,List}(\langle \operatorname{ncard}(N_{G}(v))\rangle)) = 0\).
By \paperlink{premise:4c65616e2e4f6d6567612e436f6e73747261696e742e616464457175616c6974795f736174}{Prerequisite~P\ref{premise:4c65616e2e4f6d6567612e436f6e73747261696e742e616464457175616c6974795f736174}}, we obtain \(\operatorname{sat^{\prime}}(\operatorname{Lean\mathord{:}Omega\mathord{:}Constraint\mathord{:}mk}(\operatorname{Option\mathord{:}some}(-(-2)),\allowbreak{} \operatorname{Option\mathord{:}some}(-(-2))),\allowbreak{} \langle 1\rangle,\allowbreak{} \operatorname{of\,List}(\langle \operatorname{ncard}(N_{G}(v))\rangle)) = \mathrm{true}\).
Use the hypothesis \(\neg \left(1 < \operatorname{ncard}(N_{G}(v))\right)\).
By \paperlink{premise:4e61742e6c655f6f665f6e6f745f6c74}{Prerequisite~P\ref{premise:4e61742e6c655f6f665f6e6f745f6c74}}, we obtain \(\operatorname{ncard}(N_{G}(v)) \le 1\).
By \paperlink{premise:4c65616e2e4f6d6567612e496e742e6f664e61745f6c655f6f665f6c65}{Prerequisite~P\ref{premise:4c65616e2e4f6d6567612e496e742e6f664e61745f6c655f6f665f6c65}}, we obtain \(\operatorname{ncard}(N_{G}(v)) \le 1\).
By \paperlink{premise:496e742e7375625f6e6f6e6e65675f6f665f6c65}{Prerequisite~P\ref{premise:496e742e7375625f6e6f6e6e65675f6f665f6c65}}, we obtain \(0 \le 1 - \operatorname{ncard}(N_{G}(v))\).
By definition, we obtain \(1 = 1\).
By the preceding argument, we obtain \(1 = \operatorname{eval}(\operatorname{Lean\mathord{:}Omega\mathord{:}Linear\,Combo\mathord{:}mk}(1,\allowbreak{} \langle\,\rangle),\allowbreak{} \operatorname{of\,List}(\langle \operatorname{ncard}(N_{G}(v))\rangle))\).
Using \paperlink{premise:4c65616e2e4f6d6567612e4c696e656172436f6d626f2e636f6f7264696e6174655f6576616c5f30}{Prerequisite~P\ref{premise:4c65616e2e4f6d6567612e4c696e656172436f6d626f2e636f6f7264696e6174655f6576616c5f30}} in reverse gives \(\operatorname{ncard}(N_{G}(v)) = \operatorname{eval}(\operatorname{coordinate}(0),\allowbreak{} \operatorname{of\,List}(\langle \operatorname{ncard}(N_{G}(v))\rangle))\).
By \paperlink{premise:4c65616e2e4f6d6567612e496e742e7375625f636f6e6772}{Prerequisite~P\ref{premise:4c65616e2e4f6d6567612e496e742e7375625f636f6e6772}}, we obtain \(1 - \operatorname{ncard}(N_{G}(v)) = \operatorname{eval}(\operatorname{Lean\mathord{:}Omega\mathord{:}Linear\,Combo\mathord{:}mk}(1,\allowbreak{} \langle\,\rangle),\allowbreak{} \operatorname{of\,List}(\langle \operatorname{ncard}(N_{G}(v))\rangle)) - \operatorname{eval}(\operatorname{coordinate}(0),\allowbreak{} \operatorname{of\,List}(\langle \operatorname{ncard}(N_{G}(v))\rangle))\).
Using \paperlink{premise:4c65616e2e4f6d6567612e4c696e656172436f6d626f2e7375625f6576616c}{Prerequisite~P\ref{premise:4c65616e2e4f6d6567612e4c696e656172436f6d626f2e7375625f6576616c}} with \(l_{1} = \operatorname{Lean\mathord{:}Omega\mathord{:}Linear\,Combo\mathord{:}mk}(1,\allowbreak{} \langle\,\rangle)\) and \(l_{2} = \operatorname{Lean\mathord{:}Omega\mathord{:}Linear\,Combo\mathord{:}mk}(0,\allowbreak{} \langle 1\rangle)\) and \(\operatorname{of\,List}\left(\langle \operatorname{ncard}(N_{G}(v))\rangle\right)\) in place of its parameter \(v\) in reverse gives \(\operatorname{eval}(\operatorname{Lean\mathord{:}Omega\mathord{:}Linear\,Combo\mathord{:}mk}(1,\allowbreak{} \langle\,\rangle),\allowbreak{} \operatorname{of\,List}(\langle \operatorname{ncard}(N_{G}(v))\rangle)) - \operatorname{eval}(\operatorname{Lean\mathord{:}Omega\mathord{:}Linear\,Combo\mathord{:}mk}(0,\allowbreak{} \langle 1\rangle),\allowbreak{} \operatorname{of\,List}(\langle \operatorname{ncard}(N_{G}(v))\rangle)) = \operatorname{eval}(\operatorname{Lean\mathord{:}Omega\mathord{:}Linear\,Combo\mathord{:}mk}(1,\allowbreak{} \langle\,\rangle) - \operatorname{Lean\mathord{:}Omega\mathord{:}Linear\,Combo\mathord{:}mk}(0,\allowbreak{} \langle 1\rangle),\allowbreak{} \operatorname{of\,List}(\langle \operatorname{ncard}(N_{G}(v))\rangle))\).
Chaining the equalities, we obtain \(1 - \operatorname{ncard}(N_{G}(v)) = \operatorname{eval}(\operatorname{Lean\mathord{:}Omega\mathord{:}Linear\,Combo\mathord{:}mk}(1,\allowbreak{} \langle\,\rangle) - \operatorname{Lean\mathord{:}Omega\mathord{:}Linear\,Combo\mathord{:}mk}(0,\allowbreak{} \langle 1\rangle),\allowbreak{} \operatorname{of\,List}(\langle \operatorname{ncard}(N_{G}(v))\rangle))\).
By \paperlink{premise:6c655f6f665f6c655f6f665f6571}{Prerequisite~P\ref{premise:6c655f6f665f6c655f6f665f6571}} with \(\alpha = \mathbb{Z}\), we obtain \(0 \le \operatorname{eval}(\operatorname{Lean\mathord{:}Omega\mathord{:}Linear\,Combo\mathord{:}mk}(1,\allowbreak{} \langle\,\rangle) - \operatorname{Lean\mathord{:}Omega\mathord{:}Linear\,Combo\mathord{:}mk}(0,\allowbreak{} \langle 1\rangle),\allowbreak{} \operatorname{of\,List}(\langle \operatorname{ncard}(N_{G}(v))\rangle))\).
By \paperlink{premise:4c65616e2e4f6d6567612e436f6e73747261696e742e616464496e657175616c6974795f736174}{Prerequisite~P\ref{premise:4c65616e2e4f6d6567612e436f6e73747261696e742e616464496e657175616c6974795f736174}}, we obtain \(\operatorname{sat^{\prime}}(\operatorname{Lean\mathord{:}Omega\mathord{:}Constraint\mathord{:}mk}(\operatorname{Option\mathord{:}some}(-1),\allowbreak{} \operatorname{none}),\allowbreak{} \langle -1\rangle,\allowbreak{} \operatorname{of\,List}(\langle \operatorname{ncard}(N_{G}(v))\rangle)) = \mathrm{true}\).
By \paperlink{premise:4c65616e2e4f6d6567612e636f6d626f5f73617427}{Prerequisite~P\ref{premise:4c65616e2e4f6d6567612e636f6d626f5f73617427}} with \(s = \operatorname{Lean\mathord{:}Omega\mathord{:}Constraint\mathord{:}mk}(\operatorname{Option\mathord{:}some}(2),\allowbreak{} \operatorname{Option\mathord{:}some}(2))\) and \(t = \operatorname{Lean\mathord{:}Omega\mathord{:}Constraint\mathord{:}mk}(\operatorname{Option\mathord{:}some}(-1),\allowbreak{} \operatorname{none})\) and \(1\) in place of its parameter \(a\) and \(x = \langle 1\rangle\) and \(b = 1\) and \(y = \langle -1\rangle\) and \(\operatorname{of\,List}\left(\langle \operatorname{ncard}(N_{G}(v))\rangle\right)\) in place of its parameter \(v\), we obtain \(\operatorname{sat^{\prime}}(\operatorname{Lean\mathord{:}Omega\mathord{:}Constraint\mathord{:}combo}(1,\allowbreak{} \operatorname{Lean\mathord{:}Omega\mathord{:}Constraint\mathord{:}mk}(\operatorname{Option\mathord{:}some}(2),\allowbreak{} \operatorname{Option\mathord{:}some}(2)),\allowbreak{} 1,\allowbreak{} \operatorname{Lean\mathord{:}Omega\mathord{:}Constraint\mathord{:}mk}(\operatorname{Option\mathord{:}some}(-1),\allowbreak{} \operatorname{none})),\allowbreak{} \operatorname{Lean\mathord{:}Omega\mathord{:}Coeffs\mathord{:}combo}(1,\allowbreak{} \langle 1\rangle,\allowbreak{} 1,\allowbreak{} \langle -1\rangle),\allowbreak{} \operatorname{of\,List}(\langle \operatorname{ncard}(N_{G}(v))\rangle)) = \mathrm{true}\).
By \paperlink{premise:4c65616e2e4f6d6567612e746964795f736174}{Prerequisite~P\ref{premise:4c65616e2e4f6d6567612e746964795f736174}}, we obtain \(\operatorname{sat^{\prime}}(\operatorname{tidy\,Constraint}(\operatorname{Lean\mathord{:}Omega\mathord{:}Constraint\mathord{:}mk}(\operatorname{Option\mathord{:}some}(1),\allowbreak{} \operatorname{none}),\allowbreak{} \langle 0\rangle),\allowbreak{} \operatorname{tidy\,Coeffs}(\operatorname{Lean\mathord{:}Omega\mathord{:}Constraint\mathord{:}mk}(\operatorname{Option\mathord{:}some}(1),\allowbreak{} \operatorname{none}),\allowbreak{} \langle 0\rangle),\allowbreak{} \operatorname{of\,List}(\langle \operatorname{ncard}(N_{G}(v))\rangle)) = \mathrm{true}\).
By \paperlink{premise:4c65616e2e4f6d6567612e436f6e73747261696e742e6e6f745f736174275f6f665f6973496d706f737369626c65}{Prerequisite~P\ref{premise:4c65616e2e4f6d6567612e436f6e73747261696e742e6e6f745f736174275f6f665f6973496d706f737369626c65}}, we obtain \(\mathrm{False}\).
\end{proofstep}
The contrary assumption leads to a contradiction. Hence \(1 < \operatorname{ncard}(N_{G}(v))\).
By \paperlink{premise:5365742e6578697374735f6e655f6f665f6f6e655f6c745f6e63617264}{Prerequisite~P\ref{premise:5365742e6578697374735f6e655f6f665f6f6e655f6c745f6e63617264}} with \(a = w\), we obtain \(\exists b\, :\, V,\; b \in N_{G}(v) \land b \ne w\).
Consider the following case.
\begin{proofstep}
Fix \(w^{'}\) in \(V\).
Assume \(w^{'} \in N_{G}(v) \land w^{'} \ne w\).
Choose the witness \(w^{'}\).
Use the hypothesis \(w^{'} \in N_{G}(v) \land w^{'} \ne w\).
Taking the second conjunct, we obtain \(w^{'} \ne w\).
By \paperlink{premise:4e652e73796d6d}{Prerequisite~P\ref{premise:4e652e73796d6d}}, we obtain \(w \ne w^{'}\).
Use the hypothesis \(w^{'} \in N_{G}(v) \land w^{'} \ne w\).
Taking the first conjunct, we obtain \(w^{'} \in N_{G}(v)\).
Combining the two conclusions, we obtain \(w \ne w^{'} \land w^{'} \in N_{G}(v)\).
This witness proves existence. Hence \(\exists w^{'}_{1}\, :\, V,\; w \ne w^{'}_{1} \land w^{'}_{1} \in N_{G}(v)\).
\end{proofstep}
Using the existential hypothesis, we obtain \(\exists w^{'}\, :\, V,\; w \ne w^{'} \land w^{'} \in N_{G}(v)\).
Rewriting with \paperlink{premise:66756e657874}{Prerequisite~P\ref{premise:66756e657874}} and the identity established above in the preceding result gives \(\exists w^{'}\, :\, V,\; w \ne w^{'} \land v \sim_{G} w^{'}\).
\end{proof}

\begin{lemma}[{Exists unique differs from adj}]\label{theorem:31}
Let \(V\) be a type. Let \(G\) be a simple graph on \(V\). Let \(v\) and \(w\) be elements of \(V\). Assume \(\text{every nonisolated vertex of }G\text{ has degree two}\). Assume \(v \sim_{G} w\). Then \(\exists! w^{'}\, :\, V,\; w \ne w^{'} \land v \sim_{G} w^{'}\).
\end{lemma}

\begin{proof}
By \paperlink{theorem:30}{Lemma~\ref{theorem:30}}, we obtain \(\exists w^{'}\, :\, V,\; w \ne w^{'} \land v \sim_{G} w^{'}\).
Consider the following case.
\begin{proofstep}
Fix \(w^{'}\) in \(V\).
Assume \(w \ne w^{'} \land v \sim_{G} w^{'}\).
Use the hypothesis \(w \ne w^{'} \land v \sim_{G} w^{'}\).
Consider the following case.
\begin{proofstep}
Assume \(w \ne w^{'}\).
Assume \(v \sim_{G} w^{'}\).
Choose the witness \(w^{'}\).
Use the hypothesis \(w \ne w^{'}\).
Use the hypothesis \(v \sim_{G} w^{'}\).
Combining the two conclusions, we obtain \(w \ne w^{'} \land v \sim_{G} w^{'}\).
For the required implication, argue as follows.
\begin{proofstep}
Fix \(y\) in \(V\).
Assume \(w \ne y \land v \sim_{G} y\).
Use the hypothesis \(w \ne y \land v \sim_{G} y\).
Take the witnesses from the preceding result.
\begin{proofstep}
Assume \(w \ne y\).
Assume \(v \sim_{G} y\).
By \paperlink{premise:5365742e6e636172645f65715f74776f}{Prerequisite~P\ref{premise:5365742e6e636172645f65715f74776f}}, we obtain \(\operatorname{ncard}(N_{G}(v)) = 2 \iff (\exists x\, :\, V,\; \exists y_{1}\, :\, V,\; x \ne y_{1} \land N_{G}(v) = \{x\} \cup \{y_{1}\})\).
Choose the witness \(w\).
Use the hypothesis \(v \sim_{G} w\).
This witness proves existence. Hence \(\exists x\, :\, V,\; x \in N_{G}(v)\).
Applying the hypothesis \(\text{every nonisolated vertex of }G\text{ has degree two}\), we obtain \(\operatorname{ncard}(N_{G}(v)) = 2\).
By the forward implication, we obtain \(\exists x\, :\, V,\; \exists y_{1}\, :\, V,\; x \ne y_{1} \land N_{G}(v) = \{x\} \cup \{y_{1}\}\).
Consider the following case.
\begin{proofstep}
Fix \(x\) in \(V\).
Assume \(\exists y_{1}\, :\, V,\; x \ne y_{1} \land N_{G}(v) = \{x\} \cup \{y_{1}\}\).
Use the hypothesis \(\exists y_{1}\, :\, V,\; x \ne y_{1} \land N_{G}(v) = \{x\} \cup \{y_{1}\}\).
Consider the following case.
\begin{proofstep}
Fix \(y^{'}\) in \(V\).
Assume \(x \ne y^{'} \land N_{G}(v) = \{x\} \cup \{y^{'}\}\).
Use the hypothesis \(x \ne y^{'} \land N_{G}(v) = \{x\} \cup \{y^{'}\}\).
By the preceding argument, we obtain \(\neg \left(x = y^{'}\right) \land N_{G}(v) = \{x\} \cup \{y^{'}\}\).
Consider the following case.
\begin{proofstep}
Assume \(\neg \left(x = y^{'}\right)\).
Assume \(N_{G}(v) = \{x\} \cup \{y^{'}\}\).
Use the hypothesis \(x \ne y^{'} \land N_{G}(v) = \{x\} \cup \{y^{'}\}\).
By the preceding argument, we obtain \(\neg \left(x = y^{'}\right) \land N_{G}(v) = \{x\} \cup \{y^{'}\}\).
By the preceding argument, we obtain \(\neg \left(x = y^{'}\right) \land N_{G}(v) = \{x\} \cup \{y^{'}\}\).
Taking the second conjunct gives \(N_{G}(v) = \{x\} \cup \{y^{'}\}\).
Use the hypothesis \(v \sim_{G} w\).
Rewriting with \paperlink{premise:53696d706c6547726170682e6d656d5f6e65696768626f72536574}{Prerequisite~P\ref{premise:53696d706c6547726170682e6d656d5f6e65696768626f72536574}} in the preceding result gives \(w \in N_{G}(v)\).
Rewriting with \(N_{G}(v) = \{x\} \cup \{y^{'}\}\), \paperlink{premise:5365742e6d656d5f696e736572745f696666}{Prerequisite~P\ref{premise:5365742e6d656d5f696e736572745f696666}} and \paperlink{premise:5365742e6d656d5f73696e676c65746f6e5f696666}{Prerequisite~P\ref{premise:5365742e6d656d5f73696e676c65746f6e5f696666}} in the preceding result gives \(w = x \lor w = y^{'}\).
Consider the following case.
\begin{proofstep}
Assume \(w = x\).
Use the hypothesis \(x \ne y^{'} \land N_{G}(v) = \{x\} \cup \{y^{'}\}\).
By the preceding argument, we obtain \(\neg \left(x = y^{'}\right) \land N_{G}(v) = \{x\} \cup \{y^{'}\}\).
By the preceding argument, we obtain \(\neg \left(x = y^{'}\right) \land N_{G}(v) = \{x\} \cup \{y^{'}\}\).
Taking the second conjunct gives \(N_{G}(v) = \{x\} \cup \{y^{'}\}\).
Use the hypothesis \(v \sim_{G} w^{'}\).
Rewriting with \paperlink{premise:53696d706c6547726170682e6d656d5f6e65696768626f72536574}{Prerequisite~P\ref{premise:53696d706c6547726170682e6d656d5f6e65696768626f72536574}} in the preceding result gives \(w^{'} \in N_{G}(v)\).
Rewriting with \(N_{G}(v) = \{x\} \cup \{y^{'}\}\), \paperlink{premise:5365742e6d656d5f696e736572745f696666}{Prerequisite~P\ref{premise:5365742e6d656d5f696e736572745f696666}} and \paperlink{premise:5365742e6d656d5f73696e676c65746f6e5f696666}{Prerequisite~P\ref{premise:5365742e6d656d5f73696e676c65746f6e5f696666}} in the preceding result gives \(w^{'} = x \lor w^{'} = y^{'}\).
Consider the following case.
\begin{proofstep}
Assume \(w^{'} = x\).
Use the hypothesis \(x \ne y^{'} \land N_{G}(v) = \{x\} \cup \{y^{'}\}\).
By the preceding argument, we obtain \(\neg \left(x = y^{'}\right) \land N_{G}(v) = \{x\} \cup \{y^{'}\}\).
By the preceding argument, we obtain \(\neg \left(x = y^{'}\right) \land N_{G}(v) = \{x\} \cup \{y^{'}\}\).
Taking the second conjunct gives \(N_{G}(v) = \{x\} \cup \{y^{'}\}\).
Use the hypothesis \(v \sim_{G} y\).
Rewriting with \paperlink{premise:53696d706c6547726170682e6d656d5f6e65696768626f72536574}{Prerequisite~P\ref{premise:53696d706c6547726170682e6d656d5f6e65696768626f72536574}} in the preceding result gives \(y \in N_{G}(v)\).
Rewriting with \(N_{G}(v) = \{x\} \cup \{y^{'}\}\), \paperlink{premise:5365742e6d656d5f696e736572745f696666}{Prerequisite~P\ref{premise:5365742e6d656d5f696e736572745f696666}} and \paperlink{premise:5365742e6d656d5f73696e676c65746f6e5f696666}{Prerequisite~P\ref{premise:5365742e6d656d5f73696e676c65746f6e5f696666}} in the preceding result gives \(y = x \lor y = y^{'}\).
Consider the following case.
\begin{proofstep}
Assume \(y = x\).
For the required implication, argue as follows.
\begin{proofstep}
Assume \(\neg \left(y = y^{'}\right)\).
Assume \(N_{G}(v) = \{y\} \cup \{y^{'}\}\).
Assume \(w = y\).
Assume \(w^{'} = y\).
For the required implication, argue as follows.
\begin{proofstep}
Assume \(\neg \left(w = w\right)\).
Assume \(\neg \left(w = y^{'}\right)\).
Assume \(N_{G}(v) = \{w\} \cup \{y^{'}\}\).
Assume \(w^{'} = w\).
For the required implication, argue as follows.
\begin{proofstep}
Assume \(\neg \left(w^{'} = w^{'}\right)\).
Assume \(\neg \left(w^{'} = w^{'}\right)\).
Assume \(\neg \left(w^{'} = y^{'}\right)\).
Assume \(N_{G}(v) = \{w^{'}\} \cup \{y^{'}\}\).
Use the hypothesis \(\neg \left(w^{'} = w^{'}\right)\).
Rewriting with \paperlink{premise:6e6f745f747275655f65715f66616c7365}{Prerequisite~P\ref{premise:6e6f745f747275655f65715f66616c7365}} in the preceding result gives \(\mathrm{False}\).
The contradiction rules out this case. Hence \(w^{'} = w^{'}\).
\end{proofstep}
Use the hypothesis \(w^{'} = w\).
Use the hypothesis \(w \ne w^{'}\).
By the preceding argument, we obtain \(\neg \left(w = w^{'}\right)\).
Use the hypothesis \(\neg \left(w = w\right)\).
Use the hypothesis \(\neg \left(w = y^{'}\right)\).
Use the hypothesis \(N_{G}(v) = \{w\} \cup \{y^{'}\}\).
Substituting equals, we obtain \(w = w^{'}\).
\end{proofstep}
Use the hypothesis \(w = y\).
Use the hypothesis \(w \ne y\).
By the preceding argument, we obtain \(\neg \left(w = y\right)\).
Use the hypothesis \(\neg \left(y = y^{'}\right)\).
Use the hypothesis \(N_{G}(v) = \{y\} \cup \{y^{'}\}\).
Use the hypothesis \(w^{'} = y\).
Substituting equals, we obtain \(y = w^{'}\).
\end{proofstep}
Use the hypothesis \(y = x\).
Use the hypothesis \(\neg \left(x = y^{'}\right)\).
Use the hypothesis \(N_{G}(v) = \{x\} \cup \{y^{'}\}\).
Use the hypothesis \(w = x\).
Use the hypothesis \(w^{'} = x\).
Substituting equals, we obtain \(y = w^{'}\).
\end{proofstep}
Consider the following case.
\begin{proofstep}
Assume \(y = y^{'}\).
For the required implication, argue as follows.
\begin{proofstep}
Assume \(\neg \left(w = y^{'}\right)\).
Assume \(N_{G}(v) = \{w\} \cup \{y^{'}\}\).
Assume \(w^{'} = w\).
For the required implication, argue as follows.
\begin{proofstep}
Assume \(\neg \left(w^{'} = w^{'}\right)\).
Assume \(\neg \left(w^{'} = y\right)\).
Assume \(\neg \left(w^{'} = y^{'}\right)\).
Assume \(N_{G}(v) = \{w^{'}\} \cup \{y^{'}\}\).
For the required implication, argue as follows.
\begin{proofstep}
Assume \(\neg \left(w^{'} = y\right)\).
Assume \(N_{G}(v) = \{w^{'}\} \cup \{y\}\).
Use the hypothesis \(\neg \left(w^{'} = w^{'}\right)\).
Rewriting with \paperlink{premise:6e6f745f747275655f65715f66616c7365}{Prerequisite~P\ref{premise:6e6f745f747275655f65715f66616c7365}} in the preceding result gives \(\mathrm{False}\).
The contradiction rules out this case. Hence \(y = w^{'}\).
\end{proofstep}
Use the hypothesis \(y = y^{'}\).
Use the hypothesis \(\neg \left(w^{'} = y^{'}\right)\).
Use the hypothesis \(N_{G}(v) = \{w^{'}\} \cup \{y^{'}\}\).
Substituting equals, we obtain \(y = w^{'}\).
\end{proofstep}
Use the hypothesis \(w^{'} = w\).
Use the hypothesis \(w \ne w^{'}\).
By the preceding argument, we obtain \(\neg \left(w = w^{'}\right)\).
Use the hypothesis \(w \ne y\).
By the preceding argument, we obtain \(\neg \left(w = y\right)\).
Use the hypothesis \(\neg \left(w = y^{'}\right)\).
Use the hypothesis \(N_{G}(v) = \{w\} \cup \{y^{'}\}\).
Substituting equals, we obtain \(y = w^{'}\).
\end{proofstep}
Use the hypothesis \(w = x\).
Use the hypothesis \(\neg \left(x = y^{'}\right)\).
Use the hypothesis \(N_{G}(v) = \{x\} \cup \{y^{'}\}\).
Use the hypothesis \(w^{'} = x\).
Substituting equals, we obtain \(y = w^{'}\).
\end{proofstep}
The two alternatives exhaust the cases. Hence \(y = w^{'}\).
\end{proofstep}
Consider the following case.
\begin{proofstep}
Assume \(w^{'} = y^{'}\).
Use the hypothesis \(x \ne y^{'} \land N_{G}(v) = \{x\} \cup \{y^{'}\}\).
By the preceding argument, we obtain \(\neg \left(x = y^{'}\right) \land N_{G}(v) = \{x\} \cup \{y^{'}\}\).
By the preceding argument, we obtain \(\neg \left(x = y^{'}\right) \land N_{G}(v) = \{x\} \cup \{y^{'}\}\).
Taking the second conjunct gives \(N_{G}(v) = \{x\} \cup \{y^{'}\}\).
Use the hypothesis \(v \sim_{G} y\).
Rewriting with \paperlink{premise:53696d706c6547726170682e6d656d5f6e65696768626f72536574}{Prerequisite~P\ref{premise:53696d706c6547726170682e6d656d5f6e65696768626f72536574}} in the preceding result gives \(y \in N_{G}(v)\).
Rewriting with \(N_{G}(v) = \{x\} \cup \{y^{'}\}\), \paperlink{premise:5365742e6d656d5f696e736572745f696666}{Prerequisite~P\ref{premise:5365742e6d656d5f696e736572745f696666}} and \paperlink{premise:5365742e6d656d5f73696e676c65746f6e5f696666}{Prerequisite~P\ref{premise:5365742e6d656d5f73696e676c65746f6e5f696666}} in the preceding result gives \(y = x \lor y = y^{'}\).
Consider the following case.
\begin{proofstep}
Assume \(y = x\).
For the required implication, argue as follows.
\begin{proofstep}
Assume \(\neg \left(y = y^{'}\right)\).
Assume \(N_{G}(v) = \{y\} \cup \{y^{'}\}\).
Assume \(w = y\).
For the required implication, argue as follows.
\begin{proofstep}
Assume \(\neg \left(y = w^{'}\right)\).
Assume \(N_{G}(v) = \{y\} \cup \{w^{'}\}\).
For the required implication, argue as follows.
\begin{proofstep}
Assume \(\neg \left(w = w\right)\).
Assume \(\neg \left(w = w^{'}\right)\).
Assume \(N_{G}(v) = \{w\} \cup \{w^{'}\}\).
Use the hypothesis \(\neg \left(w = w\right)\).
Rewriting with \paperlink{premise:6e6f745f747275655f65715f66616c7365}{Prerequisite~P\ref{premise:6e6f745f747275655f65715f66616c7365}} in the preceding result gives \(\mathrm{False}\).
The contradiction rules out this case. Hence \(w = w^{'}\).
\end{proofstep}
Use the hypothesis \(w = y\).
Use the hypothesis \(w \ne y\).
By the preceding argument, we obtain \(\neg \left(w = y\right)\).
Use the hypothesis \(\neg \left(y = w^{'}\right)\).
Use the hypothesis \(N_{G}(v) = \{y\} \cup \{w^{'}\}\).
Substituting equals, we obtain \(y = w^{'}\).
\end{proofstep}
Use the hypothesis \(w^{'} = y^{'}\).
Use the hypothesis \(\neg \left(y = y^{'}\right)\).
Use the hypothesis \(N_{G}(v) = \{y\} \cup \{y^{'}\}\).
Substituting equals, we obtain \(y = w^{'}\).
\end{proofstep}
Use the hypothesis \(y = x\).
Use the hypothesis \(\neg \left(x = y^{'}\right)\).
Use the hypothesis \(N_{G}(v) = \{x\} \cup \{y^{'}\}\).
Use the hypothesis \(w = x\).
Substituting equals, we obtain \(y = w^{'}\).
\end{proofstep}
Consider the following case.
\begin{proofstep}
Assume \(y = y^{'}\).
For the required implication, argue as follows.
\begin{proofstep}
Assume \(\neg \left(x = w^{'}\right)\).
Assume \(N_{G}(v) = \{x\} \cup \{w^{'}\}\).
Assume \(y = w^{'}\).
For the required implication, argue as follows.
\begin{proofstep}
Assume \(\neg \left(w = w^{'}\right)\).
Assume \(N_{G}(v) = \{w\} \cup \{w^{'}\}\).
For the required implication, argue as follows.
\begin{proofstep}
Assume \(\neg \left(w = y\right)\).
Assume \(\neg \left(w = y\right)\).
Assume \(N_{G}(v) = \{w\} \cup \{y\}\).
Simplifying by definition gives \(y = y\).
\end{proofstep}
Use the hypothesis \(y = w^{'}\).
Use the hypothesis \(w \ne w^{'}\).
By the preceding argument, we obtain \(\neg \left(w = w^{'}\right)\).
Use the hypothesis \(\neg \left(w = w^{'}\right)\).
Use the hypothesis \(N_{G}(v) = \{w\} \cup \{w^{'}\}\).
Substituting equals, we obtain \(y = w^{'}\).
\end{proofstep}
Use the hypothesis \(w = x\).
Use the hypothesis \(\neg \left(x = w^{'}\right)\).
Use the hypothesis \(N_{G}(v) = \{x\} \cup \{w^{'}\}\).
Substituting equals, we obtain \(y = w^{'}\).
\end{proofstep}
Use the hypothesis \(w^{'} = y^{'}\).
Use the hypothesis \(\neg \left(x = y^{'}\right)\).
Use the hypothesis \(N_{G}(v) = \{x\} \cup \{y^{'}\}\).
Use the hypothesis \(y = y^{'}\).
Substituting equals, we obtain \(y = w^{'}\).
\end{proofstep}
The two alternatives exhaust the cases. Hence \(y = w^{'}\).
\end{proofstep}
The two alternatives exhaust the cases. Hence \(y = w^{'}\).
\end{proofstep}
Consider the following case.
\begin{proofstep}
Assume \(w = y^{'}\).
Use the hypothesis \(x \ne y^{'} \land N_{G}(v) = \{x\} \cup \{y^{'}\}\).
By the preceding argument, we obtain \(\neg \left(x = y^{'}\right) \land N_{G}(v) = \{x\} \cup \{y^{'}\}\).
By the preceding argument, we obtain \(\neg \left(x = y^{'}\right) \land N_{G}(v) = \{x\} \cup \{y^{'}\}\).
Taking the second conjunct gives \(N_{G}(v) = \{x\} \cup \{y^{'}\}\).
Use the hypothesis \(v \sim_{G} w^{'}\).
Rewriting with \paperlink{premise:53696d706c6547726170682e6d656d5f6e65696768626f72536574}{Prerequisite~P\ref{premise:53696d706c6547726170682e6d656d5f6e65696768626f72536574}} in the preceding result gives \(w^{'} \in N_{G}(v)\).
Rewriting with \(N_{G}(v) = \{x\} \cup \{y^{'}\}\), \paperlink{premise:5365742e6d656d5f696e736572745f696666}{Prerequisite~P\ref{premise:5365742e6d656d5f696e736572745f696666}} and \paperlink{premise:5365742e6d656d5f73696e676c65746f6e5f696666}{Prerequisite~P\ref{premise:5365742e6d656d5f73696e676c65746f6e5f696666}} in the preceding result gives \(w^{'} = x \lor w^{'} = y^{'}\).
Consider the following case.
\begin{proofstep}
Assume \(w^{'} = x\).
Use the hypothesis \(x \ne y^{'} \land N_{G}(v) = \{x\} \cup \{y^{'}\}\).
By the preceding argument, we obtain \(\neg \left(x = y^{'}\right) \land N_{G}(v) = \{x\} \cup \{y^{'}\}\).
By the preceding argument, we obtain \(\neg \left(x = y^{'}\right) \land N_{G}(v) = \{x\} \cup \{y^{'}\}\).
Taking the second conjunct gives \(N_{G}(v) = \{x\} \cup \{y^{'}\}\).
Use the hypothesis \(v \sim_{G} y\).
Rewriting with \paperlink{premise:53696d706c6547726170682e6d656d5f6e65696768626f72536574}{Prerequisite~P\ref{premise:53696d706c6547726170682e6d656d5f6e65696768626f72536574}} in the preceding result gives \(y \in N_{G}(v)\).
Rewriting with \(N_{G}(v) = \{x\} \cup \{y^{'}\}\), \paperlink{premise:5365742e6d656d5f696e736572745f696666}{Prerequisite~P\ref{premise:5365742e6d656d5f696e736572745f696666}} and \paperlink{premise:5365742e6d656d5f73696e676c65746f6e5f696666}{Prerequisite~P\ref{premise:5365742e6d656d5f73696e676c65746f6e5f696666}} in the preceding result gives \(y = x \lor y = y^{'}\).
Consider the following case.
\begin{proofstep}
Assume \(y = x\).
For the required implication, argue as follows.
\begin{proofstep}
Assume \(\neg \left(y = y^{'}\right)\).
Assume \(N_{G}(v) = \{y\} \cup \{y^{'}\}\).
Assume \(w^{'} = y\).
For the required implication, argue as follows.
\begin{proofstep}
Assume \(\neg \left(w = w^{'}\right)\).
Assume \(\neg \left(w^{'} = y^{'}\right)\).
Assume \(N_{G}(v) = \{w^{'}\} \cup \{y^{'}\}\).
For the required implication, argue as follows.
\begin{proofstep}
Assume \(\neg \left(w^{'} = w\right)\).
Assume \(N_{G}(v) = \{w^{'}\} \cup \{w\}\).
Simplifying by definition gives \(w^{'} = w^{'}\).
\end{proofstep}
Use the hypothesis \(w = y^{'}\).
Use the hypothesis \(\neg \left(w^{'} = y^{'}\right)\).
Use the hypothesis \(N_{G}(v) = \{w^{'}\} \cup \{y^{'}\}\).
Substituting equals, we obtain \(w^{'} = w^{'}\).
\end{proofstep}
Use the hypothesis \(w^{'} = y\).
Use the hypothesis \(w \ne y\).
By the preceding argument, we obtain \(\neg \left(w = y\right)\).
Use the hypothesis \(\neg \left(y = y^{'}\right)\).
Use the hypothesis \(N_{G}(v) = \{y\} \cup \{y^{'}\}\).
Substituting equals, we obtain \(y = w^{'}\).
\end{proofstep}
Use the hypothesis \(y = x\).
Use the hypothesis \(\neg \left(x = y^{'}\right)\).
Use the hypothesis \(N_{G}(v) = \{x\} \cup \{y^{'}\}\).
Use the hypothesis \(w^{'} = x\).
Substituting equals, we obtain \(y = w^{'}\).
\end{proofstep}
Consider the following case.
\begin{proofstep}
Assume \(y = y^{'}\).
For the required implication, argue as follows.
\begin{proofstep}
Assume \(\neg \left(w^{'} = y^{'}\right)\).
Assume \(N_{G}(v) = \{w^{'}\} \cup \{y^{'}\}\).
For the required implication, argue as follows.
\begin{proofstep}
Assume \(w = y\).
Assume \(\neg \left(w^{'} = y\right)\).
Assume \(N_{G}(v) = \{w^{'}\} \cup \{y\}\).
For the required implication, argue as follows.
\begin{proofstep}
Assume \(\neg \left(w = w\right)\).
Assume \(\neg \left(w^{'} = w\right)\).
Assume \(N_{G}(v) = \{w^{'}\} \cup \{w\}\).
Use the hypothesis \(\neg \left(w = w\right)\).
Rewriting with \paperlink{premise:6e6f745f747275655f65715f66616c7365}{Prerequisite~P\ref{premise:6e6f745f747275655f65715f66616c7365}} in the preceding result gives \(\mathrm{False}\).
The contradiction rules out this case. Hence \(w = w^{'}\).
\end{proofstep}
Use the hypothesis \(w = y\).
Use the hypothesis \(w \ne y\).
By the preceding argument, we obtain \(\neg \left(w = y\right)\).
Use the hypothesis \(\neg \left(w^{'} = y\right)\).
Use the hypothesis \(N_{G}(v) = \{w^{'}\} \cup \{y\}\).
Substituting equals, we obtain \(y = w^{'}\).
\end{proofstep}
Use the hypothesis \(y = y^{'}\).
Use the hypothesis \(w = y^{'}\).
Use the hypothesis \(\neg \left(w^{'} = y^{'}\right)\).
Use the hypothesis \(N_{G}(v) = \{w^{'}\} \cup \{y^{'}\}\).
Substituting equals, we obtain \(y = w^{'}\).
\end{proofstep}
Use the hypothesis \(w^{'} = x\).
Use the hypothesis \(\neg \left(x = y^{'}\right)\).
Use the hypothesis \(N_{G}(v) = \{x\} \cup \{y^{'}\}\).
Substituting equals, we obtain \(y = w^{'}\).
\end{proofstep}
The two alternatives exhaust the cases. Hence \(y = w^{'}\).
\end{proofstep}
Consider the following case.
\begin{proofstep}
Assume \(w^{'} = y^{'}\).
Use the hypothesis \(x \ne y^{'} \land N_{G}(v) = \{x\} \cup \{y^{'}\}\).
By the preceding argument, we obtain \(\neg \left(x = y^{'}\right) \land N_{G}(v) = \{x\} \cup \{y^{'}\}\).
By the preceding argument, we obtain \(\neg \left(x = y^{'}\right) \land N_{G}(v) = \{x\} \cup \{y^{'}\}\).
Taking the second conjunct gives \(N_{G}(v) = \{x\} \cup \{y^{'}\}\).
Use the hypothesis \(v \sim_{G} y\).
Rewriting with \paperlink{premise:53696d706c6547726170682e6d656d5f6e65696768626f72536574}{Prerequisite~P\ref{premise:53696d706c6547726170682e6d656d5f6e65696768626f72536574}} in the preceding result gives \(y \in N_{G}(v)\).
Rewriting with \(N_{G}(v) = \{x\} \cup \{y^{'}\}\), \paperlink{premise:5365742e6d656d5f696e736572745f696666}{Prerequisite~P\ref{premise:5365742e6d656d5f696e736572745f696666}} and \paperlink{premise:5365742e6d656d5f73696e676c65746f6e5f696666}{Prerequisite~P\ref{premise:5365742e6d656d5f73696e676c65746f6e5f696666}} in the preceding result gives \(y = x \lor y = y^{'}\).
Consider the following case.
\begin{proofstep}
Assume \(y = x\).
For the required implication, argue as follows.
\begin{proofstep}
Assume \(\neg \left(x = w^{'}\right)\).
Assume \(N_{G}(v) = \{x\} \cup \{w^{'}\}\).
Assume \(w = w^{'}\).
For the required implication, argue as follows.
\begin{proofstep}
Assume \(\neg \left(y = w^{'}\right)\).
Assume \(N_{G}(v) = \{y\} \cup \{w^{'}\}\).
For the required implication, argue as follows.
\begin{proofstep}
Assume \(\neg \left(w = w\right)\).
Assume \(\neg \left(y = w\right)\).
Assume \(N_{G}(v) = \{y\} \cup \{w\}\).
Use the hypothesis \(\neg \left(w = w\right)\).
Rewriting with \paperlink{premise:6e6f745f747275655f65715f66616c7365}{Prerequisite~P\ref{premise:6e6f745f747275655f65715f66616c7365}} in the preceding result gives \(\mathrm{False}\).
The contradiction rules out this case. Hence \(y = w\).
\end{proofstep}
Use the hypothesis \(w = w^{'}\).
Use the hypothesis \(w \ne w^{'}\).
By the preceding argument, we obtain \(\neg \left(w = w^{'}\right)\).
Use the hypothesis \(\neg \left(y = w^{'}\right)\).
Use the hypothesis \(N_{G}(v) = \{y\} \cup \{w^{'}\}\).
Substituting equals, we obtain \(y = w^{'}\).
\end{proofstep}
Use the hypothesis \(y = x\).
Use the hypothesis \(\neg \left(x = w^{'}\right)\).
Use the hypothesis \(N_{G}(v) = \{x\} \cup \{w^{'}\}\).
Substituting equals, we obtain \(y = w^{'}\).
\end{proofstep}
Use the hypothesis \(w^{'} = y^{'}\).
Use the hypothesis \(\neg \left(x = y^{'}\right)\).
Use the hypothesis \(N_{G}(v) = \{x\} \cup \{y^{'}\}\).
Use the hypothesis \(w = y^{'}\).
Substituting equals, we obtain \(y = w^{'}\).
\end{proofstep}
Consider the following case.
\begin{proofstep}
Assume \(y = y^{'}\).
For the required implication, argue as follows.
\begin{proofstep}
Assume \(\neg \left(x = y\right)\).
Assume \(N_{G}(v) = \{x\} \cup \{y\}\).
Assume \(w = y\).
Assume \(w^{'} = y\).
For the required implication, argue as follows.
\begin{proofstep}
Assume \(\neg \left(w = w^{'}\right)\).
Assume \(\neg \left(x = w^{'}\right)\).
Assume \(N_{G}(v) = \{x\} \cup \{w^{'}\}\).
Assume \(w = w^{'}\).
For the required implication, argue as follows.
\begin{proofstep}
Assume \(\neg \left(w = w\right)\).
Assume \(\neg \left(w = w\right)\).
Assume \(\neg \left(x = w\right)\).
Assume \(N_{G}(v) = \{x\} \cup \{w\}\).
Use the hypothesis \(\neg \left(w = w\right)\).
Rewriting with \paperlink{premise:6e6f745f747275655f65715f66616c7365}{Prerequisite~P\ref{premise:6e6f745f747275655f65715f66616c7365}} in the preceding result gives \(\mathrm{False}\).
The contradiction rules out this case. Hence \(w = w\).
\end{proofstep}
Use the hypothesis \(w = w^{'}\).
Use the hypothesis \(w \ne w^{'}\).
By the preceding argument, we obtain \(\neg \left(w = w^{'}\right)\).
Use the hypothesis \(\neg \left(w = w^{'}\right)\).
Use the hypothesis \(\neg \left(x = w^{'}\right)\).
Use the hypothesis \(N_{G}(v) = \{x\} \cup \{w^{'}\}\).
Substituting equals, we obtain \(w^{'} = w^{'}\).
\end{proofstep}
Use the hypothesis \(w^{'} = y\).
Use the hypothesis \(w \ne y\).
By the preceding argument, we obtain \(\neg \left(w = y\right)\).
Use the hypothesis \(\neg \left(x = y\right)\).
Use the hypothesis \(N_{G}(v) = \{x\} \cup \{y\}\).
Use the hypothesis \(w = y\).
Substituting equals, we obtain \(y = w^{'}\).
\end{proofstep}
Use the hypothesis \(y = y^{'}\).
Use the hypothesis \(\neg \left(x = y^{'}\right)\).
Use the hypothesis \(N_{G}(v) = \{x\} \cup \{y^{'}\}\).
Use the hypothesis \(w = y^{'}\).
Use the hypothesis \(w^{'} = y^{'}\).
Substituting equals, we obtain \(y = w^{'}\).
\end{proofstep}
The two alternatives exhaust the cases. Hence \(y = w^{'}\).
\end{proofstep}
The two alternatives exhaust the cases. Hence \(y = w^{'}\).
\end{proofstep}
The two alternatives exhaust the cases. Hence \(y = w^{'}\).
\end{proofstep}
Using the two conjuncts, we obtain \(y = w^{'}\).
\end{proofstep}
Using the existential hypothesis, we obtain \(y = w^{'}\).
\end{proofstep}
Using the existential hypothesis, we obtain \(y = w^{'}\).
\end{proofstep}
Using the witnesses and their properties, we obtain \(y = w^{'}\).
\end{proofstep}
Combining the two conclusions, we obtain \((w \ne w^{'} \land v \sim_{G} w^{'}) \land (\forall y\, :\, V,\; w \ne y \land v \sim_{G} y \implies y = w^{'})\).
This witness proves existence. Hence \(\exists x\, :\, V,\; (w \ne x \land v \sim_{G} x) \land (\forall y\, :\, V,\; w \ne y \land v \sim_{G} y \implies y = x)\).
\end{proofstep}
Using the two conjuncts, we obtain \(\exists! w^{'}_{1}\, :\, V,\; w \ne w^{'}_{1} \land v \sim_{G} w^{'}_{1}\).
\end{proofstep}
Using the existential hypothesis, we obtain \(\exists! w^{'}\, :\, V,\; w \ne w^{'} \land v \sim_{G} w^{'}\).
\end{proof}

\begin{theorem}[{Induce supp}]\label{theorem:32}
Let \(V\) be a type. Let \(G\) be a simple graph on \(V\). Let \(c\) be an element of \(\operatorname{Connected\,Component}\left(G\right)\). Assume \(\text{every nonisolated vertex of }G\text{ has degree two}\). Then
\begin{center}
\(\displaystyle \text{every nonisolated vertex of }\operatorname{Simple\,Graph\mathord{:}spanning\,Coe}(\operatorname{Simple\,Graph\mathord{:}induce}(\operatorname{supp}(c),\allowbreak{} G))\text{ has degree two}.\)
\end{center}
\end{theorem}

\begin{proof}
Fix \(v\) in \(V\).
Assume \(N_{\operatorname{Simple\,Graph\mathord{:}spanning\,Coe}(\operatorname{Simple\,Graph\mathord{:}induce}(\operatorname{supp}(c),\allowbreak{} G))}(v)\ne\varnothing\).
Use the hypothesis \(N_{\operatorname{Simple\,Graph\mathord{:}spanning\,Coe}(\operatorname{Simple\,Graph\mathord{:}induce}(\operatorname{supp}(c),\allowbreak{} G))}(v)\ne\varnothing\).
Take the witnesses from the preceding result.
\begin{proofstep}
Fix \(w\) in \(V\).
Assume \(w \in N_{\operatorname{Simple\,Graph\mathord{:}spanning\,Coe}(\operatorname{Simple\,Graph\mathord{:}induce}(\operatorname{supp}(c),\allowbreak{} G))}(v)\).
The substitution uses the following argument.
\begin{proofstep}
Choose the witness \(w\).
Use the hypothesis \(w \in N_{\operatorname{Simple\,Graph\mathord{:}spanning\,Coe}(\operatorname{Simple\,Graph\mathord{:}induce}(\operatorname{supp}(c),\allowbreak{} G))}(v)\).
Rewriting with \paperlink{premise:53696d706c6547726170682e6d656d5f6e65696768626f72536574}{Prerequisite~P\ref{premise:53696d706c6547726170682e6d656d5f6e65696768626f72536574}} in the preceding result gives \(v \sim_{\operatorname{Simple\,Graph\mathord{:}spanning\,Coe}(\operatorname{Simple\,Graph\mathord{:}induce}(\operatorname{supp}(c),\allowbreak{} G))} w\).
Rewriting with \paperlink{premise:53696d706c6547726170682e436f6e6e6563746564436f6d706f6e656e742e61646a5f7370616e6e696e67436f655f696e647563655f73757070}{Prerequisite~P\ref{premise:53696d706c6547726170682e436f6e6e6563746564436f6d706f6e656e742e61646a5f7370616e6e696e67436f655f696e647563655f73757070}} in the preceding result gives \(v \in \operatorname{supp}\left(c\right) \land v \sim_{G} w\).
Taking the second conjunct, we obtain \(v \sim_{G} w\).
This witness proves existence. Hence \(\exists x\, :\, V,\; x \in N_{G}(v)\).
\end{proofstep}
By definition, we obtain \(\operatorname{ncard}(N_{\operatorname{Simple\,Graph\mathord{:}spanning\,Coe}(\operatorname{Simple\,Graph\mathord{:}induce}(\operatorname{supp}(c),\allowbreak{} G))}(v)) = \operatorname{ncard}(N_{\operatorname{Simple\,Graph\mathord{:}spanning\,Coe}(\operatorname{Simple\,Graph\mathord{:}induce}(\operatorname{supp}(c),\allowbreak{} G))}(v))\).
For the required implication, argue as follows.
\begin{proofstep}
Fix \(w^{'}\) in \(V\).
Use the hypothesis \(w \in N_{\operatorname{Simple\,Graph\mathord{:}spanning\,Coe}(\operatorname{Simple\,Graph\mathord{:}induce}(\operatorname{supp}(c),\allowbreak{} G))}(v)\).
Rewriting with \paperlink{premise:53696d706c6547726170682e6d656d5f6e65696768626f72536574}{Prerequisite~P\ref{premise:53696d706c6547726170682e6d656d5f6e65696768626f72536574}} in the preceding result gives \(v \sim_{\operatorname{Simple\,Graph\mathord{:}spanning\,Coe}(\operatorname{Simple\,Graph\mathord{:}induce}(\operatorname{supp}(c),\allowbreak{} G))} w\).
Rewriting with \paperlink{premise:53696d706c6547726170682e436f6e6e6563746564436f6d706f6e656e742e61646a5f7370616e6e696e67436f655f696e647563655f73757070}{Prerequisite~P\ref{premise:53696d706c6547726170682e436f6e6e6563746564436f6d706f6e656e742e61646a5f7370616e6e696e67436f655f696e647563655f73757070}} in the preceding result gives \(v \in \operatorname{supp}\left(c\right) \land v \sim_{G} w\).
Taking the first conjunct gives \(v \in \operatorname{supp}\left(c\right)\).
Rewriting with \paperlink{premise:53696d706c6547726170682e6d656d5f6e65696768626f72536574}{Prerequisite~P\ref{premise:53696d706c6547726170682e6d656d5f6e65696768626f72536574}}, \paperlink{premise:53696d706c6547726170682e436f6e6e6563746564436f6d706f6e656e742e61646a5f7370616e6e696e67436f655f696e647563655f73757070}{Prerequisite~P\ref{premise:53696d706c6547726170682e436f6e6e6563746564436f6d706f6e656e742e61646a5f7370616e6e696e67436f655f696e647563655f73757070}}, \(v \in \operatorname{supp}\left(c\right)\), \paperlink{premise:747275655f616e64}{Prerequisite~P\ref{premise:747275655f616e64}} and \paperlink{premise:6966665f73656c66}{Prerequisite~P\ref{premise:6966665f73656c66}} gives \(w^{'} \in N_{\operatorname{Simple\,Graph\mathord{:}spanning\,Coe}(\operatorname{Simple\,Graph\mathord{:}induce}(\operatorname{supp}(c),\allowbreak{} G))}(v) \iff w^{'} \in N_{G}(v)\).
\end{proofstep}
By \paperlink{premise:5365742e657874}{Prerequisite~P\ref{premise:5365742e657874}}, we obtain \(N_{\operatorname{Simple\,Graph\mathord{:}spanning\,Coe}(\operatorname{Simple\,Graph\mathord{:}induce}(\operatorname{supp}(c),\allowbreak{} G))}(v) = N_{G}(v)\).
Substituting equals, we obtain \(\operatorname{ncard}(N_{\operatorname{Simple\,Graph\mathord{:}spanning\,Coe}(\operatorname{Simple\,Graph\mathord{:}induce}(\operatorname{supp}(c),\allowbreak{} G))}(v)) = \operatorname{ncard}(N_{G}(v))\).
Rewriting with \(\text{every nonisolated vertex of }G\text{ has degree two}\) and the auxiliary result established above in the preceding result gives \(\operatorname{ncard}(N_{\operatorname{Simple\,Graph\mathord{:}spanning\,Coe}(\operatorname{Simple\,Graph\mathord{:}induce}(\operatorname{supp}(c),\allowbreak{} G))}(v)) = 2\).
\end{proofstep}
Using the witnesses and their properties, we obtain \(\operatorname{ncard}(N_{\operatorname{Simple\,Graph\mathord{:}spanning\,Coe}(\operatorname{Simple\,Graph\mathord{:}induce}(\operatorname{supp}(c),\allowbreak{} G))}(v)) = 2\).
\end{proof}

\begin{lemma}[{Is cycles spanning coercion to subgraph}]\label{theorem:33}
Let \(V\) be a type. Let \(G\) be a simple graph on \(V\). Let \(u\) be an element of \(V\). Let \(p\) be a closed walk in \(G\) based at \(u\). Assume \(p\text{ is a cycle}\). Then
\begin{center}
\(\displaystyle \text{every nonisolated vertex of }\operatorname{spanning\,graph}(\operatorname{Simple\,Graph\mathord{:}Walk\mathord{:}to\,Subgraph}(p))\text{ has degree two}.\)
\end{center}
\end{lemma}

\begin{proof}
Fix \(v\) in \(V\).
Assume \(N_{\operatorname{spanning\,graph}(\operatorname{Simple\,Graph\mathord{:}Walk\mathord{:}to\,Subgraph}(p))}(v)\ne\varnothing\).
Use the hypothesis \(p\text{ is a cycle}\).
Use the hypothesis \(N_{\operatorname{spanning\,graph}(\operatorname{Simple\,Graph\mathord{:}Walk\mathord{:}to\,Subgraph}(p))}(v)\ne\varnothing\).
Consider the following case.
\begin{proofstep}
Fix \(w\) in \(V\).
Assume \(w \in N_{\operatorname{spanning\,graph}(\operatorname{Simple\,Graph\mathord{:}Walk\mathord{:}to\,Subgraph}(p))}(v)\).
By \paperlink{premise:53696d706c6547726170682e57616c6b2e6d656d5f76657274735f746f5375626772617068}{Prerequisite~P\ref{premise:53696d706c6547726170682e57616c6b2e6d656d5f76657274735f746f5375626772617068}}, we obtain \(v \in V(\operatorname{Simple\,Graph\mathord{:}Walk\mathord{:}to\,Subgraph}(p)) \iff p\text{ visits }v\).
Use the hypothesis \(w \in N_{\operatorname{spanning\,graph}(\operatorname{Simple\,Graph\mathord{:}Walk\mathord{:}to\,Subgraph}(p))}(v)\).
By \paperlink{premise:53696d706c6547726170682e53756267726170682e656467655f76657274}{Prerequisite~P\ref{premise:53696d706c6547726170682e53756267726170682e656467655f76657274}} with \(\mathit{self} = \operatorname{Simple\,Graph\mathord{:}Walk\mathord{:}to\,Subgraph}\left(p\right)\), we obtain \(v \in V(\operatorname{Simple\,Graph\mathord{:}Walk\mathord{:}to\,Subgraph}\left(p\right))\).
By the forward implication, we obtain \(p\text{ visits }v\).
\end{proofstep}
Using the existential hypothesis, we obtain \(p\text{ visits }v\).
By \paperlink{premise:53696d706c6547726170682e57616c6b2e49734379636c652e6e636172645f6e65696768626f725365745f746f53756267726170685f65715f74776f}{Prerequisite~P\ref{premise:53696d706c6547726170682e57616c6b2e49734379636c652e6e636172645f6e65696768626f725365745f746f53756267726170685f65715f74776f}}, we obtain \(\operatorname{ncard}(\operatorname{Simple\,Graph\mathord{:}Subgraph\mathord{:}neighbor\,Set}(\operatorname{Simple\,Graph\mathord{:}Walk\mathord{:}to\,Subgraph}(p),\allowbreak{} v)) = 2\).
\end{proof}

\begin{theorem}[{Is cycles spanning coercion to subgraph sup edge}]\label{theorem:34}
Let \(V\) be a type. Let \(G\) be a simple graph on \(V\). Let \(u\) and \(v\) be elements of \(V\). Let \(p\) be a walk in \(G\) from \(u\) to \(v\). Assume \(p\text{ has no repeated vertices}\). Assume \(u \ne v\). Assume \(\operatorname{Sym2\mathord{:}mk}\left(\left(v, u\right)\right) \notin \operatorname{edges}\left(p\right)\). Then
\begin{center}
\(\displaystyle \text{every nonisolated vertex of }\operatorname{spanning\,graph}(\operatorname{Simple\,Graph\mathord{:}Walk\mathord{:}to\,Subgraph}(p)) \cup \operatorname{edge}(v,\allowbreak{} u)\text{ has degree two}.\)
\end{center}
\end{theorem}

\begin{proof}
Set \(c\) equal to \(\operatorname{Simple\,Graph\mathord{:}Walk\mathord{:}cons}\left(\operatorname{map\,Le}\left(p\right)\right)\).
Rewriting with \paperlink{premise:53696d706c6547726170682e746f705f61646a}{Prerequisite~P\ref{premise:53696d706c6547726170682e746f705f61646a}}, \paperlink{premise:65715f66616c7365}{Prerequisite~P\ref{premise:65715f66616c7365}}, \paperlink{premise:4e652e73796d6d}{Prerequisite~P\ref{premise:4e652e73796d6d}} (using \(u \ne v\)) and \paperlink{premise:6e6f745f66616c73655f65715f74727565}{Prerequisite~P\ref{premise:6e6f745f66616c73655f65715f74727565}} gives \(v \sim_{\operatorname{complete\,graph}(V)} u\).
For the required implication, argue as follows.
\begin{proofstep}
Fix \(w\) in \(V\).
For the required implication, argue as follows.
\begin{proofstep}
Fix \(x\) in \(V\).
Consider the following case.
\begin{proofstep}
Assume \(w \sim_{\operatorname{Simple\,Graph\mathord{:}Walk\mathord{:}to\,Subgraph}(p)} x \lor (w = v \land x = u \lor w = u \land x = v) \land \neg (w = x)\).
Use the hypothesis \(w \sim_{\operatorname{Simple\,Graph\mathord{:}Walk\mathord{:}to\,Subgraph}(p)} x \lor (w = v \land x = u \lor w = u \land x = v) \land \neg (w = x)\).
Consider the following case.
\begin{proofstep}
Assume \(w \sim_{\operatorname{Simple\,Graph\mathord{:}Walk\mathord{:}to\,Subgraph}\left(p\right)} x\).
Rewriting with \(w \sim_{\operatorname{Simple\,Graph\mathord{:}Walk\mathord{:}to\,Subgraph}\left(p\right)} x\) and \paperlink{premise:6f725f74727565}{Prerequisite~P\ref{premise:6f725f74727565}} gives \((v = w \land u = x \lor v = x \land u = w) \lor w \sim_{\operatorname{Simple\,Graph\mathord{:}Walk\mathord{:}to\,Subgraph}(p)} x\).
\end{proofstep}
Consider the following case.
\begin{proofstep}
Assume \(\left(w = v \land x = u \lor w = u \land x = v\right) \land \neg \left(w = x\right)\).
Use the hypothesis \(\left(w = v \land x = u \lor w = u \land x = v\right) \land \neg \left(w = x\right)\).
Consider the following case.
\begin{proofstep}
Assume \(w = v \land x = u \lor w = u \land x = v\).
Assume \(\neg \left(w = x\right)\).
Use the hypothesis \(w = v \land x = u \lor w = u \land x = v\).
Consider the following case.
\begin{proofstep}
Assume \(w = v \land x = u\).
Use the hypothesis \(w = v \land x = u\).
Taking the first conjunct gives \(w = v\).
Use the hypothesis \(w = v \land x = u\).
Taking the second conjunct gives \(x = u\).
Use the hypothesis \(w = v \land x = u\).
Taking the second conjunct gives \(x = u\).
Use the hypothesis \(w = v \land x = u\).
Taking the first conjunct gives \(w = v\).
Use the hypothesis \(w = v \land x = u\).
Taking the second conjunct gives \(x = u\).
Use the hypothesis \(w = v \land x = u\).
Taking the first conjunct gives \(w = v\).
Use the hypothesis \(w = v \land x = u\).
Taking the first conjunct gives \(w = v\).
Use the hypothesis \(w = v \land x = u\).
Taking the second conjunct gives \(x = u\).
Rewriting with \(w = v\), \(x = u\), \paperlink{premise:616e645f73656c66}{Prerequisite~P\ref{premise:616e645f73656c66}}, \paperlink{premise:65715f66616c7365}{Prerequisite~P\ref{premise:65715f66616c7365}}, \(\neg \left(w = x\right)\), \paperlink{premise:65715f66616c7365}{Prerequisite~P\ref{premise:65715f66616c7365}} (using \(u \ne v\)), \paperlink{premise:6f725f66616c7365}{Prerequisite~P\ref{premise:6f725f66616c7365}} and \paperlink{premise:747275655f6f72}{Prerequisite~P\ref{premise:747275655f6f72}} gives \((v = w \land u = x \lor v = x \land u = w) \lor w \sim_{\operatorname{Simple\,Graph\mathord{:}Walk\mathord{:}to\,Subgraph}(p)} x\).
\end{proofstep}
Consider the following case.
\begin{proofstep}
Assume \(w = u \land x = v\).
Use the hypothesis \(w = u \land x = v\).
Taking the first conjunct gives \(w = u\).
Use the hypothesis \(w = u \land x = v\).
Taking the second conjunct gives \(x = v\).
Use the hypothesis \(w = u \land x = v\).
Taking the second conjunct gives \(x = v\).
Use the hypothesis \(w = u \land x = v\).
Taking the first conjunct gives \(w = u\).
Use the hypothesis \(w = u \land x = v\).
Taking the first conjunct gives \(w = u\).
Use the hypothesis \(w = u \land x = v\).
Taking the second conjunct gives \(x = v\).
Rewriting with \(w = u\), \(x = v\), \paperlink{premise:65715f66616c7365}{Prerequisite~P\ref{premise:65715f66616c7365}} (using \(u \ne v\)), \paperlink{premise:616e645f66616c7365}{Prerequisite~P\ref{premise:616e645f66616c7365}}, \paperlink{premise:616e645f73656c66}{Prerequisite~P\ref{premise:616e645f73656c66}}, \paperlink{premise:6f725f74727565}{Prerequisite~P\ref{premise:6f725f74727565}} and \paperlink{premise:747275655f6f72}{Prerequisite~P\ref{premise:747275655f6f72}} gives \((v = w \land u = x \lor v = x \land u = w) \lor w \sim_{\operatorname{Simple\,Graph\mathord{:}Walk\mathord{:}to\,Subgraph}(p)} x\).
\end{proofstep}
The two alternatives exhaust the cases. Hence \((v = w \land u = x \lor v = x \land u = w) \lor w \sim_{\operatorname{Simple\,Graph\mathord{:}Walk\mathord{:}to\,Subgraph}(p)} x\).
\end{proofstep}
Using the two conjuncts, we obtain \((v = w \land u = x \lor v = x \land u = w) \lor w \sim_{\operatorname{Simple\,Graph\mathord{:}Walk\mathord{:}to\,Subgraph}(p)} x\).
\end{proofstep}
The two alternatives exhaust the cases. Hence \((v = w \land u = x \lor v = x \land u = w) \lor w \sim_{\operatorname{Simple\,Graph\mathord{:}Walk\mathord{:}to\,Subgraph}(p)} x\).
\end{proofstep}
Consider the following case.
\begin{proofstep}
Assume \((v = w \land u = x \lor v = x \land u = w) \lor w \sim_{\operatorname{Simple\,Graph\mathord{:}Walk\mathord{:}to\,Subgraph}(p)} x\).
Use the hypothesis \((v = w \land u = x \lor v = x \land u = w) \lor w \sim_{\operatorname{Simple\,Graph\mathord{:}Walk\mathord{:}to\,Subgraph}(p)} x\).
Consider the following case.
\begin{proofstep}
Assume \(v = w \land u = x \lor v = x \land u = w\).
Use the hypothesis \(v = w \land u = x \lor v = x \land u = w\).
Consider the following case.
\begin{proofstep}
Assume \(v = w \land u = x\).
Use the hypothesis \(v = w \land u = x\).
Taking the first conjunct gives \(v = w\).
Use the hypothesis \(v = w \land u = x\).
Taking the second conjunct gives \(u = x\).
Use the hypothesis \(v = w \land u = x\).
Taking the second conjunct gives \(u = x\).
Use the hypothesis \(v = w \land u = x\).
Taking the first conjunct gives \(v = w\).
Use the hypothesis \(v = w \land u = x\).
Consider the following case.
\begin{proofstep}
Assume \(v = w\).
Assume \(u = x\).
For the required implication, argue as follows.
\begin{proofstep}
Assume \(\operatorname{Sym2\mathord{:}mk}\left(\left(w, u\right)\right) \notin \operatorname{edges}\left(p\right)\).
For the required implication, argue as follows.
\begin{proofstep}
Assume \(\operatorname{Sym2\mathord{:}mk}\left(\left(v, u\right)\right) \notin \operatorname{edges}\left(p\right)\).
For the required implication, argue as follows.
\begin{proofstep}
Assume \(v = u\).
For the required implication, argue as follows.
\begin{proofstep}
Fix \(p_{1}\) in \(\operatorname{Walk}_{G}(v, v)\).
Assume \(p_{1}\text{ has no repeated vertices}\).
Assume \(v \ne v\).
Set \(c_{1}\) equal to \(\operatorname{Simple\,Graph\mathord{:}Walk\mathord{:}cons}(\operatorname{map\,Le}(p_{1}))\).
Rewriting with \paperlink{premise:53696d706c6547726170682e746f705f61646a}{Prerequisite~P\ref{premise:53696d706c6547726170682e746f705f61646a}}, \paperlink{premise:65715f66616c7365}{Prerequisite~P\ref{premise:65715f66616c7365}}, \paperlink{premise:4e652e73796d6d}{Prerequisite~P\ref{premise:4e652e73796d6d}} (using \(v \ne v\)) and \paperlink{premise:6e6f745f66616c73655f65715f74727565}{Prerequisite~P\ref{premise:6e6f745f66616c73655f65715f74727565}} gives \(v \sim_{\operatorname{complete\,graph}(V)} v\).
Assume \(\operatorname{Sym2\mathord{:}mk}\left(\left(v, v\right)\right) \notin \operatorname{edges}\left(p_{1}\right)\).
Set \(c_{2}\) equal to \(\operatorname{Simple\,Graph\mathord{:}Walk\mathord{:}cons}(\operatorname{map\,Le}(\operatorname{Simple\,Graph\mathord{:}Walk\mathord{:}nil}))\).
Rewriting with \paperlink{premise:53696d706c6547726170682e746f705f61646a}{Prerequisite~P\ref{premise:53696d706c6547726170682e746f705f61646a}}, \paperlink{premise:65715f66616c7365}{Prerequisite~P\ref{premise:65715f66616c7365}}, \paperlink{premise:4e652e73796d6d}{Prerequisite~P\ref{premise:4e652e73796d6d}} (using \(v \ne v\)) and \paperlink{premise:6e6f745f66616c73655f65715f74727565}{Prerequisite~P\ref{premise:6e6f745f66616c73655f65715f74727565}} gives \(v \sim_{\operatorname{complete\,graph}(V)} v\).
Simplifying by definition gives \(v = v\).
Applying the hypothesis \(v \ne v\), we obtain \(\mathrm{False}\).
Use the hypothesis \(p_{1}\text{ has no repeated vertices}\).
Rewriting with \paperlink{premise:53696d706c6547726170682e57616c6b2e6973506174685f6966665f65715f6e696c}{Prerequisite~P\ref{premise:53696d706c6547726170682e57616c6b2e6973506174685f6966665f65715f6e696c}} in the preceding result gives \(p_{1} = \operatorname{Simple\,Graph\mathord{:}Walk\mathord{:}nil}\).
Reversing the equality, we obtain \(\operatorname{Simple\,Graph\mathord{:}Walk\mathord{:}nil} = p_{1}\).
Substituting equals, we obtain \(\mathrm{False}\).
\end{proofstep}
Use the hypothesis \(v = u\).
Use the hypothesis \(p\text{ has no repeated vertices}\).
Use the hypothesis \(u \ne v\).
Use the hypothesis \(\operatorname{Sym2\mathord{:}mk}\left(\left(v, u\right)\right) \notin \operatorname{edges}\left(p\right)\).
Substituting equals, we obtain \(\mathrm{False}\).
\end{proofstep}
By \paperlink{premise:4165736f702e4275696c74696e52756c65732e6e6f745f696e74726f}{Prerequisite~P\ref{premise:4165736f702e4275696c74696e52756c65732e6e6f745f696e74726f}}, we obtain \(\neg \left(v = u\right)\).
The second alternative suffices. Hence \(v \sim_{\operatorname{Simple\,Graph\mathord{:}Walk\mathord{:}to\,Subgraph}(p)} u \lor \neg (v = u)\).
\end{proofstep}
Use the hypothesis \(v = w\).
Use the hypothesis \(\operatorname{Sym2\mathord{:}mk}\left(\left(w, u\right)\right) \notin \operatorname{edges}\left(p\right)\).
Substituting equals, we obtain \(w \sim_{\operatorname{Simple\,Graph\mathord{:}Walk\mathord{:}to\,Subgraph}(p)} u \lor \neg (w = u)\).
\end{proofstep}
Use the hypothesis \(u = x\).
Use the hypothesis \(v = w \land u = x\).
Taking the first conjunct gives \(v = w\).
Use the hypothesis \(v = w \land u = x\).
Taking the second conjunct gives \(u = x\).
Use the hypothesis \(\operatorname{Sym2\mathord{:}mk}\left(\left(v, u\right)\right) \notin \operatorname{edges}\left(p\right)\).
Rewriting with \(v = w\) and \(u = x\) in the preceding result gives \(\operatorname{Sym2\mathord{:}mk}\left(\left(w, x\right)\right) \notin \operatorname{edges}\left(p\right)\).
Substituting equals, we obtain \(w \sim_{\operatorname{Simple\,Graph\mathord{:}Walk\mathord{:}to\,Subgraph}(p)} x \lor \neg (w = x)\).
\end{proofstep}
Using the two conjuncts, we obtain \(w \sim_{\operatorname{Simple\,Graph\mathord{:}Walk\mathord{:}to\,Subgraph}(p)} x \lor \neg (w = x)\).
Rewriting with \(v = w\), \(u = x\), \paperlink{premise:616e645f73656c66}{Prerequisite~P\ref{premise:616e645f73656c66}}, \paperlink{premise:747275655f6f72}{Prerequisite~P\ref{premise:747275655f6f72}} and \paperlink{premise:747275655f616e64}{Prerequisite~P\ref{premise:747275655f616e64}} in the preceding result gives \(w \sim_{\operatorname{Simple\,Graph\mathord{:}Walk\mathord{:}to\,Subgraph}(p)} x \lor (w = v \land x = u \lor w = u \land x = v) \land \neg (w = x)\).
\end{proofstep}
Consider the following case.
\begin{proofstep}
Assume \(v = x \land u = w\).
Use the hypothesis \(v = x \land u = w\).
Taking the first conjunct gives \(v = x\).
Use the hypothesis \(v = x \land u = w\).
Taking the second conjunct gives \(u = w\).
Use the hypothesis \(v = x \land u = w\).
Taking the second conjunct gives \(u = w\).
Use the hypothesis \(v = x \land u = w\).
Taking the first conjunct gives \(v = x\).
Use the hypothesis \(v = x \land u = w\).
Consider the following case.
\begin{proofstep}
Assume \(v = x\).
Assume \(u = w\).
For the required implication, argue as follows.
\begin{proofstep}
Assume \(\operatorname{Sym2\mathord{:}mk}\left(\left(v, w\right)\right) \notin \operatorname{edges}\left(p\right)\).
For the required implication, argue as follows.
\begin{proofstep}
Assume \(\operatorname{Sym2\mathord{:}mk}\left(\left(v, u\right)\right) \notin \operatorname{edges}\left(p\right)\).
Rewriting with \paperlink{premise:65715f66616c7365}{Prerequisite~P\ref{premise:65715f66616c7365}} (using \(u \ne v\)), \paperlink{premise:6e6f745f66616c73655f65715f74727565}{Prerequisite~P\ref{premise:6e6f745f66616c73655f65715f74727565}} and \paperlink{premise:6f725f74727565}{Prerequisite~P\ref{premise:6f725f74727565}} gives \(u \sim_{\operatorname{Simple\,Graph\mathord{:}Walk\mathord{:}to\,Subgraph}(p)} v \lor \neg (u = v)\).
\end{proofstep}
Use the hypothesis \(u = w\).
Use the hypothesis \(\operatorname{Sym2\mathord{:}mk}\left(\left(v, w\right)\right) \notin \operatorname{edges}\left(p\right)\).
Substituting equals, we obtain \(w \sim_{\operatorname{Simple\,Graph\mathord{:}Walk\mathord{:}to\,Subgraph}(p)} v \lor \neg (w = v)\).
\end{proofstep}
Use the hypothesis \(v = x\).
Use the hypothesis \(v = x \land u = w\).
Taking the first conjunct gives \(v = x\).
Use the hypothesis \(v = x \land u = w\).
Taking the second conjunct gives \(u = w\).
Use the hypothesis \(\operatorname{Sym2\mathord{:}mk}\left(\left(v, u\right)\right) \notin \operatorname{edges}\left(p\right)\).
Rewriting with \(v = x\) and \(u = w\) in the preceding result gives \(\operatorname{Sym2\mathord{:}mk}\left(\left(x, w\right)\right) \notin \operatorname{edges}\left(p\right)\).
Substituting equals, we obtain \(w \sim_{\operatorname{Simple\,Graph\mathord{:}Walk\mathord{:}to\,Subgraph}(p)} x \lor \neg (w = x)\).
\end{proofstep}
Using the two conjuncts, we obtain \(w \sim_{\operatorname{Simple\,Graph\mathord{:}Walk\mathord{:}to\,Subgraph}(p)} x \lor \neg (w = x)\).
Rewriting with \(v = x\), \(u = w\), \paperlink{premise:616e645f73656c66}{Prerequisite~P\ref{premise:616e645f73656c66}}, \paperlink{premise:6f725f74727565}{Prerequisite~P\ref{premise:6f725f74727565}} and \paperlink{premise:747275655f616e64}{Prerequisite~P\ref{premise:747275655f616e64}} in the preceding result gives \(w \sim_{\operatorname{Simple\,Graph\mathord{:}Walk\mathord{:}to\,Subgraph}(p)} x \lor (w = v \land x = u \lor w = u \land x = v) \land \neg (w = x)\).
\end{proofstep}
The two alternatives exhaust the cases. Hence \(w \sim_{\operatorname{Simple\,Graph\mathord{:}Walk\mathord{:}to\,Subgraph}(p)} x \lor (w = v \land x = u \lor w = u \land x = v) \land \neg (w = x)\).
\end{proofstep}
Consider the following case.
\begin{proofstep}
Assume \(w \sim_{\operatorname{Simple\,Graph\mathord{:}Walk\mathord{:}to\,Subgraph}\left(p\right)} x\).
Rewriting with \(w \sim_{\operatorname{Simple\,Graph\mathord{:}Walk\mathord{:}to\,Subgraph}\left(p\right)} x\) and \paperlink{premise:747275655f6f72}{Prerequisite~P\ref{premise:747275655f6f72}} gives \(w \sim_{\operatorname{Simple\,Graph\mathord{:}Walk\mathord{:}to\,Subgraph}(p)} x \lor (w = v \land x = u \lor w = u \land x = v) \land \neg (w = x)\).
\end{proofstep}
The two alternatives exhaust the cases. Hence \(w \sim_{\operatorname{Simple\,Graph\mathord{:}Walk\mathord{:}to\,Subgraph}(p)} x \lor (w = v \land x = u \lor w = u \land x = v) \land \neg (w = x)\).
\end{proofstep}
Together the two implications give the equivalence. Hence \(w \sim_{\operatorname{Simple\,Graph\mathord{:}Walk\mathord{:}to\,Subgraph}(p)} x \lor (w = v \land x = u \lor w = u \land x = v) \land \neg (w = x) \iff (v = w \land u = x \lor v = x \land u = w) \lor w \sim_{\operatorname{Simple\,Graph\mathord{:}Walk\mathord{:}to\,Subgraph}(p)} x\).
Rewriting with \paperlink{premise:53796d322e6571}{Prerequisite~P\ref{premise:53796d322e6571}}, \paperlink{premise:53796d322e72656c5f69666627}{Prerequisite~P\ref{premise:53796d322e72656c5f69666627}} and \paperlink{premise:50726f642e6d6b2e696e6a4571}{Prerequisite~P\ref{premise:50726f642e6d6b2e696e6a4571}} in the preceding result gives \(w \sim_{\operatorname{Simple\,Graph\mathord{:}Walk\mathord{:}to\,Subgraph}(p)} x \lor (w = v \land x = u \lor w = u \land x = v) \land w \ne x \iff \operatorname{Sym2\mathord{:}mk}((v,\allowbreak{} u)) = \operatorname{Sym2\mathord{:}mk}((w,\allowbreak{} x)) \lor w \sim_{\operatorname{Simple\,Graph\mathord{:}Walk\mathord{:}to\,Subgraph}(p)} x\).
Rewriting with \paperlink{premise:53696d706c6547726170682e7375705f61646a}{Prerequisite~P\ref{premise:53696d706c6547726170682e7375705f61646a}}, \paperlink{premise:53696d706c6547726170682e656467655f61646a}{Prerequisite~P\ref{premise:53696d706c6547726170682e656467655f61646a}}, \paperlink{premise:53696d706c6547726170682e53756267726170682e7375705f61646a}{Prerequisite~P\ref{premise:53696d706c6547726170682e53756267726170682e7375705f61646a}}, \paperlink{premise:53696d706c6547726170682e57616c6b2e61646a5f746f53756267726170685f6d61704c65}{Prerequisite~P\ref{premise:53696d706c6547726170682e57616c6b2e61646a5f746f53756267726170685f6d61704c65}} and \paperlink{premise:4f72646572546f702e6c655f746f70}{Prerequisite~P\ref{premise:4f72646572546f702e6c655f746f70}} in the preceding result gives \(w \sim_{\operatorname{spanning\,graph}(\operatorname{Simple\,Graph\mathord{:}Walk\mathord{:}to\,Subgraph}(p)) \cup \operatorname{edge}(v,\allowbreak{} u)} x \iff w \sim_{\operatorname{spanning\,graph}(\operatorname{Simple\,Graph\mathord{:}Walk\mathord{:}to\,Subgraph}(c))} x\).
By propositional extensionality, we obtain \((w \sim_{\operatorname{spanning\,graph}(\operatorname{Simple\,Graph\mathord{:}Walk\mathord{:}to\,Subgraph}(p)) \cup \operatorname{edge}(v,\allowbreak{} u)} x) = (w \sim_{\operatorname{spanning\,graph}(\operatorname{Simple\,Graph\mathord{:}Walk\mathord{:}to\,Subgraph}(c))} x)\).
\end{proofstep}
By function extensionality, we obtain \((v_{1} \mapsto w \sim_{\operatorname{spanning\,graph}(\operatorname{Simple\,Graph\mathord{:}Walk\mathord{:}to\,Subgraph}(p)) \cup \operatorname{edge}(v,\allowbreak{} u)} v_{1}) = (v_{1} \mapsto w \sim_{\operatorname{spanning\,graph}(\operatorname{Simple\,Graph\mathord{:}Walk\mathord{:}to\,Subgraph}(c))} v_{1})\).
\end{proofstep}
By function extensionality, we obtain \(\operatorname{Adj}(\operatorname{spanning\,graph}(\operatorname{Simple\,Graph\mathord{:}Walk\mathord{:}to\,Subgraph}(p)) \cup \operatorname{edge}(v,\allowbreak{} u)) = \operatorname{Adj}(\operatorname{spanning\,graph}(\operatorname{Simple\,Graph\mathord{:}Walk\mathord{:}to\,Subgraph}(c)))\).
By \paperlink{premise:53696d706c6547726170682e657874}{Prerequisite~P\ref{premise:53696d706c6547726170682e657874}}, we obtain \(\operatorname{spanning\,graph}(\operatorname{Simple\,Graph\mathord{:}Walk\mathord{:}to\,Subgraph}(p)) \cup \operatorname{edge}(v,\allowbreak{} u) = \operatorname{spanning\,graph}(\operatorname{Simple\,Graph\mathord{:}Walk\mathord{:}to\,Subgraph}(c))\).
Rewriting with \paperlink{premise:53696d706c6547726170682e57616c6b2e636f6e735f69734379636c655f696666}{Prerequisite~P\ref{premise:53696d706c6547726170682e57616c6b2e636f6e735f69734379636c655f696666}}, \paperlink{premise:53696d706c6547726170682e746f705f61646a}{Prerequisite~P\ref{premise:53696d706c6547726170682e746f705f61646a}}, \paperlink{premise:65715f66616c7365}{Prerequisite~P\ref{premise:65715f66616c7365}}, \paperlink{premise:4e652e73796d6d}{Prerequisite~P\ref{premise:4e652e73796d6d}} (using \(u \ne v\)), \paperlink{premise:6e6f745f66616c73655f65715f74727565}{Prerequisite~P\ref{premise:6e6f745f66616c73655f65715f74727565}}, \paperlink{premise:53696d706c6547726170682e57616c6b2e6d61704c655f697350617468}{Prerequisite~P\ref{premise:53696d706c6547726170682e57616c6b2e6d61704c655f697350617468}}, \paperlink{premise:4f72646572546f702e6c655f746f70}{Prerequisite~P\ref{premise:4f72646572546f702e6c655f746f70}}, \(p\text{ has no repeated vertices}\), \paperlink{premise:53696d706c6547726170682e57616c6b2e65646765735f6d6170}{Prerequisite~P\ref{premise:53696d706c6547726170682e57616c6b2e65646765735f6d6170}}, \paperlink{premise:53796d322e6d61705f6964}{Prerequisite~P\ref{premise:53796d322e6d61705f6964}}, \paperlink{premise:4c6973742e6d61705f69645f66756e}{Prerequisite~P\ref{premise:4c6973742e6d61705f69645f66756e}}, \paperlink{premise:65715f66616c7365}{Prerequisite~P\ref{premise:65715f66616c7365}} (using \(\operatorname{Sym2\mathord{:}mk}\left(\left(v, u\right)\right) \notin \operatorname{edges}\left(p\right)\)) and \paperlink{premise:616e645f73656c66}{Prerequisite~P\ref{premise:616e645f73656c66}} gives \(\operatorname{Simple\,Graph\mathord{:}Walk\mathord{:}cons}(\operatorname{map\,Le}(p))\text{ is a cycle}\).
By \paperlink{theorem:33}{Lemma~\ref{theorem:33}}, we obtain \(\text{every nonisolated vertex of }\operatorname{spanning\,graph}(\operatorname{Simple\,Graph\mathord{:}Walk\mathord{:}to\,Subgraph}(c))\text{ has degree two}\).
Use the previously proved fact \(\operatorname{spanning\,graph}(\operatorname{Simple\,Graph\mathord{:}Walk\mathord{:}to\,Subgraph}(p)) \cup \operatorname{edge}(v,\allowbreak{} u) = \operatorname{spanning\,graph}(\operatorname{Simple\,Graph\mathord{:}Walk\mathord{:}to\,Subgraph}(c))\).
Reversing the equality, we obtain \(\operatorname{spanning\,graph}(\operatorname{Simple\,Graph\mathord{:}Walk\mathord{:}to\,Subgraph}(c)) = \operatorname{spanning\,graph}(\operatorname{Simple\,Graph\mathord{:}Walk\mathord{:}to\,Subgraph}(p)) \cup \operatorname{edge}(v,\allowbreak{} u)\).
Substituting equals, we obtain \(\text{every nonisolated vertex of }\operatorname{spanning\,graph}(\operatorname{Simple\,Graph\mathord{:}Walk\mathord{:}to\,Subgraph}(p)) \cup \operatorname{edge}(v,\allowbreak{} u)\text{ has degree two}\).
\end{proof}

\begin{theorem}[{Adj to subgraph if and only if}]\label{theorem:35}
Let \(V\) be a type. Let \(G\) be a simple graph on \(V\). Let \(v\) be an element of \(V\). Fix a structure \(\mathit{inst}\) of type \(\operatorname{Locally\,Finite}\left(G\right)\). Let \(u\) be an element of \(V\). Let \(p\) be a closed walk in \(G\) based at \(u\). Assume \(p\text{ is a cycle}\). Assume \(\text{every nonisolated vertex of }G\text{ has degree two}\). Assume \(v \in V(\operatorname{Simple\,Graph\mathord{:}Walk\mathord{:}to\,Subgraph}\left(p\right))\). Let \(w\) be an element of \(V\). Then
\begin{center}
\(\displaystyle v \sim_{\operatorname{Simple\,Graph\mathord{:}Walk\mathord{:}to\,Subgraph}(p)} w \iff v \sim_{G} w.\)
\end{center}
\end{theorem}

\begin{proof}
By \paperlink{premise:46696e6974652e6f665f66696e74797065}{Prerequisite~P\ref{premise:46696e6974652e6f665f66696e74797065}} with \(\alpha = N_{G}(v)\), we obtain \(\operatorname{Finite}\left(N_{G}(v)\right)\).
By \paperlink{premise:5365742e746f46696e697465}{Prerequisite~P\ref{premise:5365742e746f46696e697465}} with \(s = N_{G}(v)\), we obtain \(\operatorname{Set\mathord{:}Finite}\left(N_{G}(v)\right)\).
By \paperlink{premise:53696d706c6547726170682e53756267726170682e61646a5f6966665f6f665f6e65696768626f725365745f6571756976}{Prerequisite~P\ref{premise:53696d706c6547726170682e53756267726170682e61646a5f6966665f6f665f6e65696768626f725365745f6571756976}} with \(h = \operatorname{default}\), we obtain \(v \sim_{\operatorname{Simple\,Graph\mathord{:}Walk\mathord{:}to\,Subgraph}\left(p\right)} w \iff v \sim_{G} w\).
\end{proof}

\begin{theorem}[{Symm diff is cycles}]\label{theorem:36}
Let \(V\) be a type. Let \(G\) and \(G^{'}\) be simple graphs on \(V\). Let \(M\) be a subgraph of \(G\). Let \(M^{'}\) be a subgraph of \(G^{'}\). Assume \(M\text{ is a perfect matching}\). Assume \(M^{'}\text{ is a perfect matching}\). Then
\begin{center}
\(\displaystyle \text{every nonisolated vertex of }\operatorname{spanning\,graph}(M) \mathbin{\triangle} \operatorname{spanning\,graph}(M^{'})\text{ has degree two}.\)
\end{center}
\end{theorem}

\begin{proof}
Fix \(v\) in \(V\).
Use the hypothesis \(M\text{ is a perfect matching}\).
Use the hypothesis \(M\text{ is a perfect matching}\).
Taking the second conjunct, we obtain \(v \in V(M)\).
Taking the first conjunct, we obtain \(\exists! w\, :\, V,\; v \sim_{M} w\).
Consider the following case.
\begin{proofstep}
Fix \(w\) in \(V\).
Assume \(v \sim_{M} w \land \left(\forall y\, :\, V,\; v \sim_{M} y \implies y = w\right)\).
Use the hypothesis \(M^{'}\text{ is a perfect matching}\).
Use the hypothesis \(M^{'}\text{ is a perfect matching}\).
Taking the second conjunct, we obtain \(v \in V(M^{'})\).
Taking the first conjunct, we obtain \(\exists! w_{1}\, :\, V,\; v \sim_{M^{'}} w_{1}\).
Consider the following case.
\begin{proofstep}
Fix \(w^{'}\) in \(V\).
Assume \(v \sim_{M^{'}} w^{'} \land \left(\forall y\, :\, V,\; v \sim_{M^{'}} y \implies y = w^{'}\right)\).
The substitution uses the following argument.
\begin{proofstep}
Fix \(x\) in \(V\).
By \paperlink{premise:53696d706c6547726170682e6d656d5f6e65696768626f72536574}{Prerequisite~P\ref{premise:53696d706c6547726170682e6d656d5f6e65696768626f72536574}} with \((\operatorname{spanning\,graph}(M) \setminus \operatorname{spanning\,graph}(M^{'})) \cup (\operatorname{spanning\,graph}(M^{'}) \setminus \operatorname{spanning\,graph}(M))\) in place of its parameter \(G\) and \(x\) in place of its parameter \(w\), we obtain \(x \in N_{(\operatorname{spanning\,graph}(M) \setminus \operatorname{spanning\,graph}(M^{'})) \cup (\operatorname{spanning\,graph}(M^{'}) \setminus \operatorname{spanning\,graph}(M))}(v) \iff v \sim_{(\operatorname{spanning\,graph}(M) \setminus \operatorname{spanning\,graph}(M^{'})) \cup (\operatorname{spanning\,graph}(M^{'}) \setminus \operatorname{spanning\,graph}(M))} x\).
By propositional extensionality, we obtain \((x \in N_{(\operatorname{spanning\,graph}(M) \setminus \operatorname{spanning\,graph}(M^{'})) \cup (\operatorname{spanning\,graph}(M^{'}) \setminus \operatorname{spanning\,graph}(M))}(v)) = (v \sim_{(\operatorname{spanning\,graph}(M) \setminus \operatorname{spanning\,graph}(M^{'})) \cup (\operatorname{spanning\,graph}(M^{'}) \setminus \operatorname{spanning\,graph}(M))} x)\).
By \paperlink{premise:53696d706c6547726170682e7375705f61646a}{Prerequisite~P\ref{premise:53696d706c6547726170682e7375705f61646a}} with \(\operatorname{spanning\,graph}(M) \setminus \operatorname{spanning\,graph}(M^{'})\) in place of its parameter \(x\) and \(y = \operatorname{spanning\,graph}(M^{'}) \setminus \operatorname{spanning\,graph}(M)\) and \(x\) in place of its parameter \(w\), we obtain \(v \sim_{(\operatorname{spanning\,graph}(M) \setminus \operatorname{spanning\,graph}(M^{'})) \cup (\operatorname{spanning\,graph}(M^{'}) \setminus \operatorname{spanning\,graph}(M))} x \iff v \sim_{\operatorname{spanning\,graph}(M) \setminus \operatorname{spanning\,graph}(M^{'})} x \lor v \sim_{\operatorname{spanning\,graph}(M^{'}) \setminus \operatorname{spanning\,graph}(M)} x\).
By propositional extensionality, we obtain \((v \sim_{(\operatorname{spanning\,graph}(M) \setminus \operatorname{spanning\,graph}(M^{'})) \cup (\operatorname{spanning\,graph}(M^{'}) \setminus \operatorname{spanning\,graph}(M))} x) = (v \sim_{\operatorname{spanning\,graph}(M) \setminus \operatorname{spanning\,graph}(M^{'})} x \lor v \sim_{\operatorname{spanning\,graph}(M^{'}) \setminus \operatorname{spanning\,graph}(M)} x)\).
By \paperlink{premise:53696d706c6547726170682e73646966665f61646a}{Prerequisite~P\ref{premise:53696d706c6547726170682e73646966665f61646a}} with \(\operatorname{spanning\,graph}(M)\) in place of its parameter \(x\) and \(y = \operatorname{spanning\,graph}(M^{'})\) and \(x\) in place of its parameter \(w\), we obtain \(v \sim_{\operatorname{spanning\,graph}(M) \setminus \operatorname{spanning\,graph}(M^{'})} x \iff v \sim_{\operatorname{spanning\,graph}(M)} x \land \neg (v \sim_{\operatorname{spanning\,graph}(M^{'})} x)\).
By propositional extensionality, we obtain \((v \sim_{\operatorname{spanning\,graph}(M) \setminus \operatorname{spanning\,graph}(M^{'})} x) = (v \sim_{\operatorname{spanning\,graph}(M)} x \land \neg (v \sim_{\operatorname{spanning\,graph}(M^{'})} x))\).
Applying the same function to equal arguments, we obtain \((y \mapsto v \sim_{\operatorname{spanning\,graph}(M) \setminus \operatorname{spanning\,graph}(M^{'})} x \lor y) = (y \mapsto v \sim_{\operatorname{spanning\,graph}(M)} x \land \neg (v \sim_{\operatorname{spanning\,graph}(M^{'})} x) \lor y)\).
By \paperlink{premise:53696d706c6547726170682e73646966665f61646a}{Prerequisite~P\ref{premise:53696d706c6547726170682e73646966665f61646a}} with \(\operatorname{spanning\,graph}(M^{'})\) in place of its parameter \(x\) and \(y = \operatorname{spanning\,graph}(M)\) and \(x\) in place of its parameter \(w\), we obtain \(v \sim_{\operatorname{spanning\,graph}(M^{'}) \setminus \operatorname{spanning\,graph}(M)} x \iff v \sim_{\operatorname{spanning\,graph}(M^{'})} x \land \neg (v \sim_{\operatorname{spanning\,graph}(M)} x)\).
By propositional extensionality, we obtain \((v \sim_{\operatorname{spanning\,graph}(M^{'}) \setminus \operatorname{spanning\,graph}(M)} x) = (v \sim_{\operatorname{spanning\,graph}(M^{'})} x \land \neg (v \sim_{\operatorname{spanning\,graph}(M)} x))\).
By congruence, we obtain \((v \sim_{\operatorname{spanning\,graph}(M) \setminus \operatorname{spanning\,graph}(M^{'})} x \lor v \sim_{\operatorname{spanning\,graph}(M^{'}) \setminus \operatorname{spanning\,graph}(M)} x) = (v \sim_{\operatorname{spanning\,graph}(M)} x \land \neg (v \sim_{\operatorname{spanning\,graph}(M^{'})} x) \lor v \sim_{\operatorname{spanning\,graph}(M^{'})} x \land \neg (v \sim_{\operatorname{spanning\,graph}(M)} x))\).
Chaining the equalities, we obtain \((v \sim_{(\operatorname{spanning\,graph}(M) \setminus \operatorname{spanning\,graph}(M^{'})) \cup (\operatorname{spanning\,graph}(M^{'}) \setminus \operatorname{spanning\,graph}(M))} x) = (v \sim_{\operatorname{spanning\,graph}(M)} x \land \neg (v \sim_{\operatorname{spanning\,graph}(M^{'})} x) \lor v \sim_{\operatorname{spanning\,graph}(M^{'})} x \land \neg (v \sim_{\operatorname{spanning\,graph}(M)} x))\).
Chaining the equalities, we obtain \((x \in N_{(\operatorname{spanning\,graph}(M) \setminus \operatorname{spanning\,graph}(M^{'})) \cup (\operatorname{spanning\,graph}(M^{'}) \setminus \operatorname{spanning\,graph}(M))}(v)) = (v \sim_{\operatorname{spanning\,graph}(M)} x \land \neg (v \sim_{\operatorname{spanning\,graph}(M^{'})} x) \lor v \sim_{\operatorname{spanning\,graph}(M^{'})} x \land \neg (v \sim_{\operatorname{spanning\,graph}(M)} x))\).
Applying the same function to equal arguments, we obtain \((x \notin N_{(\operatorname{spanning\,graph}(M) \setminus \operatorname{spanning\,graph}(M^{'})) \cup (\operatorname{spanning\,graph}(M^{'}) \setminus \operatorname{spanning\,graph}(M))}(v)) = (\neg (v \sim_{\operatorname{spanning\,graph}(M)} x \land \neg (v \sim_{\operatorname{spanning\,graph}(M^{'})} x) \lor v \sim_{\operatorname{spanning\,graph}(M^{'})} x \land \neg (v \sim_{\operatorname{spanning\,graph}(M)} x)))\).
By \paperlink{premise:6e6f745f6f72}{Prerequisite~P\ref{premise:6e6f745f6f72}}, we obtain \(\neg (v \sim_{M} x \land \neg (v \sim_{M^{'}} x) \lor v \sim_{M^{'}} x \land \neg (v \sim_{M} x)) \iff \neg (v \sim_{M} x \land \neg (v \sim_{M^{'}} x)) \land \neg (v \sim_{M^{'}} x \land \neg (v \sim_{M} x))\).
By propositional extensionality, we obtain \((\neg (v \sim_{M} x \land \neg (v \sim_{M^{'}} x) \lor v \sim_{M^{'}} x \land \neg (v \sim_{M} x))) = (\neg (v \sim_{M} x \land \neg (v \sim_{M^{'}} x)) \land \neg (v \sim_{M^{'}} x \land \neg (v \sim_{M} x)))\).
Chaining the equalities, we obtain \((x \notin N_{(\operatorname{spanning\,graph}(M) \setminus \operatorname{spanning\,graph}(M^{'})) \cup (\operatorname{spanning\,graph}(M^{'}) \setminus \operatorname{spanning\,graph}(M))}(v)) = (\neg (v \sim_{M} x \land \neg (v \sim_{M^{'}} x)) \land \neg (v \sim_{M^{'}} x \land \neg (v \sim_{M} x)))\).
By \paperlink{premise:6e6f745f616e64}{Prerequisite~P\ref{premise:6e6f745f616e64}}, we obtain \(\neg (v \sim_{M} x \land \neg (v \sim_{M^{'}} x)) \iff v \sim_{M} x \implies \neg (\neg (v \sim_{M^{'}} x))\).
By propositional extensionality, we obtain \((\neg (v \sim_{M} x \land \neg (v \sim_{M^{'}} x))) = (v \sim_{M} x \implies \neg (\neg (v \sim_{M^{'}} x)))\).
By definition, we obtain \(\left(v \sim_{M} x\right) = \left(v \sim_{M} x\right)\).
By \paperlink{premise:436c6173736963616c2e6e6f745f6e6f74}{Prerequisite~P\ref{premise:436c6173736963616c2e6e6f745f6e6f74}}, we obtain \(\neg \left(\neg \left(v \sim_{M^{'}} x\right)\right) \iff v \sim_{M^{'}} x\).
By propositional extensionality, we obtain \(\left(\neg \left(\neg \left(v \sim_{M^{'}} x\right)\right)\right) = \left(v \sim_{M^{'}} x\right)\).
By \paperlink{premise:696d706c6965735f636f6e6772}{Prerequisite~P\ref{premise:696d706c6965735f636f6e6772}} with \(p_{1} = v \sim_{M} x\) and \(p_{2} = v \sim_{M} x\) and \(q_{1} = \neg \left(\neg \left(v \sim_{M^{'}} x\right)\right)\) and \(q_{2} = v \sim_{M^{'}} x\), we obtain \((v \sim_{M} x \implies \neg (\neg (v \sim_{M^{'}} x))) = (v \sim_{M} x \implies v \sim_{M^{'}} x)\).
By \paperlink{premise:696d705f6966665f6e6f745f6f72}{Prerequisite~P\ref{premise:696d705f6966665f6e6f745f6f72}}, we obtain \(v \sim_{M} x \implies v \sim_{M^{'}} x \iff \neg \left(v \sim_{M} x\right) \lor v \sim_{M^{'}} x\).
By propositional extensionality, we obtain \((v \sim_{M} x \implies v \sim_{M^{'}} x) = (\neg (v \sim_{M} x) \lor v \sim_{M^{'}} x)\).
Chaining the equalities, we obtain \((v \sim_{M} x \implies \neg (\neg (v \sim_{M^{'}} x))) = (\neg (v \sim_{M} x) \lor v \sim_{M^{'}} x)\).
Chaining the equalities, we obtain \((\neg (v \sim_{M} x \land \neg (v \sim_{M^{'}} x))) = (\neg (v \sim_{M} x) \lor v \sim_{M^{'}} x)\).
Applying the same function to equal arguments, we obtain \((y \mapsto \neg (v \sim_{M} x \land \neg (v \sim_{M^{'}} x)) \land y) = (y \mapsto (\neg (v \sim_{M} x) \lor v \sim_{M^{'}} x) \land y)\).
By \paperlink{premise:6e6f745f616e64}{Prerequisite~P\ref{premise:6e6f745f616e64}}, we obtain \(\neg (v \sim_{M^{'}} x \land \neg (v \sim_{M} x)) \iff v \sim_{M^{'}} x \implies \neg (\neg (v \sim_{M} x))\).
By propositional extensionality, we obtain \((\neg (v \sim_{M^{'}} x \land \neg (v \sim_{M} x))) = (v \sim_{M^{'}} x \implies \neg (\neg (v \sim_{M} x)))\).
By definition, we obtain \(\left(v \sim_{M^{'}} x\right) = \left(v \sim_{M^{'}} x\right)\).
By \paperlink{premise:436c6173736963616c2e6e6f745f6e6f74}{Prerequisite~P\ref{premise:436c6173736963616c2e6e6f745f6e6f74}}, we obtain \(\neg \left(\neg \left(v \sim_{M} x\right)\right) \iff v \sim_{M} x\).
By propositional extensionality, we obtain \(\left(\neg \left(\neg \left(v \sim_{M} x\right)\right)\right) = \left(v \sim_{M} x\right)\).
By \paperlink{premise:696d706c6965735f636f6e6772}{Prerequisite~P\ref{premise:696d706c6965735f636f6e6772}} with \(p_{1} = v \sim_{M^{'}} x\) and \(p_{2} = v \sim_{M^{'}} x\) and \(q_{1} = \neg \left(\neg \left(v \sim_{M} x\right)\right)\) and \(q_{2} = v \sim_{M} x\), we obtain \((v \sim_{M^{'}} x \implies \neg (\neg (v \sim_{M} x))) = (v \sim_{M^{'}} x \implies v \sim_{M} x)\).
By \paperlink{premise:696d705f6966665f6e6f745f6f72}{Prerequisite~P\ref{premise:696d705f6966665f6e6f745f6f72}}, we obtain \(v \sim_{M^{'}} x \implies v \sim_{M} x \iff \neg \left(v \sim_{M^{'}} x\right) \lor v \sim_{M} x\).
By propositional extensionality, we obtain \((v \sim_{M^{'}} x \implies v \sim_{M} x) = (\neg (v \sim_{M^{'}} x) \lor v \sim_{M} x)\).
Chaining the equalities, we obtain \((v \sim_{M^{'}} x \implies \neg (\neg (v \sim_{M} x))) = (\neg (v \sim_{M^{'}} x) \lor v \sim_{M} x)\).
Chaining the equalities, we obtain \((\neg (v \sim_{M^{'}} x \land \neg (v \sim_{M} x))) = (\neg (v \sim_{M^{'}} x) \lor v \sim_{M} x)\).
By congruence, we obtain \((\neg (v \sim_{M} x \land \neg (v \sim_{M^{'}} x)) \land \neg (v \sim_{M^{'}} x \land \neg (v \sim_{M} x))) = ((\neg (v \sim_{M} x) \lor v \sim_{M^{'}} x) \land (\neg (v \sim_{M^{'}} x) \lor v \sim_{M} x))\).
Chaining the equalities, we obtain \((x \notin N_{(\operatorname{spanning\,graph}(M) \setminus \operatorname{spanning\,graph}(M^{'})) \cup (\operatorname{spanning\,graph}(M^{'}) \setminus \operatorname{spanning\,graph}(M))}(v)) = ((\neg (v \sim_{M} x) \lor v \sim_{M^{'}} x) \land (\neg (v \sim_{M^{'}} x) \lor v \sim_{M} x))\).
\end{proofstep}
Distinguish whether \(w = w^{'}\) holds.
Consider the following case.
\begin{proofstep}
Assume \(w = w^{'}\).
The substitution uses the following argument.
\begin{proofstep}
Fix \(x\) in \(V\).
By \paperlink{premise:696d705f6966665f6e6f745f6f72}{Prerequisite~P\ref{premise:696d705f6966665f6e6f745f6f72}}, we obtain \(v \sim_{M} x \implies v \sim_{M^{'}} x \iff \neg \left(v \sim_{M} x\right) \lor v \sim_{M^{'}} x\).
By propositional extensionality, we obtain \((v \sim_{M} x \implies v \sim_{M^{'}} x) = (\neg (v \sim_{M} x) \lor v \sim_{M^{'}} x)\).
Reversing the equality, we obtain \((\neg (v \sim_{M} x) \lor v \sim_{M^{'}} x) = (v \sim_{M} x \implies v \sim_{M^{'}} x)\).
By definition, we obtain \(\left(v \sim_{M} x\right) = \left(v \sim_{M} x\right)\).
For the required implication, argue as follows.
\begin{proofstep}
Assume \(v \sim_{M} x\).
The substitution uses the following argument.
\begin{proofstep}
Use the hypothesis \(v \sim_{M^{'}} w^{'} \land \left(\forall y\, :\, V,\; v \sim_{M^{'}} y \implies y = w^{'}\right)\).
Taking the second conjunct gives \(\forall y\, :\, V,\; v \sim_{M^{'}} y \implies y = w^{'}\).
The substitution uses the following argument.
\begin{proofstep}
Use the hypothesis \(v \sim_{M^{'}} w^{'} \land \left(\forall y\, :\, V,\; v \sim_{M^{'}} y \implies y = w^{'}\right)\).
Taking the second conjunct gives \(\forall y\, :\, V,\; v \sim_{M^{'}} y \implies y = w^{'}\).
Use the hypothesis \(v \sim_{M^{'}} w^{'} \land \left(\forall y\, :\, V,\; v \sim_{M^{'}} y \implies y = w^{'}\right)\).
Taking the first conjunct gives \(v \sim_{M^{'}} w^{'}\).
Rewriting with \(w = w^{'}\) and \(v \sim_{M^{'}} w^{'}\) gives \(v \sim_{M^{'}} w\).
Applying this conjunct to \(w\) gives \(w = w^{'}\).
\end{proofstep}
Use the hypothesis \(v \sim_{M^{'}} w^{'} \land \left(\forall y\, :\, V,\; v \sim_{M^{'}} y \implies y = w^{'}\right)\).
Taking the first conjunct gives \(v \sim_{M^{'}} w^{'}\).
Rewriting with \(v \sim_{M^{'}} w^{'}\) and the auxiliary result established above gives \(v \sim_{M^{'}} w\).
Applying this conjunct to \(w\) gives \(w = w^{'}\).
\end{proofstep}
The substitution uses the following argument.
\begin{proofstep}
Fix \(y\) in \(V\).
By definition, we obtain \(\left(v \sim_{M} y\right) = \left(v \sim_{M} y\right)\).
For the required implication, argue as follows.
\begin{proofstep}
Assume \(v \sim_{M} y\).
Use the hypothesis \(v \sim_{M^{'}} w^{'} \land \left(\forall y\, :\, V,\; v \sim_{M^{'}} y \implies y = w^{'}\right)\).
Taking the second conjunct gives \(\forall y_{1}\, :\, V,\; v \sim_{M^{'}} y_{1} \implies y_{1} = w^{'}\).
The substitution uses the following argument.
\begin{proofstep}
Use the hypothesis \(v \sim_{M^{'}} w^{'} \land \left(\forall y\, :\, V,\; v \sim_{M^{'}} y \implies y = w^{'}\right)\).
Taking the second conjunct gives \(\forall y_{1}\, :\, V,\; v \sim_{M^{'}} y_{1} \implies y_{1} = w^{'}\).
Use the hypothesis \(v \sim_{M^{'}} w^{'} \land \left(\forall y\, :\, V,\; v \sim_{M^{'}} y \implies y = w^{'}\right)\).
Taking the first conjunct gives \(v \sim_{M^{'}} w^{'}\).
Rewriting with \(w = w^{'}\) and \(v \sim_{M^{'}} w^{'}\) gives \(v \sim_{M^{'}} w\).
Applying this conjunct to \(w\) gives \(w = w^{'}\).
\end{proofstep}
Use the hypothesis \(v \sim_{M^{'}} w^{'} \land \left(\forall y\, :\, V,\; v \sim_{M^{'}} y \implies y = w^{'}\right)\).
Taking the first conjunct gives \(v \sim_{M^{'}} w^{'}\).
Rewriting with \(v \sim_{M^{'}} w^{'}\) and the auxiliary result established above gives \(v \sim_{M^{'}} w\).
Applying this conjunct to \(w\) gives \(w = w^{'}\).
Applying the same function to equal arguments, we obtain \(\left(y = w\right) = \left(y = w^{'}\right)\).
\end{proofstep}
By \paperlink{premise:696d706c6965735f636f6e67725f637478}{Prerequisite~P\ref{premise:696d706c6965735f636f6e67725f637478}}, we obtain \(\left(v \sim_{M} y \implies y = w\right) = \left(v \sim_{M} y \implies y = w^{'}\right)\).
\end{proofstep}
Use the hypothesis \(v \sim_{M} w \land \left(\forall y\, :\, V,\; v \sim_{M} y \implies y = w\right)\).
Rewriting with \paperlink{premise:666f72616c6c5f636f6e6772}{Prerequisite~P\ref{premise:666f72616c6c5f636f6e6772}} and the two auxiliary results established above in the preceding result gives \(v \sim_{M} w^{'} \land \left(\forall y\, :\, V,\; v \sim_{M} y \implies y = w^{'}\right)\).
By the preceding argument, we obtain \(v \sim_{M} w^{'} \land \left(\forall y\, :\, V,\; v \sim_{M} y \implies y = w^{'}\right)\).
Taking the second conjunct gives \(\forall y\, :\, V,\; v \sim_{M} y \implies y = w^{'}\).
The substitution uses the following argument.
\begin{proofstep}
The substitution uses the following argument.
\begin{proofstep}
Use the hypothesis \(v \sim_{M^{'}} w^{'} \land \left(\forall y\, :\, V,\; v \sim_{M^{'}} y \implies y = w^{'}\right)\).
Taking the second conjunct gives \(\forall y\, :\, V,\; v \sim_{M^{'}} y \implies y = w^{'}\).
The substitution uses the following argument.
\begin{proofstep}
Use the hypothesis \(v \sim_{M^{'}} w^{'} \land \left(\forall y\, :\, V,\; v \sim_{M^{'}} y \implies y = w^{'}\right)\).
Taking the second conjunct gives \(\forall y\, :\, V,\; v \sim_{M^{'}} y \implies y = w^{'}\).
Use the hypothesis \(v \sim_{M^{'}} w^{'} \land \left(\forall y\, :\, V,\; v \sim_{M^{'}} y \implies y = w^{'}\right)\).
Taking the first conjunct gives \(v \sim_{M^{'}} w^{'}\).
Rewriting with \(w = w^{'}\) and \(v \sim_{M^{'}} w^{'}\) gives \(v \sim_{M^{'}} w\).
Applying this conjunct to \(w\) gives \(w = w^{'}\).
\end{proofstep}
Use the hypothesis \(v \sim_{M^{'}} w^{'} \land \left(\forall y\, :\, V,\; v \sim_{M^{'}} y \implies y = w^{'}\right)\).
Taking the first conjunct gives \(v \sim_{M^{'}} w^{'}\).
Rewriting with \(v \sim_{M^{'}} w^{'}\) and the auxiliary result established above gives \(v \sim_{M^{'}} w\).
Applying this conjunct to \(w\) gives \(w = w^{'}\).
\end{proofstep}
The substitution uses the following argument.
\begin{proofstep}
Fix \(y\) in \(V\).
By definition, we obtain \(\left(v \sim_{M} y\right) = \left(v \sim_{M} y\right)\).
For the required implication, argue as follows.
\begin{proofstep}
Assume \(v \sim_{M} y\).
Use the hypothesis \(v \sim_{M^{'}} w^{'} \land \left(\forall y\, :\, V,\; v \sim_{M^{'}} y \implies y = w^{'}\right)\).
Taking the second conjunct gives \(\forall y_{1}\, :\, V,\; v \sim_{M^{'}} y_{1} \implies y_{1} = w^{'}\).
The substitution uses the following argument.
\begin{proofstep}
Use the hypothesis \(v \sim_{M^{'}} w^{'} \land \left(\forall y\, :\, V,\; v \sim_{M^{'}} y \implies y = w^{'}\right)\).
Taking the second conjunct gives \(\forall y_{1}\, :\, V,\; v \sim_{M^{'}} y_{1} \implies y_{1} = w^{'}\).
Use the hypothesis \(v \sim_{M^{'}} w^{'} \land \left(\forall y\, :\, V,\; v \sim_{M^{'}} y \implies y = w^{'}\right)\).
Taking the first conjunct gives \(v \sim_{M^{'}} w^{'}\).
Rewriting with \(w = w^{'}\) and \(v \sim_{M^{'}} w^{'}\) gives \(v \sim_{M^{'}} w\).
Applying this conjunct to \(w\) gives \(w = w^{'}\).
\end{proofstep}
Use the hypothesis \(v \sim_{M^{'}} w^{'} \land \left(\forall y\, :\, V,\; v \sim_{M^{'}} y \implies y = w^{'}\right)\).
Taking the first conjunct gives \(v \sim_{M^{'}} w^{'}\).
Rewriting with \(v \sim_{M^{'}} w^{'}\) and the auxiliary result established above gives \(v \sim_{M^{'}} w\).
Applying this conjunct to \(w\) gives \(w = w^{'}\).
Applying the same function to equal arguments, we obtain \(\left(y = w\right) = \left(y = w^{'}\right)\).
\end{proofstep}
By \paperlink{premise:696d706c6965735f636f6e67725f637478}{Prerequisite~P\ref{premise:696d706c6965735f636f6e67725f637478}}, we obtain \(\left(v \sim_{M} y \implies y = w\right) = \left(v \sim_{M} y \implies y = w^{'}\right)\).
\end{proofstep}
Use the hypothesis \(v \sim_{M} w \land \left(\forall y\, :\, V,\; v \sim_{M} y \implies y = w\right)\).
Rewriting with \paperlink{premise:666f72616c6c5f636f6e6772}{Prerequisite~P\ref{premise:666f72616c6c5f636f6e6772}} and the two auxiliary results established above in the preceding result gives \(v \sim_{M} w^{'} \land \left(\forall y\, :\, V,\; v \sim_{M} y \implies y = w^{'}\right)\).
By the preceding argument, we obtain \(v \sim_{M} w^{'} \land \left(\forall y\, :\, V,\; v \sim_{M} y \implies y = w^{'}\right)\).
Taking the second conjunct gives \(\forall y\, :\, V,\; v \sim_{M} y \implies y = w^{'}\).
Use the hypothesis \(v \sim_{M} x\).
Applying this conjunct to \(x\) gives \(x = w^{'}\).
\end{proofstep}
The substitution uses the following argument.
\begin{proofstep}
Use the hypothesis \(v \sim_{M^{'}} w^{'} \land \left(\forall y\, :\, V,\; v \sim_{M^{'}} y \implies y = w^{'}\right)\).
Taking the second conjunct gives \(\forall y\, :\, V,\; v \sim_{M^{'}} y \implies y = w^{'}\).
The substitution uses the following argument.
\begin{proofstep}
Use the hypothesis \(v \sim_{M^{'}} w^{'} \land \left(\forall y\, :\, V,\; v \sim_{M^{'}} y \implies y = w^{'}\right)\).
Taking the second conjunct gives \(\forall y\, :\, V,\; v \sim_{M^{'}} y \implies y = w^{'}\).
Use the hypothesis \(v \sim_{M^{'}} w^{'} \land \left(\forall y\, :\, V,\; v \sim_{M^{'}} y \implies y = w^{'}\right)\).
Taking the first conjunct gives \(v \sim_{M^{'}} w^{'}\).
Rewriting with \(w = w^{'}\) and \(v \sim_{M^{'}} w^{'}\) gives \(v \sim_{M^{'}} w\).
Applying this conjunct to \(w\) gives \(w = w^{'}\).
\end{proofstep}
Use the hypothesis \(v \sim_{M^{'}} w^{'} \land \left(\forall y\, :\, V,\; v \sim_{M^{'}} y \implies y = w^{'}\right)\).
Taking the first conjunct gives \(v \sim_{M^{'}} w^{'}\).
Rewriting with \(v \sim_{M^{'}} w^{'}\) and the auxiliary result established above gives \(v \sim_{M^{'}} w\).
Applying this conjunct to \(w\) gives \(w = w^{'}\).
\end{proofstep}
The substitution uses the following argument.
\begin{proofstep}
Fix \(y\) in \(V\).
By definition, we obtain \(\left(v \sim_{M} y\right) = \left(v \sim_{M} y\right)\).
For the required implication, argue as follows.
\begin{proofstep}
Assume \(v \sim_{M} y\).
Use the hypothesis \(v \sim_{M^{'}} w^{'} \land \left(\forall y\, :\, V,\; v \sim_{M^{'}} y \implies y = w^{'}\right)\).
Taking the second conjunct gives \(\forall y_{1}\, :\, V,\; v \sim_{M^{'}} y_{1} \implies y_{1} = w^{'}\).
The substitution uses the following argument.
\begin{proofstep}
Use the hypothesis \(v \sim_{M^{'}} w^{'} \land \left(\forall y\, :\, V,\; v \sim_{M^{'}} y \implies y = w^{'}\right)\).
Taking the second conjunct gives \(\forall y_{1}\, :\, V,\; v \sim_{M^{'}} y_{1} \implies y_{1} = w^{'}\).
Use the hypothesis \(v \sim_{M^{'}} w^{'} \land \left(\forall y\, :\, V,\; v \sim_{M^{'}} y \implies y = w^{'}\right)\).
Taking the first conjunct gives \(v \sim_{M^{'}} w^{'}\).
Rewriting with \(w = w^{'}\) and \(v \sim_{M^{'}} w^{'}\) gives \(v \sim_{M^{'}} w\).
Applying this conjunct to \(w\) gives \(w = w^{'}\).
\end{proofstep}
Use the hypothesis \(v \sim_{M^{'}} w^{'} \land \left(\forall y\, :\, V,\; v \sim_{M^{'}} y \implies y = w^{'}\right)\).
Taking the first conjunct gives \(v \sim_{M^{'}} w^{'}\).
Rewriting with \(v \sim_{M^{'}} w^{'}\) and the auxiliary result established above gives \(v \sim_{M^{'}} w\).
Applying this conjunct to \(w\) gives \(w = w^{'}\).
Applying the same function to equal arguments, we obtain \(\left(y = w\right) = \left(y = w^{'}\right)\).
\end{proofstep}
By \paperlink{premise:696d706c6965735f636f6e67725f637478}{Prerequisite~P\ref{premise:696d706c6965735f636f6e67725f637478}}, we obtain \(\left(v \sim_{M} y \implies y = w\right) = \left(v \sim_{M} y \implies y = w^{'}\right)\).
\end{proofstep}
Use the hypothesis \(v \sim_{M} w \land \left(\forall y\, :\, V,\; v \sim_{M} y \implies y = w\right)\).
Rewriting with \paperlink{premise:666f72616c6c5f636f6e6772}{Prerequisite~P\ref{premise:666f72616c6c5f636f6e6772}} and the two auxiliary results established above in the preceding result gives \(v \sim_{M} w^{'} \land \left(\forall y\, :\, V,\; v \sim_{M} y \implies y = w^{'}\right)\).
By the preceding argument, we obtain \(v \sim_{M} w^{'} \land \left(\forall y\, :\, V,\; v \sim_{M} y \implies y = w^{'}\right)\).
Taking the first conjunct gives \(v \sim_{M} w^{'}\).
Rewriting with \(v \sim_{M} w^{'}\) and the auxiliary result established above gives \(v \sim_{M} x\).
Applying this conjunct to \(x\) gives \(x = w^{'}\).
Applying the same function to equal arguments, we obtain \(\left(v \sim_{M^{'}} x\right) = \left(v \sim_{M^{'}} w^{'}\right)\).
Use the hypothesis \(v \sim_{M^{'}} w^{'} \land \left(\forall y\, :\, V,\; v \sim_{M^{'}} y \implies y = w^{'}\right)\).
By the preceding argument, we obtain \(v \sim_{M^{'}} w^{'} \land \left(\forall y\, :\, V,\; v \sim_{M^{'}} y \implies y = w^{'}\right)\).
By the preceding argument, we obtain \(v \sim_{M^{'}} w^{'} \land \left(\forall y\, :\, V,\; v \sim_{M^{'}} y \implies y = w^{'}\right)\).
Taking the first conjunct gives \(v \sim_{M^{'}} w^{'}\).
By \paperlink{premise:65715f74727565}{Prerequisite~P\ref{premise:65715f74727565}}, we obtain \(\left(v \sim_{M^{'}} w^{'}\right) = \mathrm{True}\).
Chaining the equalities, we obtain \(\left(v \sim_{M^{'}} x\right) = \mathrm{True}\).
\end{proofstep}
By \paperlink{premise:696d706c6965735f636f6e67725f637478}{Prerequisite~P\ref{premise:696d706c6965735f636f6e67725f637478}}, we obtain \(\left(v \sim_{M} x \implies v \sim_{M^{'}} x\right) = \left(v \sim_{M} x \implies \mathrm{True}\right)\).
By \paperlink{premise:696d706c6965735f74727565}{Prerequisite~P\ref{premise:696d706c6965735f74727565}} with \(\alpha = v \sim_{M} x\), we obtain \(\left(v \sim_{M} x \implies \mathrm{True}\right) = \mathrm{True}\).
Chaining the equalities, we obtain \(\left(v \sim_{M} x \implies v \sim_{M^{'}} x\right) = \mathrm{True}\).
Chaining the equalities, we obtain \(\left(\neg \left(v \sim_{M} x\right) \lor v \sim_{M^{'}} x\right) = \mathrm{True}\).
Applying the same function to equal arguments, we obtain \((y \mapsto (\neg (v \sim_{M} x) \lor v \sim_{M^{'}} x) \land y) = (y \mapsto \mathrm{True} \land y)\).
By \paperlink{premise:696d705f6966665f6e6f745f6f72}{Prerequisite~P\ref{premise:696d705f6966665f6e6f745f6f72}}, we obtain \(v \sim_{M^{'}} x \implies v \sim_{M} x \iff \neg \left(v \sim_{M^{'}} x\right) \lor v \sim_{M} x\).
By propositional extensionality, we obtain \((v \sim_{M^{'}} x \implies v \sim_{M} x) = (\neg (v \sim_{M^{'}} x) \lor v \sim_{M} x)\).
Reversing the equality, we obtain \((\neg (v \sim_{M^{'}} x) \lor v \sim_{M} x) = (v \sim_{M^{'}} x \implies v \sim_{M} x)\).
By definition, we obtain \(\left(v \sim_{M^{'}} x\right) = \left(v \sim_{M^{'}} x\right)\).
For the required implication, argue as follows.
\begin{proofstep}
Assume \(v \sim_{M^{'}} x\).
The substitution uses the following argument.
\begin{proofstep}
Use the hypothesis \(v \sim_{M^{'}} w^{'} \land \left(\forall y\, :\, V,\; v \sim_{M^{'}} y \implies y = w^{'}\right)\).
Taking the second conjunct gives \(\forall y\, :\, V,\; v \sim_{M^{'}} y \implies y = w^{'}\).
The substitution uses the following argument.
\begin{proofstep}
Use the hypothesis \(v \sim_{M^{'}} w^{'} \land \left(\forall y\, :\, V,\; v \sim_{M^{'}} y \implies y = w^{'}\right)\).
Taking the second conjunct gives \(\forall y\, :\, V,\; v \sim_{M^{'}} y \implies y = w^{'}\).
Use the hypothesis \(v \sim_{M^{'}} w^{'} \land \left(\forall y\, :\, V,\; v \sim_{M^{'}} y \implies y = w^{'}\right)\).
Taking the first conjunct gives \(v \sim_{M^{'}} w^{'}\).
Rewriting with \(w = w^{'}\) and \(v \sim_{M^{'}} w^{'}\) gives \(v \sim_{M^{'}} w\).
Applying this conjunct to \(w\) gives \(w = w^{'}\).
\end{proofstep}
Use the hypothesis \(v \sim_{M^{'}} w^{'} \land \left(\forall y\, :\, V,\; v \sim_{M^{'}} y \implies y = w^{'}\right)\).
Taking the first conjunct gives \(v \sim_{M^{'}} w^{'}\).
Rewriting with \(v \sim_{M^{'}} w^{'}\) and the auxiliary result established above gives \(v \sim_{M^{'}} w\).
Applying this conjunct to \(w\) gives \(w = w^{'}\).
\end{proofstep}
The substitution uses the following argument.
\begin{proofstep}
Fix \(y\) in \(V\).
By definition, we obtain \(\left(v \sim_{M} y\right) = \left(v \sim_{M} y\right)\).
For the required implication, argue as follows.
\begin{proofstep}
Assume \(v \sim_{M} y\).
Use the hypothesis \(v \sim_{M^{'}} w^{'} \land \left(\forall y\, :\, V,\; v \sim_{M^{'}} y \implies y = w^{'}\right)\).
Taking the second conjunct gives \(\forall y_{1}\, :\, V,\; v \sim_{M^{'}} y_{1} \implies y_{1} = w^{'}\).
The substitution uses the following argument.
\begin{proofstep}
Use the hypothesis \(v \sim_{M^{'}} w^{'} \land \left(\forall y\, :\, V,\; v \sim_{M^{'}} y \implies y = w^{'}\right)\).
Taking the second conjunct gives \(\forall y_{1}\, :\, V,\; v \sim_{M^{'}} y_{1} \implies y_{1} = w^{'}\).
Use the hypothesis \(v \sim_{M^{'}} w^{'} \land \left(\forall y\, :\, V,\; v \sim_{M^{'}} y \implies y = w^{'}\right)\).
Taking the first conjunct gives \(v \sim_{M^{'}} w^{'}\).
Rewriting with \(w = w^{'}\) and \(v \sim_{M^{'}} w^{'}\) gives \(v \sim_{M^{'}} w\).
Applying this conjunct to \(w\) gives \(w = w^{'}\).
\end{proofstep}
Use the hypothesis \(v \sim_{M^{'}} w^{'} \land \left(\forall y\, :\, V,\; v \sim_{M^{'}} y \implies y = w^{'}\right)\).
Taking the first conjunct gives \(v \sim_{M^{'}} w^{'}\).
Rewriting with \(v \sim_{M^{'}} w^{'}\) and the auxiliary result established above gives \(v \sim_{M^{'}} w\).
Applying this conjunct to \(w\) gives \(w = w^{'}\).
Applying the same function to equal arguments, we obtain \(\left(y = w\right) = \left(y = w^{'}\right)\).
\end{proofstep}
By \paperlink{premise:696d706c6965735f636f6e67725f637478}{Prerequisite~P\ref{premise:696d706c6965735f636f6e67725f637478}}, we obtain \(\left(v \sim_{M} y \implies y = w\right) = \left(v \sim_{M} y \implies y = w^{'}\right)\).
\end{proofstep}
Use the hypothesis \(v \sim_{M} w \land \left(\forall y\, :\, V,\; v \sim_{M} y \implies y = w\right)\).
Rewriting with \paperlink{premise:666f72616c6c5f636f6e6772}{Prerequisite~P\ref{premise:666f72616c6c5f636f6e6772}} and the two auxiliary results established above in the preceding result gives \(v \sim_{M} w^{'} \land \left(\forall y\, :\, V,\; v \sim_{M} y \implies y = w^{'}\right)\).
By the preceding argument, we obtain \(v \sim_{M} w^{'} \land \left(\forall y\, :\, V,\; v \sim_{M} y \implies y = w^{'}\right)\).
Taking the second conjunct gives \(\forall y\, :\, V,\; v \sim_{M} y \implies y = w^{'}\).
The substitution uses the following argument.
\begin{proofstep}
Use the hypothesis \(v \sim_{M^{'}} w^{'} \land \left(\forall y\, :\, V,\; v \sim_{M^{'}} y \implies y = w^{'}\right)\).
By the preceding argument, we obtain \(v \sim_{M^{'}} w^{'} \land \left(\forall y\, :\, V,\; v \sim_{M^{'}} y \implies y = w^{'}\right)\).
By the preceding argument, we obtain \(v \sim_{M^{'}} w^{'} \land \left(\forall y\, :\, V,\; v \sim_{M^{'}} y \implies y = w^{'}\right)\).
Taking the second conjunct gives \(\forall y\, :\, V,\; v \sim_{M^{'}} y \implies y = w^{'}\).
Use the hypothesis \(v \sim_{M^{'}} x\).
Applying this conjunct to \(x\) gives \(x = w^{'}\).
\end{proofstep}
The substitution uses the following argument.
\begin{proofstep}
Use the hypothesis \(v \sim_{M^{'}} w^{'} \land \left(\forall y\, :\, V,\; v \sim_{M^{'}} y \implies y = w^{'}\right)\).
Taking the second conjunct gives \(\forall y\, :\, V,\; v \sim_{M^{'}} y \implies y = w^{'}\).
The substitution uses the following argument.
\begin{proofstep}
Use the hypothesis \(v \sim_{M^{'}} w^{'} \land \left(\forall y\, :\, V,\; v \sim_{M^{'}} y \implies y = w^{'}\right)\).
Taking the second conjunct gives \(\forall y\, :\, V,\; v \sim_{M^{'}} y \implies y = w^{'}\).
Use the hypothesis \(v \sim_{M^{'}} w^{'} \land \left(\forall y\, :\, V,\; v \sim_{M^{'}} y \implies y = w^{'}\right)\).
Taking the first conjunct gives \(v \sim_{M^{'}} w^{'}\).
Rewriting with \(w = w^{'}\) and \(v \sim_{M^{'}} w^{'}\) gives \(v \sim_{M^{'}} w\).
Applying this conjunct to \(w\) gives \(w = w^{'}\).
\end{proofstep}
Use the hypothesis \(v \sim_{M^{'}} w^{'} \land \left(\forall y\, :\, V,\; v \sim_{M^{'}} y \implies y = w^{'}\right)\).
Taking the first conjunct gives \(v \sim_{M^{'}} w^{'}\).
Rewriting with \(v \sim_{M^{'}} w^{'}\) and the auxiliary result established above gives \(v \sim_{M^{'}} w\).
Applying this conjunct to \(w\) gives \(w = w^{'}\).
\end{proofstep}
The substitution uses the following argument.
\begin{proofstep}
Fix \(y\) in \(V\).
By definition, we obtain \(\left(v \sim_{M} y\right) = \left(v \sim_{M} y\right)\).
For the required implication, argue as follows.
\begin{proofstep}
Assume \(v \sim_{M} y\).
Use the hypothesis \(v \sim_{M^{'}} w^{'} \land \left(\forall y\, :\, V,\; v \sim_{M^{'}} y \implies y = w^{'}\right)\).
Taking the second conjunct gives \(\forall y_{1}\, :\, V,\; v \sim_{M^{'}} y_{1} \implies y_{1} = w^{'}\).
The substitution uses the following argument.
\begin{proofstep}
Use the hypothesis \(v \sim_{M^{'}} w^{'} \land \left(\forall y\, :\, V,\; v \sim_{M^{'}} y \implies y = w^{'}\right)\).
Taking the second conjunct gives \(\forall y_{1}\, :\, V,\; v \sim_{M^{'}} y_{1} \implies y_{1} = w^{'}\).
Use the hypothesis \(v \sim_{M^{'}} w^{'} \land \left(\forall y\, :\, V,\; v \sim_{M^{'}} y \implies y = w^{'}\right)\).
Taking the first conjunct gives \(v \sim_{M^{'}} w^{'}\).
Rewriting with \(w = w^{'}\) and \(v \sim_{M^{'}} w^{'}\) gives \(v \sim_{M^{'}} w\).
Applying this conjunct to \(w\) gives \(w = w^{'}\).
\end{proofstep}
Use the hypothesis \(v \sim_{M^{'}} w^{'} \land \left(\forall y\, :\, V,\; v \sim_{M^{'}} y \implies y = w^{'}\right)\).
Taking the first conjunct gives \(v \sim_{M^{'}} w^{'}\).
Rewriting with \(v \sim_{M^{'}} w^{'}\) and the auxiliary result established above gives \(v \sim_{M^{'}} w\).
Applying this conjunct to \(w\) gives \(w = w^{'}\).
Applying the same function to equal arguments, we obtain \(\left(y = w\right) = \left(y = w^{'}\right)\).
\end{proofstep}
By \paperlink{premise:696d706c6965735f636f6e67725f637478}{Prerequisite~P\ref{premise:696d706c6965735f636f6e67725f637478}}, we obtain \(\left(v \sim_{M} y \implies y = w\right) = \left(v \sim_{M} y \implies y = w^{'}\right)\).
\end{proofstep}
Use the hypothesis \(v \sim_{M} w \land \left(\forall y\, :\, V,\; v \sim_{M} y \implies y = w\right)\).
Rewriting with \paperlink{premise:666f72616c6c5f636f6e6772}{Prerequisite~P\ref{premise:666f72616c6c5f636f6e6772}} and the two auxiliary results established above in the preceding result gives \(v \sim_{M} w^{'} \land \left(\forall y\, :\, V,\; v \sim_{M} y \implies y = w^{'}\right)\).
By the preceding argument, we obtain \(v \sim_{M} w^{'} \land \left(\forall y\, :\, V,\; v \sim_{M} y \implies y = w^{'}\right)\).
Taking the first conjunct gives \(v \sim_{M} w^{'}\).
Rewriting with \(v \sim_{M} w^{'}\) and the auxiliary result established above gives \(v \sim_{M} x\).
Applying this conjunct to \(x\) gives \(x = w^{'}\).
Applying the same function to equal arguments, we obtain \(\left(v \sim_{M} x\right) = \left(v \sim_{M} w^{'}\right)\).
The substitution uses the following argument.
\begin{proofstep}
Use the hypothesis \(v \sim_{M^{'}} w^{'} \land \left(\forall y\, :\, V,\; v \sim_{M^{'}} y \implies y = w^{'}\right)\).
Taking the second conjunct gives \(\forall y\, :\, V,\; v \sim_{M^{'}} y \implies y = w^{'}\).
The substitution uses the following argument.
\begin{proofstep}
Use the hypothesis \(v \sim_{M^{'}} w^{'} \land \left(\forall y\, :\, V,\; v \sim_{M^{'}} y \implies y = w^{'}\right)\).
Taking the second conjunct gives \(\forall y\, :\, V,\; v \sim_{M^{'}} y \implies y = w^{'}\).
Use the hypothesis \(v \sim_{M^{'}} w^{'} \land \left(\forall y\, :\, V,\; v \sim_{M^{'}} y \implies y = w^{'}\right)\).
Taking the first conjunct gives \(v \sim_{M^{'}} w^{'}\).
Rewriting with \(w = w^{'}\) and \(v \sim_{M^{'}} w^{'}\) gives \(v \sim_{M^{'}} w\).
Applying this conjunct to \(w\) gives \(w = w^{'}\).
\end{proofstep}
Use the hypothesis \(v \sim_{M^{'}} w^{'} \land \left(\forall y\, :\, V,\; v \sim_{M^{'}} y \implies y = w^{'}\right)\).
Taking the first conjunct gives \(v \sim_{M^{'}} w^{'}\).
Rewriting with \(v \sim_{M^{'}} w^{'}\) and the auxiliary result established above gives \(v \sim_{M^{'}} w\).
Applying this conjunct to \(w\) gives \(w = w^{'}\).
\end{proofstep}
The substitution uses the following argument.
\begin{proofstep}
Fix \(y\) in \(V\).
By definition, we obtain \(\left(v \sim_{M} y\right) = \left(v \sim_{M} y\right)\).
For the required implication, argue as follows.
\begin{proofstep}
Assume \(v \sim_{M} y\).
Use the hypothesis \(v \sim_{M^{'}} w^{'} \land \left(\forall y\, :\, V,\; v \sim_{M^{'}} y \implies y = w^{'}\right)\).
Taking the second conjunct gives \(\forall y_{1}\, :\, V,\; v \sim_{M^{'}} y_{1} \implies y_{1} = w^{'}\).
The substitution uses the following argument.
\begin{proofstep}
Use the hypothesis \(v \sim_{M^{'}} w^{'} \land \left(\forall y\, :\, V,\; v \sim_{M^{'}} y \implies y = w^{'}\right)\).
Taking the second conjunct gives \(\forall y_{1}\, :\, V,\; v \sim_{M^{'}} y_{1} \implies y_{1} = w^{'}\).
Use the hypothesis \(v \sim_{M^{'}} w^{'} \land \left(\forall y\, :\, V,\; v \sim_{M^{'}} y \implies y = w^{'}\right)\).
Taking the first conjunct gives \(v \sim_{M^{'}} w^{'}\).
Rewriting with \(w = w^{'}\) and \(v \sim_{M^{'}} w^{'}\) gives \(v \sim_{M^{'}} w\).
Applying this conjunct to \(w\) gives \(w = w^{'}\).
\end{proofstep}
Use the hypothesis \(v \sim_{M^{'}} w^{'} \land \left(\forall y\, :\, V,\; v \sim_{M^{'}} y \implies y = w^{'}\right)\).
Taking the first conjunct gives \(v \sim_{M^{'}} w^{'}\).
Rewriting with \(v \sim_{M^{'}} w^{'}\) and the auxiliary result established above gives \(v \sim_{M^{'}} w\).
Applying this conjunct to \(w\) gives \(w = w^{'}\).
Applying the same function to equal arguments, we obtain \(\left(y = w\right) = \left(y = w^{'}\right)\).
\end{proofstep}
By \paperlink{premise:696d706c6965735f636f6e67725f637478}{Prerequisite~P\ref{premise:696d706c6965735f636f6e67725f637478}}, we obtain \(\left(v \sim_{M} y \implies y = w\right) = \left(v \sim_{M} y \implies y = w^{'}\right)\).
\end{proofstep}
Use the hypothesis \(v \sim_{M} w \land \left(\forall y\, :\, V,\; v \sim_{M} y \implies y = w\right)\).
Rewriting with \paperlink{premise:666f72616c6c5f636f6e6772}{Prerequisite~P\ref{premise:666f72616c6c5f636f6e6772}} and the two auxiliary results established above in the preceding result gives \(v \sim_{M} w^{'} \land \left(\forall y\, :\, V,\; v \sim_{M} y \implies y = w^{'}\right)\).
By the preceding argument, we obtain \(v \sim_{M} w^{'} \land \left(\forall y\, :\, V,\; v \sim_{M} y \implies y = w^{'}\right)\).
Taking the first conjunct gives \(v \sim_{M} w^{'}\).
By \paperlink{premise:65715f74727565}{Prerequisite~P\ref{premise:65715f74727565}}, we obtain \(\left(v \sim_{M} w^{'}\right) = \mathrm{True}\).
Chaining the equalities, we obtain \(\left(v \sim_{M} x\right) = \mathrm{True}\).
\end{proofstep}
By \paperlink{premise:696d706c6965735f636f6e67725f637478}{Prerequisite~P\ref{premise:696d706c6965735f636f6e67725f637478}}, we obtain \(\left(v \sim_{M^{'}} x \implies v \sim_{M} x\right) = \left(v \sim_{M^{'}} x \implies \mathrm{True}\right)\).
By \paperlink{premise:696d706c6965735f74727565}{Prerequisite~P\ref{premise:696d706c6965735f74727565}} with \(\alpha = v \sim_{M^{'}} x\), we obtain \(\left(v \sim_{M^{'}} x \implies \mathrm{True}\right) = \mathrm{True}\).
Chaining the equalities, we obtain \(\left(v \sim_{M^{'}} x \implies v \sim_{M} x\right) = \mathrm{True}\).
Chaining the equalities, we obtain \(\left(\neg \left(v \sim_{M^{'}} x\right) \lor v \sim_{M} x\right) = \mathrm{True}\).
By congruence, we obtain \(((\neg (v \sim_{M} x) \lor v \sim_{M^{'}} x) \land (\neg (v \sim_{M^{'}} x) \lor v \sim_{M} x)) = (\mathrm{True} \land \mathrm{True})\).
By \paperlink{premise:616e645f73656c66}{Prerequisite~P\ref{premise:616e645f73656c66}} with \(p = \mathrm{True}\), we obtain \(\left(\mathrm{True} \land \mathrm{True}\right) = \mathrm{True}\).
Chaining the equalities, we obtain \(((\neg (v \sim_{M} x) \lor v \sim_{M^{'}} x) \land (\neg (v \sim_{M^{'}} x) \lor v \sim_{M} x)) = \mathrm{True}\).
\end{proofstep}
Rewriting with \paperlink{premise:666f72616c6c5f636f6e6772}{Prerequisite~P\ref{premise:666f72616c6c5f636f6e6772}}, \paperlink{premise:696d706c6965735f74727565}{Prerequisite~P\ref{premise:696d706c6965735f74727565}}, \paperlink{premise:747275655f6f72}{Prerequisite~P\ref{premise:747275655f6f72}} and the identity established above gives \((\forall a\, :\, V,\; (\neg (v \sim_{M} a) \lor v \sim_{M^{'}} a) \land (\neg (v \sim_{M^{'}} a) \lor v \sim_{M} a)) \lor (\exists x\, :\, V,\; \exists y\, :\, V,\; \neg (x = y) \land N_{(\operatorname{spanning\,graph}(M) \setminus \operatorname{spanning\,graph}(M^{'})) \cup (\operatorname{spanning\,graph}(M^{'}) \setminus \operatorname{spanning\,graph}(M))}(v) = \{x\} \cup \{y\})\).
\end{proofstep}
Consider the following case.
\begin{proofstep}
Assume \(\neg \left(w = w^{'}\right)\).
Choose the witness \(w\).
Choose the witness \(w^{'}\).
Use the hypothesis \(v \sim_{M} w \land \left(\forall y\, :\, V,\; v \sim_{M} y \implies y = w\right)\).
By the preceding argument, we obtain \(v \sim_{M} w \land \left(\forall y\, :\, V,\; v \sim_{M} y \implies y = w\right)\).
Consider the following case.
\begin{proofstep}
Assume \(v \sim_{M} w\).
Assume \(\forall y\, :\, V,\; v \sim_{M} y \implies y = w\).
Use the hypothesis \(v \sim_{M^{'}} w^{'} \land \left(\forall y\, :\, V,\; v \sim_{M^{'}} y \implies y = w^{'}\right)\).
By the preceding argument, we obtain \(v \sim_{M^{'}} w^{'} \land \left(\forall y\, :\, V,\; v \sim_{M^{'}} y \implies y = w^{'}\right)\).
Consider the following case.
\begin{proofstep}
Assume \(v \sim_{M^{'}} w^{'}\).
Assume \(\forall y\, :\, V,\; v \sim_{M^{'}} y \implies y = w^{'}\).
For the required implication, argue as follows.
\begin{proofstep}
Fix \(x\) in \(V\).
Consider the following case.
\begin{proofstep}
Assume \(v \sim_{M} x \land \neg (v \sim_{M^{'}} x) \lor v \sim_{M^{'}} x \land \neg (v \sim_{M} x)\).
Use the hypothesis \(v \sim_{M} x \land \neg (v \sim_{M^{'}} x) \lor v \sim_{M^{'}} x \land \neg (v \sim_{M} x)\).
Consider the following case.
\begin{proofstep}
Assume \(v \sim_{M} x \land \neg \left(v \sim_{M^{'}} x\right)\).
Use the hypothesis \(v \sim_{M} x \land \neg \left(v \sim_{M^{'}} x\right)\).
Taking the first conjunct gives \(v \sim_{M} x\).
Use the hypothesis \(v \sim_{M} x \land \neg \left(v \sim_{M^{'}} x\right)\).
Taking the first conjunct gives \(v \sim_{M} x\).
Rewriting with \(\forall y\, :\, V,\; v \sim_{M} y \implies y = w\), \(v \sim_{M} x\), \(v \sim_{M} w\), \paperlink{premise:65715f66616c7365}{Prerequisite~P\ref{premise:65715f66616c7365}} (using \(\neg \left(w = w^{'}\right)\)) and \paperlink{premise:6f725f66616c7365}{Prerequisite~P\ref{premise:6f725f66616c7365}} gives \(x = w \lor x = w^{'}\).
\end{proofstep}
Consider the following case.
\begin{proofstep}
Assume \(v \sim_{M^{'}} x \land \neg \left(v \sim_{M} x\right)\).
Use the hypothesis \(v \sim_{M^{'}} x \land \neg \left(v \sim_{M} x\right)\).
Taking the first conjunct gives \(v \sim_{M^{'}} x\).
Use the hypothesis \(v \sim_{M^{'}} x \land \neg \left(v \sim_{M} x\right)\).
Taking the first conjunct gives \(v \sim_{M^{'}} x\).
Rewriting with \(\forall y\, :\, V,\; v \sim_{M^{'}} y \implies y = w^{'}\), \(v \sim_{M^{'}} x\), \(v \sim_{M^{'}} w^{'}\) and \paperlink{premise:6f725f74727565}{Prerequisite~P\ref{premise:6f725f74727565}} gives \(x = w \lor x = w^{'}\).
\end{proofstep}
The two alternatives exhaust the cases. Hence \(x = w \lor x = w^{'}\).
\end{proofstep}
Consider the following case.
\begin{proofstep}
Assume \(x = w \lor x = w^{'}\).
Use the hypothesis \(x = w \lor x = w^{'}\).
Consider the following case.
\begin{proofstep}
Assume \(x = w\).
For the required implication, argue as follows.
\begin{proofstep}
Assume \(\neg \left(x = w^{'}\right)\).
Assume \(v \sim_{M} x\).
Assume \(\forall y\, :\, V,\; v \sim_{M} y \implies y = x\).
For the required implication, argue as follows.
\begin{proofstep}
Assume \(v \sim_{M^{'}} x\).
Use the hypothesis \(\neg \left(x = w^{'}\right)\).
Rewriting with \(\forall y\, :\, V,\; v \sim_{M^{'}} y \implies y = w^{'}\), \(v \sim_{M^{'}} x\), \(v \sim_{M^{'}} w^{'}\) and \paperlink{premise:6e6f745f747275655f65715f66616c7365}{Prerequisite~P\ref{premise:6e6f745f747275655f65715f66616c7365}} in the preceding result gives \(\mathrm{False}\).
The contradiction rules out this case. Hence \(\mathrm{False}\).
\end{proofstep}
By \paperlink{premise:4165736f702e4275696c74696e52756c65732e6e6f745f696e74726f}{Prerequisite~P\ref{premise:4165736f702e4275696c74696e52756c65732e6e6f745f696e74726f}}, we obtain \(\neg \left(v \sim_{M^{'}} x\right)\).
Rewriting with \(v \sim_{M} x\), \paperlink{premise:747275655f616e64}{Prerequisite~P\ref{premise:747275655f616e64}}, \paperlink{premise:6e6f745f747275655f65715f66616c7365}{Prerequisite~P\ref{premise:6e6f745f747275655f65715f66616c7365}}, \paperlink{premise:616e645f66616c7365}{Prerequisite~P\ref{premise:616e645f66616c7365}} and \paperlink{premise:6f725f66616c7365}{Prerequisite~P\ref{premise:6f725f66616c7365}} in the preceding result gives \(v \sim_{M} x \land \neg (v \sim_{M^{'}} x) \lor v \sim_{M^{'}} x \land \neg (v \sim_{M} x)\).
\end{proofstep}
Use the hypothesis \(x = w\).
Use the hypothesis \(\neg \left(w = w^{'}\right)\).
Use the hypothesis \(v \sim_{M} w\).
Use the hypothesis \(\forall y\, :\, V,\; v \sim_{M} y \implies y = w\).
Substituting equals, we obtain \(v \sim_{M} x \land \neg (v \sim_{M^{'}} x) \lor v \sim_{M^{'}} x \land \neg (v \sim_{M} x)\).
\end{proofstep}
Consider the following case.
\begin{proofstep}
Assume \(x = w^{'}\).
For the required implication, argue as follows.
\begin{proofstep}
Assume \(\neg \left(w = x\right)\).
Assume \(v \sim_{M^{'}} x\).
Assume \(\forall y\, :\, V,\; v \sim_{M^{'}} y \implies y = x\).
For the required implication, argue as follows.
\begin{proofstep}
Assume \(v \sim_{M} x\).
Use the hypothesis \(\neg \left(w = x\right)\).
Rewriting with \(\forall y\, :\, V,\; v \sim_{M} y \implies y = w\), \(v \sim_{M} x\), \(v \sim_{M} w\) and \paperlink{premise:6e6f745f747275655f65715f66616c7365}{Prerequisite~P\ref{premise:6e6f745f747275655f65715f66616c7365}} in the preceding result gives \(\mathrm{False}\).
The contradiction rules out this case. Hence \(\mathrm{False}\).
\end{proofstep}
By \paperlink{premise:4165736f702e4275696c74696e52756c65732e6e6f745f696e74726f}{Prerequisite~P\ref{premise:4165736f702e4275696c74696e52756c65732e6e6f745f696e74726f}}, we obtain \(\neg \left(v \sim_{M} x\right)\).
Rewriting with \(v \sim_{M^{'}} x\), \paperlink{premise:6e6f745f747275655f65715f66616c7365}{Prerequisite~P\ref{premise:6e6f745f747275655f65715f66616c7365}}, \paperlink{premise:616e645f66616c7365}{Prerequisite~P\ref{premise:616e645f66616c7365}}, \paperlink{premise:747275655f616e64}{Prerequisite~P\ref{premise:747275655f616e64}} and \paperlink{premise:66616c73655f6f72}{Prerequisite~P\ref{premise:66616c73655f6f72}} in the preceding result gives \(v \sim_{M} x \land \neg (v \sim_{M^{'}} x) \lor v \sim_{M^{'}} x \land \neg (v \sim_{M} x)\).
\end{proofstep}
Use the hypothesis \(x = w^{'}\).
Use the hypothesis \(\neg \left(w = w^{'}\right)\).
Use the hypothesis \(v \sim_{M^{'}} w^{'}\).
Use the hypothesis \(\forall y\, :\, V,\; v \sim_{M^{'}} y \implies y = w^{'}\).
Substituting equals, we obtain \(v \sim_{M} x \land \neg (v \sim_{M^{'}} x) \lor v \sim_{M^{'}} x \land \neg (v \sim_{M} x)\).
\end{proofstep}
The two alternatives exhaust the cases. Hence \(v \sim_{M} x \land \neg (v \sim_{M^{'}} x) \lor v \sim_{M^{'}} x \land \neg (v \sim_{M} x)\).
\end{proofstep}
Together the two implications give the equivalence. Hence \(v \sim_{M} x \land \neg (v \sim_{M^{'}} x) \lor v \sim_{M^{'}} x \land \neg (v \sim_{M} x) \iff x = w \lor x = w^{'}\).
Rewriting with \paperlink{premise:53696d706c6547726170682e6d656d5f6e65696768626f72536574}{Prerequisite~P\ref{premise:53696d706c6547726170682e6d656d5f6e65696768626f72536574}}, \paperlink{premise:53696d706c6547726170682e7375705f61646a}{Prerequisite~P\ref{premise:53696d706c6547726170682e7375705f61646a}}, \paperlink{premise:53696d706c6547726170682e73646966665f61646a}{Prerequisite~P\ref{premise:53696d706c6547726170682e73646966665f61646a}}, \paperlink{premise:5365742e6d656d5f696e736572745f696666}{Prerequisite~P\ref{premise:5365742e6d656d5f696e736572745f696666}} and \paperlink{premise:5365742e6d656d5f73696e676c65746f6e5f696666}{Prerequisite~P\ref{premise:5365742e6d656d5f73696e676c65746f6e5f696666}} in the preceding result gives \(x \in N_{(\operatorname{spanning\,graph}(M) \setminus \operatorname{spanning\,graph}(M^{'})) \cup (\operatorname{spanning\,graph}(M^{'}) \setminus \operatorname{spanning\,graph}(M))}(v) \iff x \in \{w\} \cup \{w^{'}\}\).
\end{proofstep}
By \paperlink{premise:5365742e657874}{Prerequisite~P\ref{premise:5365742e657874}}, we obtain \(N_{(\operatorname{spanning\,graph}(M) \setminus \operatorname{spanning\,graph}(M^{'})) \cup (\operatorname{spanning\,graph}(M^{'}) \setminus \operatorname{spanning\,graph}(M))}(v) = \{w\} \cup \{w^{'}\}\).
\end{proofstep}
Using the two conjuncts, we obtain \(N_{(\operatorname{spanning\,graph}(M) \setminus \operatorname{spanning\,graph}(M^{'})) \cup (\operatorname{spanning\,graph}(M^{'}) \setminus \operatorname{spanning\,graph}(M))}(v) = \{w\} \cup \{w^{'}\}\).
\end{proofstep}
Using the two conjuncts, we obtain \(N_{(\operatorname{spanning\,graph}(M) \setminus \operatorname{spanning\,graph}(M^{'})) \cup (\operatorname{spanning\,graph}(M^{'}) \setminus \operatorname{spanning\,graph}(M))}(v) = \{w\} \cup \{w^{'}\}\).
Rewriting with \paperlink{premise:65715f66616c7365}{Prerequisite~P\ref{premise:65715f66616c7365}} (using \(\neg \left(w = w^{'}\right)\)), \paperlink{premise:6e6f745f66616c73655f65715f74727565}{Prerequisite~P\ref{premise:6e6f745f66616c73655f65715f74727565}} and \paperlink{premise:747275655f616e64}{Prerequisite~P\ref{premise:747275655f616e64}} in the preceding result gives \(\neg (w = w^{'}) \land N_{(\operatorname{spanning\,graph}(M) \setminus \operatorname{spanning\,graph}(M^{'})) \cup (\operatorname{spanning\,graph}(M^{'}) \setminus \operatorname{spanning\,graph}(M))}(v) = \{w\} \cup \{w^{'}\}\).
This witness proves existence. Hence \(\exists y\, :\, V,\; \neg (w = y) \land N_{(\operatorname{spanning\,graph}(M) \setminus \operatorname{spanning\,graph}(M^{'})) \cup (\operatorname{spanning\,graph}(M^{'}) \setminus \operatorname{spanning\,graph}(M))}(v) = \{w\} \cup \{y\}\).
This witness proves existence. Hence \(\exists x\, :\, V,\; \exists y\, :\, V,\; \neg (x = y) \land N_{(\operatorname{spanning\,graph}(M) \setminus \operatorname{spanning\,graph}(M^{'})) \cup (\operatorname{spanning\,graph}(M^{'}) \setminus \operatorname{spanning\,graph}(M))}(v) = \{x\} \cup \{y\}\).
The second alternative suffices. Hence \((\forall x\, :\, V,\; (\neg (v \sim_{M} x) \lor v \sim_{M^{'}} x) \land (\neg (v \sim_{M^{'}} x) \lor v \sim_{M} x)) \lor (\exists x\, :\, V,\; \exists y\, :\, V,\; \neg (x = y) \land N_{(\operatorname{spanning\,graph}(M) \setminus \operatorname{spanning\,graph}(M^{'})) \cup (\operatorname{spanning\,graph}(M^{'}) \setminus \operatorname{spanning\,graph}(M))}(v) = \{x\} \cup \{y\})\).
\end{proofstep}
The proposition or its negation holds. Hence \((\forall x\, :\, V,\; (\neg (v \sim_{M} x) \lor v \sim_{M^{'}} x) \land (\neg (v \sim_{M^{'}} x) \lor v \sim_{M} x)) \lor (\exists x\, :\, V,\; \exists y\, :\, V,\; \neg (x = y) \land N_{(\operatorname{spanning\,graph}(M) \setminus \operatorname{spanning\,graph}(M^{'})) \cup (\operatorname{spanning\,graph}(M^{'}) \setminus \operatorname{spanning\,graph}(M))}(v) = \{x\} \cup \{y\})\).
Rewriting with \paperlink{premise:696d706c6965735f636f6e6772}{Prerequisite~P\ref{premise:696d706c6965735f636f6e6772}}, \paperlink{premise:5365742e6e636172645f65715f74776f}{Prerequisite~P\ref{premise:5365742e6e636172645f65715f74776f}}, \paperlink{premise:696d705f6966665f6e6f745f6f72}{Prerequisite~P\ref{premise:696d705f6966665f6e6f745f6f72}}, \paperlink{premise:5365742e6e6f745f6e6f6e656d7074795f6966665f65715f656d707479}{Prerequisite~P\ref{premise:5365742e6e6f745f6e6f6e656d7074795f6966665f65715f656d707479}}, \paperlink{premise:5365742e65715f656d7074795f6966665f666f72616c6c5f6e6f745f6d656d}{Prerequisite~P\ref{premise:5365742e65715f656d7074795f6966665f666f72616c6c5f6e6f745f6d656d}}, \paperlink{premise:666f72616c6c5f636f6e6772}{Prerequisite~P\ref{premise:666f72616c6c5f636f6e6772}} and the identity established above in the preceding result gives \(N_{\operatorname{spanning\,graph}(M) \mathbin{\triangle} \operatorname{spanning\,graph}(M^{'})}(v)\ne\varnothing \implies \operatorname{ncard}(N_{\operatorname{spanning\,graph}(M) \mathbin{\triangle} \operatorname{spanning\,graph}(M^{'})}(v)) = 2\).
\end{proofstep}
Using the existential hypothesis, we obtain \(N_{\operatorname{spanning\,graph}(M) \mathbin{\triangle} \operatorname{spanning\,graph}(M^{'})}(v)\ne\varnothing \implies \operatorname{ncard}(N_{\operatorname{spanning\,graph}(M) \mathbin{\triangle} \operatorname{spanning\,graph}(M^{'})}(v)) = 2\).
\end{proofstep}
Using the existential hypothesis, we obtain \(N_{\operatorname{spanning\,graph}(M) \mathbin{\triangle} \operatorname{spanning\,graph}(M^{'})}(v)\ne\varnothing \implies \operatorname{ncard}(N_{\operatorname{spanning\,graph}(M) \mathbin{\triangle} \operatorname{spanning\,graph}(M^{'})}(v)) = 2\).
\end{proof}

\begin{lemma}[{Snd}]\label{theorem:37}
Let \(V\) be a type. Let \(G\) be a simple graph on \(V\). Assume \(V\) is finite. Let \(v\), \(w\), and \(w^{'}\) be elements of \(V\). Assume \(\text{every nonisolated vertex of }G\text{ has degree two}\). Let \(p\) be a walk in \(G\) from \(v\) to \(w\). Assume \(w \ne w^{'}\). Assume \(p\text{ visits }w^{'}\). Assume \(p\text{ has no repeated vertices}\). Assume \(v \sim_{G} w^{'}\). Then \(\operatorname{Simple\,Graph\mathord{:}Walk\mathord{:}snd}\left(p\right) = w^{'}\).
\end{lemma}

\begin{proof}
Use the hypothesis \(p\text{ has no repeated vertices}\).
Use the hypothesis \(p\text{ visits }w^{'}\).
Rewriting with \paperlink{premise:53696d706c6547726170682e57616c6b2e6d656d5f737570706f72745f6966665f6578697374735f67657456657274}{Prerequisite~P\ref{premise:53696d706c6547726170682e57616c6b2e6d656d5f737570706f72745f6966665f6578697374735f67657456657274}} in the preceding result gives \(\exists n\, :\, \mathbb{N},\; \operatorname{get\,Vert}(p,\allowbreak{} n) = w^{'} \land n \le \operatorname{length}(p)\).
Consider the following case.
\begin{proofstep}
Fix \(n\) in \(\mathbb{N}\).
Assume \(\operatorname{get\,Vert}\left(p, n\right) = w^{'} \land n \le \operatorname{length}(p)\).
Use the hypothesis \(\operatorname{get\,Vert}\left(p, n\right) = w^{'} \land n \le \operatorname{length}(p)\).
Consider the following case.
\begin{proofstep}
Assume \(\operatorname{get\,Vert}\left(p, n\right) = w^{'}\).
Assume \(n \le \operatorname{length}(p)\).
For the required implication, argue as follows.
\begin{proofstep}
Assume \(w \ne \operatorname{get\,Vert}\left(p, n\right)\).
Assume \(v \sim_{G} \operatorname{get\,Vert}\left(p, n\right)\).
Distinguish whether \(n = 0 \lor n = \operatorname{length}(p)\) holds.
Consider the following case.
\begin{proofstep}
Assume \(n = 0 \lor n = \operatorname{length}(p)\).
Use the hypothesis \(n = 0 \lor n = \operatorname{length}(p)\).
Consider the following case.
\begin{proofstep}
Assume \(n = 0\).
For the required implication, argue as follows.
\begin{proofstep}
Assume \(0 \le \operatorname{length}(p)\).
Assume \(\neg \left(w = \operatorname{get\,Vert}\left(p, 0\right)\right)\).
Assume \(v \sim_{G} \operatorname{get\,Vert}\left(p, 0\right)\).
Use the hypothesis \(v \sim_{G} \operatorname{get\,Vert}\left(p, 0\right)\).
Rewriting with \paperlink{premise:53696d706c6547726170682e57616c6b2e676574566572745f7a65726f}{Prerequisite~P\ref{premise:53696d706c6547726170682e57616c6b2e676574566572745f7a65726f}}, \paperlink{premise:65715f66616c7365}{Prerequisite~P\ref{premise:65715f66616c7365}} and \paperlink{premise:53696d706c6547726170682e69727265666c}{Prerequisite~P\ref{premise:53696d706c6547726170682e69727265666c}} in the preceding result gives \(\mathrm{False}\).
The contradiction rules out this case. Hence \(v \sim_{\operatorname{Simple\,Graph\mathord{:}Walk\mathord{:}to\,Subgraph}(p)} \operatorname{get\,Vert}(p,\allowbreak{} 0)\).
\end{proofstep}
Use the hypothesis \(n = 0\).
Reversing the equality, we obtain \(0 = n\).
Use the hypothesis \(n \le \operatorname{length}(p)\).
Use the hypothesis \(w \ne \operatorname{get\,Vert}\left(p, n\right)\).
By the preceding argument, we obtain \(\neg \left(w = \operatorname{get\,Vert}\left(p, n\right)\right)\).
Use the hypothesis \(v \sim_{G} \operatorname{get\,Vert}\left(p, n\right)\).
Substituting equals, we obtain \(v \sim_{\operatorname{Simple\,Graph\mathord{:}Walk\mathord{:}to\,Subgraph}(p)} \operatorname{get\,Vert}(p,\allowbreak{} n)\).
\end{proofstep}
Consider the following case.
\begin{proofstep}
Assume \(n = \operatorname{length}(p)\).
For the required implication, argue as follows.
\begin{proofstep}
Assume \(\operatorname{length}(p) \le \operatorname{length}(p)\).
Assume \(\neg \left(w = \operatorname{get\,Vert}\left(p, \operatorname{length}(p)\right)\right)\).
Assume \(v \sim_{G} \operatorname{get\,Vert}\left(p, \operatorname{length}(p)\right)\).
Use the hypothesis \(\neg \left(w = \operatorname{get\,Vert}\left(p, \operatorname{length}(p)\right)\right)\).
Rewriting with \paperlink{premise:53696d706c6547726170682e57616c6b2e676574566572745f6c656e677468}{Prerequisite~P\ref{premise:53696d706c6547726170682e57616c6b2e676574566572745f6c656e677468}} and \paperlink{premise:6e6f745f747275655f65715f66616c7365}{Prerequisite~P\ref{premise:6e6f745f747275655f65715f66616c7365}} in the preceding result gives \(\mathrm{False}\).
The contradiction rules out this case. Hence \(v \sim_{\operatorname{Simple\,Graph\mathord{:}Walk\mathord{:}to\,Subgraph}(p)} \operatorname{get\,Vert}(p,\allowbreak{} \operatorname{length}(p))\).
\end{proofstep}
Use the hypothesis \(n = \operatorname{length}(p)\).
Reversing the equality, we obtain \(\operatorname{length}(p) = n\).
Use the hypothesis \(n \le \operatorname{length}(p)\).
Use the hypothesis \(w \ne \operatorname{get\,Vert}\left(p, n\right)\).
By the preceding argument, we obtain \(\neg \left(w = \operatorname{get\,Vert}\left(p, n\right)\right)\).
Use the hypothesis \(v \sim_{G} \operatorname{get\,Vert}\left(p, n\right)\).
Substituting equals, we obtain \(v \sim_{\operatorname{Simple\,Graph\mathord{:}Walk\mathord{:}to\,Subgraph}(p)} \operatorname{get\,Vert}(p,\allowbreak{} n)\).
\end{proofstep}
The two alternatives exhaust the cases. Hence \(v \sim_{\operatorname{Simple\,Graph\mathord{:}Walk\mathord{:}to\,Subgraph}(p)} \operatorname{get\,Vert}(p,\allowbreak{} n)\).
\end{proofstep}
Consider the following case.
\begin{proofstep}
Assume \(\neg \left(n = 0 \lor n = \operatorname{length}(p)\right)\).
Set \(e\) equal to \(\operatorname{of\,Nonempty}\).
The substitution uses the following argument.
\begin{proofstep}
Assume \(n = 0\).
Use the hypothesis \(\neg \left(n = 0 \lor n = \operatorname{length}(p)\right)\).
By \paperlink{premise:4c65616e2e4f6d6567612e616e645f6e6f745f6e6f745f6f665f6e6f745f6f72}{Prerequisite~P\ref{premise:4c65616e2e4f6d6567612e616e645f6e6f745f6e6f745f6f665f6e6f745f6f72}}, we obtain \(\neg \left(n = 0\right) \land \neg \left(n = \operatorname{length}(p)\right)\).
Taking the second conjunct, we obtain \(\neg \left(n = \operatorname{length}(p)\right)\).
By \paperlink{premise:4e61742e6c745f6f725f67745f6f665f6e65}{Prerequisite~P\ref{premise:4e61742e6c745f6f725f67745f6f665f6e65}}, we obtain \(n < \operatorname{length}(p) \lor n > \operatorname{length}(p)\).
Consider the following case.
\begin{proofstep}
Assume \(n < \operatorname{length}(p)\).
Use the hypothesis \(\neg \left(n = 0 \lor n = \operatorname{length}(p)\right)\).
By \paperlink{premise:4c65616e2e4f6d6567612e616e645f6e6f745f6e6f745f6f665f6e6f745f6f72}{Prerequisite~P\ref{premise:4c65616e2e4f6d6567612e616e645f6e6f745f6e6f745f6f665f6e6f745f6f72}}, we obtain \(\neg \left(n = 0\right) \land \neg \left(n = \operatorname{length}(p)\right)\).
Taking the first conjunct, we obtain \(\neg \left(n = 0\right)\).
By \paperlink{premise:4e61742e6c745f6f725f67745f6f665f6e65}{Prerequisite~P\ref{premise:4e61742e6c745f6f725f67745f6f665f6e65}}, we obtain \(n < 0 \lor n > 0\).
Consider the following case.
\begin{proofstep}
Assume \(n < 0\).
By definition, we obtain \(\mathrm{true} = \mathrm{true}\).
By the preceding argument, we obtain \(\operatorname{decide}(\operatorname{is\,Impossible}(\operatorname{Lean\mathord{:}Omega\mathord{:}Constraint\mathord{:}mk}(\operatorname{Option\mathord{:}some}(1),\allowbreak{} \operatorname{Option\mathord{:}some}(0))) = \mathrm{true}) = \mathrm{true}\).
By \paperlink{premise:6f665f6465636964655f65715f74727565}{Prerequisite~P\ref{premise:6f665f6465636964655f65715f74727565}}, we obtain \(\operatorname{is\,Impossible}(\operatorname{Lean\mathord{:}Omega\mathord{:}Constraint\mathord{:}mk}(\operatorname{Option\mathord{:}some}(1),\allowbreak{} \operatorname{Option\mathord{:}some}(0))) = \mathrm{true}\).
Using \paperlink{premise:4c65616e2e4f6d6567612e4c696e656172436f6d626f2e636f6f7264696e6174655f6576616c5f31}{Prerequisite~P\ref{premise:4c65616e2e4f6d6567612e4c696e656172436f6d626f2e636f6f7264696e6174655f6576616c5f31}} in reverse gives \(n = \operatorname{eval}(\operatorname{coordinate}(1),\allowbreak{} \operatorname{of\,List}(\langle \operatorname{length}(p)\rangle \mathbin{\Vert} \langle n\rangle))\).
By definition, we obtain \(0 = 0\).
By the preceding argument, we obtain \(0 = \operatorname{eval}(\operatorname{Lean\mathord{:}Omega\mathord{:}Linear\,Combo\mathord{:}mk}(0,\allowbreak{} \langle\,\rangle),\allowbreak{} \operatorname{of\,List}(\langle \operatorname{length}(p)\rangle \mathbin{\Vert} \langle n\rangle))\).
By \paperlink{premise:4c65616e2e4f6d6567612e496e742e7375625f636f6e6772}{Prerequisite~P\ref{premise:4c65616e2e4f6d6567612e496e742e7375625f636f6e6772}}, we obtain \(n - 0 = \operatorname{eval}(\operatorname{coordinate}(1),\allowbreak{} \operatorname{of\,List}(\langle \operatorname{length}(p)\rangle \mathbin{\Vert} \langle n\rangle)) - \operatorname{eval}(\operatorname{Lean\mathord{:}Omega\mathord{:}Linear\,Combo\mathord{:}mk}(0,\allowbreak{} \langle\,\rangle),\allowbreak{} \operatorname{of\,List}(\langle \operatorname{length}(p)\rangle \mathbin{\Vert} \langle n\rangle))\).
Using \paperlink{premise:4c65616e2e4f6d6567612e4c696e656172436f6d626f2e7375625f6576616c}{Prerequisite~P\ref{premise:4c65616e2e4f6d6567612e4c696e656172436f6d626f2e7375625f6576616c}} with \(l_{1} = \operatorname{Lean\mathord{:}Omega\mathord{:}Linear\,Combo\mathord{:}mk}(0,\allowbreak{} \langle 0\rangle \mathbin{\Vert} \langle 1\rangle)\) and \(l_{2} = \operatorname{Lean\mathord{:}Omega\mathord{:}Linear\,Combo\mathord{:}mk}(0,\allowbreak{} \langle\,\rangle)\) and \(\operatorname{of\,List}\left(\langle \operatorname{length}(p)\rangle \mathbin{\Vert} \langle n\rangle\right)\) in place of its parameter \(v\) in reverse gives \(\operatorname{eval}(\operatorname{Lean\mathord{:}Omega\mathord{:}Linear\,Combo\mathord{:}mk}(0,\allowbreak{} \langle 0\rangle \mathbin{\Vert} \langle 1\rangle),\allowbreak{} \operatorname{of\,List}(\langle \operatorname{length}(p)\rangle \mathbin{\Vert} \langle n\rangle)) - \operatorname{eval}(\operatorname{Lean\mathord{:}Omega\mathord{:}Linear\,Combo\mathord{:}mk}(0,\allowbreak{} \langle\,\rangle),\allowbreak{} \operatorname{of\,List}(\langle \operatorname{length}(p)\rangle \mathbin{\Vert} \langle n\rangle)) = \operatorname{eval}(\operatorname{Lean\mathord{:}Omega\mathord{:}Linear\,Combo\mathord{:}mk}(0,\allowbreak{} \langle 0\rangle \mathbin{\Vert} \langle 1\rangle) - \operatorname{Lean\mathord{:}Omega\mathord{:}Linear\,Combo\mathord{:}mk}(0,\allowbreak{} \langle\,\rangle),\allowbreak{} \operatorname{of\,List}(\langle \operatorname{length}(p)\rangle \mathbin{\Vert} \langle n\rangle))\).
Chaining the equalities, we obtain \(n - 0 = \operatorname{eval}(\operatorname{Lean\mathord{:}Omega\mathord{:}Linear\,Combo\mathord{:}mk}(0,\allowbreak{} \langle 0\rangle \mathbin{\Vert} \langle 1\rangle) - \operatorname{Lean\mathord{:}Omega\mathord{:}Linear\,Combo\mathord{:}mk}(0,\allowbreak{} \langle\,\rangle),\allowbreak{} \operatorname{of\,List}(\langle \operatorname{length}(p)\rangle \mathbin{\Vert} \langle n\rangle))\).
Reversing the equality, we obtain \(\operatorname{eval}(\operatorname{Lean\mathord{:}Omega\mathord{:}Linear\,Combo\mathord{:}mk}(0,\allowbreak{} \langle 0\rangle \mathbin{\Vert} \langle 1\rangle) - \operatorname{Lean\mathord{:}Omega\mathord{:}Linear\,Combo\mathord{:}mk}(0,\allowbreak{} \langle\,\rangle),\allowbreak{} \operatorname{of\,List}(\langle \operatorname{length}(p)\rangle \mathbin{\Vert} \langle n\rangle)) = n - 0\).
Use the hypothesis \(n = 0\).
By \paperlink{premise:4c65616e2e4f6d6567612e496e742e6f664e61745f636f6e6772}{Prerequisite~P\ref{premise:4c65616e2e4f6d6567612e496e742e6f664e61745f636f6e6772}}, we obtain \(n = 0\).
By \paperlink{premise:496e742e7375625f65715f7a65726f5f6f665f6571}{Prerequisite~P\ref{premise:496e742e7375625f65715f7a65726f5f6f665f6571}}, we obtain \(n - 0 = 0\).
Chaining the equalities, we obtain \(\operatorname{eval}(\operatorname{Lean\mathord{:}Omega\mathord{:}Linear\,Combo\mathord{:}mk}(0,\allowbreak{} \langle 0\rangle \mathbin{\Vert} \langle 1\rangle) - \operatorname{Lean\mathord{:}Omega\mathord{:}Linear\,Combo\mathord{:}mk}(0,\allowbreak{} \langle\,\rangle),\allowbreak{} \operatorname{of\,List}(\langle \operatorname{length}(p)\rangle \mathbin{\Vert} \langle n\rangle)) = 0\).
By \paperlink{premise:4c65616e2e4f6d6567612e436f6e73747261696e742e616464457175616c6974795f736174}{Prerequisite~P\ref{premise:4c65616e2e4f6d6567612e436f6e73747261696e742e616464457175616c6974795f736174}}, we obtain \(\operatorname{sat^{\prime}}(\operatorname{Lean\mathord{:}Omega\mathord{:}Constraint\mathord{:}mk}(\operatorname{Option\mathord{:}some}(-0),\allowbreak{} \operatorname{Option\mathord{:}some}(-0)),\allowbreak{} \langle 0\rangle \mathbin{\Vert} \langle 1\rangle,\allowbreak{} \operatorname{of\,List}(\langle \operatorname{length}(p)\rangle \mathbin{\Vert} \langle n\rangle)) = \mathrm{true}\).
Use the hypothesis \(n < 0\).
By \paperlink{premise:4c65616e2e4f6d6567612e496e742e6f664e61745f6c745f6f665f6c74}{Prerequisite~P\ref{premise:4c65616e2e4f6d6567612e496e742e6f664e61745f6c745f6f665f6c74}}, we obtain \(n < 0\).
By \paperlink{premise:496e742e6164645f6f6e655f6c655f6f665f6c74}{Prerequisite~P\ref{premise:496e742e6164645f6f6e655f6c655f6f665f6c74}}, we obtain \(n + 1 \le 0\).
By \paperlink{premise:496e742e7375625f6e6f6e6e65675f6f665f6c65}{Prerequisite~P\ref{premise:496e742e7375625f6e6f6e6e65675f6f665f6c65}}, we obtain \(0 \le 0 - \left(n + 1\right)\).
By definition, we obtain \(0 = 0\).
By the preceding argument, we obtain \(0 = \operatorname{eval}(\operatorname{Lean\mathord{:}Omega\mathord{:}Linear\,Combo\mathord{:}mk}(0,\allowbreak{} \langle\,\rangle),\allowbreak{} \operatorname{of\,List}(\langle \operatorname{length}(p)\rangle \mathbin{\Vert} \langle n\rangle))\).
Using \paperlink{premise:4c65616e2e4f6d6567612e4c696e656172436f6d626f2e636f6f7264696e6174655f6576616c5f31}{Prerequisite~P\ref{premise:4c65616e2e4f6d6567612e4c696e656172436f6d626f2e636f6f7264696e6174655f6576616c5f31}} in reverse gives \(n = \operatorname{eval}(\operatorname{coordinate}(1),\allowbreak{} \operatorname{of\,List}(\langle \operatorname{length}(p)\rangle \mathbin{\Vert} \langle n\rangle))\).
By definition, we obtain \(1 = 1\).
By the preceding argument, we obtain \(1 = \operatorname{eval}(\operatorname{Lean\mathord{:}Omega\mathord{:}Linear\,Combo\mathord{:}mk}(1,\allowbreak{} \langle\,\rangle),\allowbreak{} \operatorname{of\,List}(\langle \operatorname{length}(p)\rangle \mathbin{\Vert} \langle n\rangle))\).
By \paperlink{premise:4c65616e2e4f6d6567612e496e742e6164645f636f6e6772}{Prerequisite~P\ref{premise:4c65616e2e4f6d6567612e496e742e6164645f636f6e6772}}, we obtain \(n + 1 = \operatorname{eval}(\operatorname{coordinate}(1),\allowbreak{} \operatorname{of\,List}(\langle \operatorname{length}(p)\rangle \mathbin{\Vert} \langle n\rangle)) + \operatorname{eval}(\operatorname{Lean\mathord{:}Omega\mathord{:}Linear\,Combo\mathord{:}mk}(1,\allowbreak{} \langle\,\rangle),\allowbreak{} \operatorname{of\,List}(\langle \operatorname{length}(p)\rangle \mathbin{\Vert} \langle n\rangle))\).
Using \paperlink{premise:4c65616e2e4f6d6567612e4c696e656172436f6d626f2e6164645f6576616c}{Prerequisite~P\ref{premise:4c65616e2e4f6d6567612e4c696e656172436f6d626f2e6164645f6576616c}} with \(l_{1} = \operatorname{Lean\mathord{:}Omega\mathord{:}Linear\,Combo\mathord{:}mk}(0,\allowbreak{} \langle 0\rangle \mathbin{\Vert} \langle 1\rangle)\) and \(l_{2} = \operatorname{Lean\mathord{:}Omega\mathord{:}Linear\,Combo\mathord{:}mk}(1,\allowbreak{} \langle\,\rangle)\) and \(\operatorname{of\,List}\left(\langle \operatorname{length}(p)\rangle \mathbin{\Vert} \langle n\rangle\right)\) in place of its parameter \(v\) in reverse gives \(\operatorname{eval}(\operatorname{Lean\mathord{:}Omega\mathord{:}Linear\,Combo\mathord{:}mk}(0,\allowbreak{} \langle 0\rangle \mathbin{\Vert} \langle 1\rangle),\allowbreak{} \operatorname{of\,List}(\langle \operatorname{length}(p)\rangle \mathbin{\Vert} \langle n\rangle)) + \operatorname{eval}(\operatorname{Lean\mathord{:}Omega\mathord{:}Linear\,Combo\mathord{:}mk}(1,\allowbreak{} \langle\,\rangle),\allowbreak{} \operatorname{of\,List}(\langle \operatorname{length}(p)\rangle \mathbin{\Vert} \langle n\rangle)) = \operatorname{eval}(\operatorname{Lean\mathord{:}Omega\mathord{:}Linear\,Combo\mathord{:}mk}(0,\allowbreak{} \langle 0\rangle \mathbin{\Vert} \langle 1\rangle) + \operatorname{Lean\mathord{:}Omega\mathord{:}Linear\,Combo\mathord{:}mk}(1,\allowbreak{} \langle\,\rangle),\allowbreak{} \operatorname{of\,List}(\langle \operatorname{length}(p)\rangle \mathbin{\Vert} \langle n\rangle))\).
Chaining the equalities, we obtain \(n + 1 = \operatorname{eval}(\operatorname{Lean\mathord{:}Omega\mathord{:}Linear\,Combo\mathord{:}mk}(0,\allowbreak{} \langle 0\rangle \mathbin{\Vert} \langle 1\rangle) + \operatorname{Lean\mathord{:}Omega\mathord{:}Linear\,Combo\mathord{:}mk}(1,\allowbreak{} \langle\,\rangle),\allowbreak{} \operatorname{of\,List}(\langle \operatorname{length}(p)\rangle \mathbin{\Vert} \langle n\rangle))\).
By \paperlink{premise:4c65616e2e4f6d6567612e496e742e7375625f636f6e6772}{Prerequisite~P\ref{premise:4c65616e2e4f6d6567612e496e742e7375625f636f6e6772}}, we obtain \(0 - (n + 1) = \operatorname{eval}(\operatorname{Lean\mathord{:}Omega\mathord{:}Linear\,Combo\mathord{:}mk}(0,\allowbreak{} \langle\,\rangle),\allowbreak{} \operatorname{of\,List}(\langle \operatorname{length}(p)\rangle \mathbin{\Vert} \langle n\rangle)) - \operatorname{eval}(\operatorname{Lean\mathord{:}Omega\mathord{:}Linear\,Combo\mathord{:}mk}(0,\allowbreak{} \langle 0\rangle \mathbin{\Vert} \langle 1\rangle) + \operatorname{Lean\mathord{:}Omega\mathord{:}Linear\,Combo\mathord{:}mk}(1,\allowbreak{} \langle\,\rangle),\allowbreak{} \operatorname{of\,List}(\langle \operatorname{length}(p)\rangle \mathbin{\Vert} \langle n\rangle))\).
Using \paperlink{premise:4c65616e2e4f6d6567612e4c696e656172436f6d626f2e7375625f6576616c}{Prerequisite~P\ref{premise:4c65616e2e4f6d6567612e4c696e656172436f6d626f2e7375625f6576616c}} with \(l_{1} = \operatorname{Lean\mathord{:}Omega\mathord{:}Linear\,Combo\mathord{:}mk}(0,\allowbreak{} \langle\,\rangle)\) and \(l_{2} = \operatorname{Lean\mathord{:}Omega\mathord{:}Linear\,Combo\mathord{:}mk}(1,\allowbreak{} \langle 0\rangle \mathbin{\Vert} \langle 1\rangle)\) and \(\operatorname{of\,List}\left(\langle \operatorname{length}(p)\rangle \mathbin{\Vert} \langle n\rangle\right)\) in place of its parameter \(v\) in reverse gives \(\operatorname{eval}(\operatorname{Lean\mathord{:}Omega\mathord{:}Linear\,Combo\mathord{:}mk}(0,\allowbreak{} \langle\,\rangle),\allowbreak{} \operatorname{of\,List}(\langle \operatorname{length}(p)\rangle \mathbin{\Vert} \langle n\rangle)) - \operatorname{eval}(\operatorname{Lean\mathord{:}Omega\mathord{:}Linear\,Combo\mathord{:}mk}(1,\allowbreak{} \langle 0\rangle \mathbin{\Vert} \langle 1\rangle),\allowbreak{} \operatorname{of\,List}(\langle \operatorname{length}(p)\rangle \mathbin{\Vert} \langle n\rangle)) = \operatorname{eval}(\operatorname{Lean\mathord{:}Omega\mathord{:}Linear\,Combo\mathord{:}mk}(0,\allowbreak{} \langle\,\rangle) - \operatorname{Lean\mathord{:}Omega\mathord{:}Linear\,Combo\mathord{:}mk}(1,\allowbreak{} \langle 0\rangle \mathbin{\Vert} \langle 1\rangle),\allowbreak{} \operatorname{of\,List}(\langle \operatorname{length}(p)\rangle \mathbin{\Vert} \langle n\rangle))\).
Chaining the equalities, we obtain \(0 - (n + 1) = \operatorname{eval}(\operatorname{Lean\mathord{:}Omega\mathord{:}Linear\,Combo\mathord{:}mk}(0,\allowbreak{} \langle\,\rangle) - \operatorname{Lean\mathord{:}Omega\mathord{:}Linear\,Combo\mathord{:}mk}(1,\allowbreak{} \langle 0\rangle \mathbin{\Vert} \langle 1\rangle),\allowbreak{} \operatorname{of\,List}(\langle \operatorname{length}(p)\rangle \mathbin{\Vert} \langle n\rangle))\).
By \paperlink{premise:6c655f6f665f6c655f6f665f6571}{Prerequisite~P\ref{premise:6c655f6f665f6c655f6f665f6571}} with \(\alpha = \mathbb{Z}\), we obtain \(0 \le \operatorname{eval}(\operatorname{Lean\mathord{:}Omega\mathord{:}Linear\,Combo\mathord{:}mk}(0,\allowbreak{} \langle\,\rangle) - \operatorname{Lean\mathord{:}Omega\mathord{:}Linear\,Combo\mathord{:}mk}(1,\allowbreak{} \langle 0\rangle \mathbin{\Vert} \langle 1\rangle),\allowbreak{} \operatorname{of\,List}(\langle \operatorname{length}(p)\rangle \mathbin{\Vert} \langle n\rangle))\).
By \paperlink{premise:4c65616e2e4f6d6567612e436f6e73747261696e742e616464496e657175616c6974795f736174}{Prerequisite~P\ref{premise:4c65616e2e4f6d6567612e436f6e73747261696e742e616464496e657175616c6974795f736174}}, we obtain \(\operatorname{sat^{\prime}}(\operatorname{Lean\mathord{:}Omega\mathord{:}Constraint\mathord{:}mk}(\operatorname{Option\mathord{:}some}(-(-1)),\allowbreak{} \operatorname{none}),\allowbreak{} \langle 0\rangle \mathbin{\Vert} \langle -1\rangle,\allowbreak{} \operatorname{of\,List}(\langle \operatorname{length}(p)\rangle \mathbin{\Vert} \langle n\rangle)) = \mathrm{true}\).
By \paperlink{premise:4c65616e2e4f6d6567612e636f6d626f5f73617427}{Prerequisite~P\ref{premise:4c65616e2e4f6d6567612e636f6d626f5f73617427}} with \(s = \operatorname{Lean\mathord{:}Omega\mathord{:}Constraint\mathord{:}mk}(\operatorname{Option\mathord{:}some}(0),\allowbreak{} \operatorname{Option\mathord{:}some}(0))\) and \(t = \operatorname{Lean\mathord{:}Omega\mathord{:}Constraint\mathord{:}mk}(\operatorname{Option\mathord{:}some}(1),\allowbreak{} \operatorname{none})\) and \(1\) in place of its parameter \(a\) and \(x = \langle 0\rangle \mathbin{\Vert} \langle 1\rangle\) and \(b = 1\) and \(y = \langle 0\rangle \mathbin{\Vert} \langle -1\rangle\) and \(\operatorname{of\,List}\left(\langle \operatorname{length}(p)\rangle \mathbin{\Vert} \langle n\rangle\right)\) in place of its parameter \(v\), we obtain \(\operatorname{sat^{\prime}}(\operatorname{Lean\mathord{:}Omega\mathord{:}Constraint\mathord{:}combo}(1,\allowbreak{} \operatorname{Lean\mathord{:}Omega\mathord{:}Constraint\mathord{:}mk}(\operatorname{Option\mathord{:}some}(0),\allowbreak{} \operatorname{Option\mathord{:}some}(0)),\allowbreak{} 1,\allowbreak{} \operatorname{Lean\mathord{:}Omega\mathord{:}Constraint\mathord{:}mk}(\operatorname{Option\mathord{:}some}(1),\allowbreak{} \operatorname{none})),\allowbreak{} \operatorname{Lean\mathord{:}Omega\mathord{:}Coeffs\mathord{:}combo}(1,\allowbreak{} \langle 0\rangle \mathbin{\Vert} \langle 1\rangle,\allowbreak{} 1,\allowbreak{} \langle 0\rangle \mathbin{\Vert} \langle -1\rangle),\allowbreak{} \operatorname{of\,List}(\langle \operatorname{length}(p)\rangle \mathbin{\Vert} \langle n\rangle)) = \mathrm{true}\).
By \paperlink{premise:4c65616e2e4f6d6567612e746964795f736174}{Prerequisite~P\ref{premise:4c65616e2e4f6d6567612e746964795f736174}}, we obtain \(\operatorname{sat^{\prime}}(\operatorname{tidy\,Constraint}(\operatorname{Lean\mathord{:}Omega\mathord{:}Constraint\mathord{:}mk}(\operatorname{Option\mathord{:}some}(1),\allowbreak{} \operatorname{none}),\allowbreak{} \langle 0\rangle \mathbin{\Vert} \langle 0\rangle),\allowbreak{} \operatorname{tidy\,Coeffs}(\operatorname{Lean\mathord{:}Omega\mathord{:}Constraint\mathord{:}mk}(\operatorname{Option\mathord{:}some}(1),\allowbreak{} \operatorname{none}),\allowbreak{} \langle 0\rangle \mathbin{\Vert} \langle 0\rangle),\allowbreak{} \operatorname{of\,List}(\langle \operatorname{length}(p)\rangle \mathbin{\Vert} \langle n\rangle)) = \mathrm{true}\).
By \paperlink{premise:4c65616e2e4f6d6567612e436f6e73747261696e742e6e6f745f736174275f6f665f6973496d706f737369626c65}{Prerequisite~P\ref{premise:4c65616e2e4f6d6567612e436f6e73747261696e742e6e6f745f736174275f6f665f6973496d706f737369626c65}}, we obtain \(\mathrm{False}\).
\end{proofstep}
Consider the following case.
\begin{proofstep}
Assume \(n > 0\).
By definition, we obtain \(\mathrm{true} = \mathrm{true}\).
By the preceding argument, we obtain \(\operatorname{decide}(\operatorname{is\,Impossible}(\operatorname{Lean\mathord{:}Omega\mathord{:}Constraint\mathord{:}mk}(\operatorname{Option\mathord{:}some}(1),\allowbreak{} \operatorname{Option\mathord{:}some}(0))) = \mathrm{true}) = \mathrm{true}\).
By \paperlink{premise:6f665f6465636964655f65715f74727565}{Prerequisite~P\ref{premise:6f665f6465636964655f65715f74727565}}, we obtain \(\operatorname{is\,Impossible}(\operatorname{Lean\mathord{:}Omega\mathord{:}Constraint\mathord{:}mk}(\operatorname{Option\mathord{:}some}(1),\allowbreak{} \operatorname{Option\mathord{:}some}(0))) = \mathrm{true}\).
Using \paperlink{premise:4c65616e2e4f6d6567612e4c696e656172436f6d626f2e636f6f7264696e6174655f6576616c5f31}{Prerequisite~P\ref{premise:4c65616e2e4f6d6567612e4c696e656172436f6d626f2e636f6f7264696e6174655f6576616c5f31}} in reverse gives \(n = \operatorname{eval}(\operatorname{coordinate}(1),\allowbreak{} \operatorname{of\,List}(\langle \operatorname{length}(p)\rangle \mathbin{\Vert} \langle n\rangle))\).
By definition, we obtain \(0 = 0\).
By the preceding argument, we obtain \(0 = \operatorname{eval}(\operatorname{Lean\mathord{:}Omega\mathord{:}Linear\,Combo\mathord{:}mk}(0,\allowbreak{} \langle\,\rangle),\allowbreak{} \operatorname{of\,List}(\langle \operatorname{length}(p)\rangle \mathbin{\Vert} \langle n\rangle))\).
By \paperlink{premise:4c65616e2e4f6d6567612e496e742e7375625f636f6e6772}{Prerequisite~P\ref{premise:4c65616e2e4f6d6567612e496e742e7375625f636f6e6772}}, we obtain \(n - 0 = \operatorname{eval}(\operatorname{coordinate}(1),\allowbreak{} \operatorname{of\,List}(\langle \operatorname{length}(p)\rangle \mathbin{\Vert} \langle n\rangle)) - \operatorname{eval}(\operatorname{Lean\mathord{:}Omega\mathord{:}Linear\,Combo\mathord{:}mk}(0,\allowbreak{} \langle\,\rangle),\allowbreak{} \operatorname{of\,List}(\langle \operatorname{length}(p)\rangle \mathbin{\Vert} \langle n\rangle))\).
Using \paperlink{premise:4c65616e2e4f6d6567612e4c696e656172436f6d626f2e7375625f6576616c}{Prerequisite~P\ref{premise:4c65616e2e4f6d6567612e4c696e656172436f6d626f2e7375625f6576616c}} with \(l_{1} = \operatorname{Lean\mathord{:}Omega\mathord{:}Linear\,Combo\mathord{:}mk}(0,\allowbreak{} \langle 0\rangle \mathbin{\Vert} \langle 1\rangle)\) and \(l_{2} = \operatorname{Lean\mathord{:}Omega\mathord{:}Linear\,Combo\mathord{:}mk}(0,\allowbreak{} \langle\,\rangle)\) and \(\operatorname{of\,List}\left(\langle \operatorname{length}(p)\rangle \mathbin{\Vert} \langle n\rangle\right)\) in place of its parameter \(v\) in reverse gives \(\operatorname{eval}(\operatorname{Lean\mathord{:}Omega\mathord{:}Linear\,Combo\mathord{:}mk}(0,\allowbreak{} \langle 0\rangle \mathbin{\Vert} \langle 1\rangle),\allowbreak{} \operatorname{of\,List}(\langle \operatorname{length}(p)\rangle \mathbin{\Vert} \langle n\rangle)) - \operatorname{eval}(\operatorname{Lean\mathord{:}Omega\mathord{:}Linear\,Combo\mathord{:}mk}(0,\allowbreak{} \langle\,\rangle),\allowbreak{} \operatorname{of\,List}(\langle \operatorname{length}(p)\rangle \mathbin{\Vert} \langle n\rangle)) = \operatorname{eval}(\operatorname{Lean\mathord{:}Omega\mathord{:}Linear\,Combo\mathord{:}mk}(0,\allowbreak{} \langle 0\rangle \mathbin{\Vert} \langle 1\rangle) - \operatorname{Lean\mathord{:}Omega\mathord{:}Linear\,Combo\mathord{:}mk}(0,\allowbreak{} \langle\,\rangle),\allowbreak{} \operatorname{of\,List}(\langle \operatorname{length}(p)\rangle \mathbin{\Vert} \langle n\rangle))\).
Chaining the equalities, we obtain \(n - 0 = \operatorname{eval}(\operatorname{Lean\mathord{:}Omega\mathord{:}Linear\,Combo\mathord{:}mk}(0,\allowbreak{} \langle 0\rangle \mathbin{\Vert} \langle 1\rangle) - \operatorname{Lean\mathord{:}Omega\mathord{:}Linear\,Combo\mathord{:}mk}(0,\allowbreak{} \langle\,\rangle),\allowbreak{} \operatorname{of\,List}(\langle \operatorname{length}(p)\rangle \mathbin{\Vert} \langle n\rangle))\).
Reversing the equality, we obtain \(\operatorname{eval}(\operatorname{Lean\mathord{:}Omega\mathord{:}Linear\,Combo\mathord{:}mk}(0,\allowbreak{} \langle 0\rangle \mathbin{\Vert} \langle 1\rangle) - \operatorname{Lean\mathord{:}Omega\mathord{:}Linear\,Combo\mathord{:}mk}(0,\allowbreak{} \langle\,\rangle),\allowbreak{} \operatorname{of\,List}(\langle \operatorname{length}(p)\rangle \mathbin{\Vert} \langle n\rangle)) = n - 0\).
Use the hypothesis \(n = 0\).
By \paperlink{premise:4c65616e2e4f6d6567612e496e742e6f664e61745f636f6e6772}{Prerequisite~P\ref{premise:4c65616e2e4f6d6567612e496e742e6f664e61745f636f6e6772}}, we obtain \(n = 0\).
By \paperlink{premise:496e742e7375625f65715f7a65726f5f6f665f6571}{Prerequisite~P\ref{premise:496e742e7375625f65715f7a65726f5f6f665f6571}}, we obtain \(n - 0 = 0\).
Chaining the equalities, we obtain \(\operatorname{eval}(\operatorname{Lean\mathord{:}Omega\mathord{:}Linear\,Combo\mathord{:}mk}(0,\allowbreak{} \langle 0\rangle \mathbin{\Vert} \langle 1\rangle) - \operatorname{Lean\mathord{:}Omega\mathord{:}Linear\,Combo\mathord{:}mk}(0,\allowbreak{} \langle\,\rangle),\allowbreak{} \operatorname{of\,List}(\langle \operatorname{length}(p)\rangle \mathbin{\Vert} \langle n\rangle)) = 0\).
By \paperlink{premise:4c65616e2e4f6d6567612e436f6e73747261696e742e616464457175616c6974795f736174}{Prerequisite~P\ref{premise:4c65616e2e4f6d6567612e436f6e73747261696e742e616464457175616c6974795f736174}}, we obtain \(\operatorname{sat^{\prime}}(\operatorname{Lean\mathord{:}Omega\mathord{:}Constraint\mathord{:}mk}(\operatorname{Option\mathord{:}some}(-0),\allowbreak{} \operatorname{Option\mathord{:}some}(-0)),\allowbreak{} \langle 0\rangle \mathbin{\Vert} \langle 1\rangle,\allowbreak{} \operatorname{of\,List}(\langle \operatorname{length}(p)\rangle \mathbin{\Vert} \langle n\rangle)) = \mathrm{true}\).
Use the hypothesis \(n > 0\).
By \paperlink{premise:4c65616e2e4f6d6567612e496e742e6f664e61745f6c745f6f665f6c74}{Prerequisite~P\ref{premise:4c65616e2e4f6d6567612e496e742e6f664e61745f6c745f6f665f6c74}}, we obtain \(0 < n\).
By \paperlink{premise:496e742e6164645f6f6e655f6c655f6f665f6c74}{Prerequisite~P\ref{premise:496e742e6164645f6f6e655f6c655f6f665f6c74}}, we obtain \(0 + 1 \le n\).
By \paperlink{premise:496e742e7375625f6e6f6e6e65675f6f665f6c65}{Prerequisite~P\ref{premise:496e742e7375625f6e6f6e6e65675f6f665f6c65}}, we obtain \(0 \le n - \left(0 + 1\right)\).
Using \paperlink{premise:4c65616e2e4f6d6567612e4c696e656172436f6d626f2e636f6f7264696e6174655f6576616c5f31}{Prerequisite~P\ref{premise:4c65616e2e4f6d6567612e4c696e656172436f6d626f2e636f6f7264696e6174655f6576616c5f31}} in reverse gives \(n = \operatorname{eval}(\operatorname{coordinate}(1),\allowbreak{} \operatorname{of\,List}(\langle \operatorname{length}(p)\rangle \mathbin{\Vert} \langle n\rangle))\).
By definition, we obtain \(0 + 1 = 0 + 1\).
By the preceding argument, we obtain \(0 + 1 = \operatorname{eval}(\operatorname{Lean\mathord{:}Omega\mathord{:}Linear\,Combo\mathord{:}mk}(1,\allowbreak{} \langle\,\rangle),\allowbreak{} \operatorname{of\,List}(\langle \operatorname{length}(p)\rangle \mathbin{\Vert} \langle n\rangle))\).
By \paperlink{premise:4c65616e2e4f6d6567612e496e742e7375625f636f6e6772}{Prerequisite~P\ref{premise:4c65616e2e4f6d6567612e496e742e7375625f636f6e6772}}, we obtain \(n - (0 + 1) = \operatorname{eval}(\operatorname{coordinate}(1),\allowbreak{} \operatorname{of\,List}(\langle \operatorname{length}(p)\rangle \mathbin{\Vert} \langle n\rangle)) - \operatorname{eval}(\operatorname{Lean\mathord{:}Omega\mathord{:}Linear\,Combo\mathord{:}mk}(1,\allowbreak{} \langle\,\rangle),\allowbreak{} \operatorname{of\,List}(\langle \operatorname{length}(p)\rangle \mathbin{\Vert} \langle n\rangle))\).
Using \paperlink{premise:4c65616e2e4f6d6567612e4c696e656172436f6d626f2e7375625f6576616c}{Prerequisite~P\ref{premise:4c65616e2e4f6d6567612e4c696e656172436f6d626f2e7375625f6576616c}} with \(l_{1} = \operatorname{Lean\mathord{:}Omega\mathord{:}Linear\,Combo\mathord{:}mk}(0,\allowbreak{} \langle 0\rangle \mathbin{\Vert} \langle 1\rangle)\) and \(l_{2} = \operatorname{Lean\mathord{:}Omega\mathord{:}Linear\,Combo\mathord{:}mk}(1,\allowbreak{} \langle\,\rangle)\) and \(\operatorname{of\,List}\left(\langle \operatorname{length}(p)\rangle \mathbin{\Vert} \langle n\rangle\right)\) in place of its parameter \(v\) in reverse gives \(\operatorname{eval}(\operatorname{Lean\mathord{:}Omega\mathord{:}Linear\,Combo\mathord{:}mk}(0,\allowbreak{} \langle 0\rangle \mathbin{\Vert} \langle 1\rangle),\allowbreak{} \operatorname{of\,List}(\langle \operatorname{length}(p)\rangle \mathbin{\Vert} \langle n\rangle)) - \operatorname{eval}(\operatorname{Lean\mathord{:}Omega\mathord{:}Linear\,Combo\mathord{:}mk}(1,\allowbreak{} \langle\,\rangle),\allowbreak{} \operatorname{of\,List}(\langle \operatorname{length}(p)\rangle \mathbin{\Vert} \langle n\rangle)) = \operatorname{eval}(\operatorname{Lean\mathord{:}Omega\mathord{:}Linear\,Combo\mathord{:}mk}(0,\allowbreak{} \langle 0\rangle \mathbin{\Vert} \langle 1\rangle) - \operatorname{Lean\mathord{:}Omega\mathord{:}Linear\,Combo\mathord{:}mk}(1,\allowbreak{} \langle\,\rangle),\allowbreak{} \operatorname{of\,List}(\langle \operatorname{length}(p)\rangle \mathbin{\Vert} \langle n\rangle))\).
Chaining the equalities, we obtain \(n - (0 + 1) = \operatorname{eval}(\operatorname{Lean\mathord{:}Omega\mathord{:}Linear\,Combo\mathord{:}mk}(0,\allowbreak{} \langle 0\rangle \mathbin{\Vert} \langle 1\rangle) - \operatorname{Lean\mathord{:}Omega\mathord{:}Linear\,Combo\mathord{:}mk}(1,\allowbreak{} \langle\,\rangle),\allowbreak{} \operatorname{of\,List}(\langle \operatorname{length}(p)\rangle \mathbin{\Vert} \langle n\rangle))\).
By \paperlink{premise:6c655f6f665f6c655f6f665f6571}{Prerequisite~P\ref{premise:6c655f6f665f6c655f6f665f6571}} with \(\alpha = \mathbb{Z}\), we obtain \(0 \le \operatorname{eval}(\operatorname{Lean\mathord{:}Omega\mathord{:}Linear\,Combo\mathord{:}mk}(0,\allowbreak{} \langle 0\rangle \mathbin{\Vert} \langle 1\rangle) - \operatorname{Lean\mathord{:}Omega\mathord{:}Linear\,Combo\mathord{:}mk}(1,\allowbreak{} \langle\,\rangle),\allowbreak{} \operatorname{of\,List}(\langle \operatorname{length}(p)\rangle \mathbin{\Vert} \langle n\rangle))\).
By \paperlink{premise:4c65616e2e4f6d6567612e436f6e73747261696e742e616464496e657175616c6974795f736174}{Prerequisite~P\ref{premise:4c65616e2e4f6d6567612e436f6e73747261696e742e616464496e657175616c6974795f736174}}, we obtain \(\operatorname{sat^{\prime}}(\operatorname{Lean\mathord{:}Omega\mathord{:}Constraint\mathord{:}mk}(\operatorname{Option\mathord{:}some}(-(-1)),\allowbreak{} \operatorname{none}),\allowbreak{} \langle 0\rangle \mathbin{\Vert} \langle 1\rangle,\allowbreak{} \operatorname{of\,List}(\langle \operatorname{length}(p)\rangle \mathbin{\Vert} \langle n\rangle)) = \mathrm{true}\).
By \paperlink{premise:4c65616e2e4f6d6567612e636f6d626f5f73617427}{Prerequisite~P\ref{premise:4c65616e2e4f6d6567612e636f6d626f5f73617427}} with \(s = \operatorname{Lean\mathord{:}Omega\mathord{:}Constraint\mathord{:}mk}(\operatorname{Option\mathord{:}some}(0),\allowbreak{} \operatorname{Option\mathord{:}some}(0))\) and \(t = \operatorname{Lean\mathord{:}Omega\mathord{:}Constraint\mathord{:}mk}(\operatorname{Option\mathord{:}some}(1),\allowbreak{} \operatorname{none})\) and \(-1\) in place of its parameter \(a\) and \(x = \langle 0\rangle \mathbin{\Vert} \langle 1\rangle\) and \(b = 1\) and \(y = \langle 0\rangle \mathbin{\Vert} \langle 1\rangle\) and \(\operatorname{of\,List}\left(\langle \operatorname{length}(p)\rangle \mathbin{\Vert} \langle n\rangle\right)\) in place of its parameter \(v\), we obtain \(\operatorname{sat^{\prime}}(\operatorname{Lean\mathord{:}Omega\mathord{:}Constraint\mathord{:}combo}(-1,\allowbreak{} \operatorname{Lean\mathord{:}Omega\mathord{:}Constraint\mathord{:}mk}(\operatorname{Option\mathord{:}some}(0),\allowbreak{} \operatorname{Option\mathord{:}some}(0)),\allowbreak{} 1,\allowbreak{} \operatorname{Lean\mathord{:}Omega\mathord{:}Constraint\mathord{:}mk}(\operatorname{Option\mathord{:}some}(1),\allowbreak{} \operatorname{none})),\allowbreak{} \operatorname{Lean\mathord{:}Omega\mathord{:}Coeffs\mathord{:}combo}(-1,\allowbreak{} \langle 0\rangle \mathbin{\Vert} \langle 1\rangle,\allowbreak{} 1,\allowbreak{} \langle 0\rangle \mathbin{\Vert} \langle 1\rangle),\allowbreak{} \operatorname{of\,List}(\langle \operatorname{length}(p)\rangle \mathbin{\Vert} \langle n\rangle)) = \mathrm{true}\).
By \paperlink{premise:4c65616e2e4f6d6567612e746964795f736174}{Prerequisite~P\ref{premise:4c65616e2e4f6d6567612e746964795f736174}}, we obtain \(\operatorname{sat^{\prime}}(\operatorname{tidy\,Constraint}(\operatorname{Lean\mathord{:}Omega\mathord{:}Constraint\mathord{:}mk}(\operatorname{Option\mathord{:}some}(1),\allowbreak{} \operatorname{none}),\allowbreak{} \langle 0\rangle \mathbin{\Vert} \langle 0\rangle),\allowbreak{} \operatorname{tidy\,Coeffs}(\operatorname{Lean\mathord{:}Omega\mathord{:}Constraint\mathord{:}mk}(\operatorname{Option\mathord{:}some}(1),\allowbreak{} \operatorname{none}),\allowbreak{} \langle 0\rangle \mathbin{\Vert} \langle 0\rangle),\allowbreak{} \operatorname{of\,List}(\langle \operatorname{length}(p)\rangle \mathbin{\Vert} \langle n\rangle)) = \mathrm{true}\).
By \paperlink{premise:4c65616e2e4f6d6567612e436f6e73747261696e742e6e6f745f736174275f6f665f6973496d706f737369626c65}{Prerequisite~P\ref{premise:4c65616e2e4f6d6567612e436f6e73747261696e742e6e6f745f736174275f6f665f6973496d706f737369626c65}}, we obtain \(\mathrm{False}\).
\end{proofstep}
The two alternatives exhaust the cases. Hence \(\mathrm{False}\).
\end{proofstep}
Consider the following case.
\begin{proofstep}
Assume \(n > \operatorname{length}(p)\).
By definition, we obtain \(\mathrm{true} = \mathrm{true}\).
By the preceding argument, we obtain \(\operatorname{decide}(\operatorname{is\,Impossible}(\operatorname{Lean\mathord{:}Omega\mathord{:}Constraint\mathord{:}mk}(\operatorname{Option\mathord{:}some}(0),\allowbreak{} \operatorname{Option\mathord{:}some}(-1))) = \mathrm{true}) = \mathrm{true}\).
By \paperlink{premise:6f665f6465636964655f65715f74727565}{Prerequisite~P\ref{premise:6f665f6465636964655f65715f74727565}}, we obtain \(\operatorname{is\,Impossible}(\operatorname{Lean\mathord{:}Omega\mathord{:}Constraint\mathord{:}mk}(\operatorname{Option\mathord{:}some}(0),\allowbreak{} \operatorname{Option\mathord{:}some}(-1))) = \mathrm{true}\).
Using \paperlink{premise:4c65616e2e4f6d6567612e4c696e656172436f6d626f2e636f6f7264696e6174655f6576616c5f31}{Prerequisite~P\ref{premise:4c65616e2e4f6d6567612e4c696e656172436f6d626f2e636f6f7264696e6174655f6576616c5f31}} in reverse gives \(n = \operatorname{eval}(\operatorname{coordinate}(1),\allowbreak{} \operatorname{of\,List}(\langle \operatorname{length}(p)\rangle \mathbin{\Vert} \langle n\rangle))\).
By definition, we obtain \(0 = 0\).
By the preceding argument, we obtain \(0 = \operatorname{eval}(\operatorname{Lean\mathord{:}Omega\mathord{:}Linear\,Combo\mathord{:}mk}(0,\allowbreak{} \langle\,\rangle),\allowbreak{} \operatorname{of\,List}(\langle \operatorname{length}(p)\rangle \mathbin{\Vert} \langle n\rangle))\).
By \paperlink{premise:4c65616e2e4f6d6567612e496e742e7375625f636f6e6772}{Prerequisite~P\ref{premise:4c65616e2e4f6d6567612e496e742e7375625f636f6e6772}}, we obtain \(n - 0 = \operatorname{eval}(\operatorname{coordinate}(1),\allowbreak{} \operatorname{of\,List}(\langle \operatorname{length}(p)\rangle \mathbin{\Vert} \langle n\rangle)) - \operatorname{eval}(\operatorname{Lean\mathord{:}Omega\mathord{:}Linear\,Combo\mathord{:}mk}(0,\allowbreak{} \langle\,\rangle),\allowbreak{} \operatorname{of\,List}(\langle \operatorname{length}(p)\rangle \mathbin{\Vert} \langle n\rangle))\).
Using \paperlink{premise:4c65616e2e4f6d6567612e4c696e656172436f6d626f2e7375625f6576616c}{Prerequisite~P\ref{premise:4c65616e2e4f6d6567612e4c696e656172436f6d626f2e7375625f6576616c}} with \(l_{1} = \operatorname{Lean\mathord{:}Omega\mathord{:}Linear\,Combo\mathord{:}mk}(0,\allowbreak{} \langle 0\rangle \mathbin{\Vert} \langle 1\rangle)\) and \(l_{2} = \operatorname{Lean\mathord{:}Omega\mathord{:}Linear\,Combo\mathord{:}mk}(0,\allowbreak{} \langle\,\rangle)\) and \(\operatorname{of\,List}\left(\langle \operatorname{length}(p)\rangle \mathbin{\Vert} \langle n\rangle\right)\) in place of its parameter \(v\) in reverse gives \(\operatorname{eval}(\operatorname{Lean\mathord{:}Omega\mathord{:}Linear\,Combo\mathord{:}mk}(0,\allowbreak{} \langle 0\rangle \mathbin{\Vert} \langle 1\rangle),\allowbreak{} \operatorname{of\,List}(\langle \operatorname{length}(p)\rangle \mathbin{\Vert} \langle n\rangle)) - \operatorname{eval}(\operatorname{Lean\mathord{:}Omega\mathord{:}Linear\,Combo\mathord{:}mk}(0,\allowbreak{} \langle\,\rangle),\allowbreak{} \operatorname{of\,List}(\langle \operatorname{length}(p)\rangle \mathbin{\Vert} \langle n\rangle)) = \operatorname{eval}(\operatorname{Lean\mathord{:}Omega\mathord{:}Linear\,Combo\mathord{:}mk}(0,\allowbreak{} \langle 0\rangle \mathbin{\Vert} \langle 1\rangle) - \operatorname{Lean\mathord{:}Omega\mathord{:}Linear\,Combo\mathord{:}mk}(0,\allowbreak{} \langle\,\rangle),\allowbreak{} \operatorname{of\,List}(\langle \operatorname{length}(p)\rangle \mathbin{\Vert} \langle n\rangle))\).
Chaining the equalities, we obtain \(n - 0 = \operatorname{eval}(\operatorname{Lean\mathord{:}Omega\mathord{:}Linear\,Combo\mathord{:}mk}(0,\allowbreak{} \langle 0\rangle \mathbin{\Vert} \langle 1\rangle) - \operatorname{Lean\mathord{:}Omega\mathord{:}Linear\,Combo\mathord{:}mk}(0,\allowbreak{} \langle\,\rangle),\allowbreak{} \operatorname{of\,List}(\langle \operatorname{length}(p)\rangle \mathbin{\Vert} \langle n\rangle))\).
Reversing the equality, we obtain \(\operatorname{eval}(\operatorname{Lean\mathord{:}Omega\mathord{:}Linear\,Combo\mathord{:}mk}(0,\allowbreak{} \langle 0\rangle \mathbin{\Vert} \langle 1\rangle) - \operatorname{Lean\mathord{:}Omega\mathord{:}Linear\,Combo\mathord{:}mk}(0,\allowbreak{} \langle\,\rangle),\allowbreak{} \operatorname{of\,List}(\langle \operatorname{length}(p)\rangle \mathbin{\Vert} \langle n\rangle)) = n - 0\).
Use the hypothesis \(n = 0\).
By \paperlink{premise:4c65616e2e4f6d6567612e496e742e6f664e61745f636f6e6772}{Prerequisite~P\ref{premise:4c65616e2e4f6d6567612e496e742e6f664e61745f636f6e6772}}, we obtain \(n = 0\).
By \paperlink{premise:496e742e7375625f65715f7a65726f5f6f665f6571}{Prerequisite~P\ref{premise:496e742e7375625f65715f7a65726f5f6f665f6571}}, we obtain \(n - 0 = 0\).
Chaining the equalities, we obtain \(\operatorname{eval}(\operatorname{Lean\mathord{:}Omega\mathord{:}Linear\,Combo\mathord{:}mk}(0,\allowbreak{} \langle 0\rangle \mathbin{\Vert} \langle 1\rangle) - \operatorname{Lean\mathord{:}Omega\mathord{:}Linear\,Combo\mathord{:}mk}(0,\allowbreak{} \langle\,\rangle),\allowbreak{} \operatorname{of\,List}(\langle \operatorname{length}(p)\rangle \mathbin{\Vert} \langle n\rangle)) = 0\).
By \paperlink{premise:4c65616e2e4f6d6567612e436f6e73747261696e742e616464457175616c6974795f736174}{Prerequisite~P\ref{premise:4c65616e2e4f6d6567612e436f6e73747261696e742e616464457175616c6974795f736174}}, we obtain \(\operatorname{sat^{\prime}}(\operatorname{Lean\mathord{:}Omega\mathord{:}Constraint\mathord{:}mk}(\operatorname{Option\mathord{:}some}(-0),\allowbreak{} \operatorname{Option\mathord{:}some}(-0)),\allowbreak{} \langle 0\rangle \mathbin{\Vert} \langle 1\rangle,\allowbreak{} \operatorname{of\,List}(\langle \operatorname{length}(p)\rangle \mathbin{\Vert} \langle n\rangle)) = \mathrm{true}\).
Use the hypothesis \(n > \operatorname{length}(p)\).
By \paperlink{premise:4c65616e2e4f6d6567612e496e742e6f664e61745f6c745f6f665f6c74}{Prerequisite~P\ref{premise:4c65616e2e4f6d6567612e496e742e6f664e61745f6c745f6f665f6c74}}, we obtain \(\operatorname{length}(p) < n\).
By \paperlink{premise:496e742e6164645f6f6e655f6c655f6f665f6c74}{Prerequisite~P\ref{premise:496e742e6164645f6f6e655f6c655f6f665f6c74}}, we obtain \(\operatorname{length}(p) + 1 \le n\).
By \paperlink{premise:496e742e7375625f6e6f6e6e65675f6f665f6c65}{Prerequisite~P\ref{premise:496e742e7375625f6e6f6e6e65675f6f665f6c65}}, we obtain \(0 \le n - \left(\operatorname{length}(p) + 1\right)\).
Using \paperlink{premise:4c65616e2e4f6d6567612e4c696e656172436f6d626f2e636f6f7264696e6174655f6576616c5f31}{Prerequisite~P\ref{premise:4c65616e2e4f6d6567612e4c696e656172436f6d626f2e636f6f7264696e6174655f6576616c5f31}} in reverse gives \(n = \operatorname{eval}(\operatorname{coordinate}(1),\allowbreak{} \operatorname{of\,List}(\langle \operatorname{length}(p)\rangle \mathbin{\Vert} \langle n\rangle))\).
Using \paperlink{premise:4c65616e2e4f6d6567612e4c696e656172436f6d626f2e636f6f7264696e6174655f6576616c5f30}{Prerequisite~P\ref{premise:4c65616e2e4f6d6567612e4c696e656172436f6d626f2e636f6f7264696e6174655f6576616c5f30}} in reverse gives \(\operatorname{length}(p) = \operatorname{eval}(\operatorname{coordinate}(0),\allowbreak{} \operatorname{of\,List}(\langle \operatorname{length}(p)\rangle \mathbin{\Vert} \langle n\rangle))\).
By definition, we obtain \(1 = 1\).
By the preceding argument, we obtain \(1 = \operatorname{eval}(\operatorname{Lean\mathord{:}Omega\mathord{:}Linear\,Combo\mathord{:}mk}(1,\allowbreak{} \langle\,\rangle),\allowbreak{} \operatorname{of\,List}(\langle \operatorname{length}(p)\rangle \mathbin{\Vert} \langle n\rangle))\).
By \paperlink{premise:4c65616e2e4f6d6567612e496e742e6164645f636f6e6772}{Prerequisite~P\ref{premise:4c65616e2e4f6d6567612e496e742e6164645f636f6e6772}}, we obtain \(\operatorname{length}(p) + 1 = \operatorname{eval}(\operatorname{coordinate}(0),\allowbreak{} \operatorname{of\,List}(\langle \operatorname{length}(p)\rangle \mathbin{\Vert} \langle n\rangle)) + \operatorname{eval}(\operatorname{Lean\mathord{:}Omega\mathord{:}Linear\,Combo\mathord{:}mk}(1,\allowbreak{} \langle\,\rangle),\allowbreak{} \operatorname{of\,List}(\langle \operatorname{length}(p)\rangle \mathbin{\Vert} \langle n\rangle))\).
Using \paperlink{premise:4c65616e2e4f6d6567612e4c696e656172436f6d626f2e6164645f6576616c}{Prerequisite~P\ref{premise:4c65616e2e4f6d6567612e4c696e656172436f6d626f2e6164645f6576616c}} with \(l_{1} = \operatorname{Lean\mathord{:}Omega\mathord{:}Linear\,Combo\mathord{:}mk}(0,\allowbreak{} \langle 1\rangle)\) and \(l_{2} = \operatorname{Lean\mathord{:}Omega\mathord{:}Linear\,Combo\mathord{:}mk}(1,\allowbreak{} \langle\,\rangle)\) and \(\operatorname{of\,List}\left(\langle \operatorname{length}(p)\rangle \mathbin{\Vert} \langle n\rangle\right)\) in place of its parameter \(v\) in reverse gives \(\operatorname{eval}(\operatorname{Lean\mathord{:}Omega\mathord{:}Linear\,Combo\mathord{:}mk}(0,\allowbreak{} \langle 1\rangle),\allowbreak{} \operatorname{of\,List}(\langle \operatorname{length}(p)\rangle \mathbin{\Vert} \langle n\rangle)) + \operatorname{eval}(\operatorname{Lean\mathord{:}Omega\mathord{:}Linear\,Combo\mathord{:}mk}(1,\allowbreak{} \langle\,\rangle),\allowbreak{} \operatorname{of\,List}(\langle \operatorname{length}(p)\rangle \mathbin{\Vert} \langle n\rangle)) = \operatorname{eval}(\operatorname{Lean\mathord{:}Omega\mathord{:}Linear\,Combo\mathord{:}mk}(0,\allowbreak{} \langle 1\rangle) + \operatorname{Lean\mathord{:}Omega\mathord{:}Linear\,Combo\mathord{:}mk}(1,\allowbreak{} \langle\,\rangle),\allowbreak{} \operatorname{of\,List}(\langle \operatorname{length}(p)\rangle \mathbin{\Vert} \langle n\rangle))\).
Chaining the equalities, we obtain \(\operatorname{length}(p) + 1 = \operatorname{eval}(\operatorname{Lean\mathord{:}Omega\mathord{:}Linear\,Combo\mathord{:}mk}(0,\allowbreak{} \langle 1\rangle) + \operatorname{Lean\mathord{:}Omega\mathord{:}Linear\,Combo\mathord{:}mk}(1,\allowbreak{} \langle\,\rangle),\allowbreak{} \operatorname{of\,List}(\langle \operatorname{length}(p)\rangle \mathbin{\Vert} \langle n\rangle))\).
By \paperlink{premise:4c65616e2e4f6d6567612e496e742e7375625f636f6e6772}{Prerequisite~P\ref{premise:4c65616e2e4f6d6567612e496e742e7375625f636f6e6772}}, we obtain \(n - (\operatorname{length}(p) + 1) = \operatorname{eval}(\operatorname{coordinate}(1),\allowbreak{} \operatorname{of\,List}(\langle \operatorname{length}(p)\rangle \mathbin{\Vert} \langle n\rangle)) - \operatorname{eval}(\operatorname{Lean\mathord{:}Omega\mathord{:}Linear\,Combo\mathord{:}mk}(0,\allowbreak{} \langle 1\rangle) + \operatorname{Lean\mathord{:}Omega\mathord{:}Linear\,Combo\mathord{:}mk}(1,\allowbreak{} \langle\,\rangle),\allowbreak{} \operatorname{of\,List}(\langle \operatorname{length}(p)\rangle \mathbin{\Vert} \langle n\rangle))\).
Using \paperlink{premise:4c65616e2e4f6d6567612e4c696e656172436f6d626f2e7375625f6576616c}{Prerequisite~P\ref{premise:4c65616e2e4f6d6567612e4c696e656172436f6d626f2e7375625f6576616c}} with \(l_{1} = \operatorname{Lean\mathord{:}Omega\mathord{:}Linear\,Combo\mathord{:}mk}(0,\allowbreak{} \langle 0\rangle \mathbin{\Vert} \langle 1\rangle)\) and \(l_{2} = \operatorname{Lean\mathord{:}Omega\mathord{:}Linear\,Combo\mathord{:}mk}(1,\allowbreak{} \langle 1\rangle)\) and \(\operatorname{of\,List}\left(\langle \operatorname{length}(p)\rangle \mathbin{\Vert} \langle n\rangle\right)\) in place of its parameter \(v\) in reverse gives \(\operatorname{eval}(\operatorname{Lean\mathord{:}Omega\mathord{:}Linear\,Combo\mathord{:}mk}(0,\allowbreak{} \langle 0\rangle \mathbin{\Vert} \langle 1\rangle),\allowbreak{} \operatorname{of\,List}(\langle \operatorname{length}(p)\rangle \mathbin{\Vert} \langle n\rangle)) - \operatorname{eval}(\operatorname{Lean\mathord{:}Omega\mathord{:}Linear\,Combo\mathord{:}mk}(1,\allowbreak{} \langle 1\rangle),\allowbreak{} \operatorname{of\,List}(\langle \operatorname{length}(p)\rangle \mathbin{\Vert} \langle n\rangle)) = \operatorname{eval}(\operatorname{Lean\mathord{:}Omega\mathord{:}Linear\,Combo\mathord{:}mk}(0,\allowbreak{} \langle 0\rangle \mathbin{\Vert} \langle 1\rangle) - \operatorname{Lean\mathord{:}Omega\mathord{:}Linear\,Combo\mathord{:}mk}(1,\allowbreak{} \langle 1\rangle),\allowbreak{} \operatorname{of\,List}(\langle \operatorname{length}(p)\rangle \mathbin{\Vert} \langle n\rangle))\).
Chaining the equalities, we obtain \(n - (\operatorname{length}(p) + 1) = \operatorname{eval}(\operatorname{Lean\mathord{:}Omega\mathord{:}Linear\,Combo\mathord{:}mk}(0,\allowbreak{} \langle 0\rangle \mathbin{\Vert} \langle 1\rangle) - \operatorname{Lean\mathord{:}Omega\mathord{:}Linear\,Combo\mathord{:}mk}(1,\allowbreak{} \langle 1\rangle),\allowbreak{} \operatorname{of\,List}(\langle \operatorname{length}(p)\rangle \mathbin{\Vert} \langle n\rangle))\).
By \paperlink{premise:6c655f6f665f6c655f6f665f6571}{Prerequisite~P\ref{premise:6c655f6f665f6c655f6f665f6571}} with \(\alpha = \mathbb{Z}\), we obtain \(0 \le \operatorname{eval}(\operatorname{Lean\mathord{:}Omega\mathord{:}Linear\,Combo\mathord{:}mk}(0,\allowbreak{} \langle 0\rangle \mathbin{\Vert} \langle 1\rangle) - \operatorname{Lean\mathord{:}Omega\mathord{:}Linear\,Combo\mathord{:}mk}(1,\allowbreak{} \langle 1\rangle),\allowbreak{} \operatorname{of\,List}(\langle \operatorname{length}(p)\rangle \mathbin{\Vert} \langle n\rangle))\).
By \paperlink{premise:4c65616e2e4f6d6567612e436f6e73747261696e742e616464496e657175616c6974795f736174}{Prerequisite~P\ref{premise:4c65616e2e4f6d6567612e436f6e73747261696e742e616464496e657175616c6974795f736174}}, we obtain \(\operatorname{sat^{\prime}}(\operatorname{Lean\mathord{:}Omega\mathord{:}Constraint\mathord{:}mk}(\operatorname{Option\mathord{:}some}(-(-1)),\allowbreak{} \operatorname{none}),\allowbreak{} \langle -1\rangle \mathbin{\Vert} \langle 1\rangle,\allowbreak{} \operatorname{of\,List}(\langle \operatorname{length}(p)\rangle \mathbin{\Vert} \langle n\rangle)) = \mathrm{true}\).
By \paperlink{premise:4c65616e2e4f6d6567612e636f6d626f5f73617427}{Prerequisite~P\ref{premise:4c65616e2e4f6d6567612e636f6d626f5f73617427}} with \(s = \operatorname{Lean\mathord{:}Omega\mathord{:}Constraint\mathord{:}mk}(\operatorname{Option\mathord{:}some}(0),\allowbreak{} \operatorname{Option\mathord{:}some}(0))\) and \(t = \operatorname{Lean\mathord{:}Omega\mathord{:}Constraint\mathord{:}mk}(\operatorname{Option\mathord{:}some}(1),\allowbreak{} \operatorname{none})\) and \(-1\) in place of its parameter \(a\) and \(x = \langle 0\rangle \mathbin{\Vert} \langle 1\rangle\) and \(b = 1\) and \(y = \langle -1\rangle \mathbin{\Vert} \langle 1\rangle\) and \(\operatorname{of\,List}\left(\langle \operatorname{length}(p)\rangle \mathbin{\Vert} \langle n\rangle\right)\) in place of its parameter \(v\), we obtain \(\operatorname{sat^{\prime}}(\operatorname{Lean\mathord{:}Omega\mathord{:}Constraint\mathord{:}combo}(-1,\allowbreak{} \operatorname{Lean\mathord{:}Omega\mathord{:}Constraint\mathord{:}mk}(\operatorname{Option\mathord{:}some}(0),\allowbreak{} \operatorname{Option\mathord{:}some}(0)),\allowbreak{} 1,\allowbreak{} \operatorname{Lean\mathord{:}Omega\mathord{:}Constraint\mathord{:}mk}(\operatorname{Option\mathord{:}some}(1),\allowbreak{} \operatorname{none})),\allowbreak{} \operatorname{Lean\mathord{:}Omega\mathord{:}Coeffs\mathord{:}combo}(-1,\allowbreak{} \langle 0\rangle \mathbin{\Vert} \langle 1\rangle,\allowbreak{} 1,\allowbreak{} \langle -1\rangle \mathbin{\Vert} \langle 1\rangle),\allowbreak{} \operatorname{of\,List}(\langle \operatorname{length}(p)\rangle \mathbin{\Vert} \langle n\rangle)) = \mathrm{true}\).
By \paperlink{premise:4c65616e2e4f6d6567612e746964795f736174}{Prerequisite~P\ref{premise:4c65616e2e4f6d6567612e746964795f736174}}, we obtain \(\operatorname{sat^{\prime}}(\operatorname{tidy\,Constraint}(\operatorname{Lean\mathord{:}Omega\mathord{:}Constraint\mathord{:}mk}(\operatorname{Option\mathord{:}some}(1),\allowbreak{} \operatorname{none}),\allowbreak{} \langle -1\rangle \mathbin{\Vert} \langle 0\rangle),\allowbreak{} \operatorname{tidy\,Coeffs}(\operatorname{Lean\mathord{:}Omega\mathord{:}Constraint\mathord{:}mk}(\operatorname{Option\mathord{:}some}(1),\allowbreak{} \operatorname{none}),\allowbreak{} \langle -1\rangle \mathbin{\Vert} \langle 0\rangle),\allowbreak{} \operatorname{of\,List}(\langle \operatorname{length}(p)\rangle \mathbin{\Vert} \langle n\rangle)) = \mathrm{true}\).
Using \paperlink{premise:4c65616e2e4f6d6567612e4c696e656172436f6d626f2e636f6f7264696e6174655f6576616c5f31}{Prerequisite~P\ref{premise:4c65616e2e4f6d6567612e4c696e656172436f6d626f2e636f6f7264696e6174655f6576616c5f31}} in reverse gives \(n = \operatorname{eval}(\operatorname{coordinate}(1),\allowbreak{} \operatorname{of\,List}(\langle \operatorname{length}(p)\rangle \mathbin{\Vert} \langle n\rangle))\).
By definition, we obtain \(0 = 0\).
By the preceding argument, we obtain \(0 = \operatorname{eval}(\operatorname{Lean\mathord{:}Omega\mathord{:}Linear\,Combo\mathord{:}mk}(0,\allowbreak{} \langle\,\rangle),\allowbreak{} \operatorname{of\,List}(\langle \operatorname{length}(p)\rangle \mathbin{\Vert} \langle n\rangle))\).
By \paperlink{premise:4c65616e2e4f6d6567612e496e742e7375625f636f6e6772}{Prerequisite~P\ref{premise:4c65616e2e4f6d6567612e496e742e7375625f636f6e6772}}, we obtain \(n - 0 = \operatorname{eval}(\operatorname{coordinate}(1),\allowbreak{} \operatorname{of\,List}(\langle \operatorname{length}(p)\rangle \mathbin{\Vert} \langle n\rangle)) - \operatorname{eval}(\operatorname{Lean\mathord{:}Omega\mathord{:}Linear\,Combo\mathord{:}mk}(0,\allowbreak{} \langle\,\rangle),\allowbreak{} \operatorname{of\,List}(\langle \operatorname{length}(p)\rangle \mathbin{\Vert} \langle n\rangle))\).
Using \paperlink{premise:4c65616e2e4f6d6567612e4c696e656172436f6d626f2e7375625f6576616c}{Prerequisite~P\ref{premise:4c65616e2e4f6d6567612e4c696e656172436f6d626f2e7375625f6576616c}} with \(l_{1} = \operatorname{Lean\mathord{:}Omega\mathord{:}Linear\,Combo\mathord{:}mk}(0,\allowbreak{} \langle 0\rangle \mathbin{\Vert} \langle 1\rangle)\) and \(l_{2} = \operatorname{Lean\mathord{:}Omega\mathord{:}Linear\,Combo\mathord{:}mk}(0,\allowbreak{} \langle\,\rangle)\) and \(\operatorname{of\,List}\left(\langle \operatorname{length}(p)\rangle \mathbin{\Vert} \langle n\rangle\right)\) in place of its parameter \(v\) in reverse gives \(\operatorname{eval}(\operatorname{Lean\mathord{:}Omega\mathord{:}Linear\,Combo\mathord{:}mk}(0,\allowbreak{} \langle 0\rangle \mathbin{\Vert} \langle 1\rangle),\allowbreak{} \operatorname{of\,List}(\langle \operatorname{length}(p)\rangle \mathbin{\Vert} \langle n\rangle)) - \operatorname{eval}(\operatorname{Lean\mathord{:}Omega\mathord{:}Linear\,Combo\mathord{:}mk}(0,\allowbreak{} \langle\,\rangle),\allowbreak{} \operatorname{of\,List}(\langle \operatorname{length}(p)\rangle \mathbin{\Vert} \langle n\rangle)) = \operatorname{eval}(\operatorname{Lean\mathord{:}Omega\mathord{:}Linear\,Combo\mathord{:}mk}(0,\allowbreak{} \langle 0\rangle \mathbin{\Vert} \langle 1\rangle) - \operatorname{Lean\mathord{:}Omega\mathord{:}Linear\,Combo\mathord{:}mk}(0,\allowbreak{} \langle\,\rangle),\allowbreak{} \operatorname{of\,List}(\langle \operatorname{length}(p)\rangle \mathbin{\Vert} \langle n\rangle))\).
Chaining the equalities, we obtain \(n - 0 = \operatorname{eval}(\operatorname{Lean\mathord{:}Omega\mathord{:}Linear\,Combo\mathord{:}mk}(0,\allowbreak{} \langle 0\rangle \mathbin{\Vert} \langle 1\rangle) - \operatorname{Lean\mathord{:}Omega\mathord{:}Linear\,Combo\mathord{:}mk}(0,\allowbreak{} \langle\,\rangle),\allowbreak{} \operatorname{of\,List}(\langle \operatorname{length}(p)\rangle \mathbin{\Vert} \langle n\rangle))\).
Reversing the equality, we obtain \(\operatorname{eval}(\operatorname{Lean\mathord{:}Omega\mathord{:}Linear\,Combo\mathord{:}mk}(0,\allowbreak{} \langle 0\rangle \mathbin{\Vert} \langle 1\rangle) - \operatorname{Lean\mathord{:}Omega\mathord{:}Linear\,Combo\mathord{:}mk}(0,\allowbreak{} \langle\,\rangle),\allowbreak{} \operatorname{of\,List}(\langle \operatorname{length}(p)\rangle \mathbin{\Vert} \langle n\rangle)) = n - 0\).
Use the hypothesis \(n = 0\).
By \paperlink{premise:4c65616e2e4f6d6567612e496e742e6f664e61745f636f6e6772}{Prerequisite~P\ref{premise:4c65616e2e4f6d6567612e496e742e6f664e61745f636f6e6772}}, we obtain \(n = 0\).
By \paperlink{premise:496e742e7375625f65715f7a65726f5f6f665f6571}{Prerequisite~P\ref{premise:496e742e7375625f65715f7a65726f5f6f665f6571}}, we obtain \(n - 0 = 0\).
Chaining the equalities, we obtain \(\operatorname{eval}(\operatorname{Lean\mathord{:}Omega\mathord{:}Linear\,Combo\mathord{:}mk}(0,\allowbreak{} \langle 0\rangle \mathbin{\Vert} \langle 1\rangle) - \operatorname{Lean\mathord{:}Omega\mathord{:}Linear\,Combo\mathord{:}mk}(0,\allowbreak{} \langle\,\rangle),\allowbreak{} \operatorname{of\,List}(\langle \operatorname{length}(p)\rangle \mathbin{\Vert} \langle n\rangle)) = 0\).
By \paperlink{premise:4c65616e2e4f6d6567612e436f6e73747261696e742e616464457175616c6974795f736174}{Prerequisite~P\ref{premise:4c65616e2e4f6d6567612e436f6e73747261696e742e616464457175616c6974795f736174}}, we obtain \(\operatorname{sat^{\prime}}(\operatorname{Lean\mathord{:}Omega\mathord{:}Constraint\mathord{:}mk}(\operatorname{Option\mathord{:}some}(-0),\allowbreak{} \operatorname{Option\mathord{:}some}(-0)),\allowbreak{} \langle 0\rangle \mathbin{\Vert} \langle 1\rangle,\allowbreak{} \operatorname{of\,List}(\langle \operatorname{length}(p)\rangle \mathbin{\Vert} \langle n\rangle)) = \mathrm{true}\).
Use the hypothesis \(n \le \operatorname{length}(p)\).
By \paperlink{premise:4c65616e2e4f6d6567612e496e742e6f664e61745f6c655f6f665f6c65}{Prerequisite~P\ref{premise:4c65616e2e4f6d6567612e496e742e6f664e61745f6c655f6f665f6c65}}, we obtain \(n \le \operatorname{length}(p)\).
By \paperlink{premise:496e742e7375625f6e6f6e6e65675f6f665f6c65}{Prerequisite~P\ref{premise:496e742e7375625f6e6f6e6e65675f6f665f6c65}}, we obtain \(0 \le \operatorname{length}(p) - n\).
Using \paperlink{premise:4c65616e2e4f6d6567612e4c696e656172436f6d626f2e636f6f7264696e6174655f6576616c5f30}{Prerequisite~P\ref{premise:4c65616e2e4f6d6567612e4c696e656172436f6d626f2e636f6f7264696e6174655f6576616c5f30}} in reverse gives \(\operatorname{length}(p) = \operatorname{eval}(\operatorname{coordinate}(0),\allowbreak{} \operatorname{of\,List}(\langle \operatorname{length}(p)\rangle \mathbin{\Vert} \langle n\rangle))\).
Using \paperlink{premise:4c65616e2e4f6d6567612e4c696e656172436f6d626f2e636f6f7264696e6174655f6576616c5f31}{Prerequisite~P\ref{premise:4c65616e2e4f6d6567612e4c696e656172436f6d626f2e636f6f7264696e6174655f6576616c5f31}} in reverse gives \(n = \operatorname{eval}(\operatorname{coordinate}(1),\allowbreak{} \operatorname{of\,List}(\langle \operatorname{length}(p)\rangle \mathbin{\Vert} \langle n\rangle))\).
By \paperlink{premise:4c65616e2e4f6d6567612e496e742e7375625f636f6e6772}{Prerequisite~P\ref{premise:4c65616e2e4f6d6567612e496e742e7375625f636f6e6772}}, we obtain \(\operatorname{length}(p) - n = \operatorname{eval}(\operatorname{coordinate}(0),\allowbreak{} \operatorname{of\,List}(\langle \operatorname{length}(p)\rangle \mathbin{\Vert} \langle n\rangle)) - \operatorname{eval}(\operatorname{coordinate}(1),\allowbreak{} \operatorname{of\,List}(\langle \operatorname{length}(p)\rangle \mathbin{\Vert} \langle n\rangle))\).
Using \paperlink{premise:4c65616e2e4f6d6567612e4c696e656172436f6d626f2e7375625f6576616c}{Prerequisite~P\ref{premise:4c65616e2e4f6d6567612e4c696e656172436f6d626f2e7375625f6576616c}} with \(l_{1} = \operatorname{Lean\mathord{:}Omega\mathord{:}Linear\,Combo\mathord{:}mk}(0,\allowbreak{} \langle 1\rangle)\) and \(l_{2} = \operatorname{Lean\mathord{:}Omega\mathord{:}Linear\,Combo\mathord{:}mk}(0,\allowbreak{} \langle 0\rangle \mathbin{\Vert} \langle 1\rangle)\) and \(\operatorname{of\,List}\left(\langle \operatorname{length}(p)\rangle \mathbin{\Vert} \langle n\rangle\right)\) in place of its parameter \(v\) in reverse gives \(\operatorname{eval}(\operatorname{Lean\mathord{:}Omega\mathord{:}Linear\,Combo\mathord{:}mk}(0,\allowbreak{} \langle 1\rangle),\allowbreak{} \operatorname{of\,List}(\langle \operatorname{length}(p)\rangle \mathbin{\Vert} \langle n\rangle)) - \operatorname{eval}(\operatorname{Lean\mathord{:}Omega\mathord{:}Linear\,Combo\mathord{:}mk}(0,\allowbreak{} \langle 0\rangle \mathbin{\Vert} \langle 1\rangle),\allowbreak{} \operatorname{of\,List}(\langle \operatorname{length}(p)\rangle \mathbin{\Vert} \langle n\rangle)) = \operatorname{eval}(\operatorname{Lean\mathord{:}Omega\mathord{:}Linear\,Combo\mathord{:}mk}(0,\allowbreak{} \langle 1\rangle) - \operatorname{Lean\mathord{:}Omega\mathord{:}Linear\,Combo\mathord{:}mk}(0,\allowbreak{} \langle 0\rangle \mathbin{\Vert} \langle 1\rangle),\allowbreak{} \operatorname{of\,List}(\langle \operatorname{length}(p)\rangle \mathbin{\Vert} \langle n\rangle))\).
Chaining the equalities, we obtain \(\operatorname{length}(p) - n = \operatorname{eval}(\operatorname{Lean\mathord{:}Omega\mathord{:}Linear\,Combo\mathord{:}mk}(0,\allowbreak{} \langle 1\rangle) - \operatorname{Lean\mathord{:}Omega\mathord{:}Linear\,Combo\mathord{:}mk}(0,\allowbreak{} \langle 0\rangle \mathbin{\Vert} \langle 1\rangle),\allowbreak{} \operatorname{of\,List}(\langle \operatorname{length}(p)\rangle \mathbin{\Vert} \langle n\rangle))\).
By \paperlink{premise:6c655f6f665f6c655f6f665f6571}{Prerequisite~P\ref{premise:6c655f6f665f6c655f6f665f6571}} with \(\alpha = \mathbb{Z}\), we obtain \(0 \le \operatorname{eval}(\operatorname{Lean\mathord{:}Omega\mathord{:}Linear\,Combo\mathord{:}mk}(0,\allowbreak{} \langle 1\rangle) - \operatorname{Lean\mathord{:}Omega\mathord{:}Linear\,Combo\mathord{:}mk}(0,\allowbreak{} \langle 0\rangle \mathbin{\Vert} \langle 1\rangle),\allowbreak{} \operatorname{of\,List}(\langle \operatorname{length}(p)\rangle \mathbin{\Vert} \langle n\rangle))\).
By \paperlink{premise:4c65616e2e4f6d6567612e436f6e73747261696e742e616464496e657175616c6974795f736174}{Prerequisite~P\ref{premise:4c65616e2e4f6d6567612e436f6e73747261696e742e616464496e657175616c6974795f736174}}, we obtain \(\operatorname{sat^{\prime}}(\operatorname{Lean\mathord{:}Omega\mathord{:}Constraint\mathord{:}mk}(\operatorname{Option\mathord{:}some}(-0),\allowbreak{} \operatorname{none}),\allowbreak{} \langle 1\rangle \mathbin{\Vert} \langle -1\rangle,\allowbreak{} \operatorname{of\,List}(\langle \operatorname{length}(p)\rangle \mathbin{\Vert} \langle n\rangle)) = \mathrm{true}\).
By \paperlink{premise:4c65616e2e4f6d6567612e636f6d626f5f73617427}{Prerequisite~P\ref{premise:4c65616e2e4f6d6567612e636f6d626f5f73617427}} with \(s = \operatorname{Lean\mathord{:}Omega\mathord{:}Constraint\mathord{:}mk}(\operatorname{Option\mathord{:}some}(0),\allowbreak{} \operatorname{Option\mathord{:}some}(0))\) and \(t = \operatorname{Lean\mathord{:}Omega\mathord{:}Constraint\mathord{:}mk}(\operatorname{Option\mathord{:}some}(0),\allowbreak{} \operatorname{none})\) and \(1\) in place of its parameter \(a\) and \(x = \langle 0\rangle \mathbin{\Vert} \langle 1\rangle\) and \(b = 1\) and \(y = \langle 1\rangle \mathbin{\Vert} \langle -1\rangle\) and \(\operatorname{of\,List}\left(\langle \operatorname{length}(p)\rangle \mathbin{\Vert} \langle n\rangle\right)\) in place of its parameter \(v\), we obtain \(\operatorname{sat^{\prime}}(\operatorname{Lean\mathord{:}Omega\mathord{:}Constraint\mathord{:}combo}(1,\allowbreak{} \operatorname{Lean\mathord{:}Omega\mathord{:}Constraint\mathord{:}mk}(\operatorname{Option\mathord{:}some}(0),\allowbreak{} \operatorname{Option\mathord{:}some}(0)),\allowbreak{} 1,\allowbreak{} \operatorname{Lean\mathord{:}Omega\mathord{:}Constraint\mathord{:}mk}(\operatorname{Option\mathord{:}some}(0),\allowbreak{} \operatorname{none})),\allowbreak{} \operatorname{Lean\mathord{:}Omega\mathord{:}Coeffs\mathord{:}combo}(1,\allowbreak{} \langle 0\rangle \mathbin{\Vert} \langle 1\rangle,\allowbreak{} 1,\allowbreak{} \langle 1\rangle \mathbin{\Vert} \langle -1\rangle),\allowbreak{} \operatorname{of\,List}(\langle \operatorname{length}(p)\rangle \mathbin{\Vert} \langle n\rangle)) = \mathrm{true}\).
By \paperlink{premise:4c65616e2e4f6d6567612e436f6e73747261696e742e636f6d62696e655f73617427}{Prerequisite~P\ref{premise:4c65616e2e4f6d6567612e436f6e73747261696e742e636f6d62696e655f73617427}}, we obtain \(\operatorname{sat^{\prime}}(\operatorname{combine}(\operatorname{Lean\mathord{:}Omega\mathord{:}Constraint\mathord{:}mk}(\operatorname{none},\allowbreak{} \operatorname{Option\mathord{:}some}(-1)),\allowbreak{} \operatorname{Lean\mathord{:}Omega\mathord{:}Constraint\mathord{:}mk}(\operatorname{Option\mathord{:}some}(0),\allowbreak{} \operatorname{none})),\allowbreak{} \langle 1\rangle \mathbin{\Vert} \langle 0\rangle,\allowbreak{} \operatorname{of\,List}(\langle \operatorname{length}(p)\rangle \mathbin{\Vert} \langle n\rangle)) = \mathrm{true}\).
By \paperlink{premise:4c65616e2e4f6d6567612e436f6e73747261696e742e6e6f745f736174275f6f665f6973496d706f737369626c65}{Prerequisite~P\ref{premise:4c65616e2e4f6d6567612e436f6e73747261696e742e6e6f745f736174275f6f665f6973496d706f737369626c65}}, we obtain \(\mathrm{False}\).
\end{proofstep}
The two alternatives exhaust the cases. Hence \(\mathrm{False}\).
\end{proofstep}
The substitution uses the following argument.
\begin{proofstep}
Assume \(\neg \left(n < \operatorname{length}(p)\right)\).
Use the hypothesis \(\neg \left(n = 0 \lor n = \operatorname{length}(p)\right)\).
By \paperlink{premise:4c65616e2e4f6d6567612e616e645f6e6f745f6e6f745f6f665f6e6f745f6f72}{Prerequisite~P\ref{premise:4c65616e2e4f6d6567612e616e645f6e6f745f6e6f745f6f665f6e6f745f6f72}}, we obtain \(\neg \left(n = 0\right) \land \neg \left(n = \operatorname{length}(p)\right)\).
Taking the second conjunct, we obtain \(\neg \left(n = \operatorname{length}(p)\right)\).
By \paperlink{premise:4e61742e6c745f6f725f67745f6f665f6e65}{Prerequisite~P\ref{premise:4e61742e6c745f6f725f67745f6f665f6e65}}, we obtain \(n < \operatorname{length}(p) \lor n > \operatorname{length}(p)\).
Consider the following case.
\begin{proofstep}
Assume \(n < \operatorname{length}(p)\).
By definition, we obtain \(\mathrm{true} = \mathrm{true}\).
By the preceding argument, we obtain \(\operatorname{decide}(\operatorname{is\,Impossible}(\operatorname{Lean\mathord{:}Omega\mathord{:}Constraint\mathord{:}mk}(\operatorname{Option\mathord{:}some}(1),\allowbreak{} \operatorname{Option\mathord{:}some}(0))) = \mathrm{true}) = \mathrm{true}\).
By \paperlink{premise:6f665f6465636964655f65715f74727565}{Prerequisite~P\ref{premise:6f665f6465636964655f65715f74727565}}, we obtain \(\operatorname{is\,Impossible}(\operatorname{Lean\mathord{:}Omega\mathord{:}Constraint\mathord{:}mk}(\operatorname{Option\mathord{:}some}(1),\allowbreak{} \operatorname{Option\mathord{:}some}(0))) = \mathrm{true}\).
Use the hypothesis \(\neg \left(n < \operatorname{length}(p)\right)\).
By \paperlink{premise:4e61742e6c655f6f665f6e6f745f6c74}{Prerequisite~P\ref{premise:4e61742e6c655f6f665f6e6f745f6c74}}, we obtain \(\operatorname{length}(p) \le n\).
By \paperlink{premise:4c65616e2e4f6d6567612e496e742e6f664e61745f6c655f6f665f6c65}{Prerequisite~P\ref{premise:4c65616e2e4f6d6567612e496e742e6f664e61745f6c655f6f665f6c65}}, we obtain \(\operatorname{length}(p) \le n\).
By \paperlink{premise:496e742e7375625f6e6f6e6e65675f6f665f6c65}{Prerequisite~P\ref{premise:496e742e7375625f6e6f6e6e65675f6f665f6c65}}, we obtain \(0 \le n - \operatorname{length}(p)\).
Using \paperlink{premise:4c65616e2e4f6d6567612e4c696e656172436f6d626f2e636f6f7264696e6174655f6576616c5f31}{Prerequisite~P\ref{premise:4c65616e2e4f6d6567612e4c696e656172436f6d626f2e636f6f7264696e6174655f6576616c5f31}} in reverse gives \(n = \operatorname{eval}(\operatorname{coordinate}(1),\allowbreak{} \operatorname{of\,List}(\langle \operatorname{length}(p)\rangle \mathbin{\Vert} \langle n\rangle))\).
Using \paperlink{premise:4c65616e2e4f6d6567612e4c696e656172436f6d626f2e636f6f7264696e6174655f6576616c5f30}{Prerequisite~P\ref{premise:4c65616e2e4f6d6567612e4c696e656172436f6d626f2e636f6f7264696e6174655f6576616c5f30}} in reverse gives \(\operatorname{length}(p) = \operatorname{eval}(\operatorname{coordinate}(0),\allowbreak{} \operatorname{of\,List}(\langle \operatorname{length}(p)\rangle \mathbin{\Vert} \langle n\rangle))\).
By \paperlink{premise:4c65616e2e4f6d6567612e496e742e7375625f636f6e6772}{Prerequisite~P\ref{premise:4c65616e2e4f6d6567612e496e742e7375625f636f6e6772}}, we obtain \(n - \operatorname{length}(p) = \operatorname{eval}(\operatorname{coordinate}(1),\allowbreak{} \operatorname{of\,List}(\langle \operatorname{length}(p)\rangle \mathbin{\Vert} \langle n\rangle)) - \operatorname{eval}(\operatorname{coordinate}(0),\allowbreak{} \operatorname{of\,List}(\langle \operatorname{length}(p)\rangle \mathbin{\Vert} \langle n\rangle))\).
Using \paperlink{premise:4c65616e2e4f6d6567612e4c696e656172436f6d626f2e7375625f6576616c}{Prerequisite~P\ref{premise:4c65616e2e4f6d6567612e4c696e656172436f6d626f2e7375625f6576616c}} with \(l_{1} = \operatorname{Lean\mathord{:}Omega\mathord{:}Linear\,Combo\mathord{:}mk}(0,\allowbreak{} \langle 0\rangle \mathbin{\Vert} \langle 1\rangle)\) and \(l_{2} = \operatorname{Lean\mathord{:}Omega\mathord{:}Linear\,Combo\mathord{:}mk}(0,\allowbreak{} \langle 1\rangle)\) and \(\operatorname{of\,List}\left(\langle \operatorname{length}(p)\rangle \mathbin{\Vert} \langle n\rangle\right)\) in place of its parameter \(v\) in reverse gives \(\operatorname{eval}(\operatorname{Lean\mathord{:}Omega\mathord{:}Linear\,Combo\mathord{:}mk}(0,\allowbreak{} \langle 0\rangle \mathbin{\Vert} \langle 1\rangle),\allowbreak{} \operatorname{of\,List}(\langle \operatorname{length}(p)\rangle \mathbin{\Vert} \langle n\rangle)) - \operatorname{eval}(\operatorname{Lean\mathord{:}Omega\mathord{:}Linear\,Combo\mathord{:}mk}(0,\allowbreak{} \langle 1\rangle),\allowbreak{} \operatorname{of\,List}(\langle \operatorname{length}(p)\rangle \mathbin{\Vert} \langle n\rangle)) = \operatorname{eval}(\operatorname{Lean\mathord{:}Omega\mathord{:}Linear\,Combo\mathord{:}mk}(0,\allowbreak{} \langle 0\rangle \mathbin{\Vert} \langle 1\rangle) - \operatorname{Lean\mathord{:}Omega\mathord{:}Linear\,Combo\mathord{:}mk}(0,\allowbreak{} \langle 1\rangle),\allowbreak{} \operatorname{of\,List}(\langle \operatorname{length}(p)\rangle \mathbin{\Vert} \langle n\rangle))\).
Chaining the equalities, we obtain \(n - \operatorname{length}(p) = \operatorname{eval}(\operatorname{Lean\mathord{:}Omega\mathord{:}Linear\,Combo\mathord{:}mk}(0,\allowbreak{} \langle 0\rangle \mathbin{\Vert} \langle 1\rangle) - \operatorname{Lean\mathord{:}Omega\mathord{:}Linear\,Combo\mathord{:}mk}(0,\allowbreak{} \langle 1\rangle),\allowbreak{} \operatorname{of\,List}(\langle \operatorname{length}(p)\rangle \mathbin{\Vert} \langle n\rangle))\).
By \paperlink{premise:6c655f6f665f6c655f6f665f6571}{Prerequisite~P\ref{premise:6c655f6f665f6c655f6f665f6571}} with \(\alpha = \mathbb{Z}\), we obtain \(0 \le \operatorname{eval}(\operatorname{Lean\mathord{:}Omega\mathord{:}Linear\,Combo\mathord{:}mk}(0,\allowbreak{} \langle 0\rangle \mathbin{\Vert} \langle 1\rangle) - \operatorname{Lean\mathord{:}Omega\mathord{:}Linear\,Combo\mathord{:}mk}(0,\allowbreak{} \langle 1\rangle),\allowbreak{} \operatorname{of\,List}(\langle \operatorname{length}(p)\rangle \mathbin{\Vert} \langle n\rangle))\).
By \paperlink{premise:4c65616e2e4f6d6567612e436f6e73747261696e742e616464496e657175616c6974795f736174}{Prerequisite~P\ref{premise:4c65616e2e4f6d6567612e436f6e73747261696e742e616464496e657175616c6974795f736174}}, we obtain \(\operatorname{sat^{\prime}}(\operatorname{Lean\mathord{:}Omega\mathord{:}Constraint\mathord{:}mk}(\operatorname{Option\mathord{:}some}(-0),\allowbreak{} \operatorname{none}),\allowbreak{} \langle -1\rangle \mathbin{\Vert} \langle 1\rangle,\allowbreak{} \operatorname{of\,List}(\langle \operatorname{length}(p)\rangle \mathbin{\Vert} \langle n\rangle)) = \mathrm{true}\).
By \paperlink{premise:4c65616e2e4f6d6567612e746964795f736174}{Prerequisite~P\ref{premise:4c65616e2e4f6d6567612e746964795f736174}}, we obtain \(\operatorname{sat^{\prime}}(\operatorname{tidy\,Constraint}(\operatorname{Lean\mathord{:}Omega\mathord{:}Constraint\mathord{:}mk}(\operatorname{Option\mathord{:}some}(0),\allowbreak{} \operatorname{none}),\allowbreak{} \langle -1\rangle \mathbin{\Vert} \langle 1\rangle),\allowbreak{} \operatorname{tidy\,Coeffs}(\operatorname{Lean\mathord{:}Omega\mathord{:}Constraint\mathord{:}mk}(\operatorname{Option\mathord{:}some}(0),\allowbreak{} \operatorname{none}),\allowbreak{} \langle -1\rangle \mathbin{\Vert} \langle 1\rangle),\allowbreak{} \operatorname{of\,List}(\langle \operatorname{length}(p)\rangle \mathbin{\Vert} \langle n\rangle)) = \mathrm{true}\).
Use the hypothesis \(n \le \operatorname{length}(p)\).
By \paperlink{premise:4c65616e2e4f6d6567612e496e742e6f664e61745f6c655f6f665f6c65}{Prerequisite~P\ref{premise:4c65616e2e4f6d6567612e496e742e6f664e61745f6c655f6f665f6c65}}, we obtain \(n \le \operatorname{length}(p)\).
By \paperlink{premise:496e742e7375625f6e6f6e6e65675f6f665f6c65}{Prerequisite~P\ref{premise:496e742e7375625f6e6f6e6e65675f6f665f6c65}}, we obtain \(0 \le \operatorname{length}(p) - n\).
Using \paperlink{premise:4c65616e2e4f6d6567612e4c696e656172436f6d626f2e636f6f7264696e6174655f6576616c5f30}{Prerequisite~P\ref{premise:4c65616e2e4f6d6567612e4c696e656172436f6d626f2e636f6f7264696e6174655f6576616c5f30}} in reverse gives \(\operatorname{length}(p) = \operatorname{eval}(\operatorname{coordinate}(0),\allowbreak{} \operatorname{of\,List}(\langle \operatorname{length}(p)\rangle \mathbin{\Vert} \langle n\rangle))\).
Using \paperlink{premise:4c65616e2e4f6d6567612e4c696e656172436f6d626f2e636f6f7264696e6174655f6576616c5f31}{Prerequisite~P\ref{premise:4c65616e2e4f6d6567612e4c696e656172436f6d626f2e636f6f7264696e6174655f6576616c5f31}} in reverse gives \(n = \operatorname{eval}(\operatorname{coordinate}(1),\allowbreak{} \operatorname{of\,List}(\langle \operatorname{length}(p)\rangle \mathbin{\Vert} \langle n\rangle))\).
By \paperlink{premise:4c65616e2e4f6d6567612e496e742e7375625f636f6e6772}{Prerequisite~P\ref{premise:4c65616e2e4f6d6567612e496e742e7375625f636f6e6772}}, we obtain \(\operatorname{length}(p) - n = \operatorname{eval}(\operatorname{coordinate}(0),\allowbreak{} \operatorname{of\,List}(\langle \operatorname{length}(p)\rangle \mathbin{\Vert} \langle n\rangle)) - \operatorname{eval}(\operatorname{coordinate}(1),\allowbreak{} \operatorname{of\,List}(\langle \operatorname{length}(p)\rangle \mathbin{\Vert} \langle n\rangle))\).
Using \paperlink{premise:4c65616e2e4f6d6567612e4c696e656172436f6d626f2e7375625f6576616c}{Prerequisite~P\ref{premise:4c65616e2e4f6d6567612e4c696e656172436f6d626f2e7375625f6576616c}} with \(l_{1} = \operatorname{Lean\mathord{:}Omega\mathord{:}Linear\,Combo\mathord{:}mk}(0,\allowbreak{} \langle 1\rangle)\) and \(l_{2} = \operatorname{Lean\mathord{:}Omega\mathord{:}Linear\,Combo\mathord{:}mk}(0,\allowbreak{} \langle 0\rangle \mathbin{\Vert} \langle 1\rangle)\) and \(\operatorname{of\,List}\left(\langle \operatorname{length}(p)\rangle \mathbin{\Vert} \langle n\rangle\right)\) in place of its parameter \(v\) in reverse gives \(\operatorname{eval}(\operatorname{Lean\mathord{:}Omega\mathord{:}Linear\,Combo\mathord{:}mk}(0,\allowbreak{} \langle 1\rangle),\allowbreak{} \operatorname{of\,List}(\langle \operatorname{length}(p)\rangle \mathbin{\Vert} \langle n\rangle)) - \operatorname{eval}(\operatorname{Lean\mathord{:}Omega\mathord{:}Linear\,Combo\mathord{:}mk}(0,\allowbreak{} \langle 0\rangle \mathbin{\Vert} \langle 1\rangle),\allowbreak{} \operatorname{of\,List}(\langle \operatorname{length}(p)\rangle \mathbin{\Vert} \langle n\rangle)) = \operatorname{eval}(\operatorname{Lean\mathord{:}Omega\mathord{:}Linear\,Combo\mathord{:}mk}(0,\allowbreak{} \langle 1\rangle) - \operatorname{Lean\mathord{:}Omega\mathord{:}Linear\,Combo\mathord{:}mk}(0,\allowbreak{} \langle 0\rangle \mathbin{\Vert} \langle 1\rangle),\allowbreak{} \operatorname{of\,List}(\langle \operatorname{length}(p)\rangle \mathbin{\Vert} \langle n\rangle))\).
Chaining the equalities, we obtain \(\operatorname{length}(p) - n = \operatorname{eval}(\operatorname{Lean\mathord{:}Omega\mathord{:}Linear\,Combo\mathord{:}mk}(0,\allowbreak{} \langle 1\rangle) - \operatorname{Lean\mathord{:}Omega\mathord{:}Linear\,Combo\mathord{:}mk}(0,\allowbreak{} \langle 0\rangle \mathbin{\Vert} \langle 1\rangle),\allowbreak{} \operatorname{of\,List}(\langle \operatorname{length}(p)\rangle \mathbin{\Vert} \langle n\rangle))\).
By \paperlink{premise:6c655f6f665f6c655f6f665f6571}{Prerequisite~P\ref{premise:6c655f6f665f6c655f6f665f6571}} with \(\alpha = \mathbb{Z}\), we obtain \(0 \le \operatorname{eval}(\operatorname{Lean\mathord{:}Omega\mathord{:}Linear\,Combo\mathord{:}mk}(0,\allowbreak{} \langle 1\rangle) - \operatorname{Lean\mathord{:}Omega\mathord{:}Linear\,Combo\mathord{:}mk}(0,\allowbreak{} \langle 0\rangle \mathbin{\Vert} \langle 1\rangle),\allowbreak{} \operatorname{of\,List}(\langle \operatorname{length}(p)\rangle \mathbin{\Vert} \langle n\rangle))\).
By \paperlink{premise:4c65616e2e4f6d6567612e436f6e73747261696e742e616464496e657175616c6974795f736174}{Prerequisite~P\ref{premise:4c65616e2e4f6d6567612e436f6e73747261696e742e616464496e657175616c6974795f736174}}, we obtain \(\operatorname{sat^{\prime}}(\operatorname{Lean\mathord{:}Omega\mathord{:}Constraint\mathord{:}mk}(\operatorname{Option\mathord{:}some}(-0),\allowbreak{} \operatorname{none}),\allowbreak{} \langle 1\rangle \mathbin{\Vert} \langle -1\rangle,\allowbreak{} \operatorname{of\,List}(\langle \operatorname{length}(p)\rangle \mathbin{\Vert} \langle n\rangle)) = \mathrm{true}\).
By \paperlink{premise:4c65616e2e4f6d6567612e436f6e73747261696e742e636f6d62696e655f73617427}{Prerequisite~P\ref{premise:4c65616e2e4f6d6567612e436f6e73747261696e742e636f6d62696e655f73617427}}, we obtain \(\operatorname{sat^{\prime}}(\operatorname{combine}(\operatorname{Lean\mathord{:}Omega\mathord{:}Constraint\mathord{:}mk}(\operatorname{none},\allowbreak{} \operatorname{Option\mathord{:}some}(0)),\allowbreak{} \operatorname{Lean\mathord{:}Omega\mathord{:}Constraint\mathord{:}mk}(\operatorname{Option\mathord{:}some}(0),\allowbreak{} \operatorname{none})),\allowbreak{} \langle 1\rangle \mathbin{\Vert} \langle -1\rangle,\allowbreak{} \operatorname{of\,List}(\langle \operatorname{length}(p)\rangle \mathbin{\Vert} \langle n\rangle)) = \mathrm{true}\).
Use the hypothesis \(n < \operatorname{length}(p)\).
By \paperlink{premise:4c65616e2e4f6d6567612e496e742e6f664e61745f6c745f6f665f6c74}{Prerequisite~P\ref{premise:4c65616e2e4f6d6567612e496e742e6f664e61745f6c745f6f665f6c74}}, we obtain \(n < \operatorname{length}(p)\).
By \paperlink{premise:496e742e6164645f6f6e655f6c655f6f665f6c74}{Prerequisite~P\ref{premise:496e742e6164645f6f6e655f6c655f6f665f6c74}}, we obtain \(n + 1 \le \operatorname{length}(p)\).
By \paperlink{premise:496e742e7375625f6e6f6e6e65675f6f665f6c65}{Prerequisite~P\ref{premise:496e742e7375625f6e6f6e6e65675f6f665f6c65}}, we obtain \(0 \le \operatorname{length}(p) - \left(n + 1\right)\).
Using \paperlink{premise:4c65616e2e4f6d6567612e4c696e656172436f6d626f2e636f6f7264696e6174655f6576616c5f30}{Prerequisite~P\ref{premise:4c65616e2e4f6d6567612e4c696e656172436f6d626f2e636f6f7264696e6174655f6576616c5f30}} in reverse gives \(\operatorname{length}(p) = \operatorname{eval}(\operatorname{coordinate}(0),\allowbreak{} \operatorname{of\,List}(\langle \operatorname{length}(p)\rangle \mathbin{\Vert} \langle n\rangle))\).
Using \paperlink{premise:4c65616e2e4f6d6567612e4c696e656172436f6d626f2e636f6f7264696e6174655f6576616c5f31}{Prerequisite~P\ref{premise:4c65616e2e4f6d6567612e4c696e656172436f6d626f2e636f6f7264696e6174655f6576616c5f31}} in reverse gives \(n = \operatorname{eval}(\operatorname{coordinate}(1),\allowbreak{} \operatorname{of\,List}(\langle \operatorname{length}(p)\rangle \mathbin{\Vert} \langle n\rangle))\).
By definition, we obtain \(1 = 1\).
By the preceding argument, we obtain \(1 = \operatorname{eval}(\operatorname{Lean\mathord{:}Omega\mathord{:}Linear\,Combo\mathord{:}mk}(1,\allowbreak{} \langle\,\rangle),\allowbreak{} \operatorname{of\,List}(\langle \operatorname{length}(p)\rangle \mathbin{\Vert} \langle n\rangle))\).
By \paperlink{premise:4c65616e2e4f6d6567612e496e742e6164645f636f6e6772}{Prerequisite~P\ref{premise:4c65616e2e4f6d6567612e496e742e6164645f636f6e6772}}, we obtain \(n + 1 = \operatorname{eval}(\operatorname{coordinate}(1),\allowbreak{} \operatorname{of\,List}(\langle \operatorname{length}(p)\rangle \mathbin{\Vert} \langle n\rangle)) + \operatorname{eval}(\operatorname{Lean\mathord{:}Omega\mathord{:}Linear\,Combo\mathord{:}mk}(1,\allowbreak{} \langle\,\rangle),\allowbreak{} \operatorname{of\,List}(\langle \operatorname{length}(p)\rangle \mathbin{\Vert} \langle n\rangle))\).
Using \paperlink{premise:4c65616e2e4f6d6567612e4c696e656172436f6d626f2e6164645f6576616c}{Prerequisite~P\ref{premise:4c65616e2e4f6d6567612e4c696e656172436f6d626f2e6164645f6576616c}} with \(l_{1} = \operatorname{Lean\mathord{:}Omega\mathord{:}Linear\,Combo\mathord{:}mk}(0,\allowbreak{} \langle 0\rangle \mathbin{\Vert} \langle 1\rangle)\) and \(l_{2} = \operatorname{Lean\mathord{:}Omega\mathord{:}Linear\,Combo\mathord{:}mk}(1,\allowbreak{} \langle\,\rangle)\) and \(\operatorname{of\,List}\left(\langle \operatorname{length}(p)\rangle \mathbin{\Vert} \langle n\rangle\right)\) in place of its parameter \(v\) in reverse gives \(\operatorname{eval}(\operatorname{Lean\mathord{:}Omega\mathord{:}Linear\,Combo\mathord{:}mk}(0,\allowbreak{} \langle 0\rangle \mathbin{\Vert} \langle 1\rangle),\allowbreak{} \operatorname{of\,List}(\langle \operatorname{length}(p)\rangle \mathbin{\Vert} \langle n\rangle)) + \operatorname{eval}(\operatorname{Lean\mathord{:}Omega\mathord{:}Linear\,Combo\mathord{:}mk}(1,\allowbreak{} \langle\,\rangle),\allowbreak{} \operatorname{of\,List}(\langle \operatorname{length}(p)\rangle \mathbin{\Vert} \langle n\rangle)) = \operatorname{eval}(\operatorname{Lean\mathord{:}Omega\mathord{:}Linear\,Combo\mathord{:}mk}(0,\allowbreak{} \langle 0\rangle \mathbin{\Vert} \langle 1\rangle) + \operatorname{Lean\mathord{:}Omega\mathord{:}Linear\,Combo\mathord{:}mk}(1,\allowbreak{} \langle\,\rangle),\allowbreak{} \operatorname{of\,List}(\langle \operatorname{length}(p)\rangle \mathbin{\Vert} \langle n\rangle))\).
Chaining the equalities, we obtain \(n + 1 = \operatorname{eval}(\operatorname{Lean\mathord{:}Omega\mathord{:}Linear\,Combo\mathord{:}mk}(0,\allowbreak{} \langle 0\rangle \mathbin{\Vert} \langle 1\rangle) + \operatorname{Lean\mathord{:}Omega\mathord{:}Linear\,Combo\mathord{:}mk}(1,\allowbreak{} \langle\,\rangle),\allowbreak{} \operatorname{of\,List}(\langle \operatorname{length}(p)\rangle \mathbin{\Vert} \langle n\rangle))\).
By \paperlink{premise:4c65616e2e4f6d6567612e496e742e7375625f636f6e6772}{Prerequisite~P\ref{premise:4c65616e2e4f6d6567612e496e742e7375625f636f6e6772}}, we obtain \(\operatorname{length}(p) - (n + 1) = \operatorname{eval}(\operatorname{coordinate}(0),\allowbreak{} \operatorname{of\,List}(\langle \operatorname{length}(p)\rangle \mathbin{\Vert} \langle n\rangle)) - \operatorname{eval}(\operatorname{Lean\mathord{:}Omega\mathord{:}Linear\,Combo\mathord{:}mk}(0,\allowbreak{} \langle 0\rangle \mathbin{\Vert} \langle 1\rangle) + \operatorname{Lean\mathord{:}Omega\mathord{:}Linear\,Combo\mathord{:}mk}(1,\allowbreak{} \langle\,\rangle),\allowbreak{} \operatorname{of\,List}(\langle \operatorname{length}(p)\rangle \mathbin{\Vert} \langle n\rangle))\).
Using \paperlink{premise:4c65616e2e4f6d6567612e4c696e656172436f6d626f2e7375625f6576616c}{Prerequisite~P\ref{premise:4c65616e2e4f6d6567612e4c696e656172436f6d626f2e7375625f6576616c}} with \(l_{1} = \operatorname{Lean\mathord{:}Omega\mathord{:}Linear\,Combo\mathord{:}mk}(0,\allowbreak{} \langle 1\rangle)\) and \(l_{2} = \operatorname{Lean\mathord{:}Omega\mathord{:}Linear\,Combo\mathord{:}mk}(1,\allowbreak{} \langle 0\rangle \mathbin{\Vert} \langle 1\rangle)\) and \(\operatorname{of\,List}\left(\langle \operatorname{length}(p)\rangle \mathbin{\Vert} \langle n\rangle\right)\) in place of its parameter \(v\) in reverse gives \(\operatorname{eval}(\operatorname{Lean\mathord{:}Omega\mathord{:}Linear\,Combo\mathord{:}mk}(0,\allowbreak{} \langle 1\rangle),\allowbreak{} \operatorname{of\,List}(\langle \operatorname{length}(p)\rangle \mathbin{\Vert} \langle n\rangle)) - \operatorname{eval}(\operatorname{Lean\mathord{:}Omega\mathord{:}Linear\,Combo\mathord{:}mk}(1,\allowbreak{} \langle 0\rangle \mathbin{\Vert} \langle 1\rangle),\allowbreak{} \operatorname{of\,List}(\langle \operatorname{length}(p)\rangle \mathbin{\Vert} \langle n\rangle)) = \operatorname{eval}(\operatorname{Lean\mathord{:}Omega\mathord{:}Linear\,Combo\mathord{:}mk}(0,\allowbreak{} \langle 1\rangle) - \operatorname{Lean\mathord{:}Omega\mathord{:}Linear\,Combo\mathord{:}mk}(1,\allowbreak{} \langle 0\rangle \mathbin{\Vert} \langle 1\rangle),\allowbreak{} \operatorname{of\,List}(\langle \operatorname{length}(p)\rangle \mathbin{\Vert} \langle n\rangle))\).
Chaining the equalities, we obtain \(\operatorname{length}(p) - (n + 1) = \operatorname{eval}(\operatorname{Lean\mathord{:}Omega\mathord{:}Linear\,Combo\mathord{:}mk}(0,\allowbreak{} \langle 1\rangle) - \operatorname{Lean\mathord{:}Omega\mathord{:}Linear\,Combo\mathord{:}mk}(1,\allowbreak{} \langle 0\rangle \mathbin{\Vert} \langle 1\rangle),\allowbreak{} \operatorname{of\,List}(\langle \operatorname{length}(p)\rangle \mathbin{\Vert} \langle n\rangle))\).
By \paperlink{premise:6c655f6f665f6c655f6f665f6571}{Prerequisite~P\ref{premise:6c655f6f665f6c655f6f665f6571}} with \(\alpha = \mathbb{Z}\), we obtain \(0 \le \operatorname{eval}(\operatorname{Lean\mathord{:}Omega\mathord{:}Linear\,Combo\mathord{:}mk}(0,\allowbreak{} \langle 1\rangle) - \operatorname{Lean\mathord{:}Omega\mathord{:}Linear\,Combo\mathord{:}mk}(1,\allowbreak{} \langle 0\rangle \mathbin{\Vert} \langle 1\rangle),\allowbreak{} \operatorname{of\,List}(\langle \operatorname{length}(p)\rangle \mathbin{\Vert} \langle n\rangle))\).
By \paperlink{premise:4c65616e2e4f6d6567612e436f6e73747261696e742e616464496e657175616c6974795f736174}{Prerequisite~P\ref{premise:4c65616e2e4f6d6567612e436f6e73747261696e742e616464496e657175616c6974795f736174}}, we obtain \(\operatorname{sat^{\prime}}(\operatorname{Lean\mathord{:}Omega\mathord{:}Constraint\mathord{:}mk}(\operatorname{Option\mathord{:}some}(-(-1)),\allowbreak{} \operatorname{none}),\allowbreak{} \langle 1\rangle \mathbin{\Vert} \langle -1\rangle,\allowbreak{} \operatorname{of\,List}(\langle \operatorname{length}(p)\rangle \mathbin{\Vert} \langle n\rangle)) = \mathrm{true}\).
By \paperlink{premise:4c65616e2e4f6d6567612e636f6d626f5f73617427}{Prerequisite~P\ref{premise:4c65616e2e4f6d6567612e636f6d626f5f73617427}} with \(s = \operatorname{Lean\mathord{:}Omega\mathord{:}Constraint\mathord{:}mk}(\operatorname{Option\mathord{:}some}(0),\allowbreak{} \operatorname{Option\mathord{:}some}(0))\) and \(t = \operatorname{Lean\mathord{:}Omega\mathord{:}Constraint\mathord{:}mk}(\operatorname{Option\mathord{:}some}(1),\allowbreak{} \operatorname{none})\) and \(-1\) in place of its parameter \(a\) and \(x = \langle 1\rangle \mathbin{\Vert} \langle -1\rangle\) and \(b = 1\) and \(y = \langle 1\rangle \mathbin{\Vert} \langle -1\rangle\) and \(\operatorname{of\,List}\left(\langle \operatorname{length}(p)\rangle \mathbin{\Vert} \langle n\rangle\right)\) in place of its parameter \(v\), we obtain \(\operatorname{sat^{\prime}}(\operatorname{Lean\mathord{:}Omega\mathord{:}Constraint\mathord{:}combo}(-1,\allowbreak{} \operatorname{Lean\mathord{:}Omega\mathord{:}Constraint\mathord{:}mk}(\operatorname{Option\mathord{:}some}(0),\allowbreak{} \operatorname{Option\mathord{:}some}(0)),\allowbreak{} 1,\allowbreak{} \operatorname{Lean\mathord{:}Omega\mathord{:}Constraint\mathord{:}mk}(\operatorname{Option\mathord{:}some}(1),\allowbreak{} \operatorname{none})),\allowbreak{} \operatorname{Lean\mathord{:}Omega\mathord{:}Coeffs\mathord{:}combo}(-1,\allowbreak{} \langle 1\rangle \mathbin{\Vert} \langle -1\rangle,\allowbreak{} 1,\allowbreak{} \langle 1\rangle \mathbin{\Vert} \langle -1\rangle),\allowbreak{} \operatorname{of\,List}(\langle \operatorname{length}(p)\rangle \mathbin{\Vert} \langle n\rangle)) = \mathrm{true}\).
By \paperlink{premise:4c65616e2e4f6d6567612e746964795f736174}{Prerequisite~P\ref{premise:4c65616e2e4f6d6567612e746964795f736174}}, we obtain \(\operatorname{sat^{\prime}}(\operatorname{tidy\,Constraint}(\operatorname{Lean\mathord{:}Omega\mathord{:}Constraint\mathord{:}mk}(\operatorname{Option\mathord{:}some}(1),\allowbreak{} \operatorname{none}),\allowbreak{} \langle 0\rangle \mathbin{\Vert} \langle 0\rangle),\allowbreak{} \operatorname{tidy\,Coeffs}(\operatorname{Lean\mathord{:}Omega\mathord{:}Constraint\mathord{:}mk}(\operatorname{Option\mathord{:}some}(1),\allowbreak{} \operatorname{none}),\allowbreak{} \langle 0\rangle \mathbin{\Vert} \langle 0\rangle),\allowbreak{} \operatorname{of\,List}(\langle \operatorname{length}(p)\rangle \mathbin{\Vert} \langle n\rangle)) = \mathrm{true}\).
By \paperlink{premise:4c65616e2e4f6d6567612e436f6e73747261696e742e6e6f745f736174275f6f665f6973496d706f737369626c65}{Prerequisite~P\ref{premise:4c65616e2e4f6d6567612e436f6e73747261696e742e6e6f745f736174275f6f665f6973496d706f737369626c65}}, we obtain \(\mathrm{False}\).
\end{proofstep}
Consider the following case.
\begin{proofstep}
Assume \(n > \operatorname{length}(p)\).
By definition, we obtain \(\mathrm{true} = \mathrm{true}\).
By the preceding argument, we obtain \(\operatorname{decide}(\operatorname{is\,Impossible}(\operatorname{Lean\mathord{:}Omega\mathord{:}Constraint\mathord{:}mk}(\operatorname{Option\mathord{:}some}(1),\allowbreak{} \operatorname{Option\mathord{:}some}(0))) = \mathrm{true}) = \mathrm{true}\).
By \paperlink{premise:6f665f6465636964655f65715f74727565}{Prerequisite~P\ref{premise:6f665f6465636964655f65715f74727565}}, we obtain \(\operatorname{is\,Impossible}(\operatorname{Lean\mathord{:}Omega\mathord{:}Constraint\mathord{:}mk}(\operatorname{Option\mathord{:}some}(1),\allowbreak{} \operatorname{Option\mathord{:}some}(0))) = \mathrm{true}\).
Use the hypothesis \(\neg \left(n < \operatorname{length}(p)\right)\).
By \paperlink{premise:4e61742e6c655f6f665f6e6f745f6c74}{Prerequisite~P\ref{premise:4e61742e6c655f6f665f6e6f745f6c74}}, we obtain \(\operatorname{length}(p) \le n\).
By \paperlink{premise:4c65616e2e4f6d6567612e496e742e6f664e61745f6c655f6f665f6c65}{Prerequisite~P\ref{premise:4c65616e2e4f6d6567612e496e742e6f664e61745f6c655f6f665f6c65}}, we obtain \(\operatorname{length}(p) \le n\).
By \paperlink{premise:496e742e7375625f6e6f6e6e65675f6f665f6c65}{Prerequisite~P\ref{premise:496e742e7375625f6e6f6e6e65675f6f665f6c65}}, we obtain \(0 \le n - \operatorname{length}(p)\).
Using \paperlink{premise:4c65616e2e4f6d6567612e4c696e656172436f6d626f2e636f6f7264696e6174655f6576616c5f31}{Prerequisite~P\ref{premise:4c65616e2e4f6d6567612e4c696e656172436f6d626f2e636f6f7264696e6174655f6576616c5f31}} in reverse gives \(n = \operatorname{eval}(\operatorname{coordinate}(1),\allowbreak{} \operatorname{of\,List}(\langle \operatorname{length}(p)\rangle \mathbin{\Vert} \langle n\rangle))\).
Using \paperlink{premise:4c65616e2e4f6d6567612e4c696e656172436f6d626f2e636f6f7264696e6174655f6576616c5f30}{Prerequisite~P\ref{premise:4c65616e2e4f6d6567612e4c696e656172436f6d626f2e636f6f7264696e6174655f6576616c5f30}} in reverse gives \(\operatorname{length}(p) = \operatorname{eval}(\operatorname{coordinate}(0),\allowbreak{} \operatorname{of\,List}(\langle \operatorname{length}(p)\rangle \mathbin{\Vert} \langle n\rangle))\).
By \paperlink{premise:4c65616e2e4f6d6567612e496e742e7375625f636f6e6772}{Prerequisite~P\ref{premise:4c65616e2e4f6d6567612e496e742e7375625f636f6e6772}}, we obtain \(n - \operatorname{length}(p) = \operatorname{eval}(\operatorname{coordinate}(1),\allowbreak{} \operatorname{of\,List}(\langle \operatorname{length}(p)\rangle \mathbin{\Vert} \langle n\rangle)) - \operatorname{eval}(\operatorname{coordinate}(0),\allowbreak{} \operatorname{of\,List}(\langle \operatorname{length}(p)\rangle \mathbin{\Vert} \langle n\rangle))\).
Using \paperlink{premise:4c65616e2e4f6d6567612e4c696e656172436f6d626f2e7375625f6576616c}{Prerequisite~P\ref{premise:4c65616e2e4f6d6567612e4c696e656172436f6d626f2e7375625f6576616c}} with \(l_{1} = \operatorname{Lean\mathord{:}Omega\mathord{:}Linear\,Combo\mathord{:}mk}(0,\allowbreak{} \langle 0\rangle \mathbin{\Vert} \langle 1\rangle)\) and \(l_{2} = \operatorname{Lean\mathord{:}Omega\mathord{:}Linear\,Combo\mathord{:}mk}(0,\allowbreak{} \langle 1\rangle)\) and \(\operatorname{of\,List}\left(\langle \operatorname{length}(p)\rangle \mathbin{\Vert} \langle n\rangle\right)\) in place of its parameter \(v\) in reverse gives \(\operatorname{eval}(\operatorname{Lean\mathord{:}Omega\mathord{:}Linear\,Combo\mathord{:}mk}(0,\allowbreak{} \langle 0\rangle \mathbin{\Vert} \langle 1\rangle),\allowbreak{} \operatorname{of\,List}(\langle \operatorname{length}(p)\rangle \mathbin{\Vert} \langle n\rangle)) - \operatorname{eval}(\operatorname{Lean\mathord{:}Omega\mathord{:}Linear\,Combo\mathord{:}mk}(0,\allowbreak{} \langle 1\rangle),\allowbreak{} \operatorname{of\,List}(\langle \operatorname{length}(p)\rangle \mathbin{\Vert} \langle n\rangle)) = \operatorname{eval}(\operatorname{Lean\mathord{:}Omega\mathord{:}Linear\,Combo\mathord{:}mk}(0,\allowbreak{} \langle 0\rangle \mathbin{\Vert} \langle 1\rangle) - \operatorname{Lean\mathord{:}Omega\mathord{:}Linear\,Combo\mathord{:}mk}(0,\allowbreak{} \langle 1\rangle),\allowbreak{} \operatorname{of\,List}(\langle \operatorname{length}(p)\rangle \mathbin{\Vert} \langle n\rangle))\).
Chaining the equalities, we obtain \(n - \operatorname{length}(p) = \operatorname{eval}(\operatorname{Lean\mathord{:}Omega\mathord{:}Linear\,Combo\mathord{:}mk}(0,\allowbreak{} \langle 0\rangle \mathbin{\Vert} \langle 1\rangle) - \operatorname{Lean\mathord{:}Omega\mathord{:}Linear\,Combo\mathord{:}mk}(0,\allowbreak{} \langle 1\rangle),\allowbreak{} \operatorname{of\,List}(\langle \operatorname{length}(p)\rangle \mathbin{\Vert} \langle n\rangle))\).
By \paperlink{premise:6c655f6f665f6c655f6f665f6571}{Prerequisite~P\ref{premise:6c655f6f665f6c655f6f665f6571}} with \(\alpha = \mathbb{Z}\), we obtain \(0 \le \operatorname{eval}(\operatorname{Lean\mathord{:}Omega\mathord{:}Linear\,Combo\mathord{:}mk}(0,\allowbreak{} \langle 0\rangle \mathbin{\Vert} \langle 1\rangle) - \operatorname{Lean\mathord{:}Omega\mathord{:}Linear\,Combo\mathord{:}mk}(0,\allowbreak{} \langle 1\rangle),\allowbreak{} \operatorname{of\,List}(\langle \operatorname{length}(p)\rangle \mathbin{\Vert} \langle n\rangle))\).
By \paperlink{premise:4c65616e2e4f6d6567612e436f6e73747261696e742e616464496e657175616c6974795f736174}{Prerequisite~P\ref{premise:4c65616e2e4f6d6567612e436f6e73747261696e742e616464496e657175616c6974795f736174}}, we obtain \(\operatorname{sat^{\prime}}(\operatorname{Lean\mathord{:}Omega\mathord{:}Constraint\mathord{:}mk}(\operatorname{Option\mathord{:}some}(-0),\allowbreak{} \operatorname{none}),\allowbreak{} \langle -1\rangle \mathbin{\Vert} \langle 1\rangle,\allowbreak{} \operatorname{of\,List}(\langle \operatorname{length}(p)\rangle \mathbin{\Vert} \langle n\rangle)) = \mathrm{true}\).
By \paperlink{premise:4c65616e2e4f6d6567612e746964795f736174}{Prerequisite~P\ref{premise:4c65616e2e4f6d6567612e746964795f736174}}, we obtain \(\operatorname{sat^{\prime}}(\operatorname{tidy\,Constraint}(\operatorname{Lean\mathord{:}Omega\mathord{:}Constraint\mathord{:}mk}(\operatorname{Option\mathord{:}some}(0),\allowbreak{} \operatorname{none}),\allowbreak{} \langle -1\rangle \mathbin{\Vert} \langle 1\rangle),\allowbreak{} \operatorname{tidy\,Coeffs}(\operatorname{Lean\mathord{:}Omega\mathord{:}Constraint\mathord{:}mk}(\operatorname{Option\mathord{:}some}(0),\allowbreak{} \operatorname{none}),\allowbreak{} \langle -1\rangle \mathbin{\Vert} \langle 1\rangle),\allowbreak{} \operatorname{of\,List}(\langle \operatorname{length}(p)\rangle \mathbin{\Vert} \langle n\rangle)) = \mathrm{true}\).
Use the hypothesis \(n \le \operatorname{length}(p)\).
By \paperlink{premise:4c65616e2e4f6d6567612e496e742e6f664e61745f6c655f6f665f6c65}{Prerequisite~P\ref{premise:4c65616e2e4f6d6567612e496e742e6f664e61745f6c655f6f665f6c65}}, we obtain \(n \le \operatorname{length}(p)\).
By \paperlink{premise:496e742e7375625f6e6f6e6e65675f6f665f6c65}{Prerequisite~P\ref{premise:496e742e7375625f6e6f6e6e65675f6f665f6c65}}, we obtain \(0 \le \operatorname{length}(p) - n\).
Using \paperlink{premise:4c65616e2e4f6d6567612e4c696e656172436f6d626f2e636f6f7264696e6174655f6576616c5f30}{Prerequisite~P\ref{premise:4c65616e2e4f6d6567612e4c696e656172436f6d626f2e636f6f7264696e6174655f6576616c5f30}} in reverse gives \(\operatorname{length}(p) = \operatorname{eval}(\operatorname{coordinate}(0),\allowbreak{} \operatorname{of\,List}(\langle \operatorname{length}(p)\rangle \mathbin{\Vert} \langle n\rangle))\).
Using \paperlink{premise:4c65616e2e4f6d6567612e4c696e656172436f6d626f2e636f6f7264696e6174655f6576616c5f31}{Prerequisite~P\ref{premise:4c65616e2e4f6d6567612e4c696e656172436f6d626f2e636f6f7264696e6174655f6576616c5f31}} in reverse gives \(n = \operatorname{eval}(\operatorname{coordinate}(1),\allowbreak{} \operatorname{of\,List}(\langle \operatorname{length}(p)\rangle \mathbin{\Vert} \langle n\rangle))\).
By \paperlink{premise:4c65616e2e4f6d6567612e496e742e7375625f636f6e6772}{Prerequisite~P\ref{premise:4c65616e2e4f6d6567612e496e742e7375625f636f6e6772}}, we obtain \(\operatorname{length}(p) - n = \operatorname{eval}(\operatorname{coordinate}(0),\allowbreak{} \operatorname{of\,List}(\langle \operatorname{length}(p)\rangle \mathbin{\Vert} \langle n\rangle)) - \operatorname{eval}(\operatorname{coordinate}(1),\allowbreak{} \operatorname{of\,List}(\langle \operatorname{length}(p)\rangle \mathbin{\Vert} \langle n\rangle))\).
Using \paperlink{premise:4c65616e2e4f6d6567612e4c696e656172436f6d626f2e7375625f6576616c}{Prerequisite~P\ref{premise:4c65616e2e4f6d6567612e4c696e656172436f6d626f2e7375625f6576616c}} with \(l_{1} = \operatorname{Lean\mathord{:}Omega\mathord{:}Linear\,Combo\mathord{:}mk}(0,\allowbreak{} \langle 1\rangle)\) and \(l_{2} = \operatorname{Lean\mathord{:}Omega\mathord{:}Linear\,Combo\mathord{:}mk}(0,\allowbreak{} \langle 0\rangle \mathbin{\Vert} \langle 1\rangle)\) and \(\operatorname{of\,List}\left(\langle \operatorname{length}(p)\rangle \mathbin{\Vert} \langle n\rangle\right)\) in place of its parameter \(v\) in reverse gives \(\operatorname{eval}(\operatorname{Lean\mathord{:}Omega\mathord{:}Linear\,Combo\mathord{:}mk}(0,\allowbreak{} \langle 1\rangle),\allowbreak{} \operatorname{of\,List}(\langle \operatorname{length}(p)\rangle \mathbin{\Vert} \langle n\rangle)) - \operatorname{eval}(\operatorname{Lean\mathord{:}Omega\mathord{:}Linear\,Combo\mathord{:}mk}(0,\allowbreak{} \langle 0\rangle \mathbin{\Vert} \langle 1\rangle),\allowbreak{} \operatorname{of\,List}(\langle \operatorname{length}(p)\rangle \mathbin{\Vert} \langle n\rangle)) = \operatorname{eval}(\operatorname{Lean\mathord{:}Omega\mathord{:}Linear\,Combo\mathord{:}mk}(0,\allowbreak{} \langle 1\rangle) - \operatorname{Lean\mathord{:}Omega\mathord{:}Linear\,Combo\mathord{:}mk}(0,\allowbreak{} \langle 0\rangle \mathbin{\Vert} \langle 1\rangle),\allowbreak{} \operatorname{of\,List}(\langle \operatorname{length}(p)\rangle \mathbin{\Vert} \langle n\rangle))\).
Chaining the equalities, we obtain \(\operatorname{length}(p) - n = \operatorname{eval}(\operatorname{Lean\mathord{:}Omega\mathord{:}Linear\,Combo\mathord{:}mk}(0,\allowbreak{} \langle 1\rangle) - \operatorname{Lean\mathord{:}Omega\mathord{:}Linear\,Combo\mathord{:}mk}(0,\allowbreak{} \langle 0\rangle \mathbin{\Vert} \langle 1\rangle),\allowbreak{} \operatorname{of\,List}(\langle \operatorname{length}(p)\rangle \mathbin{\Vert} \langle n\rangle))\).
By \paperlink{premise:6c655f6f665f6c655f6f665f6571}{Prerequisite~P\ref{premise:6c655f6f665f6c655f6f665f6571}} with \(\alpha = \mathbb{Z}\), we obtain \(0 \le \operatorname{eval}(\operatorname{Lean\mathord{:}Omega\mathord{:}Linear\,Combo\mathord{:}mk}(0,\allowbreak{} \langle 1\rangle) - \operatorname{Lean\mathord{:}Omega\mathord{:}Linear\,Combo\mathord{:}mk}(0,\allowbreak{} \langle 0\rangle \mathbin{\Vert} \langle 1\rangle),\allowbreak{} \operatorname{of\,List}(\langle \operatorname{length}(p)\rangle \mathbin{\Vert} \langle n\rangle))\).
By \paperlink{premise:4c65616e2e4f6d6567612e436f6e73747261696e742e616464496e657175616c6974795f736174}{Prerequisite~P\ref{premise:4c65616e2e4f6d6567612e436f6e73747261696e742e616464496e657175616c6974795f736174}}, we obtain \(\operatorname{sat^{\prime}}(\operatorname{Lean\mathord{:}Omega\mathord{:}Constraint\mathord{:}mk}(\operatorname{Option\mathord{:}some}(-0),\allowbreak{} \operatorname{none}),\allowbreak{} \langle 1\rangle \mathbin{\Vert} \langle -1\rangle,\allowbreak{} \operatorname{of\,List}(\langle \operatorname{length}(p)\rangle \mathbin{\Vert} \langle n\rangle)) = \mathrm{true}\).
By \paperlink{premise:4c65616e2e4f6d6567612e436f6e73747261696e742e636f6d62696e655f73617427}{Prerequisite~P\ref{premise:4c65616e2e4f6d6567612e436f6e73747261696e742e636f6d62696e655f73617427}}, we obtain \(\operatorname{sat^{\prime}}(\operatorname{combine}(\operatorname{Lean\mathord{:}Omega\mathord{:}Constraint\mathord{:}mk}(\operatorname{none},\allowbreak{} \operatorname{Option\mathord{:}some}(0)),\allowbreak{} \operatorname{Lean\mathord{:}Omega\mathord{:}Constraint\mathord{:}mk}(\operatorname{Option\mathord{:}some}(0),\allowbreak{} \operatorname{none})),\allowbreak{} \langle 1\rangle \mathbin{\Vert} \langle -1\rangle,\allowbreak{} \operatorname{of\,List}(\langle \operatorname{length}(p)\rangle \mathbin{\Vert} \langle n\rangle)) = \mathrm{true}\).
Use the hypothesis \(n > \operatorname{length}(p)\).
By \paperlink{premise:4c65616e2e4f6d6567612e496e742e6f664e61745f6c745f6f665f6c74}{Prerequisite~P\ref{premise:4c65616e2e4f6d6567612e496e742e6f664e61745f6c745f6f665f6c74}}, we obtain \(\operatorname{length}(p) < n\).
By \paperlink{premise:496e742e6164645f6f6e655f6c655f6f665f6c74}{Prerequisite~P\ref{premise:496e742e6164645f6f6e655f6c655f6f665f6c74}}, we obtain \(\operatorname{length}(p) + 1 \le n\).
By \paperlink{premise:496e742e7375625f6e6f6e6e65675f6f665f6c65}{Prerequisite~P\ref{premise:496e742e7375625f6e6f6e6e65675f6f665f6c65}}, we obtain \(0 \le n - \left(\operatorname{length}(p) + 1\right)\).
Using \paperlink{premise:4c65616e2e4f6d6567612e4c696e656172436f6d626f2e636f6f7264696e6174655f6576616c5f31}{Prerequisite~P\ref{premise:4c65616e2e4f6d6567612e4c696e656172436f6d626f2e636f6f7264696e6174655f6576616c5f31}} in reverse gives \(n = \operatorname{eval}(\operatorname{coordinate}(1),\allowbreak{} \operatorname{of\,List}(\langle \operatorname{length}(p)\rangle \mathbin{\Vert} \langle n\rangle))\).
Using \paperlink{premise:4c65616e2e4f6d6567612e4c696e656172436f6d626f2e636f6f7264696e6174655f6576616c5f30}{Prerequisite~P\ref{premise:4c65616e2e4f6d6567612e4c696e656172436f6d626f2e636f6f7264696e6174655f6576616c5f30}} in reverse gives \(\operatorname{length}(p) = \operatorname{eval}(\operatorname{coordinate}(0),\allowbreak{} \operatorname{of\,List}(\langle \operatorname{length}(p)\rangle \mathbin{\Vert} \langle n\rangle))\).
By definition, we obtain \(1 = 1\).
By the preceding argument, we obtain \(1 = \operatorname{eval}(\operatorname{Lean\mathord{:}Omega\mathord{:}Linear\,Combo\mathord{:}mk}(1,\allowbreak{} \langle\,\rangle),\allowbreak{} \operatorname{of\,List}(\langle \operatorname{length}(p)\rangle \mathbin{\Vert} \langle n\rangle))\).
By \paperlink{premise:4c65616e2e4f6d6567612e496e742e6164645f636f6e6772}{Prerequisite~P\ref{premise:4c65616e2e4f6d6567612e496e742e6164645f636f6e6772}}, we obtain \(\operatorname{length}(p) + 1 = \operatorname{eval}(\operatorname{coordinate}(0),\allowbreak{} \operatorname{of\,List}(\langle \operatorname{length}(p)\rangle \mathbin{\Vert} \langle n\rangle)) + \operatorname{eval}(\operatorname{Lean\mathord{:}Omega\mathord{:}Linear\,Combo\mathord{:}mk}(1,\allowbreak{} \langle\,\rangle),\allowbreak{} \operatorname{of\,List}(\langle \operatorname{length}(p)\rangle \mathbin{\Vert} \langle n\rangle))\).
Using \paperlink{premise:4c65616e2e4f6d6567612e4c696e656172436f6d626f2e6164645f6576616c}{Prerequisite~P\ref{premise:4c65616e2e4f6d6567612e4c696e656172436f6d626f2e6164645f6576616c}} with \(l_{1} = \operatorname{Lean\mathord{:}Omega\mathord{:}Linear\,Combo\mathord{:}mk}(0,\allowbreak{} \langle 1\rangle)\) and \(l_{2} = \operatorname{Lean\mathord{:}Omega\mathord{:}Linear\,Combo\mathord{:}mk}(1,\allowbreak{} \langle\,\rangle)\) and \(\operatorname{of\,List}\left(\langle \operatorname{length}(p)\rangle \mathbin{\Vert} \langle n\rangle\right)\) in place of its parameter \(v\) in reverse gives \(\operatorname{eval}(\operatorname{Lean\mathord{:}Omega\mathord{:}Linear\,Combo\mathord{:}mk}(0,\allowbreak{} \langle 1\rangle),\allowbreak{} \operatorname{of\,List}(\langle \operatorname{length}(p)\rangle \mathbin{\Vert} \langle n\rangle)) + \operatorname{eval}(\operatorname{Lean\mathord{:}Omega\mathord{:}Linear\,Combo\mathord{:}mk}(1,\allowbreak{} \langle\,\rangle),\allowbreak{} \operatorname{of\,List}(\langle \operatorname{length}(p)\rangle \mathbin{\Vert} \langle n\rangle)) = \operatorname{eval}(\operatorname{Lean\mathord{:}Omega\mathord{:}Linear\,Combo\mathord{:}mk}(0,\allowbreak{} \langle 1\rangle) + \operatorname{Lean\mathord{:}Omega\mathord{:}Linear\,Combo\mathord{:}mk}(1,\allowbreak{} \langle\,\rangle),\allowbreak{} \operatorname{of\,List}(\langle \operatorname{length}(p)\rangle \mathbin{\Vert} \langle n\rangle))\).
Chaining the equalities, we obtain \(\operatorname{length}(p) + 1 = \operatorname{eval}(\operatorname{Lean\mathord{:}Omega\mathord{:}Linear\,Combo\mathord{:}mk}(0,\allowbreak{} \langle 1\rangle) + \operatorname{Lean\mathord{:}Omega\mathord{:}Linear\,Combo\mathord{:}mk}(1,\allowbreak{} \langle\,\rangle),\allowbreak{} \operatorname{of\,List}(\langle \operatorname{length}(p)\rangle \mathbin{\Vert} \langle n\rangle))\).
By \paperlink{premise:4c65616e2e4f6d6567612e496e742e7375625f636f6e6772}{Prerequisite~P\ref{premise:4c65616e2e4f6d6567612e496e742e7375625f636f6e6772}}, we obtain \(n - (\operatorname{length}(p) + 1) = \operatorname{eval}(\operatorname{coordinate}(1),\allowbreak{} \operatorname{of\,List}(\langle \operatorname{length}(p)\rangle \mathbin{\Vert} \langle n\rangle)) - \operatorname{eval}(\operatorname{Lean\mathord{:}Omega\mathord{:}Linear\,Combo\mathord{:}mk}(0,\allowbreak{} \langle 1\rangle) + \operatorname{Lean\mathord{:}Omega\mathord{:}Linear\,Combo\mathord{:}mk}(1,\allowbreak{} \langle\,\rangle),\allowbreak{} \operatorname{of\,List}(\langle \operatorname{length}(p)\rangle \mathbin{\Vert} \langle n\rangle))\).
Using \paperlink{premise:4c65616e2e4f6d6567612e4c696e656172436f6d626f2e7375625f6576616c}{Prerequisite~P\ref{premise:4c65616e2e4f6d6567612e4c696e656172436f6d626f2e7375625f6576616c}} with \(l_{1} = \operatorname{Lean\mathord{:}Omega\mathord{:}Linear\,Combo\mathord{:}mk}(0,\allowbreak{} \langle 0\rangle \mathbin{\Vert} \langle 1\rangle)\) and \(l_{2} = \operatorname{Lean\mathord{:}Omega\mathord{:}Linear\,Combo\mathord{:}mk}(1,\allowbreak{} \langle 1\rangle)\) and \(\operatorname{of\,List}\left(\langle \operatorname{length}(p)\rangle \mathbin{\Vert} \langle n\rangle\right)\) in place of its parameter \(v\) in reverse gives \(\operatorname{eval}(\operatorname{Lean\mathord{:}Omega\mathord{:}Linear\,Combo\mathord{:}mk}(0,\allowbreak{} \langle 0\rangle \mathbin{\Vert} \langle 1\rangle),\allowbreak{} \operatorname{of\,List}(\langle \operatorname{length}(p)\rangle \mathbin{\Vert} \langle n\rangle)) - \operatorname{eval}(\operatorname{Lean\mathord{:}Omega\mathord{:}Linear\,Combo\mathord{:}mk}(1,\allowbreak{} \langle 1\rangle),\allowbreak{} \operatorname{of\,List}(\langle \operatorname{length}(p)\rangle \mathbin{\Vert} \langle n\rangle)) = \operatorname{eval}(\operatorname{Lean\mathord{:}Omega\mathord{:}Linear\,Combo\mathord{:}mk}(0,\allowbreak{} \langle 0\rangle \mathbin{\Vert} \langle 1\rangle) - \operatorname{Lean\mathord{:}Omega\mathord{:}Linear\,Combo\mathord{:}mk}(1,\allowbreak{} \langle 1\rangle),\allowbreak{} \operatorname{of\,List}(\langle \operatorname{length}(p)\rangle \mathbin{\Vert} \langle n\rangle))\).
Chaining the equalities, we obtain \(n - (\operatorname{length}(p) + 1) = \operatorname{eval}(\operatorname{Lean\mathord{:}Omega\mathord{:}Linear\,Combo\mathord{:}mk}(0,\allowbreak{} \langle 0\rangle \mathbin{\Vert} \langle 1\rangle) - \operatorname{Lean\mathord{:}Omega\mathord{:}Linear\,Combo\mathord{:}mk}(1,\allowbreak{} \langle 1\rangle),\allowbreak{} \operatorname{of\,List}(\langle \operatorname{length}(p)\rangle \mathbin{\Vert} \langle n\rangle))\).
By \paperlink{premise:6c655f6f665f6c655f6f665f6571}{Prerequisite~P\ref{premise:6c655f6f665f6c655f6f665f6571}} with \(\alpha = \mathbb{Z}\), we obtain \(0 \le \operatorname{eval}(\operatorname{Lean\mathord{:}Omega\mathord{:}Linear\,Combo\mathord{:}mk}(0,\allowbreak{} \langle 0\rangle \mathbin{\Vert} \langle 1\rangle) - \operatorname{Lean\mathord{:}Omega\mathord{:}Linear\,Combo\mathord{:}mk}(1,\allowbreak{} \langle 1\rangle),\allowbreak{} \operatorname{of\,List}(\langle \operatorname{length}(p)\rangle \mathbin{\Vert} \langle n\rangle))\).
By \paperlink{premise:4c65616e2e4f6d6567612e436f6e73747261696e742e616464496e657175616c6974795f736174}{Prerequisite~P\ref{premise:4c65616e2e4f6d6567612e436f6e73747261696e742e616464496e657175616c6974795f736174}}, we obtain \(\operatorname{sat^{\prime}}(\operatorname{Lean\mathord{:}Omega\mathord{:}Constraint\mathord{:}mk}(\operatorname{Option\mathord{:}some}(-(-1)),\allowbreak{} \operatorname{none}),\allowbreak{} \langle -1\rangle \mathbin{\Vert} \langle 1\rangle,\allowbreak{} \operatorname{of\,List}(\langle \operatorname{length}(p)\rangle \mathbin{\Vert} \langle n\rangle)) = \mathrm{true}\).
By \paperlink{premise:4c65616e2e4f6d6567612e636f6d626f5f73617427}{Prerequisite~P\ref{premise:4c65616e2e4f6d6567612e636f6d626f5f73617427}} with \(s = \operatorname{Lean\mathord{:}Omega\mathord{:}Constraint\mathord{:}mk}(\operatorname{Option\mathord{:}some}(0),\allowbreak{} \operatorname{Option\mathord{:}some}(0))\) and \(t = \operatorname{Lean\mathord{:}Omega\mathord{:}Constraint\mathord{:}mk}(\operatorname{Option\mathord{:}some}(1),\allowbreak{} \operatorname{none})\) and \(1\) in place of its parameter \(a\) and \(x = \langle 1\rangle \mathbin{\Vert} \langle -1\rangle\) and \(b = 1\) and \(y = \langle -1\rangle \mathbin{\Vert} \langle 1\rangle\) and \(\operatorname{of\,List}\left(\langle \operatorname{length}(p)\rangle \mathbin{\Vert} \langle n\rangle\right)\) in place of its parameter \(v\), we obtain \(\operatorname{sat^{\prime}}(\operatorname{Lean\mathord{:}Omega\mathord{:}Constraint\mathord{:}combo}(1,\allowbreak{} \operatorname{Lean\mathord{:}Omega\mathord{:}Constraint\mathord{:}mk}(\operatorname{Option\mathord{:}some}(0),\allowbreak{} \operatorname{Option\mathord{:}some}(0)),\allowbreak{} 1,\allowbreak{} \operatorname{Lean\mathord{:}Omega\mathord{:}Constraint\mathord{:}mk}(\operatorname{Option\mathord{:}some}(1),\allowbreak{} \operatorname{none})),\allowbreak{} \operatorname{Lean\mathord{:}Omega\mathord{:}Coeffs\mathord{:}combo}(1,\allowbreak{} \langle 1\rangle \mathbin{\Vert} \langle -1\rangle,\allowbreak{} 1,\allowbreak{} \langle -1\rangle \mathbin{\Vert} \langle 1\rangle),\allowbreak{} \operatorname{of\,List}(\langle \operatorname{length}(p)\rangle \mathbin{\Vert} \langle n\rangle)) = \mathrm{true}\).
By \paperlink{premise:4c65616e2e4f6d6567612e746964795f736174}{Prerequisite~P\ref{premise:4c65616e2e4f6d6567612e746964795f736174}}, we obtain \(\operatorname{sat^{\prime}}(\operatorname{tidy\,Constraint}(\operatorname{Lean\mathord{:}Omega\mathord{:}Constraint\mathord{:}mk}(\operatorname{Option\mathord{:}some}(1),\allowbreak{} \operatorname{none}),\allowbreak{} \langle 0\rangle \mathbin{\Vert} \langle 0\rangle),\allowbreak{} \operatorname{tidy\,Coeffs}(\operatorname{Lean\mathord{:}Omega\mathord{:}Constraint\mathord{:}mk}(\operatorname{Option\mathord{:}some}(1),\allowbreak{} \operatorname{none}),\allowbreak{} \langle 0\rangle \mathbin{\Vert} \langle 0\rangle),\allowbreak{} \operatorname{of\,List}(\langle \operatorname{length}(p)\rangle \mathbin{\Vert} \langle n\rangle)) = \mathrm{true}\).
By \paperlink{premise:4c65616e2e4f6d6567612e436f6e73747261696e742e6e6f745f736174275f6f665f6973496d706f737369626c65}{Prerequisite~P\ref{premise:4c65616e2e4f6d6567612e436f6e73747261696e742e6e6f745f736174275f6f665f6973496d706f737369626c65}}, we obtain \(\mathrm{False}\).
\end{proofstep}
The two alternatives exhaust the cases. Hence \(\mathrm{False}\).
\end{proofstep}
Rewriting with \(\text{every nonisolated vertex of }G\text{ has degree two}\), \paperlink{premise:5365742e6e6f6e656d7074795f6f665f6d656d}{Prerequisite~P\ref{premise:5365742e6e6f6e656d7074795f6f665f6d656d}}, \paperlink{premise:53696d706c6547726170682e41646a2e73796d6d}{Prerequisite~P\ref{premise:53696d706c6547726170682e41646a2e73796d6d}} (using \(v \sim_{G} \operatorname{get\,Vert}\left(p, n\right)\)), \paperlink{premise:53696d706c6547726170682e57616c6b2e4973506174682e6e636172645f6e65696768626f725365745f746f53756267726170685f696e7465726e616c5f65715f74776f}{Prerequisite~P\ref{premise:53696d706c6547726170682e57616c6b2e4973506174682e6e636172645f6e65696768626f725365745f746f53756267726170685f696e7465726e616c5f65715f74776f}} (using \(p\text{ has no repeated vertices}\)), \paperlink{premise:446563696461626c652e6279436f6e74726164696374696f6e}{Prerequisite~P\ref{premise:446563696461626c652e6279436f6e74726164696374696f6e}}, \paperlink{premise:5365742e636173745f6e63617264}{Prerequisite~P\ref{premise:5365742e636173745f6e63617264}}, \paperlink{premise:5365742e746f46696e697465}{Prerequisite~P\ref{premise:5365742e746f46696e697465}}, \paperlink{premise:537562747970652e66696e697465}{Prerequisite~P\ref{premise:537562747970652e66696e697465}} (using \(\operatorname{Finite}\left(V\right)\)), \paperlink{premise:43617264696e616c2e6571}{Prerequisite~P\ref{premise:43617264696e616c2e6571}} and the two auxiliary results established above gives \(\operatorname{Bij}(N_{G}(\operatorname{get\,Vert}(p,\allowbreak{} n)),\allowbreak{} \operatorname{Simple\,Graph\mathord{:}Subgraph\mathord{:}neighbor\,Set}(\operatorname{Simple\,Graph\mathord{:}Walk\mathord{:}to\,Subgraph}(p),\allowbreak{} \operatorname{get\,Vert}(p,\allowbreak{} n)))\text{ is nonempty}\).
Use the hypothesis \(v \sim_{G} \operatorname{get\,Vert}\left(p, n\right)\).
By \paperlink{premise:53696d706c6547726170682e41646a2e73796d6d}{Prerequisite~P\ref{premise:53696d706c6547726170682e41646a2e73796d6d}}, we obtain \(\operatorname{get\,Vert}\left(p, n\right) \sim_{G} v\).
Rewriting with \paperlink{premise:53696d706c6547726170682e53756267726170682e61646a5f6966665f6f665f6e65696768626f725365745f6571756976}{Prerequisite~P\ref{premise:53696d706c6547726170682e53756267726170682e61646a5f6966665f6f665f6e65696768626f725365745f6571756976}}, \paperlink{premise:5365742e746f46696e697465}{Prerequisite~P\ref{premise:5365742e746f46696e697465}} and \paperlink{premise:537562747970652e66696e697465}{Prerequisite~P\ref{premise:537562747970652e66696e697465}} (using \(\operatorname{Finite}\left(V\right)\)) in the preceding result gives \(\operatorname{get\,Vert}(p,\allowbreak{} n) \sim_{\operatorname{Simple\,Graph\mathord{:}Walk\mathord{:}to\,Subgraph}(p)} v\).
Rewriting with \paperlink{premise:53696d706c6547726170682e53756267726170682e61646a5f636f6d6d}{Prerequisite~P\ref{premise:53696d706c6547726170682e53756267726170682e61646a5f636f6d6d}} in the preceding result gives \(v \sim_{\operatorname{Simple\,Graph\mathord{:}Walk\mathord{:}to\,Subgraph}(p)} \operatorname{get\,Vert}(p,\allowbreak{} n)\).
\end{proofstep}
The proposition or its negation holds. Hence \(v \sim_{\operatorname{Simple\,Graph\mathord{:}Walk\mathord{:}to\,Subgraph}(p)} \operatorname{get\,Vert}(p,\allowbreak{} n)\).
\end{proofstep}
Use the hypothesis \(\operatorname{get\,Vert}\left(p, n\right) = w^{'}\).
Use the hypothesis \(w \ne w^{'}\).
Use the hypothesis \(v \sim_{G} w^{'}\).
Substituting equals, we obtain \(v \sim_{\operatorname{Simple\,Graph\mathord{:}Walk\mathord{:}to\,Subgraph}\left(p\right)} w^{'}\).
\end{proofstep}
Using the two conjuncts, we obtain \(v \sim_{\operatorname{Simple\,Graph\mathord{:}Walk\mathord{:}to\,Subgraph}\left(p\right)} w^{'}\).
\end{proofstep}
Using the existential hypothesis, we obtain \(v \sim_{\operatorname{Simple\,Graph\mathord{:}Walk\mathord{:}to\,Subgraph}\left(p\right)} w^{'}\).
By \paperlink{premise:53696d706c6547726170682e57616c6b2e4973506174682e736e645f6f665f746f53756267726170685f61646a}{Prerequisite~P\ref{premise:53696d706c6547726170682e57616c6b2e4973506174682e736e645f6f665f746f53756267726170685f61646a}}, we obtain \(\operatorname{Simple\,Graph\mathord{:}Walk\mathord{:}snd}\left(p\right) = w^{'}\).
\end{proof}

\begin{lemma}[{Reachable sdiff to subgraph spanning coercion}]\label{theorem:38}
Let \(V\) be a type. Let \(G\) be a simple graph on \(V\). Assume \(V\) is finite. Let \(v\) and \(w\) be elements of \(V\). Assume \(\text{every nonisolated vertex of }G\text{ has degree two}\). Let \(p\) be a walk in \(G\) from \(v\) to \(w\). Assume \(p\text{ has no repeated vertices}\). Then
\begin{center}
\(\displaystyle w\text{ can reach }v\text{ in }G \setminus \operatorname{spanning\,graph}(\operatorname{Simple\,Graph\mathord{:}Walk\mathord{:}to\,Subgraph}(p)).\)
\end{center}
\end{lemma}

\begin{proof}
Set \(\mathit{this}\) equal to \(\operatorname{of\,Finite}\left(V\right)\).
Use the hypothesis \(\operatorname{Finite}\left(V\right)\).
A descending chain would make the natural-number measure \(x \mapsto \left|V\right| + 1 - \operatorname{length}(\pi_1(\pi_2(x)))\) strictly decrease at each step, so the relation is well-founded.
Proceed by well-founded induction. For an arbitrary element we may assume the claim for each strictly smaller element.
\begin{proofstep}
Fix \(x\) in \(\{(v_{1},\allowbreak{} y)\mid v_{1}\, :\, V,\; y\, :\, \{(p_{1}, y)\mid p_{1}\, :\, \operatorname{Walk}_{G}(v_{1},\allowbreak{} w),\; y\text{ proves }p_{1}\text{ has no repeated vertices}\}\}\).
Assume \(\forall y\, :\, \{(v_{1}, y)\mid v_{1}\, :\, V,\; y\, :\, \{(p_{1}, y)\mid p_{1}\, :\, \operatorname{Walk}_{G}(v_{1},\allowbreak{} w),\; y\text{ proves }p_{1}\text{ has no repeated vertices}\}\},\; |V| + 1 - \operatorname{length}(\pi_1(\pi_2(y))) < |V| + 1 - \operatorname{length}(\pi_1(\pi_2(x))) \implies w\text{ can reach }\pi_1(y)\text{ in }G \setminus \operatorname{spanning\,graph}(\operatorname{Simple\,Graph\mathord{:}Walk\mathord{:}to\,Subgraph}(\pi_1(\pi_2(y))))\).
Use the hypothesis \(\forall y\, :\, \{(v_{1}, y)\mid v_{1}\, :\, V,\; y\, :\, \{(p_{1}, y)\mid p_{1}\, :\, \operatorname{Walk}_{G}(v_{1},\allowbreak{} w),\; y\text{ proves }p_{1}\text{ has no repeated vertices}\}\},\; |V| + 1 - \operatorname{length}(\pi_1(\pi_2(y))) < |V| + 1 - \operatorname{length}(\pi_1(\pi_2(x))) \implies w\text{ can reach }\pi_1(y)\text{ in }G \setminus \operatorname{spanning\,graph}(\operatorname{Simple\,Graph\mathord{:}Walk\mathord{:}to\,Subgraph}(\pi_1(\pi_2(y))))\).
Write the pair \(x\) in terms of its two components.
\begin{proofstep}
Fix \(v_{1}\) in \(V\).
Fix \(p_{1}\) in \(\{(p_{2},\allowbreak{} h)\mid p_{2}\, :\, \operatorname{Walk}_{G}(v_{1},\allowbreak{} w),\; h\text{ proves }p_{2}\text{ has no repeated vertices}\}\).
Assume \(\forall y\, :\, \{(v_{2}, y)\mid v_{2}\, :\, V,\; y\, :\, \{(p_{2}, y)\mid p_{2}\, :\, \operatorname{Walk}_{G}(v_{2},\allowbreak{} w),\; y\text{ proves }p_{2}\text{ has no repeated vertices}\}\},\; |V| + 1 - \operatorname{length}(\pi_1(\pi_2(y))) < |V| + 1 - \operatorname{length}(\pi_1(p_{1})) \implies w\text{ can reach }\pi_1(y)\text{ in }G \setminus \operatorname{spanning\,graph}(\operatorname{Simple\,Graph\mathord{:}Walk\mathord{:}to\,Subgraph}(\pi_1(\pi_2(y))))\).
Use the hypothesis \(\forall y\, :\, \{(v_{2}, y)\mid v_{2}\, :\, V,\; y\, :\, \{(p_{2}, y)\mid p_{2}\, :\, \operatorname{Walk}_{G}(v_{2},\allowbreak{} w),\; y\text{ proves }p_{2}\text{ has no repeated vertices}\}\},\; |V| + 1 - \operatorname{length}(\pi_1(\pi_2(y))) < |V| + 1 - \operatorname{length}(\pi_1(p_{1})) \implies w\text{ can reach }\pi_1(y)\text{ in }G \setminus \operatorname{spanning\,graph}(\operatorname{Simple\,Graph\mathord{:}Walk\mathord{:}to\,Subgraph}(\pi_1(\pi_2(y))))\).
Write the pair \(p_{1}\) in terms of its two components.
\begin{proofstep}
Fix \(p_{2}\) in \(\operatorname{Walk}_{G}(v_{1}, w)\).
Assume \(p_{2}\text{ has no repeated vertices}\).
Assume \(\forall y\, :\, \{(v_{2}, y)\mid v_{2}\, :\, V,\; y\, :\, \{(p_{3}, y)\mid p_{3}\, :\, \operatorname{Walk}_{G}(v_{2},\allowbreak{} w),\; y\text{ proves }p_{3}\text{ has no repeated vertices}\}\},\; |V| + 1 - \operatorname{length}(\pi_1(\pi_2(y))) < |V| + 1 - \operatorname{length}(\pi_1(\langle p_{2},\allowbreak{} \mathit{hp}_{1}\rangle)) \implies w\text{ can reach }\pi_1(y)\text{ in }G \setminus \operatorname{spanning\,graph}(\operatorname{Simple\,Graph\mathord{:}Walk\mathord{:}to\,Subgraph}(\pi_1(\pi_2(y))))\).
Distinguish whether \(v_{1} = w\) holds.
Consider the following case.
\begin{proofstep}
Assume \(v_{1} = w\).
For the required implication, argue as follows.
\begin{proofstep}
Fix \(p_{3}\) in \(\operatorname{Walk}_{G}(v_{1}, v_{1})\).
Assume \(p_{3}\text{ has no repeated vertices}\).
Choose the witness \(\operatorname{Simple\,Graph\mathord{:}Walk\mathord{:}nil}\).
This element witnesses nonemptiness. Hence \(\operatorname{Walk}_{G \setminus \operatorname{spanning\,graph}(\operatorname{Simple\,Graph\mathord{:}Walk\mathord{:}to\,Subgraph}(p_{3}))}(v_{1},\allowbreak{} v_{1})\text{ is nonempty}\).
\end{proofstep}
Use the hypothesis \(v_{1} = w\).
Use the hypothesis \(p_{2}\text{ has no repeated vertices}\).
Substituting equals, we obtain \(w\text{ can reach }v_{1}\text{ in }G \setminus \operatorname{spanning\,graph}(\operatorname{Simple\,Graph\mathord{:}Walk\mathord{:}to\,Subgraph}(p_{2}))\).
\end{proofstep}
Consider the following case.
\begin{proofstep}
Assume \(\neg \left(v_{1} = w\right)\).
Use the hypothesis \(\neg \left(v_{1} = w\right)\).
By \paperlink{premise:53696d706c6547726170682e57616c6b2e6e6f745f6e696c5f6f665f6e65}{Prerequisite~P\ref{premise:53696d706c6547726170682e57616c6b2e6e6f745f6e696c5f6f665f6e65}}, we obtain \(\neg \operatorname{Nil}\left(p_{2}\right)\).
Use the previously proved fact \(\neg \operatorname{Nil}\left(p_{2}\right)\).
By \paperlink{premise:53696d706c6547726170682e57616c6b2e746f53756267726170685f61646a5f736e64}{Prerequisite~P\ref{premise:53696d706c6547726170682e57616c6b2e746f53756267726170685f61646a5f736e64}} with \(p_{2}\) in place of its parameter \(w\), we obtain \(v_{1} \sim_{\operatorname{Simple\,Graph\mathord{:}Walk\mathord{:}to\,Subgraph}(p_{2})} \operatorname{Simple\,Graph\mathord{:}Walk\mathord{:}snd}(p_{2})\).
By \paperlink{premise:53696d706c6547726170682e53756267726170682e41646a2e61646a5f737562}{Prerequisite~P\ref{premise:53696d706c6547726170682e53756267726170682e41646a2e61646a5f737562}}, we obtain \(v_{1} \sim_{G} \operatorname{Simple\,Graph\mathord{:}Walk\mathord{:}snd}\left(p_{2}\right)\).
By \paperlink{theorem:31}{Lemma~\ref{theorem:31}}, we obtain \(\exists! w^{'}\, :\, V,\; \operatorname{Simple\,Graph\mathord{:}Walk\mathord{:}snd}(p_{2}) \ne w^{'} \land v_{1} \sim_{G} w^{'}\).
Consider the following case.
\begin{proofstep}
Fix \(w^{'}\) in \(V\).
Assume \((\operatorname{Simple\,Graph\mathord{:}Walk\mathord{:}snd}(p_{2}) \ne w^{'} \land v_{1} \sim_{G} w^{'}) \land (\forall y\, :\, V,\; \operatorname{Simple\,Graph\mathord{:}Walk\mathord{:}snd}(p_{2}) \ne y \land v_{1} \sim_{G} y \implies y = w^{'})\).
Use the hypothesis \((\operatorname{Simple\,Graph\mathord{:}Walk\mathord{:}snd}(p_{2}) \ne w^{'} \land v_{1} \sim_{G} w^{'}) \land (\forall y\, :\, V,\; \operatorname{Simple\,Graph\mathord{:}Walk\mathord{:}snd}(p_{2}) \ne y \land v_{1} \sim_{G} y \implies y = w^{'})\).
Consider the following case.
\begin{proofstep}
Assume \(\operatorname{Simple\,Graph\mathord{:}Walk\mathord{:}snd}(p_{2}) \ne w^{'} \land v_{1} \sim_{G} w^{'}\).
Assume \(\forall y\, :\, V,\; \operatorname{Simple\,Graph\mathord{:}Walk\mathord{:}snd}(p_{2}) \ne y \land v_{1} \sim_{G} y \implies y = w^{'}\).
Use the hypothesis \(\operatorname{Simple\,Graph\mathord{:}Walk\mathord{:}snd}(p_{2}) \ne w^{'} \land v_{1} \sim_{G} w^{'}\).
Consider the following case.
\begin{proofstep}
Assume \(\operatorname{Simple\,Graph\mathord{:}Walk\mathord{:}snd}\left(p_{2}\right) \ne w^{'}\).
Assume \(v_{1} \sim_{G} w^{'}\).
Assume \(v_{1} \sim_{\operatorname{Simple\,Graph\mathord{:}Walk\mathord{:}to\,Subgraph}\left(p_{2}\right)} w^{'}\).
Use the hypothesis \(p_{2}\text{ has no repeated vertices}\).
Use the hypothesis \(v_{1} \sim_{\operatorname{Simple\,Graph\mathord{:}Walk\mathord{:}to\,Subgraph}\left(p_{2}\right)} w^{'}\).
By \paperlink{premise:53696d706c6547726170682e57616c6b2e4973506174682e736e645f6f665f746f53756267726170685f61646a}{Prerequisite~P\ref{premise:53696d706c6547726170682e57616c6b2e4973506174682e736e645f6f665f746f53756267726170685f61646a}}, we obtain \(\operatorname{Simple\,Graph\mathord{:}Walk\mathord{:}snd}\left(p_{2}\right) = w^{'}\).
Applying the hypothesis \(\operatorname{Simple\,Graph\mathord{:}Walk\mathord{:}snd}\left(p_{2}\right) \ne w^{'}\), we obtain \(\mathrm{False}\).
Distinguish whether \(w = w^{'}\) holds.
Consider the following case.
\begin{proofstep}
Assume \(w = w^{'}\).
For the required implication, argue as follows.
\begin{proofstep}
Assume \(\forall y\, :\, V,\; \operatorname{Simple\,Graph\mathord{:}Walk\mathord{:}snd}(p_{2}) \ne y \land v_{1} \sim_{G} y \implies y = w\).
Assume \(\operatorname{Simple\,Graph\mathord{:}Walk\mathord{:}snd}\left(p_{2}\right) \ne w\).
Assume \(v_{1} \sim_{G} w\).
Assume \(\neg (v_{1} \sim_{\operatorname{Simple\,Graph\mathord{:}Walk\mathord{:}to\,Subgraph}(p_{2})} w)\).
Use the hypothesis \(v_{1} \sim_{G} w\).
By \paperlink{premise:53696d706c6547726170682e41646a2e73796d6d}{Prerequisite~P\ref{premise:53696d706c6547726170682e41646a2e73796d6d}}, we obtain \(w \sim_{G} v_{1}\).
For the required implication, argue as follows.
\begin{proofstep}
Assume \(w \sim_{\operatorname{Simple\,Graph\mathord{:}Walk\mathord{:}to\,Subgraph}\left(p_{2}\right)} v_{1}\).
Use the hypothesis \(w \sim_{\operatorname{Simple\,Graph\mathord{:}Walk\mathord{:}to\,Subgraph}\left(p_{2}\right)} v_{1}\).
By \paperlink{premise:53696d706c6547726170682e53756267726170682e41646a2e73796d6d}{Prerequisite~P\ref{premise:53696d706c6547726170682e53756267726170682e41646a2e73796d6d}}, we obtain \(v_{1} \sim_{\operatorname{Simple\,Graph\mathord{:}Walk\mathord{:}to\,Subgraph}\left(p_{2}\right)} w\).
Applying the hypothesis \(\neg (v_{1} \sim_{\operatorname{Simple\,Graph\mathord{:}Walk\mathord{:}to\,Subgraph}(p_{2})} w)\), we obtain \(\mathrm{False}\).
\end{proofstep}
Combining the two conclusions, we obtain \(w \sim_{G} v_{1} \land \neg (w \sim_{\operatorname{Simple\,Graph\mathord{:}Walk\mathord{:}to\,Subgraph}(p_{2})} v_{1})\).
Rewriting with \paperlink{premise:53696d706c6547726170682e73646966665f61646a}{Prerequisite~P\ref{premise:53696d706c6547726170682e73646966665f61646a}} in the preceding result gives \(w \sim_{G \setminus \operatorname{spanning\,graph}(\operatorname{Simple\,Graph\mathord{:}Walk\mathord{:}to\,Subgraph}(p_{2}))} v_{1}\).
Use the previously proved fact \(w \sim_{G \setminus \operatorname{spanning\,graph}(\operatorname{Simple\,Graph\mathord{:}Walk\mathord{:}to\,Subgraph}(p_{2}))} v_{1}\).
By \paperlink{premise:53696d706c6547726170682e41646a2e726561636861626c65}{Prerequisite~P\ref{premise:53696d706c6547726170682e41646a2e726561636861626c65}}, we obtain \(w\text{ can reach }v_{1}\text{ in }G \setminus \operatorname{spanning\,graph}(\operatorname{Simple\,Graph\mathord{:}Walk\mathord{:}to\,Subgraph}(p_{2}))\).
\end{proofstep}
Use the hypothesis \(w = w^{'}\).
Use the hypothesis \(\forall y\, :\, V,\; \operatorname{Simple\,Graph\mathord{:}Walk\mathord{:}snd}(p_{2}) \ne y \land v_{1} \sim_{G} y \implies y = w^{'}\).
Use the hypothesis \(\operatorname{Simple\,Graph\mathord{:}Walk\mathord{:}snd}\left(p_{2}\right) \ne w^{'}\).
Use the hypothesis \(v_{1} \sim_{G} w^{'}\).
Use the previously proved fact \(\neg (v_{1} \sim_{\operatorname{Simple\,Graph\mathord{:}Walk\mathord{:}to\,Subgraph}(p_{2})} w^{'})\).
Substituting equals, we obtain \(w\text{ can reach }v_{1}\text{ in }G \setminus \operatorname{spanning\,graph}(\operatorname{Simple\,Graph\mathord{:}Walk\mathord{:}to\,Subgraph}(p_{2}))\).
\end{proofstep}
Consider the following case.
\begin{proofstep}
Assume \(\neg \left(w = w^{'}\right)\).
By \paperlink{premise:6c655f7265666c}{Prerequisite~P\ref{premise:6c655f7265666c}} with \(\alpha = \operatorname{Simple\,Graph}(V)\) and \(G\) in place of its parameter \(a\), we obtain \(G \subseteq G\).
Rewriting with \paperlink{premise:53696d706c6547726170682e53756267726170682e7370616e6e696e67436f655f73756267726170684f6641646a}{Prerequisite~P\ref{premise:53696d706c6547726170682e53756267726170682e7370616e6e696e67436f655f73756267726170684f6641646a}}, \paperlink{premise:53696d706c6547726170682e41646a2e73796d6d}{Prerequisite~P\ref{premise:53696d706c6547726170682e41646a2e73796d6d}} (using \(v_{1} \sim_{G} w^{'}\)) and \paperlink{premise:6c655f7375705f7269676874}{Prerequisite~P\ref{premise:6c655f7375705f7269676874}} gives \(\operatorname{spanning\,graph}(\operatorname{Simple\,Graph\mathord{:}Walk\mathord{:}to\,Subgraph}(p_{2})) \le \operatorname{spanning\,graph}(\operatorname{subgraph\,Of\,Adj}(G)) \cup \operatorname{spanning\,graph}(\operatorname{Simple\,Graph\mathord{:}Walk\mathord{:}to\,Subgraph}(p_{2}))\).
By \paperlink{premise:73646966665f6c655f7364696666}{Prerequisite~P\ref{premise:73646966665f6c655f7364696666}} with \(\alpha = \operatorname{Simple\,Graph}(V)\), we obtain \(G \setminus \operatorname{spanning\,graph}(\operatorname{Simple\,Graph\mathord{:}Walk\mathord{:}to\,Subgraph}(\operatorname{Simple\,Graph\mathord{:}Walk\mathord{:}cons}(p_{2}))) \subseteq G \setminus \operatorname{spanning\,graph}(\operatorname{Simple\,Graph\mathord{:}Walk\mathord{:}to\,Subgraph}(p_{2}))\).
Use the hypothesis \(p_{2}\text{ has no repeated vertices}\).
For the required implication, argue as follows.
\begin{proofstep}
Assume \(p_{2}\text{ visits }w^{'}\).
By \paperlink{premise:46696e6974652e6f665f66696e74797065}{Prerequisite~P\ref{premise:46696e6974652e6f665f66696e74797065}} with \(\alpha = V\), we obtain \(\operatorname{Finite}\left(V\right)\).
By \paperlink{theorem:37}{Lemma~\ref{theorem:37}} with \(p_{2}\) in place of its parameter \(p\), we obtain \(\operatorname{Simple\,Graph\mathord{:}Walk\mathord{:}snd}\left(p_{2}\right) = w^{'}\).
Applying the hypothesis \(\operatorname{Simple\,Graph\mathord{:}Walk\mathord{:}snd}\left(p_{2}\right) \ne w^{'}\), we obtain \(\mathrm{False}\).
\end{proofstep}
Combining the two conclusions, we obtain \(p_{2}\text{ has no repeated vertices} \land w^{'} \notin \operatorname{vertex\,sequence}(p_{2})\).
Rewriting with \paperlink{premise:53696d706c6547726170682e57616c6b2e636f6e735f6973506174685f696666}{Prerequisite~P\ref{premise:53696d706c6547726170682e57616c6b2e636f6e735f6973506174685f696666}} and \paperlink{premise:53696d706c6547726170682e41646a2e73796d6d}{Prerequisite~P\ref{premise:53696d706c6547726170682e41646a2e73796d6d}} (using \(v_{1} \sim_{G} w^{'}\)) in the preceding result gives \(\operatorname{Simple\,Graph\mathord{:}Walk\mathord{:}cons}\left(p_{2}\right)\text{ has no repeated vertices}\).
Use the hypothesis \(v_{1} \sim_{G} w^{'}\).
By \paperlink{premise:53696d706c6547726170682e41646a2e73796d6d}{Prerequisite~P\ref{premise:53696d706c6547726170682e41646a2e73796d6d}}, we obtain \(w^{'} \sim_{G} v_{1}\).
For the required implication, argue as follows.
\begin{proofstep}
Assume \(w^{'} \sim_{\operatorname{Simple\,Graph\mathord{:}Walk\mathord{:}to\,Subgraph}\left(p_{2}\right)} v_{1}\).
Use the hypothesis \(w^{'} \sim_{\operatorname{Simple\,Graph\mathord{:}Walk\mathord{:}to\,Subgraph}\left(p_{2}\right)} v_{1}\).
By \paperlink{premise:53696d706c6547726170682e53756267726170682e41646a2e73796d6d}{Prerequisite~P\ref{premise:53696d706c6547726170682e53756267726170682e41646a2e73796d6d}}, we obtain \(v_{1} \sim_{\operatorname{Simple\,Graph\mathord{:}Walk\mathord{:}to\,Subgraph}\left(p_{2}\right)} w^{'}\).
Applying the hypothesis \(\neg (v_{1} \sim_{\operatorname{Simple\,Graph\mathord{:}Walk\mathord{:}to\,Subgraph}(p_{2})} w^{'})\), we obtain \(\mathrm{False}\).
\end{proofstep}
Combining the two conclusions, we obtain \(w^{'} \sim_{G} v_{1} \land \neg (w^{'} \sim_{\operatorname{Simple\,Graph\mathord{:}Walk\mathord{:}to\,Subgraph}(p_{2})} v_{1})\).
Rewriting with \paperlink{premise:53696d706c6547726170682e73646966665f61646a}{Prerequisite~P\ref{premise:53696d706c6547726170682e73646966665f61646a}} in the preceding result gives \(w^{'} \sim_{G \setminus \operatorname{spanning\,graph}(\operatorname{Simple\,Graph\mathord{:}Walk\mathord{:}to\,Subgraph}(p_{2}))} v_{1}\).
Choose the witness \(\operatorname{map\,Le}(\operatorname{Nonempty\mathord{:}some}) \mathbin{\Vert} \operatorname{to\,Walk}\).
This element witnesses nonemptiness. Hence \(w\text{ can reach }v_{1}\text{ in }G \setminus \operatorname{spanning\,graph}(\operatorname{Simple\,Graph\mathord{:}Walk\mathord{:}to\,Subgraph}(p_{2}))\).
\end{proofstep}
The proposition or its negation holds. Hence \(w\text{ can reach }v_{1}\text{ in }G \setminus \operatorname{spanning\,graph}(\operatorname{Simple\,Graph\mathord{:}Walk\mathord{:}to\,Subgraph}(p_{2}))\).
\end{proofstep}
Using the two conjuncts, we obtain \(w\text{ can reach }v_{1}\text{ in }G \setminus \operatorname{spanning\,graph}(\operatorname{Simple\,Graph\mathord{:}Walk\mathord{:}to\,Subgraph}(p_{2}))\).
\end{proofstep}
Using the two conjuncts, we obtain \(w\text{ can reach }v_{1}\text{ in }G \setminus \operatorname{spanning\,graph}(\operatorname{Simple\,Graph\mathord{:}Walk\mathord{:}to\,Subgraph}(p_{2}))\).
\end{proofstep}
Using the existential hypothesis, we obtain \(w\text{ can reach }v_{1}\text{ in }G \setminus \operatorname{spanning\,graph}(\operatorname{Simple\,Graph\mathord{:}Walk\mathord{:}to\,Subgraph}(p_{2}))\).
\end{proofstep}
The proposition or its negation holds. Hence \(w\text{ can reach }v_{1}\text{ in }G \setminus \operatorname{spanning\,graph}(\operatorname{Simple\,Graph\mathord{:}Walk\mathord{:}to\,Subgraph}(p_{2}))\).
\end{proofstep}
This proves \(w\text{ can reach }v_{1}\text{ in }G \setminus \operatorname{spanning\,graph}(\operatorname{Simple\,Graph\mathord{:}Walk\mathord{:}to\,Subgraph}(\pi_1(p_{1})))\).
\end{proofstep}
This proves \(w\text{ can reach }\pi_1(x)\text{ in }G \setminus \operatorname{spanning\,graph}(\operatorname{Simple\,Graph\mathord{:}Walk\mathord{:}to\,Subgraph}(\pi_1(\pi_2(x))))\).
\end{proofstep}
The induction gives \(w\text{ can reach }v\text{ in }G \setminus \operatorname{spanning\,graph}(\operatorname{Simple\,Graph\mathord{:}Walk\mathord{:}to\,Subgraph}(p))\).
\end{proof}

\begin{lemma}[{Reachable delete edges}]\label{theorem:39}
Let \(V\) be a type. Let \(G\) be a simple graph on \(V\). Let \(v\) and \(w\) be elements of \(V\). Assume \(V\) is finite. Assume \(v \sim_{G} w\). Assume \(\text{every nonisolated vertex of }G\text{ has degree two}\). Then
\begin{center}
\(\displaystyle v\text{ can reach }w\text{ in }\operatorname{delete\,Edges}(G,\allowbreak{} \{\operatorname{Sym2\mathord{:}mk}((v,\allowbreak{} w))\}).\)
\end{center}
\end{lemma}

\begin{proof}
Using \paperlink{premise:53696d706c6547726170682e53756267726170682e7370616e6e696e67436f655f73756267726170684f6641646a}{Prerequisite~P\ref{premise:53696d706c6547726170682e53756267726170682e7370616e6e696e67436f655f73756267726170684f6641646a}} in reverse under \(v \sim_{G} w\) gives \(\operatorname{from\,Edge\,Set}(\{\operatorname{Sym2\mathord{:}mk}((v,\allowbreak{} w))\}) = \operatorname{spanning\,graph}(\operatorname{subgraph\,Of\,Adj}(G))\).
Rewriting with \paperlink{premise:7375705f6f665f6c655f6c656674}{Prerequisite~P\ref{premise:7375705f6f665f6c655f6c656674}}, \paperlink{premise:53696d706c6547726170682e73696e676c65746f6e53756267726170685f6c655f696666}{Prerequisite~P\ref{premise:53696d706c6547726170682e73696e676c65746f6e53756267726170685f6c655f696666}}, \paperlink{premise:5365742e6d656d5f696e736572745f696666}{Prerequisite~P\ref{premise:5365742e6d656d5f696e736572745f696666}}, \paperlink{premise:5365742e6d656d5f73696e676c65746f6e5f696666}{Prerequisite~P\ref{premise:5365742e6d656d5f73696e676c65746f6e5f696666}} and \paperlink{premise:6f725f74727565}{Prerequisite~P\ref{premise:6f725f74727565}} in the preceding result gives \(\operatorname{from\,Edge\,Set}(\{\operatorname{Sym2\mathord{:}mk}((v,\allowbreak{} w))\}) = \operatorname{spanning\,graph}(\operatorname{Simple\,Graph\mathord{:}Walk\mathord{:}to\,Subgraph}(\operatorname{to\,Walk}))\).
The substitution uses the following argument.
\begin{proofstep}
By definition, we obtain \(\operatorname{delete\,Edges}(G,\allowbreak{} \{\operatorname{Sym2\mathord{:}mk}((v,\allowbreak{} w))\}) = G \setminus \operatorname{from\,Edge\,Set}(\{\operatorname{Sym2\mathord{:}mk}((v,\allowbreak{} w))\})\).
Use the previously proved fact \(\operatorname{delete\,Edges}(G,\allowbreak{} \{\operatorname{Sym2\mathord{:}mk}((v,\allowbreak{} w))\}) = G \setminus \operatorname{from\,Edge\,Set}(\{\operatorname{Sym2\mathord{:}mk}((v,\allowbreak{} w))\})\).
\end{proofstep}
Use the hypothesis \(v \sim_{G} w\).
By \paperlink{premise:53696d706c6547726170682e57616c6b2e4973506174682e6f665f61646a}{Prerequisite~P\ref{premise:53696d706c6547726170682e57616c6b2e4973506174682e6f665f61646a}}, we obtain \(\operatorname{to\,Walk}\text{ has no repeated vertices}\).
By \paperlink{theorem:38}{Lemma~\ref{theorem:38}} with \(p = \operatorname{to\,Walk}\), we obtain \(w\text{ can reach }v\text{ in }G \setminus \operatorname{spanning\,graph}(\operatorname{Simple\,Graph\mathord{:}Walk\mathord{:}to\,Subgraph}(\operatorname{to\,Walk}))\).
By \paperlink{premise:53696d706c6547726170682e526561636861626c652e73796d6d}{Prerequisite~P\ref{premise:53696d706c6547726170682e526561636861626c652e73796d6d}}, we obtain \(v\text{ can reach }w\text{ in }G \setminus \operatorname{spanning\,graph}(\operatorname{Simple\,Graph\mathord{:}Walk\mathord{:}to\,Subgraph}(\operatorname{to\,Walk}))\).
Use the previously proved fact \(\operatorname{from\,Edge\,Set}(\{\operatorname{Sym2\mathord{:}mk}((v,\allowbreak{} w))\}) = \operatorname{spanning\,graph}(\operatorname{Simple\,Graph\mathord{:}Walk\mathord{:}to\,Subgraph}(\operatorname{to\,Walk}))\).
Reversing the equality, we obtain \(\operatorname{spanning\,graph}(\operatorname{Simple\,Graph\mathord{:}Walk\mathord{:}to\,Subgraph}(\operatorname{to\,Walk})) = \operatorname{from\,Edge\,Set}(\{\operatorname{Sym2\mathord{:}mk}((v,\allowbreak{} w))\})\).
Substituting equals, we obtain \(v\text{ can reach }w\text{ in }G \setminus \operatorname{from\,Edge\,Set}(\{\operatorname{Sym2\mathord{:}mk}((v,\allowbreak{} w))\})\).
Rewriting with the identity established above in the preceding result gives \(v\text{ can reach }w\text{ in }\operatorname{delete\,Edges}(G,\allowbreak{} \{\operatorname{Sym2\mathord{:}mk}((v,\allowbreak{} w))\})\).
\end{proof}

\begin{theorem}[{Exists cycle to subgraph verts equals connected component supp}]\label{theorem:40}
Let \(V\) be a type. Let \(G\) be a simple graph on \(V\). Let \(v\) be an element of \(V\). Assume \(V\) is finite. Let \(c\) be an element of \(\operatorname{Connected\,Component}\left(G\right)\). Assume \(\text{every nonisolated vertex of }G\text{ has degree two}\). Assume \(v \in \operatorname{supp}\left(c\right)\). Assume \(N_{G}(v)\ne\varnothing\). Then
\begin{center}
\(\displaystyle \exists p\, :\, \operatorname{Walk}_{G}(v,\allowbreak{} v),\; p\text{ is a cycle} \land V(\operatorname{Simple\,Graph\mathord{:}Walk\mathord{:}to\,Subgraph}(p)) = \operatorname{supp}(c).\)
\end{center}
\end{theorem}

\begin{proof}
Use the hypothesis \(N_{G}(v)\ne\varnothing\).
Consider the following case.
\begin{proofstep}
Fix \(w\) in \(V\).
Assume \(w \in N_{G}(v)\).
By \paperlink{premise:53696d706c6547726170682e61646a5f616e645f726561636861626c655f64656c6574655f65646765735f6966665f6578697374735f6379636c65}{Prerequisite~P\ref{premise:53696d706c6547726170682e61646a5f616e645f726561636861626c655f64656c6574655f65646765735f6966665f6578697374735f6379636c65}}, we obtain \(v \sim_{G} w \land v\text{ can reach }w\text{ in }G \setminus \operatorname{from\,Edge\,Set}(\{\operatorname{Sym2\mathord{:}mk}((v,\allowbreak{} w))\}) \iff (\exists u\, :\, V,\; \exists p\, :\, \operatorname{Walk}_{G}(u,\allowbreak{} u),\; p\text{ is a cycle} \land \operatorname{Sym2\mathord{:}mk}((v,\allowbreak{} w)) \in \operatorname{edges}(p))\).
Use the hypothesis \(w \in N_{G}(v)\).
By \paperlink{theorem:39}{Lemma~\ref{theorem:39}}, we obtain \(v\text{ can reach }w\text{ in }\operatorname{delete\,Edges}(G,\allowbreak{} \{\operatorname{Sym2\mathord{:}mk}((v,\allowbreak{} w))\})\).
Combining the two conclusions, we obtain \(v \sim_{G} w \land v\text{ can reach }w\text{ in }G \setminus \operatorname{from\,Edge\,Set}(\{\operatorname{Sym2\mathord{:}mk}((v,\allowbreak{} w))\})\).
By the forward implication, we obtain \(\exists u\, :\, V,\; \exists p\, :\, \operatorname{Walk}_{G}(u,\allowbreak{} u),\; p\text{ is a cycle} \land \operatorname{Sym2\mathord{:}mk}((v,\allowbreak{} w)) \in \operatorname{edges}(p)\).
Consider the following case.
\begin{proofstep}
Fix \(u\) in \(V\).
Assume \(\exists p\, :\, \operatorname{Walk}_{G}(u,\allowbreak{} u),\; p\text{ is a cycle} \land \operatorname{Sym2\mathord{:}mk}((v,\allowbreak{} w)) \in \operatorname{edges}(p)\).
Use the hypothesis \(\exists p\, :\, \operatorname{Walk}_{G}(u,\allowbreak{} u),\; p\text{ is a cycle} \land \operatorname{Sym2\mathord{:}mk}((v,\allowbreak{} w)) \in \operatorname{edges}(p)\).
Consider the following case.
\begin{proofstep}
Fix \(p\) in \(\operatorname{Walk}_{G}(u, u)\).
Assume \(p\text{ is a cycle} \land \operatorname{Sym2\mathord{:}mk}((v,\allowbreak{} w)) \in \operatorname{edges}(p)\).
Use the hypothesis \(p\text{ is a cycle} \land \operatorname{Sym2\mathord{:}mk}((v,\allowbreak{} w)) \in \operatorname{edges}(p)\).
Taking the second conjunct, we obtain \(\operatorname{Sym2\mathord{:}mk}\left(\left(v, w\right)\right) \in \operatorname{edges}\left(p\right)\).
By \paperlink{premise:53696d706c6547726170682e57616c6b2e6673745f6d656d5f737570706f72745f6f665f6d656d5f6564676573}{Prerequisite~P\ref{premise:53696d706c6547726170682e57616c6b2e6673745f6d656d5f737570706f72745f6f665f6d656d5f6564676573}}, we obtain \(p\text{ visits }v\).
By \paperlink{premise:53696d706c6547726170682e57616c6b2e746f53756267726170685f636f6e6e6563746564}{Prerequisite~P\ref{premise:53696d706c6547726170682e57616c6b2e746f53756267726170685f636f6e6e6563746564}}, we obtain \(\operatorname{Connected}(\operatorname{Simple\,Graph\mathord{:}Walk\mathord{:}to\,Subgraph}(p))\).
For the required implication, argue as follows.
\begin{proofstep}
Fix \(v_{1}\) in \(V\).
Assume \(v_{1} \in V(\operatorname{Simple\,Graph\mathord{:}Walk\mathord{:}to\,Subgraph}\left(p\right))\).
Fix \(w_{1}\) in \(V\).
Assume \(v_{1} \sim_{G} w_{1}\).
Use the hypothesis \(\operatorname{Finite}\left(V\right)\).
By \paperlink{premise:537562747970652e66696e697465}{Prerequisite~P\ref{premise:537562747970652e66696e697465}}, we obtain \(\operatorname{Finite}\left(N_{G}(v_{1})\right)\).
By \paperlink{premise:5365742e746f46696e697465}{Prerequisite~P\ref{premise:5365742e746f46696e697465}} with \(s = N_{G}(v_{1})\), we obtain \(\operatorname{Set\mathord{:}Finite}\left(N_{G}(v_{1})\right)\).
By \paperlink{premise:53696d706c6547726170682e53756267726170682e61646a5f6966665f6f665f6e65696768626f725365745f6571756976}{Prerequisite~P\ref{premise:53696d706c6547726170682e53756267726170682e61646a5f6966665f6f665f6e65696768626f725365745f6571756976}} with \(\operatorname{let\,Fun}\left(\mathit{this} \mapsto \operatorname{of\,Nonempty}\right)\) in place of its parameter \(h\), we obtain \(v_{1} \sim_{\operatorname{Simple\,Graph\mathord{:}Walk\mathord{:}to\,Subgraph}(p)} w_{1} \iff v_{1} \sim_{G} w_{1}\).
Use the hypothesis \(v_{1} \sim_{G} w_{1}\).
By the reverse implication, we obtain \(v_{1} \sim_{\operatorname{Simple\,Graph\mathord{:}Walk\mathord{:}to\,Subgraph}\left(p\right)} w_{1}\).
\end{proofstep}
By \paperlink{premise:53696d706c6547726170682e53756267726170682e436f6e6e65637465642e6578697374735f76657274735f65715f636f6e6e6563746564436f6d706f6e656e7453757070}{Prerequisite~P\ref{premise:53696d706c6547726170682e53756267726170682e436f6e6e65637465642e6578697374735f76657274735f65715f636f6e6e6563746564436f6d706f6e656e7453757070}}, we obtain \(\exists c_{1}\, :\, \operatorname{Connected\,Component}(G),\; V(\operatorname{Simple\,Graph\mathord{:}Walk\mathord{:}to\,Subgraph}(p)) = \operatorname{supp}(c_{1})\).
Consider the following case.
\begin{proofstep}
Fix \(c^{'}\) in \(\operatorname{Connected\,Component}\left(G\right)\).
Assume \(V(\operatorname{Simple\,Graph\mathord{:}Walk\mathord{:}to\,Subgraph}(p)) = \operatorname{supp}(c^{'})\).
Use the previously proved fact \(p\text{ visits }v\).
Rewriting with \paperlink{premise:53696d706c6547726170682e57616c6b2e6d656d5f76657274735f746f5375626772617068}{Prerequisite~P\ref{premise:53696d706c6547726170682e57616c6b2e6d656d5f76657274735f746f5375626772617068}} in the preceding result gives \(v \in V(\operatorname{Simple\,Graph\mathord{:}Walk\mathord{:}to\,Subgraph}\left(p\right))\).
Rewriting with \(V(\operatorname{Simple\,Graph\mathord{:}Walk\mathord{:}to\,Subgraph}(p)) = \operatorname{supp}(c^{'})\) in the preceding result gives \(v \in \operatorname{supp}\left(c^{'}\right)\).
Rewriting with \paperlink{premise:53696d706c6547726170682e436f6e6e6563746564436f6d706f6e656e742e6d656d5f737570705f696666}{Prerequisite~P\ref{premise:53696d706c6547726170682e436f6e6e6563746564436f6d706f6e656e742e6d656d5f737570705f696666}}, \(v \in \operatorname{supp}\left(c\right)\) and \(v \in \operatorname{supp}\left(c^{'}\right)\) gives \(\operatorname{supp}\left(c^{'}\right) = \operatorname{supp}\left(c\right)\).
Rewriting with \(V(\operatorname{Simple\,Graph\mathord{:}Walk\mathord{:}to\,Subgraph}(p)) = \operatorname{supp}(c^{'})\) in the preceding result gives \(V(\operatorname{Simple\,Graph\mathord{:}Walk\mathord{:}to\,Subgraph}(p)) = \operatorname{supp}(c)\).
\end{proofstep}
Using the existential hypothesis, we obtain \(V(\operatorname{Simple\,Graph\mathord{:}Walk\mathord{:}to\,Subgraph}(p)) = \operatorname{supp}(c)\).
Choose the witness \(\operatorname{Simple\,Graph\mathord{:}Walk\mathord{:}rotate}\left(p\right)\).
Use the previously proved fact \(p\text{ visits }v\).
Use the hypothesis \(p\text{ is a cycle} \land \operatorname{Sym2\mathord{:}mk}((v,\allowbreak{} w)) \in \operatorname{edges}(p)\).
Taking the first conjunct, we obtain \(p\text{ is a cycle}\).
Use the previously proved fact \(p\text{ visits }v\).
By \paperlink{premise:53696d706c6547726170682e57616c6b2e49734379636c652e726f74617465}{Prerequisite~P\ref{premise:53696d706c6547726170682e57616c6b2e49734379636c652e726f74617465}}, we obtain \(\operatorname{Simple\,Graph\mathord{:}Walk\mathord{:}rotate}\left(p\right)\text{ is a cycle}\).
Rewriting with \paperlink{premise:53696d706c6547726170682e57616c6b2e746f53756267726170685f726f74617465}{Prerequisite~P\ref{premise:53696d706c6547726170682e57616c6b2e746f53756267726170685f726f74617465}} (using \(p\text{ visits }v\)), \paperlink{premise:53696d706c6547726170682e57616c6b2e76657274735f746f5375626772617068}{Prerequisite~P\ref{premise:53696d706c6547726170682e57616c6b2e76657274735f746f5375626772617068}}, \(V(\operatorname{Simple\,Graph\mathord{:}Walk\mathord{:}to\,Subgraph}(p)) = \operatorname{supp}(c)\) and substitution of equals gives \(V(\operatorname{Simple\,Graph\mathord{:}Walk\mathord{:}to\,Subgraph}(\operatorname{Simple\,Graph\mathord{:}Walk\mathord{:}rotate}(p))) = V(\operatorname{Simple\,Graph\mathord{:}Walk\mathord{:}to\,Subgraph}(p))\).
Combining the two conclusions, we obtain \(\operatorname{Simple\,Graph\mathord{:}Walk\mathord{:}rotate}(p)\text{ is a cycle} \land V(\operatorname{Simple\,Graph\mathord{:}Walk\mathord{:}to\,Subgraph}(\operatorname{Simple\,Graph\mathord{:}Walk\mathord{:}rotate}(p))) = V(\operatorname{Simple\,Graph\mathord{:}Walk\mathord{:}to\,Subgraph}(p))\).
Rewriting with \(V(\operatorname{Simple\,Graph\mathord{:}Walk\mathord{:}to\,Subgraph}(p)) = \operatorname{supp}(c)\) in the preceding result gives \(\operatorname{Simple\,Graph\mathord{:}Walk\mathord{:}rotate}(p)\text{ is a cycle} \land V(\operatorname{Simple\,Graph\mathord{:}Walk\mathord{:}to\,Subgraph}(\operatorname{Simple\,Graph\mathord{:}Walk\mathord{:}rotate}(p))) = \operatorname{supp}(c)\).
This witness proves existence. Hence \(\exists p_{1}\, :\, \operatorname{Walk}_{G}(v,\allowbreak{} v),\; p_{1}\text{ is a cycle} \land V(\operatorname{Simple\,Graph\mathord{:}Walk\mathord{:}to\,Subgraph}(p_{1})) = \operatorname{supp}(c)\).
\end{proofstep}
Using the existential hypothesis, we obtain \(\exists p\, :\, \operatorname{Walk}_{G}(v,\allowbreak{} v),\; p\text{ is a cycle} \land V(\operatorname{Simple\,Graph\mathord{:}Walk\mathord{:}to\,Subgraph}(p)) = \operatorname{supp}(c)\).
\end{proofstep}
Using the existential hypothesis, we obtain \(\exists p\, :\, \operatorname{Walk}_{G}(v,\allowbreak{} v),\; p\text{ is a cycle} \land V(\operatorname{Simple\,Graph\mathord{:}Walk\mathord{:}to\,Subgraph}(p)) = \operatorname{supp}(c)\).
\end{proofstep}
Using the existential hypothesis, we obtain \(\exists p\, :\, \operatorname{Walk}_{G}(v,\allowbreak{} v),\; p\text{ is a cycle} \land V(\operatorname{Simple\,Graph\mathord{:}Walk\mathord{:}to\,Subgraph}(p)) = \operatorname{supp}(c)\).
\end{proof}

\begin{theorem}[{Mono}]\label{theorem:41}
Let \(V\) be a type. Let \(G\), \(G^{'}\), and \(G^{''}\) be simple graphs on \(V\). Assume \(G\text{ alternates with respect to }G^{'}\). Assume \(G^{''} \subseteq G\). Then \(G^{''}\text{ alternates with respect to }G^{'}\).
\end{theorem}

\begin{proof}
Fix \(x\) in \(V\).
Fix \(x_{1}\) in \(V\).
Fix \(x_{2}\) in \(V\).
Assume \(x_{1} \ne x_{2}\).
Assume \(x \sim_{G^{''}} x_{1}\).
Assume \(x \sim_{G^{''}} x_{2}\).
Use the hypothesis \(x_{1} \ne x_{2}\).
Use the hypothesis \(x \sim_{G^{''}} x_{1}\).
Applying the hypothesis \(G^{''} \subseteq G\), we obtain \(x \sim_{G} x_{1}\).
Use the hypothesis \(x \sim_{G^{''}} x_{2}\).
Applying the hypothesis \(G^{''} \subseteq G\), we obtain \(x \sim_{G} x_{2}\).
Applying the hypothesis \(G\text{ alternates with respect to }G^{'}\), we obtain \(x \sim_{G^{'}} x_{1} \iff \neg \left(x \sim_{G^{'}} x_{2}\right)\).
\end{proof}

\begin{theorem}[{Spanning coercion}]\label{theorem:42}
Let \(V\) be a type. Let \(G\) and \(G^{'}\) be simple graphs on \(V\). Assume \(G\text{ alternates with respect to }G^{'}\). Let \(H\) be a subgraph of \(G\). Then \(\operatorname{spanning\,graph}(H)\text{ alternates with respect to }G^{'}\).
\end{theorem}

\begin{proof}
Fix \(v\) in \(V\).
Fix \(w\) in \(V\).
Fix \(w^{'}\) in \(V\).
Assume \(w \ne w^{'}\).
Assume \(v \sim_{\operatorname{spanning\,graph}(H)} w\).
Assume \(v \sim_{\operatorname{spanning\,graph}(H)} w^{'}\).
Use the hypothesis \(w \ne w^{'}\).
Use the hypothesis \(v \sim_{\operatorname{spanning\,graph}(H)} w\).
By \paperlink{premise:53696d706c6547726170682e53756267726170682e41646a2e61646a5f737562}{Prerequisite~P\ref{premise:53696d706c6547726170682e53756267726170682e41646a2e61646a5f737562}}, we obtain \(v \sim_{G} w\).
Use the hypothesis \(v \sim_{\operatorname{spanning\,graph}(H)} w^{'}\).
By \paperlink{premise:53696d706c6547726170682e53756267726170682e41646a2e61646a5f737562}{Prerequisite~P\ref{premise:53696d706c6547726170682e53756267726170682e41646a2e61646a5f737562}}, we obtain \(v \sim_{G} w^{'}\).
Applying the hypothesis \(G\text{ alternates with respect to }G^{'}\), we obtain \(v \sim_{G^{'}} w \iff \neg \left(v \sim_{G^{'}} w^{'}\right)\).
\end{proof}

\begin{theorem}[{Sup edge}]\label{theorem:43}
Let \(V\) be a type. Let \(G\) and \(G^{'}\) be simple graphs on \(V\). Let \(u\) and \(x\) be elements of \(V\). Assume \(G\text{ alternates with respect to }G^{'}\). Assume \(\neg \left(u \sim_{G^{'}} x\right)\). Assume \(\forall u^{'}\, :\, V,\; u^{'} \ne u \implies x \sim_{G} u^{'} \implies x \sim_{G^{'}} u^{'}\). Assume \(\forall x^{'}\, :\, V,\; x^{'} \ne x \implies x^{'} \sim_{G} u \implies x^{'} \sim_{G^{'}} u\). Then \(G \cup \operatorname{edge}\left(u, x\right)\text{ alternates with respect to }G^{'}\).
\end{theorem}

\begin{proof}
Distinguish whether \(u \sim_{G} x\) holds.
Consider the following case.
\begin{proofstep}
Assume \(u \sim_{G} x\).
Use the hypothesis \(G\text{ alternates with respect to }G^{'}\).
Rewriting with \paperlink{premise:53696d706c6547726170682e7375705f656467655f6f665f61646a}{Prerequisite~P\ref{premise:53696d706c6547726170682e7375705f656467655f6f665f61646a}} (using \(u \sim_{G} x\)) in the preceding result gives \(G \cup \operatorname{edge}\left(u, x\right)\text{ alternates with respect to }G^{'}\).
\end{proofstep}
Consider the following case.
\begin{proofstep}
Assume \(\neg \left(u \sim_{G} x\right)\).
Fix \(v\) in \(V\).
Fix \(w\) in \(V\).
Fix \(w^{'}\) in \(V\).
Assume \(w \ne w^{'}\).
Assume \(v \sim_{G \cup \operatorname{edge}\left(u, x\right)} w\).
Assume \(v \sim_{G \cup \operatorname{edge}\left(u, x\right)} w^{'}\).
Use the hypothesis \(v \sim_{G \cup \operatorname{edge}\left(u, x\right)} w\).
Rewriting with \paperlink{premise:53696d706c6547726170682e7375705f61646a}{Prerequisite~P\ref{premise:53696d706c6547726170682e7375705f61646a}} and \paperlink{premise:53696d706c6547726170682e656467655f61646a}{Prerequisite~P\ref{premise:53696d706c6547726170682e656467655f61646a}} in the preceding result gives \(v \sim_{G} w \lor \left(v = u \land w = x \lor v = x \land w = u\right) \land v \ne w\).
Consider the following case.
\begin{proofstep}
Assume \(v \sim_{G} w\).
Use the hypothesis \(v \sim_{G \cup \operatorname{edge}\left(u, x\right)} w^{'}\).
Rewriting with \paperlink{premise:53696d706c6547726170682e7375705f61646a}{Prerequisite~P\ref{premise:53696d706c6547726170682e7375705f61646a}} and \paperlink{premise:53696d706c6547726170682e656467655f61646a}{Prerequisite~P\ref{premise:53696d706c6547726170682e656467655f61646a}} in the preceding result gives \(v \sim_{G} w^{'} \lor \left(v = u \land w^{'} = x \lor v = x \land w^{'} = u\right) \land v \ne w^{'}\).
Consider the following case.
\begin{proofstep}
Assume \(v \sim_{G} w^{'}\).
Use the hypothesis \(w \ne w^{'}\).
Use the hypothesis \(v \sim_{G} w\).
Use the hypothesis \(v \sim_{G} w^{'}\).
Applying the hypothesis \(G\text{ alternates with respect to }G^{'}\), we obtain \(v \sim_{G^{'}} w \iff \neg \left(v \sim_{G^{'}} w^{'}\right)\).
\end{proofstep}
Consider the following case.
\begin{proofstep}
Assume \(\left(v = u \land w^{'} = x \lor v = x \land w^{'} = u\right) \land v \ne w^{'}\).
Use the hypothesis \(\left(v = u \land w^{'} = x \lor v = x \land w^{'} = u\right) \land v \ne w^{'}\).
By the preceding argument, we obtain \(\left(v = u \land w^{'} = x \lor v = x \land w^{'} = u\right) \land \neg \left(v = w^{'}\right)\).
By the preceding argument, we obtain \(\left(v = u \land w^{'} = x \lor v = x \land w^{'} = u\right) \land \neg \left(v = w^{'}\right)\).
Taking the first conjunct gives \(v = u \land w^{'} = x \lor v = x \land w^{'} = u\).
Use the hypothesis \(\left(v = u \land w^{'} = x \lor v = x \land w^{'} = u\right) \land v \ne w^{'}\).
Taking the first conjunct, we obtain \(v = u \land w^{'} = x \lor v = x \land w^{'} = u\).
Consider the following case.
\begin{proofstep}
Assume \(v = u \land w^{'} = x\).
Use the hypothesis \(v = u \land w^{'} = x\).
Consider the following case.
\begin{proofstep}
Assume \(v = u\).
Assume \(w^{'} = x\).
For the required implication, argue as follows.
\begin{proofstep}
Assume \(\neg \left(v \sim_{G^{'}} x\right)\).
Assume \(\forall u^{'}\, :\, V,\; u^{'} \ne v \implies x \sim_{G} u^{'} \implies x \sim_{G^{'}} u^{'}\).
Assume \(\forall x^{'}\, :\, V,\; x^{'} \ne x \implies x^{'} \sim_{G} v \implies x^{'} \sim_{G^{'}} v\).
Assume \(\neg \left(v \sim_{G} x\right)\).
Assume \(\left(v = v \land w^{'} = x \lor v = x \land w^{'} = v\right) \land v \ne w^{'}\).
For the required implication, argue as follows.
\begin{proofstep}
Assume \(\neg \left(v \sim_{G^{'}} w^{'}\right)\).
Assume \(\forall u^{'}\, :\, V,\; u^{'} \ne v \implies w^{'} \sim_{G} u^{'} \implies w^{'} \sim_{G^{'}} u^{'}\).
Assume \(\forall x^{'}\, :\, V,\; x^{'} \ne w^{'} \implies x^{'} \sim_{G} v \implies x^{'} \sim_{G^{'}} v\).
Assume \(\neg \left(v \sim_{G} w^{'}\right)\).
Assume \(\left(v = v \land w^{'} = w^{'} \lor v = w^{'} \land w^{'} = v\right) \land v \ne w^{'}\).
Use the hypothesis \(w \ne w^{'}\).
Use the hypothesis \(v \sim_{G} w\).
By \paperlink{premise:53696d706c6547726170682e41646a2e73796d6d}{Prerequisite~P\ref{premise:53696d706c6547726170682e41646a2e73796d6d}}, we obtain \(w \sim_{G} v\).
Applying the hypothesis \(\forall x^{'}\, :\, V,\; x^{'} \ne w^{'} \implies x^{'} \sim_{G} v \implies x^{'} \sim_{G^{'}} v\), we obtain \(w \sim_{G^{'}} v\).
By \paperlink{premise:53696d706c6547726170682e41646a2e73796d6d}{Prerequisite~P\ref{premise:53696d706c6547726170682e41646a2e73796d6d}}, we obtain \(v \sim_{G^{'}} w\).
\end{proofstep}
Use the hypothesis \(w^{'} = x\).
Use the hypothesis \(\neg \left(v \sim_{G^{'}} x\right)\).
Use the hypothesis \(\forall u^{'}\, :\, V,\; u^{'} \ne v \implies x \sim_{G} u^{'} \implies x \sim_{G^{'}} u^{'}\).
Use the hypothesis \(\forall x^{'}\, :\, V,\; x^{'} \ne x \implies x^{'} \sim_{G} v \implies x^{'} \sim_{G^{'}} v\).
Use the hypothesis \(\neg \left(v \sim_{G} x\right)\).
Use the hypothesis \(\left(v = v \land w^{'} = x \lor v = x \land w^{'} = v\right) \land v \ne w^{'}\).
Substituting equals, we obtain \(v \sim_{G^{'}} w\).
\end{proofstep}
Use the hypothesis \(v = u\).
Use the hypothesis \(\neg \left(u \sim_{G^{'}} x\right)\).
Use the hypothesis \(\forall u^{'}\, :\, V,\; u^{'} \ne u \implies x \sim_{G} u^{'} \implies x \sim_{G^{'}} u^{'}\).
Use the hypothesis \(\forall x^{'}\, :\, V,\; x^{'} \ne x \implies x^{'} \sim_{G} u \implies x^{'} \sim_{G^{'}} u\).
Use the hypothesis \(\neg \left(u \sim_{G} x\right)\).
Use the hypothesis \(\left(v = u \land w^{'} = x \lor v = x \land w^{'} = u\right) \land v \ne w^{'}\).
Substituting equals, we obtain \(v \sim_{G^{'}} w\).
\end{proofstep}
Using the two conjuncts, we obtain \(v \sim_{G^{'}} w\).
\end{proofstep}
Consider the following case.
\begin{proofstep}
Assume \(v = x \land w^{'} = u\).
Use the hypothesis \(v = x \land w^{'} = u\).
Consider the following case.
\begin{proofstep}
Assume \(v = x\).
Assume \(w^{'} = u\).
For the required implication, argue as follows.
\begin{proofstep}
Assume \(\neg \left(w^{'} \sim_{G^{'}} x\right)\).
Assume \(\forall u^{'}\, :\, V,\; u^{'} \ne w^{'} \implies x \sim_{G} u^{'} \implies x \sim_{G^{'}} u^{'}\).
Assume \(\forall x^{'}\, :\, V,\; x^{'} \ne x \implies x^{'} \sim_{G} w^{'} \implies x^{'} \sim_{G^{'}} w^{'}\).
Assume \(\neg \left(w^{'} \sim_{G} x\right)\).
Assume \(\left(v = w^{'} \land w^{'} = x \lor v = x \land w^{'} = w^{'}\right) \land v \ne w^{'}\).
For the required implication, argue as follows.
\begin{proofstep}
Assume \(\neg \left(w^{'} \sim_{G^{'}} v\right)\).
Assume \(\forall u^{'}\, :\, V,\; u^{'} \ne w^{'} \implies v \sim_{G} u^{'} \implies v \sim_{G^{'}} u^{'}\).
Assume \(\forall x^{'}\, :\, V,\; x^{'} \ne v \implies x^{'} \sim_{G} w^{'} \implies x^{'} \sim_{G^{'}} w^{'}\).
Assume \(\neg \left(w^{'} \sim_{G} v\right)\).
Assume \(\left(v = w^{'} \land w^{'} = v \lor v = v \land w^{'} = w^{'}\right) \land v \ne w^{'}\).
Rewriting with \(\forall u^{'}\, :\, V,\; u^{'} \ne w^{'} \implies v \sim_{G} u^{'} \implies v \sim_{G^{'}} u^{'}\), \paperlink{premise:65715f66616c7365}{Prerequisite~P\ref{premise:65715f66616c7365}}, \(w \ne w^{'}\), \paperlink{premise:6e6f745f66616c73655f65715f74727565}{Prerequisite~P\ref{premise:6e6f745f66616c73655f65715f74727565}} and \(v \sim_{G} w\) gives \(v \sim_{G^{'}} w\).
\end{proofstep}
Use the hypothesis \(v = x\).
Use the hypothesis \(\neg \left(w^{'} \sim_{G^{'}} x\right)\).
Use the hypothesis \(\forall u^{'}\, :\, V,\; u^{'} \ne w^{'} \implies x \sim_{G} u^{'} \implies x \sim_{G^{'}} u^{'}\).
Use the hypothesis \(\forall x^{'}\, :\, V,\; x^{'} \ne x \implies x^{'} \sim_{G} w^{'} \implies x^{'} \sim_{G^{'}} w^{'}\).
Use the hypothesis \(\neg \left(w^{'} \sim_{G} x\right)\).
Use the hypothesis \(\left(v = w^{'} \land w^{'} = x \lor v = x \land w^{'} = w^{'}\right) \land v \ne w^{'}\).
Substituting equals, we obtain \(v \sim_{G^{'}} w\).
\end{proofstep}
Use the hypothesis \(w^{'} = u\).
Use the hypothesis \(\neg \left(u \sim_{G^{'}} x\right)\).
Use the hypothesis \(\forall u^{'}\, :\, V,\; u^{'} \ne u \implies x \sim_{G} u^{'} \implies x \sim_{G^{'}} u^{'}\).
Use the hypothesis \(\forall x^{'}\, :\, V,\; x^{'} \ne x \implies x^{'} \sim_{G} u \implies x^{'} \sim_{G^{'}} u\).
Use the hypothesis \(\neg \left(u \sim_{G} x\right)\).
Use the hypothesis \(\left(v = u \land w^{'} = x \lor v = x \land w^{'} = u\right) \land v \ne w^{'}\).
Substituting equals, we obtain \(v \sim_{G^{'}} w\).
\end{proofstep}
Using the two conjuncts, we obtain \(v \sim_{G^{'}} w\).
\end{proofstep}
The two alternatives exhaust the cases. Hence \(v \sim_{G^{'}} w\).
Rewriting with \paperlink{premise:65715f66616c7365}{Prerequisite~P\ref{premise:65715f66616c7365}} (using \(\neg \left(u \sim_{G^{'}} x\right)\)), \paperlink{premise:6e6f745f66616c73655f65715f74727565}{Prerequisite~P\ref{premise:6e6f745f66616c73655f65715f74727565}} and \paperlink{premise:6966665f74727565}{Prerequisite~P\ref{premise:6966665f74727565}} in the preceding result gives \(v \sim_{G^{'}} w \iff \neg \left(u \sim_{G^{'}} x\right)\).
Rewriting with \paperlink{premise:53696d706c6547726170682e61646a5f636f6e67725f6f665f73796d32}{Prerequisite~P\ref{premise:53696d706c6547726170682e61646a5f636f6e67725f6f665f73796d32}}, \paperlink{premise:53796d322e6571}{Prerequisite~P\ref{premise:53796d322e6571}}, \paperlink{premise:53796d322e72656c5f69666627}{Prerequisite~P\ref{premise:53796d322e72656c5f69666627}}, \paperlink{premise:50726f642e6d6b2e696e6a4571}{Prerequisite~P\ref{premise:50726f642e6d6b2e696e6a4571}} and \(v = u \land w^{'} = x \lor v = x \land w^{'} = u\) in the preceding result gives \(v \sim_{G^{'}} w \iff \neg \left(v \sim_{G^{'}} w^{'}\right)\).
\end{proofstep}
The two alternatives exhaust the cases. Hence \(v \sim_{G^{'}} w \iff \neg \left(v \sim_{G^{'}} w^{'}\right)\).
\end{proofstep}
Consider the following case.
\begin{proofstep}
Assume \(\left(v = u \land w = x \lor v = x \land w = u\right) \land v \ne w\).
Use the hypothesis \(v \sim_{G \cup \operatorname{edge}\left(u, x\right)} w^{'}\).
Rewriting with \paperlink{premise:53696d706c6547726170682e7375705f61646a}{Prerequisite~P\ref{premise:53696d706c6547726170682e7375705f61646a}} and \paperlink{premise:53696d706c6547726170682e656467655f61646a}{Prerequisite~P\ref{premise:53696d706c6547726170682e656467655f61646a}} in the preceding result gives \(v \sim_{G} w^{'} \lor \left(v = u \land w^{'} = x \lor v = x \land w^{'} = u\right) \land v \ne w^{'}\).
Consider the following case.
\begin{proofstep}
Assume \(v \sim_{G} w^{'}\).
Use the hypothesis \(\left(v = u \land w = x \lor v = x \land w = u\right) \land v \ne w\).
By the preceding argument, we obtain \(\left(v = u \land w = x \lor v = x \land w = u\right) \land \neg \left(v = w\right)\).
By the preceding argument, we obtain \(\left(v = u \land w = x \lor v = x \land w = u\right) \land \neg \left(v = w\right)\).
Taking the first conjunct gives \(v = u \land w = x \lor v = x \land w = u\).
Use the hypothesis \(\left(v = u \land w = x \lor v = x \land w = u\right) \land v \ne w\).
Taking the first conjunct, we obtain \(v = u \land w = x \lor v = x \land w = u\).
Consider the following case.
\begin{proofstep}
Assume \(v = u \land w = x\).
Use the hypothesis \(v = u \land w = x\).
Consider the following case.
\begin{proofstep}
Assume \(v = u\).
Assume \(w = x\).
For the required implication, argue as follows.
\begin{proofstep}
Assume \(\neg \left(v \sim_{G^{'}} x\right)\).
Assume \(\forall u^{'}\, :\, V,\; u^{'} \ne v \implies x \sim_{G} u^{'} \implies x \sim_{G^{'}} u^{'}\).
Assume \(\forall x^{'}\, :\, V,\; x^{'} \ne x \implies x^{'} \sim_{G} v \implies x^{'} \sim_{G^{'}} v\).
Assume \(\neg \left(v \sim_{G} x\right)\).
Assume \(\left(v = v \land w = x \lor v = x \land w = v\right) \land v \ne w\).
For the required implication, argue as follows.
\begin{proofstep}
Assume \(\neg \left(v \sim_{G^{'}} w\right)\).
Assume \(\forall u^{'}\, :\, V,\; u^{'} \ne v \implies w \sim_{G} u^{'} \implies w \sim_{G^{'}} u^{'}\).
Assume \(\forall x^{'}\, :\, V,\; x^{'} \ne w \implies x^{'} \sim_{G} v \implies x^{'} \sim_{G^{'}} v\).
Assume \(\neg \left(v \sim_{G} w\right)\).
Assume \(\left(v = v \land w = w \lor v = w \land w = v\right) \land v \ne w\).
Use the hypothesis \(w \ne w^{'}\).
By \paperlink{premise:4e652e73796d6d}{Prerequisite~P\ref{premise:4e652e73796d6d}}, we obtain \(w^{'} \ne w\).
Use the hypothesis \(v \sim_{G} w^{'}\).
By \paperlink{premise:53696d706c6547726170682e41646a2e73796d6d}{Prerequisite~P\ref{premise:53696d706c6547726170682e41646a2e73796d6d}}, we obtain \(w^{'} \sim_{G} v\).
Applying the hypothesis \(\forall x^{'}\, :\, V,\; x^{'} \ne w \implies x^{'} \sim_{G} v \implies x^{'} \sim_{G^{'}} v\), we obtain \(w^{'} \sim_{G^{'}} v\).
By \paperlink{premise:53696d706c6547726170682e41646a2e73796d6d}{Prerequisite~P\ref{premise:53696d706c6547726170682e41646a2e73796d6d}}, we obtain \(v \sim_{G^{'}} w^{'}\).
\end{proofstep}
Use the hypothesis \(w = x\).
Use the hypothesis \(\neg \left(v \sim_{G^{'}} x\right)\).
Use the hypothesis \(\forall u^{'}\, :\, V,\; u^{'} \ne v \implies x \sim_{G} u^{'} \implies x \sim_{G^{'}} u^{'}\).
Use the hypothesis \(\forall x^{'}\, :\, V,\; x^{'} \ne x \implies x^{'} \sim_{G} v \implies x^{'} \sim_{G^{'}} v\).
Use the hypothesis \(\neg \left(v \sim_{G} x\right)\).
Use the hypothesis \(\left(v = v \land w = x \lor v = x \land w = v\right) \land v \ne w\).
Substituting equals, we obtain \(v \sim_{G^{'}} w^{'}\).
\end{proofstep}
Use the hypothesis \(v = u\).
Use the hypothesis \(\neg \left(u \sim_{G^{'}} x\right)\).
Use the hypothesis \(\forall u^{'}\, :\, V,\; u^{'} \ne u \implies x \sim_{G} u^{'} \implies x \sim_{G^{'}} u^{'}\).
Use the hypothesis \(\forall x^{'}\, :\, V,\; x^{'} \ne x \implies x^{'} \sim_{G} u \implies x^{'} \sim_{G^{'}} u\).
Use the hypothesis \(\neg \left(u \sim_{G} x\right)\).
Use the hypothesis \(\left(v = u \land w = x \lor v = x \land w = u\right) \land v \ne w\).
Substituting equals, we obtain \(v \sim_{G^{'}} w^{'}\).
\end{proofstep}
Using the two conjuncts, we obtain \(v \sim_{G^{'}} w^{'}\).
\end{proofstep}
Consider the following case.
\begin{proofstep}
Assume \(v = x \land w = u\).
Use the hypothesis \(v = x \land w = u\).
Consider the following case.
\begin{proofstep}
Assume \(v = x\).
Assume \(w = u\).
For the required implication, argue as follows.
\begin{proofstep}
Assume \(\neg \left(w \sim_{G^{'}} x\right)\).
Assume \(\forall u^{'}\, :\, V,\; u^{'} \ne w \implies x \sim_{G} u^{'} \implies x \sim_{G^{'}} u^{'}\).
Assume \(\forall x^{'}\, :\, V,\; x^{'} \ne x \implies x^{'} \sim_{G} w \implies x^{'} \sim_{G^{'}} w\).
Assume \(\neg \left(w \sim_{G} x\right)\).
Assume \(\left(v = w \land w = x \lor v = x \land w = w\right) \land v \ne w\).
For the required implication, argue as follows.
\begin{proofstep}
Assume \(\neg \left(w \sim_{G^{'}} v\right)\).
Assume \(\forall u^{'}\, :\, V,\; u^{'} \ne w \implies v \sim_{G} u^{'} \implies v \sim_{G^{'}} u^{'}\).
Assume \(\forall x^{'}\, :\, V,\; x^{'} \ne v \implies x^{'} \sim_{G} w \implies x^{'} \sim_{G^{'}} w\).
Assume \(\neg \left(w \sim_{G} v\right)\).
Assume \(\left(v = w \land w = v \lor v = v \land w = w\right) \land v \ne w\).
Use the hypothesis \(\forall u^{'}\, :\, V,\; u^{'} \ne w \implies v \sim_{G} u^{'} \implies v \sim_{G^{'}} u^{'}\).
For the required implication, argue as follows.
\begin{proofstep}
Assume \(w^{'} = w\).
For the required implication, argue as follows.
\begin{proofstep}
Assume \(\neg \left(w^{'} = w^{'}\right)\).
Assume \(\neg \left(w^{'} \sim_{G^{'}} v\right)\).
Assume \(\forall u^{'}\, :\, V,\; \neg (u^{'} = w^{'}) \implies v \sim_{G} u^{'} \implies v \sim_{G^{'}} u^{'}\).
Assume \(\forall x^{'}\, :\, V,\; \neg (x^{'} = v) \implies x^{'} \sim_{G} w^{'} \implies x^{'} \sim_{G^{'}} w^{'}\).
Assume \(\neg \left(w^{'} \sim_{G} v\right)\).
Assume \(\neg \left(v = w^{'}\right)\).
Use the hypothesis \(\neg \left(w^{'} = w^{'}\right)\).
Rewriting with \paperlink{premise:6e6f745f747275655f65715f66616c7365}{Prerequisite~P\ref{premise:6e6f745f747275655f65715f66616c7365}} in the preceding result gives \(\mathrm{False}\).
The contradiction rules out this case. Hence \(\mathrm{False}\).
\end{proofstep}
Use the hypothesis \(w^{'} = w\).
Use the hypothesis \(w \ne w^{'}\).
By the preceding argument, we obtain \(\neg \left(w = w^{'}\right)\).
Use the hypothesis \(\neg \left(w \sim_{G^{'}} v\right)\).
Use the hypothesis \(\forall u^{'}\, :\, V,\; u^{'} \ne w \implies v \sim_{G} u^{'} \implies v \sim_{G^{'}} u^{'}\).
By the preceding argument, we obtain \(\forall u^{'}\, :\, V,\; \neg \left(u^{'} = w\right) \implies v \sim_{G} u^{'} \implies v \sim_{G^{'}} u^{'}\).
Use the hypothesis \(\forall x^{'}\, :\, V,\; x^{'} \ne v \implies x^{'} \sim_{G} w \implies x^{'} \sim_{G^{'}} w\).
By the preceding argument, we obtain \(\forall x^{'}\, :\, V,\; \neg \left(x^{'} = v\right) \implies x^{'} \sim_{G} w \implies x^{'} \sim_{G^{'}} w\).
Use the hypothesis \(\neg \left(w \sim_{G} v\right)\).
Use the hypothesis \(\left(v = w \land w = v \lor v = v \land w = w\right) \land v \ne w\).
Rewriting with \paperlink{premise:616e645f73656c66}{Prerequisite~P\ref{premise:616e645f73656c66}}, \paperlink{premise:6f725f74727565}{Prerequisite~P\ref{premise:6f725f74727565}} and \paperlink{premise:747275655f616e64}{Prerequisite~P\ref{premise:747275655f616e64}} in the preceding result gives \(\neg \left(v = w\right)\).
Substituting equals, we obtain \(\mathrm{False}\).
\end{proofstep}
By \paperlink{premise:4165736f702e4275696c74696e52756c65732e6e6f745f696e74726f}{Prerequisite~P\ref{premise:4165736f702e4275696c74696e52756c65732e6e6f745f696e74726f}}, we obtain \(\neg \left(w^{'} = w\right)\).
Use the hypothesis \(v \sim_{G} w^{'}\).
By the preceding argument, we obtain \(v \sim_{G^{'}} w^{'}\).
\end{proofstep}
Use the hypothesis \(v = x\).
Use the hypothesis \(\neg \left(w \sim_{G^{'}} x\right)\).
Use the hypothesis \(\forall u^{'}\, :\, V,\; u^{'} \ne w \implies x \sim_{G} u^{'} \implies x \sim_{G^{'}} u^{'}\).
Use the hypothesis \(\forall x^{'}\, :\, V,\; x^{'} \ne x \implies x^{'} \sim_{G} w \implies x^{'} \sim_{G^{'}} w\).
Use the hypothesis \(\neg \left(w \sim_{G} x\right)\).
Use the hypothesis \(\left(v = w \land w = x \lor v = x \land w = w\right) \land v \ne w\).
Substituting equals, we obtain \(v \sim_{G^{'}} w^{'}\).
\end{proofstep}
Use the hypothesis \(w = u\).
Use the hypothesis \(\neg \left(u \sim_{G^{'}} x\right)\).
Use the hypothesis \(\forall u^{'}\, :\, V,\; u^{'} \ne u \implies x \sim_{G} u^{'} \implies x \sim_{G^{'}} u^{'}\).
Use the hypothesis \(\forall x^{'}\, :\, V,\; x^{'} \ne x \implies x^{'} \sim_{G} u \implies x^{'} \sim_{G^{'}} u\).
Use the hypothesis \(\neg \left(u \sim_{G} x\right)\).
Use the hypothesis \(\left(v = u \land w = x \lor v = x \land w = u\right) \land v \ne w\).
Substituting equals, we obtain \(v \sim_{G^{'}} w^{'}\).
\end{proofstep}
Using the two conjuncts, we obtain \(v \sim_{G^{'}} w^{'}\).
\end{proofstep}
The two alternatives exhaust the cases. Hence \(v \sim_{G^{'}} w^{'}\).
Rewriting with \paperlink{premise:65715f66616c7365}{Prerequisite~P\ref{premise:65715f66616c7365}} (using \(\neg \left(u \sim_{G^{'}} x\right)\)), \paperlink{premise:66616c73655f696666}{Prerequisite~P\ref{premise:66616c73655f696666}} and \paperlink{premise:436c6173736963616c2e6e6f745f6e6f74}{Prerequisite~P\ref{premise:436c6173736963616c2e6e6f745f6e6f74}} in the preceding result gives \(u \sim_{G^{'}} x \iff \neg \left(v \sim_{G^{'}} w^{'}\right)\).
Rewriting with \paperlink{premise:53696d706c6547726170682e61646a5f636f6e67725f6f665f73796d32}{Prerequisite~P\ref{premise:53696d706c6547726170682e61646a5f636f6e67725f6f665f73796d32}}, \paperlink{premise:53796d322e6571}{Prerequisite~P\ref{premise:53796d322e6571}}, \paperlink{premise:53796d322e72656c5f69666627}{Prerequisite~P\ref{premise:53796d322e72656c5f69666627}}, \paperlink{premise:50726f642e6d6b2e696e6a4571}{Prerequisite~P\ref{premise:50726f642e6d6b2e696e6a4571}} and \(v = u \land w = x \lor v = x \land w = u\) in the preceding result gives \(v \sim_{G^{'}} w \iff \neg \left(v \sim_{G^{'}} w^{'}\right)\).
\end{proofstep}
Consider the following case.
\begin{proofstep}
Assume \(\left(v = u \land w^{'} = x \lor v = x \land w^{'} = u\right) \land v \ne w^{'}\).
Use the hypothesis \(\left(v = u \land w = x \lor v = x \land w = u\right) \land v \ne w\).
By the preceding argument, we obtain \(\left(v = u \land w = x \lor v = x \land w = u\right) \land \neg \left(v = w\right)\).
Consider the following case.
\begin{proofstep}
Assume \(v = u \land w = x \lor v = x \land w = u\).
Assume \(\neg \left(v = w\right)\).
Use the hypothesis \(\left(v = u \land w^{'} = x \lor v = x \land w^{'} = u\right) \land v \ne w^{'}\).
By the preceding argument, we obtain \(\left(v = u \land w^{'} = x \lor v = x \land w^{'} = u\right) \land \neg \left(v = w^{'}\right)\).
Consider the following case.
\begin{proofstep}
Assume \(v = u \land w^{'} = x \lor v = x \land w^{'} = u\).
Assume \(\neg \left(v = w^{'}\right)\).
Use the hypothesis \(v = u \land w = x \lor v = x \land w = u\).
Consider the following case.
\begin{proofstep}
Assume \(v = u \land w = x\).
Use the hypothesis \(v = u \land w^{'} = x \lor v = x \land w^{'} = u\).
Consider the following case.
\begin{proofstep}
Assume \(v = u \land w^{'} = x\).
Use the hypothesis \(v = u \land w = x\).
Taking the second conjunct gives \(w = x\).
Use the hypothesis \(v = u \land w^{'} = x\).
Taking the second conjunct gives \(w^{'} = x\).
Use the hypothesis \(w \ne w^{'}\).
By the preceding argument, we obtain \(\neg \left(w = w^{'}\right)\).
Rewriting with \(w = x\), \(w^{'} = x\) and \paperlink{premise:6e6f745f747275655f65715f66616c7365}{Prerequisite~P\ref{premise:6e6f745f747275655f65715f66616c7365}} in the preceding result gives \(\mathrm{False}\).
The contradiction rules out this case. Hence \(v \sim_{G^{'}} w \iff \neg \left(v \sim_{G^{'}} w^{'}\right)\).
\end{proofstep}
Consider the following case.
\begin{proofstep}
Assume \(v = x \land w^{'} = u\).
Use the hypothesis \(v = u \land w = x\).
Taking the first conjunct gives \(v = u\).
Use the hypothesis \(v = x \land w^{'} = u\).
Taking the second conjunct gives \(w^{'} = u\).
Use the hypothesis \(\neg \left(v = w^{'}\right)\).
Rewriting with \(v = u\), \(w^{'} = u\) and \paperlink{premise:6e6f745f747275655f65715f66616c7365}{Prerequisite~P\ref{premise:6e6f745f747275655f65715f66616c7365}} in the preceding result gives \(\mathrm{False}\).
The contradiction rules out this case. Hence \(v \sim_{G^{'}} w \iff \neg \left(v \sim_{G^{'}} w^{'}\right)\).
\end{proofstep}
The two alternatives exhaust the cases. Hence \(v \sim_{G^{'}} w \iff \neg \left(v \sim_{G^{'}} w^{'}\right)\).
\end{proofstep}
Consider the following case.
\begin{proofstep}
Assume \(v = x \land w = u\).
Use the hypothesis \(v = u \land w^{'} = x \lor v = x \land w^{'} = u\).
Consider the following case.
\begin{proofstep}
Assume \(v = u \land w^{'} = x\).
Use the hypothesis \(v = x \land w = u\).
Taking the first conjunct gives \(v = x\).
Use the hypothesis \(v = u \land w^{'} = x\).
Taking the second conjunct gives \(w^{'} = x\).
Use the hypothesis \(\neg \left(v = w^{'}\right)\).
Rewriting with \(v = x\), \(w^{'} = x\) and \paperlink{premise:6e6f745f747275655f65715f66616c7365}{Prerequisite~P\ref{premise:6e6f745f747275655f65715f66616c7365}} in the preceding result gives \(\mathrm{False}\).
The contradiction rules out this case. Hence \(v \sim_{G^{'}} w \iff \neg \left(v \sim_{G^{'}} w^{'}\right)\).
\end{proofstep}
Consider the following case.
\begin{proofstep}
Assume \(v = x \land w^{'} = u\).
Use the hypothesis \(v = x \land w = u\).
Taking the second conjunct gives \(w = u\).
Use the hypothesis \(v = x \land w^{'} = u\).
Taking the second conjunct gives \(w^{'} = u\).
Use the hypothesis \(w \ne w^{'}\).
By the preceding argument, we obtain \(\neg \left(w = w^{'}\right)\).
Rewriting with \(w = u\), \(w^{'} = u\) and \paperlink{premise:6e6f745f747275655f65715f66616c7365}{Prerequisite~P\ref{premise:6e6f745f747275655f65715f66616c7365}} in the preceding result gives \(\mathrm{False}\).
The contradiction rules out this case. Hence \(v \sim_{G^{'}} w \iff \neg \left(v \sim_{G^{'}} w^{'}\right)\).
\end{proofstep}
The two alternatives exhaust the cases. Hence \(v \sim_{G^{'}} w \iff \neg \left(v \sim_{G^{'}} w^{'}\right)\).
\end{proofstep}
The two alternatives exhaust the cases. Hence \(v \sim_{G^{'}} w \iff \neg \left(v \sim_{G^{'}} w^{'}\right)\).
\end{proofstep}
Using the two conjuncts, we obtain \(v \sim_{G^{'}} w \iff \neg \left(v \sim_{G^{'}} w^{'}\right)\).
\end{proofstep}
Using the two conjuncts, we obtain \(v \sim_{G^{'}} w \iff \neg \left(v \sim_{G^{'}} w^{'}\right)\).
\end{proofstep}
The two alternatives exhaust the cases. Hence \(v \sim_{G^{'}} w \iff \neg \left(v \sim_{G^{'}} w^{'}\right)\).
\end{proofstep}
The two alternatives exhaust the cases. Hence \(v \sim_{G^{'}} w \iff \neg \left(v \sim_{G^{'}} w^{'}\right)\).
\end{proofstep}
The proposition or its negation holds. Hence \(G \cup \operatorname{edge}\left(u, x\right)\text{ alternates with respect to }G^{'}\).
\end{proof}

\begin{theorem}[{Symm diff}]\label{theorem:44}
Let \(V\) be a type. Let \(G\) and \(G^{'}\) be simple graphs on \(V\). Let \(M\) be a subgraph of \(G\). Assume \(M\text{ is a perfect matching}\). Assume \(G^{'}\text{ alternates with respect to }\operatorname{spanning\,graph}(M)\). Assume \(\text{every nonisolated vertex of }G^{'}\text{ has degree two}\). Then
\begin{center}
\(\displaystyle \operatorname{full\,subgraph}(\operatorname{spanning\,graph}(M) \mathbin{\triangle} G^{'})\text{ is a perfect matching}.\)
\end{center}
\end{theorem}

\begin{proof}
Fix \(v\) in \(V\).
Use the hypothesis \(M\text{ is a perfect matching}\).
Use the hypothesis \(M\text{ is a perfect matching}\).
Taking the second conjunct, we obtain \(v \in V(M)\).
Taking the first conjunct, we obtain \(\exists! w\, :\, V,\; v \sim_{M} w\).
Consider the following case.
\begin{proofstep}
Fix \(w\) in \(V\).
Assume \(v \sim_{M} w \land \left(\forall y\, :\, V,\; v \sim_{M} y \implies y = w\right)\).
Distinguish whether \(v \sim_{G^{'}} w\) holds.
Consider the following case.
\begin{proofstep}
Assume \(v \sim_{G^{'}} w\).
By \paperlink{theorem:30}{Lemma~\ref{theorem:30}}, we obtain \(\exists w^{'}\, :\, V,\; w \ne w^{'} \land v \sim_{G^{'}} w^{'}\).
Consider the following case.
\begin{proofstep}
Fix \(w^{'}\) in \(V\).
Assume \(w \ne w^{'} \land v \sim_{G^{'}} w^{'}\).
Use the hypothesis \(w \ne w^{'} \land v \sim_{G^{'}} w^{'}\).
Taking the first conjunct, we obtain \(w \ne w^{'}\).
Use the hypothesis \(v \sim_{G^{'}} w\).
Use the hypothesis \(w \ne w^{'} \land v \sim_{G^{'}} w^{'}\).
Taking the second conjunct, we obtain \(v \sim_{G^{'}} w^{'}\).
Applying the hypothesis \(G^{'}\text{ alternates with respect to }\operatorname{spanning\,graph}(M)\), we obtain \(v \sim_{M} w \iff \neg \left(v \sim_{M} w^{'}\right)\).
Choose the witness \(w^{'}\).
The substitution uses the following argument.
\begin{proofstep}
Fix \(y\) in \(V\).
By \paperlink{premise:53696d706c6547726170682e53756267726170682e746f705f61646a}{Prerequisite~P\ref{premise:53696d706c6547726170682e53756267726170682e746f705f61646a}}, we obtain \(v \sim_{\operatorname{full\,subgraph}((\operatorname{spanning\,graph}(M) \setminus G^{'}) \cup (G^{'} \setminus \operatorname{spanning\,graph}(M)))} y \iff v \sim_{(\operatorname{spanning\,graph}(M) \setminus G^{'}) \cup (G^{'} \setminus \operatorname{spanning\,graph}(M))} y\).
By propositional extensionality, we obtain \((v \sim_{\operatorname{full\,subgraph}((\operatorname{spanning\,graph}(M) \setminus G^{'}) \cup (G^{'} \setminus \operatorname{spanning\,graph}(M)))} y) = (v \sim_{(\operatorname{spanning\,graph}(M) \setminus G^{'}) \cup (G^{'} \setminus \operatorname{spanning\,graph}(M))} y)\).
By \paperlink{premise:53696d706c6547726170682e7375705f61646a}{Prerequisite~P\ref{premise:53696d706c6547726170682e7375705f61646a}} with \(x = \operatorname{spanning\,graph}(M) \setminus G^{'}\) and \(G^{'} \setminus \operatorname{spanning\,graph}(M)\) in place of its parameter \(y\) and \(y\) in place of its parameter \(w\), we obtain \(v \sim_{(\operatorname{spanning\,graph}(M) \setminus G^{'}) \cup (G^{'} \setminus \operatorname{spanning\,graph}(M))} y \iff v \sim_{\operatorname{spanning\,graph}(M) \setminus G^{'}} y \lor v \sim_{G^{'} \setminus \operatorname{spanning\,graph}(M)} y\).
By propositional extensionality, we obtain \((v \sim_{(\operatorname{spanning\,graph}(M) \setminus G^{'}) \cup (G^{'} \setminus \operatorname{spanning\,graph}(M))} y) = (v \sim_{\operatorname{spanning\,graph}(M) \setminus G^{'}} y \lor v \sim_{G^{'} \setminus \operatorname{spanning\,graph}(M)} y)\).
By \paperlink{premise:53696d706c6547726170682e73646966665f61646a}{Prerequisite~P\ref{premise:53696d706c6547726170682e73646966665f61646a}} with \(x = \operatorname{spanning\,graph}(M)\) and \(G^{'}\) in place of its parameter \(y\) and \(y\) in place of its parameter \(w\), we obtain \(v \sim_{\operatorname{spanning\,graph}(M) \setminus G^{'}} y \iff v \sim_{\operatorname{spanning\,graph}(M)} y \land \neg (v \sim_{G^{'}} y)\).
By propositional extensionality, we obtain \((v \sim_{\operatorname{spanning\,graph}(M) \setminus G^{'}} y) = (v \sim_{\operatorname{spanning\,graph}(M)} y \land \neg (v \sim_{G^{'}} y))\).
Applying the same function to equal arguments, we obtain \((y_{1} \mapsto v \sim_{\operatorname{spanning\,graph}(M) \setminus G^{'}} y \lor y_{1}) = (y_{1} \mapsto v \sim_{\operatorname{spanning\,graph}(M)} y \land \neg (v \sim_{G^{'}} y) \lor y_{1})\).
By \paperlink{premise:53696d706c6547726170682e73646966665f61646a}{Prerequisite~P\ref{premise:53696d706c6547726170682e73646966665f61646a}} with \(x = G^{'}\) and \(\operatorname{spanning\,graph}(M)\) in place of its parameter \(y\) and \(y\) in place of its parameter \(w\), we obtain \(v \sim_{G^{'} \setminus \operatorname{spanning\,graph}(M)} y \iff v \sim_{G^{'}} y \land \neg (v \sim_{\operatorname{spanning\,graph}(M)} y)\).
By propositional extensionality, we obtain \((v \sim_{G^{'} \setminus \operatorname{spanning\,graph}(M)} y) = (v \sim_{G^{'}} y \land \neg (v \sim_{\operatorname{spanning\,graph}(M)} y))\).
By congruence, we obtain \((v \sim_{\operatorname{spanning\,graph}(M) \setminus G^{'}} y \lor v \sim_{G^{'} \setminus \operatorname{spanning\,graph}(M)} y) = (v \sim_{\operatorname{spanning\,graph}(M)} y \land \neg (v \sim_{G^{'}} y) \lor v \sim_{G^{'}} y \land \neg (v \sim_{\operatorname{spanning\,graph}(M)} y))\).
Chaining the equalities, we obtain \((v \sim_{(\operatorname{spanning\,graph}(M) \setminus G^{'}) \cup (G^{'} \setminus \operatorname{spanning\,graph}(M))} y) = (v \sim_{\operatorname{spanning\,graph}(M)} y \land \neg (v \sim_{G^{'}} y) \lor v \sim_{G^{'}} y \land \neg (v \sim_{\operatorname{spanning\,graph}(M)} y))\).
Chaining the equalities, we obtain \((v \sim_{\operatorname{full\,subgraph}((\operatorname{spanning\,graph}(M) \setminus G^{'}) \cup (G^{'} \setminus \operatorname{spanning\,graph}(M)))} y) = (v \sim_{\operatorname{spanning\,graph}(M)} y \land \neg (v \sim_{G^{'}} y) \lor v \sim_{G^{'}} y \land \neg (v \sim_{\operatorname{spanning\,graph}(M)} y))\).
By definition, we obtain \(\left(y = w^{'}\right) = \left(y = w^{'}\right)\).
By \paperlink{premise:696d706c6965735f636f6e6772}{Prerequisite~P\ref{premise:696d706c6965735f636f6e6772}} with \(p_{1} = v \sim_{\operatorname{full\,subgraph}((\operatorname{spanning\,graph}(M) \setminus G^{'}) \cup (G^{'} \setminus \operatorname{spanning\,graph}(M)))} y\) and \(p_{2} = v \sim_{\operatorname{spanning\,graph}(M)} y \land \neg (v \sim_{G^{'}} y) \lor v \sim_{G^{'}} y \land \neg (v \sim_{\operatorname{spanning\,graph}(M)} y)\) and \(q_{1} = y = w^{'}\) and \(q_{2} = y = w^{'}\), we obtain \((v \sim_{\operatorname{full\,subgraph}((\operatorname{spanning\,graph}(M) \setminus G^{'}) \cup (G^{'} \setminus \operatorname{spanning\,graph}(M)))} y \implies y = w^{'}) = (v \sim_{\operatorname{spanning\,graph}(M)} y \land \neg (v \sim_{G^{'}} y) \lor v \sim_{G^{'}} y \land \neg (v \sim_{\operatorname{spanning\,graph}(M)} y) \implies y = w^{'})\).
\end{proofstep}
Fix \(y\) in \(V\).
Assume \(v \sim_{M} y \land \neg (v \sim_{G^{'}} y) \lor v \sim_{G^{'}} y \land \neg (v \sim_{M} y)\).
Use the hypothesis \(v \sim_{M} y \land \neg (v \sim_{G^{'}} y) \lor v \sim_{G^{'}} y \land \neg (v \sim_{M} y)\).
Consider the following case.
\begin{proofstep}
Assume \(v \sim_{M} y \land \neg \left(v \sim_{G^{'}} y\right)\).
The substitution uses the following argument.
\begin{proofstep}
Use the hypothesis \(v \sim_{M} w \land \left(\forall y\, :\, V,\; v \sim_{M} y \implies y = w\right)\).
Rewriting with \(v \sim_{M} w \iff \neg \left(v \sim_{M} w^{'}\right)\) in the preceding result gives \(\neg (v \sim_{M} w^{'}) \land (\forall y_{1}\, :\, V,\; v \sim_{M} y_{1} \implies y_{1} = w)\).
By the preceding argument, we obtain \(\neg (v \sim_{M} w^{'}) \land (\forall y_{1}\, :\, V,\; v \sim_{M} y_{1} \implies y_{1} = w)\).
Taking the second conjunct gives \(\forall y_{1}\, :\, V,\; v \sim_{M} y_{1} \implies y_{1} = w\).
The substitution uses the following argument.
\begin{proofstep}
Use the hypothesis \(v \sim_{M} w \land \left(\forall y\, :\, V,\; v \sim_{M} y \implies y = w\right)\).
Rewriting with \(v \sim_{M} w \iff \neg \left(v \sim_{M} w^{'}\right)\) in the preceding result gives \(\neg (v \sim_{M} w^{'}) \land (\forall y_{1}\, :\, V,\; v \sim_{M} y_{1} \implies y_{1} = w)\).
By the preceding argument, we obtain \(\neg (v \sim_{M} w^{'}) \land (\forall y_{1}\, :\, V,\; v \sim_{M} y_{1} \implies y_{1} = w)\).
Taking the second conjunct gives \(\forall y_{1}\, :\, V,\; v \sim_{M} y_{1} \implies y_{1} = w\).
Use the hypothesis \(v \sim_{M} y \land \neg \left(v \sim_{G^{'}} y\right)\).
Taking the first conjunct gives \(v \sim_{M} y\).
Applying this conjunct to \(y\) gives \(y = w\).
\end{proofstep}
Use the hypothesis \(v \sim_{M} w \land \left(\forall y\, :\, V,\; v \sim_{M} y \implies y = w\right)\).
Rewriting with \(v \sim_{M} w \iff \neg \left(v \sim_{M} w^{'}\right)\) in the preceding result gives \(\neg (v \sim_{M} w^{'}) \land (\forall y_{1}\, :\, V,\; v \sim_{M} y_{1} \implies y_{1} = w)\).
By the preceding argument, we obtain \(\neg (v \sim_{M} w^{'}) \land (\forall y_{1}\, :\, V,\; v \sim_{M} y_{1} \implies y_{1} = w)\).
Taking the first conjunct gives \(\neg \left(v \sim_{M} w^{'}\right)\).
Rewriting with \paperlink{premise:65715f66616c7365}{Prerequisite~P\ref{premise:65715f66616c7365}}, \(\neg \left(v \sim_{M} w^{'}\right)\), \paperlink{premise:6e6f745f66616c73655f65715f74727565}{Prerequisite~P\ref{premise:6e6f745f66616c73655f65715f74727565}}, \paperlink{premise:6966665f74727565}{Prerequisite~P\ref{premise:6966665f74727565}}, \(v \sim_{M} w \iff \neg \left(v \sim_{M} w^{'}\right)\) and the auxiliary result established above gives \(v \sim_{M} y\).
Applying this conjunct to \(y\) gives \(y = w\).
\end{proofstep}
Use the hypothesis \(w \ne w^{'} \land v \sim_{G^{'}} w^{'}\).
By the preceding argument, we obtain \(\neg \left(w = w^{'}\right) \land v \sim_{G^{'}} w^{'}\).
By the preceding argument, we obtain \(\neg \left(w = w^{'}\right) \land v \sim_{G^{'}} w^{'}\).
Taking the first conjunct gives \(\neg \left(w = w^{'}\right)\).
Use the hypothesis \(v \sim_{M} w \land \left(\forall y\, :\, V,\; v \sim_{M} y \implies y = w\right)\), which simplifies to \((w_{1} \mapsto v \sim_{M} w_{1})(w) \land (\forall y_{1}\, :\, V,\; v \sim_{M} y_{1} \implies y_{1} = w)\).
Rewriting with \(v \sim_{M} w \iff \neg \left(v \sim_{M} w^{'}\right)\) in the preceding result gives \(\neg (v \sim_{M} w^{'}) \land (\forall y_{1}\, :\, V,\; v \sim_{M} y_{1} \implies y_{1} = w)\).
Consider the following case.
\begin{proofstep}
Assume \(\neg \left(v \sim_{M} w^{'}\right)\).
Assume \(\forall y_{1}\, :\, V,\; v \sim_{M} y_{1} \implies y_{1} = w\).
Use the hypothesis \(w \ne w^{'} \land v \sim_{G^{'}} w^{'}\).
By the preceding argument, we obtain \(\neg \left(w = w^{'}\right) \land v \sim_{G^{'}} w^{'}\).
Consider the following case.
\begin{proofstep}
Assume \(\neg \left(w = w^{'}\right)\).
Assume \(v \sim_{G^{'}} w^{'}\).
Use the hypothesis \(v \sim_{M} y \land \neg \left(v \sim_{G^{'}} y\right)\).
Consider the following case.
\begin{proofstep}
Assume \(v \sim_{M} y\).
Assume \(\neg \left(v \sim_{G^{'}} y\right)\).
Use the hypothesis \(v \sim_{M} w \land \left(\forall y\, :\, V,\; v \sim_{M} y \implies y = w\right)\).
Rewriting with \(v \sim_{M} w \iff \neg \left(v \sim_{M} w^{'}\right)\) in the preceding result gives \(\neg (v \sim_{M} w^{'}) \land (\forall y_{1}\, :\, V,\; v \sim_{M} y_{1} \implies y_{1} = w)\).
By the preceding argument, we obtain \(\neg (v \sim_{M} w^{'}) \land (\forall y_{1}\, :\, V,\; v \sim_{M} y_{1} \implies y_{1} = w)\).
Taking the first conjunct gives \(\neg \left(v \sim_{M} w^{'}\right)\).
Use the hypothesis \(\neg \left(v \sim_{G^{'}} y\right)\).
Rewriting with \(\forall y_{1}\, :\, V,\; v \sim_{M} y_{1} \implies y_{1} = w\), \(v \sim_{M} y\), \paperlink{premise:65715f66616c7365}{Prerequisite~P\ref{premise:65715f66616c7365}}, \(\neg \left(v \sim_{M} w^{'}\right)\), \paperlink{premise:6e6f745f66616c73655f65715f74727565}{Prerequisite~P\ref{premise:6e6f745f66616c73655f65715f74727565}}, \paperlink{premise:6966665f74727565}{Prerequisite~P\ref{premise:6966665f74727565}}, \(v \sim_{M} w \iff \neg \left(v \sim_{M} w^{'}\right)\), \(v \sim_{G^{'}} w\) and \paperlink{premise:6e6f745f747275655f65715f66616c7365}{Prerequisite~P\ref{premise:6e6f745f747275655f65715f66616c7365}} in the preceding result gives \(\mathrm{False}\).
The contradiction rules out this case. Hence \(\mathrm{False}\).
\end{proofstep}
Using the two conjuncts, we obtain \(\mathrm{False}\).
\end{proofstep}
Using the two conjuncts, we obtain \(\mathrm{False}\).
\end{proofstep}
Using the two conjuncts, we obtain \(\mathrm{False}\).
Rewriting with \paperlink{premise:65715f66616c7365}{Prerequisite~P\ref{premise:65715f66616c7365}}, \(\neg \left(w = w^{'}\right)\) and the auxiliary result established above in the preceding result gives \(y = w^{'}\).
\end{proofstep}
Consider the following case.
\begin{proofstep}
Assume \(v \sim_{G^{'}} y \land \neg \left(v \sim_{M} y\right)\).
Use the hypothesis \(v \sim_{G^{'}} y \land \neg \left(v \sim_{M} y\right)\).
Taking the first conjunct, we obtain \(v \sim_{G^{'}} y\).
By \paperlink{theorem:30}{Lemma~\ref{theorem:30}}, we obtain \(\exists w^{'}_{1}\, :\, V,\; y \ne w^{'}_{1} \land v \sim_{G^{'}} w^{'}_{1}\).
Consider the following case.
\begin{proofstep}
Fix \(w^{''}\) in \(V\).
Assume \(y \ne w^{''} \land v \sim_{G^{'}} w^{''}\).
Consider the following case.
\begin{proofstep}
Assume \(\neg \left(y = w^{'}\right)\).
The substitution uses the following argument.
\begin{proofstep}
Use the hypothesis \(v \sim_{M} w \land \left(\forall y\, :\, V,\; v \sim_{M} y \implies y = w\right)\).
Rewriting with \(v \sim_{M} w \iff \neg \left(v \sim_{M} w^{'}\right)\) in the preceding result gives \(\neg (v \sim_{M} w^{'}) \land (\forall y_{1}\, :\, V,\; v \sim_{M} y_{1} \implies y_{1} = w)\).
By the preceding argument, we obtain \(\neg (v \sim_{M} w^{'}) \land (\forall y_{1}\, :\, V,\; v \sim_{M} y_{1} \implies y_{1} = w)\).
Taking the second conjunct gives \(\forall y_{1}\, :\, V,\; v \sim_{M} y_{1} \implies y_{1} = w\).
The substitution uses the following argument.
\begin{proofstep}
Use the hypothesis \(v \sim_{M} w \land \left(\forall y\, :\, V,\; v \sim_{M} y \implies y = w\right)\).
Rewriting with \(v \sim_{M} w \iff \neg \left(v \sim_{M} w^{'}\right)\) in the preceding result gives \(\neg (v \sim_{M} w^{'}) \land (\forall y_{1}\, :\, V,\; v \sim_{M} y_{1} \implies y_{1} = w)\).
By the preceding argument, we obtain \(\neg (v \sim_{M} w^{'}) \land (\forall y_{1}\, :\, V,\; v \sim_{M} y_{1} \implies y_{1} = w)\).
Taking the second conjunct gives \(\forall y_{1}\, :\, V,\; v \sim_{M} y_{1} \implies y_{1} = w\).
The substitution uses the following argument.
\begin{proofstep}
Use the hypothesis \(\neg \left(y = w^{'}\right)\).
Use the hypothesis \(v \sim_{G^{'}} y \land \neg \left(v \sim_{M} y\right)\).
Taking the first conjunct, we obtain \(v \sim_{G^{'}} y\).
Use the hypothesis \(w \ne w^{'} \land v \sim_{G^{'}} w^{'}\).
Taking the second conjunct, we obtain \(v \sim_{G^{'}} w^{'}\).
Applying the hypothesis \(G^{'}\text{ alternates with respect to }\operatorname{spanning\,graph}(M)\), we obtain \(v \sim_{M} y \iff \neg \left(v \sim_{M} w^{'}\right)\).
Use the previously proved fact \(v \sim_{M} y \iff \neg \left(v \sim_{M} w^{'}\right)\).
\end{proofstep}
Use the hypothesis \(v \sim_{M} w \land \left(\forall y\, :\, V,\; v \sim_{M} y \implies y = w\right)\).
Rewriting with \(v \sim_{M} w \iff \neg \left(v \sim_{M} w^{'}\right)\) in the preceding result gives \(\neg (v \sim_{M} w^{'}) \land (\forall y_{1}\, :\, V,\; v \sim_{M} y_{1} \implies y_{1} = w)\).
By the preceding argument, we obtain \(\neg (v \sim_{M} w^{'}) \land (\forall y_{1}\, :\, V,\; v \sim_{M} y_{1} \implies y_{1} = w)\).
Taking the first conjunct gives \(\neg \left(v \sim_{M} w^{'}\right)\).
Rewriting with \paperlink{premise:65715f66616c7365}{Prerequisite~P\ref{premise:65715f66616c7365}}, \(\neg \left(v \sim_{M} w^{'}\right)\), \paperlink{premise:6e6f745f66616c73655f65715f74727565}{Prerequisite~P\ref{premise:6e6f745f66616c73655f65715f74727565}} and the auxiliary result established above gives \(v \sim_{M} y\).
Applying this conjunct to \(y\) gives \(y = w\).
\end{proofstep}
Use the hypothesis \(v \sim_{M} w \land \left(\forall y\, :\, V,\; v \sim_{M} y \implies y = w\right)\).
Rewriting with \(v \sim_{M} w \iff \neg \left(v \sim_{M} w^{'}\right)\) in the preceding result gives \(\neg (v \sim_{M} w^{'}) \land (\forall y_{1}\, :\, V,\; v \sim_{M} y_{1} \implies y_{1} = w)\).
By the preceding argument, we obtain \(\neg (v \sim_{M} w^{'}) \land (\forall y_{1}\, :\, V,\; v \sim_{M} y_{1} \implies y_{1} = w)\).
Taking the first conjunct gives \(\neg \left(v \sim_{M} w^{'}\right)\).
Rewriting with \paperlink{premise:65715f66616c7365}{Prerequisite~P\ref{premise:65715f66616c7365}}, \(\neg \left(v \sim_{M} w^{'}\right)\), \paperlink{premise:6e6f745f66616c73655f65715f74727565}{Prerequisite~P\ref{premise:6e6f745f66616c73655f65715f74727565}}, \paperlink{premise:6966665f74727565}{Prerequisite~P\ref{premise:6966665f74727565}}, \(v \sim_{M} w \iff \neg \left(v \sim_{M} w^{'}\right)\) and the auxiliary result established above gives \(v \sim_{M} y\).
Applying this conjunct to \(y\) gives \(y = w\).
\end{proofstep}
The substitution uses the following argument.
\begin{proofstep}
Use the hypothesis \(v \sim_{M} w \land \left(\forall y\, :\, V,\; v \sim_{M} y \implies y = w\right)\).
Rewriting with \(v \sim_{M} w \iff \neg \left(v \sim_{M} w^{'}\right)\) in the preceding result gives \(\neg (v \sim_{M} w^{'}) \land (\forall y_{1}\, :\, V,\; v \sim_{M} y_{1} \implies y_{1} = w)\).
By the preceding argument, we obtain \(\neg (v \sim_{M} w^{'}) \land (\forall y_{1}\, :\, V,\; v \sim_{M} y_{1} \implies y_{1} = w)\).
Taking the second conjunct gives \(\forall y_{1}\, :\, V,\; v \sim_{M} y_{1} \implies y_{1} = w\).
The substitution uses the following argument.
\begin{proofstep}
Use the hypothesis \(v \sim_{M} w \land \left(\forall y\, :\, V,\; v \sim_{M} y \implies y = w\right)\).
Rewriting with \(v \sim_{M} w \iff \neg \left(v \sim_{M} w^{'}\right)\) in the preceding result gives \(\neg (v \sim_{M} w^{'}) \land (\forall y_{1}\, :\, V,\; v \sim_{M} y_{1} \implies y_{1} = w)\).
By the preceding argument, we obtain \(\neg (v \sim_{M} w^{'}) \land (\forall y_{1}\, :\, V,\; v \sim_{M} y_{1} \implies y_{1} = w)\).
Taking the second conjunct gives \(\forall y_{1}\, :\, V,\; v \sim_{M} y_{1} \implies y_{1} = w\).
The substitution uses the following argument.
\begin{proofstep}
Use the hypothesis \(\neg \left(y = w^{'}\right)\).
Use the hypothesis \(v \sim_{G^{'}} y \land \neg \left(v \sim_{M} y\right)\).
Taking the first conjunct, we obtain \(v \sim_{G^{'}} y\).
Use the hypothesis \(w \ne w^{'} \land v \sim_{G^{'}} w^{'}\).
Taking the second conjunct, we obtain \(v \sim_{G^{'}} w^{'}\).
Applying the hypothesis \(G^{'}\text{ alternates with respect to }\operatorname{spanning\,graph}(M)\), we obtain \(v \sim_{M} y \iff \neg \left(v \sim_{M} w^{'}\right)\).
Use the previously proved fact \(v \sim_{M} y \iff \neg \left(v \sim_{M} w^{'}\right)\).
\end{proofstep}
Use the hypothesis \(v \sim_{M} w \land \left(\forall y\, :\, V,\; v \sim_{M} y \implies y = w\right)\).
Rewriting with \(v \sim_{M} w \iff \neg \left(v \sim_{M} w^{'}\right)\) in the preceding result gives \(\neg (v \sim_{M} w^{'}) \land (\forall y_{1}\, :\, V,\; v \sim_{M} y_{1} \implies y_{1} = w)\).
By the preceding argument, we obtain \(\neg (v \sim_{M} w^{'}) \land (\forall y_{1}\, :\, V,\; v \sim_{M} y_{1} \implies y_{1} = w)\).
Taking the first conjunct gives \(\neg \left(v \sim_{M} w^{'}\right)\).
Rewriting with \paperlink{premise:65715f66616c7365}{Prerequisite~P\ref{premise:65715f66616c7365}}, \(\neg \left(v \sim_{M} w^{'}\right)\), \paperlink{premise:6e6f745f66616c73655f65715f74727565}{Prerequisite~P\ref{premise:6e6f745f66616c73655f65715f74727565}} and the auxiliary result established above gives \(v \sim_{M} y\).
Applying this conjunct to \(y\) gives \(y = w\).
\end{proofstep}
Use the hypothesis \(v \sim_{M} w \land \left(\forall y\, :\, V,\; v \sim_{M} y \implies y = w\right)\).
Rewriting with \(v \sim_{M} w \iff \neg \left(v \sim_{M} w^{'}\right)\) in the preceding result gives \(\neg (v \sim_{M} w^{'}) \land (\forall y_{1}\, :\, V,\; v \sim_{M} y_{1} \implies y_{1} = w)\).
By the preceding argument, we obtain \(\neg (v \sim_{M} w^{'}) \land (\forall y_{1}\, :\, V,\; v \sim_{M} y_{1} \implies y_{1} = w)\).
Taking the first conjunct gives \(\neg \left(v \sim_{M} w^{'}\right)\).
Rewriting with \paperlink{premise:65715f66616c7365}{Prerequisite~P\ref{premise:65715f66616c7365}}, \(\neg \left(v \sim_{M} w^{'}\right)\), \paperlink{premise:6e6f745f66616c73655f65715f74727565}{Prerequisite~P\ref{premise:6e6f745f66616c73655f65715f74727565}}, \paperlink{premise:6966665f74727565}{Prerequisite~P\ref{premise:6966665f74727565}}, \(v \sim_{M} w \iff \neg \left(v \sim_{M} w^{'}\right)\) and the auxiliary result established above gives \(v \sim_{M} y\).
Applying this conjunct to \(y\) gives \(y = w\).
\end{proofstep}
Use the hypothesis \(v \sim_{M} w \land \left(\forall y\, :\, V,\; v \sim_{M} y \implies y = w\right)\).
Rewriting with \(v \sim_{M} w \iff \neg \left(v \sim_{M} w^{'}\right)\) in the preceding result gives \(\neg (v \sim_{M} w^{'}) \land (\forall y_{1}\, :\, V,\; v \sim_{M} y_{1} \implies y_{1} = w)\).
By the preceding argument, we obtain \(\neg (v \sim_{M} w^{'}) \land (\forall y_{1}\, :\, V,\; v \sim_{M} y_{1} \implies y_{1} = w)\).
Taking the first conjunct gives \(\neg \left(v \sim_{M} w^{'}\right)\).
Use the hypothesis \(v \sim_{G^{'}} y \land \neg \left(v \sim_{M} y\right)\).
Rewriting with \(v \sim_{G^{'}} w\), \paperlink{premise:65715f66616c7365}{Prerequisite~P\ref{premise:65715f66616c7365}}, \(\neg \left(v \sim_{M} w^{'}\right)\), \paperlink{premise:6e6f745f66616c73655f65715f74727565}{Prerequisite~P\ref{premise:6e6f745f66616c73655f65715f74727565}}, \paperlink{premise:6966665f74727565}{Prerequisite~P\ref{premise:6966665f74727565}}, \(v \sim_{M} w \iff \neg \left(v \sim_{M} w^{'}\right)\), \paperlink{premise:6e6f745f747275655f65715f66616c7365}{Prerequisite~P\ref{premise:6e6f745f747275655f65715f66616c7365}}, \paperlink{premise:616e645f66616c7365}{Prerequisite~P\ref{premise:616e645f66616c7365}} and the two auxiliary results established above in the preceding result gives \(\mathrm{False}\).
The contradiction rules out this case. Hence \(\mathrm{False}\).
\end{proofstep}
The contrary assumption leads to a contradiction. Hence \(y = w^{'}\).
\end{proofstep}
Using the existential hypothesis, we obtain \(y = w^{'}\).
\end{proofstep}
The two alternatives exhaust the cases. Hence \(y = w^{'}\).
Rewriting with \paperlink{premise:53696d706c6547726170682e53756267726170682e746f705f61646a}{Prerequisite~P\ref{premise:53696d706c6547726170682e53756267726170682e746f705f61646a}}, \paperlink{premise:53696d706c6547726170682e7375705f61646a}{Prerequisite~P\ref{premise:53696d706c6547726170682e7375705f61646a}}, \paperlink{premise:53696d706c6547726170682e73646966665f61646a}{Prerequisite~P\ref{premise:53696d706c6547726170682e73646966665f61646a}}, \paperlink{premise:65715f66616c7365}{Prerequisite~P\ref{premise:65715f66616c7365}}, \paperlink{premise:4966662e6d70}{Prerequisite~P\ref{premise:4966662e6d70}} (using \(v \sim_{M} w \iff \neg \left(v \sim_{M} w^{'}\right)\)), \paperlink{premise:416e642e6c656674}{Prerequisite~P\ref{premise:416e642e6c656674}} (using \(v \sim_{M} w \land \left(\forall y\, :\, V,\; v \sim_{M} y \implies y = w\right)\)), \paperlink{premise:416e642e7269676874}{Prerequisite~P\ref{premise:416e642e7269676874}} (using \(w \ne w^{'} \land v \sim_{G^{'}} w^{'}\)), \paperlink{premise:6e6f745f747275655f65715f66616c7365}{Prerequisite~P\ref{premise:6e6f745f747275655f65715f66616c7365}}, \paperlink{premise:616e645f73656c66}{Prerequisite~P\ref{premise:616e645f73656c66}}, \paperlink{premise:6e6f745f66616c73655f65715f74727565}{Prerequisite~P\ref{premise:6e6f745f66616c73655f65715f74727565}}, \paperlink{premise:6f725f74727565}{Prerequisite~P\ref{premise:6f725f74727565}}, \paperlink{premise:666f72616c6c5f636f6e6772}{Prerequisite~P\ref{premise:666f72616c6c5f636f6e6772}}, \paperlink{premise:747275655f616e64}{Prerequisite~P\ref{premise:747275655f616e64}} and the identity established above in the preceding result gives \(v \sim_{\operatorname{full\,subgraph}((\operatorname{spanning\,graph}(M) \setminus G^{'}) \cup (G^{'} \setminus \operatorname{spanning\,graph}(M)))} w^{'} \land (\forall y\, :\, V,\; v \sim_{\operatorname{full\,subgraph}((\operatorname{spanning\,graph}(M) \setminus G^{'}) \cup (G^{'} \setminus \operatorname{spanning\,graph}(M)))} y \implies y = w^{'})\).
This witness proves existence. Hence \(\exists! w_{1}\, :\, V,\; v \sim_{\operatorname{full\,subgraph}((\operatorname{spanning\,graph}(M) \setminus G^{'}) \cup (G^{'} \setminus \operatorname{spanning\,graph}(M)))} w_{1}\).
\end{proofstep}
Using the existential hypothesis, we obtain \(\exists! w_{1}\, :\, V,\; v \sim_{\operatorname{full\,subgraph}((\operatorname{spanning\,graph}(M) \setminus G^{'}) \cup (G^{'} \setminus \operatorname{spanning\,graph}(M)))} w_{1}\).
\end{proofstep}
Consider the following case.
\begin{proofstep}
Assume \(\neg \left(v \sim_{G^{'}} w\right)\).
Choose the witness \(w\).
The substitution uses the following argument.
\begin{proofstep}
Fix \(y\) in \(V\).
By \paperlink{premise:53696d706c6547726170682e53756267726170682e746f705f61646a}{Prerequisite~P\ref{premise:53696d706c6547726170682e53756267726170682e746f705f61646a}}, we obtain \(v \sim_{\operatorname{full\,subgraph}((\operatorname{spanning\,graph}(M) \setminus G^{'}) \cup (G^{'} \setminus \operatorname{spanning\,graph}(M)))} y \iff v \sim_{(\operatorname{spanning\,graph}(M) \setminus G^{'}) \cup (G^{'} \setminus \operatorname{spanning\,graph}(M))} y\).
By propositional extensionality, we obtain \((v \sim_{\operatorname{full\,subgraph}((\operatorname{spanning\,graph}(M) \setminus G^{'}) \cup (G^{'} \setminus \operatorname{spanning\,graph}(M)))} y) = (v \sim_{(\operatorname{spanning\,graph}(M) \setminus G^{'}) \cup (G^{'} \setminus \operatorname{spanning\,graph}(M))} y)\).
By \paperlink{premise:53696d706c6547726170682e7375705f61646a}{Prerequisite~P\ref{premise:53696d706c6547726170682e7375705f61646a}} with \(x = \operatorname{spanning\,graph}(M) \setminus G^{'}\) and \(G^{'} \setminus \operatorname{spanning\,graph}(M)\) in place of its parameter \(y\) and \(y\) in place of its parameter \(w\), we obtain \(v \sim_{(\operatorname{spanning\,graph}(M) \setminus G^{'}) \cup (G^{'} \setminus \operatorname{spanning\,graph}(M))} y \iff v \sim_{\operatorname{spanning\,graph}(M) \setminus G^{'}} y \lor v \sim_{G^{'} \setminus \operatorname{spanning\,graph}(M)} y\).
By propositional extensionality, we obtain \((v \sim_{(\operatorname{spanning\,graph}(M) \setminus G^{'}) \cup (G^{'} \setminus \operatorname{spanning\,graph}(M))} y) = (v \sim_{\operatorname{spanning\,graph}(M) \setminus G^{'}} y \lor v \sim_{G^{'} \setminus \operatorname{spanning\,graph}(M)} y)\).
By \paperlink{premise:53696d706c6547726170682e73646966665f61646a}{Prerequisite~P\ref{premise:53696d706c6547726170682e73646966665f61646a}} with \(x = \operatorname{spanning\,graph}(M)\) and \(G^{'}\) in place of its parameter \(y\) and \(y\) in place of its parameter \(w\), we obtain \(v \sim_{\operatorname{spanning\,graph}(M) \setminus G^{'}} y \iff v \sim_{\operatorname{spanning\,graph}(M)} y \land \neg (v \sim_{G^{'}} y)\).
By propositional extensionality, we obtain \((v \sim_{\operatorname{spanning\,graph}(M) \setminus G^{'}} y) = (v \sim_{\operatorname{spanning\,graph}(M)} y \land \neg (v \sim_{G^{'}} y))\).
Applying the same function to equal arguments, we obtain \((y_{1} \mapsto v \sim_{\operatorname{spanning\,graph}(M) \setminus G^{'}} y \lor y_{1}) = (y_{1} \mapsto v \sim_{\operatorname{spanning\,graph}(M)} y \land \neg (v \sim_{G^{'}} y) \lor y_{1})\).
By \paperlink{premise:53696d706c6547726170682e73646966665f61646a}{Prerequisite~P\ref{premise:53696d706c6547726170682e73646966665f61646a}} with \(x = G^{'}\) and \(\operatorname{spanning\,graph}(M)\) in place of its parameter \(y\) and \(y\) in place of its parameter \(w\), we obtain \(v \sim_{G^{'} \setminus \operatorname{spanning\,graph}(M)} y \iff v \sim_{G^{'}} y \land \neg (v \sim_{\operatorname{spanning\,graph}(M)} y)\).
By propositional extensionality, we obtain \((v \sim_{G^{'} \setminus \operatorname{spanning\,graph}(M)} y) = (v \sim_{G^{'}} y \land \neg (v \sim_{\operatorname{spanning\,graph}(M)} y))\).
By congruence, we obtain \((v \sim_{\operatorname{spanning\,graph}(M) \setminus G^{'}} y \lor v \sim_{G^{'} \setminus \operatorname{spanning\,graph}(M)} y) = (v \sim_{\operatorname{spanning\,graph}(M)} y \land \neg (v \sim_{G^{'}} y) \lor v \sim_{G^{'}} y \land \neg (v \sim_{\operatorname{spanning\,graph}(M)} y))\).
Chaining the equalities, we obtain \((v \sim_{(\operatorname{spanning\,graph}(M) \setminus G^{'}) \cup (G^{'} \setminus \operatorname{spanning\,graph}(M))} y) = (v \sim_{\operatorname{spanning\,graph}(M)} y \land \neg (v \sim_{G^{'}} y) \lor v \sim_{G^{'}} y \land \neg (v \sim_{\operatorname{spanning\,graph}(M)} y))\).
Chaining the equalities, we obtain \((v \sim_{\operatorname{full\,subgraph}((\operatorname{spanning\,graph}(M) \setminus G^{'}) \cup (G^{'} \setminus \operatorname{spanning\,graph}(M)))} y) = (v \sim_{\operatorname{spanning\,graph}(M)} y \land \neg (v \sim_{G^{'}} y) \lor v \sim_{G^{'}} y \land \neg (v \sim_{\operatorname{spanning\,graph}(M)} y))\).
By definition, we obtain \(\left(y = w\right) = \left(y = w\right)\).
By \paperlink{premise:696d706c6965735f636f6e6772}{Prerequisite~P\ref{premise:696d706c6965735f636f6e6772}} with \(p_{1} = v \sim_{\operatorname{full\,subgraph}((\operatorname{spanning\,graph}(M) \setminus G^{'}) \cup (G^{'} \setminus \operatorname{spanning\,graph}(M)))} y\) and \(p_{2} = v \sim_{\operatorname{spanning\,graph}(M)} y \land \neg (v \sim_{G^{'}} y) \lor v \sim_{G^{'}} y \land \neg (v \sim_{\operatorname{spanning\,graph}(M)} y)\) and \(q_{1} = y = w\) and \(q_{2} = y = w\), we obtain \((v \sim_{\operatorname{full\,subgraph}((\operatorname{spanning\,graph}(M) \setminus G^{'}) \cup (G^{'} \setminus \operatorname{spanning\,graph}(M)))} y \implies y = w) = (v \sim_{\operatorname{spanning\,graph}(M)} y \land \neg (v \sim_{G^{'}} y) \lor v \sim_{G^{'}} y \land \neg (v \sim_{\operatorname{spanning\,graph}(M)} y) \implies y = w)\).
\end{proofstep}
Fix \(y\) in \(V\).
Assume \(v \sim_{M} y \land \neg (v \sim_{G^{'}} y) \lor v \sim_{G^{'}} y \land \neg (v \sim_{M} y)\).
Use the hypothesis \(v \sim_{M} y \land \neg (v \sim_{G^{'}} y) \lor v \sim_{G^{'}} y \land \neg (v \sim_{M} y)\).
Consider the following case.
\begin{proofstep}
Assume \(v \sim_{M} y \land \neg \left(v \sim_{G^{'}} y\right)\).
Use the hypothesis \(v \sim_{M} w \land \left(\forall y\, :\, V,\; v \sim_{M} y \implies y = w\right)\).
Use the hypothesis \(v \sim_{M} y \land \neg \left(v \sim_{G^{'}} y\right)\).
Taking the first conjunct, we obtain \(v \sim_{M} y\).
Taking the second conjunct, we obtain \(y = w\).
\end{proofstep}
Consider the following case.
\begin{proofstep}
Assume \(v \sim_{G^{'}} y \land \neg \left(v \sim_{M} y\right)\).
Use the hypothesis \(v \sim_{G^{'}} y \land \neg \left(v \sim_{M} y\right)\).
Taking the first conjunct, we obtain \(v \sim_{G^{'}} y\).
By \paperlink{theorem:30}{Lemma~\ref{theorem:30}}, we obtain \(\exists w^{'}\, :\, V,\; y \ne w^{'} \land v \sim_{G^{'}} w^{'}\).
Take the witnesses from the preceding result.
\begin{proofstep}
Fix \(w^{'}\) in \(V\).
Assume \(y \ne w^{'} \land v \sim_{G^{'}} w^{'}\).
The substitution uses the following argument.
\begin{proofstep}
Use the hypothesis \(v \sim_{M} w \land \left(\forall y\, :\, V,\; v \sim_{M} y \implies y = w\right)\).
By the preceding argument, we obtain \(v \sim_{M} w \land \left(\forall y_{1}\, :\, V,\; v \sim_{M} y_{1} \implies y_{1} = w\right)\).
By the preceding argument, we obtain \(v \sim_{M} w \land \left(\forall y_{1}\, :\, V,\; v \sim_{M} y_{1} \implies y_{1} = w\right)\).
Taking the second conjunct gives \(\forall y_{1}\, :\, V,\; v \sim_{M} y_{1} \implies y_{1} = w\).
The substitution uses the following argument.
\begin{proofstep}
Use the hypothesis \(v \sim_{M} w \land \left(\forall y\, :\, V,\; v \sim_{M} y \implies y = w\right)\).
By the preceding argument, we obtain \(v \sim_{M} w \land \left(\forall y_{1}\, :\, V,\; v \sim_{M} y_{1} \implies y_{1} = w\right)\).
By the preceding argument, we obtain \(v \sim_{M} w \land \left(\forall y_{1}\, :\, V,\; v \sim_{M} y_{1} \implies y_{1} = w\right)\).
Taking the second conjunct gives \(\forall y_{1}\, :\, V,\; v \sim_{M} y_{1} \implies y_{1} = w\).
The substitution uses the following argument.
\begin{proofstep}
Use the hypothesis \(y \ne w^{'} \land v \sim_{G^{'}} w^{'}\).
Taking the first conjunct, we obtain \(y \ne w^{'}\).
Use the hypothesis \(v \sim_{G^{'}} y \land \neg \left(v \sim_{M} y\right)\).
Taking the first conjunct, we obtain \(v \sim_{G^{'}} y\).
Use the hypothesis \(y \ne w^{'} \land v \sim_{G^{'}} w^{'}\).
Taking the second conjunct, we obtain \(v \sim_{G^{'}} w^{'}\).
Applying the hypothesis \(G^{'}\text{ alternates with respect to }\operatorname{spanning\,graph}(M)\), we obtain \(v \sim_{M} y \iff \neg \left(v \sim_{M} w^{'}\right)\).
Use the previously proved fact \(v \sim_{M} y \iff \neg \left(v \sim_{M} w^{'}\right)\).
\end{proofstep}
Use the hypothesis \(v \sim_{G^{'}} y \land \neg \left(v \sim_{M} y\right)\).
Rewriting with \paperlink{premise:436c6173736963616c2e6e6f745f6e6f74}{Prerequisite~P\ref{premise:436c6173736963616c2e6e6f745f6e6f74}} and the auxiliary result established above in the preceding result gives \(v \sim_{G^{'}} y \land v \sim_{M} w^{'}\).
By the preceding argument, we obtain \(v \sim_{G^{'}} y \land v \sim_{M} w^{'}\).
Taking the second conjunct gives \(v \sim_{M} w^{'}\).
Applying this conjunct to \(w^{'}\) gives \(w^{'} = w\).
\end{proofstep}
Use the hypothesis \(v \sim_{M} w \land \left(\forall y\, :\, V,\; v \sim_{M} y \implies y = w\right)\).
By the preceding argument, we obtain \(v \sim_{M} w \land \left(\forall y_{1}\, :\, V,\; v \sim_{M} y_{1} \implies y_{1} = w\right)\).
By the preceding argument, we obtain \(v \sim_{M} w \land \left(\forall y_{1}\, :\, V,\; v \sim_{M} y_{1} \implies y_{1} = w\right)\).
Taking the first conjunct gives \(v \sim_{M} w\).
Rewriting with \(v \sim_{M} w\) and the auxiliary result established above gives \(v \sim_{M} w^{'}\).
Applying this conjunct to \(w^{'}\) gives \(w^{'} = w\).
\end{proofstep}
The substitution uses the following argument.
\begin{proofstep}
Use the hypothesis \(v \sim_{M} w \land \left(\forall y\, :\, V,\; v \sim_{M} y \implies y = w\right)\).
By the preceding argument, we obtain \(v \sim_{M} w \land \left(\forall y_{1}\, :\, V,\; v \sim_{M} y_{1} \implies y_{1} = w\right)\).
By the preceding argument, we obtain \(v \sim_{M} w \land \left(\forall y_{1}\, :\, V,\; v \sim_{M} y_{1} \implies y_{1} = w\right)\).
Taking the second conjunct gives \(\forall y_{1}\, :\, V,\; v \sim_{M} y_{1} \implies y_{1} = w\).
The substitution uses the following argument.
\begin{proofstep}
Use the hypothesis \(v \sim_{M} w \land \left(\forall y\, :\, V,\; v \sim_{M} y \implies y = w\right)\).
By the preceding argument, we obtain \(v \sim_{M} w \land \left(\forall y_{1}\, :\, V,\; v \sim_{M} y_{1} \implies y_{1} = w\right)\).
By the preceding argument, we obtain \(v \sim_{M} w \land \left(\forall y_{1}\, :\, V,\; v \sim_{M} y_{1} \implies y_{1} = w\right)\).
Taking the second conjunct gives \(\forall y_{1}\, :\, V,\; v \sim_{M} y_{1} \implies y_{1} = w\).
The substitution uses the following argument.
\begin{proofstep}
Use the hypothesis \(y \ne w^{'} \land v \sim_{G^{'}} w^{'}\).
Taking the first conjunct, we obtain \(y \ne w^{'}\).
Use the hypothesis \(v \sim_{G^{'}} y \land \neg \left(v \sim_{M} y\right)\).
Taking the first conjunct, we obtain \(v \sim_{G^{'}} y\).
Use the hypothesis \(y \ne w^{'} \land v \sim_{G^{'}} w^{'}\).
Taking the second conjunct, we obtain \(v \sim_{G^{'}} w^{'}\).
Applying the hypothesis \(G^{'}\text{ alternates with respect to }\operatorname{spanning\,graph}(M)\), we obtain \(v \sim_{M} y \iff \neg \left(v \sim_{M} w^{'}\right)\).
Use the previously proved fact \(v \sim_{M} y \iff \neg \left(v \sim_{M} w^{'}\right)\).
\end{proofstep}
Use the hypothesis \(v \sim_{G^{'}} y \land \neg \left(v \sim_{M} y\right)\).
Rewriting with \paperlink{premise:436c6173736963616c2e6e6f745f6e6f74}{Prerequisite~P\ref{premise:436c6173736963616c2e6e6f745f6e6f74}} and the auxiliary result established above in the preceding result gives \(v \sim_{G^{'}} y \land v \sim_{M} w^{'}\).
By the preceding argument, we obtain \(v \sim_{G^{'}} y \land v \sim_{M} w^{'}\).
Taking the second conjunct gives \(v \sim_{M} w^{'}\).
Applying this conjunct to \(w^{'}\) gives \(w^{'} = w\).
\end{proofstep}
Use the hypothesis \(v \sim_{M} w \land \left(\forall y\, :\, V,\; v \sim_{M} y \implies y = w\right)\).
By the preceding argument, we obtain \(v \sim_{M} w \land \left(\forall y_{1}\, :\, V,\; v \sim_{M} y_{1} \implies y_{1} = w\right)\).
By the preceding argument, we obtain \(v \sim_{M} w \land \left(\forall y_{1}\, :\, V,\; v \sim_{M} y_{1} \implies y_{1} = w\right)\).
Taking the first conjunct gives \(v \sim_{M} w\).
Rewriting with \(v \sim_{M} w\) and the auxiliary result established above gives \(v \sim_{M} w^{'}\).
Applying this conjunct to \(w^{'}\) gives \(w^{'} = w\).
\end{proofstep}
Use the hypothesis \(y \ne w^{'} \land v \sim_{G^{'}} w^{'}\).
Rewriting with \paperlink{premise:65715f66616c7365}{Prerequisite~P\ref{premise:65715f66616c7365}} (using \(\neg \left(v \sim_{G^{'}} w\right)\)), \paperlink{premise:616e645f66616c7365}{Prerequisite~P\ref{premise:616e645f66616c7365}} and the two auxiliary results established above in the preceding result gives \(\mathrm{False}\).
The contradiction rules out this case. Hence \(y = w\).
\end{proofstep}
Using the witnesses and their properties, we obtain \(y = w\).
\end{proofstep}
The two alternatives exhaust the cases. Hence \(y = w\).
Rewriting with \paperlink{premise:53696d706c6547726170682e53756267726170682e746f705f61646a}{Prerequisite~P\ref{premise:53696d706c6547726170682e53756267726170682e746f705f61646a}}, \paperlink{premise:53696d706c6547726170682e7375705f61646a}{Prerequisite~P\ref{premise:53696d706c6547726170682e7375705f61646a}}, \paperlink{premise:53696d706c6547726170682e73646966665f61646a}{Prerequisite~P\ref{premise:53696d706c6547726170682e73646966665f61646a}}, \paperlink{premise:416e642e6c656674}{Prerequisite~P\ref{premise:416e642e6c656674}} (using \(v \sim_{M} w \land \left(\forall y\, :\, V,\; v \sim_{M} y \implies y = w\right)\)), \paperlink{premise:65715f66616c7365}{Prerequisite~P\ref{premise:65715f66616c7365}} (using \(\neg \left(v \sim_{G^{'}} w\right)\)), \paperlink{premise:6e6f745f66616c73655f65715f74727565}{Prerequisite~P\ref{premise:6e6f745f66616c73655f65715f74727565}}, \paperlink{premise:616e645f73656c66}{Prerequisite~P\ref{premise:616e645f73656c66}}, \paperlink{premise:6e6f745f747275655f65715f66616c7365}{Prerequisite~P\ref{premise:6e6f745f747275655f65715f66616c7365}}, \paperlink{premise:6f725f66616c7365}{Prerequisite~P\ref{premise:6f725f66616c7365}}, \paperlink{premise:666f72616c6c5f636f6e6772}{Prerequisite~P\ref{premise:666f72616c6c5f636f6e6772}}, \paperlink{premise:747275655f616e64}{Prerequisite~P\ref{premise:747275655f616e64}} and the identity established above in the preceding result gives \(v \sim_{\operatorname{full\,subgraph}((\operatorname{spanning\,graph}(M) \setminus G^{'}) \cup (G^{'} \setminus \operatorname{spanning\,graph}(M)))} w \land (\forall y\, :\, V,\; v \sim_{\operatorname{full\,subgraph}((\operatorname{spanning\,graph}(M) \setminus G^{'}) \cup (G^{'} \setminus \operatorname{spanning\,graph}(M)))} y \implies y = w)\).
This witness proves existence. Hence \(\exists x\, :\, V,\; v \sim_{\operatorname{full\,subgraph}((\operatorname{spanning\,graph}(M) \setminus G^{'}) \cup (G^{'} \setminus \operatorname{spanning\,graph}(M)))} x \land (\forall y\, :\, V,\; v \sim_{\operatorname{full\,subgraph}((\operatorname{spanning\,graph}(M) \setminus G^{'}) \cup (G^{'} \setminus \operatorname{spanning\,graph}(M)))} y \implies y = x)\).
\end{proofstep}
The proposition or its negation holds. Hence \(\exists! w_{1}\, :\, V,\; v \sim_{\operatorname{full\,subgraph}((\operatorname{spanning\,graph}(M) \setminus G^{'}) \cup (G^{'} \setminus \operatorname{spanning\,graph}(M)))} w_{1}\).
\end{proofstep}
Using the existential hypothesis, we obtain \(\exists! w\, :\, V,\; v \sim_{\operatorname{full\,subgraph}((\operatorname{spanning\,graph}(M) \setminus G^{'}) \cup (G^{'} \setminus \operatorname{spanning\,graph}(M)))} w\).
By the preceding argument, we obtain \(\exists! w\, :\, V,\; v \sim_{\operatorname{full\,subgraph}(\operatorname{spanning\,graph}(M) \mathbin{\triangle} G^{'})} w\).
Rewriting with \paperlink{theorem:19}{Lemma~\ref{theorem:19}} in the preceding result gives \(\operatorname{full\,subgraph}(\operatorname{spanning\,graph}(M) \mathbin{\triangle} G^{'})\text{ is a perfect matching}\).
\end{proof}

\begin{lemma}[{Is alternating symm diff left}]\label{theorem:45}
Let \(V\) be a type. Let \(G\) and \(G^{'}\) be simple graphs on \(V\). Let \(M\) be a subgraph of \(G\). Let \(M^{'}\) be a subgraph of \(G^{'}\). Assume \(M\text{ is a perfect matching}\). Assume \(M^{'}\text{ is a perfect matching}\). Then
\begin{center}
\(\displaystyle \operatorname{spanning\,graph}(M) \mathbin{\triangle} \operatorname{spanning\,graph}(M^{'})\text{ alternates with respect to }\operatorname{spanning\,graph}(M).\)
\end{center}
\end{lemma}

\begin{proof}
Fix \(v\) in \(V\).
Fix \(w\) in \(V\).
Fix \(w^{'}\) in \(V\).
Assume \(w \ne w^{'}\).
Assume \(v \sim_{\operatorname{spanning\,graph}(M) \mathbin{\triangle} \operatorname{spanning\,graph}(M^{'})} w\).
Assume \(v \sim_{\operatorname{spanning\,graph}(M) \mathbin{\triangle} \operatorname{spanning\,graph}(M^{'})} w^{'}\).
Use the hypothesis \(M\text{ is a perfect matching}\).
Use the hypothesis \(M\text{ is a perfect matching}\).
Taking the second conjunct, we obtain \(v \in V(M)\).
Taking the first conjunct, we obtain \(\exists! w_{1}\, :\, V,\; v \sim_{M} w_{1}\).
Consider the following case.
\begin{proofstep}
Fix \(\mathit{v1}\) in \(V\).
Assume \(v \sim_{M} \mathit{v1} \land \left(\forall y\, :\, V,\; v \sim_{M} y \implies y = \mathit{v1}\right)\).
Use the hypothesis \(v \sim_{M} \mathit{v1} \land \left(\forall y\, :\, V,\; v \sim_{M} y \implies y = \mathit{v1}\right)\).
Consider the following case.
\begin{proofstep}
Assume \(v \sim_{M} \mathit{v1}\).
Assume \(\forall y\, :\, V,\; v \sim_{M} y \implies y = \mathit{v1}\).
Use the hypothesis \(M^{'}\text{ is a perfect matching}\).
Use the hypothesis \(M^{'}\text{ is a perfect matching}\).
Taking the second conjunct, we obtain \(v \in V(M^{'})\).
Taking the first conjunct, we obtain \(\exists! w_{1}\, :\, V,\; v \sim_{M^{'}} w_{1}\).
Consider the following case.
\begin{proofstep}
Fix \(\mathit{v2}\) in \(V\).
Assume \(v \sim_{M^{'}} \mathit{v2} \land \left(\forall y\, :\, V,\; v \sim_{M^{'}} y \implies y = \mathit{v2}\right)\).
Use the hypothesis \(v \sim_{M^{'}} \mathit{v2} \land \left(\forall y\, :\, V,\; v \sim_{M^{'}} y \implies y = \mathit{v2}\right)\).
Consider the following case.
\begin{proofstep}
Assume \(v \sim_{M^{'}} \mathit{v2}\).
Assume \(\forall y\, :\, V,\; v \sim_{M^{'}} y \implies y = \mathit{v2}\).
Use the hypothesis \(v \sim_{\operatorname{spanning\,graph}(M) \mathbin{\triangle} \operatorname{spanning\,graph}(M^{'})} w\), which simplifies to \(v \sim_{(\operatorname{spanning\,graph}(M) \setminus \operatorname{spanning\,graph}(M^{'})) \cup (\operatorname{spanning\,graph}(M^{'}) \setminus \operatorname{spanning\,graph}(M))} w\).
Rewriting with \paperlink{premise:53696d706c6547726170682e7375705f61646a}{Prerequisite~P\ref{premise:53696d706c6547726170682e7375705f61646a}} and \paperlink{premise:53696d706c6547726170682e73646966665f61646a}{Prerequisite~P\ref{premise:53696d706c6547726170682e73646966665f61646a}} in the preceding result gives \(v \sim_{M} w \land \neg (v \sim_{M^{'}} w) \lor v \sim_{M^{'}} w \land \neg (v \sim_{M} w)\).
Consider the following case.
\begin{proofstep}
Assume \(v \sim_{M} w \land \neg \left(v \sim_{M^{'}} w\right)\).
Use the hypothesis \(v \sim_{\operatorname{spanning\,graph}(M) \mathbin{\triangle} \operatorname{spanning\,graph}(M^{'})} w^{'}\), which simplifies to \(v \sim_{(\operatorname{spanning\,graph}(M) \setminus \operatorname{spanning\,graph}(M^{'})) \cup (\operatorname{spanning\,graph}(M^{'}) \setminus \operatorname{spanning\,graph}(M))} w^{'}\).
Rewriting with \paperlink{premise:53696d706c6547726170682e7375705f61646a}{Prerequisite~P\ref{premise:53696d706c6547726170682e7375705f61646a}} and \paperlink{premise:53696d706c6547726170682e73646966665f61646a}{Prerequisite~P\ref{premise:53696d706c6547726170682e73646966665f61646a}} in the preceding result gives \(v \sim_{M} w^{'} \land \neg (v \sim_{M^{'}} w^{'}) \lor v \sim_{M^{'}} w^{'} \land \neg (v \sim_{M} w^{'})\).
Consider the following case.
\begin{proofstep}
Assume \(v \sim_{M} w^{'} \land \neg \left(v \sim_{M^{'}} w^{'}\right)\).
Use the hypothesis \(v \sim_{M} w \land \neg \left(v \sim_{M^{'}} w\right)\).
Taking the first conjunct gives \(v \sim_{M} w\).
Use the hypothesis \(v \sim_{M} w^{'} \land \neg \left(v \sim_{M^{'}} w^{'}\right)\).
Taking the first conjunct gives \(v \sim_{M} w^{'}\).
Use the hypothesis \(w \ne w^{'}\).
By the preceding argument, we obtain \(\neg \left(w = w^{'}\right)\).
Rewriting with \(\forall y\, :\, V,\; v \sim_{M} y \implies y = \mathit{v1}\), \(v \sim_{M} w\), \(v \sim_{M} \mathit{v1}\), \(v \sim_{M} w^{'}\) and \paperlink{premise:6e6f745f747275655f65715f66616c7365}{Prerequisite~P\ref{premise:6e6f745f747275655f65715f66616c7365}} in the preceding result gives \(\mathrm{False}\).
The contradiction rules out this case. Hence \(v \sim_{M} w \iff \neg \left(v \sim_{M} w^{'}\right)\).
\end{proofstep}
Consider the following case.
\begin{proofstep}
Assume \(v \sim_{M^{'}} w^{'} \land \neg \left(v \sim_{M} w^{'}\right)\).
Use the hypothesis \(v \sim_{M} w \land \neg \left(v \sim_{M^{'}} w\right)\).
Taking the first conjunct gives \(v \sim_{M} w\).
Use the hypothesis \(v \sim_{M^{'}} w^{'} \land \neg \left(v \sim_{M} w^{'}\right)\).
Taking the first conjunct gives \(v \sim_{M^{'}} w^{'}\).
Use the hypothesis \(v \sim_{M} w \land \neg \left(v \sim_{M^{'}} w\right)\).
Consider the following case.
\begin{proofstep}
Assume \(v \sim_{M} w\).
Assume \(\neg \left(v \sim_{M^{'}} w\right)\).
Use the hypothesis \(v \sim_{M^{'}} w^{'} \land \neg \left(v \sim_{M} w^{'}\right)\).
Consider the following case.
\begin{proofstep}
Assume \(v \sim_{M^{'}} w^{'}\).
Assume \(\neg \left(v \sim_{M} w^{'}\right)\).
Rewriting with \paperlink{premise:65715f66616c7365}{Prerequisite~P\ref{premise:65715f66616c7365}}, \(\forall y\, :\, V,\; v \sim_{M^{'}} y \implies y = \mathit{v2}\), \paperlink{premise:6f665f65715f74727565}{Prerequisite~P\ref{premise:6f665f65715f74727565}}, \(v \sim_{M^{'}} w^{'}\), \paperlink{premise:65715f74727565}{Prerequisite~P\ref{premise:65715f74727565}}, \(v \sim_{M^{'}} \mathit{v2}\), \(\neg \left(v \sim_{M} w^{'}\right)\), substitution of equals and \paperlink{premise:6e6f745f66616c73655f65715f74727565}{Prerequisite~P\ref{premise:6e6f745f66616c73655f65715f74727565}} gives \(\neg \left(v \sim_{M} \mathit{v2}\right)\).
\end{proofstep}
Using the two conjuncts, we obtain \(\neg \left(v \sim_{M} \mathit{v2}\right)\).
\end{proofstep}
Using the two conjuncts, we obtain \(\neg \left(v \sim_{M} \mathit{v2}\right)\).
Rewriting with \(\forall y\, :\, V,\; v \sim_{M} y \implies y = \mathit{v1}\), \(v \sim_{M} w\), \(v \sim_{M} \mathit{v1}\), \(\forall y\, :\, V,\; v \sim_{M^{'}} y \implies y = \mathit{v2}\), \(v \sim_{M^{'}} w^{'}\), \(v \sim_{M^{'}} \mathit{v2}\) and \paperlink{premise:747275655f696666}{Prerequisite~P\ref{premise:747275655f696666}} in the preceding result gives \(v \sim_{M} w \iff \neg \left(v \sim_{M} w^{'}\right)\).
\end{proofstep}
The two alternatives exhaust the cases. Hence \(v \sim_{M} w \iff \neg \left(v \sim_{M} w^{'}\right)\).
\end{proofstep}
Consider the following case.
\begin{proofstep}
Assume \(v \sim_{M^{'}} w \land \neg \left(v \sim_{M} w\right)\).
Use the hypothesis \(v \sim_{\operatorname{spanning\,graph}(M) \mathbin{\triangle} \operatorname{spanning\,graph}(M^{'})} w^{'}\), which simplifies to \(v \sim_{(\operatorname{spanning\,graph}(M) \setminus \operatorname{spanning\,graph}(M^{'})) \cup (\operatorname{spanning\,graph}(M^{'}) \setminus \operatorname{spanning\,graph}(M))} w^{'}\).
Rewriting with \paperlink{premise:53696d706c6547726170682e7375705f61646a}{Prerequisite~P\ref{premise:53696d706c6547726170682e7375705f61646a}} and \paperlink{premise:53696d706c6547726170682e73646966665f61646a}{Prerequisite~P\ref{premise:53696d706c6547726170682e73646966665f61646a}} in the preceding result gives \(v \sim_{M} w^{'} \land \neg (v \sim_{M^{'}} w^{'}) \lor v \sim_{M^{'}} w^{'} \land \neg (v \sim_{M} w^{'})\).
Consider the following case.
\begin{proofstep}
Assume \(v \sim_{M} w^{'} \land \neg \left(v \sim_{M^{'}} w^{'}\right)\).
Use the hypothesis \(v \sim_{M^{'}} w \land \neg \left(v \sim_{M} w\right)\).
Taking the first conjunct gives \(v \sim_{M^{'}} w\).
Use the hypothesis \(v \sim_{M} w^{'} \land \neg \left(v \sim_{M^{'}} w^{'}\right)\).
Taking the first conjunct gives \(v \sim_{M} w^{'}\).
Use the hypothesis \(v \sim_{M^{'}} w \land \neg \left(v \sim_{M} w\right)\).
Consider the following case.
\begin{proofstep}
Assume \(v \sim_{M^{'}} w\).
Assume \(\neg \left(v \sim_{M} w\right)\).
Use the hypothesis \(v \sim_{M} w^{'} \land \neg \left(v \sim_{M^{'}} w^{'}\right)\).
Consider the following case.
\begin{proofstep}
Assume \(v \sim_{M} w^{'}\).
Assume \(\neg \left(v \sim_{M^{'}} w^{'}\right)\).
Rewriting with \paperlink{premise:65715f66616c7365}{Prerequisite~P\ref{premise:65715f66616c7365}}, \(\forall y\, :\, V,\; v \sim_{M^{'}} y \implies y = \mathit{v2}\), \paperlink{premise:6f665f65715f74727565}{Prerequisite~P\ref{premise:6f665f65715f74727565}}, \(v \sim_{M^{'}} w\), \paperlink{premise:65715f74727565}{Prerequisite~P\ref{premise:65715f74727565}}, \(v \sim_{M^{'}} \mathit{v2}\), \(\neg \left(v \sim_{M} w\right)\), substitution of equals and \paperlink{premise:6e6f745f66616c73655f65715f74727565}{Prerequisite~P\ref{premise:6e6f745f66616c73655f65715f74727565}} gives \(\neg \left(v \sim_{M} \mathit{v2}\right)\).
\end{proofstep}
Using the two conjuncts, we obtain \(\neg \left(v \sim_{M} \mathit{v2}\right)\).
\end{proofstep}
Using the two conjuncts, we obtain \(\neg \left(v \sim_{M} \mathit{v2}\right)\).
Rewriting with \(\forall y\, :\, V,\; v \sim_{M^{'}} y \implies y = \mathit{v2}\), \(v \sim_{M^{'}} w\), \(v \sim_{M^{'}} \mathit{v2}\), \(\forall y\, :\, V,\; v \sim_{M} y \implies y = \mathit{v1}\), \(v \sim_{M} w^{'}\), \(v \sim_{M} \mathit{v1}\), \paperlink{premise:6e6f745f747275655f65715f66616c7365}{Prerequisite~P\ref{premise:6e6f745f747275655f65715f66616c7365}} and \paperlink{premise:6966665f66616c7365}{Prerequisite~P\ref{premise:6966665f66616c7365}} in the preceding result gives \(v \sim_{M} w \iff \neg \left(v \sim_{M} w^{'}\right)\).
\end{proofstep}
Consider the following case.
\begin{proofstep}
Assume \(v \sim_{M^{'}} w^{'} \land \neg \left(v \sim_{M} w^{'}\right)\).
Use the hypothesis \(v \sim_{M^{'}} w \land \neg \left(v \sim_{M} w\right)\).
Taking the first conjunct gives \(v \sim_{M^{'}} w\).
Use the hypothesis \(v \sim_{M^{'}} w^{'} \land \neg \left(v \sim_{M} w^{'}\right)\).
Taking the first conjunct gives \(v \sim_{M^{'}} w^{'}\).
Use the hypothesis \(w \ne w^{'}\).
By the preceding argument, we obtain \(\neg \left(w = w^{'}\right)\).
Rewriting with \(\forall y\, :\, V,\; v \sim_{M^{'}} y \implies y = \mathit{v2}\), \(v \sim_{M^{'}} w\), \(v \sim_{M^{'}} \mathit{v2}\), \(v \sim_{M^{'}} w^{'}\) and \paperlink{premise:6e6f745f747275655f65715f66616c7365}{Prerequisite~P\ref{premise:6e6f745f747275655f65715f66616c7365}} in the preceding result gives \(\mathrm{False}\).
The contradiction rules out this case. Hence \(v \sim_{M} w \iff \neg \left(v \sim_{M} w^{'}\right)\).
\end{proofstep}
The two alternatives exhaust the cases. Hence \(v \sim_{M} w \iff \neg \left(v \sim_{M} w^{'}\right)\).
\end{proofstep}
The two alternatives exhaust the cases. Hence \(v \sim_{M} w \iff \neg \left(v \sim_{M} w^{'}\right)\).
By the preceding argument, we obtain \(v \sim_{\operatorname{spanning\,graph}(M)} w \iff \neg (v \sim_{\operatorname{spanning\,graph}(M)} w^{'})\).
\end{proofstep}
Using the two conjuncts, we obtain \(v \sim_{\operatorname{spanning\,graph}(M)} w \iff \neg (v \sim_{\operatorname{spanning\,graph}(M)} w^{'})\).
\end{proofstep}
Using the existential hypothesis, we obtain \(v \sim_{\operatorname{spanning\,graph}(M)} w \iff \neg (v \sim_{\operatorname{spanning\,graph}(M)} w^{'})\).
\end{proofstep}
Using the two conjuncts, we obtain \(v \sim_{\operatorname{spanning\,graph}(M)} w \iff \neg (v \sim_{\operatorname{spanning\,graph}(M)} w^{'})\).
\end{proofstep}
Using the existential hypothesis, we obtain \(v \sim_{\operatorname{spanning\,graph}(M)} w \iff \neg (v \sim_{\operatorname{spanning\,graph}(M)} w^{'})\).
\end{proof}

\begin{corollary}[{Is alternating symm diff right}]\label{theorem:46}
Let \(V\) be a type. Let \(G\) and \(G^{'}\) be simple graphs on \(V\). Let \(M\) be a subgraph of \(G\). Let \(M^{'}\) be a subgraph of \(G^{'}\). Assume \(M\text{ is a perfect matching}\). Assume \(M^{'}\text{ is a perfect matching}\). Then
\begin{center}
\(\displaystyle \operatorname{spanning\,graph}(M) \mathbin{\triangle} \operatorname{spanning\,graph}(M^{'})\text{ alternates with respect to }\operatorname{spanning\,graph}(M^{'}).\)
\end{center}
\end{corollary}

\begin{proof}
By \paperlink{premise:73796d6d446966665f636f6d6d}{Prerequisite~P\ref{premise:73796d6d446966665f636f6d6d}} with \(\alpha = \operatorname{Simple\,Graph}(V)\) and \(a = \operatorname{spanning\,graph}(M^{'})\) and \(b = \operatorname{spanning\,graph}(M)\), we obtain \(\operatorname{symmetric\,difference}_{\operatorname{to\,Max},\operatorname{to\,SDiff}}(\operatorname{spanning\,graph}(M^{'}),\allowbreak{} \operatorname{spanning\,graph}(M)) = \operatorname{symmetric\,difference}_{\operatorname{to\,Max},\operatorname{to\,SDiff}}(\operatorname{spanning\,graph}(M),\allowbreak{} \operatorname{spanning\,graph}(M^{'}))\).
Applying the same function to equal arguments, we obtain \((\operatorname{symmetric\,difference}_{\operatorname{to\,Max},\operatorname{to\,SDiff}}(\operatorname{spanning\,graph}(M^{'}),\allowbreak{} \operatorname{spanning\,graph}(M))\text{ alternates with respect to }\operatorname{spanning\,graph}(M^{'})) = (\operatorname{symmetric\,difference}_{\operatorname{to\,Max},\operatorname{to\,SDiff}}(\operatorname{spanning\,graph}(M),\allowbreak{} \operatorname{spanning\,graph}(M^{'}))\text{ alternates with respect to }\operatorname{spanning\,graph}(M^{'}))\).
By \paperlink{theorem:45}{Lemma~\ref{theorem:45}}, we obtain \(\operatorname{spanning\,graph}(M^{'}) \mathbin{\triangle} \operatorname{spanning\,graph}(M)\text{ alternates with respect to }\operatorname{spanning\,graph}(M^{'})\).
Substituting along the equality, we obtain \(\operatorname{spanning\,graph}(M) \mathbin{\triangle} \operatorname{spanning\,graph}(M^{'})\text{ alternates with respect to }\operatorname{spanning\,graph}(M^{'})\).
\end{proof}

\end{document}
